%% Generated by Sphinx.
\def\sphinxdocclass{report}
\documentclass[a4paper,10pt,turkish]{sphinxmanual}
\ifdefined\pdfpxdimen
   \let\sphinxpxdimen\pdfpxdimen\else\newdimen\sphinxpxdimen
\fi \sphinxpxdimen=.75bp\relax
\ifdefined\pdfimageresolution
    \pdfimageresolution= \numexpr \dimexpr1in\relax/\sphinxpxdimen\relax
\fi
\newdimen\sphinxremdimen\sphinxremdimen = 10pt
%% let collapsible pdf bookmarks panel have high depth per default
\PassOptionsToPackage{bookmarksdepth=5}{hyperref}

\PassOptionsToPackage{booktabs}{sphinx}
\PassOptionsToPackage{colorrows}{sphinx}

\PassOptionsToPackage{warn}{textcomp}
\usepackage[utf8]{inputenc}
\ifdefined\DeclareUnicodeCharacter
% support both utf8 and utf8x syntaxes
  \ifdefined\DeclareUnicodeCharacterAsOptional
    \def\sphinxDUC#1{\DeclareUnicodeCharacter{"#1}}
  \else
    \let\sphinxDUC\DeclareUnicodeCharacter
  \fi
  \sphinxDUC{00A0}{\nobreakspace}
  \sphinxDUC{2500}{\sphinxunichar{2500}}
  \sphinxDUC{2502}{\sphinxunichar{2502}}
  \sphinxDUC{2514}{\sphinxunichar{2514}}
  \sphinxDUC{251C}{\sphinxunichar{251C}}
  \sphinxDUC{2572}{\textbackslash}
\fi
\usepackage{cmap}
\usepackage[T1]{fontenc}
\usepackage{amsmath,amssymb,amstext}
\usepackage{babel}



\usepackage{tgtermes}
\usepackage{tgheros}
\renewcommand{\ttdefault}{txtt}



\usepackage[Sonny]{fncychap}
\ChNameVar{\Large\normalfont\sffamily}
\ChTitleVar{\Large\normalfont\sffamily}
\usepackage{sphinx}

\fvset{fontsize=auto}
\usepackage{geometry}


% Include hyperref last.
\usepackage{hyperref}
% Fix anchor placement for figures with captions.
\usepackage{hypcap}% it must be loaded after hyperref.
% Set up styles of URL: it should be placed after hyperref.
\urlstyle{same}

\addto\captionsturkish{\renewcommand{\contentsname}{İçindekiler}}

\usepackage{sphinxmessages}
\setcounter{tocdepth}{1}


        \usepackage[
            a4paper,
            top=1cm,
            bottom=1cm,
            left=1cm,
            right=1cm
        ]{geometry}
    

\title{LinuxKernelBook}
\date{25 May 2026}
\release{1.0.0}
\author{Kaan Aslan}
\newcommand{\sphinxlogo}{\vbox{}}
\renewcommand{\releasename}{Yayım}
\makeindex
\begin{document}

\ifdefined\shorthandoff
  \ifnum\catcode`\=\string=\active\shorthandoff{=}\fi
  \ifnum\catcode`\"=\active\shorthandoff{"}\fi
\fi

\pagestyle{empty}
\sphinxmaketitle
\pagestyle{plain}
\sphinxtableofcontents
\pagestyle{normal}
\phantomsection\label{\detokenize{index::doc}}


\sphinxAtStartPar
Bu kitap Kaan ASLAN tarafından \sphinxstyleemphasis{C ve Sistem Programcıları Derneğindeki} \sphinxstylestrong{Linux Kernel İşletim Sistemlerinin Tasarımı
ve Gerçekleştirilmesi} kursundaki kurs notları temel alınarak oluşturulmuştur.

\sphinxAtStartPar
Kitaptaki içerik belli bir kaynak referans alınarak oluşturulmamıştır ve içeriğin oluşturulmasında yapay zeka araçlarından
faydalanılmamıştır. Kitap ile benzer içeriklere sahip kaynakların listesi \sphinxstyleemphasis{Giriş} bölümünde listelenmiştir.


\bigskip\hrule\bigskip

\subsubsection*{Ön Gereksinimler}

\sphinxAtStartPar
Okuyucunun C programlama diline hâkim olduğu ve temel Linux komut satırını
kullanabildiği varsayılmaktadır. Çekirdeği derleyip denemek için sanallaştırma
ortamına kurulu bir Debian türevi Linux dağıtımı önerilmektedir.


\bigskip\hrule\bigskip


\sphinxstepscope


\chapter{Giriş}
\label{\detokenize{introduction:giris}}\label{\detokenize{introduction:id1}}\label{\detokenize{introduction::doc}}

\section{Gereksinimler}
\label{\detokenize{introduction:gereksinimler}}
\sphinxAtStartPar
Kursumuz için bir sanal makineye herhangi bir Linux dağıtımının kurulması gerekmektedir.
Dağıtım minimalist biçimde kurulabilir. Ancak Linux içerisinde C derleyicisinin (gcc ya da
clang) ve binary utility araçlarının bulunuyor olması gerekir. Kursumuzda Linux kaynak
kodları üzerinde değişiklikler yapıp çekirdeği yeniden derleyerek çeşitli denemeler
yapacağız. Bu nedenle kurduğunuz sistemin bir kopyasını da saklamanızı salık veriyoruz.

\sphinxAtStartPar
Kursumuzda Debian türevi bir Linux dağıtımının kurulmuş olduğu varsayılacaktır. Biz bazı
konularda açıklamalar yaparken Debian dağıtımının \sphinxstyleemphasis{apt} paket sistemini kullanacağız.


\section{Kaynak Listesi}
\label{\detokenize{introduction:kaynak-listesi}}
\sphinxAtStartPar
Linux çekirdeği konusunda yararlanabileceğiniz ek kaynakların listesini aşağıda veriyoruz. Bunların arasında en ayrıntılı
olanı \sphinxstyleemphasis{Understandibg the Linux Kernel} kitabıdır. Bu kitabın üç baskısı yapılmıştır. Birincisi çekirdeğin 2.2’li versiyonlarına,
ikincisi 2.4 versiyonlarına ve üçüncüsü de 2.6’lı versiyonlarına ilişkindir. Ancak 2.6’dan sonra çekirdekte çeşitli alt sistemlerde
değişiklikler yapıldığı için bu kitap Linux çekirdeğinin güncel versiyonlarını tam olarak yansıtmamaktadır.

\sphinxAtStartPar
Biz kursumuzda belli bir kitabın programını izlemeyeceğiz. Kurs konuları Kaan ASLAN’ın kendi deneyimlerinden ve bilgi birikiminden
damıtılarak oluşturulmuştur.


\subsection{Understanding the Linux Kernel}
\label{\detokenize{introduction:understanding-the-linux-kernel}}\begin{enumerate}
\sphinxsetlistlabels{\arabic}{enumi}{enumii}{}{.}%
\item {} 
\sphinxAtStartPar
{[}1{]} D. P. Bovet and M. Cesati, \sphinxstyleemphasis{Understanding the Linux Kernel}, 1st ed.
Sebastopol, CA, USA: O’Reilly Media, 2000.
ISBN: 978\sphinxhyphen{}0\sphinxhyphen{}596\sphinxhyphen{}00002\sphinxhyphen{}5.

\item {} 
\sphinxAtStartPar
{[}2{]} D. P. Bovet and M. Cesati, \sphinxstyleemphasis{Understanding the Linux Kernel}, 2nd ed.
Sebastopol, CA, USA: O’Reilly Media, 2002.
ISBN: 978\sphinxhyphen{}0\sphinxhyphen{}596\sphinxhyphen{}00213\sphinxhyphen{}5.

\item {} 
\sphinxAtStartPar
{[}3{]} D. P. Bovet and M. Cesati, \sphinxstyleemphasis{Understanding the Linux Kernel}, 3rd ed.
Sebastopol, CA, USA: O’Reilly Media, 2005.
ISBN: 978\sphinxhyphen{}0\sphinxhyphen{}596\sphinxhyphen{}00565\sphinxhyphen{}5.

\end{enumerate}

\begin{sphinxadmonition}{note}{Not:}
\sphinxAtStartPar
3. baskı Linux 2.6.11 çekirdeğini kapsar ve serinin son baskısıdır.
Linux 5.x / 6.x için kernel.org resmi dokümantasyonu ve LWN.net
makaleleri tamamlayıcı kaynak olarak kullanılmalıdır.
\end{sphinxadmonition}


\bigskip\hrule\bigskip



\subsection{Linux Device Drivers}
\label{\detokenize{introduction:linux-device-drivers}}\begin{enumerate}
\sphinxsetlistlabels{\arabic}{enumi}{enumii}{}{.}%
\item {} 
\sphinxAtStartPar
{[}4{]} A. Rubini and J. Corbet, \sphinxstyleemphasis{Linux Device Drivers}, 1st ed.
Sebastopol, CA, USA: O’Reilly Media, 1998.

\item {} 
\sphinxAtStartPar
{[}5{]} A. Rubini and J. Corbet, \sphinxstyleemphasis{Linux Device Drivers}, 2nd ed.
Sebastopol, CA, USA: O’Reilly Media, 2001.

\item {} 
\sphinxAtStartPar
{[}6{]} J. Corbet, A. Rubini, and G. Kroah\sphinxhyphen{}Hartman, \sphinxstyleemphasis{Linux Device Drivers},
3rd ed. Sebastopol, CA, USA: O’Reilly Media, 2005.
ISBN: 978\sphinxhyphen{}0\sphinxhyphen{}596\sphinxhyphen{}00590\sphinxhyphen{}7.
Çevrimiçi (ücretsiz): \sphinxurl{https://lwn.net/Kernel/LDD3/}

\end{enumerate}

\begin{sphinxadmonition}{note}{Not:}
\sphinxAtStartPar
3. baskı Linux 2.6 çekirdeğini hedefler ve O’Reilly tarafından
ücretsiz olarak çevrimiçi yayınlanmaktadır. 4. baskı henüz
yayınlanmamıştır.
\end{sphinxadmonition}


\bigskip\hrule\bigskip



\subsection{Kernel İnternals ve Genel Mimari}
\label{\detokenize{introduction:kernel-internals-ve-genel-mimari}}\begin{enumerate}
\sphinxsetlistlabels{\arabic}{enumi}{enumii}{}{.}%
\item {} 
\sphinxAtStartPar
{[}7{]} R. Love, \sphinxstyleemphasis{Linux Kernel Development}, 3rd ed.
Upper Saddle River, NJ, USA: Addison\sphinxhyphen{}Wesley, 2010.
ISBN: 978\sphinxhyphen{}0\sphinxhyphen{}672\sphinxhyphen{}32946\sphinxhyphen{}3.

\item {} 
\sphinxAtStartPar
{[}8{]} M. Gorman, \sphinxstyleemphasis{Understanding the Linux Virtual Memory Manager}.
Upper Saddle River, NJ, USA: Prentice Hall, 2004.
ISBN: 978\sphinxhyphen{}0\sphinxhyphen{}131\sphinxhyphen{}45348\sphinxhyphen{}5.
Çevrimiçi (ücretsiz): \sphinxurl{https://www.kernel.org/doc/gorman/}

\item {} 
\sphinxAtStartPar
{[}9{]} W. Mauerer, \sphinxstyleemphasis{Professional Linux Kernel Architecture}.
Indianapolis, IN, USA: Wrox Press, 2008.
ISBN: 978\sphinxhyphen{}0\sphinxhyphen{}470\sphinxhyphen{}34343\sphinxhyphen{}2.

\item {} \begin{description}
\sphinxlineitem{{[}10{]} G. Kroah\sphinxhyphen{}Hartman, \sphinxstyleemphasis{Linux Kernel in a Nutshell}.}
\sphinxAtStartPar
Sebastopol, CA, USA: O’Reilly Media, 2006.
Çevrimiçi (ücretsiz): \sphinxurl{http://www.kroah.com/lkn/}

\end{description}

\end{enumerate}


\bigskip\hrule\bigskip



\subsection{Linux Kernel Programming}
\label{\detokenize{introduction:linux-kernel-programming}}\begin{enumerate}
\sphinxsetlistlabels{\arabic}{enumi}{enumii}{}{.}%
\item {} \begin{description}
\sphinxlineitem{{[}11{]} K. N. Billimoria, \sphinxstyleemphasis{Linux Kernel Programming: A Comprehensive Guide to Kernel Internals, Writing Kernel Modules, and Kernel Synchronization}, 1st ed.}
\sphinxAtStartPar
Birmingham, UK: Packt Publishing, 2021.

\end{description}

\item {} \begin{description}
\sphinxlineitem{{[}12{]} K. N. Billimoria, \sphinxstyleemphasis{Linux Kernel Programming Part 2: Writing Character Device Drivers: Learn to Work with User\sphinxhyphen{}Kernel Interfaces, Handle Peripheral I/O and Hardware Interrupts}, 1st ed.}
\sphinxAtStartPar
Birmingham, UK: Packt Publishing, 2021.

\end{description}

\end{enumerate}


\bigskip\hrule\bigskip



\subsection{Linux Kernel Debugging}
\label{\detokenize{introduction:linux-kernel-debugging}}\begin{enumerate}
\sphinxsetlistlabels{\arabic}{enumi}{enumii}{}{.}%
\item {} \begin{description}
\sphinxlineitem{{[}13{]} K. N. Billimoria, \sphinxstyleemphasis{Linux Kernel Debugging: Leverage Proven Tools and Advanced Techniques to Effectively Debug Linux Kernels and Kernel Modules}, 1st ed.}
\sphinxAtStartPar
Birmingham, UK: Packt Publishing, 2023.

\end{description}

\end{enumerate}

\begin{sphinxadmonition}{note}{Not:}
\sphinxAtStartPar
Billimoria’nın üç kitabı ({[}11{]}, {[}12{]}, {[}13{]}) birbirini tamamlar
niteliktedir: kernel programlama \(\rightarrow\) karakter sürücüleri \(\rightarrow\) kernel hata
ayıklama şeklinde doğal bir okuma sırası oluşturur. {[}13{]} numaralı
\sphinxstyleemphasis{Linux Kernel Debugging} kitabı \sphinxcode{\sphinxupquote{ftrace}}, \sphinxcode{\sphinxupquote{kprobes}}, \sphinxcode{\sphinxupquote{KASAN}},
\sphinxcode{\sphinxupquote{lockdep}} ve \sphinxcode{\sphinxupquote{kdump}} araçlarını kaynak kod düzeyinde ele alır.
\end{sphinxadmonition}


\bigskip\hrule\bigskip



\subsection{Sistem Programlama ve Kullanıcı Alanı Kaynakları}
\label{\detokenize{introduction:sistem-programlama-ve-kullanici-alani-kaynaklari}}\begin{enumerate}
\sphinxsetlistlabels{\arabic}{enumi}{enumii}{}{.}%
\item {} \begin{description}
\sphinxlineitem{{[}14{]} M. Kerrisk, \sphinxstyleemphasis{The Linux Programming Interface}.}
\sphinxAtStartPar
San Francisco, CA, USA: No Starch Press, 2010.
ISBN: 978\sphinxhyphen{}1\sphinxhyphen{}593\sphinxhyphen{}27220\sphinxhyphen{}3.

\end{description}

\item {} \begin{description}
\sphinxlineitem{{[}15{]} W. R. Stevens and S. A. Rago, \sphinxstyleemphasis{Advanced Programming in the UNIX Environment}, 3rd ed.}
\sphinxAtStartPar
Upper Saddle River, NJ, USA: Addison\sphinxhyphen{}Wesley, 2013.
ISBN: 978\sphinxhyphen{}0\sphinxhyphen{}321\sphinxhyphen{}63773\sphinxhyphen{}4.

\end{description}

\item {} \begin{description}
\sphinxlineitem{{[}16{]} B. Ward, \sphinxstyleemphasis{How Linux Works: What Every Superuser Should Know},}
\sphinxAtStartPar
3rd ed.
San Francisco, CA, USA: No Starch Press, 2021.

\end{description}

\item {} \begin{description}
\sphinxlineitem{{[}17{]} J. Levine, \sphinxstyleemphasis{Linkers and Loaders}.}
\sphinxAtStartPar
San Francisco, CA, USA: Morgan Kaufmann, 2000.

\end{description}

\end{enumerate}


\bigskip\hrule\bigskip



\subsection{Çevrimiçi Kaynaklar}
\label{\detokenize{introduction:cevrimici-kaynaklar}}\begin{enumerate}
\sphinxsetlistlabels{\arabic}{enumi}{enumii}{}{.}%
\item {} \begin{description}
\sphinxlineitem{{[}18{]} The Linux Kernel Organization, “The Linux Kernel Documentation,”}
\sphinxAtStartPar
{[}Çevrimiçi{]}. Erişim: \sphinxurl{https://www.kernel.org/doc/html/latest/}
{[}Erişim tarihi: Mayıs 2026{]}.

\end{description}

\item {} \begin{description}
\sphinxlineitem{{[}19{]} M. Gorman, “Understanding the Linux Virtual Memory Manager,”}
\sphinxAtStartPar
{[}Çevrimiçi{]}. Erişim: \sphinxurl{https://www.kernel.org/doc/gorman/}
{[}Erişim tarihi: Mayıs 2026{]}.

\end{description}

\item {} \begin{description}
\sphinxlineitem{{[}20{]} LWN.net, “Kernel development news and articles,” {[}Çevrimiçi{]}.}
\sphinxAtStartPar
Erişim: \sphinxurl{https://lwn.net}
{[}Erişim tarihi: Mayıs 2026{]}.

\end{description}

\item {} \begin{description}
\sphinxlineitem{{[}21{]} LWN.net, “Kernel Index — Memory Management,” {[}Çevrimiçi{]}.}
\sphinxAtStartPar
Erişim: \sphinxurl{https://lwn.net/Kernel/Index/\#Memory\_management}
{[}Erişim tarihi: Mayıs 2026{]}.

\end{description}

\item {} \begin{description}
\sphinxlineitem{{[}22{]} Bootlin, “Elixir Cross Referencer — Linux Kernel Source,”}
\sphinxAtStartPar
{[}Çevrimiçi{]}. Erişim: \sphinxurl{https://elixir.bootlin.com/linux/latest/source}
{[}Erişim tarihi: Mayıs 2026{]}.

\end{description}

\item {} \begin{description}
\sphinxlineitem{{[}23{]} KernelNewbies, “Linux Kernel Newbies,” {[}Çevrimiçi{]}.}
\sphinxAtStartPar
Erişim: \sphinxurl{https://kernelnewbies.org}
{[}Erişim tarihi: Mayıs 2026{]}.

\end{description}

\end{enumerate}


\bigskip\hrule\bigskip



\section{Çekirdek Kodları Üzerinde Gezinme}
\label{\detokenize{introduction:cekirdek-kodlari-uzerinde-gezinme}}
\sphinxAtStartPar
Kursumuzda Linux işletim sisteminin kaynak kodlarında gezinmek için
\sphinxurl{https://elixir.bootlin.com/linux/} sitesini kullanacağız. Bu sitede Linux’un 0.01
versiyonundan günümüzdeki en son versiyonuna kadar bütün resmi versiyonlarının kaynak
kodları bulunmaktadır ve bu kaynak kodlar üzerinde gezintiler (navigation) yapılabilmektedir.


\section{İşletim Sistemlerine Giriş}
\label{\detokenize{introduction:isletim-sistemlerine-giris}}
\sphinxAtStartPar
Bu bölümde işletim sistemleri hakkında genel bilgiler vereceğiz.


\section{İşletim Sisteminin Tanımı ve Yapısı}
\label{\detokenize{introduction:isletim-sisteminin-tanimi-ve-yapisi}}
\sphinxAtStartPar
İşletim sistemleri bilgisayar donanımının kaynaklarını yöneten, bilgisayar donanımı ile
kullanıcı arasında arayüz oluşturan sistem programlarıdır. Bilgisayar bilimlerinin akademik
öncülerinin çoğu işletim sistemlerini bir kaynak yöneticisi (resource manager) olarak
tanımlamıştır.

\begin{sphinxVerbatim}[commandchars=\\\{\}]
    +\PYGZhy{}\PYGZhy{}\PYGZhy{}\PYGZhy{}\PYGZhy{}\PYGZhy{}\PYGZhy{}\PYGZhy{}\PYGZhy{}\PYGZhy{}\PYGZhy{}\PYGZhy{}\PYGZhy{}\PYGZhy{}\PYGZhy{}\PYGZhy{}\PYGZhy{}\PYGZhy{}\PYGZhy{}\PYGZhy{}\PYGZhy{}\PYGZhy{}+
    | Uygulama Programları |
    +\PYGZhy{}\PYGZhy{}\PYGZhy{}\PYGZhy{}\PYGZhy{}\PYGZhy{}\PYGZhy{}\PYGZhy{}\PYGZhy{}\PYGZhy{}\PYGZhy{}\PYGZhy{}\PYGZhy{}\PYGZhy{}\PYGZhy{}\PYGZhy{}\PYGZhy{}\PYGZhy{}\PYGZhy{}\PYGZhy{}\PYGZhy{}\PYGZhy{}+
+\PYGZhy{}\PYGZhy{}\PYGZhy{}\PYGZhy{}\PYGZhy{}\PYGZhy{}\PYGZhy{}\PYGZhy{}\PYGZhy{}\PYGZhy{}\PYGZhy{}\PYGZhy{}\PYGZhy{}\PYGZhy{}\PYGZhy{}\PYGZhy{}\PYGZhy{}\PYGZhy{}\PYGZhy{}\PYGZhy{}\PYGZhy{}\PYGZhy{}\PYGZhy{}\PYGZhy{}\PYGZhy{}\PYGZhy{}\PYGZhy{}\PYGZhy{}\PYGZhy{}+
|       İşletim Sistemi       |
+\PYGZhy{}\PYGZhy{}\PYGZhy{}\PYGZhy{}\PYGZhy{}\PYGZhy{}\PYGZhy{}\PYGZhy{}\PYGZhy{}\PYGZhy{}\PYGZhy{}\PYGZhy{}\PYGZhy{}\PYGZhy{}\PYGZhy{}\PYGZhy{}\PYGZhy{}\PYGZhy{}\PYGZhy{}\PYGZhy{}\PYGZhy{}\PYGZhy{}\PYGZhy{}\PYGZhy{}\PYGZhy{}\PYGZhy{}\PYGZhy{}\PYGZhy{}\PYGZhy{}+
+\PYGZhy{}\PYGZhy{}\PYGZhy{}\PYGZhy{}\PYGZhy{}\PYGZhy{}\PYGZhy{}\PYGZhy{}\PYGZhy{}\PYGZhy{}\PYGZhy{}\PYGZhy{}\PYGZhy{}\PYGZhy{}\PYGZhy{}\PYGZhy{}\PYGZhy{}\PYGZhy{}\PYGZhy{}\PYGZhy{}\PYGZhy{}\PYGZhy{}\PYGZhy{}\PYGZhy{}\PYGZhy{}\PYGZhy{}\PYGZhy{}\PYGZhy{}\PYGZhy{}\PYGZhy{}\PYGZhy{}\PYGZhy{}\PYGZhy{}\PYGZhy{}\PYGZhy{}+
|        Bilgisayar Donanımı        |
+\PYGZhy{}\PYGZhy{}\PYGZhy{}\PYGZhy{}\PYGZhy{}\PYGZhy{}\PYGZhy{}\PYGZhy{}\PYGZhy{}\PYGZhy{}\PYGZhy{}\PYGZhy{}\PYGZhy{}\PYGZhy{}\PYGZhy{}\PYGZhy{}\PYGZhy{}\PYGZhy{}\PYGZhy{}\PYGZhy{}\PYGZhy{}\PYGZhy{}\PYGZhy{}\PYGZhy{}\PYGZhy{}\PYGZhy{}\PYGZhy{}\PYGZhy{}\PYGZhy{}\PYGZhy{}\PYGZhy{}\PYGZhy{}\PYGZhy{}\PYGZhy{}\PYGZhy{}+
\end{sphinxVerbatim}

\sphinxAtStartPar
İşletim sistemlerinin yönettiği kaynakların en önemlileri şunlardır:
\begin{itemize}
\item {} 
\sphinxAtStartPar
\sphinxstylestrong{CPU:} İşletim sistemi hangi programın ne zaman, ne kadar süre için CPU’ya atanacağına
karar verip bu işlemleri gerçekleştirmektedir.

\item {} 
\sphinxAtStartPar
\sphinxstylestrong{Ana Bellek (Main Memory / RAM):} İşletim sistemi programların ana belleğin neresine
yükleneceğine karar verir ve ana bellek kullanımını düzenler.

\item {} 
\sphinxAtStartPar
\sphinxstylestrong{İkincil Bellekler:} İşletim sistemi bir dosya sistemi (file system) oluşturarak
dosyaların parçalarını ikincil belleklerde etkin bir biçimde tutar ve kullanıcılara bir
dosya kavramıyla sunar.

\item {} 
\sphinxAtStartPar
\sphinxstylestrong{Çevre Birimleri (klavye, fare, yazıcı vb.):} İşletim sistemi fare, klavye, yazıcı
gibi çevre birimlerini yöneterek onları kullanıma hazır hale getirir. Yardımcı
işlemcileri (denetleyicileri) programlayarak onların işlev görmesini sağlamaktadır.

\item {} 
\sphinxAtStartPar
\sphinxstylestrong{Ağ İşlemleri:} İşletim sistemi ağa ilişkin donanım birimlerini yöneterek dışarıdan
gelen bilgileri onları talep eden programlara iletir.

\end{itemize}

\sphinxAtStartPar
İşletim sistemleri kaynak yönetimine göre alt sistemlere ayrılarak da incelenebilmektedir.
Örneğin işletim sisteminin \sphinxstyleemphasis{çizelgeleyici (scheduler)} alt sistemi demekle CPU yönetimini
sağlayan alt sistemi kastedilmektedir. Ana bellek yönetimi (memory management) yine
soyutlanarak incelenen önemli alt sistemlerden biridir. İşletim sistemlerinin ikincil bellek
yönetimine \sphinxstyleemphasis{dosya sistemi (file system)} da denilmektedir. Tabii bütün bu sistemler
birbirinden kopuk olarak değil birbirleriyle ilişkili bir biçimde işlev görmektedir. Bu
durumu insanın \sphinxstyleemphasis{solunum sistemi}, \sphinxstyleemphasis{dolaşım sistemi}, \sphinxstyleemphasis{sinir sistemi}, \sphinxstyleemphasis{boşaltım sistemi}
gibi alt sistemlerine benzetebiliriz. Bu alt sistemlerin birinde bile çalışma bozukluğu
oluşsa insan yaşamını yitirebilmektedir.

\sphinxAtStartPar
İşletim sistemleri yapı olarak iki kısımdan oluşmaktadır: Çekirdek (kernel) ve kabuk
(shell). Çekirdek işletim sisteminin donanımı kontrol eden ve kaynakları yöneten motor
kısmıdır. Aslında işletim sistemi denildiğinde akla çekirdek gelmektedir. Kabuk ise
işletim sisteminin kullanıcı ile arayüz oluşturan önyüzüdür. Örneğin UNIX/Linux
sistemlerinde bash gibi komut satırı, GNOME, KDE gibi pencere yöneticileri, Windows’taki
masaüstü (Explorer), macOS’teki masaüstü (Aqua) bu işletim sistemlerinin kabuk kısımlarını
oluşturmaktadır.

\begin{sphinxVerbatim}[commandchars=\\\{\}]
+\PYGZhy{}\PYGZhy{}\PYGZhy{}\PYGZhy{}\PYGZhy{}\PYGZhy{}\PYGZhy{}\PYGZhy{}\PYGZhy{}\PYGZhy{}\PYGZhy{}\PYGZhy{}\PYGZhy{}\PYGZhy{}\PYGZhy{}\PYGZhy{}\PYGZhy{}\PYGZhy{}\PYGZhy{}\PYGZhy{}\PYGZhy{}\PYGZhy{}\PYGZhy{}\PYGZhy{}\PYGZhy{}\PYGZhy{}\PYGZhy{}\PYGZhy{}\PYGZhy{}\PYGZhy{}\PYGZhy{}\PYGZhy{}\PYGZhy{}\PYGZhy{}\PYGZhy{}+
|           Kabuk (Shell)           |
|   +\PYGZhy{}\PYGZhy{}\PYGZhy{}\PYGZhy{}\PYGZhy{}\PYGZhy{}\PYGZhy{}\PYGZhy{}\PYGZhy{}\PYGZhy{}\PYGZhy{}\PYGZhy{}\PYGZhy{}\PYGZhy{}\PYGZhy{}\PYGZhy{}\PYGZhy{}\PYGZhy{}\PYGZhy{}\PYGZhy{}\PYGZhy{}\PYGZhy{}\PYGZhy{}\PYGZhy{}\PYGZhy{}\PYGZhy{}\PYGZhy{}+   |
|   |                           |   |
|   |    Çekirdek (Kernel)      |   |
|   |                           |   |
|   +\PYGZhy{}\PYGZhy{}\PYGZhy{}\PYGZhy{}\PYGZhy{}\PYGZhy{}\PYGZhy{}\PYGZhy{}\PYGZhy{}\PYGZhy{}\PYGZhy{}\PYGZhy{}\PYGZhy{}\PYGZhy{}\PYGZhy{}\PYGZhy{}\PYGZhy{}\PYGZhy{}\PYGZhy{}\PYGZhy{}\PYGZhy{}\PYGZhy{}\PYGZhy{}\PYGZhy{}\PYGZhy{}\PYGZhy{}\PYGZhy{}+   |
+\PYGZhy{}\PYGZhy{}\PYGZhy{}\PYGZhy{}\PYGZhy{}\PYGZhy{}\PYGZhy{}\PYGZhy{}\PYGZhy{}\PYGZhy{}\PYGZhy{}\PYGZhy{}\PYGZhy{}\PYGZhy{}\PYGZhy{}\PYGZhy{}\PYGZhy{}\PYGZhy{}\PYGZhy{}\PYGZhy{}\PYGZhy{}\PYGZhy{}\PYGZhy{}\PYGZhy{}\PYGZhy{}\PYGZhy{}\PYGZhy{}\PYGZhy{}\PYGZhy{}\PYGZhy{}\PYGZhy{}\PYGZhy{}\PYGZhy{}\PYGZhy{}\PYGZhy{}+
\end{sphinxVerbatim}

\sphinxAtStartPar
Peki işletim sistemi bu kadar temel donanım yönetimini sağlıyorsa işletim sistemi olmadan
programlama yapılabilir mi? İşletim sistemi olmadan programlama faaliyetine halk arasında
\sphinxstyleemphasis{bare metal programlama} denilmektedir. Bare metal programlama tipik olarak gömülü
sistemlerde, mikrodenetleyicilerin kullanıldığı uygulamalarda kullanılmaktadır. Bare metal
programlama genellikle özel bir amaca hizmet edecek biçimde yapılmaktadır. Amaçlar
fazlalaştığı zaman ve sistem karmaşıklaştığı zaman artık işletim sistemlerine gereksinim
duyulmaktadır.

\sphinxAtStartPar
Bazı kontrol yazılımları işletim sistemlerinin bazı etkinliklerini de sağlamaktadır. Bir
kontrol yazılımının işletim sistemi olarak isimlendirilmesi için yukarıda açıkladığımız
kaynak yönetimlerinin önemli bir bölümünü sağlıyor olması gerekir. Bu kaynak yönetimlerinin
çoğunu sağlamayan kontrol yazılımlarına genel olarak \sphinxstyleemphasis{firmware} de denilmektedir.


\section{İşletim Sistemlerinin Sınıflandırılması}
\label{\detokenize{introduction:isletim-sistemlerinin-siniflandirilmasi}}
\sphinxAtStartPar
İşletim sistemleri çeşitli biçimlerde sınıflandırılabilmektedir.


\subsection{Proses Yönetimine Göre}
\label{\detokenize{introduction:proses-yonetimine-gore}}
\sphinxAtStartPar
Aynı anda tek bir programı çalıştıran işletim sistemlerine \sphinxstyleemphasis{tek prosesli (single
processing)}, aynı anda birden fazla programı çalıştırabilen işletim sistemlerine ise
\sphinxstyleemphasis{çok prosesli (multiprocessing) işletim sistemleri} denilmektedir. Örneğin DOS işletim
sistemi tek prosesli bir sistemdi. Biz bu işletim sisteminde bir programı çalıştırırdık
ancak çalıştırdığımız program sonlanınca başka bir programı çalıştırabilirdik. Halbuki
Windows, UNIX/Linux, macOS gibi işletim sistemleri çok prosesli işletim sistemleridir.


\subsection{Kullanıcı Sayısına Göre}
\label{\detokenize{introduction:kullanici-sayisina-gore}}
\sphinxAtStartPar
Birden fazla farklı kullanıcının çalışabildiği sistemlere \sphinxstyleemphasis{çok kullanıcılı (multiuser)},
tek bir kullanıcının çalışabildiği sistemlere \sphinxstyleemphasis{tek kullanıcılı (single user)} sistemler
denilmektedir. Genellikle çok prosesli işletim sistemleri aynı zamanda çok kullanıcılı
sistemlerdir. Birden fazla kullanıcının söz konusu olduğu sistemlerde kullanıcıların
yetkilerinin ayarlanması, kullanıcıların birbirlerinin alanlarına erişmesinin engellenmesi,
sistem kaynaklarını belli oranlarda bölüşmesi gerekebilmektedir. Örneğin DOS tek
kullanıcılı bir sistemdi. Halbuki Windows, UNIX/Linux ve macOS sistemleri çok kullanıcılı
sistemlerdir.


\subsection{Çekirdek Yapısına Göre}
\label{\detokenize{introduction:cekirdek-yapisina-gore}}
\sphinxAtStartPar
İşletim sistemleri çekirdek yapısına göre \sphinxstyleemphasis{tek parçalı çekirdekli (monolithic kernel)} ve
\sphinxstyleemphasis{mikro çekirdekli (microkernel)} olmak üzere ikiye ayrılmaktadır. Tek parçalı çekirdekli
işletim sisteminin büyük kısmı çekirdek modunda çalışır. Mikro çekirdekli sistemlerde ise
çekirdek modunda çalışan kısım minimize edilmeye çalışılmıştır. Aslında tek parçalı ve
mikro çekirdekli tasarımları bir spektrum olarak düşünebiliriz. (Örneğin bu spektrumda
bazı çekirdekler tek parçalı tarafa yakın, bazıları ise mikro tarafa yakın olabilmektedir.)


\subsection{Dışsal Olaylarla Yanıt Verebilme Özelliğine Göre}
\label{\detokenize{introduction:dissal-olaylarla-yanit-verebilme-ozelligine-gore}}
\sphinxAtStartPar
İşletim sistemleri dışsal olaylara yanıt verme bakımından gerçek zamanlı olan (real\sphinxhyphen{}time)
ve gerçek zamanlı olmayan (non\sphinxhyphen{}real\sphinxhyphen{}time) sistemler olmak üzere ikiye ayrılabilir. Dışsal
olaylara hızlı bir biçimde yanıt verebilecek çekirdek yapısına sahip olan işletim
sistemlerine \sphinxstyleemphasis{gerçek zamanlı (real\sphinxhyphen{}time) işletim sistemleri} denilmektedir. Gerçek zamanlı
işletim sistemleri de kendi aralarında \sphinxstyleemphasis{katı (hard real\sphinxhyphen{}time)} ve \sphinxstyleemphasis{gevşek (soft real\sphinxhyphen{}time)}
işletim sistemleri olmak üzere ikiye ayrılabilmektedir. Katı gerçek zamanlı sistemler
dışsal olaylara yanıt verme bakımından çok güvenilir olma iddiasındadır. Gevşek gerçek
zamanlı sistemler ise bu konuda daha toleranslıdır.


\subsection{Dağıtıklık Durumuna Göre}
\label{\detokenize{introduction:dagitiklik-durumuna-gore}}
\sphinxAtStartPar
İşletim sistemleri dağıtıklık durumuna göre \sphinxstyleemphasis{dağıtık olan (distributed)} ve \sphinxstyleemphasis{dağıtık
olmayan (non\sphinxhyphen{}distributed)} sistemler biçiminde ikiye ayrılabilmektedir. Dağıtık işletim
sistemlerinde sistem birden fazla bilgisayardan oluşan tek bir sistem gibi davranmaktadır.
Örneğin 10 tane makineyi tek bir sistem olarak düşünebilirsiniz. Bu durumda bu
bilgisayarların kaynakları (örneğin diskleri ve CPU’ları) bu 10 makine tarafından
paylaşılmaktadır. Windows, UNIX/Linux ve macOS dağıtık işletim sistemleri değildir. Ancak
bu sistemlerde dağıtık uygulamalar yapılabilmektedir.


\subsection{Donanım Özelliğine Göre}
\label{\detokenize{introduction:donanim-ozelligine-gore}}
\sphinxAtStartPar
Neredeyse her yaygın masaüstü işletim sisteminin bir mobil versiyonu da oluşturulmuştur.
iOS (iPhone Operating System) ve iPadOS Apple firmasının (yani macOS sistemlerinin) mobil
işletim sistemleridir. Android bir çeşit mobil Linux sistemi olarak değerlendirilebilir.
Android projesinde Linux çekirdeği alınmış, biraz özelleştirilmiş, bazı parçaları atılmış,
buna bir mobil arayüz giydirilmiş ve sistem akıllı telefonlara ve tabletlere uygun hale
getirilmiştir. Nokia eskiden Symbian sistemlerinde büyük bir pazar payına sahipti. Ancak
bu firma akıllı telefon geçişini iyi yönetemedi. MeeGo ve Maemo gibi işletim sistemlerini
denedi. Sonra ekonomik sıkıntılar sonucunca büyük ölçüde Microsoft tarafından satın alındı.
Windows’un mobil versiyonuna genel olarak Windows CE denilmektedir. Windows CE’nin akıllı
telefonlar ve tabletler için özelleştirilmiş biçimine ise Windows Mobile ve Windows Phone
denilmektedir. Ancak Microsoft 2010 yılında Windows Mobile işletim sistemini, 2017’de de
Windows Phone işletim sistemini sonlandırmıştır ve bu alandaki rekabetten tamamen
çekilmiştir. Windows CE ise Windows IoT Core ismiyle farklı bir tasarımla evrimleşerek
devam ettirilmektedir.


\subsection{Kaynak Kod Lisansına Göre}
\label{\detokenize{introduction:kaynak-kod-lisansina-gore}}
\sphinxAtStartPar
Kaynak kod lisansına göre işletim sistemlerini kabaca açık kaynak kodlu (open source) ve
mülkiyete bağlı (proprietary) olmak üzere ikiye ayırabiliriz. Açık kaynak kodlu işletim
sistemleri değişik açık kaynak kod lisanslarına sahip olabilmektedir. Bunların kaynak
kodları indirilip üzerinde değişiklikler yapılabilmektedir. Örneğin Windows işletim sistemi
mülkiyete sahiptir. Oysa Linux, BSD sistemleri, Solaris, Android gibi sistemler açık kaynak
kodludur. macOS sistemlerinin ise çekirdeği açık, diğer kısımları (örneğin kabuk kısmı ve
diğer katmanları) kapalıdır.


\subsection{Kaynak Kodun Özgünlüğüne Göre}
\label{\detokenize{introduction:kaynak-kodun-ozgunlugune-gore}}
\sphinxAtStartPar
Bazı işletim sistemleri bazı işletim sistemlerinin kodları alınıp değiştirilerek
oluşturulmuştur (örneğin Android ve macOS’ta olduğu gibi). Bazı işletim sistemlerinin
kodları ise sıfırdan yazılmıştır. Kodları sıfırdan yazılan, yani orijinal kod temeline
dayanan işletim sistemlerinden bazıları şunlardır:
\begin{itemize}
\item {} 
\sphinxAtStartPar
AT\&T UNIX

\item {} 
\sphinxAtStartPar
DOS

\item {} 
\sphinxAtStartPar
Windows

\item {} 
\sphinxAtStartPar
Linux

\item {} 
\sphinxAtStartPar
BSD’ler (belli bir yıldan sonra)

\item {} 
\sphinxAtStartPar
Solaris

\item {} 
\sphinxAtStartPar
XENIX

\item {} 
\sphinxAtStartPar
VMS

\end{itemize}

\sphinxAtStartPar
Burada orijinal mimari ile orijinal kod tabanını birbirine karıştırmamak gerekiyor. Linux
UNIX işletim sisteminin mimarisini temel almıştır. Ancak tüm kodları sıfırdan yazılmıştır.
Yani orijinal AT\&T UNIX sistemindeki kaynak kodların bir bölümü kopyalanarak
kullanılmamıştır.


\subsection{GUI Çalışma Desteğine Göre}
\label{\detokenize{introduction:gui-calisma-destegine-gore}}
\sphinxAtStartPar
Bazı işletim sistemleri GUI çalışma modelini doğrudan desteklerken bazıları
desteklememektedir. Örneğin Windows sistemleri çekirdekle entegre edilmiş bir GUI çalışma
modeli sunmaktadır. UNIX/Linux sistemleri de X Window (ya da X11) ve Wayland katmanlarıyla
benzer bir model sunmaktadır. Fakat örneğin DOS işletim sisteminin böyle bir doğal GUI
desteği yoktu.


\subsection{Ağ Üzerinde Hizmet Alıp Verme Rollerine Göre}
\label{\detokenize{introduction:ag-uzerinde-hizmet-alip-verme-rollerine-gore}}
\sphinxAtStartPar
İşletim sistemlerini ağ altında hizmet alıp verme rollerine göre \sphinxstyleemphasis{istemci (client) ve
sunucu (server)} biçiminde de iki gruba ayırabiliriz. Bazı işletim sistemlerinin istemci
versiyonları birbirlerinden ayrılmıştır. Bazılarında ise bu ayrım yapılmamıştır. Örneğin
Windows 7, 8, 10, 11 sistemleri bu bakımdan istemci (client) sistemleridir. Halbuki Windows
Server 2016, 2019 sunucu sistemleri olarak piyasaya sürülmüştür. Eskiden Mac OS X’in
istemci ve sunucu versiyonları farklıydı. Fakat Mac OS X 10.7 (Lion) ile birlikte istemci
ve sunucu versiyonları birleştirildi. Linux dağıtımlarının çoğu da hem istemci hem de
sunucu olarak kullanılabilmektedir. Ancak bazı dağıtımların istemci ve sunucu versiyonları
farklıdır.

\sphinxAtStartPar
Peki işletim sistemlerinin istemci ve sunucu versiyonları arasındaki farklılıklar nelerdir?
Kabaca iki tür farklılığın olduğunu söyleyebiliriz. Birincisi çekirdekle ilgili
farklılıklardır. Genellikle sunucu sistemlerinde çizelgeleyici alt sistemde istemci
sistemlerine göre farklılıklar bulunmaktadır. İkincisi ise barındırdıkları yardımcı
yazılımlardır. İşletim sistemlerinin sunucu versiyonları hazır bazı sunucu programlarını
da içermektedir.


\section{İşletim Sistemlerinin Tarihsel Gelişimi}
\label{\detokenize{introduction:isletim-sistemlerinin-tarihsel-gelisimi}}
\sphinxAtStartPar
1940’lı yıllarda ilk elektronik bilgisayarlar yapıldığında henüz bir işletim sistemi
kavramı yoktu. Bu bilgisayarlara program yazacak olanlar işletim sistemi faaliyetlerini de
kendileri yapmak zorunda kalıyordu. (Yani şimdi mikrodenetleyicilere bare metal kod
yazanlarda olduğu gibi.) Transistör bulunduktan sonra 1950’li yıllarda artık elektronik
bilgisayarlar yavaş yavaş transistörlerle yapılmaya başlandı. Transistörlerin ortaya
çıkması hem bilgisayarların kapasitelerini ve güvenilirliklerini artırmış, hem de güç
harcamalarını düşürmüştür.

\sphinxAtStartPar
1950’li yıllarda IBM gibi pek çok bilgisayar üreticisi firma yalnızca donanım satıyordu.
İşletim sistemi gibi programları yazmak kullanıcıların yapması gereken bir işti. Böylece
donanımı satın alan her kurum işletim sistemine benzeyen programları da kendisi yazıyordu.
Bu anlamda standart bir işletim sistemi yoktu. Bugünkü anlamda ilk işletim sisteminin
General Motors’un 1956 yılında IBM’in 701 sistemi için yazdığı NAA IO (North American
Aviation Input Output System) olduğu söylenebilir.

\sphinxAtStartPar
1960’lara gelindiğinde IBM, System/360 isminde yeni bir bilgisayar donanımı geliştirme
işine girişti ve artık donanımla işletim sistemini birlikte satma fikrini benimsedi. Bu
donanım 1964 yılında duyuruldu ve 1965 yılında gerçekleştirildi. İlk System/360 Model 30
bilgisayarları o zamanın \sphinxstyleemphasis{Solid Logic Technology (SLD)} teknolojisiyle üretilmişti. Hem
öncekilerden daha güçlüydü hem de daha az yer kaplıyordu. Saniyede 34500 işlem
yapabiliyordu ve 8K ile 64K ana belleğe sahipti. 1967 yılında System/360’ın Model 60’ı
piyasaya sürüldü. Bu model saniyede 16.6 milyon komut çalıştırabiliyordu ve ana belleği
tipik olarak 512K, 768K ve 1 MB idi. IBM Sistem 360 donanımları için 1964 yılında ilk kez
OS/360 işletim sistemini geliştirdi. IBM daha sonra 1967 yılında OS/360 Model 67 için
OS/360’ın TSS 360 isminde zaman paylaşımlı (time sharing system) bir versiyonunu daha
geliştirmiştir. IBM’in System/360 makineleri ve işletim sistemleri önemli ticari başarı
kazandı. System/360’ı System/370 izledi. System/360 ve System/370 için başka kurumlar da
işletim sistemleri geliştirmiştir. Michigan Terminal System (MTS) ve MUSIC/SP bunlar
arasında önemli olanlardandır.

\sphinxAtStartPar
1960’lı yıllarda başka firmalar da işletim sistemleri geliştirmiştir. Örneğin Control Data
Corporation firmasının SCOPE işletim sistemi batch işlemler yapabiliyordu. Aynı firma MACE
isminde bu işletim sisteminin zaman paylaşımlı bir versiyonunu da yazmıştır. Firma bu
çalışmalarını 1970’li yıllarda Kronos işletim sistemiyle devam ettirmiştir. Burroughs
firması 1961 yılında MCP işletim sistemi ile B5000 bilgisayarlarını, GE firması da 1962
yılında GECOS işletim sistemiyle GE\sphinxhyphen{}600 serisi bilgisayarlarını piyasaya sürdü. UNIVAC
dünyanın ilk ticari bilgisayarlarını üreten firmadır. Bu firma da 1962 yılında UNIVAC 1107
için EXEC I işletim sistemini yazdı. Bu işletim sistemini sırasıyla Exec 2 ve Exec 8 izledi.

\sphinxAtStartPar
DEC (Digital Equipment Corporation) eskilerin en önemli bilgisayar üretici firmalarından
biriydi. (DEC 1998 yılında Compaq firması tarafından, Compaq firması da 2002 yılında HP
firması tarafından satın alındı.) Firmanın en önemli ürünleri PDP (Programmed Data
Processor) isimli bilgisayarlarıdır. Firma PDP\sphinxhyphen{}1’den (1959) başlayarak PDP\sphinxhyphen{}16’ya
(1971\sphinxhyphen{}1972) kadar PDP makinelerinin 16 versiyonunu piyasaya sürmüştür. DEC’in PDP\sphinxhyphen{}8’inin
mini bilgisayar devrimini başlattığı söylenebilir. Bu model 50.000’in üzerinde satışa
ulaşmıştır. UNIX işletim sistemi 1969 yılında ilk kez DEC’in PDP\sphinxhyphen{}7 modeli üzerinde
yazılmıştır. 1965 yılında piyasaya sürülen DEC PDP\sphinxhyphen{}7 18 bitlik bir makineydi. Makine
DECsys denilen işletim sistemi benzeri bir yönetici programla beraber satılıyordu. DEC’in
1966 yılında çıkardığı PDP\sphinxhyphen{}10 26 bitlik bir makineydi. DEC bu modelle birlikte işletim
sistemi olarak TOPS\sphinxhyphen{}10 isimli bir sisteme geçti.

\sphinxAtStartPar
1960’lı yılların sonuna kadar işletim sistemleri ağırlıklı olarak sembolik makine diliyle
yazılıyordu. 1960’lı yılların sonlarında AT\&T Bell Lab. tarafından UNIX işletim sistemi
geliştirildiğinde önemli bir devrim yaşandı. UNIX işletim sistemi 1973 yılında C ile yeniden
yazılmıştır. Böylece artık işletim sistemlerinin yüksek seviyeli dillerle de yazılabildiği
görülmüştür. PDP\sphinxhyphen{}11’i 16 bitlik PDP\sphinxhyphen{}12 izledi. PDP\sphinxhyphen{}12 Intel’in x86 ve Motorola’nın 6800
işlemcileri için ilham kaynağı olmuştur.

\sphinxAtStartPar
1970’li yılların ikinci yarısında entegre devrelerin de geliştirilmesiyle \sphinxstyleemphasis{ev bilgisayarları
(home computer)} ortaya çıkmaya başladı. Bunlarda genellikle BASIC yorumlayıcıları ile iç
içe geçmiş CP/M ya da GEOS işletim sistemleri kullanılıyordu. 1970’li yıllarda pek çok
firma farklı ev bilgisayarları üretmiştir. BBC Micro, Commodore 64, Apple II, Atari,
Amstrad, ZX Spectrum dönemin en ünlü ev bilgisayarlarındandı. Bu makinelerde kullanılan
işlemciler Intel’in 8080’i, Zilog’un Z80’i, Motorola’nın 6800’ü gibi 8 bitlik
işlemcilerdi.

\sphinxAtStartPar
DEC firması 1977 yılında VAX isimli bilgisayarı ve 32 bitlik işlemci birimini piyasaya
sürdü. VAX ailesi makineler o yıllarda önemli bir ticari başarı kazanmıştır. DEC VAX
makineleri için VAX/VMS isimli bir işletim sistemi yazmıştı. DEC bu işletim sisteminin
ismini 1992 yılında OpenVMS olarak değiştirdi. DEC 1992 yılında 64 bitlik RISC tasarımı
olan Alpha işlemcilerini piyasaya sürdü ve OpenVMS Alpha işlemcilerine port edildi. OpenVMS
hala kullanılmaya devam etmektedir. Itanium ve X86\sphinxhyphen{}64 portları da vardır.

\sphinxAtStartPar
Apple firması 1976 yılında kuruldu. Apple’ın ilk bilgisayarı Apple I idi. Bunu 1977’de
Apple II, 1980’de de Apple III izledi. Bu ilk Apple bilgisayarlarında AppleDOS isimli
işletim sistemleri kullanılıyordu. Daha sonra Apple 1983’te Lisa modelini piyasaya sürdü.
1983’ün sonlarında da ilk Macintosh bilgisayarını çıkardı. Lisa ile birlikte Apple grafik
tabanlı işletim sistemlerine geçiş yaptı. Lisa ve sonraki Apple bilgisayarlarının hepsi
grafik bir arayüze sahiptir. Macintosh markası daha sonra Mac olarak telaffuz edilmeye
başlandı. Lisa bilgisayarlarında kullanılan işletim sistemi LisaOS ismindeydi. Apple daha
sonra Macintosh bilgisayarlarının değişik versiyonlarını piyasaya sürdü. Bunlardaki işletim
sistemini \sphinxstyleemphasis{System Software 1 (1984), System Software 2 (1985), System Software 3 (1986),
System Software 4 (1987), System Software 5 (1987), System Software 6 (1988) ve System
Software 7 (1991)} olarak isimlendirdi. Apple \sphinxstyleemphasis{System Software} 7.5’ten sonra işletim
sisteminin ismini \sphinxstyleemphasis{System Software} yerine Mac OS olarak değiştirdi ve System Software 7.6
versiyonu Mac OS 7.6 ismiyle çıktı. Daha sonra Apple 1997 yılında Mac OS 8’i, 1999 yılında
da Mac OS 9’u çıkarmıştır.

\sphinxAtStartPar
1980’li yıllarda Mac bilgisayarlarının fiyatı çok yüksekti ve satışları da iyi gitmiyordu.
Çünkü Steve Jobs bilgisayarların program yazmak için değil kullanmak için satın alınması
gerektiğini düşünüyordu. Nihayet Apple’daki çalkantılar sonucunda Steve Jobs 1985 yılında
Apple’dan ayrılmak zorunda kaldı (kovuldu da denebilir) ve NeXT firmasını kurdu. NeXT
firması NeXT isimli bilgisayarları geliştirdi. Bu bilgisayarlarda NeXTSTEP isimli işletim
sistemi kullanılıyordu. Daha sonra bu sistem açık hale getirildi ve OPENSTEP ismini aldı.
Dünyanın ilk Web tarayıcısı Tim Berners Lee tarafından CERN’de NeXT bilgisayarları üzerinde
gerçekleştirilmiştir.

\sphinxAtStartPar
Steve Jobs 1997 yılında Apple’a geri döndü. Apple da NeXT firmasını 200 milyon dolara satın
aldı. Sonra piyasaya iMac ve Power Mac serileri çıktı. Daha sonra Steve Jobs Mac’lerin
çekirdeklerini tamamen değiştirme kararı aldı. Mac’ler Mac OS’un 10 versiyonu ile birlikte
yeni bir çekirdeğe geçtiler. Mac OS işletim sistemlerinin 10’lu versiyonları Roma rakamıyla
Mac OS X biçiminde isimlendirilmiştir. Apple Mac OS X ismini 2012 yılında Mountain Lion
(10.8) sürümü ile OS X olarak, 2016 yılında da Sierra (10.12) sürümüyle birlikte de macOS
olarak değiştirmiştir.

\sphinxAtStartPar
DOS işletim sistemi text ekranda çalışıyordu. Microsoft da geleceğin grafik tabanlı işletim
sistemlerinde olduğunu gördü ve yavaş yavaş DOS’u bırakarak grafik tabanlı bir sisteme
geçmeyi planladı. Bunun için Windows isimli grafik arayüzün birinci versiyonunu 1985’te
çıkardı. Bunu 1987’de Windows 2, 1990’da Windows 3.0 ve 1992’de de Windows 3.1 izledi.
Bu 16 bit Windows sistemleri işletim sistemi değildi; DOS üzerinden çalıştırılan birer
grafik arayüz gibiydi. Microsoft daha sonra Windows’u Windows NT 3.1 ile bağımsız bir
işletim sistemi haline getirdi. Microsoft bundan sonra sırasıyla 1994 yılında Windows NT
3.5’i, 1995 yılında Windows NT 3.51’i ve Windows 95’i, 1998 yılında Windows 98’i, 2000
yılında Windows 2000 ve Windows ME’yi, 2001 yılında Windows XP’yi, 2006 yılında Windows
Vista’yı, 2012 yılında Windows 8’i, 2015 yılında Windows 10’u ve nihayet 2021 yılında da
Windows 11’i çıkarmıştır.

\sphinxAtStartPar
Linux işletim sistemi 1992 yılında bir dağıtım biçiminde piyasaya çıkmıştır. Linux işletim
sisteminin hikayesi daha geniş olarak izleyen derste ele alınmaktadır.


\section{UNIX Türevi Sistemlerin Tarihsel Gelişimi}
\label{\detokenize{introduction:unix-turevi-sistemlerin-tarihsel-gelisimi}}\label{\detokenize{introduction:ders03}}
\sphinxAtStartPar
UNIX işletim sistemi AT\&T Bell Laboratuvarlarında 1969\sphinxhyphen{}1971 yılları arasında geliştirildi.
Proje ekibinin lideri Ken Thompson’du. Çalışma ekibinde Dennis Ritchie, Brian Kernighan gibi
önemli isimler de vardı. Ekip daha önce General Electric’in GE\sphinxhyphen{}645 mainframe bilgisayarı
için Multics işletim sistemi üzerinde çalışıyordu. (Multics işletim sisteminin
geliştirilmesine 1964 yılında başlandı. Projede General Electric, MIT ve Bell Lab birlikte
çalışıyordu. Sonra proje Honeywell şirketi tarafından devralınmıştır.)

\sphinxAtStartPar
AT\&T 1969 yılında bu projeden çekilerek kendi işletim sistemini geliştirmek istemiştir.
Geliştirme çalışmasına DEC’in PDP\sphinxhyphen{}7 makinelerinde başlanmıştır. UNIX ismi 1970 yılında
Brian Kernighan tarafından Multics’ten kelime oyunu yapılarak uydurulmuştur. Proje ekibi
AT\&T’yi DEC PDP\sphinxhyphen{}11 almaya ikna etti ve böylece geliştirme çalışmaları PDP\sphinxhyphen{}11 ile devam
etti. UNIX’in resmi olarak ilk sürümü Ekim 1971’de, ikinci sürümü Aralık 1972’de, üçüncü
ve dördüncü sürümleri de 1973 yılında yayınlanmıştır.

\sphinxAtStartPar
UNIX işletim sistemi büyük ölçüde PDP’nin sembolik makine dili ve Ken Thompson’ın B isimli
programlama diliyle geliştirilmiştir. B programlama dili fonksiyonları alıp DEC’in makine
diline dönüştürüyordu. Bu bakımdan B bir yorumlayıcı değil derleyiciydi. İşte 1972 yılında
Dennis Ritchie, Ken Thompson’ın B programlama dilinden hareketle C Programlama Dilini
geliştirmiştir. UNIX işletim sisteminin dördüncü sürümü 1973 yılında yeniden C Programlama
Dili ile yazılmıştır.

\sphinxAtStartPar
1974 yılında UNIX’in beşinci sürümü oluşturuldu. Bu sürümlerin hepsi araştırma amaçlıydı
ve \sphinxstyleemphasis{educational license} ismiyle lisanslanmıştı. UNIX işletim sistemi bir araştırma projesi
olarak organize edilmişti. Bu nedenle AT\&T kaynak kodlarını araştırma kuruluşlarına ücretsiz
dağıtmıştır. 1975 yılında UNIX’in altıncı sürümü şirketlere yönelik hazırlandı. UNIX’in
altıncı versiyonunun kaynak kodları 20.000 dolara (bugünün yaklaşık 120.000 doları) şirketlere
sunuldu. 1977 yılında Bell Lab, UNIX’i Interdata 7/32 isimli 32 bit mimariye port etti.
Bunu 1978’de VAX portu izledi.

\sphinxAtStartPar
1974 yılında California Üniversitesi (Berkeley) işletim sisteminin kopyasını Bell Lab’tan
aldı. 1978 yılında \sphinxstyleemphasis{Berkeley Software Distribution (1BSD)} ismiyle AT\&T dışındaki ilk UNIX
dağıtımını gerçekleştirdi. Bu dağıtım hayatını hala FreeBSD, OpenBSD ve NetBSD olarak
devam ettirmektedir. 1979’da BSD’nin ikinci versiyonu (2BSD) ve 1979’un sonlarına doğru da
üçüncü versiyonu (3BSD) piyasaya sürüldü. Bunu 1980 yılında versiyon 4 (4BSD) izlemiştir.
1991 yılında BSD UNIX’ten AT\&T kodları tamamen arındırılmış ve kod bakımından özgün hale
getirilmiştir. BSD’nin son versiyonu 1995’te 4.4BSD Lite Release 2 olarak çıkmıştır.

\sphinxAtStartPar
1980’li yıllarda pek çok kurum ve ticari firma UNIX kodlarını lisans ücreti ödeyerek
AT\&T’den satın alıp kendilerine yönelik UNIX sistemleri oluşturmuştur. Bunların önemli
olanları şunlardır:
\begin{description}
\sphinxlineitem{AIX}
\sphinxAtStartPar
IBM tarafından geliştirilmiş olan UNIX türevi sistemlerdir. İlk kez 1986 yılında
piyasaya sürülmüştür. IBM AIX’i System/370, RS/6000 ve PS2 bilgisayarlarında
kullanıyordu. Bu sistemler AT\&T UNIX System 5 kodları temel alınarak geliştirilmiştir.
AIX hala kullanılmaktadır. Son sürümü 2021 yılında 7.3 olarak piyasaya sürülmüştür.
AIX PowerPC ve x86 işlemcileri için de port edilmiştir.

\sphinxlineitem{IRIX}
\sphinxAtStartPar
SGI firması tarafından AT\&T ve BSD kodları değiştirilerek 1988’de oluşturulmuştur.
2006’da bırakılmıştır.

\sphinxlineitem{HP\sphinxhyphen{}UX}
\sphinxAtStartPar
HP 9000 bilgisayarları için 1982’de oluşturulmuştur. Motorola 68000 ve Itanium
işlemcileri için yazılmıştır. Hala devam ettirilmektedir.

\sphinxlineitem{ULTRIX}
\sphinxAtStartPar
DEC firmasının PDP\sphinxhyphen{}7, PDP\sphinxhyphen{}11 ve VAX donanımları için geliştirdiği UNIX sistemiydi.
İlk versiyonu 1984 yılında çıktı. 1995 yılında piyasadan çekildi.

\sphinxlineitem{XENIX}
\sphinxAtStartPar
Microsoft tarafından 1980 yılında geliştirilmeye başlanmıştır. İlk versiyonu 1980’in
sonlarına doğru çıkmıştır. Daha sonra SCO firması Microsoft’la bu konuda iş birliği
yapmış, 1987 yılında da Microsoft sistemi tamamen SCO’ya devretmiştir.

\sphinxlineitem{SCO\sphinxhyphen{}UNIX}
\sphinxAtStartPar
SCO firması XENIX’i Microsoft’tan alınca bunu SCO\sphinxhyphen{}UNIX olarak devam ettirdi. SCO\sphinxhyphen{}UNIX’in
ilk versiyonu 1989 yılında çıktı. SCO sonra bunu OpenServer ismiyle devam ettirmiştir.

\sphinxlineitem{FreeBSD, NetBSD ve OpenBSD}
\sphinxAtStartPar
4.3BSD sistemleri temel alınarak geliştirilmiştir. FreeBSD ve NetBSD 1993 yılında,
OpenBSD ise 1996 yılında piyasaya çıkmıştır. Sürdürülmeye devam etmektedir. FreeBSD
genel amaçlı client ve server işletim sistemi olma niyetindedir. NetBSD daha taşınabilir
ve geniş bir port’a sahiptir; daha çok bilimsel çalışmalarda tercih edilmektedir.
OpenBSD güvenliğin önemli olduğu alanlarda tercih edilmektedir.

\sphinxlineitem{SunOS / Solaris}
\sphinxAtStartPar
Sun firmasının BSD kodlarıyla oluşturduğu UNIX türevi işletim sistemiydi. İlk versiyonu
1982 yılında çıktı. SunOS işletim sistemi 5.2 versiyonundan sonra (1992) Solaris ismiyle
pazarlanmaya başlamıştır. Solaris daha sonra OpenSolaris biçiminde açık kaynak kodlu
olarak bir süre varlığını devam ettirdi. Oracle firmasının Sun firmasını 2010’da satın
almasından sonra bu proje de durduruldu. Bu proje Illumos ismiyle başka bir ekip
tarafından devam ettirilmektedir.

\sphinxlineitem{Linux}
\sphinxAtStartPar
Linus Torvalds’ın öncülüğünde geliştirilmiş en yaygın UNIX türevi işletim sistemidir.
İlk versiyonu 1991 yılında çıkmıştır. Hala devam ettirilmektedir. Linux’un tarihsel
gelişimi aşağıdaki bölümde ayrıntılı bir biçimde açıklanmaktadır.

\sphinxlineitem{Mac OS X / OS X / macOS}
\sphinxAtStartPar
Carnegie Mellon üniversitesinin Mach isimli çekirdeği ile BSD UNIX sisteminin bir araya
getirilmesiyle oluşturulmuştur. İlk versiyonu 2001 yılında piyasaya sürülmüştür.

\end{description}


\section{Mac OS X Türevi Sistemler}
\label{\detokenize{introduction:mac-os-x-turevi-sistemler}}
\sphinxAtStartPar
İşletim sistemlerinin tarihsel gelişimini ele aldığımız önceki bölümde de belirttiğimiz gibi
Apple firmasının Mac bilgisayarları Mac OS’un 10 versiyonu ile birlikte yeni bir çekirdeğe
geçtiler. Mac OS işletim sistemlerinin 10’lu versiyonları Roma rakamıyla Mac OS X biçiminde
isimlendirildi. Apple Mac OS X ismini 2012 yılında Mountain Lion (10.8) sürümü ile OS X
olarak, 2016 yılında da Sierra (10.12) sürümüyle birlikte de macOS olarak değiştirdi. Biz
Mac OS X, OS X ve macOS sistemlerine bu bölümde \sphinxstyleemphasis{Mac OS X türevi sistemler} de diyeceğiz.

\sphinxAtStartPar
Mac OS X türevi işletim sistemleri UNIX türevi sistemlerdir. Bu işletim sistemlerinin
çekirdeğine Darwin denilmektedir. Darwin açık kaynak kodlu bir işletim sistemdir. Ancak
Mac OS X türevi sistemler tam anlamıyla açık sistemler değildir. Bu sistemlerin çekirdeği
açık olsa da geri kalan kısımları mülkiyete sahip (proprietary) biçimdedir.

\sphinxAtStartPar
Darwin’in hikayesi 1989 yılında NeXT’in NeXTSTEP işletim sistemiyle başladı. NeXTSTEP
daha sonra OpenStep ismiyle API düzeyinde standart hale getirildi. 1996’nın sonunda,
1997’nin başında Steve Jobs Apple’a dönerken Apple da NeXT firmasını satın aldı ve sonraki
işletim sistemini OpenStep üzerine kuracağını açıkladı. Bundan sonra Apple 1997’de OpenStep
üzerine kurulu olan Rapsody’yi çıkardı. 1998’de de yeni işletim sisteminin Mac OS X
olacağını açıkladı. Daha sonra 2000 yılında Apple Rapsody’den Darwin projesini türetti.
Darwin her ne kadar Mac sistemlerinin çekirdeği olarak tasarlanmışsa da ayrı bir işletim
sistemi olarak da yüklenebilmektedir. Ancak Darwin grafik arayüzü olmadığı için Mac
programlarını çalıştıramamaktadır. Daha sonra Darwin’i bağımsız bir işletim sistemi haline
getirmek amacıyla Darwin’den de çeşitli projeler türetilmiştir. Bunlardan biri Apple
tarafından 2002’de başlatılan OpenDarwin’dir. Bu proje 2006’da sonlandırılmıştır. 2007’de
PureDarwin projesi başlatılmıştır.

\sphinxAtStartPar
Darwin’in çekirdeği XNU üzerine oturtulmuştur. XNU, NeXT firması tarafından NEXTSTEP
işletim sisteminde kullanılmak üzere geliştirilmiş bir çekirdektir. XNU, Carnegie Mellon
(\sphinxstyleemphasis{Karnegi} diye okunuyor) üniversitesinin Mach 3 mikrokernel çekirdeği ile 4.3BSD karışımı
hibrit bir sistemdir.

\sphinxAtStartPar
Mac OS X türevi sistemlerin versiyonları şunlardır:
\begin{itemize}
\item {} 
\sphinxAtStartPar
Mac OS X 10.0 (Cheetah, 2001)

\item {} 
\sphinxAtStartPar
Mac OS X 10.1 (Puma, 2001)

\item {} 
\sphinxAtStartPar
Mac OS X 10.2 (Jaguar, 2002)

\item {} 
\sphinxAtStartPar
Mac OS X 10.3 (Panther, 2003)

\item {} 
\sphinxAtStartPar
Mac OS X 10.4 (Tiger, 2005)

\item {} 
\sphinxAtStartPar
Mac OS X 10.5 (Leopard, 2007)

\item {} 
\sphinxAtStartPar
Mac OS X 10.6 (Snow Leopard, 2009)

\item {} 
\sphinxAtStartPar
Mac OS X 10.7 (Lion, 2011)

\item {} 
\sphinxAtStartPar
OS X 10.8 (Mountain Lion, 2012)

\item {} 
\sphinxAtStartPar
OS X 10.9 (Mavericks, 2013)

\item {} 
\sphinxAtStartPar
OS X 10.10 (Yosemite, 2014)

\item {} 
\sphinxAtStartPar
OS X 10.11 (El Capitan, 2015)

\item {} 
\sphinxAtStartPar
macOS 10.12 (Sierra, 2017)

\item {} 
\sphinxAtStartPar
macOS 10.13 (High Sierra, 2017)

\item {} 
\sphinxAtStartPar
macOS 10.14 (Mojave, 2018)

\item {} 
\sphinxAtStartPar
macOS 10.15 (Catalina, 2019)

\item {} 
\sphinxAtStartPar
macOS 11 (Big Sur, 2020)

\item {} 
\sphinxAtStartPar
macOS 12 (Monterey, 2021)

\item {} 
\sphinxAtStartPar
macOS 13 (Ventura, 2022)

\item {} 
\sphinxAtStartPar
macOS 14 (Sonoma, 2023)

\item {} 
\sphinxAtStartPar
macOS 15 (Sequoia, 2024)

\end{itemize}

\sphinxAtStartPar
macOS büyük ölçüde POSIX uyumlu bir sistemdir.


\section{GNU Projesi ve Özgür Yazılım}
\label{\detokenize{introduction:gnu-projesi-ve-ozgur-yazilim}}
\sphinxAtStartPar
UNIX/Linux dünyasında önemli bir yeri olan GNU Projesi, özgür yazılım ve açık kaynak kod
akımları üzerinde durmak istiyoruz.

\sphinxAtStartPar
1970’lerdeki mikro bilgisayarlar devrimine kadar yazılımda bir telif anlayışı yoktu. Yani
yazılımın dağıtılması konusunda sözleşmeler ve hukuki yaptırımlara gerek duyulmamıştı.
Yazılım zaten donanımla birlikte satılıyordu ya da kuruma özel yapılıyordu. 1969 yılında
IBM yazılımı donanımla birlikte verdiği için rekabet kurallarına uymadığı gerekçesiyle
mahkemeye verilmiştir ve cezaya çarptırılmıştır. 1970’li yıllarda yazılım maliyetleri
artmış, yazılım sektörü genişlemiş ve lisanslama politikaları da uygulamaya sokulmuştur.
Pek çok yazılım bu yıllarda özel lisanslarla piyasaya sürülmeye başlanmıştır. 1980’li
yıllarda bu lisanslama faaliyetleri hız kazanmıştır.

\sphinxAtStartPar
1980’li yıllarda tüm UNIX türevi sistemler çeşitli biçimlerde sınırlandırıcı lisanslara
sahipti. Bu nedenle bedava ve sınırlamasız UNIX türevi bir işletim sistemine gereksinim
duyulmaya başlandı. İşte durumdan vazife çıkaran ünlü Emacs editörünün yazarı Richard
Stallman 1983 yılının sonlarına doğru GNU projesini başlattı ve özgür yazılım (free
software) fikrini ortaya attı. GNU projesinin amacı açık kaynak kodlu UNIX benzeri bir
işletim sistemini ve geliştirme araçlarını yazmaktı. Proje fiilen 1984 yılında başlamıştır.

\sphinxAtStartPar
Stallman 1985 yılında özgür yazılım kavramını yaygınlaştırmak amacıyla Free Software
Foundation (\sphinxurl{https://www.fsf.org}) isimli kurumu kurdu ve artık GNU projesi bu kurum
tarafından yürütülmeye başlandı. FSF özgür yazılım modeli için GPL (GNU Public License)
denilen bir lisans da geliştirdi. Özgür yazılım akımında oluşturulan bir yazılım
istenildiği gibi çalıştırılabilir, kopyalanabilir, incelenebilir, dağıtılabilir,
değiştirilebilir ve iyileştirilebilir. Özgür yazılım tipik olarak aşağıdaki dört özgürlükle
tanımlanmıştır:
\begin{description}
\sphinxlineitem{Özgürlük 0}
\sphinxAtStartPar
Programı her türlü amaç için çalıştırma özgürlüğü.

\sphinxlineitem{Özgürlük 1}
\sphinxAtStartPar
Programın kaynak kodunu inceleme ve değiştirebilme özgürlüğü.

\sphinxlineitem{Özgürlük 2}
\sphinxAtStartPar
Programın kopyalarını çıkartabilme ve yeniden dağıtabilme özgürlüğü.

\sphinxlineitem{Özgürlük 3}
\sphinxAtStartPar
Programı iyileştirebilme ve iyileştirilmiş programı yayınlama özgürlüğü.

\end{description}

\sphinxAtStartPar
GNU projesi bağlamında pek çok temel araç (gcc derleyici, ld bağlayıcı vb.) geliştirilmiştir.
Fakat hedeflenen çekirdek bir türlü oluşturulamamıştır.

\sphinxAtStartPar
Aslında özgür yazılım (free software) ile açık kaynak kod (open source) akımları arasında
bazı farklar olmakla birlikte her iki akımın da hedefleri benzerdir. Özgür yazılım bir
sosyal harekete benzetilirken açık kaynak kod akımı bir geliştirme metodolojisine
benzetilmektedir. Biz kursumuzda tüm bu akımları \sphinxstyleemphasis{açık kaynak kod (open source)} olarak
nitelendireceğiz. Özgür yazılım akımının temel lisansı GPL’dir (GNU Public Licence). Bunun
yumuşatılmış LGPL (Lesser GPL) biçiminde bir versiyonu da oluşturulmuştur. Ayrıca Apache,
MIT, BSD gibi açık kaynak kodlu başka lisanslar da vardır.


\section{Linux’un Tarihsel Gelişimi}
\label{\detokenize{introduction:linux-un-tarihsel-gelisimi}}
\sphinxAtStartPar
Linus Torvalds Helsinki Üniversitesinde öğrenciyken bir işletim sistemi yazmaya
niyetlenmiştir. O zamanlarda telif uygulanmayan UNIX türevi bir işletim sistemi kalmamıştı
ve GNU projesinin işletim sistemi de (GNU Hurd) bitirilememişti. MINIX isimli bir işletim
sistemi Andrew Tanenbaum tarafından yazılmıştı ancak bu sisteminin lisansı yalnızca akademik
kullanımlara izin verecek biçimde sınırlandırılmıştı. Linus Torvalds projesini USENET haber
gruplarında duyurdu ve zamanla kendisine gönüllü yardım edecek sistem programcıları buldu.
Yazılım dünyasında bu tür girişimlerle sık karşılaşıldığı halde başarı olasılığı nispeten
düşük olmaktadır. Linus Torvalds’ın bu girişimi başarıya ulaşmıştır.

\sphinxAtStartPar
1991 yılında Linux’un 0.01 sürümü oluşturuldu. 1994 yılında stabil bir biçimde Linux 1.0
versiyonu dağıtılmaya başlandı. Bunu 1996 yılında Linux 2.0, 1999 yılında 2.2, 2000 yılında
2.4 ve 2003 yılında 2.6 izledi. Daha sonra Linux versiyon numaralandırma sistemi
değiştirilmiştir. 2011 yılında 3.0, 2015 yılında 4.0, 2019 yılında 5.0, 2022 yılında da
6.0 versiyonu çıkmıştır.

\sphinxAtStartPar
Linux çekirdeklerinin versiyonları ve bu versiyonlarda eklenen önemli özellikler aşağıdaki
tabloda verilmiştir.


\begin{savenotes}\sphinxattablestart
\sphinxthistablewithglobalstyle
\centering
\sphinxcapstartof{table}
\sphinxthecaptionisattop
\sphinxcaption{Linux Çekirdeği Sürüm Geçmişi}\label{\detokenize{introduction:id2}}
\sphinxaftertopcaption
\begin{tabular}[t]{\X{10}{100}\X{18}{100}\X{72}{100}}
\sphinxtoprule
\begin{varwidth}[t]{\sphinxcolwidth{1}{3}}
\sphinxstyletheadfamily \sphinxAtStartPar
Sürüm
\sphinxbeforeendvarwidth
\end{varwidth}%
&\begin{varwidth}[t]{\sphinxcolwidth{1}{3}}
\sphinxstyletheadfamily \sphinxAtStartPar
Tarih
\sphinxbeforeendvarwidth
\end{varwidth}%
&\begin{varwidth}[t]{\sphinxcolwidth{1}{3}}
\sphinxstyletheadfamily \sphinxAtStartPar
Önemli Yenilikler
\sphinxbeforeendvarwidth
\end{varwidth}%
\\
\sphinxmidrule
\sphinxtableatstartofbodyhook\begin{varwidth}[t]{\sphinxcolwidth{1}{3}}
\sphinxAtStartPar
0.01
\sphinxbeforeendvarwidth
\end{varwidth}%
&\begin{varwidth}[t]{\sphinxcolwidth{1}{3}}
\sphinxAtStartPar
Ağustos 1991
\sphinxbeforeendvarwidth
\end{varwidth}%
&\begin{varwidth}[t]{\sphinxcolwidth{1}{3}}
\sphinxAtStartPar
Linus Torvalds tarafından duyurulan ilk sürüm; sadece temel fonksiyonlara sahipti.
\sphinxbeforeendvarwidth
\end{varwidth}%
\\
\sphinxhline\begin{varwidth}[t]{\sphinxcolwidth{1}{3}}
\sphinxAtStartPar
1.0
\sphinxbeforeendvarwidth
\end{varwidth}%
&\begin{varwidth}[t]{\sphinxcolwidth{1}{3}}
\sphinxAtStartPar
Mart 1994
\sphinxbeforeendvarwidth
\end{varwidth}%
&\begin{varwidth}[t]{\sphinxcolwidth{1}{3}}
\sphinxAtStartPar
İlk resmi sürüm; çoklu işlemci desteği yoktu. Ağ üzerinden TCP/IP desteği sağlandı.
\sphinxbeforeendvarwidth
\end{varwidth}%
\\
\sphinxhline\begin{varwidth}[t]{\sphinxcolwidth{1}{3}}
\sphinxAtStartPar
1.2
\sphinxbeforeendvarwidth
\end{varwidth}%
&\begin{varwidth}[t]{\sphinxcolwidth{1}{3}}
\sphinxAtStartPar
Mart 1995
\sphinxbeforeendvarwidth
\end{varwidth}%
&\begin{varwidth}[t]{\sphinxcolwidth{1}{3}}
\sphinxAtStartPar
x86 dışı mimarilere (Alpha, MIPS) ilk destekler geldi.
\sphinxbeforeendvarwidth
\end{varwidth}%
\\
\sphinxhline\begin{varwidth}[t]{\sphinxcolwidth{1}{3}}
\sphinxAtStartPar
2.0
\sphinxbeforeendvarwidth
\end{varwidth}%
&\begin{varwidth}[t]{\sphinxcolwidth{1}{3}}
\sphinxAtStartPar
Haziran 1996
\sphinxbeforeendvarwidth
\end{varwidth}%
&\begin{varwidth}[t]{\sphinxcolwidth{1}{3}}
\sphinxAtStartPar
SMP (Simetrik Çoklu İşlemci) desteği eklendi. Daha fazla mimari desteği sunuldu.
\sphinxbeforeendvarwidth
\end{varwidth}%
\\
\sphinxhline\begin{varwidth}[t]{\sphinxcolwidth{1}{3}}
\sphinxAtStartPar
2.2
\sphinxbeforeendvarwidth
\end{varwidth}%
&\begin{varwidth}[t]{\sphinxcolwidth{1}{3}}
\sphinxAtStartPar
Ocak 1999
\sphinxbeforeendvarwidth
\end{varwidth}%
&\begin{varwidth}[t]{\sphinxcolwidth{1}{3}}
\sphinxAtStartPar
Gelişmiş ağ yığını, IPv6 desteği, daha fazla SMP ölçeklenebilirliği.
\sphinxbeforeendvarwidth
\end{varwidth}%
\\
\sphinxhline\begin{varwidth}[t]{\sphinxcolwidth{1}{3}}
\sphinxAtStartPar
2.4
\sphinxbeforeendvarwidth
\end{varwidth}%
&\begin{varwidth}[t]{\sphinxcolwidth{1}{3}}
\sphinxAtStartPar
Ocak 2001
\sphinxbeforeendvarwidth
\end{varwidth}%
&\begin{varwidth}[t]{\sphinxcolwidth{1}{3}}
\sphinxAtStartPar
USB, PCMCIA ve Bluetooth desteği; 64 GB RAM’e kadar bellek desteği (PAE ile).
\sphinxbeforeendvarwidth
\end{varwidth}%
\\
\sphinxhline\begin{varwidth}[t]{\sphinxcolwidth{1}{3}}
\sphinxAtStartPar
2.6
\sphinxbeforeendvarwidth
\end{varwidth}%
&\begin{varwidth}[t]{\sphinxcolwidth{1}{3}}
\sphinxAtStartPar
Aralık 2003
\sphinxbeforeendvarwidth
\end{varwidth}%
&\begin{varwidth}[t]{\sphinxcolwidth{1}{3}}
\sphinxAtStartPar
Yeni scheduler (O(1)), udev ile dinamik aygıt yönetimi, sysfs, Native POSIX Thread
Library (NPTL), preemptive kernel.
\sphinxbeforeendvarwidth
\end{varwidth}%
\\
\sphinxhline\begin{varwidth}[t]{\sphinxcolwidth{1}{3}}
\sphinxAtStartPar
3.0
\sphinxbeforeendvarwidth
\end{varwidth}%
&\begin{varwidth}[t]{\sphinxcolwidth{1}{3}}
\sphinxAtStartPar
Temmuz 2011
\sphinxbeforeendvarwidth
\end{varwidth}%
&\begin{varwidth}[t]{\sphinxcolwidth{1}{3}}
\sphinxAtStartPar
Büyük bir teknik değişiklik yok; sadece sürüm numarası sadeleştirildi (2.6.x’lerin
devamı).
\sphinxbeforeendvarwidth
\end{varwidth}%
\\
\sphinxhline\begin{varwidth}[t]{\sphinxcolwidth{1}{3}}
\sphinxAtStartPar
4.0
\sphinxbeforeendvarwidth
\end{varwidth}%
&\begin{varwidth}[t]{\sphinxcolwidth{1}{3}}
\sphinxAtStartPar
Nisan 2015
\sphinxbeforeendvarwidth
\end{varwidth}%
&\begin{varwidth}[t]{\sphinxcolwidth{1}{3}}
\sphinxAtStartPar
Canlı kernel güncelleme (live patching) özelliği eklendi.
\sphinxbeforeendvarwidth
\end{varwidth}%
\\
\sphinxhline\begin{varwidth}[t]{\sphinxcolwidth{1}{3}}
\sphinxAtStartPar
5.0
\sphinxbeforeendvarwidth
\end{varwidth}%
&\begin{varwidth}[t]{\sphinxcolwidth{1}{3}}
\sphinxAtStartPar
Mart 2019
\sphinxbeforeendvarwidth
\end{varwidth}%
&\begin{varwidth}[t]{\sphinxcolwidth{1}{3}}
\sphinxAtStartPar
Yeni donanım desteği, enerji verimliliği iyileştirmeleri, Adiantum şifreleme
algoritması.
\sphinxbeforeendvarwidth
\end{varwidth}%
\\
\sphinxhline\begin{varwidth}[t]{\sphinxcolwidth{1}{3}}
\sphinxAtStartPar
5.4
\sphinxbeforeendvarwidth
\end{varwidth}%
&\begin{varwidth}[t]{\sphinxcolwidth{1}{3}}
\sphinxAtStartPar
Kasım 2019
\sphinxbeforeendvarwidth
\end{varwidth}%
&\begin{varwidth}[t]{\sphinxcolwidth{1}{3}}
\sphinxAtStartPar
Lockdown modu, fs\sphinxhyphen{}verity desteği.
\sphinxbeforeendvarwidth
\end{varwidth}%
\\
\sphinxhline\begin{varwidth}[t]{\sphinxcolwidth{1}{3}}
\sphinxAtStartPar
5.10
\sphinxbeforeendvarwidth
\end{varwidth}%
&\begin{varwidth}[t]{\sphinxcolwidth{1}{3}}
\sphinxAtStartPar
Aralık 2020
\sphinxbeforeendvarwidth
\end{varwidth}%
&\begin{varwidth}[t]{\sphinxcolwidth{1}{3}}
\sphinxAtStartPar
EXT4 ve Btrfs iyileştirmeleri, AMD GPU desteği.
\sphinxbeforeendvarwidth
\end{varwidth}%
\\
\sphinxhline\begin{varwidth}[t]{\sphinxcolwidth{1}{3}}
\sphinxAtStartPar
5.15
\sphinxbeforeendvarwidth
\end{varwidth}%
&\begin{varwidth}[t]{\sphinxcolwidth{1}{3}}
\sphinxAtStartPar
Kasım 2021
\sphinxbeforeendvarwidth
\end{varwidth}%
&\begin{varwidth}[t]{\sphinxcolwidth{1}{3}}
\sphinxAtStartPar
NTFS3 dosya sistemi, yeni I/O kontrolcüsü.
\sphinxbeforeendvarwidth
\end{varwidth}%
\\
\sphinxhline\begin{varwidth}[t]{\sphinxcolwidth{1}{3}}
\sphinxAtStartPar
6.0
\sphinxbeforeendvarwidth
\end{varwidth}%
&\begin{varwidth}[t]{\sphinxcolwidth{1}{3}}
\sphinxAtStartPar
Ekim 2022
\sphinxbeforeendvarwidth
\end{varwidth}%
&\begin{varwidth}[t]{\sphinxcolwidth{1}{3}}
\sphinxAtStartPar
Rust diline ilk çekirdek içi destek, Scheduler ve TCP performans iyileştirmeleri.
\sphinxbeforeendvarwidth
\end{varwidth}%
\\
\sphinxhline\begin{varwidth}[t]{\sphinxcolwidth{1}{3}}
\sphinxAtStartPar
6.1
\sphinxbeforeendvarwidth
\end{varwidth}%
&\begin{varwidth}[t]{\sphinxcolwidth{1}{3}}
\sphinxAtStartPar
Aralık 2022
\sphinxbeforeendvarwidth
\end{varwidth}%
&\begin{varwidth}[t]{\sphinxcolwidth{1}{3}}
\sphinxAtStartPar
Rust desteği genişletildi, Intel AMX desteği.
\sphinxbeforeendvarwidth
\end{varwidth}%
\\
\sphinxhline\begin{varwidth}[t]{\sphinxcolwidth{1}{3}}
\sphinxAtStartPar
6.6
\sphinxbeforeendvarwidth
\end{varwidth}%
&\begin{varwidth}[t]{\sphinxcolwidth{1}{3}}
\sphinxAtStartPar
Kasım 2023
\sphinxbeforeendvarwidth
\end{varwidth}%
&\begin{varwidth}[t]{\sphinxcolwidth{1}{3}}
\sphinxAtStartPar
Apple Silicon (M1/M2) desteği, yeni enerji yönetimi özellikleri.
\sphinxbeforeendvarwidth
\end{varwidth}%
\\
\sphinxbottomrule
\end{tabular}
\sphinxtableafterendhook\par
\sphinxattableend\end{savenotes}

\sphinxAtStartPar
Linux monolithic bir çekirdek yapısına sahiptir. Büyük ölçüde POSIX uyumu bulunmaktadır.


\section{Linux Dağıtımları}
\label{\detokenize{introduction:linux-dagitimlari}}
\sphinxAtStartPar
Açık kaynak kodlu yazılımlar bir araya getirilip paketlenerek istenildiği gibi
dağıtılabilmektedir. Dağıtım (distribution) bu anlamda kullanılan genel bir terimdir ve
her türlü açık kaynak kodlu yazılım için dağıtım söz konusu olabilir. Ancak biz burada
Linux dağıtımları üzerinde duracağız.

\sphinxAtStartPar
Linux temel olarak bir çekirdek geliştirme projesidir. Linux kaynak kodlarına baktığınızda
tüm kodların çekirdekle ilgili olduğunu görürsünüz. Çekirdeğin dışındaki tüm yazılımlar
(örneğin init prosesinden başlayarak, kabuk yazılımları, paket yöneticileri, pencere
yöneticileri vb.) hep başka proje grupları tarafından gerçekleştirilmiş açık kaynak kodlu
yazılımlardır. İşte tüm bu açık kaynak kodlu yazılımların Linux çekirdeği temelinde bir
araya getirilmesi ve doğrudan kullanıcının install edip çalıştırabileceği biçimde
paketlenmesine Linux dağıtımları denilmektedir. Linux dağıtımları pencere yöneticileri
(KDE, GNOME gibi), paket yöneticileri (APT, RPM, YUM, DPKG, PACMAN, ZYPPER gibi) ve
diğer yararlı uygulama programları bakımından farklılıklar gösterebilmektedir.

\sphinxAtStartPar
Toplamda iki yüzün üzerinde Linux dağıtımının olduğu söylenebilir. Ancak bunlar arasında
az sayıda dağıtım çok popüler olmuştur. Bazı dağıtımlar bazı dağıtımlardan fork edilerek
oluşturulmuştur. Aşağıda en çok kullanılan dağıtımlara ilişkin dağıtım ağacı verilmektedir:

\begin{sphinxVerbatim}[commandchars=\\\{\}]
Linux
├── Debian
│   ├── Ubuntu
│   │   ├── Linux Mint
│   │   ├── Pop!\PYGZus{}OS
│   │   ├── elementary OS
│   │   └── Zorin OS
│   ├── Devuan       \PYGZsh{} Systemd olmayan Debian
│   └── Kali Linux   \PYGZsh{} Güvenlik test amaçlı
├── Red Hat Linux (eski)
│   ├── Fedora       \PYGZsh{} Topluluk temelli, RHEL\PYGZsq{}in test yatağı
│   │   └── RHEL (Red Hat Enterprise Linux)
│   │       ├── CentOS (\(\rightarrow\) 2021 sonrası CentOS Stream)
│   │       ├── AlmaLinux
│   │       └── Rocky Linux
├── Slackware
│   └── Slax         \PYGZsh{} Hafif sürüm
├── Arch Linux
│   ├── Manjaro
│   └── EndeavourOS
├── Gentoo
│   └── Calculate Linux
├── SUSE Linux
│   ├── openSUSE Leap
│   └── openSUSE Tumbleweed
├── Android          \PYGZsh{} Mobil, Linux çekirdeğine dayalı
├── Alpine Linux     \PYGZsh{} Minimal, güvenli, konteyner dostu
└── Chrome OS
    └── Chromium OS  \PYGZsh{} Açık kaynak tabanı
\end{sphinxVerbatim}

\sphinxAtStartPar
En çok kullanılan Linux dağıtımları aşağıda özetlenmiştir.

\sphinxAtStartPar
\sphinxstylestrong{Debian:} En önemli ve en eski Linux dağıtımlarından biridir. Knoppix, Mint ve Ubuntu
dağıtımları Debian türevi dağıtımlardır.

\sphinxAtStartPar
\sphinxstylestrong{Fedora:} Red Hat firması tarafından çıkarılmış olan dağıtımdır. İlk kez 2003 yılında
oluşturulmuştur. RPM paket yöneticisini kullanır. CentOS ve Scientific Linux en önemli
Fedora türevi dağıtımlardır. 2000 yılında ilk sürümü yapılan Red Hat Enterprise Linux
(RHEL) en önemli Fedora türevidir. Ondan da CentOS, Scientific Linux gibi dağıtımlar
türetilmiştir. CentOS server makinelerde en yaygın kullanılan Linux versiyonudur.

\sphinxAtStartPar
\sphinxstylestrong{OpenSUSE:} Alman SUSE firmasının desteklediği dağıtımdır. SUSE Linux Enterprise isminde
ticari bir versiyonu da vardır. ZYpp, YaST ve RPM paket yöneticilerini kullanmaktadır.

\sphinxAtStartPar
\sphinxstylestrong{Slackware:} En eski Linux dağıtımıdır. 1993 yılında oluşturulmuştur. Sürdürümü yavaş
olmakla birlikte hala devam etmektedir.


\section{POSIX Standartları}
\label{\detokenize{introduction:posix-standartlari}}
\sphinxAtStartPar
1980’li yıllarda AT\&T ya da BSD kodlarından türetilmiş olan ve çoğunluğu şirketlere ait
olan pek çok UNIX türevi işletim sistemi oluşturuldu. Bu işletim sistemleri birbirlerine
çok benzemekle birlikte aralarında bazı farklılıklara da sahipti. İşte IEEE durumdan vazife
çıkartarak bu UNIX türevi sistemleri standardize etmek için kolları sıvadı ve bunun sonucu
olarak da POSIX standartlarını oluşturdu.

\sphinxAtStartPar
POSIX sözcüğü Richard Stallman tarafından önerilmiştir. \sphinxstyleemphasis{Portable Operating System Interface
for UNIX} sözcüklerinden kısaltılarak uydurulmuştur ve \sphinxstyleemphasis{poziks} biçiminde okunmaktadır.
POSIX standartları üzerinde çalışmalar 1985 yılında başlamıştır ve ilk standartlar 1988
yılında \sphinxstyleemphasis{IEEE Std 1003.1\sphinxhyphen{}1988} kod numarasıyla oluşturulmuştur. POSIX her ne kadar UNIX
türevi sistemler için düşünülmüşse de UNIX türevi mimariye sahip olmayan sistemler için de
kullanılabilecek bir standarttır. (Örneğin Windows sistemleri Interix denilen alt sistemle
POSIX uyumlu olarak da kullanılabilmekteydi. Interix alt sistemi daha sonra Windows 8 ile
birlikte Windows’tan kaldırılmıştır.)

\sphinxAtStartPar
POSIX standartları 4 bölümden oluşmaktadır:
\begin{enumerate}
\sphinxsetlistlabels{\arabic}{enumi}{enumii}{}{.}%
\item {} 
\sphinxAtStartPar
\sphinxstylestrong{Base Definitions:} Bu bölümde temel tanımlamalar bulunmaktadır.

\item {} 
\sphinxAtStartPar
\sphinxstylestrong{Shell \& Utilities:} Bu bölümde kabuk komutları ve standart utility programlar ele
alınmaktadır.

\item {} 
\sphinxAtStartPar
\sphinxstylestrong{System Interfaces:} Bu bölümde C programcıları için hazır bulunan POSIX fonksiyonları
açıklanmaktadır.

\item {} 
\sphinxAtStartPar
\sphinxstylestrong{Rationale:} Çeşitli kuralların ve özelliklerin gerekçeleri bu bölümde açıklanmaktadır.

\end{enumerate}

\sphinxAtStartPar
POSIX standartlarının zaman içerisinde çeşitli versiyonları çıkartılmıştır. Bu versiyonlarda
hem yeni POSIX fonksiyonları kütüphaneye eklenmiş hem de standartlardaki bazı bozukluklar
ve uyumsuzluklar düzeltilmiştir. Standardın önemli versiyonları şu senelerde yayınlanmıştır:
1992, 1993, 1995, 1997, 2001, 2004, 2008, 2017. Standartlardaki en önemli değişim 1993
yılında \sphinxstyleemphasis{POSIX 1.b} diye de isimlendirilen \sphinxstyleemphasis{Realtime\sphinxhyphen{}extensions} ile 1995 yılında \sphinxstyleemphasis{POSIX 1.c}
diye isimlendirilen \sphinxstyleemphasis{Thread\sphinxhyphen{}extensions} isimli eklemelerdir. Bu eklemelerle POSIX’e gerçek
zamanlı işlemler için çeşitli özelliklerle thread kullanımı eklenmiştir.

\sphinxAtStartPar
Single UNIX Specification, UNIX türevi sistemler için oluşturulmuş diğer önemli standarttır.
Bir sistemin UNIX olarak değerlendirilebilmesi için bu standartlara uygun olması
gerekmektedir. Standartlar Austin Group isimli toplulukla Open Group isimli dernek tarafından
geliştirilmiştir. Sürdürümü Open Group tarafından yapılmaktadır. Open Group hali hazırda
UNIX sistemlerinin isim haklarını elinde bulundurmaktadır.

\sphinxAtStartPar
POSIX standartları ile Single UNIX Specification standartları arasında eskiden daha fazla
farklılıklar vardı. Ancak bugün itibariyle bu iki standart birbirlerine yaklaştırılmış ve
son versiyonlarla aynı hale getirilmiştir. Single UNIX Specification dokümanlarına
İnternet’ten Open Group’un web sitesinden erişilebilir:

\sphinxAtStartPar
\sphinxurl{https://pubs.opengroup.org/onlinepubs/9799919799/}


\section{Linux Çekirdeğinin Kaynak Kod Yapısı}
\label{\detokenize{introduction:linux-cekirdeginin-kaynak-kod-yapisi}}
\sphinxAtStartPar
Linux çekirdeğinin kaynak kod ağacı üzerinde temel dizinler aynı kalmak üzere zaman
içerisinde çeşitli değişiklikler yapılmıştır. Biz burada çekirdeğin son versiyonlarını
dikkate alarak açıklamalar yapacağız. Linux kaynak kod ağacındaki temel dizinler şunlardır:

\begin{sphinxVerbatim}[commandchars=\\\{\}]
linux\PYGZhy{}6.x/
├── arch/          \(\rightarrow\) Mimarilere özel kodlar (x86, ARM, RISC\PYGZhy{}V, vb.)
├── block/         \(\rightarrow\) Blok aygıt altyapısı
├── certs/         \(\rightarrow\) Kernel modül imzalama için sertifikalar
├── crypto/        \(\rightarrow\) Kriptografik algoritmalar ve API\PYGZsq{}ler
├── Documentation/ \(\rightarrow\) Belgeler (API, özellikler, davranışlar)
├── drivers/       \(\rightarrow\) Aygıt sürücüleri (net, usb, gpu, vb.)
├── fs/            \(\rightarrow\) Dosya sistemi sürücüleri (ext4, btrfs, vb.)
├── include/       \(\rightarrow\) Global başlık dosyaları (headers)
├── io\PYGZus{}uring/      \(\rightarrow\) 5.1 çekirdeği ile eklenen asenkron io\PYGZus{}uring sistem fonksiyonları
├── init/          \(\rightarrow\) Kernelin ilk başlatma kodları
├── ipc/           \(\rightarrow\) IPC mekanizmaları (semaphore, message queue, vb.)
├── kernel/        \(\rightarrow\) Temel kernel işlevleri (zamanlayıcı, process, vb.)
├── lib/           \(\rightarrow\) Genel amaçlı yardımcı fonksiyonlar
├── mm/            \(\rightarrow\) Bellek yönetimi
├── net/           \(\rightarrow\) Ağ yığını (TCP/IP, socket, protokoller)
├── rust/          \(\rightarrow\) Rust dili ile yazılmış çekirdek kodları (6.x ile)
├── samples/       \(\rightarrow\) Örnek kodlar (eBPF, modül kodları, vb.)
├── scripts/       \(\rightarrow\) Derleme sürecine yardımcı betikler
├── security/      \(\rightarrow\) Güvenlik altyapısı (LSM, SELinux, AppArmor, vb.)
├── sound/         \(\rightarrow\) Ses sürücüleri (ALSA, vb.)
├── tools/         \(\rightarrow\) Kullanıcı uzayı araçları (perf, bpftool, vb.)
├── usr/           \(\rightarrow\) Yerleşik initramfs oluşturmak için
├── virt/          \(\rightarrow\) Sanallaştırma (KVM, Xen, vb.)
├── MAINTAINERS    \(\rightarrow\) Dosyaların kim tarafından korunduğu bilgisi
├── Makefile       \(\rightarrow\) Ana derleme talimatı
└── Kconfig        \(\rightarrow\) Kernel yapılandırma seçenekleri
\end{sphinxVerbatim}

\sphinxAtStartPar
Dizinlerin açıklamaları aşağıda verilmiştir.
\begin{description}
\sphinxlineitem{\sphinxcode{\sphinxupquote{arch}}}
\sphinxAtStartPar
Dizin adı \sphinxstyleemphasis{architecture} sözcüğünden kısaltılmıştır. Bu dizinde işlemci mimarisine göre
değişebilen çekirdek kodları, her bir işlemci ailesi ayrı bir dizinde olacak biçimde
bulundurulmaktadır.

\sphinxlineitem{\sphinxcode{\sphinxupquote{block}}}
\sphinxAtStartPar
Eskiden bu dizin \sphinxstyleemphasis{drivers} dizini içerisindeydi. Burada işletim sisteminin blok düzeyinde
işlemler yapan kodları bulundurulmuştur.

\sphinxlineitem{\sphinxcode{\sphinxupquote{certs}}}
\sphinxAtStartPar
Bu dizinde çekirdek modüllerine ilişkin sertifikasyonlarla ilgili kodlar bulunmaktadır.

\sphinxlineitem{\sphinxcode{\sphinxupquote{crypto}}}
\sphinxAtStartPar
Çekirdeğin içerisinde kullanılan şifreleme işlemlerine yönelik kaynak kodlar
bulunmaktadır.

\sphinxlineitem{\sphinxcode{\sphinxupquote{Documentation}}}
\sphinxAtStartPar
Burada çeşitli alt sistemlere ilişkin dokümanlar bulunmaktadır.

\sphinxlineitem{\sphinxcode{\sphinxupquote{drivers}}}
\sphinxAtStartPar
Burada aygıt sürücülere ilişkin çekirdek kodları bulunmaktadır. Tabii bu aygıt sürücüler
isteğe göre seçilerek yalnızca bir grubu derlenmektedir.

\sphinxlineitem{\sphinxcode{\sphinxupquote{fs}}}
\sphinxAtStartPar
Bu dizinde dosya sistemlerine yönelik çekirdek kodları bulundurulmaktadır.

\sphinxlineitem{\sphinxcode{\sphinxupquote{include}}}
\sphinxAtStartPar
Çekirdek kodlarında include edilen tüm başlık dosyaları bu dizinin içerisinde
bulundurulmaktadır. Tabii bu dizinde de alt dizinler vardır.

\sphinxlineitem{\sphinxcode{\sphinxupquote{init}}}
\sphinxAtStartPar
Çekirdeğin çalışabilir hale gelmesi için yapılan ilk işlemlere ilişkin kodlar burada
bulundurulmaktadır.

\sphinxlineitem{\sphinxcode{\sphinxupquote{io\_uring}}}
\sphinxAtStartPar
5.1 çekirdekleri ile eklenen yüksek performanslı asenkron IO işlemleri için kullanılan
sistem fonksiyonlarının gerçekleştirimine ilişkin kodlar bu dizinde bulunmaktadır.

\sphinxlineitem{\sphinxcode{\sphinxupquote{ipc}}}
\sphinxAtStartPar
Borular, mesaj kuyrukları gibi IPC nesnelerine ilişkin kaynak kodlar bu dizin
içerisindedir.

\sphinxlineitem{\sphinxcode{\sphinxupquote{kernel}}}
\sphinxAtStartPar
Bu dizin çekirdeğin en temel işlevlerine ilişkin kaynak kodları barındırmaktadır.
(Çekirdek kodlarının içerisinde ayrıca bir \sphinxcode{\sphinxupquote{kernel}} dizinin bulunması biraz tuhaf olsa
da bu aslında adeta \sphinxstyleemphasis{çekirdeğin çekirdeği} gibi bir anlama gelmektedir.)

\sphinxlineitem{\sphinxcode{\sphinxupquote{lib}}}
\sphinxAtStartPar
Bu dizinde çekirdeğin değişik yerlerinde kullanılan genel amaçlı utility fonksiyonların
kodları bulundurulmaktadır. Burada \sphinxstyleemphasis{lib} sözcüğünün statik ve dinamik kütüphane ile bir
ilgisi yoktur. Zaten Linux derlenirken bu biçimde kütüphane dosyaları
oluşturulmamaktadır.

\sphinxlineitem{\sphinxcode{\sphinxupquote{mm}}}
\sphinxAtStartPar
Çekirdeğin tüm ana bellek yönetim kodları bu dizinde bulunmaktadır.

\sphinxlineitem{\sphinxcode{\sphinxupquote{net}}}
\sphinxAtStartPar
Çekirdeğin tüm ağ (network) alt sistemine ilişkin kodları bu dizinde bulunmaktadır.

\sphinxlineitem{\sphinxcode{\sphinxupquote{rust}}}
\sphinxAtStartPar
Linux çekirdeğinde 6’lı versiyonlarla Rust Programlama Dili de kullanılmaya başlanmıştır.
Bu dizinde çekirdeğe ilişkin Rust kodları bulunmaktadır. Kursun yapıldığı tarihteki
Linux çekirdeklerinde Rust kodları çok az yer kaplamaktadır.

\sphinxlineitem{\sphinxcode{\sphinxupquote{samples}}}
\sphinxAtStartPar
Burada alt sistemlere ilişkin örnek kodlar bulunmaktadır.

\sphinxlineitem{\sphinxcode{\sphinxupquote{scripts}}}
\sphinxAtStartPar
Bu dizinde çekirdek derlemesi sırasında faydalanılan script dosyaları
bulundurulmuştur.

\sphinxlineitem{\sphinxcode{\sphinxupquote{security}}}
\sphinxAtStartPar
Çekirdeğin sistem güvenliği ile ilgili kodları bu dizinde bulunmaktadır.

\sphinxlineitem{\sphinxcode{\sphinxupquote{sound}}}
\sphinxAtStartPar
Ses işlemlerine yönelik çekirdek kodları bu dizinde bulunmaktadır.

\sphinxlineitem{\sphinxcode{\sphinxupquote{tools}}}
\sphinxAtStartPar
Bu dizinde çekirdek geliştirmesinde ve analizinde kullanılan yardımcı programlar
bulundurulmaktadır. Aslında bu dizindeki programların çekirdekle bir ilgisi yoktur.
Bunlar yardımcı araçlardır; çekirdek kodlarının içerisinde yer almazlar.

\sphinxlineitem{\sphinxcode{\sphinxupquote{usr}}}
\sphinxAtStartPar
Bu dizinde \sphinxstyleemphasis{geçici kök dosya sistemine (initial ramdisk)} ilişkin kodlar
bulunmaktadır. Bunların bazı sistemlerde konfigürasyon yoluyla çekirdek kodlarına
aktarılmaktadır.

\sphinxlineitem{\sphinxcode{\sphinxupquote{virt}}}
\sphinxAtStartPar
Burada çekirdeğin sanallaştırmaya hizmet eden kodları bulundurulmaktadır.

\end{description}


\section{Linux Sistemlerinin Boot Edilmesi}
\label{\detokenize{introduction:linux-sistemlerinin-boot-edilmesi}}
\sphinxAtStartPar
Kursumuzun bu bölümünde Linux sistemlerinin boot edilmesi süreci ele alınacaktır. İşletim sisteminin otomatik olarak
yüklenerek çalışır hale getirilmesi sürecine \sphinxstyleemphasis{boot} işlemi denilmektedir. (\sphinxstyleemphasis{Boot} terimi İngilizce \sphinxstyleemphasis{askeri bottan
(çizmeden)} gelmektedir.) Biz bilgisayar sistemine güç verdiğimizde bir süre sonra işletim sisteminin otomatik
yüklendiğini görmekteyiz. Aslında işletim sisteminin yüklenmesi biraz karmaşık bir süreçle gerçekleşmektedir.


\subsection{Reset Vektörü ve BIOS/Firmware}
\label{\detokenize{introduction:reset-vektoru-ve-bios-firmware}}
\sphinxAtStartPar
Mikroişlemciler güç uygulandığında (reset edildiğinde) belli bir adresten itibaren çalışacak biçimde tasarlanmıştır.
Buna işlemcilerin \sphinxstyleemphasis{reset vektörü} denilmektedir. İşlemcilerin reset vektörleri genellikle fiziksel belleğin başında
ya da sonunda bulunmaktadır. Örneğin Intel işlemcilerinde reset vektörü belleğin sonunda, ARM işlemcilerinde ise
belleğin başında bulunmaktadır. İşlemci reset edildiğinde belli bir adresten çalışmaya başladığına göre orada
çalışmaya hazır bir kodun bulunuyor olması gerekir. Tabii bu kod RAM (DRAM) bellekte bulunamaz. Çünkü sisteme güç
uygulandığında RAM belleğin de içeriği sıfırlanmaktadır, ayrıca bugünkü DRAM belleklerin kullanılabilmesi için de
onların donanımsal olarak programlanması gerekmektedir. İşte bilgisayar sistemlerinde işlemcinin reset vektöründe
tipik olarak ROM bellekler bulundurulur. Eskiden bu amaçla EPROM bellekler kullanılıyordu. Artık uzunca bir süredir
EEPROM (flash EPROM) bellekler kullanılmaktadır.

\sphinxAtStartPar
Bugün kullandığımız DRAM bellekler güç kaynağı verilir verilmez kullanıma hazır hale getirilememektedir. Onların da
kullanıma hazır hale getirilmesi (initialize edilmesi) gerekir. Benzer biçimde yine bazı donanım aygıtlarının
kullanılmadan önce kullanıma hazır hale getirilmesi de gerekmektedir. İşte reset vektörüne yerleştirilmiş olan
kodlar bilgisayar sistemini kullanıma hazır hale getirecek kodları barındırmaktadır. Tabii burada donanımdan
donanıma değişen bazı ayrıntılar da söz konusu olabilmektedir. Reset vektöründeki kodlar genellikle donanımı üreten
firmalar tarafından yazılmaktadır. Bunlar çok temel kodlar olduğu için bunlara \sphinxstyleemphasis{BIOS} ya da \sphinxstyleemphasis{firmware} gibi isimler
verilmiştir. Bu tür temel BIOS ya da firmware kodlarının donanımı kullanıma hazır hale getirmek için yaptığı
işlemlerden bazıları şunlardır:
\begin{itemize}
\item {} 
\sphinxAtStartPar
DRAM bellekleri (bildiğimiz RAM) kullanıma hazır getirilmesi

\item {} 
\sphinxAtStartPar
Çok çekirdekli sistemlerde çekirdeklerin kullanıma hazır hale getirilmesi

\item {} 
\sphinxAtStartPar
Donanım aygıtlarının ve çevre birimlerinin tespit edilmesi ve bunlar hakkındaki bilgilerin oluşturulması
(örneğin ACPI tablosunun oluşturulması)

\item {} 
\sphinxAtStartPar
Çevre birimlerinin (yani yardımcı işlemcilerin) kullanıma hazır hale getirilmesi

\item {} 
\sphinxAtStartPar
SSD ve HDD gibi depolama birimlerinin kullanıma hazır hale getirilmesi

\item {} 
\sphinxAtStartPar
Klavye, fare gibi giriş çıkış aygıtlarının kullanıma hazır hale getirilmesi

\item {} 
\sphinxAtStartPar
Ekran kartı gibi görüntü aygıtlarının kullanıma hazır hale getirilmesi

\end{itemize}

\sphinxAtStartPar
Reset vektöründeki kodlar çalıştırılınca artık donanım genel olarak çalışmaya hazır bir duruma getirilmiştir.
Bundan sonra akışın bir biçimde sistem programcısına devredilmesi gerekir. Akışın devredilmesi tipik olarak
ikincil bellekteki belli bir disk bloğunun RAM’e yüklenmesi ve oradaki kodun çalıştırılması yoluyla yapılmaktadır.
Tersten gidersek sistem programcısı programını ikincil belleğin (diskin) bir bloğuna yerleştirir ve bu süreçler
sonucunda bu program çalışır hale gelir. Peki işletim sistemi nasıl yüklenmektedir? ROM’daki BIOS ya da firmware
kodunun işletim sistemini yüklemesi genel olarak mümkün değildir. Bunun temel nedenleri şunlardır:
\begin{itemize}
\item {} 
\sphinxAtStartPar
İşletim sistemleri çok büyük olabilir. ROM’daki kod bunu yapabilecek yeterlilikte olmayabilir.

\item {} 
\sphinxAtStartPar
İşletim sistemlerinin yüklenmesi sistemden sisteme değişebilen ve nispeten karmaşık bir süreçtir. ROM’daki küçük
kodun bunu yapması genellikle mümkün olamamaktadır.

\item {} 
\sphinxAtStartPar
İkincil bellekte birden fazla işletim sistemi bulunuyor olabilir. Bu durumda bunların hangisinin nasıl
yükleneceğine karar verilmesi ve karar verilen işletim sisteminin yüklenmesi ROM’daki küçük program tarafından
genellikle yapılamamaktadır.

\end{itemize}

\sphinxAtStartPar
O halde tek çıkar yol ROM’daki programın diskteki bir önyükleyiciyi yüklemesi ve işletim sisteminin yüklenmesinin
bu program tarafından yapılmasıdır.


\subsection{Önyükleyici (Bootloader) Programlar}
\label{\detokenize{introduction:onyukleyici-bootloader-programlar}}
\sphinxAtStartPar
İşletim sistemlerini yükleyen bu tür programlara \sphinxstyleemphasis{önyükleyici (bootloader)} denilmektedir. (Biz \sphinxstyleemphasis{bootloader}
terimi yerine bunun Türkçesi olan \sphinxstyleemphasis{önyükleyici} terimini de kullanacağız.) Yani yukarıda da belirttiğimiz gibi
sistem programcısı diskteki önceden tespit edilmiş alana kendi kodunu yerleştirir. (Genellikle BIOS ya da firmware
kodları küçük bir disk bloğunu yükleyip çalıştırmaktadır.) İşte sistem programcısının özel disk bloğuna
yerleştirdiği bu program da önyükleyici (bootloader) programını yüklemektedir.

\sphinxAtStartPar
Bugün değişik platformlarda kullanılan değişik \sphinxstyleemphasis{önyükleyici} programlar bulunmaktadır. Örneğin Microsoft Windows
sistemleri için kendi \sphinxstyleemphasis{önyükleyici} programını kullanmaktadır. Buna \sphinxstyleemphasis{Windows Boot Manager (bootmgr)} denilmektedir.
UNIX/Linux sistemlerinde değişik proje grupları tarafından yazılmış olan değişik önyükleyici programlar
kullanılabilmektedir. Örneğin Linux sistemlerinde eskiden \sphinxstyleemphasis{LILO} isimli önyükleyici yoğun biçimde kullanılıyordu.
Daha sonra \sphinxstyleemphasis{GRUB} isimli önyükleyici yaygın biçimde kullanılmaya başlandı. Son yıllarda SYSLINUX isimli bootloader
paketi de belli bir ölçüde kullanım bulmuştur. SYSLINUX önyükleyicisi minimalist bir tasarıma sahiptir ve daha çok
küçük sistemlerde tercih edilmektedir. Gömülü Linux sistemlerinde ise bugün en çok tercih edilen \sphinxstyleemphasis{Das U\sphinxhyphen{}Boot} ya da
kısaca \sphinxstyleemphasis{U\sphinxhyphen{}Boot (Universal Bootloader)} denilen önyükleyici programdır.

\sphinxAtStartPar
Peki ROM’daki program önyükleyici (bootloader) programını ikincil bellekte nasıl arayıp bulmaktadır? Bunun için
birkaç teknik kullanılabilmektedir. Birincisi ROM’daki programın doğrudan ikincil belleğin önceden belirlenmiş
bloklarını yüklemesidir. Örneğin klasik PC sistemlerinde ROM’daki (BIOS’taki) program diskin 0’ıncı bloğunu
(yani ilk 512 byte’ı içeren bloğu) yüklemektedir. (Ancak daha sonra PC sistemlerinde \sphinxstyleemphasis{UEFI BIOS} ismi altında eski
klasik BIOS kodları oldukça geliştirilmiştir. Artık bu UEFI BIOS kodları bazı dosya sistemlerini de tanımaktadır.)
Örneğin Raspberry Pi’da ROM’daki kodlar FAT dosya sistemini tanıyabilmektedir. FAT dosya sistemi Microsoft’un DOS
sistemlerinde kullandığı karmaşık olmayan sade bir dosya sistemidir. İşte ROM’daki kodlar FAT gibi bir dosya
sistemini tanıyorsa FAT dosya sisteminin kök dizininde belli isimdeki dosyayı bulup RAM’e de yükleyebilmektedir.
Bazı sistemlerde (örneğin Beaglebone Black’lerde) ROM’daki kod önce küçük bir önyükleyiciyi yüklemekte, bu küçük
önyükleyici de asıl önyükleyiciyi yükleyebilmektedir.

\sphinxAtStartPar
Burada birkaç soru aklınıza gelebilir. Bunlardan biri ROM’daki reset vektöründe bulunan kodların oraya kimin
tarafından yerleştirilmiş olduğudur. ROM’daki reset vektöründe bulunan kodlar çok aşağı seviyeli kodlar olduğu
için bunlar genellikle donanımı tasarlayan kurumlar tarafından yazılıp oraya yerleştirilmektedir. Örneğin bugün
kullandığımız PC sistemlerinde ROM’daki bu kodlara BIOS (Basic Input Output System) denilmektedir. BIOS kodları PC
donanımını tasarlayan IBM tarafından yazılmıştı. Örneğin Beaglebone Black kartlarındaki ROM programı Texas
Instruments (TI) firması tarafından, Raspberry Pi kartlarındaki ROM programı ise Broadcom firması tarafından
yazılmıştır. Özetle ROM’da bulunan bu aşağı seviyeli kodlar ilgili donanımı tasarlayan kurumlardaki sistem
programcıları tarafından yazılmaktadır. Belli bir süredir bu tür BIOS alanlarında EEPROM teknolojisi kullanıldığı
için BIOS güncellemeleri de yapılabilmektedir.


\subsection{Genel Boot Akışı}
\label{\detokenize{introduction:genel-boot-akisi}}
\sphinxAtStartPar
İşletim sisteminin yüklenmesi pek çok donanım ve platformda aşağıdaki aşamalardan geçilerek yapılmaktadır:

\sphinxincludegraphics[]{graphviz-2eb08623d9f98bf468681086a9dc8433f58d030f.pdf}


\subsection{Boot Aşamaları Terminolojisi}
\label{\detokenize{introduction:boot-asamalari-terminolojisi}}
\sphinxAtStartPar
Sistem reset edildiğinde tüm boot işlemi bir bütünün parçaları gibi düşünülürse burada devreye giren programlara
birer aşama numarası da verilebilmektedir. Örneğin boot işleminin reset vektöründe bulunan kısmına \sphinxstyleemphasis{birinci aşama
önyükleyici (stage\sphinxhyphen{}1 / first stage bootloader)}, bu programın yüklediği programa \sphinxstyleemphasis{ikinci aşama önyükleyici
(stage\sphinxhyphen{}2 / second stage bootloader)} ve bunun da yüklediği programa (yani işletim sistemini yükleyen GRUB gibi
programlara da) \sphinxstyleemphasis{üçüncü aşama önyükleyici (stage\sphinxhyphen{}3 / third stage bootloader)} denilebilmektedir. Ancak bazı
kaynaklar bu ROM’daki kodu boot sürecinin bir parçası olarak ele almamaktadır. Dolayısıyla bu kaynaklar ROM
kodunun kendisine değil, onun yüklediği programa \sphinxstyleemphasis{birinci aşama önyükleyici (stage\sphinxhyphen{}1 / first stage bootloader)},
işletim sistemini yükleyen programa (yani GRUB gibi) ise \sphinxstyleemphasis{ikinci aşama önyükleyici (stage\sphinxhyphen{}2 / second stage
bootloader)} demektedir. Örneğin Wikipedia boot işleminde donanım bileşenlerini hazır hale getiren reset
vektöründeki programı \sphinxstyleemphasis{birinci aşama önyükleyici}, işletim sistemini yükleyen programı ise (GRUB gibi) \sphinxstyleemphasis{ikinci
aşama önyükleyici} olarak isimlendirmiştir. Biz kursumuzda izleyen paragraflarda bu terminolojiyi kullanacağız.

\sphinxAtStartPar
Bugün pek çok masaüstü Linux dağıtımında GRUB isimli önyükleyici program kullanılmaktadır. Bu önyükleyici program
Linux çekirdek imajını belleğe yükleyerek çalıştırmaktadır. Linux’un çekirdek parametreleri de bu önyükleyici
tarafından kullanıcıdan alınıp Linux’a verilmektedir. Biz kursumuzda çekirdek derlemesi yaparken ve sistemi
çekirdekle boot ederken önyükleyicinin GRUB olduğunu varsayacağız.


\subsection{PC Sistemlerinde Boot Süreci (Klasik BIOS)}
\label{\detokenize{introduction:pc-sistemlerinde-boot-sureci-klasik-bios}}
\sphinxAtStartPar
Kursumuzdaki örnekler x86 tabanlı masaüstü bilgisayar kullanılarak verilmektedir. (Burada \sphinxstyleemphasis{masaüstü (desktop)}
terimi yalnızca kasalı bilgisayarlar için değil notebook’lar için de kullanılan genel bir terimdir.) x86 tabanlı
masaüstü bilgisayarlara tarihsel bakımdan \sphinxstyleemphasis{PC (Personal Computer)} de denilmektedir. (2000’lerin başında Apple
Intel tabanlı PC mimarisine geçtiyse de kullandığı PC donanımında bazı farklılıklar da oluşturmuştur. Sonra
Apple’ın Intel mimarisini de bırakıp ARM mimarisine geçtiğini biliyorsunuz. Ancak PC denildiğinde Intel ve ARM
tabanlı macOS sistemleri kastedilmemektedir.)

\sphinxAtStartPar
Bugün kullandığımız PC sistemlerinin temel donanımları 70’li yılların sonlarında ve 80’li yılların başlarında IBM
tarafından tasarlanmıştır. Bu bilgisayarlar 1980’in Aralık ayında IBM’in tasarladığı donanım ve Microsoft’un
tasarladığı DOS işletim sistemiyle piyasaya sürüldü. Zaman ilerledikçe bu PC donanımlarında bazı iyileştirmeler
yapıldıysa da temel mimari büyük ölçüde aynı kalmıştır.

\sphinxAtStartPar
Eskiden PC’lerde klasik BIOS (\sphinxstyleemphasis{legacy BIOS}) kullanılıyordu. Ancak 2010’lardan sonra modern UEFI BIOS’lar
yaygınlaştı. Artık bugün ağırlıklı olarak UEFI BIOS’lar kullanılıyor. (Ancak bu modern UEFI BIOS’lar eski klasik
BIOS (\sphinxstyleemphasis{legacy BIOS}) gibi de çalışabilecek özelliklere sahip olabilmektedir.) Biz burada klasik eski BIOS’a sahip
PC’lerin boot sürecinden bahsedeceğiz.

\sphinxAtStartPar
Eski klasik BIOS’lara (\sphinxstyleemphasis{legacy BIOS}) sahip PC’ler reset edildiğinde BIOS kodları DRAM belleği ve pek çok donanım
birimini programladıktan sonra CMOS setup’ta belirtilen \sphinxstyleemphasis{boot sırasına (boot sequence)} göre ilgili medyanın
0’ıncı sektörünü belleğe yükleyip akışı oraya devretmektedir. (PC mimarisinde diskten okunabilecek ya da diske
yazılabilecek en küçük birime \sphinxstyleemphasis{sektör (sector)} denilmektedir.) İkincil belleklerdeki ilk sektöre \sphinxstyleemphasis{MBR (Master
Boot Record)} denilmektedir. MBR’de toplam 512 byte’lık bir program bulunur. MBR’nin sonunda 64 byte’lık \sphinxstyleemphasis{Disk
Bölümleme Tablosu (Disk Partition Table)} vardır. MBR’deki program duruma göre ya bir önyükleyici programını
yüklemekte ya da default durumda aktif disk bölümünün 0’ıncı sektörünü RAM’e yükleyip akışı oraya devretmektedir.
(PC sistemlerinde her disk bölümünün ilk sektörüne (0’ıncı sektörüne) \sphinxstyleemphasis{boot sektör} denilmektedir. Boot sektör
ilgili işletim sisteminin yüklenmesinden sorumludur.) Tabii MBR’deki program daha büyük bir önyükleyici programını
da yükleyebilmektedir. Bu durumda hangi işletim sisteminin yükleneceği önyükleyici program tarafından bir menü
yoluyla kullanıcıya da sorulabilmektedir.

\sphinxAtStartPar
UEFI BIOS öncesindeki eski BIOS sistemini kullanan (\sphinxstyleemphasis{legacy BIOS}) PC sistemlerindeki boot sürecini aşağıdaki
gibi özetleyebiliriz:

\sphinxincludegraphics[]{graphviz-5dfbb4b6a4f89c35bcc81e4bf8240691fc0b7529.pdf}

\sphinxAtStartPar
Kursumuzda UEFI BIOS sistemlerinin işlevi ve temel çalışma mekanizması başka bir bölümde ele alınacaktır.

\sphinxAtStartPar
Yukarıda da belirttiğimiz gibi bugünkü Linux yüklü olan Intel x86 tabanlı PC sistemlerinde genellikle önyükleyici
olarak GRUB tercih edilmektedir.

\sphinxstepscope


\chapter{Çekirdeğin Derlenmesi}
\label{\detokenize{kernelcompilation:cekirdegin-derlenmesi}}\label{\detokenize{kernelcompilation::doc}}
\sphinxAtStartPar
Bu bölümde ayrıntılı bir biçimde Linux çekireğinin derlenmesi sürecini ele alacapız.


\section{Çekirdek Davranışını Değiştirme Yöntemleri}
\label{\detokenize{kernelcompilation:cekirdek-davranisini-degistirme-yontemleri}}
\sphinxAtStartPar
Linux çekirdeğinin kodlarına dokunmadan onunla ilgili bazı davranış değişikliklerinin yapılması
temelde beş biçimde sağlanabilmektedir:

\sphinxAtStartPar
\sphinxstylestrong{1. Çekirdek modülleri ve aygıt sürücüler yoluyla}

\sphinxAtStartPar
Çekirdek modunda çalışan \sphinxstyleemphasis{çekirdek modülleri (kernel modules)} ve \sphinxstyleemphasis{aygıt sürücüler (device drivers)}
yoluyla davranış değişikliği oluşturulabilmektedir. İşletim sistemlerinin çoğu çekirdeğin bir parçası
gibi işlev görecek kodların yüklenmesine ve çalıştırılmasına olanak sağlamaktadır. Bu yöntemde
çekirdeğin yeniden derlenmesine gerek yoktur. Zaten çekirdek modülleri ve aygıt sürücüleri çalışmakta
olan çekirdek içerisine yüklenmektedir. Çekirdek modüllerini ve aygıt sürücüleri \sphinxstyleemphasis{kasalı bilgisayarlardaki
genişleme yuvalarına takılan kartlara benzetebiliriz.} Nasıl takılan bu kartlar donanımın bir parçası
haline geliyorsa çekirdek modülleri ve aygıt sürücüler de çekirdeğin bir parçası haline gelmektedir.
Linux’taki çekirdek modüllerinin ve aygıt sürücülerinin yazımı kursumuzun konularına dahil değildir.

\sphinxAtStartPar
\sphinxstylestrong{2. Çekirdek komut satırı parametreleri yoluyla}

\sphinxAtStartPar
Çekirdek yüklenip başlatılırken (initialize ederken) ismine \sphinxstyleemphasis{çekirdek komut satırı parametreleri
(kernel command line parameters)} denilen parametreler yoluyla çekirdeğin davranışı
değiştirilebilmektedir. Çekirdek komut satırı parametreleri \sphinxstyleemphasis{önyükleyici tarafından} çekirdeğe
iletilmektedir. Linux çekirdeğinin pek çok komut satırı parametresi vardır. Bu parametreler yoluyla
çekirdekte bazı davranış değişiklikleri oluşturulabilmektedir. Bunun için de çekirdeğin yeniden
derlenmesi gerekmez.

\sphinxAtStartPar
\sphinxstylestrong{3. Konfigürasyon parametreleri yoluyla}

\sphinxAtStartPar
Çekirdek derlenirken çekirdek kodlarına hiç dokunmadan bazı konfigürasyon parametreleri ile oynanarak
çekirdekte davranış değişiklikleri oluşturulabilmektedir. Çekirdek konfigüre edilirken bazı alt
sistemler çekirdekten çıkartılabilmekte, bazı alt sistemler üzerinde ince ayarlar yapılabilmektedir.
Tabii çekirdek konfigüre edilirken her ne kadar biz çekirdek kodlarını değiştirmiyor olsak da aslında
sembolik sabitler yoluyla arka planda derleme işlemine sokulan kodlar üzerinde de değişikler yapılmış
olmaktadır. O halde çekirdeğin konfigüre edilmesinin iki amacı vardır:
\begin{itemize}
\item {} 
\sphinxAtStartPar
Çekirdek içerisindeki alt sistemlere ilişkin bileşenlerin çekirdeğe eklenmesini ve çekirdekten
çıkartılmasını sağlamak.

\item {} 
\sphinxAtStartPar
Çekirdeğin üzerinde bazı davranış değişikliklerini oluşturmak.

\end{itemize}

\sphinxAtStartPar
Çekirdek konfigüre edildikten sonra yeniden derlenmelidir. Yani bu yöntem çekirdeğin yeniden
derlenmesini gerektirmektedir.

\sphinxAtStartPar
\sphinxstylestrong{4. Çekirdek kodlarında doğrudan değişiklik yaparak}

\sphinxAtStartPar
Çekirdek kodlarında doğrudan değişiklikler yapılıp çekirdek yeniden derlenebilir. Bu çekirdekte
davranış değişikliği oluşturmak için kullanılan en aşağı seviyeli yöntemdir.

\sphinxAtStartPar
\sphinxstylestrong{5. Çekirdek çalışırken komutlar ve mekanizmalar yoluyla}

\sphinxAtStartPar
Nihayet çekirdek çalışırken de çekirdeğin davranışı bazı komutlarla (bu komutlar bazı mekanizmaları
ve sistem fonksiyonlarını kullanmaktadır), konfigürasyon dosyalarıyla ve bazı mekanizmalarla da
(örneğin \sphinxstyleemphasis{proc} ve \sphinxstyleemphasis{sys} dosya sistemi yoluyla) çekirdek davranışları değiştirilebilmektedir. Tabii
bunun için de çekirdeğin yeniden derlenmesi gerekmez. Örneğin sistem çalışırken bir prosesin
açabileceği dosya sayısını \sphinxstyleemphasis{proc} dosya sistemi yoluyla şöyle değiştirebiliriz:

\begin{sphinxVerbatim}[commandchars=\\\{\}]
\PYGZdl{}\PYG{+w}{ }\PYG{n+nb}{echo}\PYG{+w}{ }\PYG{l+m}{2048}\PYG{+w}{ }\PYG{p}{|}\PYG{+w}{ }sudo\PYG{+w}{ }tee\PYG{+w}{ }/proc/sys/fs/file\PYGZhy{}max
\end{sphinxVerbatim}


\section{Çekirdek Komut Satırı Parametreleri}
\label{\detokenize{kernelcompilation:cekirdek-komut-satiri-parametreleri}}
\sphinxAtStartPar
Linux’ta çekirdek komut satırı parametreleri birbirinden boşluklarla ayrılmış yazılardan
oluşmaktadır. Bazı parametrelerin argümanları yoktur, bazılarının vardır. Eğer parametrenin bir
argümanı varsa \sphinxstyleemphasis{parametre=değer} biçiminde (\sphinxcode{\sphinxupquote{=}} karakteri boşluksuz olarak kullanılmalıdır),
yoksa yalnızca \sphinxstyleemphasis{parametre} biçiminde belirtilmektedir. Çekirdek komut satırı parametreleri tek
bir yazı biçiminde çekirdeğe aktarılmaktadır. Çekirdek bu komut satırı parametrelerini kendini
kullanıma hazır hale getirmenin (kendini initialize etmenin) ön aşamalarında parse eder ve bu
değerleri yapılandırma amacıyla kullanır. Tabii çekirdeğin komut satırı parametreleri tipik olarak
önyükleyici (bootloader) tarafından çekirdeğe iletilmektedir. Örneğin çekirdek komut satırı
parametreleri aşağıdaki gibi bir görünümde olabilir:

\begin{sphinxVerbatim}[commandchars=\\\{\}]
console=serial0,115200 console=tty1 root=PARTUUID=382d6f16\PYGZhy{}02 rootfstype=ext4 fsck.repair=yes
rootwait quiet splash plymouth.ignore\PYGZhy{}serial\PYGZhy{}consoles cfg80211.ieee80211\PYGZus{}regdom=TR
\end{sphinxVerbatim}

\sphinxAtStartPar
Linux çekirdeğinin onlarca farklı komut satırı parametresi vardır. Bunların çoğu spesifik konulara
ilişkindir ve ancak çekirdeğin yapısını iyi bilen kişiler tarafından anlamlandırılabilmektedir. Tabii
bazı parametrelerin anlamları herkes tarafından anlaşılacak kadar açıktır. Çekirdek parametrelerinin
dokümantasyonuna aşağıdaki bağlantıdan erişebilirsiniz:
\begin{quote}

\sphinxAtStartPar
\sphinxurl{https://www.kernel.org/doc/html/v4.14/admin-guide/kernel-parameters.html}
\end{quote}

\sphinxAtStartPar
Linux çekirdeği komut satırı parametrelerini parse ederken gerçekte olmayan (yani çekirdek tarafından
kullanılmayan) bir parametre gördüğünde onu dikkate almamaktadır. Ancak biz kendi parametrelerimizin de
çekirdek tarafından saklanmasını sağlayabiliriz. Bu konunun bazı ayrıntıları vardır.

\sphinxAtStartPar
Biz kursumuzda çeşitli bölümlerde çekirdeğin bazı komut satırı parametreleri hakkında bazı
açıklamalarda bulunacağız.


\section{Çekirdeği Yeniden Derlemenin Gerektiği Durumlar}
\label{\detokenize{kernelcompilation:cekirdegi-yeniden-derlemenin-gerektigi-durumlar}}
\sphinxAtStartPar
Bazı durumlarda çekirdeğin sıfırdan derlenmesi gerekebilmektedir. Çekirdeğin yeniden derlenmesinin
gerektiği tipik durumlar şunlardır:
\begin{itemize}
\item {} 
\sphinxAtStartPar
Bazı çekirdek modüllerinin ve aygıt sürücülerin çekirdek imajından çıkartılması ve dolayısıyla
çekirdeğin küçültülmesi için.

\item {} 
\sphinxAtStartPar
Yeni birtakım modüllerin ve aygıt sürücülerin çekirdek imajına eklenmesi için.

\item {} 
\sphinxAtStartPar
Çekirdeğe tamamen başka birtakım özelliklerin ve alt sistemlerin eklenmesi için.

\item {} 
\sphinxAtStartPar
Çekirdek üzerinde çekirdek komut satırı parametreleriyle sağlanamayacak bazı konfigürasyon
değişikliklerinin yapılabilmesi için.

\item {} 
\sphinxAtStartPar
Çekirdek kodlarında yapılan değişikliklerin etkin hale getirilmesi için.

\item {} 
\sphinxAtStartPar
Çekirdeğe yama (patch) yapılması için.

\item {} 
\sphinxAtStartPar
Yeni çıkan çekirdek kodlarının kullanılabilir hale getirilmesi için ya da eski çekirdeklerin
kullanımını sağlamak için.

\end{itemize}


\section{Çekirdek Derlemesi Hakkında}
\label{\detokenize{kernelcompilation:cekirdek-derlemesi-hakkinda}}
\sphinxAtStartPar
Bu bölümde çekirdek kodlarının nasıl derleneceği ve derlenmiş olan çekirdekle sistemin nasıl boot
edileceği üzerinde duracağız.


\subsection{Çapraz Derleme}
\label{\detokenize{kernelcompilation:capraz-derleme}}
\sphinxAtStartPar
Çekirdek derlemesi o anda çalışılan platform için yapılabileceği gibi başka bir platform için de
yapılabilmektedir. Aşağı seviyeli yazılım terminolojisinde başka bir sistem için yapılan derleme
işlemlerine \sphinxstyleemphasis{çapraz derleme (cross compile)} denilmektedir. Örneğin Intel tabanlı PC’lerde ARM
işlemcilerinin bulunduğu gömülü aygıtlar için Linux çekirdeğini derlemek isteyebiliriz. Bu bir
çapraz derleme faaliyetidir. Çapraz derleme işlemleri için \sphinxstyleemphasis{çapraz araç zincirlerinin (cross
toolchains)} kullanılması gerekmektedir. Linux çekirdeğinin derlenmesinde yalnızca C derleyicisi
değil pek çok yardımcı programlara da gereksinim duyulmaktadır.


\subsection{Çekirdek Kaynak Kodlarının İndirilmesi}
\label{\detokenize{kernelcompilation:cekirdek-kaynak-kodlarinin-indirilmesi}}
\sphinxAtStartPar
Çekirdeğin derlenmesi için öncelikle çekirdek kaynak kodlarının derleme yapılacak bilgisayara
indirilmesi gerekir. Bazı dağıtımlar default durumda çekirdeğin kaynak kodlarını da kurulum sırasında
makineye çekmektedir. Çekirdek kodları resmi olarak \sphinxcode{\sphinxupquote{kernel.org}} sitesinde bulundurulmaktadır.
Tarayıcıdan \sphinxcode{\sphinxupquote{kernel.org}} sitesine girip \sphinxcode{\sphinxupquote{pub/linux/kernel}} dizinine geçtiğinizde tüm yayınlanmış
çekirdek kodlarını göreceksiniz. İndirmeyi tarayıcıdan doğrudan yapabilirsiniz. Eğer indirmeyi komut
satırından \sphinxstyleemphasis{wget} programıyla yapmak istiyorsanız aşağıdaki URL’yi kullanabilirsiniz:

\begin{sphinxVerbatim}[commandchars=\\\{\}]
https://cdn.kernel.org/pub/linux/kernel/[MAJOR\PYGZus{}VERSION].x/linux\PYGZhy{}[VERSION].tar.xz
\end{sphinxVerbatim}

\sphinxAtStartPar
Buradaki \sphinxcode{\sphinxupquote{MAJOR\_VERSION}} 3, 4, 5, 6 gibi tek bir sayıyı belirtmektedir. \sphinxcode{\sphinxupquote{VERSION}} ise çekirdeğin
büyük ve küçük numaralarını belirtmektedir. Örneğin biz çekirdeğin 6.9.2 versiyonunu komut satırından
şöyle indirebiliriz:

\begin{sphinxVerbatim}[commandchars=\\\{\}]
\PYGZdl{}\PYG{+w}{ }wget\PYG{+w}{ }https://cdn.kernel.org/pub/linux/kernel/v6.x/linux\PYGZhy{}6.9.2.tar.xz
\end{sphinxVerbatim}


\subsubsection{Sıkıştırma Formatları ve tar Komutu}
\label{\detokenize{kernelcompilation:sikistirma-formatlari-ve-tar-komutu}}
\sphinxAtStartPar
Çekirdek kodları indirildikten sonra açılması gerekir. Açma işlemi \sphinxstyleemphasis{tar} komutuyla aşağıdaki gibi
yapılabilir:

\begin{sphinxVerbatim}[commandchars=\\\{\}]
\PYGZdl{}\PYG{+w}{ }tar\PYG{+w}{ }\PYGZhy{}xvJf\PYG{+w}{ }linux\PYGZhy{}5.15.12.tar.xz
\end{sphinxVerbatim}

\sphinxAtStartPar
UNIX/Linux sistemlerinde kullanılan \sphinxstyleemphasis{tar} programı sıkıştırma yapmamaktadır. Yalnızca dosyaları uç uca
ekleyip onları tek bir dosya haline getirmektedir. Bu sayede dosyalar dosya sisteminde daha az yer
kaplar hale getirilmektedir. Aynı zamanda onların iletilmesi ve kopyalanması da kolaylaştırılmaktadır.
Sıkıştırma programları ise aslında tek bir dosyayı sıkıştırmaktadır. O halde UNIX/Linux sistemlerinde
kullanıcılar bir grup dosyayı sıkıştırmak için önce onları tar’layıp tek dosya haline getirirler, sonra
bu dosyayı sıkıştırırlar. Bu işlemler sonucunda dosya uzantısı \sphinxcode{\sphinxupquote{.tar.gz}} gibi, \sphinxcode{\sphinxupquote{.tar.xz}} gibi
dosyanın hem tar’lanmış hem de sıkıştırılmış olduğunu belirten biçimde olur. Açım sırasında da önce
sıkıştırılan dosyalar açılır, buradan \sphinxcode{\sphinxupquote{.tar}} dosyası elde edilir; sonra da bu \sphinxcode{\sphinxupquote{.tar}} dosyasının
içerisindeki dosyalar dışarı çıkartılır.

\sphinxAtStartPar
Linux’ta çeşitli sıkıştırma formatları kullanılmaktadır:
\begin{itemize}
\item {} 
\sphinxAtStartPar
\sphinxstylestrong{gzip} — uzantı: \sphinxcode{\sphinxupquote{.gz}}

\item {} 
\sphinxAtStartPar
\sphinxstylestrong{bz2} — uzantı: \sphinxcode{\sphinxupquote{.bz2}}

\item {} 
\sphinxAtStartPar
\sphinxstylestrong{xz} — uzantı: \sphinxcode{\sphinxupquote{.xz}}

\item {} 
\sphinxAtStartPar
\sphinxstylestrong{zip} — uzantı: \sphinxcode{\sphinxupquote{.zip}} ya da \sphinxcode{\sphinxupquote{.z}}

\end{itemize}

\sphinxAtStartPar
Bu formatlar sıkıştırma bakımından farklı performans göstermektedir. Ancak formatın sıkıştırma
performansı ne kadar yüksekse işlem yapma süresi de o kadar uzun olmaktadır. Tipik sıkıştırma
performans sıralaması şöyledir:

\begin{sphinxVerbatim}[commandchars=\\\{\}]
xz  \PYGZgt{}  bz2  \PYGZgt{}  gzip ≈ zip
\end{sphinxVerbatim}

\sphinxAtStartPar
Yani en iyi sıkıştırma \sphinxstyleemphasis{xz} formatında, daha sonra \sphinxstyleemphasis{bz2} formatında, daha sonra da \sphinxstyleemphasis{gzip}
formatındadır. \sphinxstyleemphasis{gzip} ile \sphinxstyleemphasis{zip} aynı algoritmaları kullandığından performansları birbirine benzerdir.

\sphinxAtStartPar
Her format için ayrı araçlar mevcuttur:


\begin{savenotes}\sphinxattablestart
\sphinxthistablewithglobalstyle
\centering
\begin{tabular}[t]{\X{20}{70}\X{25}{70}\X{25}{70}}
\sphinxtoprule
\begin{varwidth}[t]{\sphinxcolwidth{1}{3}}
\sphinxstyletheadfamily \sphinxAtStartPar
Format
\sphinxbeforeendvarwidth
\end{varwidth}%
&\begin{varwidth}[t]{\sphinxcolwidth{1}{3}}
\sphinxstyletheadfamily \sphinxAtStartPar
Sıkıştırma
\sphinxbeforeendvarwidth
\end{varwidth}%
&\begin{varwidth}[t]{\sphinxcolwidth{1}{3}}
\sphinxstyletheadfamily \sphinxAtStartPar
Açma
\sphinxbeforeendvarwidth
\end{varwidth}%
\\
\sphinxmidrule
\sphinxtableatstartofbodyhook\begin{varwidth}[t]{\sphinxcolwidth{1}{3}}
\sphinxAtStartPar
gzip
\sphinxbeforeendvarwidth
\end{varwidth}%
&\begin{varwidth}[t]{\sphinxcolwidth{1}{3}}
\sphinxAtStartPar
\sphinxcode{\sphinxupquote{gzip}}
\sphinxbeforeendvarwidth
\end{varwidth}%
&\begin{varwidth}[t]{\sphinxcolwidth{1}{3}}
\sphinxAtStartPar
\sphinxcode{\sphinxupquote{gunzip}}
\sphinxbeforeendvarwidth
\end{varwidth}%
\\
\sphinxhline\begin{varwidth}[t]{\sphinxcolwidth{1}{3}}
\sphinxAtStartPar
bz2
\sphinxbeforeendvarwidth
\end{varwidth}%
&\begin{varwidth}[t]{\sphinxcolwidth{1}{3}}
\sphinxAtStartPar
\sphinxcode{\sphinxupquote{bzip2}}
\sphinxbeforeendvarwidth
\end{varwidth}%
&\begin{varwidth}[t]{\sphinxcolwidth{1}{3}}
\sphinxAtStartPar
\sphinxcode{\sphinxupquote{bunzip2}}
\sphinxbeforeendvarwidth
\end{varwidth}%
\\
\sphinxhline\begin{varwidth}[t]{\sphinxcolwidth{1}{3}}
\sphinxAtStartPar
xz
\sphinxbeforeendvarwidth
\end{varwidth}%
&\begin{varwidth}[t]{\sphinxcolwidth{1}{3}}
\sphinxAtStartPar
\sphinxcode{\sphinxupquote{xz}}
\sphinxbeforeendvarwidth
\end{varwidth}%
&\begin{varwidth}[t]{\sphinxcolwidth{1}{3}}
\sphinxAtStartPar
\sphinxcode{\sphinxupquote{unxz}}
\sphinxbeforeendvarwidth
\end{varwidth}%
\\
\sphinxhline\begin{varwidth}[t]{\sphinxcolwidth{1}{3}}
\sphinxAtStartPar
zip
\sphinxbeforeendvarwidth
\end{varwidth}%
&\begin{varwidth}[t]{\sphinxcolwidth{1}{3}}
\sphinxAtStartPar
\sphinxcode{\sphinxupquote{zip}}
\sphinxbeforeendvarwidth
\end{varwidth}%
&\begin{varwidth}[t]{\sphinxcolwidth{1}{3}}
\sphinxAtStartPar
\sphinxcode{\sphinxupquote{unzip}}
\sphinxbeforeendvarwidth
\end{varwidth}%
\\
\sphinxbottomrule
\end{tabular}
\sphinxtableafterendhook\par
\sphinxattableend\end{savenotes}

\sphinxAtStartPar
\sphinxstyleemphasis{tar} programının en çok kullanılan seçenekleri şunlardır: \sphinxcode{\sphinxupquote{\sphinxhyphen{}c}} (tar’lamak için), \sphinxcode{\sphinxupquote{\sphinxhyphen{}x}} (açmak
için), \sphinxcode{\sphinxupquote{\sphinxhyphen{}v}} (ayrıntılı çıktı), \sphinxcode{\sphinxupquote{\sphinxhyphen{}f}} (\sphinxcode{\sphinxupquote{.tar}} dosyasını belirtmek için; seçenek listesinin
sonunda olmalıdır).

\begin{sphinxadmonition}{warning}{Uyarı:}
\sphinxAtStartPar
\sphinxcode{\sphinxupquote{\sphinxhyphen{}f}} seçeneğinin seçenek listesinde en sonda yer alması gerektiğine dikkat ediniz.
\end{sphinxadmonition}

\sphinxAtStartPar
\sphinxstylestrong{gzip ile örnekler:}

\begin{sphinxVerbatim}[commandchars=\\\{\}]
\PYG{c+c1}{\PYGZsh{} Yalnızca tar\PYGZsq{}lama:}
\PYGZdl{}\PYG{+w}{ }tar\PYG{+w}{ }\PYGZhy{}cf\PYG{+w}{ }test.tar\PYG{+w}{ }x.txt\PYG{+w}{ }y.txt\PYG{+w}{ }z.txt

\PYG{c+c1}{\PYGZsh{} Tek adımda tar\PYGZsq{}la + gzip sıkıştır (\PYGZhy{}z seçeneği):}
\PYGZdl{}\PYG{+w}{ }tar\PYG{+w}{ }\PYGZhy{}cvzf\PYG{+w}{ }test.tar.gz\PYG{+w}{ }x.txt\PYG{+w}{ }y.txt

\PYG{c+c1}{\PYGZsh{} Aç:}
\PYGZdl{}\PYG{+w}{ }tar\PYG{+w}{ }\PYGZhy{}xvzf\PYG{+w}{ }test.tar.gz
\end{sphinxVerbatim}

\begin{sphinxadmonition}{note}{Not:}
\sphinxAtStartPar
\sphinxstyleemphasis{gzip} ve \sphinxstyleemphasis{gunzip} programlarının kaynak dosyayı sildiğine dikkat ediniz.
\end{sphinxadmonition}

\sphinxAtStartPar
\sphinxstylestrong{bzip2 ile örnekler:}

\begin{sphinxVerbatim}[commandchars=\\\{\}]
\PYG{c+c1}{\PYGZsh{} Tek adımda tar\PYGZsq{}la + bzip2 sıkıştır (\PYGZhy{}j seçeneği):}
\PYGZdl{}\PYG{+w}{ }tar\PYG{+w}{ }\PYGZhy{}cvjf\PYG{+w}{ }test.tar.bz2\PYG{+w}{ }x.txt\PYG{+w}{ }y.txt

\PYG{c+c1}{\PYGZsh{} Aç:}
\PYGZdl{}\PYG{+w}{ }tar\PYG{+w}{ }\PYGZhy{}xvjf\PYG{+w}{ }test.tar.bz2
\end{sphinxVerbatim}

\sphinxAtStartPar
\sphinxstylestrong{xz ile örnekler:}

\begin{sphinxVerbatim}[commandchars=\\\{\}]
\PYG{c+c1}{\PYGZsh{} Tek adımda tar\PYGZsq{}la + xz sıkıştır (\PYGZhy{}J seçeneği, büyük harf):}
\PYGZdl{}\PYG{+w}{ }tar\PYG{+w}{ }\PYGZhy{}cvJf\PYG{+w}{ }test.tar.xz\PYG{+w}{ }x.txt\PYG{+w}{ }y.txt

\PYG{c+c1}{\PYGZsh{} Aç:}
\PYGZdl{}\PYG{+w}{ }tar\PYG{+w}{ }\PYGZhy{}xvJf\PYG{+w}{ }test.tar.xz
\end{sphinxVerbatim}

\begin{sphinxadmonition}{note}{Not:}
\sphinxAtStartPar
\sphinxstyleemphasis{tar} seçeneklerinde \sphinxcode{\sphinxupquote{\sphinxhyphen{}z}} gzip için, \sphinxcode{\sphinxupquote{\sphinxhyphen{}j}} bzip2 için, \sphinxcode{\sphinxupquote{\sphinxhyphen{}J}} (büyük harf) xz için
kullanılmaktadır.
\end{sphinxadmonition}


\subsection{Dağıtıma Özgü Çekirdek Kaynak Kodları}
\label{\detokenize{kernelcompilation:dagitima-ozgu-cekirdek-kaynak-kodlari}}
\sphinxAtStartPar
Dağıtımlar (Ubuntu, Mint, Fedora gibi) çekirdek kodlarında küçük değişiklikler ve kendilerine özgü
özelleştirmeler ve yamamalar da yapabilmektedir. Debian tabanlı sistemlerde o anda makinede yüklü olan
mevcut çekirdeğin ilgili dağıtıma ilişkin kaynak kodlarını indirmek için aşağıdaki komutu
kullanabilirsiniz:

\begin{sphinxVerbatim}[commandchars=\\\{\}]
\PYGZdl{}\PYG{+w}{ }sudo\PYG{+w}{ }apt\PYGZhy{}get\PYG{+w}{ }install\PYG{+w}{ }linux\PYGZhy{}source
\end{sphinxVerbatim}

\sphinxAtStartPar
Burada yükleme \sphinxcode{\sphinxupquote{/usr/src}} dizinine yapılacaktır. Ancak bu komut doğrudan sıkıştırılmış dosyayı
indirmektedir; açım yapmamaktadır. Ayrıca belirli bir versiyona ilişkin kaynak kodlar da indirilebilir.
Bunun için komutta \sphinxcode{\sphinxupquote{linux\sphinxhyphen{}source}} argümanına istenilen versiyonun majör, minör ve patch numarası
\sphinxcode{\sphinxupquote{\sphinxhyphen{}majör.minör.patch}} biçiminde eklenir. Örneğin:

\begin{sphinxVerbatim}[commandchars=\\\{\}]
\PYGZdl{}\PYG{+w}{ }sudo\PYG{+w}{ }apt\PYGZhy{}get\PYG{+w}{ }install\PYG{+w}{ }linux\PYGZhy{}source\PYGZhy{}6.8.0
\end{sphinxVerbatim}

\begin{sphinxadmonition}{note}{Not:}
\sphinxAtStartPar
Bu biçimde indirilen çekirdek kaynak kodları Debian ya da Ubuntu depolarından indirilmektedir.
Bunlar bu dağıtımlar tarafından yamanmış kodlardır. Örneğin Mint’te çalışıyorsanız indirdiğiniz
kodlar Ubuntu için yamanmış kodlar olacaktır.
\end{sphinxadmonition}

\sphinxAtStartPar
\sphinxstylestrong{BBB (BeagleBone Black) için:}

\sphinxAtStartPar
BBB için bazı özelleştirmelerin de yapılmış olduğu kaynak kodların indirilip derlenmesi birtakım
kolaylıklar sağlamaktadır:

\begin{sphinxVerbatim}[commandchars=\\\{\}]
\PYGZdl{}\PYG{+w}{ }git\PYG{+w}{ }clone\PYG{+w}{ }https://github.com/beagleboard/linux.git
\end{sphinxVerbatim}

\sphinxAtStartPar
\sphinxstylestrong{Raspberry Pi için:}

\sphinxAtStartPar
Raspberry Pi’a özgü daha güncel aygıt sürücüler ve aygıt ağacı dosyalarını içeren Linux kaynak
kodlarının projenin kendi sitesinden indirilmesi daha uygun olur:

\begin{sphinxVerbatim}[commandchars=\\\{\}]
\PYGZdl{}\PYG{+w}{ }git\PYG{+w}{ }clone\PYG{+w}{ }\PYGZhy{}\PYGZhy{}depth\PYG{o}{=}\PYG{l+m}{1}\PYG{+w}{ }https://github.com/raspberrypi/linux.git
\end{sphinxVerbatim}


\section{Çekirdek Versiyon Numaralandırması}
\label{\detokenize{kernelcompilation:cekirdek-versiyon-numaralandirmasi}}
\sphinxAtStartPar
Linux kaynak kodlarının versiyonlanması eskiden daha farklıydı. Çekirdeğin 2.6 versiyonlarından sonra
versiyon numaralandırma sistemi değiştirildi. Eskiden (2.6 ve öncesinde) versiyon numaraları çok yavaş
ilerletiliyordu. 2.6 sonrasındaki yeni versiyonlamada versiyon numaraları daha hızlı ilerletilmeye
başlanmıştır. Bugün kullanılan Linux versiyonları nokta ile ayrılmış üç sayıdan oluşmaktadır:

\begin{sphinxVerbatim}[commandchars=\\\{\}]
Majör.Minör.Patch
\end{sphinxVerbatim}

\sphinxAtStartPar
Buradaki \sphinxstyleemphasis{majör numara} büyük ilerlemeleri, \sphinxstyleemphasis{minör numara} ise küçük ilerlemeleri belirtmektedir.
Eskiden (2.6 ve öncesinde) tek sayı olan minör numaralar \sphinxstyleemphasis{geliştirme versiyonlarını (ya da beta
versiyonlarını)}, çift olanlar ise stabil hale getirilmiş dağıtılan versiyonları belirtiyordu. Ancak
2.6 sonrasında artık tek ve çift minör numaralar arasında böyle bir anlam farklılığı kalmamıştır.
\sphinxstyleemphasis{Patch numarası} birtakım böceklerin giderildiği ya da çok küçük yeniliklerin çekirdeğe dahil edildiği
versiyonları belirtmektedir.

\sphinxAtStartPar
Linux kaynak kodları konfigüre edilip derlendiğinde çekirdek imajının ismine bir alan daha
eklenebilmektedir. Bu alana biz \sphinxstyleemphasis{Extra} alanı diyeceğiz. Bu durumda çekirdek imaj ismi şu hale
gelecektir:

\begin{sphinxVerbatim}[commandchars=\\\{\}]
Majör.Minör.Patch\PYGZhy{}Extra
\end{sphinxVerbatim}

\sphinxAtStartPar
\sphinxstyleemphasis{Extra} için \sphinxcode{\sphinxupquote{\sphinxhyphen{}rcX}}, \sphinxcode{\sphinxupquote{\sphinxhyphen{}stable}}, \sphinxcode{\sphinxupquote{\sphinxhyphen{}custom}}, \sphinxcode{\sphinxupquote{\sphinxhyphen{}generic}} gibi sözcükler kullanılabilir. Bu
\sphinxstyleemphasis{extra} alanı tamamen derleme işlemini yapan kişi ya da kurumun isteğine göre belirlenmiş bir
alandır:
\begin{itemize}
\item {} 
\sphinxAtStartPar
\sphinxstyleemphasis{rcX} — \sphinxstyleemphasis{release candidate} (yayın adayı) anlamındadır; \sphinxcode{\sphinxupquote{X}} bir sayı belirtir.

\item {} 
\sphinxAtStartPar
\sphinxstyleemphasis{stable} — dağıtılan sürümün \sphinxstyleemphasis{kararlı sürüm} olduğunu belirtir.

\item {} 
\sphinxAtStartPar
\sphinxstyleemphasis{custom} — sistem programcısının çekirdekte birtakım değişiklikler yaptığını belirtir.

\item {} 
\sphinxAtStartPar
\sphinxstyleemphasis{generic} — çekirdeğin \sphinxstyleemphasis{genel kullanım için konfigüre edilmiş} olduğunu belirtir; dağıtımlar
tarafından sıkça kullanılır.

\item {} 
\sphinxAtStartPar
\sphinxstyleemphasis{realtime} — genellikle gerçek zamanlı bir yapılandırmada kullanılmaktadır.

\end{itemize}

\sphinxAtStartPar
\sphinxstyleemphasis{generic} ve \sphinxstyleemphasis{realtime} sözcüklerinin öncesinde \sphinxcode{\sphinxupquote{\sphinxhyphen{}N\sphinxhyphen{}}} biçiminde bir sayı da bulunabilmektedir. Bu
sayı dağıtıma özgü yama ya da derleme numarasını belirtmektedir. Örneğin:

\begin{sphinxVerbatim}[commandchars=\\\{\}]
6.8.0\PYGZhy{}51\PYGZhy{}generic
\end{sphinxVerbatim}

\sphinxAtStartPar
Burada \sphinxcode{\sphinxupquote{\sphinxhyphen{}51}} yapılandırmaya özgü bir numarayı, \sphinxcode{\sphinxupquote{\sphinxhyphen{}generic}} ise yapılandırmanın genel kullanım için
olduğunu belirtmektedir.

\sphinxAtStartPar
Çalışmakta olan Linux sistemi hakkında bilgiler \sphinxstyleemphasis{“uname \sphinxhyphen{}a”} komutu ile elde edilebilir:

\begin{sphinxVerbatim}[commandchars=\\\{\}]
\PYGZdl{}\PYG{+w}{ }uname\PYG{+w}{ }\PYGZhy{}a
Linux\PYG{+w}{ }kaan\PYGZhy{}virtual\PYGZhy{}machine\PYG{+w}{ }\PYG{l+m}{5}.15.0\PYGZhy{}91\PYGZhy{}generic\PYG{+w}{ }Ubuntu\PYG{+w}{ }SMP\PYG{+w}{ }Tue\PYG{+w}{ }Nov\PYG{+w}{ }\PYG{l+m}{14}\PYG{+w}{ }\PYG{l+m}{13}:30:08\PYG{+w}{ }UTC\PYG{+w}{ }\PYG{l+m}{2023}\PYG{+w}{ }x86\PYGZus{}64\PYG{+w}{ }x86\PYGZus{}64\PYG{+w}{ }x86\PYGZus{}64\PYG{+w}{ }GNU/Linux
\end{sphinxVerbatim}

\sphinxAtStartPar
Bu bilgi içerisinden yalnızca çekirdek versiyonu görüntülenmek isteniyorsa \sphinxstyleemphasis{“uname \sphinxhyphen{}r”} seçeneği
kullanılmalıdır:

\begin{sphinxVerbatim}[commandchars=\\\{\}]
\PYGZdl{}\PYG{+w}{ }uname\PYG{+w}{ }\PYGZhy{}r
\PYG{l+m}{6}.8.0\PYGZhy{}51\PYGZhy{}generic
\end{sphinxVerbatim}

\sphinxAtStartPar
Buradan biz çekirdeğin \sphinxcode{\sphinxupquote{6.8.0}} sürümünün kullanıldığını anlıyoruz. Burada genel yapılandırılmış bir
çekirdek söz konusudur. \sphinxcode{\sphinxupquote{51}} sayısı dağıtıma özgü yama ya da derleme numarasını belirtir.

\sphinxAtStartPar
Daha önceden de belirttiğimiz gibi \sphinxstyleemphasis{uname} komutu bu bilgileri \sphinxstyleemphasis{proc} dosya sisteminin içerisinden
almaktadır. Örneğin:

\begin{sphinxVerbatim}[commandchars=\\\{\}]
\PYGZdl{}\PYG{+w}{ }cat\PYG{+w}{ }/proc/version
Linux\PYG{+w}{ }version\PYG{+w}{ }\PYG{l+m}{5}.15.0\PYGZhy{}91\PYGZhy{}generic\PYG{+w}{ }\PYG{o}{(}buildd@lcy02\PYGZhy{}amd64\PYGZhy{}045\PYG{o}{)}\PYG{+w}{ }\PYG{o}{(}gcc\PYG{+w}{ }\PYG{o}{(}Ubuntu\PYG{+w}{ }\PYG{l+m}{11}.4.0\PYGZhy{}1ubuntu1\PYGZti{}22.04\PYG{o}{)}\PYG{+w}{ }\PYG{l+m}{11}.4.0,
GNU\PYG{+w}{ }ld\PYG{+w}{ }\PYG{o}{(}GNU\PYG{+w}{ }Binutils\PYG{+w}{ }\PYG{k}{for}\PYG{+w}{ }Ubuntu\PYG{o}{)}\PYG{+w}{ }\PYG{l+m}{2}.38\PYG{o}{)}\PYG{+w}{ }\PYG{c+c1}{\PYGZsh{}101\PYGZhy{}Ubuntu SMP Tue Nov 14 13:30:08 UTC 2023}
\end{sphinxVerbatim}


\section{Derleme Araçları}
\label{\detokenize{kernelcompilation:derleme-araclari}}
\sphinxAtStartPar
Bilindiği gibi büyük projelerin derlenmesi için \sphinxstyleemphasis{build otomasyon araçları (build automation tools)}
denilen araçlar kullanılmaktadır. Bunların en yaygın kullanılanı \sphinxstyleemphasis{make} denilen araçtır. Linux
çekirdeklerinin derlenmesi de \sphinxstyleemphasis{make} aracı ile yapılmaktadır. Ancak Linux çekirdeklerinin derlenmesinde
projeye özgü bazı yapılar ve yöntemler de kullanılmıştır. Buna \sphinxstyleemphasis{KConfig sistemi} ya da \sphinxstyleemphasis{KBuild sistemi}
denilmektedir. (\sphinxstyleemphasis{KConfig} sistemi ya da \sphinxstyleemphasis{KBuild} sistemi Linux çekirdeğinin dışında aşağı seviyeli pek
çok proje tarafından da kullanılmaktadır.) Biz önce çekirdek derleme işleminin hangi adımlardan
geçilerek yapılacağını göreceğiz. Sonra çekirdeğin önemli konfigürasyon parametreleri üzerinde biraz
duracağız. Sonra da çekirdekte bazı değişiklikler yapıp sistemin değiştirilmiş çekirdekle açılmasını
sağlayacağız.


\section{Çekirdek Derleme Adımları}
\label{\detokenize{kernelcompilation:cekirdek-derleme-adimlari}}
\sphinxAtStartPar
Linux’ta çekirdek derlemesi tipik olarak aşağıdaki aşamalardan geçilerek gerçekleştirilmektedir.


\subsection{Adım 1: Gerekli Araçların Kurulması}
\label{\detokenize{kernelcompilation:adim-1-gerekli-araclarin-kurulmasi}}
\sphinxAtStartPar
Derleme öncesinde derlemenin yapılacağı makinede bazı programların yüklenmiş olması gerekmektedir.
Çünkü \sphinxstyleemphasis{KBuild} sistemi yalnızca binary araçları değil bazı başka kütüphaneleri ve araçları da
kullanmaktadır. Çekirdeğin derlenmesi için gerekebilecek programları şöyle yükleyebilirsiniz
(burada Debian türevi bir dağıtımın kullanıldığını varsayacağız):

\begin{sphinxVerbatim}[commandchars=\\\{\}]
\PYGZdl{}\PYG{+w}{ }sudo\PYG{+w}{ }apt\PYG{+w}{ }update
\PYGZdl{}\PYG{+w}{ }sudo\PYG{+w}{ }apt\PYG{+w}{ }install\PYG{+w}{ }build\PYGZhy{}essential\PYG{+w}{ }libncurses\PYGZhy{}dev\PYG{+w}{ }bison\PYG{+w}{ }flex\PYG{+w}{ }libssl\PYGZhy{}dev\PYG{+w}{ }libelf\PYGZhy{}dev\PYG{+w}{ }dwarves
\end{sphinxVerbatim}


\subsection{Adım 2: Kaynak Kodların İndirilmesi ve Açılması}
\label{\detokenize{kernelcompilation:adim-2-kaynak-kodlarin-indirilmesi-ve-acilmasi}}
\sphinxAtStartPar
Çekirdek kodları indirilir ve açılır. Biz bu konuyu yukarıda ele almıştık. İndirmeyi şöyle
yapabiliriz:

\begin{sphinxVerbatim}[commandchars=\\\{\}]
\PYGZdl{}\PYG{+w}{ }wget\PYG{+w}{ }https://cdn.kernel.org/pub/linux/kernel/v6.x/linux\PYGZhy{}6.9.2.tar.xz
\PYGZdl{}\PYG{+w}{ }tar\PYG{+w}{ }\PYGZhy{}xvJf\PYG{+w}{ }linux\PYGZhy{}6.9.2.tar.xz
\end{sphinxVerbatim}

\sphinxAtStartPar
Bu işlemden sonra \sphinxcode{\sphinxupquote{linux\sphinxhyphen{}6.9.2}} isminde bir dizin oluşturulacaktır.


\subsection{Adım 3: Çekirdeğin Konfigüre Edilmesi}
\label{\detokenize{kernelcompilation:adim-3-cekirdegin-konfigure-edilmesi}}
\sphinxAtStartPar
Çekirdek derlenmeden önce konfigüre edilmelidir. Çekirdeğin konfigüre edilmesi birtakım çekirdek
özelliklerinin belirlenmesi anlamına gelmektedir. Bu belirmeler konfigürasyon dosyası yoluyla
yapılmaktadır. Konfigürasyon dosyası kaynak kod ağacının kök dizininde (örneğimizde \sphinxcode{\sphinxupquote{linux\sphinxhyphen{}6.9.2}}
dizini) \sphinxcode{\sphinxupquote{.config}} ismiyle bulunmalıdır.

\sphinxAtStartPar
Bu \sphinxcode{\sphinxupquote{.config}} dosyası default durumda kaynak dosyaların kök dizininde yoktur; derlemeyi yapacak
kişi tarafından oluşturulması gerekmektedir. Çekirdek konfigürasyon parametreleri oldukça fazladır
ve bunların bazılarının anlamlandırılması özel bilgi gerektirmektedir. Çekirdek konfigürasyon
parametreleri çok fazla olduğu için bazı genel amaçları karşılayacak biçimde default değerlerle
önceden oluşturulmuş konfigürasyon dosyaları kaynak kod ağacında \sphinxcode{\sphinxupquote{arch/\textless{}mimari\textgreater{}/configs}} dizininin
içerisinde bulunmaktadır. Örneğin Intel x86 mimarisi için bu default konfigürasyon dosyaları şöyledir:

\begin{sphinxVerbatim}[commandchars=\\\{\}]
\PYGZdl{}\PYG{+w}{ }ls\PYG{+w}{ }arch/x86/configs
hardening.config\PYG{+w}{  }i386\PYGZus{}defconfig\PYG{+w}{  }tiny.config\PYG{+w}{  }x86\PYGZus{}64\PYGZus{}defconfig\PYG{+w}{  }xen.config
\end{sphinxVerbatim}

\sphinxAtStartPar
64 bit Linux sistemleri için \sphinxcode{\sphinxupquote{x86\_64\_defconfig}} dosyasını kullanabiliriz. O halde bu dosyayı
kaynak dosyaların bulunduğu kök dizine \sphinxcode{\sphinxupquote{.config}} ismiyle kopyalayabiliriz (bütün işlemlerde
çekirdek kaynak kodlarının kök dizininde bulunduğumuzu varsayıyoruz):

\begin{sphinxVerbatim}[commandchars=\\\{\}]
\PYGZdl{}\PYG{+w}{ }cp\PYG{+w}{ }arch/x86/configs/x86\PYGZus{}64\PYGZus{}defconfig\PYG{+w}{ }.config
\end{sphinxVerbatim}

\begin{sphinxadmonition}{note}{Not:}
\sphinxAtStartPar
Linux kaynak kodlarındaki default konfigürasyon dosyaları minimalist biçimde konfigüre
edilmiştir. Bu nedenle pek çok modül bu default konfigürasyon dosyalarında işaretlenmiş değildir.
Bu tür denemeleri zaten çalışan çekirdeğin derlenmesinde kullanılan konfigürasyon dosyalarından
hareketle yaparsanız daha fazla modül dosyası oluşturulabilir ancak daha az zahmet çekebilirsiniz.
\end{sphinxadmonition}

\sphinxAtStartPar
Burada bir noktaya dikkatinizi çekmek istiyoruz. Çekirdek kaynak kodlarındaki
\sphinxcode{\sphinxupquote{arch/\textless{}platform\textgreater{}/configs}} dizinindeki \sphinxcode{\sphinxupquote{x86\_64\_defconfig}} konfigürasyon dosyası \sphinxcode{\sphinxupquote{.config}}
ismiyle kopyalandıktan sonra ayrıca \sphinxstyleemphasis{“make menuconfig”} ya da \sphinxstyleemphasis{“make oldconfig”} gibi bir işlemle
onun satırlarına eklenecek çekirdekte bulunan yeni birtakım özelliklere ilişkin bazı default
değerlerin de eklenmesi gerekir.

\sphinxAtStartPar
Linux sistemlerinde genel olarak \sphinxcode{\sphinxupquote{/boot}} dizini içerisinde \sphinxcode{\sphinxupquote{config\sphinxhyphen{}\textless{}çekirdek\_sürümü\textgreater{}}} ismiyle
mevcut çekirdeğe ilişkin konfigürasyon dosyası bulundurulmaktadır.

\sphinxAtStartPar
\sphinxcode{\sphinxupquote{.config}} dosyasını oluşturmanın alternatif yolları şöyle özetlenebilir:
\begin{description}
\sphinxlineitem{\sphinxcode{\sphinxupquote{make defconfig}}}
\sphinxAtStartPar
Çalışılan sisteme uygun olan konfigürasyon dosyasını temel alarak mevcut donanım bileşenlerini
de gözden geçirerek sistemin açılması için gerekli minimal bir konfigürasyon dosyasını \sphinxcode{\sphinxupquote{.config}}
ismiyle oluşturur. Örneğin 64 bit Intel sisteminde \sphinxcode{\sphinxupquote{make defconfig}} çalıştırıldığında
\sphinxcode{\sphinxupquote{arch/x86/configs/x86\_64\_defconfig}} dosyası temel alınır.

\sphinxlineitem{\sphinxcode{\sphinxupquote{make oldconfig}}}
\sphinxAtStartPar
Bu seçeneği kullanmak için kaynak kök dizinde \sphinxcode{\sphinxupquote{.config}} dosyasının bulunuyor olması gerekir.
\sphinxstyleemphasis{KConfig} dosyalarındaki ve kaynak dosya ağacındaki diğer değişiklikleri de göz önüne alarak bu
eski \sphinxcode{\sphinxupquote{.config}} dosyasını yeni versiyona uyumlandırmaktadır. Başka bir deyişle \sphinxstyleemphasis{“make oldconfig”}
eski bir konfigürasyon dosyasını yeni çekirdekler için uyumlandırmaktadır.

\sphinxlineitem{\sphinxcode{\sphinxupquote{make \textless{}platform\textgreater{}\_defconfig}}}
\sphinxAtStartPar
Belli bir platformun default konfigürasyon dosyasını \sphinxcode{\sphinxupquote{.config}} dosyası olarak kaydeder. Örneğin
Intel makinelerinde çalışırken BBB için \sphinxcode{\sphinxupquote{make bb.org\_defconfig}} komutu uygulanabilir.

\sphinxlineitem{\sphinxcode{\sphinxupquote{make modules}}}
\sphinxAtStartPar
Yalnızca modülleri (aygıt sürücü dosyaları) derler; çekirdek derlemesi yapmaz. Yalnızca \sphinxstyleemphasis{make}
işlemi zaten aynı zamanda bu işlemi de yapmaktadır.

\sphinxlineitem{\sphinxcode{\sphinxupquote{make uninstall}}}
\sphinxAtStartPar
\sphinxstyleemphasis{“make install”} işlemi ile yapılanları geri alır.

\end{description}

\sphinxAtStartPar
Aşağıdaki \sphinxstyleemphasis{make xxxconfig} komutları seyrek kullanılmaktadır:
\begin{description}
\sphinxlineitem{\sphinxcode{\sphinxupquote{make allnoconfig}}}
\sphinxAtStartPar
Tüm seçenekleri \sphinxstyleemphasis{hayır (no)} olarak ayarlar (minimal özellikler).

\sphinxlineitem{\sphinxcode{\sphinxupquote{make allyesconfig}}}
\sphinxAtStartPar
Tüm seçenekleri \sphinxstyleemphasis{evet (yes)} olarak ayarlar (maksimum özellikler).

\sphinxlineitem{\sphinxcode{\sphinxupquote{make allmodconfig}}}
\sphinxAtStartPar
Tüm aygıt sürücülerin çekirdeğin dışında modül biçiminde derleneceğini belirtir.

\sphinxlineitem{\sphinxcode{\sphinxupquote{make localmodconfig}}}
\sphinxAtStartPar
Sistemde o anda yüklü modüllere dayalı bir yapılandırma dosyası oluşturur.

\sphinxlineitem{\sphinxcode{\sphinxupquote{make silentoldconfig}}}
\sphinxAtStartPar
Yeni seçenekler için onları görmezden gelir; yeni özellikler \sphinxcode{\sphinxupquote{.config}} dosyasına yansıtılmaz.

\sphinxlineitem{\sphinxcode{\sphinxupquote{make dtbs}}}
\sphinxAtStartPar
Kaynak kod ağacında \sphinxcode{\sphinxupquote{/arch/platform/boot/dts}} dizinindeki aygıt ağacı kaynak dosyalarını
derler ve \sphinxcode{\sphinxupquote{dtb}} dosyalarını elde eder. Gömülü sistemlerde bu işlemin yapılması ve her çekirdek
versiyonuyla o versiyonun \sphinxcode{\sphinxupquote{dtb}} dosyasının kullanılması tavsiye edilir.

\end{description}

\sphinxAtStartPar
Pek çok dağıtım o anda yüklü olan çekirdeğe ilişkin konfigürasyon dosyasını \sphinxcode{\sphinxupquote{/boot}} dizini
içerisinde \sphinxcode{\sphinxupquote{config\sphinxhyphen{}\$(uname \sphinxhyphen{}r)}} ismiyle bulundurmaktadır. Örneğin kursun yapılmakta olduğu Mint
dağıtımında \sphinxcode{\sphinxupquote{/boot}} dizininin içeriği şöyledir:

\begin{sphinxVerbatim}[commandchars=\\\{\}]
\PYGZdl{} ls /boot
config\PYGZhy{}6.8.0\PYGZhy{}51\PYGZhy{}generic      initrd.img.old
System.map\PYGZhy{}6.8.0\PYGZhy{}51\PYGZhy{}generic  efi
grub                         vmlinuz
initrd.img                   vmlinuz\PYGZhy{}6.8.0\PYGZhy{}51\PYGZhy{}generic
initrd.img\PYGZhy{}6.8.0\PYGZhy{}51\PYGZhy{}generic
\end{sphinxVerbatim}

\sphinxAtStartPar
Buradaki \sphinxcode{\sphinxupquote{config\sphinxhyphen{}6.8.0\sphinxhyphen{}51\sphinxhyphen{}generic}} dosyası çalışmakta olduğumuz çekirdekte kullanılan
konfigürasyon dosyasıdır. Bu dosya sistem açılırken herhangi bir biçimde kullanılmamaktadır (yani bu
dosyayı silseniz hiçbir sorun oluşmaz). Bu dosya o çekirdeği yeniden derleyecek kişiler için
kolaylık sağlamak amacıyla bulundurulmaktadır.

\sphinxAtStartPar
Eğer çalışılan sistemdeki konfigürasyon dosyasını temel alacaksanız bu dosyayı Linux kaynak
kodlarının bulunduğu kök dizine \sphinxcode{\sphinxupquote{.config}} ismiyle kopyalayabilirsiniz:

\begin{sphinxVerbatim}[commandchars=\\\{\}]
\PYGZdl{}\PYG{+w}{ }cp\PYG{+w}{ }/boot/config\PYGZhy{}\PYG{k}{\PYGZdl{}(}uname\PYG{+w}{ }\PYGZhy{}r\PYG{k}{)}\PYG{+w}{ }.config
\end{sphinxVerbatim}

\sphinxAtStartPar
Eski bir konfigürasyon dosyasını yeni bir çekirdekle kullanmak için ayrıca \sphinxstyleemphasis{“make oldconfig”}
işleminin de yapılması gerekmektedir. Bir sonraki adımda göreceğimiz \sphinxstyleemphasis{“make menuconfig”} işlemi
aynı zamanda \sphinxstyleemphasis{“make oldconfig”} işlemini de kendi içerisinde barındırmaktadır.


\subsection{Adım 4: Konfigürasyon Değişikliklerinin Yapılması}
\label{\detokenize{kernelcompilation:adim-4-konfigurasyon-degisikliklerinin-yapilmasi}}
\sphinxAtStartPar
Elimizde pek çok değerin set edilmiş olduğu \sphinxcode{\sphinxupquote{.config}} isimli bir konfigürasyon dosyası vardır.
Artık bu konfigürasyon dosyasından hareketle yalnızca istediğimiz özellikleri değiştirebiliriz.
Bunun için \sphinxstyleemphasis{“make menuconfig”} komutunu kullanabiliriz:

\begin{sphinxVerbatim}[commandchars=\\\{\}]
\PYGZdl{}\PYG{+w}{ }make\PYG{+w}{ }menuconfig
\end{sphinxVerbatim}

\sphinxAtStartPar
Bu komutla birlikte text ekranda konfigürasyon seçenekleri listelenecektir. Tabii buradaki seçenekler
\sphinxcode{\sphinxupquote{.config}} dosyasındaki içerikten hareketle oluşturulmuş durumdadır. Bunların üzerinde değişiklikler
yaparak \sphinxcode{\sphinxupquote{.config}} dosyasını yeniden save edebiliriz. Aslında \sphinxstyleemphasis{“make menuconfig”} işlemi hiç
\sphinxcode{\sphinxupquote{.config}} dosyası oluşturulmadan doğrudan da yapılabilmektedir. Bu durumda hangi sistemde
çalışılıyorsa o sisteme özgü default konfigürasyon dosyası temel alınmaktadır. Biz en azından
\sphinxcode{\sphinxupquote{General Setup/Local version \sphinxhyphen{} append to kernel release}} seçeneğine \sphinxcode{\sphinxupquote{\sphinxhyphen{}custom}} gibi bir sonek
girmenizi, böylece yeni çekirdeğe \sphinxcode{\sphinxupquote{\sphinxhyphen{}custom}} soneki iliştirmenizi tavsiye ederiz.

\begin{sphinxadmonition}{note}{Not:}
\sphinxAtStartPar
\sphinxstyleemphasis{“make menuconfig”} işlemi zaten \sphinxstyleemphasis{“make oldconfig”} işlemini de kendi içerisinde
barındırmaktadır. Dolayısıyla \sphinxstyleemphasis{“make menuconfig”} yapıyorsanız ayrıca \sphinxstyleemphasis{“make oldconfig”}
işlemini yapmanıza gerek yoktur.
\end{sphinxadmonition}

\sphinxAtStartPar
Peki biz hazır bir \sphinxcode{\sphinxupquote{.config}} dosyasını kaynak kod ağacının kök dizinine kopyaladıktan sonra hiç
\sphinxstyleemphasis{“make menuconfig”} ya da \sphinxstyleemphasis{“make oldconfig”} yazmazsak ne olur? Bu durumda sorun çıkmayabilir. Eğer
kaynak kod çekirdeği yeniyse \sphinxstyleemphasis{make} işlemi sırasında \sphinxstyleemphasis{“make oldconfig”} gibi bir işlem de
yapılmaktadır. Ancak \sphinxcode{\sphinxupquote{.config}} dosyasını kök dizine kopyaladıktan sonra \sphinxstyleemphasis{“make menuconfig”} ya da
\sphinxstyleemphasis{“make oldconfig”} işlemini yapmanızı salık veriyoruz.

\sphinxAtStartPar
\sphinxcode{\sphinxupquote{.config}} dosyası elde edildiğinde çekirdek imzalamasını ortadan kaldırmak için dosyayı açıp
aşağıdaki özellikleri aşağıda gösterildiği gibi değiştirebilirsiniz (bunların bazıları zaten default
durumda aşağıdaki gibi olabilir):

\begin{sphinxVerbatim}[commandchars=\\\{\}]
CONFIG\PYGZus{}SYSTEM\PYGZus{}TRUSTED\PYGZus{}KEYS=\PYGZdq{}\PYGZdq{}
CONFIG\PYGZus{}SYSTEM\PYGZus{}REVOCATION\PYGZus{}KEYS=\PYGZdq{}\PYGZdq{}
CONFIG\PYGZus{}SYSTEM\PYGZus{}TRUSTED\PYGZus{}KEYRING=n
CONFIG\PYGZus{}SECONDARY\PYGZus{}TRUSTED\PYGZus{}KEYRING=n

CONFIG\PYGZus{}MODULE\PYGZus{}SIG=n
CONFIG\PYGZus{}MODULE\PYGZus{}SIG\PYGZus{}ALL=n
CONFIG\PYGZus{}MODULE\PYGZus{}SIG\PYGZus{}KEY=\PYGZdq{}\PYGZdq{}
\end{sphinxVerbatim}

\sphinxAtStartPar
Çekirdek imzalaması konusu daha ileride ele alınacaktır.

\sphinxAtStartPar
Derlenecek çekirdeklere yerel bir versiyon belirteci ve numarası da atanabilmektedir. Bu işlem
\sphinxstyleemphasis{“make menuconfig”} menüsünde \sphinxcode{\sphinxupquote{General Setup/Local version \sphinxhyphen{} append custom release}} seçeneği
kullanılarak ya da \sphinxcode{\sphinxupquote{.config}} dosyasında \sphinxcode{\sphinxupquote{CONFIG\_LOCALVERSION}} satırı düzenlenerek yapılabilir.
Örneğin:

\begin{sphinxVerbatim}[commandchars=\\\{\}]
CONFIG\PYGZus{}LOCALVERSION=\PYGZdq{}\PYGZhy{}custom\PYGZdq{}
\end{sphinxVerbatim}

\sphinxAtStartPar
Bu işlemle artık çekirdek sürümüne \sphinxcode{\sphinxupquote{\sphinxhyphen{}custom}} sonekini eklemiş olduk.


\subsection{Adım 5: Derleme}
\label{\detokenize{kernelcompilation:adim-5-derleme}}
\sphinxAtStartPar
Derleme işlemi için \sphinxstyleemphasis{make} komutu kullanılmaktadır:

\begin{sphinxVerbatim}[commandchars=\\\{\}]
\PYGZdl{}\PYG{+w}{ }make
\end{sphinxVerbatim}

\sphinxAtStartPar
Eğer derleme işleminin birden fazla CPU ya da çekirdek ile yapılmasını istiyorsanız
\sphinxcode{\sphinxupquote{\sphinxhyphen{}j\textless{}cpu\_sayısı\textgreater{}}} seçeneğini komuta dahil edebilirsiniz. Çalışılan sistemdeki CPU sayısının \sphinxstyleemphasis{nproc}
komutuyla elde edildiğini anımsayınız. O halde derleme için \sphinxstyleemphasis{make} komutunu şöyle kullanabiliriz:

\begin{sphinxVerbatim}[commandchars=\\\{\}]
\PYGZdl{}\PYG{+w}{ }make\PYG{+w}{ }\PYGZhy{}j\PYG{k}{\PYGZdl{}(}nproc\PYG{k}{)}
\end{sphinxVerbatim}

\sphinxAtStartPar
Derleme işlemi bittiğinde ürün olarak çekirdek imajı, çekirdek tarafından yüklenecek olan modül
dosyaları ve diğer bazı dosyalar elde edilmiş olur. Derleme işleminden sonra oluşturulan dosyalar
ve yerleri şöyledir (buradaki \sphinxcode{\sphinxupquote{\textless{}çekirdek\_sürümü\textgreater{}}}, \sphinxstyleemphasis{“uname \sphinxhyphen{}r”} ile elde edilecek yazıyı
belirtiyor):
\begin{itemize}
\item {} 
\sphinxAtStartPar
\sphinxstylestrong{Sıkıştırılmış Çekirdek İmajı:} \sphinxcode{\sphinxupquote{arch/\textless{}platform\textgreater{}/boot}} dizininde \sphinxcode{\sphinxupquote{bzImage}} ismiyle
oluşturulmaktadır. Denemeyi yaptığımız Intel makinede dosyanın yol ifadesi
\sphinxcode{\sphinxupquote{arch/x86\_64/boot/bzImage}} biçimindedir. (Ancak buradaki dosya x86\_64 platformu için
\sphinxcode{\sphinxupquote{arch/x86/boot/bzImage}} dosyasına sembolik link de yapılmış olabilir.)

\item {} 
\sphinxAtStartPar
\sphinxstylestrong{Çekirdeğin Sıkıştırılmamış ELF İmajı:} Kaynak kök dizininde \sphinxcode{\sphinxupquote{vmlinux}} ismiyle
oluşturulmaktadır.

\item {} 
\sphinxAtStartPar
\sphinxstylestrong{Çekirdek Modülleri (Aygıt Sürücü Dosyaları):} \sphinxcode{\sphinxupquote{drivers}}, \sphinxcode{\sphinxupquote{fs}} ve \sphinxcode{\sphinxupquote{net}} dizinlerinin
altındaki dizinlerde bulunur. Ancak \sphinxstyleemphasis{“make modules\_install”} ile bunların hepsi belirli bir
dizine çekilebilir.

\item {} 
\sphinxAtStartPar
\sphinxstylestrong{Çekirdek Sembol Tablosu:} Kaynak kök dizininde \sphinxcode{\sphinxupquote{System.map}} ismiyle bulunur. Çekirdek sembol
tablosundan yalnızca çekirdek debug edilirken faydalanılmaktadır. Bu dosya silinse bile sistemin
çalışmasında bir sorun oluşmaz.

\end{itemize}

\begin{sphinxadmonition}{note}{Not:}
\sphinxAtStartPar
Derleme süresini uzatan en önemli etken çekirdek konfigüre edilirken seçilen modül (aygıt sürücü)
sayısıdır. Pek çok dağıtım “belki ileride lazım olur” gerekçesiyle konfigürasyon dosyalarında
çok sayıda modülü dahil etmektedir. Bu nedenle bir dağıtımın konfigürasyon dosyasını kullandığınız
zaman çekirdek derlemesi uzayacaktır. Çekirdek kodlarındaki platforma özgü default konfigürasyon
dosyaları daha minimalist biçimde oluşturulmuş durumdadır.

\sphinxAtStartPar
Tabii çekirdek bütünsel olarak bir kez derlendikten sonra çekirdek kodlarında değişiklik yapıp
çekirdeği yeniden derlemek istediğimizde artık derleme süresi bütünsel derleme kadar uzun
olmayacaktır.
\end{sphinxadmonition}


\subsection{Adım 6: Modüllerin Kurulması}
\label{\detokenize{kernelcompilation:adim-6-modullerin-kurulmasi}}
\sphinxAtStartPar
Farklı dizinlerde oluşturulmuş olan aygıt sürücü dosyalarını (modülleri) belli bir dizine kopyalamak
için \sphinxstyleemphasis{“make modules\_install”} komutu kullanılmaktadır. Bu komut seçeneksiz kullanılırsa default olarak
\sphinxcode{\sphinxupquote{/lib/modules/\textless{}çekirdek\_sürümü\textgreater{}}} dizinine kopyalama yapar:

\begin{sphinxVerbatim}[commandchars=\\\{\}]
\PYGZdl{}\PYG{+w}{ }sudo\PYG{+w}{ }make\PYG{+w}{ }modules\PYGZus{}install
\end{sphinxVerbatim}

\begin{sphinxadmonition}{note}{Not:}
\sphinxAtStartPar
Eskiden \sphinxstyleemphasis{make} işlemi çekirdek modüllerinin derlenmesini sağlamıyordu. Çekirdek modüllerinin
derlenmesi için \sphinxstyleemphasis{“make modules”} işlemi gerekiyordu. Ancak çekirdeğin 2.6 versiyonundan itibaren
\sphinxstyleemphasis{make} işlemi zaten çekirdek modüllerinin de derlenmesini sağlamaktadır.
\end{sphinxadmonition}

\sphinxAtStartPar
Ayrıca \sphinxstyleemphasis{“make modules\_install”} komutunun modül dosyalarını istediğimiz bir dizine kopyalamasını da
sağlayabiliriz. Bunun için \sphinxcode{\sphinxupquote{INSTALL\_MOD\_PATH}} çevre değişkeni kullanılmaktadır. Örneğin:

\begin{sphinxVerbatim}[commandchars=\\\{\}]
\PYGZdl{}\PYG{+w}{ }sudo\PYG{+w}{ }\PYG{n+nv}{INSTALL\PYGZus{}MOD\PYGZus{}PATH}\PYG{o}{=}modules\PYG{+w}{ }make\PYG{+w}{ }modules\PYGZus{}install
\end{sphinxVerbatim}

\sphinxAtStartPar
Burada aygıt sürücü dosyaları \sphinxcode{\sphinxupquote{/lib/modules/\textless{}çekirdek\_sürümü\textgreater{}}} dizinine değil, bulunulan yerdeki
\sphinxcode{\sphinxupquote{modules}} dizinine kopyalanacaktır.


\subsection{Adım 7: Aygıt Ağacı Dosyalarının Derlenmesi}
\label{\detokenize{kernelcompilation:adim-7-aygit-agaci-dosyalarinin-derlenmesi}}
\sphinxAtStartPar
Eğer gömülü sistemler için derleme yapıyorsanız kaynak kod ağacında
\sphinxcode{\sphinxupquote{arch/\textless{}platform\textgreater{}/boot/dts}} dizini içerisindeki aygıt ağacı kaynak dosyalarını da derlemelisiniz.
Aygıt ağacı kaynak dosyalarını derlemek için \sphinxstyleemphasis{“make dtbs”} komutunu kullanabilirsiniz:

\begin{sphinxVerbatim}[commandchars=\\\{\}]
\PYGZdl{}\PYG{+w}{ }make\PYG{+w}{ }dtbs
\end{sphinxVerbatim}

\sphinxAtStartPar
Derlenmiş aygıt ağacı dosyaları \sphinxcode{\sphinxupquote{arch/\textless{}platform\textgreater{}/boot/dts}} dizininde ya da bu dizinin altındaki
ilgili vendor dizininde oluşturulacaktır.

\begin{sphinxadmonition}{note}{Not:}
\sphinxAtStartPar
Aygıt ağaçları gömülü sistemlerde kullanılmaktadır. Intel tabanlı PC’lerde donanım birimlerinin
tespit edilmesi otomatik olarak ACPI protokolü yoluyla yapıldığı için aygıt ağacı dosyaları bu
platformda kullanılmamaktadır. Ancak ARM platformunu kullanan gömülü sistemlerde ya da BBB ve
Raspberry Pi gibi SBC’lerde aygıt ağaçları kullanılmaktadır.
\end{sphinxadmonition}


\subsection{Adım 8: Kurulum}
\label{\detokenize{kernelcompilation:adim-8-kurulum}}
\sphinxAtStartPar
Bizim çekirdek imajını, geçici kök dosya sistemine ilişkin dosyayı ve aygıt ağacı dosyasını \sphinxcode{\sphinxupquote{/boot}}
dizinine kopyalamamız gerekir. Ancak aslında bu işlem de \sphinxstyleemphasis{“make install”} komutuyla otomatik olarak
yapılabilmektedir. \sphinxstyleemphasis{“make install”} komutu bu dosyaları \sphinxcode{\sphinxupquote{/boot}} dizinine kopyalamanın yanı sıra
aynı zamanda GRUB önyükleyici programın konfigürasyon dosyalarında da güncellemeler yapıp yeni
çekirdeğin GRUB menüsü içerisinde görünmesini de sağlamaktadır:

\begin{sphinxVerbatim}[commandchars=\\\{\}]
\PYGZdl{}\PYG{+w}{ }sudo\PYG{+w}{ }make\PYG{+w}{ }install
\end{sphinxVerbatim}


\subsubsection{Geçici Kök Dosya Sistemi (initrd)}
\label{\detokenize{kernelcompilation:gecici-kok-dosya-sistemi-initrd}}
\sphinxAtStartPar
Burada \sphinxstyleemphasis{geçici kök dosya sistemi} (\sphinxstyleemphasis{“initial ramdisk”} ya da \sphinxstyleemphasis{“initrd”}) diye bir terim kullandık.
Geçici kök dosya sistemi diskteki asıl kök dosya sistemi mount edilene kadar geçici bir süre sanki
diskteki dosya sistemiymiş gibi işlev gören bir RAM disk imajıdır. Tipik Linux sistemlerinde önce
geçici kök dosya sistemi mount edilerek temel dosyalara oradan erişilir. Sonra bu geçici dosya sistemi
RAM’den atılıp diskteki gerçek kök dosya sistemi mount edilmektedir.

\sphinxAtStartPar
Geçici kök dosya sistemine gereksinimi basit bir örnekle anlayabiliriz. Diyelim ki diske erişmekte
kullanılan aygıt sürücüsü çekirdeğin içerisine yerleştirilmemiş olsun, yani dışarıda
\sphinxcode{\sphinxupquote{lib/modules/\$(uname \sphinxhyphen{}r)}} dizininde bir dosya biçiminde bulunuyor olsun. Şimdi çekirdeğin diske
erişebilmesi için bu aygıt sürücüye ihtiyacı olacaktır. Ancak aygıt sürücü de disktedir. İşte böyle
bir durumda bu aygıt sürücüleri de barındıran bir RAM disk dosya sistemi oluşturulmakta ve önyükleyici
tarafından (örneğin GRUB önyükleyicisi) bu dosya da RAM’e yüklenmektedir. Böylece çekirdek artık
diske erişebilir hale gelir.

\begin{sphinxadmonition}{note}{Not:}
\sphinxAtStartPar
Eğer çekirdeğin gereksinim duyacağı bütün aygıt sürücü dosyaları konfigürasyon aşamasında
çekirdeğin içerisine gömülmüşse geçici kök dosya sistemi oluşturmadan da sistem boot edilebilir.
Ancak bu sırada çözülmesi gereken problemlerle de karşılaşılabilmektedir. Özellikle masaüstü
sistemlerinde geçici kök dosya sistemi olmadan sistemi boot etmek oldukça zahmetlidir.
\end{sphinxadmonition}

\sphinxAtStartPar
Geçici kök dosya sistemi aynı zamanda işletim sistemini \sphinxstyleemphasis{güvenli kipte (safe mode)} açmak için ve
sistem güncellemelerinde de kullanılmaktadır.

\sphinxAtStartPar
Debian türevi dağıtımlarda geçici kök dosya sistemini oluşturmak için \sphinxstyleemphasis{“update\sphinxhyphen{}initramfs”} komutu
kullanılmaktadır. Bu komut da \sphinxstyleemphasis{mkinitramfs} komutunu çalıştırmaktadır. Geçici kök dosya sistemi
oluşturulurken bu komutlar işlemlerini kabaca aşağıdaki adımlarla gerçekleştirmektedir:

\sphinxincludegraphics[]{graphviz-43a178dd74d17c866e26ea957698c636aec067ab.pdf}

\sphinxAtStartPar
\sphinxstyleemphasis{“make install”} komutu masaüstü sistemlerde sırasıyla şunları yapmaktadır:
\begin{itemize}
\item {} 
\sphinxAtStartPar
\sphinxcode{\sphinxupquote{arch/\textless{}platform\textgreater{}/boot}} dizinindeki \sphinxcode{\sphinxupquote{bzImage}} çekirdek imajı alınarak hedef \sphinxcode{\sphinxupquote{/boot}} dizinine
\sphinxcode{\sphinxupquote{vmlinuz\sphinxhyphen{}\textless{}çekirdek\_sürümü\textgreater{}}} ismiyle kopyalanır.

\item {} 
\sphinxAtStartPar
\sphinxcode{\sphinxupquote{System.map}} dosyası hedef \sphinxcode{\sphinxupquote{/boot}} dizinine \sphinxcode{\sphinxupquote{System.map\sphinxhyphen{}\textless{}çekirdek\_sürümü\textgreater{}}} ismiyle kopyalanır.

\item {} 
\sphinxAtStartPar
\sphinxcode{\sphinxupquote{.config}} dosyası \sphinxcode{\sphinxupquote{/boot}} dizinine \sphinxcode{\sphinxupquote{config\sphinxhyphen{}\textless{}çekirdek\_sürümü\textgreater{}}} ismiyle kopyalanır.

\item {} 
\sphinxAtStartPar
Geçici kök dosya sistemine ilişkin dosya oluşturulur ve hedef \sphinxcode{\sphinxupquote{/boot}} dizinine
\sphinxcode{\sphinxupquote{initrd.img\sphinxhyphen{}\textless{}çekirdek\_sürümü\textgreater{}}} ismiyle kopyalanır.

\item {} 
\sphinxAtStartPar
Eğer GRUB önyükleyicisi kullanılıyorsa GRUB konfigürasyonu güncellenir ve GRUB menüsüne yeni
girişler eklenir.

\end{itemize}

\sphinxAtStartPar
Böylece sistemin otomatik olarak yeni çekirdekle açılması sağlanır.

\sphinxAtStartPar
\sphinxstyleemphasis{“make install”} komutu uygulandığında eğer çekirdeğin versiyon bilgisi aynı ise \sphinxcode{\sphinxupquote{/boot}}
dizinindeki bir önceki kurulumun dosyaları \sphinxcode{\sphinxupquote{.old}} uzantısıyla saklanmaktadır. Böylece son
\sphinxstyleemphasis{“make install”} komutundan önceki kuruluma manuel olarak geri dönebilirsiniz. Örneğin yeniden
aynı çekirdek versiyonunu \sphinxstyleemphasis{“make install”} yaptığımızda \sphinxcode{\sphinxupquote{/boot}} dizininin içeriği şöyle
olacaktır:

\begin{sphinxVerbatim}[commandchars=\\\{\}]
kaan@kaan\PYGZhy{}Huawei:\PYGZti{}/Study/LinuxKernel/linux\PYGZhy{}6.9.2\PYGZdl{} ls \PYGZhy{}l /boot
total 655480
\PYGZhy{}rw\PYGZhy{}r\PYGZhy{}\PYGZhy{}r\PYGZhy{}\PYGZhy{} 1 root root    287375 Haz  7  2024 config\PYGZhy{}6.8.0\PYGZhy{}38\PYGZhy{}generic
\PYGZhy{}rw\PYGZhy{}r\PYGZhy{}\PYGZhy{}r\PYGZhy{}\PYGZhy{} 1 root root    287766 Ağu 15 11:57 config\PYGZhy{}6.9.2\PYGZhy{}custom
\PYGZhy{}rw\PYGZhy{}r\PYGZhy{}\PYGZhy{}r\PYGZhy{}\PYGZhy{} 1 root root    287766 Ağu 15 11:55 config\PYGZhy{}6.9.2\PYGZhy{}custom.old
drwx\PYGZhy{}\PYGZhy{}\PYGZhy{}\PYGZhy{}\PYGZhy{}\PYGZhy{} 3 root root      4096 Oca  1  1970 efi
drwxr\PYGZhy{}xr\PYGZhy{}x 6 root root      4096 Ağu 15 11:57 grub
lrwxrwxrwx 1 root root        23 Ağu 10 11:20 initrd.img \PYGZhy{}\PYGZgt{} initrd.img\PYGZhy{}6.9.2\PYGZhy{}custom
\PYGZhy{}rw\PYGZhy{}r\PYGZhy{}\PYGZhy{}r\PYGZhy{}\PYGZhy{} 1 root root  73052139 Kas 19  2024 initrd.img\PYGZhy{}6.8.0\PYGZhy{}38\PYGZhy{}generic
\PYGZhy{}rw\PYGZhy{}r\PYGZhy{}\PYGZhy{}r\PYGZhy{}\PYGZhy{} 1 root root 526888758 Ağu 15 11:57 initrd.img\PYGZhy{}6.9.2\PYGZhy{}custom
\PYGZhy{}rw\PYGZhy{}\PYGZhy{}\PYGZhy{}\PYGZhy{}\PYGZhy{}\PYGZhy{}\PYGZhy{} 1 root root   9055262 Haz  7  2024 System.map\PYGZhy{}6.8.0\PYGZhy{}38\PYGZhy{}generic
\PYGZhy{}rw\PYGZhy{}r\PYGZhy{}\PYGZhy{}r\PYGZhy{}\PYGZhy{} 1 root root   8365115 Ağu 15 11:57 System.map\PYGZhy{}6.9.2\PYGZhy{}custom
\PYGZhy{}rw\PYGZhy{}r\PYGZhy{}\PYGZhy{}r\PYGZhy{}\PYGZhy{} 1 root root   8365115 Ağu 15 11:55 System.map\PYGZhy{}6.9.2\PYGZhy{}custom.old
lrwxrwxrwx 1 root root        20 Ağu 15 11:57 vmlinuz \PYGZhy{}\PYGZgt{} vmlinuz\PYGZhy{}6.9.2\PYGZhy{}custom
\PYGZhy{}rw\PYGZhy{}r\PYGZhy{}\PYGZhy{}r\PYGZhy{}\PYGZhy{} 1 root root  14944648 Haz  7  2024 vmlinuz\PYGZhy{}6.8.0\PYGZhy{}38\PYGZhy{}generic
\PYGZhy{}rw\PYGZhy{}r\PYGZhy{}\PYGZhy{}r\PYGZhy{}\PYGZhy{} 1 root root  14819840 Ağu 15 11:57 vmlinuz\PYGZhy{}6.9.2\PYGZhy{}custom
\PYGZhy{}rw\PYGZhy{}r\PYGZhy{}\PYGZhy{}r\PYGZhy{}\PYGZhy{} 1 root root  14819840 Ağu 15 11:55 vmlinuz\PYGZhy{}6.9.2\PYGZhy{}custom.old
lrwxrwxrwx 1 root root        24 Ağu 15 11:57 vmlinuz.old \PYGZhy{}\PYGZgt{} vmlinuz\PYGZhy{}6.9.2\PYGZhy{}custom.old
\end{sphinxVerbatim}

\begin{sphinxadmonition}{note}{Not:}
\sphinxAtStartPar
\sphinxstyleemphasis{“make install”} komutu masaüstü sistemler için kullanılmaktadır. Gömülü sistemlerde tüm
dosyaların toparlanıp hedef sisteme aktarılması gerekir. Gömülü sistemlerde \sphinxstyleemphasis{“make install”}
tarafından yapılan işlemleri siz manuel bir biçimde yapabilirsiniz. Tabii gömülü sistemlerde
GRUB önyükleyicisi yerine ağırlıklı olarak \sphinxstyleemphasis{U\sphinxhyphen{}Boot} denilen önyükleyici kullanılmaktadır.
Dolayısıyla gömülü sistemlerde sistemin yeni çekirdekle açılmasını sağlamak için \sphinxstyleemphasis{U\sphinxhyphen{}Boot}
önyükleyicisini de ayarlamanız gerekecektir.
\end{sphinxadmonition}


\section{make modules\_install Komutu ile İlgili Ayrıntılar}
\label{\detokenize{kernelcompilation:make-modules-install-komutu-ile-ilgili-ayrintilar}}
\sphinxAtStartPar
\sphinxstyleemphasis{“make modules\_install”} komutu yalnızca modül dosyalarını hedef dizine kopyalamakla kalmaz; aynı
zamanda bazı dosyaları da oluşturup onları da hedef dizine kopyalar. Hedef dizinin default
\sphinxcode{\sphinxupquote{/lib/modules/\textless{}çekirdek\_sürümü\textgreater{}}} dizini olduğu varsayımıyla komut sırasıyla şunları yapmaktadır:
\begin{itemize}
\item {} 
\sphinxAtStartPar
Modül dosyalarını \sphinxcode{\sphinxupquote{/lib/modules/\textless{}çekirdek\_sürümü\textgreater{}}} dizinine kopyalar.

\item {} 
\sphinxAtStartPar
\sphinxcode{\sphinxupquote{modules.dep}} isimli dosyayı oluşturur ve bunu \sphinxcode{\sphinxupquote{/lib/modules/\textless{}çekirdek\_sürümü\textgreater{}}} dizinine
kopyalar.

\item {} 
\sphinxAtStartPar
\sphinxcode{\sphinxupquote{modules.alias}} isimli dosyayı oluşturur ve bunu \sphinxcode{\sphinxupquote{/lib/modules/\textless{}çekirdek\_sürümü\textgreater{}}} dizinine
kopyalar.

\item {} 
\sphinxAtStartPar
\sphinxcode{\sphinxupquote{modules.order}} isimli dosyayı oluşturur ve \sphinxcode{\sphinxupquote{/lib/modules/\textless{}çekirdek\_sürümü\textgreater{}}} dizinine kopyalar.

\item {} 
\sphinxAtStartPar
\sphinxcode{\sphinxupquote{modules.builtin}} isimli dosyayı \sphinxcode{\sphinxupquote{/lib/modules/\textless{}çekirdek\_sürümü\textgreater{}}} dizinine kopyalar.

\end{itemize}

\sphinxAtStartPar
Aslında burada oluşturulan dosyaların bazıları mutlak anlamda bulundurulmak zorunda değildir. Ancak
sistemin öngörüldüğü gibi işlev göstermesi için bu dosyaların ilgili dizinde bulundurulması uygundur.


\subsection{modules.dep}
\label{\detokenize{kernelcompilation:modules-dep}}
\sphinxAtStartPar
Bir aygıt sürücü başka aygıt sürücüleri de kullanıyor olabilir. Yani bir aygıt sürücünün
çalışabilmesi için başka aygıt sürücülerin de yüklü olması gerekebilmektedir. İşte \sphinxcode{\sphinxupquote{modules.dep}}
dosyası bir aygıt sürücünün yüklenmesi için başka hangi aygıt sürücülerin yüklenmesi gerektiği
bilgisini tutmaktadır. \sphinxcode{\sphinxupquote{modules.dep}} bir text dosyadır. Satırların içeriği şöyledir:

\begin{sphinxVerbatim}[commandchars=\\\{\}]
\PYGZlt{}modül\PYGZus{}yolu\PYGZgt{}: \PYGZlt{}bağımlılık1\PYGZgt{} \PYGZlt{}bağımlılık2\PYGZgt{} ...
\end{sphinxVerbatim}

\sphinxAtStartPar
Dosyanın içeriğine örnek verebiliriz:

\begin{sphinxVerbatim}[commandchars=\\\{\}]
kernel/arch/x86/crypto/nhpoly1305\PYGZhy{}sse2.ko.zst: kernel/crypto/nhpoly1305.ko.zst kernel/lib/crypto/libpoly1305.ko.zst
kernel/arch/x86/crypto/nhpoly1305\PYGZhy{}avx2.ko.zst: kernel/crypto/nhpoly1305.ko.zst kernel/lib/crypto/libpoly1305.ko.zst
kernel/arch/x86/crypto/curve25519\PYGZhy{}x86\PYGZus{}64.ko.zst: kernel/lib/crypto/libcurve25519\PYGZhy{}generic.ko.zst
\end{sphinxVerbatim}

\sphinxAtStartPar
Eğer bu \sphinxcode{\sphinxupquote{modules.dep}} dosyası olmazsa \sphinxstyleemphasis{modprobe} komutu çalışmaz ve çekirdek modülleri yüklenirken
eksik yükleme yapılabilir. Dolayısıyla sistem düzgün bir biçimde çalışmayabilir. Eğer bu dosya
elimizde yoksa ya da bir biçimde silinmişse \sphinxstyleemphasis{“depmod \sphinxhyphen{}a”} komutu ile yeniden oluşturulabilir:

\begin{sphinxVerbatim}[commandchars=\\\{\}]
\PYGZdl{}\PYG{+w}{ }sudo\PYG{+w}{ }depmod\PYG{+w}{ }\PYGZhy{}a
\end{sphinxVerbatim}

\sphinxAtStartPar
Siz yüklü olan başka bir çekirdek sürümü için \sphinxcode{\sphinxupquote{modules.dep}} dosyasını oluşturmak istiyorsanız
çekirdek sürümünü de argüman olarak vermelisiniz:

\begin{sphinxVerbatim}[commandchars=\\\{\}]
\PYGZdl{}\PYG{+w}{ }sudo\PYG{+w}{ }depmod\PYG{+w}{ }\PYGZhy{}a\PYG{+w}{ }\PYGZlt{}çekirdek\PYGZus{}sürümü\PYGZgt{}
\end{sphinxVerbatim}

\begin{sphinxadmonition}{note}{Not:}
\sphinxAtStartPar
\sphinxstyleemphasis{depmod} komutunun çalışabilmesi için \sphinxcode{\sphinxupquote{/lib/modules/\textless{}çekirdek\_sürümü\textgreater{}}} dizininde modül
dosyalarının bulunuyor olması gerekir. Çünkü bu komut bu dizindeki modül dosyalarını tek tek bulup
ELF formatının ilgili bölümlerine bakarak modülün hangi modülleri kullandığını tespit ederek
\sphinxcode{\sphinxupquote{modules.dep}} dosyasını oluşturmaktadır.
\end{sphinxadmonition}


\subsection{modules.alias}
\label{\detokenize{kernelcompilation:modules-alias}}
\sphinxAtStartPar
\sphinxcode{\sphinxupquote{modules.alias}} dosyası belli bir isim ya da id ile aygıt sürücü dosyasını eşleştiren bir text
dosyadır. Bu dosyanın bulunmaması bazı durumlarda sorunlara yol açmayabilir. Ancak örneğin USB porta
bir aygıt takıldığında bu aygıta ilişkin aygıt sürücünün hangisi olduğu bilgisi bu dosyada
tutulmaktadır. Bu durumda bu dosyanın olmayışı aygıt sürücünün yüklenememesine neden olabilir.
Dosyanın içeriği aşağıdaki formata uygun satırlardan oluşmaktadır:

\begin{sphinxVerbatim}[commandchars=\\\{\}]
alias \PYGZlt{}tanımlayıcı\PYGZgt{} \PYGZlt{}modül\PYGZus{}adı\PYGZgt{}
\end{sphinxVerbatim}

\sphinxAtStartPar
Örnek bir içerik:

\begin{sphinxVerbatim}[commandchars=\\\{\}]
alias usb:v05ACp*d*dc*dsc*dp*ic*isc*ip*in* apple\PYGZus{}mfi\PYGZus{}fastcharge
alias usb:v8086p0B63d*dc*dsc*dp*ic*isc*ip*in* usb\PYGZus{}ljca
alias usb:v0681p0010d*dc*dsc*dp*ic*isc*ip*in* idmouse
alias usb:v0681p0005d*dc*dsc*dp*ic*isc*ip*in* idmouse
alias usb:v07C0p1506d*dc*dsc*dp*ic*isc*ip*in* iowarrior
alias usb:v07C0p1505d*dc*dsc*dp*ic*isc*ip*in* iowarrior
\end{sphinxVerbatim}

\sphinxAtStartPar
Bu dosya silinirse yine \sphinxstyleemphasis{“depmod \sphinxhyphen{}a”} komutu ile oluşturulabilir.


\subsection{modules.order}
\label{\detokenize{kernelcompilation:modules-order}}
\sphinxAtStartPar
\sphinxcode{\sphinxupquote{modules.order}} dosyası aygıt sürücü dosyalarının yüklenme sırasını barındıran bir text dosyadır.
Bu dosyanın her satırında bir çekirdek aygıt sürücüsünün dosya yol ifadesi bulunur. Daha önce
yazılmış aygıt sürücüler daha sonra yazılanlardan daha önce yüklenir. Bu dosyanın olmaması genellikle
bir soruna yol açmaz. Ancak modüllerin belli sırada yüklenmemesi bazı durumlarda bozukluklara da
neden olabilmektedir. Bu dosyanın silinmesi durumunda yine \sphinxstyleemphasis{“depmod \sphinxhyphen{}a”} komutuyla oluşturulabilmektedir.


\section{Manuel Kurulum}
\label{\detokenize{kernelcompilation:manuel-kurulum}}
\sphinxAtStartPar
Biz \sphinxstyleemphasis{“make install”} komutu ile yapılan işlemleri manuel olarak da yapabiliriz. Aslında çekirdek
imajı ve geçici kök dosya sistemi dosyaları default yerlerin dışında başka yerlerde de
bulundurulabilir; önyükleyiciye bu konuda bilgi verilebilir. Ancak yukarıdaki dosyaların hedef
sistemde bulundurulduğu default yerler şöyledir:


\begin{savenotes}\sphinxattablestart
\sphinxthistablewithglobalstyle
\centering
\begin{tabular}[t]{\X{40}{100}\X{60}{100}}
\sphinxtoprule
\begin{varwidth}[t]{\sphinxcolwidth{1}{2}}
\sphinxstyletheadfamily \sphinxAtStartPar
Dosya
\sphinxbeforeendvarwidth
\end{varwidth}%
&\begin{varwidth}[t]{\sphinxcolwidth{1}{2}}
\sphinxstyletheadfamily \sphinxAtStartPar
Hedef Dizin
\sphinxbeforeendvarwidth
\end{varwidth}%
\\
\sphinxmidrule
\sphinxtableatstartofbodyhook\begin{varwidth}[t]{\sphinxcolwidth{1}{2}}
\sphinxAtStartPar
Çekirdek İmajı
\sphinxbeforeendvarwidth
\end{varwidth}%
&\begin{varwidth}[t]{\sphinxcolwidth{1}{2}}
\sphinxAtStartPar
\sphinxcode{\sphinxupquote{/boot}}
\sphinxbeforeendvarwidth
\end{varwidth}%
\\
\sphinxhline\begin{varwidth}[t]{\sphinxcolwidth{1}{2}}
\sphinxAtStartPar
Çekirdek Sembol Tablosu
\sphinxbeforeendvarwidth
\end{varwidth}%
&\begin{varwidth}[t]{\sphinxcolwidth{1}{2}}
\sphinxAtStartPar
\sphinxcode{\sphinxupquote{/boot}}
\sphinxbeforeendvarwidth
\end{varwidth}%
\\
\sphinxhline\begin{varwidth}[t]{\sphinxcolwidth{1}{2}}
\sphinxAtStartPar
Modül Dosyaları
\sphinxbeforeendvarwidth
\end{varwidth}%
&\begin{varwidth}[t]{\sphinxcolwidth{1}{2}}
\sphinxAtStartPar
\sphinxcode{\sphinxupquote{/lib/modules/\textless{}çekirdek\_sürümü\textgreater{}/kernel}}
\sphinxbeforeendvarwidth
\end{varwidth}%
\\
\sphinxhline\begin{varwidth}[t]{\sphinxcolwidth{1}{2}}
\sphinxAtStartPar
Geçici Kök Dosya Sistemi Dosyası
\sphinxbeforeendvarwidth
\end{varwidth}%
&\begin{varwidth}[t]{\sphinxcolwidth{1}{2}}
\sphinxAtStartPar
\sphinxcode{\sphinxupquote{/boot}}
\sphinxbeforeendvarwidth
\end{varwidth}%
\\
\sphinxbottomrule
\end{tabular}
\sphinxtableafterendhook\par
\sphinxattableend\end{savenotes}

\sphinxAtStartPar
İsteğe bağlı olarak aşağıdaki dosyalar da hedef sisteme konuşlandırılabilir:


\begin{savenotes}\sphinxattablestart
\sphinxthistablewithglobalstyle
\centering
\begin{tabular}[t]{\X{40}{100}\X{60}{100}}
\sphinxtoprule
\begin{varwidth}[t]{\sphinxcolwidth{1}{2}}
\sphinxstyletheadfamily \sphinxAtStartPar
Dosya
\sphinxbeforeendvarwidth
\end{varwidth}%
&\begin{varwidth}[t]{\sphinxcolwidth{1}{2}}
\sphinxstyletheadfamily \sphinxAtStartPar
Hedef Dizin
\sphinxbeforeendvarwidth
\end{varwidth}%
\\
\sphinxmidrule
\sphinxtableatstartofbodyhook\begin{varwidth}[t]{\sphinxcolwidth{1}{2}}
\sphinxAtStartPar
Konfigürasyon Dosyası
\sphinxbeforeendvarwidth
\end{varwidth}%
&\begin{varwidth}[t]{\sphinxcolwidth{1}{2}}
\sphinxAtStartPar
\sphinxcode{\sphinxupquote{/boot}}
\sphinxbeforeendvarwidth
\end{varwidth}%
\\
\sphinxhline\begin{varwidth}[t]{\sphinxcolwidth{1}{2}}
\sphinxAtStartPar
Modüllere İlişkin Bazı Dosyalar
\sphinxbeforeendvarwidth
\end{varwidth}%
&\begin{varwidth}[t]{\sphinxcolwidth{1}{2}}
\sphinxAtStartPar
\sphinxcode{\sphinxupquote{/lib/modules/\textless{}çekirdek\_sürümü\textgreater{}}}
\sphinxbeforeendvarwidth
\end{varwidth}%
\\
\sphinxbottomrule
\end{tabular}
\sphinxtableafterendhook\par
\sphinxattableend\end{savenotes}

\sphinxAtStartPar
Peki yukarıda belirttiğimiz dosyalar hedef sistemdeki ilgili dizinlere hangi isimlerle
kopyalanmalıdır? Tipik isimlendirme şöyle olmalıdır (buradaki \sphinxcode{\sphinxupquote{\textless{}çekirdek\_sürümü\textgreater{}}}, \sphinxstyleemphasis{“uname \sphinxhyphen{}r”}
komutuyla elde edilecek olan yazıdır):
\begin{itemize}
\item {} 
\sphinxAtStartPar
\sphinxstylestrong{Çekirdek İmajı:} \sphinxcode{\sphinxupquote{/boot/vmlinuz\sphinxhyphen{}\textless{}çekirdek\_sürümü\textgreater{}}} — örneğin \sphinxcode{\sphinxupquote{vmlinuz\sphinxhyphen{}6.9.2\sphinxhyphen{}custom}}

\item {} 
\sphinxAtStartPar
\sphinxstylestrong{Çekirdek Sembol Tablosu:} \sphinxcode{\sphinxupquote{/boot/System.map\sphinxhyphen{}\textless{}çekirdek\_sürümü\textgreater{}}} — örneğin
\sphinxcode{\sphinxupquote{System.map\sphinxhyphen{}6.9.2\sphinxhyphen{}custom}}

\item {} 
\sphinxAtStartPar
\sphinxstylestrong{Modüllere İlişkin Dosyalar:} \sphinxcode{\sphinxupquote{/lib/modules/\textless{}çekirdek\_sürümü\textgreater{}}} dizininin içerisine

\item {} 
\sphinxAtStartPar
\sphinxstylestrong{Konfigürasyon Dosyası:} \sphinxcode{\sphinxupquote{/boot/config\sphinxhyphen{}\textless{}çekirdek\_sürümü\textgreater{}}} — örneğin \sphinxcode{\sphinxupquote{config\sphinxhyphen{}6.9.2\sphinxhyphen{}custom}}

\item {} 
\sphinxAtStartPar
\sphinxstylestrong{Geçici Kök Dosya Sistemi:} \sphinxcode{\sphinxupquote{/boot/initrd.img\sphinxhyphen{}\textless{}çekirdek\_sürümü\textgreater{}}} — örneğin
\sphinxcode{\sphinxupquote{initrd.img\sphinxhyphen{}6.9.2\sphinxhyphen{}custom}}; \sphinxstyleemphasis{“update\sphinxhyphen{}initramfs”} programıyla oluşturulabilir.

\end{itemize}

\sphinxAtStartPar
Ayrıca bazı dağıtımlarda \sphinxcode{\sphinxupquote{/boot}} dizini içerisindeki \sphinxcode{\sphinxupquote{vmlinuz}} dosyası default olan
\sphinxcode{\sphinxupquote{vmlinuz\sphinxhyphen{}\textless{}çekirdek\_sürümü\textgreater{}}} dosyasına, \sphinxcode{\sphinxupquote{initrd.img}} dosyası da
\sphinxcode{\sphinxupquote{/boot/initrd.img\sphinxhyphen{}\textless{}çekirdek\_sürümü\textgreater{}}} dosyasına sembolik link yapılmış durumda olabilir. Ancak bu
sembolik bağlantıları GRUB kullanmamaktadır. Aşağıda Intel sistemindeki 6.8.0 çekirdeğinin yüklü
olduğu \sphinxcode{\sphinxupquote{/boot}} dizininin default içeriğini görüyorsunuz:

\begin{sphinxVerbatim}[commandchars=\\\{\}]
\PYGZdl{} ls \PYGZhy{}l /boot
\PYGZhy{}rw\PYGZhy{}r\PYGZhy{}\PYGZhy{}r\PYGZhy{}\PYGZhy{} 1 root root    287375 Haz  7  2024 config\PYGZhy{}6.8.0\PYGZhy{}38\PYGZhy{}generic
drwx\PYGZhy{}\PYGZhy{}\PYGZhy{}\PYGZhy{}\PYGZhy{}\PYGZhy{} 3 root root      4096 Oca  1  1970 efi
drwxr\PYGZhy{}xr\PYGZhy{}x 6 root root      4096 Ağu 15 11:57 grub
lrwxrwxrwx 1 root root        23 Ağu 10 11:20 initrd.img \PYGZhy{}\PYGZgt{} initrd.img\PYGZhy{}6.8.0\PYGZhy{}38\PYGZhy{}generic
\PYGZhy{}rw\PYGZhy{}r\PYGZhy{}\PYGZhy{}r\PYGZhy{}\PYGZhy{} 1 root root  73052139 Kas 19  2024 initrd.img\PYGZhy{}6.8.0\PYGZhy{}38\PYGZhy{}generic
\PYGZhy{}rw\PYGZhy{}\PYGZhy{}\PYGZhy{}\PYGZhy{}\PYGZhy{}\PYGZhy{}\PYGZhy{} 1 root root   9055262 Haz  7  2024 System.map\PYGZhy{}6.8.0\PYGZhy{}38\PYGZhy{}generic
lrwxrwxrwx 1 root root        20 Ağu 15 11:57 vmlinuz \PYGZhy{}\PYGZgt{} vmlinuz\PYGZhy{}6.8.0\PYGZhy{}38\PYGZhy{}generic
\PYGZhy{}rw\PYGZhy{}r\PYGZhy{}\PYGZhy{}r\PYGZhy{}\PYGZhy{} 1 root root  14944648 Haz  7  2024 vmlinuz\PYGZhy{}6.8.0\PYGZhy{}38\PYGZhy{}generic
\end{sphinxVerbatim}

\sphinxAtStartPar
Geçici kök dosya sisteminin içerisindeki dosyalar \sphinxstyleemphasis{cpio} denilen bir arşiv formatıyla
arşivlenmektedir. \sphinxstyleemphasis{cpio} arşivi tıpkı \sphinxstyleemphasis{tar} arşivinde olduğu gibi yalnızca dosyaların uç uca
eklenmesiyle oluşturulmaktadır. İsterseniz geçici kök dosya sistemine ilişkin \sphinxcode{\sphinxupquote{initrd\sphinxhyphen{}xxx}}
dosyalarını açabilirsiniz. Ancak bu dosyaların içeriği dağıtımların versiyonlarına göre
değişebilmektedir. Yeni dağıtımlarda bu \sphinxcode{\sphinxupquote{initrd\sphinxhyphen{}xxx}} arşiv dosyası birkaç bölümden oluşmaktadır.
Eğer biz bu dosyayı \sphinxstyleemphasis{cpio} programıyla açmaya çalışırsak bu program yalnızca ilk bölümü açacaktır.
Tüm bölümleri açmak için \sphinxstyleemphasis{unmkinitramfs} programından faydalanabilirsiniz:

\begin{sphinxVerbatim}[commandchars=\\\{\}]
unmkinitramfs\PYG{+w}{ }/boot/initrd.img\PYGZhy{}6.9.2\PYGZhy{}custom\PYG{+w}{ }initrd
\end{sphinxVerbatim}

\sphinxAtStartPar
Bu komutla geçici kök dosya sistemine ilişkin arşiv dosyası \sphinxcode{\sphinxupquote{initrd}} isimli dizinin altına
açılacaktır.


\subsection{Çekirdek Sürüm Yazısının Belirlenmesi}
\label{\detokenize{kernelcompilation:cekirdek-surum-yazisinin-belirlenmesi}}
\sphinxAtStartPar
Derleme sonucunda elde ettiğimiz dosyaları manuel isimlendirirken çekirdek sürüm yazısını nasıl
bileceğiz? Bunun için \sphinxstyleemphasis{“uname \sphinxhyphen{}r”} komutunu kullanamayız. Çünkü bu komut bize o anda çalışmakta
olan çekirdeğin sürüm yazısını verir. Sürüm yazısı çekirdek imajının içerisine de yazılmaktadır ve
bizim bazı dosyalara verdiğimiz isimlerin çekirdek içerisindeki bu yazıyla uyumlu olması gerekir.

\sphinxAtStartPar
Default olarak \sphinxcode{\sphinxupquote{kernel.org}} sitesinden indirilen kaynak kodlar derlendiğinde çekirdek sürümü
\sphinxcode{\sphinxupquote{6.9.2}} gibi üç haneli bir sayılardan oluşmaktadır. Yani yazının sonunda \sphinxcode{\sphinxupquote{\sphinxhyphen{}generic}} ya da
\sphinxcode{\sphinxupquote{\sphinxhyphen{}custom}} gibi sonekler yoktur. Tabii çekirdeği derlemeden önce \sphinxcode{\sphinxupquote{.config}} dosyasında
\sphinxcode{\sphinxupquote{CONFIG\_LOCALVERSION}} özelliğine bu sürüm numarasından sonra eklenecek bilgiyi girebilirsiniz.

\sphinxAtStartPar
Sürüm yazısını bulmak için sıkıştırılmamış çekirdek dosyası içerisindeki (kaynak kök dizindeki
\sphinxcode{\sphinxupquote{vmlinux}} dosyası) string tablosunda \sphinxcode{\sphinxupquote{Linux version}} yazısını aramak yeterlidir:

\begin{sphinxVerbatim}[commandchars=\\\{\}]
\PYGZdl{}\PYG{+w}{ }strings\PYG{+w}{ }vmlinux\PYG{+w}{ }\PYG{p}{|}\PYG{+w}{ }grep\PYG{+w}{ }\PYG{l+s+s2}{\PYGZdq{}Linux version\PYGZdq{}}
Linux\PYG{+w}{ }version\PYG{+w}{ }\PYG{l+m}{6}.9.2\PYGZhy{}custom\PYG{+w}{ }\PYG{o}{(}kaan@kaan\PYGZhy{}virtual\PYGZhy{}machine\PYG{o}{)}\PYG{+w}{ }\PYG{o}{(}gcc\PYG{+w}{ }\PYG{o}{(}Ubuntu\PYG{+w}{ }\PYG{l+m}{11}.4.0\PYGZhy{}1ubuntu1\PYGZti{}22.04\PYG{o}{)}\PYG{+w}{ }\PYG{l+m}{11}.4.0,
GNU\PYG{+w}{ }ld\PYG{+w}{ }\PYG{o}{(}GNU\PYG{+w}{ }Binutils\PYG{+w}{ }\PYG{k}{for}\PYG{+w}{ }Ubuntu\PYG{o}{)}\PYG{+w}{ }\PYG{l+m}{2}.38\PYG{o}{)}\PYG{+w}{ }\PYG{c+c1}{\PYGZsh{} SMP PREEMPT\PYGZus{}DYNAMIC}
Linux\PYG{+w}{ }version\PYG{+w}{ }\PYG{l+m}{6}.9.2\PYGZhy{}custom\PYG{+w}{ }\PYG{o}{(}kaan@kaan\PYGZhy{}virtual\PYGZhy{}machine\PYG{o}{)}\PYG{+w}{ }\PYG{o}{(}gcc\PYG{+w}{ }\PYG{o}{(}Ubuntu\PYG{+w}{ }\PYG{l+m}{11}.4.0\PYGZhy{}1ubuntu1\PYGZti{}22.04\PYG{o}{)}\PYG{+w}{ }\PYG{l+m}{11}.4.0,
GNU\PYG{+w}{ }ld\PYG{+w}{ }\PYG{o}{(}GNU\PYG{+w}{ }Binutils\PYG{+w}{ }\PYG{k}{for}\PYG{+w}{ }Ubuntu\PYG{o}{)}\PYG{+w}{ }\PYG{l+m}{2}.38\PYG{o}{)}\PYG{+w}{ }\PYG{c+c1}{\PYGZsh{}2 SMP PREEMPT\PYGZus{}DYNAMIC Thu Dec 5 17:55:14 +03 2024}
\end{sphinxVerbatim}

\sphinxAtStartPar
Buradan sürüm yazısının \sphinxcode{\sphinxupquote{6.9.2\sphinxhyphen{}custom}} olduğu görülmektedir. O halde derleme sonucunda elde
ettiğimiz dosyaları manuel biçimde kopyalarken sürüm bilgisi olarak \sphinxcode{\sphinxupquote{6.9.2\sphinxhyphen{}custom}} yazısını
kullanmamız gerekir. Çekirdek imajının \sphinxcode{\sphinxupquote{/boot}} dizinine manuel kopyalanması işlemi şöyle
yapılabilir (kaynak kök dizinde bulunduğumuzu varsayıyoruz):

\begin{sphinxVerbatim}[commandchars=\\\{\}]
\PYGZdl{}\PYG{+w}{ }sudo\PYG{+w}{ }cp\PYG{+w}{ }arch/x86\PYGZus{}64/boot/bzImage\PYG{+w}{ }/boot/vmlinuz\PYGZhy{}6.9.2\PYGZhy{}custom
\PYGZdl{}\PYG{+w}{ }sudo\PYG{+w}{ }cp\PYG{+w}{ }.config\PYG{+w}{ }/boot/config\PYGZhy{}6.9.2\PYGZhy{}custom
\end{sphinxVerbatim}

\begin{sphinxadmonition}{note}{Not:}
\sphinxAtStartPar
Çekirdek modüllerinin kopyalanması biraz zahmetli bir işlemdir çünkü bunlar derlediğimiz
çekirdekte farklı dizinlerde bulunmaktadır. Bu kopyalamanın en etkin yolu \sphinxstyleemphasis{“make modules\_install”}
komutunu kullanmaktır. Benzer biçimde çekirdek dosyalarının ve gerekli diğer dosyaların uygun
yerlere kopyalanması için de en etkin yöntem \sphinxstyleemphasis{“make install”} komutudur.
\end{sphinxadmonition}


\section{GRUB Yapılandırması}
\label{\detokenize{kernelcompilation:grub-yapilandirmasi}}
\sphinxAtStartPar
Normal olarak \sphinxstyleemphasis{“make install”} yaptığımızda eğer sistemimizde GRUB önyükleyicisi varsa komut GRUB
konfigürasyon dosyalarında da güncellemeler yaparak sistemin yeni çekirdekle açılmasını sağlamaktadır.
Böylece kullanıcı bir menü yoluyla sistemin kendi istediği çekirdekle açılmasını da sağlayabilmektedir.
GRUB menüsü otomatik olarak görüntülenmemektedir. Boot işlemi sırasında ESC tuşuna basılırsa menü
görüntülenir. Eğer GRUB menüsünün her zaman görüntülenmesi isteniyorsa \sphinxcode{\sphinxupquote{/etc/default/grub}}
dosyasındaki iki satır aşağıdaki gibi değiştirilmelidir:

\begin{sphinxVerbatim}[commandchars=\\\{\}]
GRUB\PYGZus{}TIMEOUT\PYGZus{}STYLE=menu
GRUB\PYGZus{}TIMEOUT=5
\end{sphinxVerbatim}

\sphinxAtStartPar
Buradaki \sphinxcode{\sphinxupquote{GRUB\_TIMEOUT}} satırı eğer müdahale yapılmamışsa menünün en fazla 5 saniye
görüntüleneceğini belirtmektedir.

\sphinxAtStartPar
Bu işlemden sonra \sphinxstyleemphasis{update\sphinxhyphen{}grub} programı da çalıştırılmalıdır:

\begin{sphinxVerbatim}[commandchars=\\\{\}]
\PYGZdl{}\PYG{+w}{ }sudo\PYG{+w}{ }update\PYGZhy{}grub
\end{sphinxVerbatim}

\sphinxAtStartPar
Bu tür denemeler yapılırken GRUB menüleri bozulabilmektedir. Düzeltme işlemleri bazı konfigürasyon
dosyalarının düzenlenmesiyle manuel biçimde yapılabilir. Konfigürasyon dosyaları güncellendikten
sonra \sphinxstyleemphasis{update\sphinxhyphen{}grub} programı mutlaka çalıştırılmalıdır. Ancak eğer GRUB konfigürasyon dosyaları
konusunda yeterli bilgiye sahip değilseniz GRUB işlemlerini görsel bir biçimde \sphinxstyleemphasis{grub\sphinxhyphen{}customizer}
isimli programla da yapabilirsiniz. Bu program Debian depolarında olmadığı için önce aşağıdaki gibi
programın bulunduğu yerin \sphinxcode{\sphinxupquote{apt}} kayıtlarına eklenmesi gerekmektedir:

\begin{sphinxVerbatim}[commandchars=\\\{\}]
\PYGZdl{}\PYG{+w}{ }sudo\PYG{+w}{ }add\PYGZhy{}apt\PYGZhy{}repository\PYG{+w}{ }ppa:danielrichter2007/grub\PYGZhy{}customizer
\PYGZdl{}\PYG{+w}{ }sudo\PYG{+w}{ }apt\PYGZhy{}get\PYG{+w}{ }update
\PYGZdl{}\PYG{+w}{ }sudo\PYG{+w}{ }apt\PYGZhy{}get\PYG{+w}{ }install\PYG{+w}{ }grub\PYGZhy{}customizer
\end{sphinxVerbatim}


\section{Çekirdek Derleme Süreci Özeti}
\label{\detokenize{kernelcompilation:cekirdek-derleme-sureci-ozeti}}
\sphinxAtStartPar
Yukarıdaki adımların özeti şöyledir:
\begin{enumerate}
\sphinxsetlistlabels{\arabic}{enumi}{enumii}{}{.}%
\item {} 
\sphinxAtStartPar
Çekirdek derlemesi için gerekli olan araçlar kurulur.

\item {} 
\sphinxAtStartPar
Çekirdek kodları indirilir ve açılır.

\item {} 
\sphinxAtStartPar
Zaten hazır olan konfigürasyon dosyası \sphinxcode{\sphinxupquote{/boot}} dizininden alınarak \sphinxcode{\sphinxupquote{.config}} ismiyle kaynak
kök dizine kopyalanır.

\item {} 
\sphinxAtStartPar
Konfigürasyon dosyası üzerinde \sphinxstyleemphasis{“make menuconfig”} komutu ile değişiklikler yapılır.

\item {} 
\sphinxAtStartPar
Eğer çekirdeğin imzalanması istenmiyorsa \sphinxcode{\sphinxupquote{.config}} dosyasındaki ilgili satırlar üzerinde
değişiklikler yapılır.

\item {} 
\sphinxAtStartPar
Çekirdek derlemesi \sphinxstyleemphasis{“make \sphinxhyphen{}j\$(nproc)”} komutu ile gerçekleştirilir.

\item {} 
\sphinxAtStartPar
Modüller ve ilgili dosyalar hedefe \sphinxstyleemphasis{“sudo make modules\_install”} komutu ile konuşlandırılır.

\item {} 
\sphinxAtStartPar
Çekirdek imajı ve ilgili dosyalar masaüstü sistemlerde \sphinxstyleemphasis{“sudo make install”} komutu ile hedefe
konuşlandırılır.

\end{enumerate}


\section{Çekirdeğin Sistemden Kaldırılması}
\label{\detokenize{kernelcompilation:cekirdegin-sistemden-kaldirilmasi}}
\sphinxAtStartPar
Yeni çekirdeği derleyip sisteme dahil ettikten sonra onu sistemden tamamen çıkartmak için yapılan
işlemlerin tersini yapmak gerekir. Kaldırma işlemi manuel biçimde şöyle yapılabilir:
\begin{itemize}
\item {} 
\sphinxAtStartPar
\sphinxcode{\sphinxupquote{/lib/modules/\textless{}çekirdek\_sürümü\textgreater{}}} dizini tamamen silinir.

\item {} 
\sphinxAtStartPar
\sphinxcode{\sphinxupquote{/boot}} dizinindeki çekirdek sürümüne ilişkin dosyalar silinir.

\item {} 
\sphinxAtStartPar
\sphinxcode{\sphinxupquote{/boot}} dizininden çekirdek sürümüne ilişkin dosyalar silindikten sonra \sphinxstyleemphasis{update\sphinxhyphen{}grub} programı
\sphinxcode{\sphinxupquote{sudo}} ile çalıştırılmalıdır. Bu program \sphinxcode{\sphinxupquote{/boot}} dizinini inceleyip otomatik olarak ilgili
girişleri GRUB menüsünden siler. Yani aslında GRUB konfigürasyon dosyaları üzerinde manuel
değişiklik yapmaya gerek yoktur.

\end{itemize}

\begin{sphinxadmonition}{note}{Not:}
\sphinxAtStartPar
\sphinxstyleemphasis{grub\sphinxhyphen{}customizer} programı ile de görsel silme yapılabilir. Ancak bu program \sphinxcode{\sphinxupquote{/boot}} dizini
içerisindeki dosyaları ve modül dosyalarını silmez; yalnızca ilgili girişleri GRUB menüsünden
çıkartmaktadır.
\end{sphinxadmonition}


\section{Çekirdek Kodları Üzerinde Değişiklik Yapma Yöntemleri}
\label{\detokenize{kernelcompilation:cekirdek-kodlari-uzerinde-degisiklik-yapma-yontemleri}}
\sphinxAtStartPar
Çekirdeği yeniden derlemenin gerekçelerinden bahsetmiştik. Bunlardan biri de çekirdek kodları üzerinde
değişikliklerin yapılmış olmasıydı. Peki çekirdek kodları üzerinde değişiklikler nasıl yapılabilir?
Çekirdek kodları üzerinde değişiklikler tipik olarak dört yolla yapılmaktadır:
\begin{enumerate}
\sphinxsetlistlabels{\arabic}{enumi}{enumii}{}{.}%
\item {} 
\sphinxAtStartPar
Çekirdek kodlarındaki zaten var olan bir dosya içerisinde bulunan fonksiyon kodlarında değişiklik
yapılması.

\item {} 
\sphinxAtStartPar
Çekirdek kodlarındaki zaten var olan bir dosya içerisine yeni bir fonksiyon eklenmesi.

\item {} 
\sphinxAtStartPar
Çekirdek kodlarındaki bir dizin içerisine yeni bir C kaynak dosyası eklenmesi.

\item {} 
\sphinxAtStartPar
Çekirdek kodlarındaki bir dizin içerisine yeni bir dizin ve bu dizinin içerisine de çok sayıda C
kaynak dosyalarının eklenmesi.

\end{enumerate}

\sphinxAtStartPar
Eğer biz birinci ve ikinci maddedeki gibi çekirdek kodlarına yeni bir dosya eklemiyorsak çekirdeğin
derlenmesini sağlayan make dosyalarında bir değişiklik yapmamıza gerek yoktur. Ancak çekirdeğe yeni bir
kaynak dosya ya da dizin ekleyeceksek bu eklemeyi yaptığımız dizindeki make dosyasında izleyen paragraflarda
açıklayacağımız biçimde bazı güncellemelerin yapılması gerekir. Böylece çekirdek yeniden derlendiğinde bu
dosyalar da çekirdek imajının içerisine eklenmiş olacaktır. Eğer kaynak kod ağacında bir dizinin altına yeni
bir dizin eklemek istiyorsak bu durumda o dizini yine üst dizine ilişkin make dosyasında belirtmemiz ve o
dizinde ayrı bir \sphinxcode{\sphinxupquote{Makefile}} dosyası oluşturmamız gerekir.


\section{KBuild Sistemi ve GNU Make}
\label{\detokenize{kernelcompilation:kbuild-sistemi-ve-gnu-make}}
\sphinxAtStartPar
GNU Make aracı oldukça ayrıntılı özelliklere sahip bir build aracıdır. Bu aracın ayrıntılarını öğrenmek
ayrı bir çabayı gerektirmektedir. Make dili aslında oldukça aşağı seviyeli bir build dilidir. Bu nedenle
özellikle son yirmi yıldır programcılar doğrudan GNU Make aracını kullanmak yerine daha üst düzey make
araçlarını kullanmayı tercih etmektedir. Bunlardan en yaygın olanlardan biri \sphinxstyleemphasis{CMake} denilen araçtır.
Microsoft ise \sphinxstyleemphasis{MSBuild} isimli kendi tasarladığı build aracını kullanmaktadır.

\sphinxAtStartPar
Make dilinde değişkenler oluşturulabilmektedir. Örneğin:

\begin{sphinxVerbatim}[commandchars=\\\{\}]
\PYG{n+nv}{obj\PYGZhy{}y}\PYG{+w}{ }\PYG{o}{=}\PYG{+w}{ }a.o
\end{sphinxVerbatim}

\sphinxAtStartPar
Burada \sphinxcode{\sphinxupquote{obj\sphinxhyphen{}y}} isimli değişken \sphinxcode{\sphinxupquote{a.o}} bilgisini tutmaktadır. Değişkenleri bir çeşit makro gibi
düşünebilirsiniz. Bir değişkene ekleme yapmak için Make dilinde \sphinxcode{\sphinxupquote{+=}} operatörü kullanılmaktadır. Örneğin:

\begin{sphinxVerbatim}[commandchars=\\\{\}]
\PYG{n+nv}{obj\PYGZhy{}y}\PYG{+w}{ }\PYG{o}{=}\PYG{+w}{ }a.o
\PYG{n+nv}{obj\PYGZhy{}y}\PYG{+w}{ }\PYG{o}{+=}\PYG{+w}{ }b.o
\PYG{n+nv}{obj\PYGZhy{}y}\PYG{+w}{ }\PYG{o}{+=}\PYG{+w}{ }c.o
\end{sphinxVerbatim}

\sphinxAtStartPar
Burada artık \sphinxcode{\sphinxupquote{obj\sphinxhyphen{}y}} değişkeni \sphinxcode{\sphinxupquote{a.o b.o c.o}} biçiminde olacaktır.

\sphinxAtStartPar
Linux çekirdeğinde özyinelemeli bir make yöntemi kullanılmaktadır. Her dizinde bir \sphinxcode{\sphinxupquote{Makefile}} dosyası
vardır. Bunun içerisindeki \sphinxcode{\sphinxupquote{obj\sphinxhyphen{}y}} gibi, \sphinxcode{\sphinxupquote{obj\sphinxhyphen{}m}} gibi bazı değişkenler \sphinxcode{\sphinxupquote{+=}} operatörüyle eklenerek
biriktirilmektedir. Bunlar da derleme ve bağlama işlemine sokulmaktadır. Yukarıda da belirttiğimiz gibi
Linux’taki bu build sistemine \sphinxstyleemphasis{KBuild} ya da \sphinxstyleemphasis{KConfig} sistemi denilmektedir.

\sphinxAtStartPar
Bizim Linux’ta \sphinxcode{\sphinxupquote{Makefile}} dosyaları üzerinde gerekli güncellemeleri yapmak için çok fazla bilgiye sahip
olmamız gerekmez. Bazı yönergeleri uygun bir biçimde yerine getirirsek hedefimize ulaşabiliriz.


\section{Makefile Dosyalarının Güncellenmesi}
\label{\detokenize{kernelcompilation:makefile-dosyalarinin-guncellenmesi}}
\sphinxAtStartPar
Linux kaynak kod ağacında dizinlerin altında \sphinxcode{\sphinxupquote{Makefile}} isimli make dosyaları bulunur. Eğer bir dizinin
altına yeni bir dosya eklenecekse o dizinin içerisinde bulunan \sphinxcode{\sphinxupquote{Makefile}} dosyasının içerisine aşağıdaki
gibi bir satırın eklenmesi gerekir:

\begin{sphinxVerbatim}[commandchars=\\\{\}]
\PYG{n+nv}{obj\PYGZhy{}y}\PYG{+w}{ }\PYG{o}{+=}\PYG{+w}{ }dosya\PYGZus{}ismi.o
\end{sphinxVerbatim}

\sphinxAtStartPar
Buradaki \sphinxcode{\sphinxupquote{+=}} operatörü \sphinxcode{\sphinxupquote{obj\sphinxhyphen{}y}} isimli hedefe ekleme yapma anlamına gelmektedir. \sphinxcode{\sphinxupquote{obj}} sözcüğünün
yanındaki \sphinxcode{\sphinxupquote{\sphinxhyphen{}y}} harfi ilgili dosyanın çekirdeğin bir parçası biçiminde çekirdek imajının içerisine
gömüleceğini belirtmektedir. Make dosyalarının bazı satırlarında \sphinxcode{\sphinxupquote{obj\sphinxhyphen{}y}} yerine \sphinxcode{\sphinxupquote{obj\sphinxhyphen{}m}} de görebilirsiniz.
Bu da ilgili dosyanın ayrı bir modül biçiminde derleneceği anlamına gelmektedir. Eklemeler genellikle
çekirdek imajının içine yapıldığı için biz de genellikle \sphinxcode{\sphinxupquote{obj\sphinxhyphen{}y}} kullanırız. Eğer bir dosyanın (aygıt
sürücüler için bu durum söz konusudur) çekirdek imajının içine gömülmesi yerine ayrı bir çekirdek modülü
olarak derlenmesi isteniyorsa bu durumda dosyanın yerleştirildiği dizinin \sphinxcode{\sphinxupquote{Makefile}} dosyasına aşağıdaki
gibi bir eklemenin yapılması gerekir:

\begin{sphinxVerbatim}[commandchars=\\\{\}]
\PYG{n+nv}{obj\PYGZhy{}m}\PYG{+w}{ }\PYG{o}{+=}\PYG{+w}{ }dosya\PYGZus{}ismi.o
\end{sphinxVerbatim}

\sphinxAtStartPar
Eğer çekirdek kaynak kodlarına tümden bir dizinin eklenmesi isteniyorsa bu durumda önce o dizinin
oluşturulduğu dizindeki \sphinxcode{\sphinxupquote{Makefile}} dosyasına aşağıdaki gibi bir satır eklenmelidir (dizin isminden
sonra \sphinxcode{\sphinxupquote{/}} karakterini unutmayınız):

\begin{sphinxVerbatim}[commandchars=\\\{\}]
\PYG{n+nv}{obj\PYGZhy{}y}\PYG{+w}{ }\PYG{o}{+=}\PYG{+w}{ }dizin\PYGZus{}ismi/
\end{sphinxVerbatim}

\sphinxAtStartPar
Tabii bu ekleme bir modül biçiminde de olabilirdi:

\begin{sphinxVerbatim}[commandchars=\\\{\}]
\PYG{n+nv}{obj\PYGZhy{}m}\PYG{+w}{ }\PYG{o}{+=}\PYG{+w}{ }dizin\PYGZus{}ismi/
\end{sphinxVerbatim}

\sphinxAtStartPar
Fakat bu ekleme tek başına yetmemektedir. Bu ekleme yapıldıktan sonra ayrıca yaratılan dizinde \sphinxcode{\sphinxupquote{Makefile}}
isimli bir dosyanın oluşturulması ve o dosyanın içerisinde o dizindeki kaynak dosyaların da belirtilmesi
gerekmektedir. Örneğin biz \sphinxcode{\sphinxupquote{drivers}} dizininin altında \sphinxcode{\sphinxupquote{mydriver}} isimli bir dizin oluşturup onun da
içerisine \sphinxcode{\sphinxupquote{a.c}}, \sphinxcode{\sphinxupquote{b.c}} ve \sphinxcode{\sphinxupquote{c.c}} dosyalarını eklemiş olalım. Bu durumda önce \sphinxcode{\sphinxupquote{drivers}} dizini
içerisindeki \sphinxcode{\sphinxupquote{Makefile}} dosyasına aşağıdaki gibi bir satır ekleriz:

\begin{sphinxVerbatim}[commandchars=\\\{\}]
\PYG{n+nv}{obj\PYGZhy{}y}\PYG{+w}{ }\PYG{o}{+=}\PYG{+w}{ }mydriver/
\end{sphinxVerbatim}

\sphinxAtStartPar
Sonra da \sphinxcode{\sphinxupquote{mydriver}} dizini içerisinde \sphinxcode{\sphinxupquote{Makefile}} isimli bir dosya oluşturup bu dosyanın içerisinde
de dizin içerisindeki dosyaları belirtiriz. Örneğin:

\begin{sphinxVerbatim}[commandchars=\\\{\}]
\PYG{n+nv}{obj\PYGZhy{}y}\PYG{+w}{ }\PYG{o}{+=}\PYG{+w}{ }a.o
\PYG{n+nv}{obj\PYGZhy{}y}\PYG{+w}{ }\PYG{o}{+=}\PYG{+w}{ }b.o
\PYG{n+nv}{obj\PYGZhy{}y}\PYG{+w}{ }\PYG{o}{+=}\PYG{+w}{ }c.o
\end{sphinxVerbatim}


\section{Kconfig Dosyaları}
\label{\detokenize{kernelcompilation:kconfig-dosyalari}}
\sphinxAtStartPar
Kaynak kod ağacında \sphinxcode{\sphinxupquote{Makefile}} dosyasının dışında build sistemiyle ilgili \sphinxcode{\sphinxupquote{Kconfig}} isimli dosyalar
da bulunmaktadır. Bu dosyaların içerisinde ilgili dosyaların ya da dizinlerin konfigürasyon dosyasına
yansıtılması için gerekli bilgiler bulundurulmaktadır. Örneğin biz eklediğimiz \sphinxcode{\sphinxupquote{mydriver}} dizinindeki
dosyaların çekirdek kodlarına dahil edilip edilmeyeceğini çekirdeği derleyenin konfigürasyon aşamasında
belirlemesini sağlayabiliriz. Bunun için bu \sphinxcode{\sphinxupquote{Kconfig}} dosyasına bir giriş eklememiz gerekir. Böylece
bu giriş de \sphinxstyleemphasis{“make menuconfig”} yapıldığında bir seçenek olarak karşımıza gelecektir. Tabii ekleyeceğimiz
dosya ve dizinleri \sphinxcode{\sphinxupquote{Kconfig}} dosyasında belirtmek zorunda değiliz.

\sphinxAtStartPar
\sphinxstyleemphasis{“make menuconfig”} menüsünde seçenekler için birkaç seçme biçimi bulunmaktadır:


\begin{savenotes}\sphinxattablestart
\sphinxthistablewithglobalstyle
\centering
\begin{tabular}[t]{\X{15}{100}\X{85}{100}}
\sphinxtoprule
\begin{varwidth}[t]{\sphinxcolwidth{1}{2}}
\sphinxstyletheadfamily \sphinxAtStartPar
Gösterim
\sphinxbeforeendvarwidth
\end{varwidth}%
&\begin{varwidth}[t]{\sphinxcolwidth{1}{2}}
\sphinxstyletheadfamily \sphinxAtStartPar
Anlamı
\sphinxbeforeendvarwidth
\end{varwidth}%
\\
\sphinxmidrule
\sphinxtableatstartofbodyhook\begin{varwidth}[t]{\sphinxcolwidth{1}{2}}
\sphinxAtStartPar
\sphinxcode{\sphinxupquote{{[} {]}}}
\sphinxbeforeendvarwidth
\end{varwidth}%
&\begin{varwidth}[t]{\sphinxcolwidth{1}{2}}
\sphinxAtStartPar
Seçilebilir ya da seçilmeyebilir. Seçildiğinde \sphinxcode{\sphinxupquote{{[}*{]}}} biçiminde gösterilir.
\sphinxbeforeendvarwidth
\end{varwidth}%
\\
\sphinxhline\begin{varwidth}[t]{\sphinxcolwidth{1}{2}}
\sphinxAtStartPar
\sphinxcode{\sphinxupquote{\textless{} \textgreater{}}}
\sphinxbeforeendvarwidth
\end{varwidth}%
&\begin{varwidth}[t]{\sphinxcolwidth{1}{2}}
\sphinxAtStartPar
Üç konumlu seçenek: boş (seçilmedi), \sphinxcode{\sphinxupquote{\textless{}M\textgreater{}}} (modül) ya da \sphinxcode{\sphinxupquote{\textless{}*\textgreater{}}} (çekirdek içine gömülü).
\sphinxbeforeendvarwidth
\end{varwidth}%
\\
\sphinxhline\begin{varwidth}[t]{\sphinxcolwidth{1}{2}}
\sphinxAtStartPar
\sphinxcode{\sphinxupquote{\sphinxhyphen{}*\sphinxhyphen{}}}
\sphinxbeforeendvarwidth
\end{varwidth}%
&\begin{varwidth}[t]{\sphinxcolwidth{1}{2}}
\sphinxAtStartPar
Seçilememezlik yapılamaz; ilgili özellik mutlaka çekirdek kodlarında bulunmalıdır.
\sphinxbeforeendvarwidth
\end{varwidth}%
\\
\sphinxhline\begin{varwidth}[t]{\sphinxcolwidth{1}{2}}
\sphinxAtStartPar
değer
\sphinxbeforeendvarwidth
\end{varwidth}%
&\begin{varwidth}[t]{\sphinxcolwidth{1}{2}}
\sphinxAtStartPar
İlgili özellik için doğrudan bir değer girilmektedir.
\sphinxbeforeendvarwidth
\end{varwidth}%
\\
\sphinxbottomrule
\end{tabular}
\sphinxtableafterendhook\par
\sphinxattableend\end{savenotes}

\sphinxAtStartPar
Tabii \sphinxstyleemphasis{“make menuconfig”} menüsünde yapılan her şey aslında \sphinxcode{\sphinxupquote{.config}} dosyasına yansıtılmaktadır.


\section{Konfigürasyon Seçeneklerinin Kaynak Kodlara Yansıtılması}
\label{\detokenize{kernelcompilation:konfigurasyon-seceneklerinin-kaynak-kodlara-yansitilmasi}}
\sphinxAtStartPar
Peki çekirdeğin konfigüre edilmesi aşamasında \sphinxstyleemphasis{“make menuconfig”} işleminde belirlediğimiz seçenekler
kaynak kodlara nasıl yansıtılmaktadır? Örneğin biz \sphinxstyleemphasis{“make menuconfig”} işleminde bir modülün çekirdek
kodlarına dahil edilmesini ilgili girişi \sphinxcode{\sphinxupquote{*}} ile seçerek sağlayabilmekteyiz. Benzer biçimde biz
konfigürasyon aşamasında bazı çekirdek parametrelerini de değiştirebilmekteyiz. Örneğin \sphinxstyleemphasis{timer tick}
frekansı \sphinxstyleemphasis{“make menuconfig”} menüsünde bir sayı biçiminde belirlenebilmektedir. Peki buradaki belirlemeler
çekirdek kodlarına ve build sistemine nasıl yansıtılmaktadır?

\sphinxAtStartPar
Anımsanacağı gibi \sphinxstyleemphasis{“make menuconfig”} ve diğer config menülerinde yapılan seçimler \sphinxcode{\sphinxupquote{.config}} isimli
bir metin dosyaya kaydedilmektedir. Bu \sphinxcode{\sphinxupquote{.config}} dosyası \sphinxstyleemphasis{özellik=değer} biçiminde satırlardan
oluşmaktadır. Aşağıda dosyanın birkaç satırını görüyorsunuz:

\begin{sphinxVerbatim}[commandchars=\\\{\}]
CONFIG\PYGZus{}CC\PYGZus{}HAS\PYGZus{}ASM\PYGZus{}INLINE=y
CONFIG\PYGZus{}CC\PYGZus{}HAS\PYGZus{}NO\PYGZus{}PROFILE\PYGZus{}FN\PYGZus{}ATTR=y
CONFIG\PYGZus{}PAHOLE\PYGZus{}VERSION=125
CONFIG\PYGZus{}IRQ\PYGZus{}WORK=y
CONFIG\PYGZus{}BUILDTIME\PYGZus{}TABLE\PYGZus{}SORT=y
CONFIG\PYGZus{}THREAD\PYGZus{}INFO\PYGZus{}IN\PYGZus{}TASK=y
\end{sphinxVerbatim}

\sphinxAtStartPar
Çekirdek derlenirken ilk aşamada bu \sphinxcode{\sphinxupquote{.config}} dosyasının içeriği
\sphinxcode{\sphinxupquote{include/generated/autoconf.h}} dosyasının içerisine \sphinxcode{\sphinxupquote{\#define}} önişlemci komutları biçiminde
aktarılmaktadır. Eğer ilgili konfigürasyon dosyasındaki değer \sphinxcode{\sphinxupquote{y}} ya da \sphinxcode{\sphinxupquote{m}} ise bu \sphinxcode{\sphinxupquote{autoconf.h}}
dosyası içerisinde buna ilişkin sembolik sabit 1 olarak görünür. (Başka bir deyişle \sphinxstyleemphasis{“make menuconfig”}
menüsünde \sphinxcode{\sphinxupquote{{[}*{]}}} ya da \sphinxcode{\sphinxupquote{\textless{}*\textgreater{}}} ya da \sphinxcode{\sphinxupquote{\textless{}M\textgreater{}}} seçenekleri için sembolik sabit 1 olur.) Eğer konfigürasyon
dosyasında ilgili seçenek gerçekten bir değer belirtiyorsa \sphinxcode{\sphinxupquote{autoconf.h}} dosyası içerisinde bu sembolik
sabit o değerde olur. Eğer konfigürasyon dosyasında ilgili seçenek \sphinxcode{\sphinxupquote{n}} biçiminde seçilmişse (yani
\sphinxstyleemphasis{“make menuconfig”} menüsünde ilgili seçenek \sphinxcode{\sphinxupquote{{[} {]}}} ya da \sphinxcode{\sphinxupquote{\textless{} \textgreater{}}} biçiminde seçilmişse) bu durumda
ilgili sembolik sabit hiç \sphinxcode{\sphinxupquote{\#define}} edilmemiş hale gelir.

\sphinxAtStartPar
Özetle aslında \sphinxcode{\sphinxupquote{.config}} dosyası içerisindeki satırlardan C dilinde anlamlı olan \sphinxcode{\sphinxupquote{\#define}} önişlemci
komutları oluşturulmaktadır. Aşağıda üretilmiş olan \sphinxcode{\sphinxupquote{autoconf.h}} dosyasının birkaç satırını görüyorsunuz:

\begin{sphinxVerbatim}[commandchars=\\\{\}]
\PYG{c+cp}{\PYGZsh{}}\PYG{c+cp}{define CONFIG\PYGZus{}IGB\PYGZus{}HWMON 1}
\PYG{c+cp}{\PYGZsh{}}\PYG{c+cp}{define CONFIG\PYGZus{}ACPI\PYGZus{}HOTPLUG\PYGZus{}CPU 1}
\PYG{c+cp}{\PYGZsh{}}\PYG{c+cp}{define CONFIG\PYGZus{}DEV\PYGZus{}DAX\PYGZus{}KMEM\PYGZus{}MODULE 1}
\PYG{c+cp}{\PYGZsh{}}\PYG{c+cp}{define CONFIG\PYGZus{}RIONET\PYGZus{}RX\PYGZus{}SIZE 128}
\PYG{c+cp}{\PYGZsh{}}\PYG{c+cp}{define CONFIG\PYGZus{}USB\PYGZus{}SERIAL\PYGZus{}KEYSPAN\PYGZus{}PDA\PYGZus{}MODULE 1}
\PYG{c+cp}{\PYGZsh{}}\PYG{c+cp}{define CONFIG\PYGZus{}BOOTTIME\PYGZus{}TRACING 1}
\end{sphinxVerbatim}

\sphinxAtStartPar
Çekirdeğe ilişkin bir C kodu içerisinde “ilgili seçenek seçilmişse” bir kod parçasını derlemeye dahil
etmek için \sphinxcode{\sphinxupquote{\#ifdef}} önişlemci komutundan faydalanabilirsiniz. Örneğin:

\begin{sphinxVerbatim}[commandchars=\\\{\}]
\PYG{c+cp}{\PYGZsh{}}\PYG{c+cp}{ifdef CONFIG\PYGZus{}XXX}
\PYG{p}{.}\PYG{p}{.}\PYG{p}{.}
\PYG{c+cp}{\PYGZsh{}}\PYG{c+cp}{endif}
\end{sphinxVerbatim}

\sphinxAtStartPar
Tabii üretilen bu \sphinxcode{\sphinxupquote{autoconf.h}} dosyası çekirdek kaynak kodlarındaki çeşitli başlık dosyalarında
doğrudan ya da dolaylı bir biçimde eklenmiş durumdadır. Linux kaynak kodları da bu sembolik sabitleri
kullanacak biçimde yazılmıştır. Linux’un kaynak kodlarında \sphinxcode{\sphinxupquote{autoconf.h}} dosyasını bulamazsınız.
Çünkü bu dosya make işlemi sırasında oluşturulmaktadır.


\section{Konfigürasyon Seçeneklerinin Makefile Dosyalarına Yansıtılması}
\label{\detokenize{kernelcompilation:konfigurasyon-seceneklerinin-makefile-dosyalarina-yansitilmasi}}
\sphinxAtStartPar
Biz yukarıdaki paragrafta \sphinxcode{\sphinxupquote{.config}} dosyasındaki konfigürasyon parametrelerinin nasıl C’ye sembolik
sabitler biçiminde yansıtıldığını gördük. Peki bu konfigürasyon seçenekleri \sphinxcode{\sphinxupquote{Makefile}} dosyalarına
nasıl yansıtılmaktadır? Aşağıda çekirdeğin bir \sphinxcode{\sphinxupquote{Makefile}} içeriğini görüyorsunuz:

\begin{sphinxVerbatim}[commandchars=\\\{\}]
\PYG{n+nv}{obj\PYGZhy{}\PYGZdl{}(CONFIG\PYGZus{}I8254)}\PYG{+w}{             }\PYG{o}{+=}\PYG{+w}{ }i8254.o
\PYG{n+nv}{obj\PYGZhy{}\PYGZdl{}(CONFIG\PYGZus{}104\PYGZus{}QUAD\PYGZus{}8)}\PYG{+w}{        }\PYG{o}{+=}\PYG{+w}{ }\PYG{l+m}{104}\PYGZhy{}quad\PYGZhy{}8.o
\PYG{n+nv}{obj\PYGZhy{}\PYGZdl{}(CONFIG\PYGZus{}INTERRUPT\PYGZus{}CNT)}\PYG{+w}{     }\PYG{o}{+=}\PYG{+w}{ }interrupt\PYGZhy{}cnt.o
\PYG{n+nv}{obj\PYGZhy{}\PYGZdl{}(CONFIG\PYGZus{}RZ\PYGZus{}MTU3\PYGZus{}CNT)}\PYG{+w}{       }\PYG{o}{+=}\PYG{+w}{ }rz\PYGZhy{}mtu3\PYGZhy{}cnt.o
\PYG{n+nv}{obj\PYGZhy{}\PYGZdl{}(CONFIG\PYGZus{}STM32\PYGZus{}TIMER\PYGZus{}CNT)}\PYG{+w}{   }\PYG{o}{+=}\PYG{+w}{ }stm32\PYGZhy{}timer\PYGZhy{}cnt.o
\end{sphinxVerbatim}

\sphinxAtStartPar
Bu satırlar aslında “ilgili konfigürasyon seçenekleri seçilmişse ilgili dosyaların derlemeye dahil
edileceğini” belirtmektedir.

\sphinxAtStartPar
İşte \sphinxstyleemphasis{KBuild} sistemi aynı zamanda bu \sphinxcode{\sphinxupquote{.config}} dosyasından hareketle make programı için anlamlı
olan değişkenler de oluşturmaktadır. Bu değişkenleri yukarıda açıkladığımız sembolik sabitlerle
karıştırmayınız. Bu değişkenler make dili için anlamlı olan make dilinin değişkenleridir. Örneğin eğer
\sphinxcode{\sphinxupquote{.config}} dosyasında bir seçenek \sphinxcode{\sphinxupquote{y}} olarak belirtilmişse (başka bir deyişle \sphinxstyleemphasis{“make menuconfig”}
menüsünde seçenek \sphinxcode{\sphinxupquote{{[}*{]}}} ya \sphinxcode{\sphinxupquote{\textless{}*\textgreater{}}} biçiminde seçilmişse) bu konfigürasyon seçeneği için make dilinde
\sphinxcode{\sphinxupquote{y}} değeri, eğer \sphinxcode{\sphinxupquote{m}} olarak belirtilmişse de (başka bir deyişle \sphinxstyleemphasis{“make menuconfig”} menüsünde
seçenek \sphinxcode{\sphinxupquote{\textless{}M\textgreater{}}} biçiminde seçilmişse) \sphinxcode{\sphinxupquote{m}} değeri oluşturulmaktadır.

\sphinxAtStartPar
Böylece aslında yukarıdaki make satırları ilgili seçenek \sphinxcode{\sphinxupquote{y}} olarak seçilmişse \sphinxcode{\sphinxupquote{obj\sphinxhyphen{}y}} biçimine,
\sphinxcode{\sphinxupquote{m}} olarak seçilmişse \sphinxcode{\sphinxupquote{obj\sphinxhyphen{}m}} biçimine dönüştürülmektedir. Örneğin:

\begin{sphinxVerbatim}[commandchars=\\\{\}]
\PYG{n+nv}{obj\PYGZhy{}\PYGZdl{}(CONFIG\PYGZus{}I8254)}\PYG{+w}{ }\PYG{o}{+=}\PYG{+w}{ }i8254.o
\end{sphinxVerbatim}

\sphinxAtStartPar
Burada \sphinxcode{\sphinxupquote{CONFIG\_I8254}} seçeneği \sphinxcode{\sphinxupquote{y}} ise ilgili dosya \sphinxcode{\sphinxupquote{obj\sphinxhyphen{}y}} değişkenine dahil olacaktır. Yani
çekirdeğin içerisinde bulunacaktır. Ancak bu \sphinxcode{\sphinxupquote{CONFIG\_I8254}} seçeneği \sphinxcode{\sphinxupquote{m}} ise ilgili dosya \sphinxcode{\sphinxupquote{obj\sphinxhyphen{}m}}
değişkenine dahil olacaktır. Eğer bu seçenek \sphinxcode{\sphinxupquote{n}} ise (yani hiç seçilmemişse) çekirdek build sistemi
ya bu \sphinxcode{\sphinxupquote{CONFIG\_I8254}} değişkenini hiç tanımlamamakta ya da bunu \sphinxcode{\sphinxupquote{n}} olarak tanımlamaktadır. Her iki
durumda da artık bu dosya herhangi bir biçimde derlemeye dahil edilmeyecektir.

\sphinxAtStartPar
Çekirdeğe birtakım kodlar ekleyenler eğer eklemeleri \sphinxcode{\sphinxupquote{Kconfig}} dosyası yoluyla konfigürasyona
yansıtmışlarsa bu durumda kendi \sphinxcode{\sphinxupquote{Makefile}} dosyasına bu eklemeleri yukarıdaki gibi girebilirler.
Örneğin biz \sphinxcode{\sphinxupquote{mymodule}} ile temsil ettiğimiz bir modül dosyası oluşturup bu modül dosyasının çekirdek
kodlarına eklenip eklenmeyeceğini konfigürasyonda \sphinxcode{\sphinxupquote{Kconfig}} dosyası yoluyla belirtebiliriz. Bu durumda
\sphinxcode{\sphinxupquote{Makefile}} içerisindeki girişi aşağıdaki gibi de oluşturabiliriz:

\begin{sphinxVerbatim}[commandchars=\\\{\}]
\PYG{n+nv}{obj\PYGZhy{}\PYGZdl{}(CONFIG\PYGZus{}MYMODULE)}\PYG{+w}{ }\PYG{o}{+=}\PYG{+w}{ }mymodule.o
\end{sphinxVerbatim}

\sphinxAtStartPar
Görüldüğü gibi burada aslında konfigüre eden nasıl seçmişse biz onun seçimini yansıtmış olmaktayız.
\sphinxcode{\sphinxupquote{.config}} dosyasında bir özellik \sphinxcode{\sphinxupquote{n}} ise bu durumda ilgili \sphinxcode{\sphinxupquote{Makefile}} satırı şu hale gelecektir:

\begin{sphinxVerbatim}[commandchars=\\\{\}]
\PYG{n+nv}{obj\PYGZhy{}n}\PYG{+w}{ }\PYG{o}{+=}\PYG{+w}{ }mymodule.o
\end{sphinxVerbatim}

\sphinxAtStartPar
\sphinxcode{\sphinxupquote{obj\sphinxhyphen{}n}} biçiminde bir birikim yapılmadığı için zaten bu satır derleme aşamasında dikkate alınmayacaktır.
Ancak bazen sistem programcıları \sphinxcode{\sphinxupquote{y}} durumu için aşağıdaki gibi bir kontrol ile modülü koşullu bir
biçimde de derleme sürecine ekleyebilmektedir:

\begin{sphinxVerbatim}[commandchars=\\\{\}]
\PYG{c+cp}{ifeq (\PYGZdl{}(CONFIG\PYGZus{}MYSYSCALL), y)}
\PYG{+w}{    }obj\PYGZhy{}y\PYG{+w}{ }+\PYG{o}{=}\PYG{+w}{ }mysyscall.o
\PYG{c+cp}{endif}
\end{sphinxVerbatim}

\sphinxAtStartPar
Burada eğer konfigürasyon yapılırken ilgili seçenek \sphinxcode{\sphinxupquote{y}} biçiminde (yani \sphinxcode{\sphinxupquote{{[}*{]}}} ya da \sphinxcode{\sphinxupquote{\textless{}*\textgreater{}}}
biçiminde) geçilmişse bu durumda biz de ilgili dosyayı derlemeye dahil etmiş olduk.
\subsubsection*{── 9. Ders · 16 Ağustos 2025, Cumartesi ──}


\bigskip\hrule\bigskip



\section{Yeni Dizin için Makefile ve Kconfig Dosyaları}
\label{\detokenize{kernelcompilation:yeni-dizin-icin-makefile-ve-kconfig-dosyalari}}
\sphinxAtStartPar
Bir C dosyasını ya da dizini çekirdek kodlarına ekledikten sonra onun konfigürasyon sırasında (örneğin
\sphinxstyleemphasis{“make menuconfig”} işlemi sırasında) görünebilirliğini sağlamak için \sphinxcode{\sphinxupquote{Kconfig}} dosyalarının
kullanıldığını belirtmiştik. Yani \sphinxcode{\sphinxupquote{Kconfig}} dosyaları yaptığımız değişikliklerin konfigüre
edilebilirliğini sağlamak için kullanılmaktadır. \sphinxcode{\sphinxupquote{Kconfig}} dosyalarının genel formatı için aşağıdaki
bağlantıya başvurabilirsiniz:
\begin{quote}

\sphinxAtStartPar
\sphinxurl{https://docs.kernel.org/kbuild/kconfig-language.html}
\end{quote}

\sphinxAtStartPar
\sphinxcode{\sphinxupquote{Kconfig}} dosyaları tıpkı \sphinxcode{\sphinxupquote{Makefile}} dosyalarında olduğu gibi özyinelemeli biçimde işletilmektedir.
Yani biz çekirdek kaynak kod ağacında bir dizin yaratmayıp zaten var olan bir dizinin içerisine bir \sphinxcode{\sphinxupquote{.c}}
dosyası yerleştiriyorsak \sphinxcode{\sphinxupquote{Makefile}} ve \sphinxcode{\sphinxupquote{Kconfig}} dosyaları oluşturmamıza gerek yoktur. Gerekli
işlemleri zaten dizin içerisinde var olan bir \sphinxcode{\sphinxupquote{Makefile}} ve \sphinxcode{\sphinxupquote{Kconfig}} dosyaları üzerinde yapabiliriz.
Ancak eğer biz bir dizin oluşturup onun içerisine dosyalar yerleştireceksek o dizin için bir tane
\sphinxcode{\sphinxupquote{Makefile}} ve bir tane de \sphinxcode{\sphinxupquote{Kconfig}} dosyası oluşturmamız gerekir.

\sphinxAtStartPar
Bir dizin yaratıp onun içerisine dosyalar yerleştirirken o dizin için \sphinxcode{\sphinxupquote{Makefile}} ve \sphinxcode{\sphinxupquote{Kconfig}}
dosyalarının yazılması gerektiğini belirtmiştik. (Tabii aslında \sphinxcode{\sphinxupquote{Kconfig}} dosyasının bulundurulması
zorunlu değildir. Ancak eklenen özelliğin konfigüre edilebilirliğinin sağlanması için gerekmektedir.)
Bu dosyalar oluşturulduktan sonra dış dizindeki \sphinxcode{\sphinxupquote{Makefile}} ve \sphinxcode{\sphinxupquote{Kconfig}} dosyalarında aşağıda
belirtilen işlemler de yapılmalıdır:
\begin{enumerate}
\sphinxsetlistlabels{\arabic}{enumi}{enumii}{}{.}%
\item {} 
\sphinxAtStartPar
Dış dizindeki \sphinxcode{\sphinxupquote{Makefile}} dosyasında alt dizinin dikkate alınacağı aşağıdaki gibi bir satırla
belirtilmelidir:

\begin{sphinxVerbatim}[commandchars=\\\{\}]
\PYG{n+nv}{obj\PYGZhy{}y}\PYG{+w}{ }\PYG{o}{+=}\PYG{+w}{ }\PYGZlt{}dizin\PYGZus{}ismi\PYGZgt{}/
\end{sphinxVerbatim}

\item {} 
\sphinxAtStartPar
Dış dizinin \sphinxcode{\sphinxupquote{Kconfig}} dosyasında iç dizindeki \sphinxcode{\sphinxupquote{Kconfig}} dosyasının dikkate alınması aşağıdaki
gibi bir satırın eklenmesiyle sağlanmaktadır (buradaki yol ifadesi çekirdek kodlarının kök dizinine
göreli olmalıdır):

\begin{sphinxVerbatim}[commandchars=\\\{\}]
source \PYGZdq{}drivers/mydriver/Kconfig\PYGZdq{}
\end{sphinxVerbatim}

\end{enumerate}

\sphinxAtStartPar
Biz \sphinxcode{\sphinxupquote{Kconfig}} dosyasına yukarıdaki gibi bir giriş yerleştirdiğimizde artık \sphinxstyleemphasis{“make menuconfig”} gibi
konfigürasyon menülerinde eklediğimiz \sphinxcode{\sphinxupquote{Kconfig}} elemanı bir menü seçeneği biçiminde karşımıza
çıkacaktır.

\sphinxAtStartPar
Örneğin biz bir aygıt sürücümüzün dosyalarını çekirdeğin kaynak kod ağacında \sphinxcode{\sphinxupquote{drivers}} dizininin
altına \sphinxcode{\sphinxupquote{mydriver}} dizinini açarak eklemek isteyelim. Bu durumda şunları yapmamız gerekir:
\begin{enumerate}
\sphinxsetlistlabels{\arabic}{enumi}{enumii}{}{.}%
\item {} 
\sphinxAtStartPar
\sphinxcode{\sphinxupquote{drivers}} dizini içerisinde \sphinxcode{\sphinxupquote{mydriver}} dizinini yaratıp içerisine \sphinxcode{\sphinxupquote{mydriver.c}} dosyasını
(belki de \sphinxcode{\sphinxupquote{mydriver.h}} gibi bir başlık dosyasını da) yerleştirmeliyiz.

\item {} 
\sphinxAtStartPar
\sphinxcode{\sphinxupquote{drivers/mydriver}} dizininde aşağıdaki gibi bir \sphinxcode{\sphinxupquote{Kconfig}} dosyasını oluşturmalıyız. Buradaki
\sphinxcode{\sphinxupquote{config MYDRIVER}} satırı aslında make dilinde \sphinxcode{\sphinxupquote{CONFIG\_MYDRIVER}} değişkeninin oluşturulmasına yol
açmaktadır:

\begin{sphinxVerbatim}[commandchars=\\\{\}]
config MYDRIVER
    tristate \PYGZdq{}My Driver\PYGZdq{}
    default y
    help
      Enable this option to include support for My Device Driver.
      It can either be built as a module or statically linked into the kernel.
\end{sphinxVerbatim}

\sphinxAtStartPar
Eğer ilgili konfigürasyon seçeneği yes/no biçimindeyse (yani \sphinxcode{\sphinxupquote{{[}*{]}}} biçimindeyse) yukarıdaki config
direktifinde \sphinxstyleemphasis{tristate} yerine \sphinxstyleemphasis{bool} kullanılmalıdır. Örneğin:

\begin{sphinxVerbatim}[commandchars=\\\{\}]
config MYDRIVER
    bool \PYGZdq{}My Driver\PYGZdq{}
    default y
    help
      Enable this option to include support for My Device Driver.
\end{sphinxVerbatim}

\item {} 
\sphinxAtStartPar
Üst dizindeki (\sphinxcode{\sphinxupquote{drivers}} dizinindeki) \sphinxcode{\sphinxupquote{Kconfig}} dosyasına aşağıdaki satırı yerleştirmeliyiz:

\begin{sphinxVerbatim}[commandchars=\\\{\}]
source \PYGZdq{}drivers/mydriver/Kconfig\PYGZdq{}
\end{sphinxVerbatim}

\item {} 
\sphinxAtStartPar
\sphinxcode{\sphinxupquote{drivers/mydriver}} dizinindeki \sphinxcode{\sphinxupquote{Makefile}} dosyası içerisine aşağıdaki gibi bir satır eklemeliyiz:

\begin{sphinxVerbatim}[commandchars=\\\{\}]
\PYG{n+nv}{obj\PYGZhy{}\PYGZdl{}(CONFIG\PYGZus{}MYDRIVER)}\PYG{+w}{ }\PYG{o}{+=}\PYG{+w}{ }mydriver.o
\end{sphinxVerbatim}

\item {} 
\sphinxAtStartPar
Üst dizindeki (yani \sphinxcode{\sphinxupquote{drivers}} dizinindeki) \sphinxcode{\sphinxupquote{Makefile}} içerisine aşağıdaki gibi bir satır
eklemeliyiz:

\begin{sphinxVerbatim}[commandchars=\\\{\}]
\PYG{n+nv}{obj\PYGZhy{}\PYGZdl{}(CONFIG\PYGZus{}MYDRIVER)}\PYG{+w}{ }\PYG{o}{+=}\PYG{+w}{ }mydriver/
\end{sphinxVerbatim}

\end{enumerate}

\sphinxAtStartPar
Tabii eğer siz bir dizin yaratmayıp zaten çekirdek kaynak kod ağacındaki bir dizine bir dosya
yerleştiriyorsanız bu durumda ayrı bir \sphinxcode{\sphinxupquote{Kconfig}} dosyasını oluşturmanıza gerek yoktur. Bu durumda
doğrudan yukarıdaki config içeriğini dizinde zaten var olan \sphinxcode{\sphinxupquote{Kconfig}} dosyasına yerleştirebilirsiniz.


\section{Örnek: Çekirdeğe Yeni Modül Ekleme}
\label{\detokenize{kernelcompilation:ornek-cekirdege-yeni-modul-ekleme}}
\sphinxAtStartPar
Şimdi çekirdeğe bazı kodlar ekleyip onu yeniden derleyerek bir deneme yapalım. Örneğin çekirdeğe yeni
bir çekirdek modülü ekleyelim ve çekirdeğin o modül gömülü olarak başlatılmasını sağlayalım. Aynı
zamanda bu çekirdek modülünün \sphinxstyleemphasis{“make menuconfig”} ile seçilebilmesini de sağlayalım. Çekirdek
modüllerinin nasıl yazılacağını bilmediğinizi varsayıyoruz. Ancak biz yine de örneğimizde “hiçbir şey
yapmayan iskelet bir çekirdek modülü” oluşturacağız. Bu işlem şu adımlardan geçilerek yapılabilir
(kaynak kod ağacının kök dizininde bulunduğumuzu varsayıyoruz):

\sphinxAtStartPar
\sphinxstylestrong{Adım 1:} \sphinxcode{\sphinxupquote{drivers/mydriver}} dizini yaratılır.

\sphinxAtStartPar
\sphinxstylestrong{Adım 2:} İskelet bir çekirdek modülü \sphinxcode{\sphinxupquote{mydriver.c}} biçiminde \sphinxcode{\sphinxupquote{drivers/mydriver}} dizininde
aşağıdaki gibi oluşturulur:

\begin{sphinxVerbatim}[commandchars=\\\{\}]
\PYG{c+cm}{/* mydriver.c */}

\PYG{c+cp}{\PYGZsh{}}\PYG{c+cp}{include}\PYG{+w}{ }\PYG{c+cpf}{\PYGZlt{}linux/module.h\PYGZgt{}}
\PYG{c+cp}{\PYGZsh{}}\PYG{c+cp}{include}\PYG{+w}{ }\PYG{c+cpf}{\PYGZlt{}linux/kernel.h\PYGZgt{}}

\PYG{n}{MODULE\PYGZus{}LICENSE}\PYG{p}{(}\PYG{l+s}{\PYGZdq{}}\PYG{l+s}{GPL}\PYG{l+s}{\PYGZdq{}}\PYG{p}{)}\PYG{p}{;}
\PYG{n}{MODULE\PYGZus{}AUTHOR}\PYG{p}{(}\PYG{l+s}{\PYGZdq{}}\PYG{l+s}{Kaan Aslan}\PYG{l+s}{\PYGZdq{}}\PYG{p}{)}\PYG{p}{;}
\PYG{n}{MODULE\PYGZus{}DESCRIPTION}\PYG{p}{(}\PYG{l+s}{\PYGZdq{}}\PYG{l+s}{General Device Driver}\PYG{l+s}{\PYGZdq{}}\PYG{p}{)}\PYG{p}{;}

\PYG{k}{static}\PYG{+w}{ }\PYG{k+kt}{int}\PYG{+w}{ }\PYG{n}{\PYGZus{}\PYGZus{}init}\PYG{+w}{ }\PYG{n+nf}{mydriver\PYGZus{}init}\PYG{p}{(}\PYG{k+kt}{void}\PYG{p}{)}
\PYG{p}{\PYGZob{}}
\PYG{+w}{    }\PYG{n}{printk}\PYG{p}{(}\PYG{n}{KERN\PYGZus{}INFO}\PYG{+w}{ }\PYG{l+s}{\PYGZdq{}}\PYG{l+s}{Hello World...}\PYG{l+s+se}{\PYGZbs{}n}\PYG{l+s}{\PYGZdq{}}\PYG{p}{)}\PYG{p}{;}
\PYG{+w}{    }\PYG{k}{return}\PYG{+w}{ }\PYG{l+m+mi}{0}\PYG{p}{;}
\PYG{p}{\PYGZcb{}}

\PYG{k}{static}\PYG{+w}{ }\PYG{k+kt}{void}\PYG{+w}{ }\PYG{n}{\PYGZus{}\PYGZus{}exit}\PYG{+w}{ }\PYG{n+nf}{mydriver\PYGZus{}exit}\PYG{p}{(}\PYG{k+kt}{void}\PYG{p}{)}
\PYG{p}{\PYGZob{}}
\PYG{+w}{    }\PYG{n}{printk}\PYG{p}{(}\PYG{n}{KERN\PYGZus{}INFO}\PYG{+w}{ }\PYG{l+s}{\PYGZdq{}}\PYG{l+s}{Goodbye World...}\PYG{l+s+se}{\PYGZbs{}n}\PYG{l+s}{\PYGZdq{}}\PYG{p}{)}\PYG{p}{;}
\PYG{p}{\PYGZcb{}}

\PYG{n}{module\PYGZus{}init}\PYG{p}{(}\PYG{n}{mydriver\PYGZus{}init}\PYG{p}{)}\PYG{p}{;}
\PYG{n}{module\PYGZus{}exit}\PYG{p}{(}\PYG{n}{mydriver\PYGZus{}exit}\PYG{p}{)}\PYG{p}{;}
\end{sphinxVerbatim}

\sphinxAtStartPar
\sphinxstylestrong{Adım 3:} \sphinxcode{\sphinxupquote{drivers/mydriver}} dizininde \sphinxcode{\sphinxupquote{Kconfig}} dosyası aşağıdaki gibi oluşturulmalıdır.
Burada konfigürasyon makrosunun ismi \sphinxcode{\sphinxupquote{CONFIG\_MYDRIVER}} biçiminde olacaktır:

\begin{sphinxVerbatim}[commandchars=\\\{\}]
config MYDRIVER
    tristate \PYGZdq{}My Character Device Driver\PYGZdq{}
    default y
    help
      Enable this option to include support for My Device Driver.
      It can either be built as a module or statically linked into the kernel.
\end{sphinxVerbatim}

\sphinxAtStartPar
\sphinxstylestrong{Adım 4:} Üst dizinin (yani \sphinxcode{\sphinxupquote{drivers}} dizininin) \sphinxcode{\sphinxupquote{Kconfig}} dosyasına aşağıdaki ekleme
yapılmalıdır:

\begin{sphinxVerbatim}[commandchars=\\\{\}]
source \PYGZdq{}drivers/mydriver/Kconfig\PYGZdq{}
\end{sphinxVerbatim}

\sphinxAtStartPar
\sphinxstylestrong{Adım 5:} \sphinxcode{\sphinxupquote{drivers/mydriver}} dizininde aşağıdaki içeriğe sahip bir \sphinxcode{\sphinxupquote{Makefile}} dosyası
oluşturulmalıdır:

\begin{sphinxVerbatim}[commandchars=\\\{\}]
\PYG{n+nv}{obj\PYGZhy{}\PYGZdl{}(CONFIG\PYGZus{}MYDRIVER)}\PYG{+w}{ }\PYG{o}{+=}\PYG{+w}{ }my\PYGZus{}driver.o
\end{sphinxVerbatim}

\sphinxAtStartPar
\sphinxstylestrong{Adım 6:} Üst dizindeki (\sphinxcode{\sphinxupquote{drivers}} dizinindeki) \sphinxcode{\sphinxupquote{Makefile}} dosyasına aşağıdaki satır
eklenmelidir:

\begin{sphinxVerbatim}[commandchars=\\\{\}]
\PYG{n+nv}{obj\PYGZhy{}\PYGZdl{}(CONFIG\PYGZus{}MYDRIVER)}\PYG{+w}{ }\PYG{o}{+=}\PYG{+w}{ }mydriver/
\end{sphinxVerbatim}

\sphinxAtStartPar
Artık çekirdeği derleyebiliriz. \sphinxstyleemphasis{“make menuconfig”} menüsünde kendi aygıt sürücümüze ilişkin seçenek
de çıkacaktır.

\sphinxAtStartPar
Çekirdek imzalamasını devre dışı bırakmak için konfigürasyon dosyasındaki satırlarda aşağıdaki
değişiklikleri yapmayı unutmayınız:

\begin{sphinxVerbatim}[commandchars=\\\{\}]
CONFIG\PYGZus{}SYSTEM\PYGZus{}TRUSTED\PYGZus{}KEYS=\PYGZdq{}\PYGZdq{}
CONFIG\PYGZus{}SYSTEM\PYGZus{}REVOCATION\PYGZus{}KEYS=\PYGZdq{}\PYGZdq{}
CONFIG\PYGZus{}SYSTEM\PYGZus{}TRUSTED\PYGZus{}KEYRING=n
CONFIG\PYGZus{}SECONDARY\PYGZus{}TRUSTED\PYGZus{}KEYRING=n

CONFIG\PYGZus{}MODULE\PYGZus{}SIG=n
CONFIG\PYGZus{}MODULE\PYGZus{}SIG\PYGZus{}ALL=n
CONFIG\PYGZus{}MODULE\PYGZus{}SIG\PYGZus{}KEY=\PYGZdq{}\PYGZdq{}
\end{sphinxVerbatim}

\sphinxAtStartPar
Yeni çekirdeğimize \sphinxcode{\sphinxupquote{\sphinxhyphen{}custom}} ismini de ekleyebiliriz. Daha önceden de belirttiğimiz gibi eğer
çekirdeğin eski versiyonundan konfigürasyon dosyası alınacaksa \sphinxstyleemphasis{“make oldconfig”} uygulanıp o
versiyondan sonra eklenmiş olan özelliklerin gözden geçirilmesi sağlanmalıdır. Ancak \sphinxstyleemphasis{“make menuconfig”}
işlemi zaten \sphinxstyleemphasis{“make oldconfig”} işlemini de içermektedir.

\sphinxAtStartPar
\sphinxstylestrong{Adım 7:} Artık çekirdek derlemesi aşağıdaki gibi yapılabilir:

\begin{sphinxVerbatim}[commandchars=\\\{\}]
\PYGZdl{}\PYG{+w}{ }make\PYG{+w}{ }\PYGZhy{}j\PYG{k}{\PYGZdl{}(}nproc\PYG{k}{)}
\end{sphinxVerbatim}

\sphinxAtStartPar
\sphinxstylestrong{Adım 8:} Derleme işlemi bittikten sonra önce çekirdek modüllerini \sphinxstyleemphasis{“make modules\_install”} ile
sonra da çekirdeğin kendisini \sphinxstyleemphasis{“make install”} ile kurabilirsiniz:

\begin{sphinxVerbatim}[commandchars=\\\{\}]
\PYGZdl{}\PYG{+w}{ }sudo\PYG{+w}{ }make\PYG{+w}{ }modules\PYGZus{}install
\PYGZdl{}\PYG{+w}{ }sudo\PYG{+w}{ }make\PYG{+w}{ }install
\end{sphinxVerbatim}

\sphinxAtStartPar
Anımsanacağı gibi \sphinxstyleemphasis{“make install”} komutu artık sistemin yeni çekirdekle açılmasını sağlayacaktır.
\sphinxstyleemphasis{“make install”} aynı zamanda geçici kök dosya sistemini \sphinxstyleemphasis{update\sphinxhyphen{}initramfs} komutu ile oluşturup
\sphinxcode{\sphinxupquote{/boot}} dizinine yerleştirmektedir. Tabii \sphinxstyleemphasis{update\sphinxhyphen{}initramfs} programını siz de gerektiğinde
kullanabilirsiniz. Programın tipik kullanımı şöyledir:

\begin{sphinxVerbatim}[commandchars=\\\{\}]
\PYGZdl{}\PYG{+w}{ }sudo\PYG{+w}{ }update\PYGZhy{}initramfs\PYG{+w}{ }\PYGZhy{}c\PYG{+w}{ }\PYGZhy{}k\PYG{+w}{ }\PYGZlt{}çekirdek\PYGZus{}sürümü\PYGZgt{}
\end{sphinxVerbatim}

\sphinxAtStartPar
Buradaki \sphinxstyleemphasis{çekirdek\_sürümü} yalnızca çekirdeğin numarasını değil ona verdiğiniz ekleri de içermelidir.
(Örneğin \sphinxcode{\sphinxupquote{6.9.2\sphinxhyphen{}custom}} gibi.) Bu komut geçici kök dosya sistemini o anda çalışmakta olan sistemin
konfigürasyonunu da dikkate alarak oluşturur ve \sphinxcode{\sphinxupquote{/boot}} dizinine kopyalar. Yukarıda da belirttiğimiz
gibi \sphinxstyleemphasis{“make install”} zaten bu programı çalıştırarak geçici kök dosya sistemini \sphinxcode{\sphinxupquote{/boot}} dizininde
oluşturmaktadır.

\sphinxAtStartPar
Peki çekirdeğin kaynak kodlarına yaptığımız eklemenin gerçekten yapılmış olduğunu nasıl anlayabiliriz?
Bizim yazdığımız iskelet aygıt sürücü kodlarında çekirdek aygıt sürücümüz yüklendiğinde
\sphinxcode{\sphinxupquote{mydriver\_init}} fonksiyonu çağrılacaktır. Bu fonksiyonun içinde de \sphinxstyleemphasis{printk} isimli çekirdek fonksiyonu
ile biz bir log mesajı yazdırdık. Bu log mesajları \sphinxstyleemphasis{kernel ring buffer} denilen bir kuyruk sistemine
yazılmaktadır. \sphinxstyleemphasis{dmesg} komutuyla bu kuyruk sistemi görüntülenebilir. Eğer \sphinxstyleemphasis{dmesg} yaptığımızda biz bu
mesajları görürsek aygıt sürücümüzün yüklenmiş olduğu sonucunu çıkartabiliriz. Örneğin:

\begin{sphinxVerbatim}[commandchars=\\\{\}]
\PYGZdl{}\PYG{+w}{ }dmesg\PYG{+w}{ }\PYG{p}{|}\PYG{+w}{ }grep\PYG{+w}{ }\PYG{l+s+s2}{\PYGZdq{}Hello\PYGZdq{}}
Hello\PYG{+w}{ }World...
\end{sphinxVerbatim}

\sphinxAtStartPar
Çekirdeğin içerisine gömülmüş olan modüller \sphinxcode{\sphinxupquote{/proc/modules}} dosyasında görünmezler, dolayısıyla da
\sphinxstyleemphasis{lsmod} komutu ile de bunları göremeyiz. Bunlar için \sphinxcode{\sphinxupquote{/sys/module}} dizininde de bir giriş
oluşturulmamaktadır. \sphinxstyleemphasis{modinfo} komutu ise çekirdeğe ilişkin bazı dosyalara da baktığı için bize bu
konuda bilgi verebilmektedir.


\section{Yeniden Derleme Sonrası Güncellemeler}
\label{\detokenize{kernelcompilation:yeniden-derleme-sonrasi-guncellemeler}}
\sphinxAtStartPar
Peki çekirdek kodlarında küçük değişiklikler yaptıktan sonra yeniden \sphinxstyleemphasis{“make modules\_install”} ve
\sphinxstyleemphasis{“make install”} işlemlerine gerek var mı? Aslında küçük değişiklikler için bu işlemler yapılmazsa
genellikle bir sorun ortaya çıkmaz. Yeni oluşturulan çekirdek imajı doğrudan eskisinin üzerine
kopyalanabilir. Ancak değişikliğin yerine ve kapsamına göre çekirdeğin sembol tabloları değişebileceği
için genel olarak her derlemeden sonra \sphinxstyleemphasis{“make install”} yapabilirsiniz.

\sphinxAtStartPar
\sphinxcode{\sphinxupquote{drivers}} dizininde \sphinxcode{\sphinxupquote{obj\sphinxhyphen{}m}} biçiminde değişiklikler yapılmışsa \sphinxstyleemphasis{“make modules\_install”} yapılmalıdır.
Yukarıdaki örnekte biz \sphinxcode{\sphinxupquote{drivers}} dizininin içerisine \sphinxcode{\sphinxupquote{obj\sphinxhyphen{}y}} ile eklemeler yaptık; bu durumda aslında
\sphinxstyleemphasis{“make modules\_install”} yapmaya gerek yoktur. Ancak aygıt sürücüler \sphinxcode{\sphinxupquote{obj\sphinxhyphen{}m}} biçiminde ekleniyorsa
\sphinxstyleemphasis{“make modules\_install”} komutu uygulanmalıdır. Çekirdeğin modüllerle ilgili olmayan kısımlarında yapılan
değişiklikler için \sphinxstyleemphasis{“make modules\_install”} yapılmasına gerek olmadığını bir kez daha belirtmek istiyoruz.
\sphinxstyleemphasis{“make modules\_install”} işleminden önce eski \sphinxcode{\sphinxupquote{/lib/modules/\textless{}çekirdek\_sürümü\textgreater{}}} dizinini \sphinxstyleemphasis{“rm \sphinxhyphen{}r”} komutu
ile silmek daha güvenli bir yaklaşımdır.


\section{Çekirdek ve Modül İmzalama}
\label{\detokenize{kernelcompilation:cekirdek-ve-modul-imzalama}}
\sphinxAtStartPar
Biz yukarıdaki çekirdek derlemesi sürecinde imzalama (signing) işlemlerini devre dışı bırakmıştık.
Çekirdek kodları ve özellikle de aygıt sürücüler belli imzalara sahip olacak biçimde derlenebilmektedir.
Böylece onlar üzerinde birtakım istenmeyen değişikliklerin yapılmış olduğu değiştirilmiş çekirdeklerin
ya da aygıt sürücülerin yüklenmesi engellenmiş olur. Yukarıda da gördüğünüz gibi çekirdek kodları ve
aygıt sürücülerde bu imzalama işlemi devre dışı da bırakılabilmektedir. Ancak imzalama süreci sistem
güvenliğini artırmaktadır. Bu tür imzalama işlemleri yalnızca Linux sistemlerinde değil diğer UNIX
türevi sistemlerde, Windows ve macOS sistemlerinde de bulunmaktadır.

\sphinxAtStartPar
Çekirdeğin imza kontrolü temel olarak UEFI BIOS (eğer \sphinxstyleemphasis{secure boot} seçeneği aktif ise) ve önyükleyiciler
(örneğin GRUB gibi, U\sphinxhyphen{}Boot gibi önyükleyiciler) tarafından yapılmaktadır. Ancak Linux çekirdeği de aygıt
sürücüler ve modüller yüklenirken imza kontrolü uygulayabilmektedir. Biz burada önce modül imzalama
işleminin ve sonra da çekirdek imzalama işleminin nasıl yapılacağı üzerinde duracağız.

\sphinxAtStartPar
İmzalama işlemi tipik olarak şu adımlardan geçilerek yapılmaktadır:

\sphinxAtStartPar
\sphinxstylestrong{Adım 1:} İmzalama işlemi için öncelikle \sphinxstyleemphasis{openssl} kütüphanesinin yüklenmiş olması gerekir. Yükleme
işlemi aşağıdaki gibi yapılabilir:

\begin{sphinxVerbatim}[commandchars=\\\{\}]
\PYGZdl{}\PYG{+w}{ }sudo\PYG{+w}{ }apt\PYGZhy{}get\PYG{+w}{ }install\PYG{+w}{ }openssl
\end{sphinxVerbatim}

\sphinxAtStartPar
Daha sonra \sphinxstyleemphasis{openssl} programı ile aşağıdaki gibi bir anahtar çifti ve sertifika dosyası üretilir.
Buradaki \sphinxcode{\sphinxupquote{.key}} dosyası özel anahtarı (private key), \sphinxcode{\sphinxupquote{.crt}} dosyası ise sertifika dosyasını
belirtmektedir. Bu dosyaların uzantıları \sphinxcode{\sphinxupquote{.key}}, \sphinxcode{\sphinxupquote{.crt}}, \sphinxcode{\sphinxupquote{.pem}} olsa da içeriği PEM (Privacy
Enhanced Mail) formatındadır. Dolayısıyla aslında burada verdiğiniz dosyaların uzantısı herhangi bir
biçimde olabilir:

\begin{sphinxVerbatim}[commandchars=\\\{\}]
\PYGZdl{}\PYG{+w}{ }openssl\PYG{+w}{ }req\PYG{+w}{ }\PYGZhy{}new\PYG{+w}{ }\PYGZhy{}x509\PYG{+w}{ }\PYGZhy{}newkey\PYG{+w}{ }rsa:2048\PYG{+w}{ }\PYGZhy{}sha256\PYG{+w}{ }\PYG{l+s+se}{\PYGZbs{}}
\PYG{+w}{    }\PYGZhy{}keyout\PYG{+w}{ }signing\PYGZus{}key.key\PYG{+w}{ }\PYGZhy{}out\PYG{+w}{ }certs/signing\PYGZus{}key.crt\PYG{+w}{ }\PYG{l+s+se}{\PYGZbs{}}
\PYG{+w}{    }\PYGZhy{}nodes\PYG{+w}{ }\PYGZhy{}days\PYG{+w}{ }\PYG{l+m}{36500}\PYG{+w}{ }\PYGZhy{}subj\PYG{+w}{ }\PYG{l+s+s2}{\PYGZdq{}/CN=Local Kernel Module Key/\PYGZdq{}}
\end{sphinxVerbatim}

\sphinxAtStartPar
Bu iki dosyanın kaynak kod ağacına aşağıdaki gibi kopyalanması gerekir:

\begin{sphinxVerbatim}[commandchars=\\\{\}]
\PYGZdl{}\PYG{+w}{ }cp\PYG{+w}{ }signing\PYGZus{}key.key\PYG{+w}{ }certs/
\PYGZdl{}\PYG{+w}{ }cp\PYG{+w}{ }signing\PYGZus{}key.crt\PYG{+w}{ }certs/
\end{sphinxVerbatim}

\sphinxAtStartPar
Daha sonra konfigürasyon dosyasında imzalama için aşağıdaki değişiklikler yapılmalıdır:

\begin{sphinxVerbatim}[commandchars=\\\{\}]
CONFIG\PYGZus{}MODULE\PYGZus{}SIG=y
CONFIG\PYGZus{}MODULE\PYGZus{}SIG\PYGZus{}ALL=y
CONFIG\PYGZus{}MODULE\PYGZus{}SIG\PYGZus{}SHA256=y
CONFIG\PYGZus{}SYSTEM\PYGZus{}TRUSTED\PYGZus{}KEYRING=y
CONFIG\PYGZus{}MODULE\PYGZus{}SIG\PYGZus{}KEY=\PYGZdq{}certs/signing\PYGZus{}key.pem\PYGZdq{}
CONFIG\PYGZus{}SYSTEM\PYGZus{}TRUSTED\PYGZus{}KEYS=\PYGZdq{}certs/signing\PYGZus{}key.crt\PYGZdq{}
\end{sphinxVerbatim}

\sphinxAtStartPar
Aslında \sphinxcode{\sphinxupquote{.key}} ve \sphinxcode{\sphinxupquote{.crt}} dosyaları tek bir dosyada da birleştirilebilir:

\begin{sphinxVerbatim}[commandchars=\\\{\}]
\PYGZdl{}\PYG{+w}{ }cat\PYG{+w}{ }certs/signing\PYGZus{}key.key\PYG{+w}{ }certs/signing\PYGZus{}key.crt\PYG{+w}{ }\PYGZgt{}\PYG{+w}{ }certs/signing\PYGZus{}key.pem
\end{sphinxVerbatim}

\sphinxAtStartPar
Tabii artık \sphinxcode{\sphinxupquote{.config}} dosyasındaki isimleri de şöyle değiştirmeliyiz:

\begin{sphinxVerbatim}[commandchars=\\\{\}]
CONFIG\PYGZus{}MODULE\PYGZus{}SIG=y
CONFIG\PYGZus{}MODULE\PYGZus{}SIG\PYGZus{}ALL=y
CONFIG\PYGZus{}MODULE\PYGZus{}SIG\PYGZus{}SHA256=y
CONFIG\PYGZus{}SYSTEM\PYGZus{}TRUSTED\PYGZus{}KEYRING=y
CONFIG\PYGZus{}MODULE\PYGZus{}SIG\PYGZus{}KEY=\PYGZdq{}certs/signing\PYGZus{}key.pem\PYGZdq{}
CONFIG\PYGZus{}SYSTEM\PYGZus{}TRUSTED\PYGZus{}KEYS=\PYGZdq{}certs/signing\PYGZus{}key.pem\PYGZdq{}
\end{sphinxVerbatim}

\sphinxAtStartPar
Biz bu işlemlerle aygıt sürücüleri imzalamış olduk. Ancak eğer \sphinxcode{\sphinxupquote{.config}} dosyasında
\sphinxcode{\sphinxupquote{CONFIG\_MODULE\_SIG\_FORCE=n}} ise (default durumda genellikle böyledir) bu durum çekirdek log amaçlı
bir uyarı oluştursa da imzalanmamış modülleri yine de yükler. Eğer imzalanmamış modülleri yüklemek
istemiyorsanız \sphinxcode{\sphinxupquote{CONFIG\_MODULE\_SIG\_FORCE=y}} yapmalısınız. (Bu işlem \sphinxstyleemphasis{“make menuconfig”} menüsünde
\sphinxstyleemphasis{Enable loadable module support / Module signature verification / Require modules to be validly signed}
seçeneğinden de yapılabilmektedir.) Bu durumda çekirdek kendi oluşturduğumuz imzayla imzalanmamış olan
modülleri artık yüklemeyecektir.

\sphinxAtStartPar
Eğer biz başkalarının yazdığı bir aygıt sürücüyü yüklerken çekirdeğin uyarı vermesini ya da yüklemeyi
reddetmesini istemiyorsak o aygıt sürücüyü de kendi ürettiğimiz anahtarla (ya da dağıtımın public
anahtarıyla) imzalamalıyız. Bu işlem çekirdek kodlarındaki \sphinxcode{\sphinxupquote{scripts}} dizini içerisinde bulunan
\sphinxstyleemphasis{sign\sphinxhyphen{}file} betiği ile yapılmaktadır. Bu betiğin tipik kullanımı şöyledir:

\begin{sphinxVerbatim}[commandchars=\\\{\}]
\PYGZdl{}\PYG{+w}{ }scripts/sign\PYGZhy{}file\PYG{+w}{ }\PYGZlt{}hash\PYGZus{}alg\PYGZgt{}\PYG{+w}{ }\PYGZlt{}private\PYGZus{}key.pem\PYGZgt{}\PYG{+w}{ }\PYGZlt{}public\PYGZus{}cert.pem\PYGZgt{}\PYG{+w}{ }\PYGZlt{}module.ko\PYGZgt{}
\end{sphinxVerbatim}

\sphinxAtStartPar
Örneğin:

\begin{sphinxVerbatim}[commandchars=\\\{\}]
\PYGZdl{}\PYG{+w}{ }scripts/sign\PYGZhy{}file\PYG{+w}{ }sha256\PYG{+w}{ }signing\PYGZus{}key.pem\PYG{+w}{ }signing\PYGZus{}key.pem\PYG{+w}{ }mydriver.ko
\end{sphinxVerbatim}

\sphinxAtStartPar
Peki Ubuntu, Mint gibi dağıtımlar çekirdek imzalaması uygulamakta mıdır? Evet, genel olarak dağıtımlar
kendi özel anahtarlarıyla (private keys) çekirdeği ve aygıt sürücüleri imzalamaktadır. Ancak aygıt
sürücülerin yüklenmesinde imza kontrolünü zorunlu hale getirmemektedir. İmzalama için kullanılacak
public anahtarlar \sphinxcode{\sphinxupquote{/proc/keys}} dosyasında belirtilmektedir.


\section{Çekirdek İmjasının İmzalanması}
\label{\detokenize{kernelcompilation:cekirdek-imjasinin-imzalanmasi}}
\sphinxAtStartPar
Biz yukarıdaki işlemleri yaptığımızda yalnızca aygıt sürücüleri imzalamış oluruz. Çekirdeğin kendisinin
imzalanması ayrıca yapılmalıdır. Yukarıda da belirttiğimiz gibi UEFI BIOS’lar ve GRUB gibi önyükleyiciler
çekirdek imajını yüklemeden önce ayarlar uygun biçime getirildiyse çekirdek imzasına bakmaktadır. Eğer
çekirdek imzası yanlışsa (çekirdek dışarıdan kasti ya da yanlışlıkla bozulmuş olabilir) çekirdeği hiç
yüklememektedir.

\sphinxAtStartPar
Çekirdeğin imzalanması için önce anahtar ve sertifikasyon dosyaları aşağıdaki gibi oluşturulur:

\begin{sphinxVerbatim}[commandchars=\\\{\}]
\PYGZdl{}\PYG{+w}{ }openssl\PYG{+w}{ }req\PYG{+w}{ }\PYGZhy{}new\PYG{+w}{ }\PYGZhy{}x509\PYG{+w}{ }\PYGZhy{}newkey\PYG{+w}{ }rsa:2048\PYG{+w}{ }\PYGZhy{}sha256\PYG{+w}{ }\PYG{l+s+se}{\PYGZbs{}}
\PYG{+w}{    }\PYGZhy{}keyout\PYG{+w}{ }certs/sb\PYGZhy{}signing.key\PYG{+w}{ }\PYG{l+s+se}{\PYGZbs{}}
\PYG{+w}{    }\PYGZhy{}out\PYG{+w}{ }certs/sb\PYGZhy{}signing.crt\PYG{+w}{ }\PYG{l+s+se}{\PYGZbs{}}
\PYG{+w}{    }\PYGZhy{}nodes\PYG{+w}{ }\PYGZhy{}days\PYG{+w}{ }\PYG{l+m}{36500}\PYG{+w}{ }\PYG{l+s+se}{\PYGZbs{}}
\PYG{+w}{    }\PYGZhy{}subj\PYG{+w}{ }\PYG{l+s+s2}{\PYGZdq{}/CN=Secure Boot Signing Key/\PYGZdq{}}
\end{sphinxVerbatim}

\sphinxAtStartPar
Sonra da \sphinxstyleemphasis{sbsign} programı ile imzalama aşağıdaki gibi yapılır:

\begin{sphinxVerbatim}[commandchars=\\\{\}]
\PYGZdl{}\PYG{+w}{ }sbsign\PYG{+w}{ }\PYGZhy{}\PYGZhy{}key\PYG{+w}{ }db.key\PYG{+w}{ }\PYGZhy{}\PYGZhy{}cert\PYG{+w}{ }signing\PYGZus{}key.pem\PYG{+w}{ }\PYGZhy{}\PYGZhy{}output\PYG{+w}{ }bzImage.signed\PYG{+w}{ }arch/x86/boot/bzImage
\end{sphinxVerbatim}

\sphinxAtStartPar
Eğer bu işlemden sonra \sphinxstyleemphasis{“make install”} yapacaksanız imzalanmış çekirdeği eski ismiyle bulundurmalısınız
(\sphinxstyleemphasis{mv} komutu hedef dosya varsa onu ezerek işlemini yapmaktadır):

\begin{sphinxVerbatim}[commandchars=\\\{\}]
\PYGZdl{}\PYG{+w}{ }mv\PYG{+w}{ }arch/x86/boot/bzImage.signed\PYG{+w}{ }arch/x86/boot/bzImage
\end{sphinxVerbatim}

\sphinxAtStartPar
Genellikle bu biçimde bir çekirdek imzalaması seyrek olarak yapılmaktadır. Önceki paragrafta biz aygıt
sürücü dosyalarının imzalandığını belirtmiştik. Aynı makinede aygıt sürücüyü derlerken (build ederken)
oluşturulan imza bilgisi de kullanılmaktadır. Yani biz aynı makinede bir aygıt sürücü derlediğimizde
aygıt sürücümüz de zaten imzalanmış olacaktır.


\section{Kök Dosya Sisteminin Oluşturulması}
\label{\detokenize{kernelcompilation:kok-dosya-sisteminin-olusturulmasi}}
\sphinxAtStartPar
Bir Linux sisteminin düzgün bir biçimde açılması için belli dizinlerin kök dosya sisteminde bulunuyor
olması gerekir. Biz bir dağıtımı kurduğumuzda zaten bu kök dosya sistemi de oluşturulmaktadır. Peki
sıfırdan dağıtımı tamamen kurmadan kök dosya sistemini nasıl oluşturulabiliriz? Bu işlem tamamen manuel
biçimde yapılabilir. Yani uygulamacı kök dizin içerisindeki gerekli dizinleri elle yaratır. Sonra
gerekli programları kaynak kodlarından hareketle hedef makine için derler ve onları konuşlandırır. Sonra
yine gerekli birtakım konfigürasyon dosyalarını elle oluşturur. Ancak bu manuel yöntem oldukça zahmetlidir.

\sphinxAtStartPar
Bunun yerine bu işlemi pratik bir biçimde yapan araçlar geliştirilmiştir. Örneğin gömülü sistemlerde
\sphinxstyleemphasis{BusyBox} denilen araç bu amaçla sıkça kullanılmaktadır. Kullanımı da oldukça kolaydır. Gömülü sistemler
için \sphinxstyleemphasis{Buildroot} ve \sphinxstyleemphasis{Yocto} gibi projeler daha genel amaçlar için gerçekleştirilmiştir ancak bunlarla kök
dosya sistemi de oluşturulabilmektedir. Bazı dağıtımların bu işi yapan özel yardımcı programları da vardır.
Örneğin \sphinxstyleemphasis{debootstrap} programı Debian tabanlı kök dosya sistemini İnternet’ten indirerek yerel makinede
oluşturabilmektedir. Ancak bu araçların bazıları esnek değildir. Özellikle gömülü sistemlerde düşük bir
sistem kaynağının olduğu dikkate alındığında bu araçların bazıları minimalist bir kurulum
sağlayamamaktadır.


\subsection{\sphinxstyleliteralintitle{\sphinxupquote{debootstrap}} ile Kök Dosya Sistemi Oluşturma}
\label{\detokenize{kernelcompilation:debootstrap-ile-kok-dosya-sistemi-olusturma}}
\sphinxAtStartPar
\sphinxstyleemphasis{debootstrap} programı default olarak sisteminizde yüklü değildir. Bunu aşağıdaki gibi kurabilirsiniz:

\begin{sphinxVerbatim}[commandchars=\\\{\}]
\PYGZdl{}\PYG{+w}{ }sudo\PYG{+w}{ }apt\PYGZhy{}get\PYG{+w}{ }install\PYG{+w}{ }debootstrap
\end{sphinxVerbatim}

\sphinxAtStartPar
\sphinxstyleemphasis{debootstrap} programının pek çok komut satırı argümanı vardır. Biz burada en önemli birkaç argüman
üzerinde duracağız:
\begin{itemize}
\item {} 
\sphinxAtStartPar
\sphinxcode{\sphinxupquote{\sphinxhyphen{}\sphinxhyphen{}arch}}: Hedef CPU mimarisini belirtmektedir. Bu argüman girilmezse o andaki platform temel alınır.
64 bit Intel platformu için \sphinxcode{\sphinxupquote{amd64}}, BBB gibi 32 bit ARM platformu için \sphinxcode{\sphinxupquote{armhf}}, 64 bit ARM platformu
için \sphinxcode{\sphinxupquote{arm64}} girilmelidir.

\item {} 
\sphinxAtStartPar
İlk seçeneksiz argüman Debian sisteminin varyantını belirtmektedir (örneğin \sphinxcode{\sphinxupquote{bullseye}}, \sphinxcode{\sphinxupquote{buster}}).

\item {} 
\sphinxAtStartPar
İkinci komut satırı argümanı hedef kök dosya sisteminin oluşturulacağı dizini belirtmektedir.

\item {} 
\sphinxAtStartPar
Üçüncü komut satırı argümanı ise paketlerin indirileceği depoyu (repository) belirtmektedir.

\end{itemize}

\sphinxAtStartPar
Örneğin:

\begin{sphinxVerbatim}[commandchars=\\\{\}]
\PYGZdl{}\PYG{+w}{ }sudo\PYG{+w}{ }debootstrap\PYG{+w}{ }\PYGZhy{}\PYGZhy{}arch\PYG{o}{=}amd64\PYG{+w}{ }\PYGZhy{}\PYGZhy{}include\PYG{o}{=}systemd\PYG{+w}{ }bullseye\PYG{+w}{ }myrootfs\PYG{+w}{ }http://deb.debian.org/debian/
\end{sphinxVerbatim}

\sphinxAtStartPar
\sphinxcode{\sphinxupquote{\sphinxhyphen{}\sphinxhyphen{}arch}} seçeneği girilmemişse programın çalıştırıldığı makine için kök dosya sistemi indirilip
kurulmaktadır. Default durumda \sphinxstyleemphasis{debootstrap} pek çok paketi kök dosya sistemine dahil ettiği için
paketlerin indirilmesi ve kök dosya sisteminin oluşturulması biraz zaman alacaktır.

\sphinxAtStartPar
Uygulamacı isterse \sphinxcode{\sphinxupquote{\sphinxhyphen{}\sphinxhyphen{}include}} ve \sphinxcode{\sphinxupquote{\sphinxhyphen{}\sphinxhyphen{}exclude}} komut satırı seçenekleriyle birtakım paketleri dahil
edebilir ya da dışlayabilir. Örneğin biz \sphinxstyleemphasis{systemd} dışında \sphinxstyleemphasis{sudo} ve \sphinxstyleemphasis{gcc} paketlerini de aşağıdaki gibi
kuruluma dahil edebiliriz:

\begin{sphinxVerbatim}[commandchars=\\\{\}]
\PYGZdl{}\PYG{+w}{ }sudo\PYG{+w}{ }debootstrap\PYG{+w}{ }\PYGZhy{}\PYGZhy{}arch\PYG{o}{=}amd64\PYG{+w}{ }\PYGZhy{}\PYGZhy{}include\PYG{o}{=}systemd,sudo,gcc\PYG{+w}{ }bullseye\PYG{+w}{ }myrootfs\PYG{+w}{ }http://deb.debian.org/debian/
\end{sphinxVerbatim}

\sphinxAtStartPar
\sphinxstyleemphasis{debootstrap} programı ile eğer host makineyle aynı platform için indirme işlemi yapılıyorsa \sphinxstyleemphasis{debootstrap}
önce İnternet’ten gerekli paketleri indirip yerel makinede bir dizin içerisinde kök dosya sistemini
oluşturur. Eğer host makineden farklı bir sistem için indirme yapılıyorsa (örneğin host sistem Intel
tabanlı bir makineyse ve ARM tabanlı bir Debian kök dosya sistemi oluşturulmak isteniyorsa) \sphinxstyleemphasis{debootstrap}
yüklemesini yapmadan önce aşağıdaki gibi \sphinxstyleemphasis{qemu} emülatör paketinin statik versiyonu ve \sphinxstyleemphasis{binfmt} destek
paketi kurulmalıdır:

\begin{sphinxVerbatim}[commandchars=\\\{\}]
\PYGZdl{}\PYG{+w}{ }sudo\PYG{+w}{ }apt\PYG{+w}{ }install\PYG{+w}{ }qemu\PYGZhy{}user\PYGZhy{}static\PYG{+w}{ }binfmt\PYGZhy{}support
\end{sphinxVerbatim}

\sphinxAtStartPar
Bu işlemden sonra \sphinxstyleemphasis{debootstrap} programı yukarıda belirttiğimiz biçimde çalıştırılabilir:

\begin{sphinxVerbatim}[commandchars=\\\{\}]
\PYGZdl{}\PYG{+w}{ }sudo\PYG{+w}{ }debootstrap\PYG{+w}{ }\PYGZhy{}\PYGZhy{}include\PYG{o}{=}systemd\PYG{+w}{ }\PYGZhy{}\PYGZhy{}arch\PYG{+w}{ }armhf\PYG{+w}{ }buster\PYG{+w}{ }myrootfs\PYG{+w}{ }http://deb.debian.org/debian/
\end{sphinxVerbatim}

\sphinxAtStartPar
Bu komut hem birinci aşama hem de ikinci aşama işlemleri yapıp bitirecektir. Artık istenildiği zaman
\sphinxstyleemphasis{chroot} işlemi de yapılabilir:

\begin{sphinxVerbatim}[commandchars=\\\{\}]
\PYGZdl{}\PYG{+w}{ }sudo\PYG{+w}{ }chroot\PYG{+w}{ }myrootfs
\end{sphinxVerbatim}


\subsection{Geçici Kök Dosya Sisteminin Oluşturulması}
\label{\detokenize{kernelcompilation:gecici-kok-dosya-sisteminin-olusturulmasi}}
\sphinxAtStartPar
\sphinxstyleemphasis{debootstrap} programı ile biz Debian kök dosya sistemi için geçici kök dosya sistemi de oluşturabiliriz.
Bunun en pratik yolu kök dosya sistemini kurduktan sonra \sphinxstyleemphasis{chroot} yapıp \sphinxstyleemphasis{update\sphinxhyphen{}initramfs} programı ile
geçici kök dosya sistemini oluşturmaktır. Ancak bunun için \sphinxcode{\sphinxupquote{/boot}} ve \sphinxcode{\sphinxupquote{/lib/modules}} dizinlerinin
uygun biçimde oluşturulmuş olması gerekir. \sphinxstyleemphasis{update\sphinxhyphen{}initramfs} programı bu dizinlerdeki içerikten
faydalanmaktadır. \sphinxstyleemphasis{update\sphinxhyphen{}initramfs} programı \sphinxstyleemphasis{initramfs\sphinxhyphen{}tools} isimli pakettedir. \sphinxstyleemphasis{chroot} yaptıktan
sonra öncelikle bu paketi aşağıdaki gibi kurmalısınız:

\begin{sphinxVerbatim}[commandchars=\\\{\}]
\PYGZdl{}\PYG{+w}{ }sudo\PYG{+w}{ }apt\PYGZhy{}get\PYG{+w}{ }install\PYG{+w}{ }initramfs\PYGZhy{}tools
\end{sphinxVerbatim}

\sphinxAtStartPar
Bundan sonra geçici kök dosya sistemini aşağıdaki gibi oluşturabilirsiniz (Debian kök dosya sisteminin
kökünde olduğumuzu varsayıyoruz):

\begin{sphinxVerbatim}[commandchars=\\\{\}]
\PYGZdl{}\PYG{+w}{ }update\PYGZhy{}initramfs\PYG{+w}{ }\PYGZhy{}c\PYG{+w}{ }\PYGZhy{}k\PYG{+w}{ }\PYG{l+m}{6}.9.2\PYGZhy{}custom\PYG{+w}{ }\PYGZhy{}b\PYG{+w}{ }.
\end{sphinxVerbatim}

\sphinxAtStartPar
Burada \sphinxcode{\sphinxupquote{6.9.2\sphinxhyphen{}custom}} çekirdeğin sürüm ismidir. Geçici kök dosya sistemi \sphinxcode{\sphinxupquote{initrd.img\sphinxhyphen{}6.9.2\sphinxhyphen{}custom}}
ismiyle bulunulan dizinde oluşturulacaktır. \sphinxcode{\sphinxupquote{\sphinxhyphen{}b .}} seçeneği oluşturulacak dosyanın dizinini
belirtmektedir.

\begin{sphinxadmonition}{note}{Not:}
\sphinxAtStartPar
Bu tür konuların ayrıntılarına bu kursta girmeyeceğiz. Bu konular daha çok \sphinxstyleemphasis{Gömülü Linux Sistemleri \sphinxhyphen{}
Geliştirme ve Uygulama} kursunun konularını oluşturmaktadır.
\end{sphinxadmonition}


\section{Çekirdek Başlık Dosyalarının Kurulumu}
\label{\detokenize{kernelcompilation:cekirdek-baslik-dosyalarinin-kurulumu}}
\sphinxAtStartPar
Aygıt sürücü geliştirirken çekirdeğin kaynak kodlarına gereksinim duyulmaz. Ancak çekirdeğin başlık
dosyalarının geliştirmenin yapıldığı bilgisayarda yüklü olması gerekir. Çekirdek kodlarının kendisini
değil de yalnızca başlık dosyalarını indirmek için aşağıdaki komut kullanılabilir:

\begin{sphinxVerbatim}[commandchars=\\\{\}]
\PYGZdl{}\PYG{+w}{ }sudo\PYG{+w}{ }apt\PYG{+w}{ }install\PYG{+w}{ }linux\PYGZhy{}headers\PYGZhy{}\PYG{k}{\PYGZdl{}(}uname\PYG{+w}{ }\PYGZhy{}r\PYG{k}{)}
\end{sphinxVerbatim}

\sphinxAtStartPar
Buradaki \sphinxcode{\sphinxupquote{\$(uname \sphinxhyphen{}r)}} çalışılmakta olan makinedeki çekirdek sürümünü belirtmektedir. Tabii biz
istediğimiz çekirdek sürümünün başlık dosyalarını da indirebiliriz. İndirilen dosyalar \sphinxcode{\sphinxupquote{/usr/src}}
dizininin altına \sphinxcode{\sphinxupquote{\$(uname \sphinxhyphen{}r)}} isimli dizin içerisine yerleştirilmektedir.

\sphinxstepscope


\chapter{Bağlı Listelerin ve Hash Tablolarının Çekirdek Gerçekleştirimleri}
\label{\detokenize{kernellinkedlistandhashtable:bagli-listelerin-ve-hash-tablolarinin-cekirdek-gerceklestirimleri}}\label{\detokenize{kernellinkedlistandhashtable::doc}}
\sphinxAtStartPar
Bu bölümde Linux çekirdeğindeki bağlı listelerin ve hash tablolarının gerçekleştirimleri üzerinde duracağız.


\section{Bağlı Listeler}
\label{\detokenize{kernellinkedlistandhashtable:bagli-listeler}}
\sphinxAtStartPar
Çekirdeğin en önemli veri yapılarından biri bağlı listelerdir. Genel olarak çekirdekteki neredeyse tüm
bağlı listeler \sphinxstyleemphasis{çift bağlı (doubly linked)} biçimde kullanılmaktadır. Önce bağlı listelerin Linux
çekirdeğindeki gerçekleştirimleri üzerinde duracağız.

\sphinxAtStartPar
Aralarında öncelik\sphinxhyphen{}sonralık ilişkisi olan veri yapılarına \sphinxstyleemphasis{liste (list)} tarzı veri yapıları denilmektedir.
Örneğin bu tanıma göre diziler de “liste tarzı” veri yapılarıdır. Liste tarzı veri yapılarının en yaygın
kullanılanlarından biri \sphinxstyleemphasis{bağlı liste (linked list)} denilen veri yapısıdır. Önceki elemanın sonraki
elemanın yerini gösterdiği dolayısıyla elemanların ardışıl olma zorunluluğunun ortadan kaldırıldığı
listelere “bağlı liste” denilmektedir. Dizi elemanlarının bellekte fiziksel olarak ardışıl biçimde
bulunduğunu anımsayınız. Bağlı listeler adeta “elemanları bellekte ardışıl olmak zorunda olmayan
diziler” gibidir.

\sphinxAtStartPar
Bağlı listelerin her elemanına \sphinxstyleemphasis{düğüm (node)} denilmektedir. Bağlı listelerde her düğüm sonraki
düğümün yerini tuttuğuna göre ilk elemanın yeri biliniyorsa liste elemanlarının hepsine erişilebilmektedir:

\sphinxincludegraphics[]{graphviz-17a828b3ed6668e1720b754b4350973a641c94c8.pdf}

\sphinxAtStartPar
Her düğümün yalnızca sonraki düğümün değil aynı zamanda önceki düğümün yerini de tuttuğu bağlı
listelere \sphinxstyleemphasis{çift bağlı listeler (doubly linked lists)} denilmektedir. Çift bağlı listelerde belli bir
düğümün adresini biliyorsak yalnızca ileriye doğru değil, geriye doğru da gidebiliriz:

\sphinxincludegraphics[]{graphviz-eae5046ec12b575d8f0bc91fffe7f7f1f965785f.pdf}

\sphinxAtStartPar
Çift bağlı listelere ilişkin bir düğümün bellekte daha fazla yer kaplayacağına dikkat ediniz. Çift bağlı
listelerin tek bağlı listelere göre en önemli özelliği “adresi bilinen bir düğümün” silinebilmesidir.
Tek bağlı listelerde bu durum mümkün değildir. Uygulamalarda ve özellikle çekirdek kodlarında buna çok
sık gereksinim duyulmaktadır.

\sphinxAtStartPar
Eğer bir bağlı listede son eleman da ilk elemanı gösteriyorsa bu tür bağlı listelere \sphinxstyleemphasis{döngüsel bağlı
listeler (circular linked lists)} denilmektedir.


\section{Bağlı Listelere Neden Gereksinim Duyulur?}
\label{\detokenize{kernellinkedlistandhashtable:bagli-listelere-neden-gereksinim-duyulur}}
\sphinxAtStartPar
Peki bağlı listelere neden gereksinim duyulmaktadır? Diziler varken bağlı listelere gerek var mıdır?
Dizilerle bağlı listeler arasındaki farklılıkları, benzerlikleri ve bağlı listelere neden gereksinim
duyulduğunu birkaç maddede açıklayabiliriz:
\begin{enumerate}
\sphinxsetlistlabels{\arabic}{enumi}{enumii}{}{.}%
\item {} 
\sphinxAtStartPar
\sphinxstylestrong{Bellek parçalanması sorunu:} Diziler ardışıl alana gereksinim duymaktadır. Ancak belleğin
bölündüğü (fragmente olduğu) durumlarda bellekte yeteri kadar küçük boş alanlar olduğu halde bunlar
ardışıl olmadığı için dizi tahsisatı mümkün olamamaktadır. Bu tür durumlarda ardışıllık gereksinimi
olmayan bağlı listeler kullanılabilir. Özellikle heap gibi bir alanda çok sayıda dinamik dizi bellek
kullanımı bakımından verimsizliğe yol açabilmektedir. Bu dinamik diziler zamanla büyüdükçe birbirini
engeller hale gelebilmektedir. İşte uzunluğu baştan belli olmayan çok sayıda dizinin oluşturulacağı
durumlarda dinamik diziler yerine bağlı listeler bellek kullanımı bakımından toplamda daha iyi
performans gösterebilmektedir. Dinamik dizilerde büyütme sırasında bloklar yer değiştirebileceği için
bu işlem yavaştır.

\item {} 
\sphinxAtStartPar
\sphinxstylestrong{Araya ekleme ve silme verimliliği:} Dizilerde araya eleman ekleme (insert etme) ve aradaki bir
elemanı silme dizinin kaydırılmasına yol açacağından yavaş bir işlemdir. Teknik olarak dizilerde eleman
insert etme ve eleman silme O(N) karmaşıklıkta bir işlemdir. Halbuki bağlı listelerde eğer düğümün yeri
biliniyorsa bu işlem O(1) karmaşıklıkta (yani döngü olmadan tekil işlemlerle) yapılabilmektedir. O halde
araya eleman eklemenin ve aradan eleman silmenin çok yapıldığı sistemlerde diziler yerine bağlı listeler
tercih edilebilmektedir. Çekirdek veri yapılarında araya eleman ekleme ve aradan eleman silme gibi
işlemler çok yoğun yapılmaktadır.

\item {} 
\sphinxAtStartPar
\sphinxstylestrong{İndeksle erişim yavaşlığı:} Bağlı listelerde belli bir indeksteki elemana erişmek O(N) karmaşıklıkta
bir işlemdir. Halbuki dizilerde elemana erişim O(1) karmaşıklıkta yani çok hızlıdır. O halde belli bir
indeks değeri ile elemana erişimin yoğun yapıldığı durumlarda bağlı listeler yerine diziler tercih
edilmelidir.

\item {} 
\sphinxAtStartPar
\sphinxstylestrong{Ek bellek maliyeti:} Bağlı listeler toplamda bellekte daha fazla yer kaplama eğilimindedir. Çünkü
bağlı listenin her düğümü sonraki (ve duruma göre önceki) elemanın yerini de tutmaktadır.

\end{enumerate}

\sphinxAtStartPar
O halde bağlı listeler tipik olarak şu durumlarda dizilere tercih edilmelidir:
\begin{itemize}
\item {} 
\sphinxAtStartPar
Eleman insert etmenin ve eleman silmenin sık yapıldığı durumlarda.

\item {} 
\sphinxAtStartPar
Uzunluğu baştan belli olmayan çok sayıda dizinin kullanıldığı durumlarda.

\item {} 
\sphinxAtStartPar
İndeks yoluyla erişimin az yapıldığı durumlarda.

\item {} 
\sphinxAtStartPar
Toplam bellek miktarının yeteri kadar fazla olduğu sistemlerde.

\end{itemize}

\sphinxAtStartPar
İşte tüm yukarıdaki nedenlerden dolayı bağlı listeler çekirdek için en önemli veri yapılarından biridir.


\section{Linux Çekirdeğindeki Bağlı Liste Gerçekleştirimi}
\label{\detokenize{kernellinkedlistandhashtable:linux-cekirdegindeki-bagli-liste-gerceklestirimi}}
\sphinxAtStartPar
Linux çekirdeğinde bağlı liste gerçekleştirimi \sphinxcode{\sphinxupquote{include/linux/list.h}} dosyasında yapılmıştır. Bu dosya
içerisindeki fonksiyonlar \sphinxcode{\sphinxupquote{static inline}} biçiminde tanımlanmıştır. Çekirdek derlemesi sırasında
derleyicinin optimizasyon seçeneği ayarlandığı için derleyici buradaki inline fonksiyonları sanki bir makro
gibi koda açmaktadır.

\sphinxAtStartPar
Linux çekirdeğindeki bağlı listelerde düğümler \sphinxcode{\sphinxupquote{list\_head}} isimli bir yapıyla temsil edilmektedir:

\begin{sphinxVerbatim}[commandchars=\\\{\}]
\PYG{k}{struct}\PYG{+w}{ }\PYG{n+nc}{list\PYGZus{}head}\PYG{+w}{ }\PYG{p}{\PYGZob{}}
\PYG{+w}{    }\PYG{k}{struct}\PYG{+w}{ }\PYG{n+nc}{list\PYGZus{}head}\PYG{+w}{ }\PYG{o}{*}\PYG{n}{next}\PYG{p}{;}
\PYG{+w}{    }\PYG{k}{struct}\PYG{+w}{ }\PYG{n+nc}{list\PYGZus{}head}\PYG{+w}{ }\PYG{o}{*}\PYG{n}{prev}\PYG{p}{;}
\PYG{p}{\PYGZcb{}}\PYG{p}{;}
\end{sphinxVerbatim}

\sphinxAtStartPar
Aslında belli yapılar değil, bu \sphinxcode{\sphinxupquote{list\_head}} yapıları bağlı liste içerisinde birbirine bağlanmaktadır:

\sphinxincludegraphics[]{graphviz-410375c161ca832e3bfd80174451cf3f3d4379af.pdf}

\sphinxAtStartPar
Tabii eğer bu \sphinxcode{\sphinxupquote{list\_head}} yapıları başka bir yapının (buna asıl yapı diyelim) içerisindeyse bu durumda
aslında bir \sphinxcode{\sphinxupquote{list\_head}} yapısının adresi asıl yapının bir elemanının adresi haline gelmektedir. Biz C’de
bir yapının bir elemanının adresini biliyorsak kolaylıkla o yapının başlangıç adresini elde edebiliriz.
Çünkü yapı elemanları ardışıldır ve standart \sphinxcode{\sphinxupquote{offsetof}} makrosuyla belli bir yapı elemanının yapının
başlangıcından itibaren hangi offset’te bulunduğu bilgisi elde edilebilmektedir. \sphinxcode{\sphinxupquote{offsetof}} makrosu
\sphinxcode{\sphinxupquote{\textless{}stddef.h\textgreater{}}} içerisinde aşağıdakine benzer biçimde tanımlanmıştır:

\begin{sphinxVerbatim}[commandchars=\\\{\}]
\PYG{c+cp}{\PYGZsh{}}\PYG{c+cp}{define offsetof(type, member) ((size\PYGZus{}t)\PYGZam{}(((type *)0)\PYGZhy{}\PYGZgt{}member))}
\end{sphinxVerbatim}

\sphinxAtStartPar
\sphinxcode{\sphinxupquote{offsetof}} makrosunun birinci parametresi yapının tür ismini, ikinci parametresi ilgili yapı elemanının
ismini almaktadır. Tabii \sphinxcode{\sphinxupquote{offsetof}} makrosu yalnızca \sphinxcode{\sphinxupquote{size\_t}} türünden bir byte offset’i vermektedir.
Elemanın adresinin makroyla elde edilen bu değerden çıkartılarak tür dönüştürmesinin yapılması gerekir.
İşte bunun için Linux çekirdeğinde \sphinxcode{\sphinxupquote{container\_of}} isimli bir makro bulundurulmuştur. Bu makronun basit
yazımı şöyle yapılabilir:

\begin{sphinxVerbatim}[commandchars=\\\{\}]
\PYG{c+cp}{\PYGZsh{}}\PYG{c+cp}{define container\PYGZus{}of(ptr, type, member) ((type *)((char *)ptr \PYGZhy{} myoffsetof(type, member)))}
\end{sphinxVerbatim}

\sphinxAtStartPar
\sphinxcode{\sphinxupquote{container\_of}} makrosunun birinci parametresi yapı elemanının adresini, ikinci parametresi yapının tür
ismini ve üçüncü parametresi de ilgili yapı elemanının ismini almaktadır. Bu makro ikinci parametresine
girilmiş olan yapı türünden bir adres vermektedir.

\sphinxAtStartPar
\sphinxcode{\sphinxupquote{container\_of}} makrosunu \sphinxcode{\sphinxupquote{list\_entry}} ismiyle de kullanabilirsiniz:

\begin{sphinxVerbatim}[commandchars=\\\{\}]
\PYG{c+cp}{\PYGZsh{}}\PYG{c+cp}{define list\PYGZus{}entry(ptr, type, member) container\PYGZus{}of(ptr, type, member)}
\end{sphinxVerbatim}


\subsection{Bağlı Liste Başlatma}
\label{\detokenize{kernellinkedlistandhashtable:bagli-liste-baslatma}}
\sphinxAtStartPar
Linux çekirdeğindeki bağlı listeler kullanılırken önce bir başlangıç düğümünün oluşturulması gerekir.
Başlangıç düğümünde \sphinxcode{\sphinxupquote{next}} ve \sphinxcode{\sphinxupquote{prev}} göstericilerinin başlangıç düğümünün kendisini göstermesi
gerekir. Örneğin:

\begin{sphinxVerbatim}[commandchars=\\\{\}]
\PYG{k}{struct}\PYG{+w}{ }\PYG{n+nc}{list\PYGZus{}head}\PYG{+w}{ }\PYG{n}{head}\PYG{+w}{ }\PYG{o}{=}\PYG{+w}{ }\PYG{p}{\PYGZob{}}\PYG{o}{\PYGZam{}}\PYG{n}{head}\PYG{p}{,}\PYG{+w}{ }\PYG{o}{\PYGZam{}}\PYG{n}{head}\PYG{p}{\PYGZcb{}}\PYG{p}{;}
\end{sphinxVerbatim}

\sphinxAtStartPar
Bu ilkdeğer verme C’de geçerlidir. Çünkü C’de bir değişken faaliyet alanına \sphinxcode{\sphinxupquote{=}} atomu ile ilkdeğer
verme kısmından önce sokulmaktadır. \sphinxcode{\sphinxupquote{list.h}} dosyasında bu işlemi yapan aşağıdaki gibi bir makro da
bulundurulmuştur:

\begin{sphinxVerbatim}[commandchars=\\\{\}]
\PYG{c+cp}{\PYGZsh{}}\PYG{c+cp}{define LIST\PYGZus{}HEAD\PYGZus{}INIT(name) \PYGZob{}\PYGZam{}(name), \PYGZam{}(name)\PYGZcb{}}
\end{sphinxVerbatim}

\sphinxAtStartPar
Bu durumda başlangıç düğümü bu makroyla şöyle de oluşturulabilir:

\begin{sphinxVerbatim}[commandchars=\\\{\}]
\PYG{k}{struct}\PYG{+w}{ }\PYG{n+nc}{list\PYGZus{}head}\PYG{+w}{ }\PYG{n}{head}\PYG{+w}{ }\PYG{o}{=}\PYG{+w}{ }\PYG{n}{LIST\PYGZus{}HEAD\PYGZus{}INIT}\PYG{p}{(}\PYG{n}{head}\PYG{p}{)}\PYG{p}{;}
\end{sphinxVerbatim}

\sphinxAtStartPar
Aslında bu tanımlamanın tamamını yapan aşağıdaki gibi bir makro da bulundurulmuştur:

\begin{sphinxVerbatim}[commandchars=\\\{\}]
\PYG{c+cp}{\PYGZsh{}}\PYG{c+cp}{define LIST\PYGZus{}HEAD(name) \PYGZbs{}}
\PYG{c+cp}{    struct list\PYGZus{}head name = LIST\PYGZus{}HEAD\PYGZus{}INIT(name)}
\end{sphinxVerbatim}

\sphinxAtStartPar
O halde biz başlangıç düğümünü basit bir biçimde şöyle oluşturabiliriz:

\begin{sphinxVerbatim}[commandchars=\\\{\}]
\PYG{n}{LIST\PYGZus{}HEAD}\PYG{p}{(}\PYG{n}{head}\PYG{p}{)}\PYG{p}{;}
\end{sphinxVerbatim}

\sphinxAtStartPar
Linux’un bağlı liste gerçekleştiriminde \sphinxcode{\sphinxupquote{next}} ve \sphinxcode{\sphinxupquote{prev}} göstericilerinde \sphinxcode{\sphinxupquote{NULL}} adres
kullanılmamıştır. Son düğümün \sphinxcode{\sphinxupquote{next}} göstericisi \sphinxcode{\sphinxupquote{NULL}} yerine başlangıç düğümünü, ilk düğümün
\sphinxcode{\sphinxupquote{prev}} göstericisi de \sphinxcode{\sphinxupquote{NULL}} yerine son düğümü göstermektedir. Yani aslında bağlı liste
gerçekleştirimi döngüsel (circular) gibidir. Herhangi bir düğümden başlanarak başlangıç düğümü de
atlanırsa tam dolaşım sağlanabilir. Bağlı listenin başlangıç düğümünün \sphinxcode{\sphinxupquote{prev}} göstericisi de son
düğümü göstermektedir.


\subsection{Düğüm Ekleme}
\label{\detokenize{kernellinkedlistandhashtable:dugum-ekleme}}
\sphinxAtStartPar
Bağlı listenin başına \sphinxcode{\sphinxupquote{list\_head}} düğümünü eklemek için \sphinxcode{\sphinxupquote{list\_add}} fonksiyonu kullanılmaktadır:

\begin{sphinxVerbatim}[commandchars=\\\{\}]
\PYG{k}{static}\PYG{+w}{ }\PYG{k+kr}{inline}\PYG{+w}{ }\PYG{k+kt}{void}\PYG{+w}{ }\PYG{n+nf}{list\PYGZus{}add}\PYG{p}{(}\PYG{k}{struct}\PYG{+w}{ }\PYG{n+nc}{list\PYGZus{}head}\PYG{+w}{ }\PYG{o}{*}\PYG{n}{new}\PYG{p}{,}\PYG{+w}{ }\PYG{k}{struct}\PYG{+w}{ }\PYG{n+nc}{list\PYGZus{}head}\PYG{+w}{ }\PYG{o}{*}\PYG{n}{head}\PYG{p}{)}
\PYG{p}{\PYGZob{}}
\PYG{+w}{    }\PYG{n}{\PYGZus{}\PYGZus{}list\PYGZus{}add}\PYG{p}{(}\PYG{n}{new}\PYG{p}{,}\PYG{+w}{ }\PYG{n}{head}\PYG{p}{,}\PYG{+w}{ }\PYG{n}{head}\PYG{o}{\PYGZhy{}}\PYG{o}{\PYGZgt{}}\PYG{n}{next}\PYG{p}{)}\PYG{p}{;}
\PYG{p}{\PYGZcb{}}
\end{sphinxVerbatim}

\sphinxAtStartPar
Fonksiyonun birinci parametresi eklenecek düğümün adresini, ikinci parametresi başlangıç düğümünün
adresini almaktadır.

\sphinxAtStartPar
Çekirdek kodlarında iki alt tire ile başlayan fonksiyonlar aşağı seviyeli fonksiyonlardır. Yani doğrudan
çağrılmayan ancak başka fonksiyonların içerisinden dolaylı biçimde çağrılan fonksiyonlardır. Buradaki
\sphinxcode{\sphinxupquote{\_\_list\_add}} fonksiyonu şöyle yazılmıştır:

\begin{sphinxVerbatim}[commandchars=\\\{\}]
\PYG{k}{static}\PYG{+w}{ }\PYG{k+kr}{inline}\PYG{+w}{ }\PYG{k+kt}{void}\PYG{+w}{ }\PYG{n+nf}{\PYGZus{}\PYGZus{}list\PYGZus{}add}\PYG{p}{(}\PYG{k}{struct}\PYG{+w}{ }\PYG{n+nc}{list\PYGZus{}head}\PYG{+w}{ }\PYG{o}{*}\PYG{n}{new}\PYG{p}{,}
\PYG{+w}{              }\PYG{k}{struct}\PYG{+w}{ }\PYG{n+nc}{list\PYGZus{}head}\PYG{+w}{ }\PYG{o}{*}\PYG{n}{prev}\PYG{p}{,}
\PYG{+w}{              }\PYG{k}{struct}\PYG{+w}{ }\PYG{n+nc}{list\PYGZus{}head}\PYG{+w}{ }\PYG{o}{*}\PYG{n}{next}\PYG{p}{)}
\PYG{p}{\PYGZob{}}
\PYG{+w}{    }\PYG{k}{if}\PYG{+w}{ }\PYG{p}{(}\PYG{o}{!}\PYG{n}{\PYGZus{}\PYGZus{}list\PYGZus{}add\PYGZus{}valid}\PYG{p}{(}\PYG{n}{new}\PYG{p}{,}\PYG{+w}{ }\PYG{n}{prev}\PYG{p}{,}\PYG{+w}{ }\PYG{n}{next}\PYG{p}{)}\PYG{p}{)}
\PYG{+w}{        }\PYG{k}{return}\PYG{p}{;}

\PYG{+w}{    }\PYG{n}{next}\PYG{o}{\PYGZhy{}}\PYG{o}{\PYGZgt{}}\PYG{n}{prev}\PYG{+w}{ }\PYG{o}{=}\PYG{+w}{ }\PYG{n}{new}\PYG{p}{;}
\PYG{+w}{    }\PYG{n}{new}\PYG{o}{\PYGZhy{}}\PYG{o}{\PYGZgt{}}\PYG{n}{next}\PYG{+w}{ }\PYG{o}{=}\PYG{+w}{ }\PYG{n}{next}\PYG{p}{;}
\PYG{+w}{    }\PYG{n}{new}\PYG{o}{\PYGZhy{}}\PYG{o}{\PYGZgt{}}\PYG{n}{prev}\PYG{+w}{ }\PYG{o}{=}\PYG{+w}{ }\PYG{n}{prev}\PYG{p}{;}
\PYG{+w}{    }\PYG{n}{WRITE\PYGZus{}ONCE}\PYG{p}{(}\PYG{n}{prev}\PYG{o}{\PYGZhy{}}\PYG{o}{\PYGZgt{}}\PYG{n}{next}\PYG{p}{,}\PYG{+w}{ }\PYG{n}{new}\PYG{p}{)}\PYG{p}{;}
\PYG{p}{\PYGZcb{}}
\end{sphinxVerbatim}

\sphinxAtStartPar
Buradaki \sphinxcode{\sphinxupquote{WRITE\_ONCE}} makrosu atama işlemini volatile erişimle yapmaktadır. \sphinxcode{\sphinxupquote{WRITE\_ONCE(a, b)}}
çağrısını \sphinxcode{\sphinxupquote{a = b}} gibi düşünebilirsiniz.

\sphinxAtStartPar
Bağlı listenin sonuna düğüm eklemek için \sphinxcode{\sphinxupquote{list\_add\_tail}} fonksiyonu kullanılmaktadır:

\begin{sphinxVerbatim}[commandchars=\\\{\}]
\PYG{k}{static}\PYG{+w}{ }\PYG{k+kr}{inline}\PYG{+w}{ }\PYG{k+kt}{void}\PYG{+w}{ }\PYG{n+nf}{list\PYGZus{}add\PYGZus{}tail}\PYG{p}{(}\PYG{k}{struct}\PYG{+w}{ }\PYG{n+nc}{list\PYGZus{}head}\PYG{+w}{ }\PYG{o}{*}\PYG{n}{new}\PYG{p}{,}\PYG{+w}{ }\PYG{k}{struct}\PYG{+w}{ }\PYG{n+nc}{list\PYGZus{}head}\PYG{+w}{ }\PYG{o}{*}\PYG{n}{head}\PYG{p}{)}
\PYG{p}{\PYGZob{}}
\PYG{+w}{    }\PYG{n}{\PYGZus{}\PYGZus{}list\PYGZus{}add}\PYG{p}{(}\PYG{n}{new}\PYG{p}{,}\PYG{+w}{ }\PYG{n}{head}\PYG{o}{\PYGZhy{}}\PYG{o}{\PYGZgt{}}\PYG{n}{prev}\PYG{p}{,}\PYG{+w}{ }\PYG{n}{head}\PYG{p}{)}\PYG{p}{;}
\PYG{p}{\PYGZcb{}}
\end{sphinxVerbatim}


\subsection{Liste Dolaşımı}
\label{\detokenize{kernellinkedlistandhashtable:liste-dolasimi}}
\sphinxAtStartPar
Peki biz bağlı listeyi nasıl dolaşabiliriz? Başlangıç düğümünü geçerek sürekli ileriye gidersek son
düğüm de başlangıç düğümünü gösterdiğine göre dolaşımı aşağıdaki gibi bir for döngüsüyle yapabiliriz:

\begin{sphinxVerbatim}[commandchars=\\\{\}]
\PYG{k}{struct}\PYG{+w}{ }\PYG{n+nc}{list\PYGZus{}head}\PYG{+w}{ }\PYG{o}{*}\PYG{n}{lh}\PYG{p}{;}

\PYG{k}{for}\PYG{+w}{ }\PYG{p}{(}\PYG{n}{lh}\PYG{+w}{ }\PYG{o}{=}\PYG{+w}{ }\PYG{n}{head}\PYG{p}{.}\PYG{n}{next}\PYG{p}{;}\PYG{+w}{ }\PYG{n}{lh}\PYG{+w}{ }\PYG{o}{!}\PYG{o}{=}\PYG{+w}{ }\PYG{o}{\PYGZam{}}\PYG{n}{head}\PYG{p}{;}\PYG{+w}{ }\PYG{n}{lh}\PYG{+w}{ }\PYG{o}{=}\PYG{+w}{ }\PYG{n}{lh}\PYG{o}{\PYGZhy{}}\PYG{o}{\PYGZgt{}}\PYG{n}{next}\PYG{p}{)}\PYG{+w}{ }\PYG{p}{\PYGZob{}}
\PYG{+w}{    }\PYG{c+cm}{/* ... */}
\PYG{p}{\PYGZcb{}}
\end{sphinxVerbatim}

\sphinxAtStartPar
Tabii böyle bir döngüde dolaşım yapılırken aslında biz her yinelemede ilgili yapının başlangıç adresini
değil \sphinxcode{\sphinxupquote{list\_head}} yapısının başlangıç adresini elde ederiz. \sphinxcode{\sphinxupquote{container\_of}} makrosu ile bu adresin
yapı adresine dönüştürülmesi gerekir. Örneğin biz \sphinxcode{\sphinxupquote{struct SAMPLE}} türünden yapı nesnelerini birbirine
bağlamış olalım:

\begin{sphinxVerbatim}[commandchars=\\\{\}]
\PYG{n}{LIST\PYGZus{}HEAD}\PYG{p}{(}\PYG{n}{head}\PYG{p}{)}\PYG{p}{;}

\PYG{k}{struct}\PYG{+w}{ }\PYG{n+nc}{SAMPLE}\PYG{+w}{ }\PYG{p}{\PYGZob{}}
\PYG{+w}{    }\PYG{k+kt}{int}\PYG{+w}{ }\PYG{n}{a}\PYG{p}{;}
\PYG{+w}{    }\PYG{k}{struct}\PYG{+w}{ }\PYG{n+nc}{list\PYGZus{}head}\PYG{+w}{ }\PYG{n}{link}\PYG{p}{;}
\PYG{p}{\PYGZcb{}}\PYG{p}{;}
\end{sphinxVerbatim}

\sphinxAtStartPar
Eklemeleri şöyle yapmış olalım:

\begin{sphinxVerbatim}[commandchars=\\\{\}]
\PYG{k}{struct}\PYG{+w}{ }\PYG{n+nc}{SAMPLE}\PYG{+w}{ }\PYG{o}{*}\PYG{n}{ps}\PYG{p}{;}

\PYG{k}{for}\PYG{+w}{ }\PYG{p}{(}\PYG{k+kt}{int}\PYG{+w}{ }\PYG{n}{i}\PYG{+w}{ }\PYG{o}{=}\PYG{+w}{ }\PYG{l+m+mi}{0}\PYG{p}{;}\PYG{+w}{ }\PYG{n}{i}\PYG{+w}{ }\PYG{o}{\PYGZlt{}}\PYG{+w}{ }\PYG{l+m+mi}{10}\PYG{p}{;}\PYG{+w}{ }\PYG{o}{+}\PYG{o}{+}\PYG{n}{i}\PYG{p}{)}\PYG{+w}{ }\PYG{p}{\PYGZob{}}
\PYG{+w}{    }\PYG{k}{if}\PYG{+w}{ }\PYG{p}{(}\PYG{p}{(}\PYG{n}{ps}\PYG{+w}{ }\PYG{o}{=}\PYG{+w}{ }\PYG{p}{(}\PYG{k}{struct}\PYG{+w}{ }\PYG{n+nc}{SAMPLE}\PYG{+w}{ }\PYG{o}{*}\PYG{p}{)}\PYG{n}{malloc}\PYG{p}{(}\PYG{k}{sizeof}\PYG{p}{(}\PYG{k}{struct}\PYG{+w}{ }\PYG{n+nc}{SAMPLE}\PYG{p}{)}\PYG{p}{)}\PYG{p}{)}\PYG{+w}{ }\PYG{o}{=}\PYG{o}{=}\PYG{+w}{ }\PYG{n+nb}{NULL}\PYG{p}{)}\PYG{+w}{ }\PYG{p}{\PYGZob{}}
\PYG{+w}{        }\PYG{n}{fprintf}\PYG{p}{(}\PYG{n}{stderr}\PYG{p}{,}\PYG{+w}{ }\PYG{l+s}{\PYGZdq{}}\PYG{l+s}{cannot allocate memory!...}\PYG{l+s+se}{\PYGZbs{}n}\PYG{l+s}{\PYGZdq{}}\PYG{p}{)}\PYG{p}{;}
\PYG{+w}{        }\PYG{n}{exit}\PYG{p}{(}\PYG{n}{EXIT\PYGZus{}FAILURE}\PYG{p}{)}\PYG{p}{;}
\PYG{+w}{    }\PYG{p}{\PYGZcb{}}
\PYG{+w}{    }\PYG{n}{ps}\PYG{o}{\PYGZhy{}}\PYG{o}{\PYGZgt{}}\PYG{n}{a}\PYG{+w}{ }\PYG{o}{=}\PYG{+w}{ }\PYG{n}{i}\PYG{p}{;}
\PYG{+w}{    }\PYG{n}{list\PYGZus{}add\PYGZus{}tail}\PYG{p}{(}\PYG{o}{\PYGZam{}}\PYG{n}{ps}\PYG{o}{\PYGZhy{}}\PYG{o}{\PYGZgt{}}\PYG{n}{link}\PYG{p}{,}\PYG{+w}{ }\PYG{o}{\PYGZam{}}\PYG{n}{head}\PYG{p}{)}\PYG{p}{;}
\PYG{p}{\PYGZcb{}}
\end{sphinxVerbatim}

\begin{sphinxadmonition}{note}{Not:}
\sphinxAtStartPar
Çekirdek kodlarında \sphinxstyleemphasis{malloc} gibi kullanıcı modu fonksiyonları kullanılamaz. Ancak burada yalnızca
anlaşılır bir test örneği verilmektedir.
\end{sphinxadmonition}

\sphinxAtStartPar
\sphinxcode{\sphinxupquote{struct SAMPLE}} yapımızdaki \sphinxcode{\sphinxupquote{link}} elemanı \sphinxcode{\sphinxupquote{list\_head}} yapılarını birbirine bağlamak için
bulundurulmuştur. Örneğimizdeki \sphinxcode{\sphinxupquote{list\_add\_tail}} fonksiyonunda ekleme yapılacak düğümün \sphinxcode{\sphinxupquote{SAMPLE}}
yapısının içerisindeki \sphinxcode{\sphinxupquote{link}} elemanının adresi olduğuna dikkat ediniz. Şimdi dolaşımı şöyle
yapabiliriz:

\begin{sphinxVerbatim}[commandchars=\\\{\}]
\PYG{k}{struct}\PYG{+w}{ }\PYG{n+nc}{SAMPLE}\PYG{+w}{ }\PYG{o}{*}\PYG{n}{ps}\PYG{p}{;}
\PYG{k}{struct}\PYG{+w}{ }\PYG{n+nc}{list\PYGZus{}head}\PYG{+w}{ }\PYG{o}{*}\PYG{n}{lh}\PYG{p}{;}

\PYG{k}{for}\PYG{+w}{ }\PYG{p}{(}\PYG{n}{lh}\PYG{+w}{ }\PYG{o}{=}\PYG{+w}{ }\PYG{n}{head}\PYG{p}{.}\PYG{n}{next}\PYG{p}{;}\PYG{+w}{ }\PYG{n}{lh}\PYG{+w}{ }\PYG{o}{!}\PYG{o}{=}\PYG{+w}{ }\PYG{o}{\PYGZam{}}\PYG{n}{head}\PYG{p}{;}\PYG{+w}{ }\PYG{n}{lh}\PYG{+w}{ }\PYG{o}{=}\PYG{+w}{ }\PYG{n}{lh}\PYG{o}{\PYGZhy{}}\PYG{o}{\PYGZgt{}}\PYG{n}{next}\PYG{p}{)}\PYG{+w}{ }\PYG{p}{\PYGZob{}}
\PYG{+w}{    }\PYG{n}{ps}\PYG{+w}{ }\PYG{o}{=}\PYG{+w}{ }\PYG{n}{container\PYGZus{}of}\PYG{p}{(}\PYG{n}{lh}\PYG{p}{,}\PYG{+w}{ }\PYG{k}{struct}\PYG{+w}{ }\PYG{n+nc}{SAMPLE}\PYG{p}{,}\PYG{+w}{ }\PYG{n}{link}\PYG{p}{)}\PYG{p}{;}
\PYG{+w}{    }\PYG{c+cm}{/* ... */}
\PYG{p}{\PYGZcb{}}
\end{sphinxVerbatim}

\sphinxAtStartPar
Burada artık \sphinxcode{\sphinxupquote{ps}} göstericisi \sphinxcode{\sphinxupquote{list\_head}} nesnesini değil \sphinxcode{\sphinxupquote{struct SAMPLE}} nesnesini göstermektedir.
Aslında \sphinxcode{\sphinxupquote{list.h}} dosyası içerisinde buradaki döngüyü oluşturan \sphinxcode{\sphinxupquote{list\_for\_each}} isimli bir makro da
bulundurulmuştur:

\begin{sphinxVerbatim}[commandchars=\\\{\}]
\PYG{k}{static}\PYG{+w}{ }\PYG{k+kr}{inline}\PYG{+w}{ }\PYG{k+kt}{int}\PYG{+w}{ }\PYG{n+nf}{list\PYGZus{}is\PYGZus{}head}\PYG{p}{(}\PYG{k}{const}\PYG{+w}{ }\PYG{k}{struct}\PYG{+w}{ }\PYG{n+nc}{list\PYGZus{}head}\PYG{+w}{ }\PYG{o}{*}\PYG{n}{list}\PYG{p}{,}
\PYG{+w}{                               }\PYG{k}{const}\PYG{+w}{ }\PYG{k}{struct}\PYG{+w}{ }\PYG{n+nc}{list\PYGZus{}head}\PYG{+w}{ }\PYG{o}{*}\PYG{n}{head}\PYG{p}{)}
\PYG{p}{\PYGZob{}}
\PYG{+w}{    }\PYG{k}{return}\PYG{+w}{ }\PYG{n}{list}\PYG{+w}{ }\PYG{o}{=}\PYG{o}{=}\PYG{+w}{ }\PYG{n}{head}\PYG{p}{;}
\PYG{p}{\PYGZcb{}}

\PYG{c+cp}{\PYGZsh{}}\PYG{c+cp}{define list\PYGZus{}for\PYGZus{}each(pos, head) \PYGZbs{}}
\PYG{c+cp}{    for (pos = (head)\PYGZhy{}\PYGZgt{}next; !list\PYGZus{}is\PYGZus{}head(pos, (head)); pos = pos\PYGZhy{}\PYGZgt{}next)}
\end{sphinxVerbatim}

\sphinxAtStartPar
Makronun birinci parametresi \sphinxcode{\sphinxupquote{list\_head}} türünden bir göstericiyi, ikinci parametresi başlangıç
düğümünün adresini almaktadır. Döngünün her yinelenmesinde bu göstericiye sonraki düğümün \sphinxcode{\sphinxupquote{list\_head}}
adresi atanmaktadır. Örneğin:

\begin{sphinxVerbatim}[commandchars=\\\{\}]
\PYG{k}{struct}\PYG{+w}{ }\PYG{n+nc}{list\PYGZus{}head}\PYG{+w}{ }\PYG{o}{*}\PYG{n}{lh}\PYG{p}{;}

\PYG{n}{list\PYGZus{}for\PYGZus{}each}\PYG{p}{(}\PYG{n}{lh}\PYG{p}{,}\PYG{+w}{ }\PYG{o}{\PYGZam{}}\PYG{n}{head}\PYG{p}{)}\PYG{+w}{ }\PYG{p}{\PYGZob{}}
\PYG{+w}{    }\PYG{n}{ps}\PYG{+w}{ }\PYG{o}{=}\PYG{+w}{ }\PYG{n}{container\PYGZus{}of}\PYG{p}{(}\PYG{n}{lh}\PYG{p}{,}\PYG{+w}{ }\PYG{k}{struct}\PYG{+w}{ }\PYG{n+nc}{SAMPLE}\PYG{p}{,}\PYG{+w}{ }\PYG{n}{link}\PYG{p}{)}\PYG{p}{;}
\PYG{+w}{    }\PYG{n}{printf}\PYG{p}{(}\PYG{l+s}{\PYGZdq{}}\PYG{l+s}{\PYGZpc{}d}\PYG{l+s+se}{\PYGZbs{}n}\PYG{l+s}{\PYGZdq{}}\PYG{p}{,}\PYG{+w}{ }\PYG{n}{ps}\PYG{o}{\PYGZhy{}}\PYG{o}{\PYGZgt{}}\PYG{n}{a}\PYG{p}{)}\PYG{p}{;}
\PYG{p}{\PYGZcb{}}
\end{sphinxVerbatim}

\sphinxAtStartPar
Aslında çekirdekteki \sphinxcode{\sphinxupquote{list.h}} dosyası içerisinde yukarıdaki dolaşımı tek hamlede yapan
\sphinxcode{\sphinxupquote{list\_for\_each\_entry}} isimli bir makro da bulunmaktadır:

\begin{sphinxVerbatim}[commandchars=\\\{\}]
\PYG{c+cp}{\PYGZsh{}}\PYG{c+cp}{define list\PYGZus{}for\PYGZus{}each\PYGZus{}entry(pos, head, member)              \PYGZbs{}}
\PYG{c+cp}{    for (pos = list\PYGZus{}first\PYGZus{}entry(head, typeof(*pos), member); \PYGZbs{}}
\PYG{c+cp}{         !list\PYGZus{}entry\PYGZus{}is\PYGZus{}head(pos, head, member);             \PYGZbs{}}
\PYG{c+cp}{         pos = list\PYGZus{}next\PYGZus{}entry(pos, member))}
\end{sphinxVerbatim}

\sphinxAtStartPar
Makronun birinci parametresi asıl yapı türünden bir göstericidir; makro her yinelemede bu göstericiye
sonraki düğümün adresini yerleştirmektedir. Makronun ikinci parametresi başlangıç düğümünün adresini,
üçüncü parametresi ise yapı içerisindeki \sphinxcode{\sphinxupquote{list\_head}} elemanın ismini belirtmektedir. Bu makronun
eşdeğeri bazı ayrıntılar ihmal edilerek şöyle de yazılabilir:

\begin{sphinxVerbatim}[commandchars=\\\{\}]
\PYG{c+cp}{\PYGZsh{}}\PYG{c+cp}{define my\PYGZus{}list\PYGZus{}for\PYGZus{}each\PYGZus{}entry(pos, head, member)                                            \PYGZbs{}}
\PYG{c+cp}{    for (pos = container\PYGZus{}of((head)\PYGZhy{}\PYGZgt{}next, typeof(*pos), member); \PYGZam{}(pos)\PYGZhy{}\PYGZgt{}member != head;     \PYGZbs{}}
\PYG{c+cp}{        pos = container\PYGZus{}of((pos)\PYGZhy{}\PYGZgt{}member.next, typeof(*pos), member))}
\end{sphinxVerbatim}

\begin{sphinxadmonition}{note}{Not:}
\sphinxAtStartPar
C standartlarında bir ifadenin türünü veren bir tür belirleyicisi yoktu. C++’a C++11 ile birlikte
\sphinxcode{\sphinxupquote{decltype}} ismi ile böyle bir belirleyici eklendi. gcc derleyicileri uzun süredir bu işi yapan
\sphinxcode{\sphinxupquote{typeof}} isimli belirleyiciyi desteklemektedir. Nihayet C’ye de C23 ile birlikte resmi olarak
\sphinxcode{\sphinxupquote{typeof}} belirleyicisi eklenmiştir. Makroda türün \sphinxcode{\sphinxupquote{typeof(*pos)}} ifadesiyle elde edildiğine dikkat
ediniz.
\end{sphinxadmonition}

\sphinxAtStartPar
Bu makro ile dolaşım şöyle sağlanabilir:

\begin{sphinxVerbatim}[commandchars=\\\{\}]
\PYG{k}{struct}\PYG{+w}{ }\PYG{n+nc}{SAMPLE}\PYG{+w}{ }\PYG{o}{*}\PYG{n}{ps}\PYG{p}{;}

\PYG{n}{list\PYGZus{}for\PYGZus{}each\PYGZus{}entry}\PYG{p}{(}\PYG{n}{ps}\PYG{p}{,}\PYG{+w}{ }\PYG{o}{\PYGZam{}}\PYG{n}{head}\PYG{p}{,}\PYG{+w}{ }\PYG{n}{link}\PYG{p}{)}\PYG{+w}{ }\PYG{p}{\PYGZob{}}
\PYG{+w}{    }\PYG{c+cm}{/* ... */}
\PYG{p}{\PYGZcb{}}
\end{sphinxVerbatim}

\begin{sphinxadmonition}{tip}{Tüyo:}
\sphinxAtStartPar
Bağlı liste makrolarının ve fonksiyonlarının isimlerinin sonunda \sphinxstyleemphasis{entry} sözcüğü varsa bu makro ya da
fonksiyon asıl yapıya ilişkin adres, yoksa \sphinxcode{\sphinxupquote{list\_head}} türünden adres verip almaktadır.
\end{sphinxadmonition}

\sphinxAtStartPar
Ayrıca \sphinxcode{\sphinxupquote{list.h}} içerisinde \sphinxcode{\sphinxupquote{list\_for\_each\_entry\_safe}} isimli bir döngü makrosu da
bulundurulmuştur. Bu makro eğer liste boşsa hiç dolaşım yapmamaktadır. Makro aşağıdaki gibi
yazılmıştır:

\begin{sphinxVerbatim}[commandchars=\\\{\}]
\PYG{c+cp}{\PYGZsh{}}\PYG{c+cp}{define list\PYGZus{}for\PYGZus{}each\PYGZus{}entry\PYGZus{}safe(pos, n, head, member)          \PYGZbs{}}
\PYG{c+cp}{    for (pos = list\PYGZus{}first\PYGZus{}entry(head, typeof(*pos), member),    \PYGZbs{}}
\PYG{c+cp}{         n = list\PYGZus{}next\PYGZus{}entry(pos, member);                       \PYGZbs{}}
\PYG{c+cp}{         !list\PYGZus{}entry\PYGZus{}is\PYGZus{}head(pos, head, member);                 \PYGZbs{}}
\PYG{c+cp}{         pos = n, n = list\PYGZus{}next\PYGZus{}entry(n, member))}
\end{sphinxVerbatim}

\sphinxAtStartPar
Makronun birinci elemanı dolaşımda kullanılacak asıl yapı türünden göstericiyi almaktadır. Makronun
ikinci parametresi de asıl yapı göstericidir. Üçüncü ve dördüncü parametreler sırasıyla kök düğümün
adresi ve bağ düğümünün ismini almaktadır.


\subsection{Düğüm Silme}
\label{\detokenize{kernellinkedlistandhashtable:dugum-silme}}
\sphinxAtStartPar
Bağlı listeden bir düğümü silmek için \sphinxcode{\sphinxupquote{list\_del}} fonksiyonu kullanılmaktadır. Fonksiyon şöyle
tanımlanmıştır:

\begin{sphinxVerbatim}[commandchars=\\\{\}]
\PYG{k}{static}\PYG{+w}{ }\PYG{k+kr}{inline}\PYG{+w}{ }\PYG{k+kt}{void}\PYG{+w}{ }\PYG{n+nf}{\PYGZus{}\PYGZus{}list\PYGZus{}del}\PYG{p}{(}\PYG{k}{struct}\PYG{+w}{ }\PYG{n+nc}{list\PYGZus{}head}\PYG{+w}{ }\PYG{o}{*}\PYG{n}{prev}\PYG{p}{,}\PYG{+w}{ }\PYG{k}{struct}\PYG{+w}{ }\PYG{n+nc}{list\PYGZus{}head}\PYG{+w}{ }\PYG{o}{*}\PYG{n}{next}\PYG{p}{)}
\PYG{p}{\PYGZob{}}
\PYG{+w}{    }\PYG{n}{next}\PYG{o}{\PYGZhy{}}\PYG{o}{\PYGZgt{}}\PYG{n}{prev}\PYG{+w}{ }\PYG{o}{=}\PYG{+w}{ }\PYG{n}{prev}\PYG{p}{;}
\PYG{+w}{    }\PYG{n}{WRITE\PYGZus{}ONCE}\PYG{p}{(}\PYG{n}{prev}\PYG{o}{\PYGZhy{}}\PYG{o}{\PYGZgt{}}\PYG{n}{next}\PYG{p}{,}\PYG{+w}{ }\PYG{n}{next}\PYG{p}{)}\PYG{p}{;}
\PYG{p}{\PYGZcb{}}

\PYG{k}{static}\PYG{+w}{ }\PYG{k+kr}{inline}\PYG{+w}{ }\PYG{k+kt}{void}\PYG{+w}{ }\PYG{n+nf}{\PYGZus{}\PYGZus{}list\PYGZus{}del\PYGZus{}entry}\PYG{p}{(}\PYG{k}{struct}\PYG{+w}{ }\PYG{n+nc}{list\PYGZus{}head}\PYG{+w}{ }\PYG{o}{*}\PYG{n}{entry}\PYG{p}{)}
\PYG{p}{\PYGZob{}}
\PYG{+w}{    }\PYG{k}{if}\PYG{+w}{ }\PYG{p}{(}\PYG{o}{!}\PYG{n}{\PYGZus{}\PYGZus{}list\PYGZus{}del\PYGZus{}entry\PYGZus{}valid}\PYG{p}{(}\PYG{n}{entry}\PYG{p}{)}\PYG{p}{)}
\PYG{+w}{        }\PYG{k}{return}\PYG{p}{;}

\PYG{+w}{    }\PYG{n}{\PYGZus{}\PYGZus{}list\PYGZus{}del}\PYG{p}{(}\PYG{n}{entry}\PYG{o}{\PYGZhy{}}\PYG{o}{\PYGZgt{}}\PYG{n}{prev}\PYG{p}{,}\PYG{+w}{ }\PYG{n}{entry}\PYG{o}{\PYGZhy{}}\PYG{o}{\PYGZgt{}}\PYG{n}{next}\PYG{p}{)}\PYG{p}{;}
\PYG{p}{\PYGZcb{}}

\PYG{k}{static}\PYG{+w}{ }\PYG{k+kr}{inline}\PYG{+w}{ }\PYG{k+kt}{void}\PYG{+w}{ }\PYG{n+nf}{list\PYGZus{}del}\PYG{p}{(}\PYG{k}{struct}\PYG{+w}{ }\PYG{n+nc}{list\PYGZus{}head}\PYG{+w}{ }\PYG{o}{*}\PYG{n}{entry}\PYG{p}{)}
\PYG{p}{\PYGZob{}}
\PYG{+w}{    }\PYG{n}{\PYGZus{}\PYGZus{}list\PYGZus{}del\PYGZus{}entry}\PYG{p}{(}\PYG{n}{entry}\PYG{p}{)}\PYG{p}{;}
\PYG{+w}{    }\PYG{n}{entry}\PYG{o}{\PYGZhy{}}\PYG{o}{\PYGZgt{}}\PYG{n}{next}\PYG{+w}{ }\PYG{o}{=}\PYG{+w}{ }\PYG{n}{LIST\PYGZus{}POISON1}\PYG{p}{;}
\PYG{+w}{    }\PYG{n}{entry}\PYG{o}{\PYGZhy{}}\PYG{o}{\PYGZgt{}}\PYG{n}{prev}\PYG{+w}{ }\PYG{o}{=}\PYG{+w}{ }\PYG{n}{LIST\PYGZus{}POISON2}\PYG{p}{;}
\PYG{p}{\PYGZcb{}}
\end{sphinxVerbatim}

\sphinxAtStartPar
Fonksiyonun parametresi silinecek düğüme ilişkin \sphinxcode{\sphinxupquote{list\_head}} nesnesinin adresini almaktadır. Burada
silinen düğümdeki \sphinxcode{\sphinxupquote{next}} ve \sphinxcode{\sphinxupquote{prev}} göstericilerine özel bazı değerlerin atandığını görüyorsunuz. Bu
değerler silinmiş düğümün kullanılması durumunda \sphinxstyleemphasis{page fault} oluşmasına yol açmaktadır; dolayısıyla bir
çeşit debug mekanizması oluşturmak amacıyla atandıklarını söyleyebiliriz.


\subsection{Araya Düğüm Ekleme}
\label{\detokenize{kernellinkedlistandhashtable:araya-dugum-ekleme}}
\sphinxAtStartPar
Bağlı listenin arasına düğüm ekleyen ayrı bir insert fonksiyonu yoktur. Zaten \sphinxcode{\sphinxupquote{list\_add}} fonksiyonu
araya ekleme işlemini de yapmaktadır. Örneğin biz araya şöyle ekleme yapabiliriz:

\begin{sphinxVerbatim}[commandchars=\\\{\}]
\PYG{n}{list\PYGZus{}add}\PYG{p}{(}\PYG{o}{\PYGZam{}}\PYG{n}{ps}\PYG{o}{\PYGZhy{}}\PYG{o}{\PYGZgt{}}\PYG{n}{link}\PYG{p}{,}\PYG{+w}{ }\PYG{o}{\PYGZam{}}\PYG{n}{ps\PYGZus{}insert}\PYG{o}{\PYGZhy{}}\PYG{o}{\PYGZgt{}}\PYG{n}{link}\PYG{p}{)}\PYG{p}{;}
\end{sphinxVerbatim}

\sphinxAtStartPar
Burada \sphinxcode{\sphinxupquote{ps}} eklenecek düğümün asıl yapı adresini, \sphinxcode{\sphinxupquote{ps\_insert}} ise önüne eklemenin yapılacağı asıl
yapı adresini belirtmektedir.

\sphinxAtStartPar
Bir bağlı listeyi tümden serbest bırakmak için düğümlerin tek tek serbest bırakılması gerekmektedir.
Buradaki düğümler aslında başka yapıların içerisinde olduğuna göre liste içerisindeki o yapı nesnelerinin
serbest bırakılması gerekir.


\section{Tam Örnek: Kullanıcı Alanında Bağlı Liste Kullanımı}
\label{\detokenize{kernellinkedlistandhashtable:tam-ornek-kullanici-alaninda-bagli-liste-kullanimi}}
\sphinxAtStartPar
Aşağıda kullanıcı alanında çekirdekteki bağlı liste kullanımına bir örnek verilmiştir.

\begin{sphinxVerbatim}[commandchars=\\\{\}]
\PYG{c+cp}{\PYGZsh{}}\PYG{c+cp}{include}\PYG{+w}{ }\PYG{c+cpf}{\PYGZlt{}stdio.h\PYGZgt{}}
\PYG{c+cp}{\PYGZsh{}}\PYG{c+cp}{include}\PYG{+w}{ }\PYG{c+cpf}{\PYGZlt{}stdlib.h\PYGZgt{}}
\PYG{c+cp}{\PYGZsh{}}\PYG{c+cp}{include}\PYG{+w}{ }\PYG{c+cpf}{\PYGZlt{}stddef.h\PYGZgt{}}

\PYG{k}{struct}\PYG{+w}{ }\PYG{n+nc}{list\PYGZus{}head}\PYG{+w}{ }\PYG{p}{\PYGZob{}}
\PYG{+w}{    }\PYG{k}{struct}\PYG{+w}{ }\PYG{n+nc}{list\PYGZus{}head}\PYG{+w}{ }\PYG{o}{*}\PYG{n}{next}\PYG{p}{,}\PYG{+w}{ }\PYG{o}{*}\PYG{n}{prev}\PYG{p}{;}
\PYG{p}{\PYGZcb{}}\PYG{p}{;}

\PYG{c+cp}{\PYGZsh{}}\PYG{c+cp}{define LIST\PYGZus{}HEAD\PYGZus{}INIT(name) \PYGZob{} \PYGZam{}(name), \PYGZam{}(name) \PYGZcb{}}

\PYG{c+cp}{\PYGZsh{}}\PYG{c+cp}{define LIST\PYGZus{}HEAD(name) \PYGZbs{}}
\PYG{c+cp}{    struct list\PYGZus{}head name = LIST\PYGZus{}HEAD\PYGZus{}INIT(name)}

\PYG{c+cp}{\PYGZsh{}}\PYG{c+cp}{define container\PYGZus{}of(ptr, type, member) (\PYGZob{}          \PYGZbs{}}
\PYG{c+cp}{        void *\PYGZus{}\PYGZus{}mptr = (void *)(ptr);               \PYGZbs{}}
\PYG{c+cp}{        ((type *)(\PYGZus{}\PYGZus{}mptr \PYGZhy{} offsetof(type, member))); \PYGZcb{})}

\PYG{c+cp}{\PYGZsh{}}\PYG{c+cp}{define list\PYGZus{}entry(ptr, type, member) \PYGZbs{}}
\PYG{c+cp}{    container\PYGZus{}of(ptr, type, member)}

\PYG{c+cp}{\PYGZsh{}}\PYG{c+cp}{define list\PYGZus{}entry\PYGZus{}is\PYGZus{}head(pos, head, member) \PYGZbs{}}
\PYG{c+cp}{    (\PYGZam{}pos\PYGZhy{}\PYGZgt{}member == (head))}

\PYG{c+cp}{\PYGZsh{}}\PYG{c+cp}{define list\PYGZus{}first\PYGZus{}entry(ptr, type, member) \PYGZbs{}}
\PYG{c+cp}{    list\PYGZus{}entry((ptr)\PYGZhy{}\PYGZgt{}next, type, member)}

\PYG{c+cp}{\PYGZsh{}}\PYG{c+cp}{define list\PYGZus{}next\PYGZus{}entry(pos, member) \PYGZbs{}}
\PYG{c+cp}{    list\PYGZus{}entry((pos)\PYGZhy{}\PYGZgt{}member.next, typeof(*(pos)), member)}

\PYG{c+cp}{\PYGZsh{}}\PYG{c+cp}{define list\PYGZus{}for\PYGZus{}each(pos, head) \PYGZbs{}}
\PYG{c+cp}{    for (pos = (head)\PYGZhy{}\PYGZgt{}next; !list\PYGZus{}is\PYGZus{}head(pos, (head)); pos = pos\PYGZhy{}\PYGZgt{}next)}

\PYG{c+cp}{\PYGZsh{}}\PYG{c+cp}{define list\PYGZus{}for\PYGZus{}each\PYGZus{}entry(pos, head, member)              \PYGZbs{}}
\PYG{c+cp}{    for (pos = list\PYGZus{}first\PYGZus{}entry(head, typeof(*pos), member); \PYGZbs{}}
\PYG{c+cp}{        !list\PYGZus{}entry\PYGZus{}is\PYGZus{}head(pos, head, member);              \PYGZbs{}}
\PYG{c+cp}{        pos = list\PYGZus{}next\PYGZus{}entry(pos, member))}

\PYG{c+cp}{\PYGZsh{}}\PYG{c+cp}{define list\PYGZus{}for\PYGZus{}each\PYGZus{}entry\PYGZus{}safe(pos, n, head, member)          \PYGZbs{}}
\PYG{c+cp}{    for (pos = list\PYGZus{}first\PYGZus{}entry(head, typeof(*pos), member),    \PYGZbs{}}
\PYG{c+cp}{         n = list\PYGZus{}next\PYGZus{}entry(pos, member);                       \PYGZbs{}}
\PYG{c+cp}{         !list\PYGZus{}entry\PYGZus{}is\PYGZus{}head(pos, head, member);                 \PYGZbs{}}
\PYG{c+cp}{         pos = n, n = list\PYGZus{}next\PYGZus{}entry(n, member))}

\PYG{c+cp}{\PYGZsh{}}\PYG{c+cp}{define my\PYGZus{}list\PYGZus{}for\PYGZus{}each\PYGZus{}entry(pos, head, member)                                        \PYGZbs{}}
\PYG{c+cp}{    for (pos = container\PYGZus{}of((head)\PYGZhy{}\PYGZgt{}next, typeof(*pos), member); \PYGZam{}(pos)\PYGZhy{}\PYGZgt{}member != head; \PYGZbs{}}
\PYG{c+cp}{            pos = container\PYGZus{}of((pos)\PYGZhy{}\PYGZgt{}member.next, typeof(*pos), member))}

\PYG{k}{static}\PYG{+w}{ }\PYG{k+kr}{inline}\PYG{+w}{ }\PYG{k+kt}{int}\PYG{+w}{ }\PYG{n+nf}{list\PYGZus{}is\PYGZus{}head}\PYG{p}{(}\PYG{k}{const}\PYG{+w}{ }\PYG{k}{struct}\PYG{+w}{ }\PYG{n+nc}{list\PYGZus{}head}\PYG{+w}{ }\PYG{o}{*}\PYG{n}{list}\PYG{p}{,}
\PYG{+w}{                               }\PYG{k}{const}\PYG{+w}{ }\PYG{k}{struct}\PYG{+w}{ }\PYG{n+nc}{list\PYGZus{}head}\PYG{+w}{ }\PYG{o}{*}\PYG{n}{head}\PYG{p}{)}
\PYG{p}{\PYGZob{}}
\PYG{+w}{    }\PYG{k}{return}\PYG{+w}{ }\PYG{n}{list}\PYG{+w}{ }\PYG{o}{=}\PYG{o}{=}\PYG{+w}{ }\PYG{n}{head}\PYG{p}{;}
\PYG{p}{\PYGZcb{}}

\PYG{k}{static}\PYG{+w}{ }\PYG{k+kr}{inline}\PYG{+w}{ }\PYG{k+kt}{void}\PYG{+w}{ }\PYG{n+nf}{\PYGZus{}\PYGZus{}list\PYGZus{}add}\PYG{p}{(}\PYG{k}{struct}\PYG{+w}{ }\PYG{n+nc}{list\PYGZus{}head}\PYG{+w}{ }\PYG{o}{*}\PYG{n}{new}\PYG{p}{,}
\PYG{+w}{                }\PYG{k}{struct}\PYG{+w}{ }\PYG{n+nc}{list\PYGZus{}head}\PYG{+w}{ }\PYG{o}{*}\PYG{n}{prev}\PYG{p}{,}
\PYG{+w}{                }\PYG{k}{struct}\PYG{+w}{ }\PYG{n+nc}{list\PYGZus{}head}\PYG{+w}{ }\PYG{o}{*}\PYG{n}{next}\PYG{p}{)}
\PYG{p}{\PYGZob{}}
\PYG{+w}{    }\PYG{n}{next}\PYG{o}{\PYGZhy{}}\PYG{o}{\PYGZgt{}}\PYG{n}{prev}\PYG{+w}{ }\PYG{o}{=}\PYG{+w}{ }\PYG{n}{new}\PYG{p}{;}
\PYG{+w}{    }\PYG{n}{new}\PYG{o}{\PYGZhy{}}\PYG{o}{\PYGZgt{}}\PYG{n}{next}\PYG{+w}{ }\PYG{o}{=}\PYG{+w}{ }\PYG{n}{next}\PYG{p}{;}
\PYG{+w}{    }\PYG{n}{new}\PYG{o}{\PYGZhy{}}\PYG{o}{\PYGZgt{}}\PYG{n}{prev}\PYG{+w}{ }\PYG{o}{=}\PYG{+w}{ }\PYG{n}{prev}\PYG{p}{;}
\PYG{+w}{    }\PYG{n}{prev}\PYG{o}{\PYGZhy{}}\PYG{o}{\PYGZgt{}}\PYG{n}{next}\PYG{+w}{ }\PYG{o}{=}\PYG{+w}{ }\PYG{n}{new}\PYG{p}{;}
\PYG{p}{\PYGZcb{}}

\PYG{k}{static}\PYG{+w}{ }\PYG{k+kr}{inline}\PYG{+w}{ }\PYG{k+kt}{void}\PYG{+w}{ }\PYG{n+nf}{list\PYGZus{}add}\PYG{p}{(}\PYG{k}{struct}\PYG{+w}{ }\PYG{n+nc}{list\PYGZus{}head}\PYG{+w}{ }\PYG{o}{*}\PYG{n}{new}\PYG{p}{,}\PYG{+w}{ }\PYG{k}{struct}\PYG{+w}{ }\PYG{n+nc}{list\PYGZus{}head}\PYG{+w}{ }\PYG{o}{*}\PYG{n}{head}\PYG{p}{)}
\PYG{p}{\PYGZob{}}
\PYG{+w}{    }\PYG{n}{\PYGZus{}\PYGZus{}list\PYGZus{}add}\PYG{p}{(}\PYG{n}{new}\PYG{p}{,}\PYG{+w}{ }\PYG{n}{head}\PYG{p}{,}\PYG{+w}{ }\PYG{n}{head}\PYG{o}{\PYGZhy{}}\PYG{o}{\PYGZgt{}}\PYG{n}{next}\PYG{p}{)}\PYG{p}{;}
\PYG{p}{\PYGZcb{}}

\PYG{k}{static}\PYG{+w}{ }\PYG{k+kr}{inline}\PYG{+w}{ }\PYG{k+kt}{void}\PYG{+w}{ }\PYG{n+nf}{list\PYGZus{}add\PYGZus{}tail}\PYG{p}{(}\PYG{k}{struct}\PYG{+w}{ }\PYG{n+nc}{list\PYGZus{}head}\PYG{+w}{ }\PYG{o}{*}\PYG{n}{new}\PYG{p}{,}\PYG{+w}{ }\PYG{k}{struct}\PYG{+w}{ }\PYG{n+nc}{list\PYGZus{}head}\PYG{+w}{ }\PYG{o}{*}\PYG{n}{head}\PYG{p}{)}
\PYG{p}{\PYGZob{}}
\PYG{+w}{    }\PYG{n}{\PYGZus{}\PYGZus{}list\PYGZus{}add}\PYG{p}{(}\PYG{n}{new}\PYG{p}{,}\PYG{+w}{ }\PYG{n}{head}\PYG{o}{\PYGZhy{}}\PYG{o}{\PYGZgt{}}\PYG{n}{prev}\PYG{p}{,}\PYG{+w}{ }\PYG{n}{head}\PYG{p}{)}\PYG{p}{;}
\PYG{p}{\PYGZcb{}}

\PYG{k}{static}\PYG{+w}{ }\PYG{k+kr}{inline}\PYG{+w}{ }\PYG{k+kt}{void}\PYG{+w}{ }\PYG{n+nf}{\PYGZus{}\PYGZus{}list\PYGZus{}del}\PYG{p}{(}\PYG{k}{struct}\PYG{+w}{ }\PYG{n+nc}{list\PYGZus{}head}\PYG{+w}{ }\PYG{o}{*}\PYG{n}{prev}\PYG{p}{,}\PYG{+w}{ }\PYG{k}{struct}\PYG{+w}{ }\PYG{n+nc}{list\PYGZus{}head}\PYG{+w}{ }\PYG{o}{*}\PYG{n}{next}\PYG{p}{)}
\PYG{p}{\PYGZob{}}
\PYG{+w}{    }\PYG{n}{next}\PYG{o}{\PYGZhy{}}\PYG{o}{\PYGZgt{}}\PYG{n}{prev}\PYG{+w}{ }\PYG{o}{=}\PYG{+w}{ }\PYG{n}{prev}\PYG{p}{;}
\PYG{+w}{    }\PYG{n}{prev}\PYG{o}{\PYGZhy{}}\PYG{o}{\PYGZgt{}}\PYG{n}{next}\PYG{+w}{ }\PYG{o}{=}\PYG{+w}{ }\PYG{n}{next}\PYG{p}{;}
\PYG{p}{\PYGZcb{}}

\PYG{k}{static}\PYG{+w}{ }\PYG{k+kr}{inline}\PYG{+w}{ }\PYG{k+kt}{void}\PYG{+w}{ }\PYG{n+nf}{\PYGZus{}\PYGZus{}list\PYGZus{}del\PYGZus{}entry}\PYG{p}{(}\PYG{k}{struct}\PYG{+w}{ }\PYG{n+nc}{list\PYGZus{}head}\PYG{+w}{ }\PYG{o}{*}\PYG{n}{entry}\PYG{p}{)}
\PYG{p}{\PYGZob{}}
\PYG{+w}{    }\PYG{n}{\PYGZus{}\PYGZus{}list\PYGZus{}del}\PYG{p}{(}\PYG{n}{entry}\PYG{o}{\PYGZhy{}}\PYG{o}{\PYGZgt{}}\PYG{n}{prev}\PYG{p}{,}\PYG{+w}{ }\PYG{n}{entry}\PYG{o}{\PYGZhy{}}\PYG{o}{\PYGZgt{}}\PYG{n}{next}\PYG{p}{)}\PYG{p}{;}
\PYG{p}{\PYGZcb{}}

\PYG{k}{static}\PYG{+w}{ }\PYG{k+kr}{inline}\PYG{+w}{ }\PYG{k+kt}{void}\PYG{+w}{ }\PYG{n+nf}{list\PYGZus{}del}\PYG{p}{(}\PYG{k}{struct}\PYG{+w}{ }\PYG{n+nc}{list\PYGZus{}head}\PYG{+w}{ }\PYG{o}{*}\PYG{n}{entry}\PYG{p}{)}
\PYG{p}{\PYGZob{}}
\PYG{+w}{    }\PYG{n}{\PYGZus{}\PYGZus{}list\PYGZus{}del\PYGZus{}entry}\PYG{p}{(}\PYG{n}{entry}\PYG{p}{)}\PYG{p}{;}
\PYG{p}{\PYGZcb{}}

\PYG{n}{LIST\PYGZus{}HEAD}\PYG{p}{(}\PYG{n}{head}\PYG{p}{)}\PYG{p}{;}

\PYG{k}{struct}\PYG{+w}{ }\PYG{n+nc}{SAMPLE}\PYG{+w}{ }\PYG{p}{\PYGZob{}}
\PYG{+w}{    }\PYG{k+kt}{int}\PYG{+w}{ }\PYG{n}{a}\PYG{p}{;}
\PYG{+w}{    }\PYG{k}{struct}\PYG{+w}{ }\PYG{n+nc}{list\PYGZus{}head}\PYG{+w}{ }\PYG{n}{link}\PYG{p}{;}
\PYG{p}{\PYGZcb{}}\PYG{p}{;}

\PYG{k+kt}{int}\PYG{+w}{ }\PYG{n+nf}{main}\PYG{p}{(}\PYG{k+kt}{void}\PYG{p}{)}
\PYG{p}{\PYGZob{}}
\PYG{+w}{    }\PYG{k}{struct}\PYG{+w}{ }\PYG{n+nc}{SAMPLE}\PYG{+w}{ }\PYG{o}{*}\PYG{n}{ps}\PYG{p}{,}\PYG{+w}{ }\PYG{o}{*}\PYG{n}{ps\PYGZus{}node}\PYG{p}{,}\PYG{+w}{ }\PYG{o}{*}\PYG{n}{ps\PYGZus{}temp}\PYG{p}{,}\PYG{+w}{ }\PYG{o}{*}\PYG{n}{ps\PYGZus{}del}\PYG{p}{,}\PYG{+w}{ }\PYG{o}{*}\PYG{n}{ps\PYGZus{}insert}\PYG{p}{;}
\PYG{+w}{    }\PYG{k}{struct}\PYG{+w}{ }\PYG{n+nc}{list\PYGZus{}head}\PYG{+w}{ }\PYG{o}{*}\PYG{n}{lh}\PYG{p}{;}

\PYG{+w}{    }\PYG{k}{for}\PYG{+w}{ }\PYG{p}{(}\PYG{k+kt}{int}\PYG{+w}{ }\PYG{n}{i}\PYG{+w}{ }\PYG{o}{=}\PYG{+w}{ }\PYG{l+m+mi}{0}\PYG{p}{;}\PYG{+w}{ }\PYG{n}{i}\PYG{+w}{ }\PYG{o}{\PYGZlt{}}\PYG{+w}{ }\PYG{l+m+mi}{10}\PYG{p}{;}\PYG{+w}{ }\PYG{o}{+}\PYG{o}{+}\PYG{n}{i}\PYG{p}{)}\PYG{+w}{ }\PYG{p}{\PYGZob{}}
\PYG{+w}{        }\PYG{k}{if}\PYG{+w}{ }\PYG{p}{(}\PYG{p}{(}\PYG{n}{ps}\PYG{+w}{ }\PYG{o}{=}\PYG{+w}{ }\PYG{p}{(}\PYG{k}{struct}\PYG{+w}{ }\PYG{n+nc}{SAMPLE}\PYG{+w}{ }\PYG{o}{*}\PYG{p}{)}\PYG{n}{malloc}\PYG{p}{(}\PYG{k}{sizeof}\PYG{p}{(}\PYG{k}{struct}\PYG{+w}{ }\PYG{n+nc}{SAMPLE}\PYG{p}{)}\PYG{p}{)}\PYG{p}{)}\PYG{+w}{ }\PYG{o}{=}\PYG{o}{=}\PYG{+w}{ }\PYG{n+nb}{NULL}\PYG{p}{)}\PYG{+w}{ }\PYG{p}{\PYGZob{}}
\PYG{+w}{            }\PYG{n}{fprintf}\PYG{p}{(}\PYG{n}{stderr}\PYG{p}{,}\PYG{+w}{ }\PYG{l+s}{\PYGZdq{}}\PYG{l+s}{cannot allocate memory!...}\PYG{l+s+se}{\PYGZbs{}n}\PYG{l+s}{\PYGZdq{}}\PYG{p}{)}\PYG{p}{;}
\PYG{+w}{            }\PYG{n}{exit}\PYG{p}{(}\PYG{n}{EXIT\PYGZus{}FAILURE}\PYG{p}{)}\PYG{p}{;}
\PYG{+w}{        }\PYG{p}{\PYGZcb{}}
\PYG{+w}{        }\PYG{n}{ps}\PYG{o}{\PYGZhy{}}\PYG{o}{\PYGZgt{}}\PYG{n}{a}\PYG{+w}{ }\PYG{o}{=}\PYG{+w}{ }\PYG{n}{i}\PYG{p}{;}
\PYG{+w}{        }\PYG{n}{list\PYGZus{}add\PYGZus{}tail}\PYG{p}{(}\PYG{o}{\PYGZam{}}\PYG{n}{ps}\PYG{o}{\PYGZhy{}}\PYG{o}{\PYGZgt{}}\PYG{n}{link}\PYG{p}{,}\PYG{+w}{ }\PYG{o}{\PYGZam{}}\PYG{n}{head}\PYG{p}{)}\PYG{p}{;}

\PYG{+w}{        }\PYG{k}{if}\PYG{+w}{ }\PYG{p}{(}\PYG{n}{i}\PYG{+w}{ }\PYG{o}{=}\PYG{o}{=}\PYG{+w}{ }\PYG{l+m+mi}{5}\PYG{p}{)}
\PYG{+w}{            }\PYG{n}{ps\PYGZus{}del}\PYG{+w}{ }\PYG{o}{=}\PYG{+w}{ }\PYG{n}{ps}\PYG{p}{;}

\PYG{+w}{        }\PYG{k}{if}\PYG{+w}{ }\PYG{p}{(}\PYG{n}{i}\PYG{+w}{ }\PYG{o}{=}\PYG{o}{=}\PYG{+w}{ }\PYG{l+m+mi}{7}\PYG{p}{)}
\PYG{+w}{            }\PYG{n}{ps\PYGZus{}insert}\PYG{+w}{ }\PYG{o}{=}\PYG{+w}{ }\PYG{n}{ps}\PYG{p}{;}
\PYG{+w}{    }\PYG{p}{\PYGZcb{}}

\PYG{+w}{    }\PYG{n}{list\PYGZus{}for\PYGZus{}each}\PYG{p}{(}\PYG{n}{lh}\PYG{p}{,}\PYG{+w}{ }\PYG{o}{\PYGZam{}}\PYG{n}{head}\PYG{p}{)}\PYG{+w}{ }\PYG{p}{\PYGZob{}}
\PYG{+w}{        }\PYG{n}{ps}\PYG{+w}{ }\PYG{o}{=}\PYG{+w}{ }\PYG{n}{container\PYGZus{}of}\PYG{p}{(}\PYG{n}{lh}\PYG{p}{,}\PYG{+w}{ }\PYG{k}{struct}\PYG{+w}{ }\PYG{n+nc}{SAMPLE}\PYG{p}{,}\PYG{+w}{ }\PYG{n}{link}\PYG{p}{)}\PYG{p}{;}
\PYG{+w}{        }\PYG{n}{printf}\PYG{p}{(}\PYG{l+s}{\PYGZdq{}}\PYG{l+s}{\PYGZpc{}d }\PYG{l+s}{\PYGZdq{}}\PYG{p}{,}\PYG{+w}{ }\PYG{n}{ps}\PYG{o}{\PYGZhy{}}\PYG{o}{\PYGZgt{}}\PYG{n}{a}\PYG{p}{)}\PYG{p}{;}
\PYG{+w}{    }\PYG{p}{\PYGZcb{}}
\PYG{+w}{    }\PYG{n}{printf}\PYG{p}{(}\PYG{l+s}{\PYGZdq{}}\PYG{l+s+se}{\PYGZbs{}n}\PYG{l+s}{\PYGZdq{}}\PYG{p}{)}\PYG{p}{;}

\PYG{+w}{    }\PYG{n}{list\PYGZus{}del}\PYG{p}{(}\PYG{o}{\PYGZam{}}\PYG{n}{ps\PYGZus{}del}\PYG{o}{\PYGZhy{}}\PYG{o}{\PYGZgt{}}\PYG{n}{link}\PYG{p}{)}\PYG{p}{;}

\PYG{+w}{    }\PYG{n}{list\PYGZus{}for\PYGZus{}each}\PYG{p}{(}\PYG{n}{lh}\PYG{p}{,}\PYG{+w}{ }\PYG{o}{\PYGZam{}}\PYG{n}{head}\PYG{p}{)}\PYG{+w}{ }\PYG{p}{\PYGZob{}}
\PYG{+w}{        }\PYG{n}{ps}\PYG{+w}{ }\PYG{o}{=}\PYG{+w}{ }\PYG{n}{container\PYGZus{}of}\PYG{p}{(}\PYG{n}{lh}\PYG{p}{,}\PYG{+w}{ }\PYG{k}{struct}\PYG{+w}{ }\PYG{n+nc}{SAMPLE}\PYG{p}{,}\PYG{+w}{ }\PYG{n}{link}\PYG{p}{)}\PYG{p}{;}
\PYG{+w}{        }\PYG{n}{printf}\PYG{p}{(}\PYG{l+s}{\PYGZdq{}}\PYG{l+s}{\PYGZpc{}d }\PYG{l+s}{\PYGZdq{}}\PYG{p}{,}\PYG{+w}{ }\PYG{n}{ps}\PYG{o}{\PYGZhy{}}\PYG{o}{\PYGZgt{}}\PYG{n}{a}\PYG{p}{)}\PYG{p}{;}
\PYG{+w}{    }\PYG{p}{\PYGZcb{}}
\PYG{+w}{    }\PYG{n}{printf}\PYG{p}{(}\PYG{l+s}{\PYGZdq{}}\PYG{l+s+se}{\PYGZbs{}n}\PYG{l+s}{\PYGZdq{}}\PYG{p}{)}\PYG{p}{;}

\PYG{+w}{    }\PYG{k}{if}\PYG{+w}{ }\PYG{p}{(}\PYG{p}{(}\PYG{n}{ps}\PYG{+w}{ }\PYG{o}{=}\PYG{+w}{ }\PYG{p}{(}\PYG{k}{struct}\PYG{+w}{ }\PYG{n+nc}{SAMPLE}\PYG{+w}{ }\PYG{o}{*}\PYG{p}{)}\PYG{n}{malloc}\PYG{p}{(}\PYG{k}{sizeof}\PYG{p}{(}\PYG{k}{struct}\PYG{+w}{ }\PYG{n+nc}{SAMPLE}\PYG{p}{)}\PYG{p}{)}\PYG{p}{)}\PYG{+w}{ }\PYG{o}{=}\PYG{o}{=}\PYG{+w}{ }\PYG{n+nb}{NULL}\PYG{p}{)}\PYG{+w}{ }\PYG{p}{\PYGZob{}}
\PYG{+w}{        }\PYG{n}{fprintf}\PYG{p}{(}\PYG{n}{stderr}\PYG{p}{,}\PYG{+w}{ }\PYG{l+s}{\PYGZdq{}}\PYG{l+s}{cannot allocate memory!...}\PYG{l+s+se}{\PYGZbs{}n}\PYG{l+s}{\PYGZdq{}}\PYG{p}{)}\PYG{p}{;}
\PYG{+w}{        }\PYG{n}{exit}\PYG{p}{(}\PYG{n}{EXIT\PYGZus{}FAILURE}\PYG{p}{)}\PYG{p}{;}
\PYG{+w}{    }\PYG{p}{\PYGZcb{}}
\PYG{+w}{    }\PYG{n}{ps}\PYG{o}{\PYGZhy{}}\PYG{o}{\PYGZgt{}}\PYG{n}{a}\PYG{+w}{ }\PYG{o}{=}\PYG{+w}{ }\PYG{l+m+mi}{100}\PYG{p}{;}

\PYG{+w}{    }\PYG{n}{list\PYGZus{}add}\PYG{p}{(}\PYG{o}{\PYGZam{}}\PYG{n}{ps}\PYG{o}{\PYGZhy{}}\PYG{o}{\PYGZgt{}}\PYG{n}{link}\PYG{p}{,}\PYG{+w}{ }\PYG{o}{\PYGZam{}}\PYG{n}{ps\PYGZus{}insert}\PYG{o}{\PYGZhy{}}\PYG{o}{\PYGZgt{}}\PYG{n}{link}\PYG{p}{)}\PYG{p}{;}

\PYG{+w}{    }\PYG{n}{list\PYGZus{}for\PYGZus{}each}\PYG{p}{(}\PYG{n}{lh}\PYG{p}{,}\PYG{+w}{ }\PYG{o}{\PYGZam{}}\PYG{n}{head}\PYG{p}{)}\PYG{+w}{ }\PYG{p}{\PYGZob{}}
\PYG{+w}{        }\PYG{n}{ps}\PYG{+w}{ }\PYG{o}{=}\PYG{+w}{ }\PYG{n}{container\PYGZus{}of}\PYG{p}{(}\PYG{n}{lh}\PYG{p}{,}\PYG{+w}{ }\PYG{k}{struct}\PYG{+w}{ }\PYG{n+nc}{SAMPLE}\PYG{p}{,}\PYG{+w}{ }\PYG{n}{link}\PYG{p}{)}\PYG{p}{;}
\PYG{+w}{        }\PYG{n}{printf}\PYG{p}{(}\PYG{l+s}{\PYGZdq{}}\PYG{l+s}{\PYGZpc{}d }\PYG{l+s}{\PYGZdq{}}\PYG{p}{,}\PYG{+w}{ }\PYG{n}{ps}\PYG{o}{\PYGZhy{}}\PYG{o}{\PYGZgt{}}\PYG{n}{a}\PYG{p}{)}\PYG{p}{;}
\PYG{+w}{    }\PYG{p}{\PYGZcb{}}
\PYG{+w}{    }\PYG{n}{printf}\PYG{p}{(}\PYG{l+s}{\PYGZdq{}}\PYG{l+s+se}{\PYGZbs{}n}\PYG{l+s}{\PYGZdq{}}\PYG{p}{)}\PYG{p}{;}

\PYG{+w}{    }\PYG{n}{ps\PYGZus{}temp}\PYG{+w}{ }\PYG{o}{=}\PYG{+w}{ }\PYG{n+nb}{NULL}\PYG{p}{;}
\PYG{+w}{    }\PYG{n}{list\PYGZus{}for\PYGZus{}each\PYGZus{}entry\PYGZus{}safe}\PYG{p}{(}\PYG{n}{ps}\PYG{p}{,}\PYG{+w}{ }\PYG{n}{ps\PYGZus{}node}\PYG{p}{,}\PYG{+w}{ }\PYG{o}{\PYGZam{}}\PYG{n}{head}\PYG{p}{,}\PYG{+w}{ }\PYG{n}{link}\PYG{p}{)}\PYG{+w}{ }\PYG{p}{\PYGZob{}}
\PYG{+w}{        }\PYG{n}{free}\PYG{p}{(}\PYG{n}{ps\PYGZus{}temp}\PYG{p}{)}\PYG{p}{;}
\PYG{+w}{        }\PYG{n}{ps\PYGZus{}temp}\PYG{+w}{ }\PYG{o}{=}\PYG{+w}{ }\PYG{n}{ps}\PYG{p}{;}
\PYG{+w}{    }\PYG{p}{\PYGZcb{}}

\PYG{+w}{    }\PYG{k}{return}\PYG{+w}{ }\PYG{l+m+mi}{0}\PYG{p}{;}
\PYG{p}{\PYGZcb{}}
\end{sphinxVerbatim}


\section{RCU Desteği}
\label{\detokenize{kernellinkedlistandhashtable:rcu-destegi}}
\sphinxAtStartPar
Belli bir süreden sonra \sphinxcode{\sphinxupquote{list.h}} dosyasına RCU (Read\sphinxhyphen{}Copy\sphinxhyphen{}Update) mekanizmasını destekleyecek biçimde
dolaşım yapan \sphinxcode{\sphinxupquote{list\_for\_each\_rcu}} isimli makro da eklenmiştir. Bu makro bir yazıcı varsa okuyucuları
bekletmeden işlem yapabilmeyi sağlamaktadır. \sphinxstyleemphasis{Read\sphinxhyphen{}Copy\sphinxhyphen{}Update} mekanizması ileride ele alınacaktır.


\section{Eski Çekirdek Sürümlerindeki list.h}
\label{\detokenize{kernellinkedlistandhashtable:eski-cekirdek-surumlerindeki-list-h}}
\sphinxAtStartPar
Çekirdeğin eski sürümlerinde \sphinxcode{\sphinxupquote{list.h}} dosyası oldukça sade idi. Sonra bu dosyanın içeriğinde de bazı
faydalı değişiklikler yapıldı. Ancak bu değişiklikler veri yapısının anlaşılmasını da biraz zorlaştırmıştır.
Örneğin 2.2.26 çekirdeğindeki \sphinxcode{\sphinxupquote{list.h}} dosyası oldukça sade bir biçimde aşağıdaki gibi oluşturulmuştur:

\begin{sphinxVerbatim}[commandchars=\\\{\}]
\PYG{c+cm}{/* 2.2.26 \PYGZlt{}list.h\PYGZgt{} dosyasının içeriği */}

\PYG{c+cp}{\PYGZsh{}}\PYG{c+cp}{ifndef \PYGZus{}LINUX\PYGZus{}LIST\PYGZus{}H}
\PYG{c+cp}{\PYGZsh{}}\PYG{c+cp}{define \PYGZus{}LINUX\PYGZus{}LIST\PYGZus{}H}

\PYG{c+cp}{\PYGZsh{}}\PYG{c+cp}{ifdef \PYGZus{}\PYGZus{}KERNEL\PYGZus{}\PYGZus{}}

\PYG{c+cm}{/*}
\PYG{c+cm}{ * Simple doubly linked list implementation.}
\PYG{c+cm}{ *}
\PYG{c+cm}{ * Some of the internal functions (\PYGZdq{}\PYGZus{}\PYGZus{}xxx\PYGZdq{}) are useful when}
\PYG{c+cm}{ * manipulating whole lists rather than single entries, as}
\PYG{c+cm}{ * sometimes we already know the next/prev entries and we can}
\PYG{c+cm}{ * generate better code by using them directly rather than}
\PYG{c+cm}{ * using the generic single\PYGZhy{}entry routines.}
\PYG{c+cm}{ */}

\PYG{k}{struct}\PYG{+w}{ }\PYG{n+nc}{list\PYGZus{}head}\PYG{+w}{ }\PYG{p}{\PYGZob{}}
\PYG{+w}{    }\PYG{k}{struct}\PYG{+w}{ }\PYG{n+nc}{list\PYGZus{}head}\PYG{+w}{ }\PYG{o}{*}\PYG{n}{next}\PYG{p}{,}\PYG{+w}{ }\PYG{o}{*}\PYG{n}{prev}\PYG{p}{;}
\PYG{p}{\PYGZcb{}}\PYG{p}{;}

\PYG{c+cp}{\PYGZsh{}}\PYG{c+cp}{define LIST\PYGZus{}HEAD\PYGZus{}INIT(name) \PYGZob{} \PYGZam{}(name), \PYGZam{}(name) \PYGZcb{}}

\PYG{c+cp}{\PYGZsh{}}\PYG{c+cp}{define LIST\PYGZus{}HEAD(name) \PYGZbs{}}
\PYG{c+cp}{    struct list\PYGZus{}head name = \PYGZob{} \PYGZam{}name, \PYGZam{}name \PYGZcb{}}

\PYG{c+cp}{\PYGZsh{}}\PYG{c+cp}{define INIT\PYGZus{}LIST\PYGZus{}HEAD(ptr) do \PYGZob{} \PYGZbs{}}
\PYG{c+cp}{    (ptr)\PYGZhy{}\PYGZgt{}next = (ptr); (ptr)\PYGZhy{}\PYGZgt{}prev = (ptr); \PYGZbs{}}
\PYG{c+cp}{\PYGZcb{} while (0)}

\PYG{k}{static}\PYG{+w}{ }\PYG{n}{\PYGZus{}\PYGZus{}inline\PYGZus{}\PYGZus{}}\PYG{+w}{ }\PYG{k+kt}{void}\PYG{+w}{ }\PYG{n}{\PYGZus{}\PYGZus{}list\PYGZus{}add}\PYG{p}{(}\PYG{k}{struct}\PYG{+w}{ }\PYG{n+nc}{list\PYGZus{}head}\PYG{+w}{ }\PYG{o}{*}\PYG{n}{new}\PYG{p}{,}
\PYG{+w}{    }\PYG{k}{struct}\PYG{+w}{ }\PYG{n+nc}{list\PYGZus{}head}\PYG{+w}{ }\PYG{o}{*}\PYG{n}{prev}\PYG{p}{,}
\PYG{+w}{    }\PYG{k}{struct}\PYG{+w}{ }\PYG{n+nc}{list\PYGZus{}head}\PYG{+w}{ }\PYG{o}{*}\PYG{n}{next}\PYG{p}{)}
\PYG{p}{\PYGZob{}}
\PYG{+w}{    }\PYG{n}{next}\PYG{o}{\PYGZhy{}}\PYG{o}{\PYGZgt{}}\PYG{n}{prev}\PYG{+w}{ }\PYG{o}{=}\PYG{+w}{ }\PYG{n}{new}\PYG{p}{;}
\PYG{+w}{    }\PYG{n}{new}\PYG{o}{\PYGZhy{}}\PYG{o}{\PYGZgt{}}\PYG{n}{next}\PYG{+w}{ }\PYG{o}{=}\PYG{+w}{ }\PYG{n}{next}\PYG{p}{;}
\PYG{+w}{    }\PYG{n}{new}\PYG{o}{\PYGZhy{}}\PYG{o}{\PYGZgt{}}\PYG{n}{prev}\PYG{+w}{ }\PYG{o}{=}\PYG{+w}{ }\PYG{n}{prev}\PYG{p}{;}
\PYG{+w}{    }\PYG{n}{prev}\PYG{o}{\PYGZhy{}}\PYG{o}{\PYGZgt{}}\PYG{n}{next}\PYG{+w}{ }\PYG{o}{=}\PYG{+w}{ }\PYG{n}{new}\PYG{p}{;}
\PYG{p}{\PYGZcb{}}

\PYG{k}{static}\PYG{+w}{ }\PYG{n}{\PYGZus{}\PYGZus{}inline\PYGZus{}\PYGZus{}}\PYG{+w}{ }\PYG{k+kt}{void}\PYG{+w}{ }\PYG{n}{list\PYGZus{}add}\PYG{p}{(}\PYG{k}{struct}\PYG{+w}{ }\PYG{n+nc}{list\PYGZus{}head}\PYG{+w}{ }\PYG{o}{*}\PYG{n}{new}\PYG{p}{,}\PYG{+w}{ }\PYG{k}{struct}\PYG{+w}{ }\PYG{n+nc}{list\PYGZus{}head}\PYG{+w}{ }\PYG{o}{*}\PYG{n}{head}\PYG{p}{)}
\PYG{p}{\PYGZob{}}
\PYG{+w}{    }\PYG{n}{\PYGZus{}\PYGZus{}list\PYGZus{}add}\PYG{p}{(}\PYG{n}{new}\PYG{p}{,}\PYG{+w}{ }\PYG{n}{head}\PYG{p}{,}\PYG{+w}{ }\PYG{n}{head}\PYG{o}{\PYGZhy{}}\PYG{o}{\PYGZgt{}}\PYG{n}{next}\PYG{p}{)}\PYG{p}{;}
\PYG{p}{\PYGZcb{}}

\PYG{k}{static}\PYG{+w}{ }\PYG{n}{\PYGZus{}\PYGZus{}inline\PYGZus{}\PYGZus{}}\PYG{+w}{ }\PYG{k+kt}{void}\PYG{+w}{ }\PYG{n}{list\PYGZus{}add\PYGZus{}tail}\PYG{p}{(}\PYG{k}{struct}\PYG{+w}{ }\PYG{n+nc}{list\PYGZus{}head}\PYG{+w}{ }\PYG{o}{*}\PYG{n}{new}\PYG{p}{,}\PYG{+w}{ }\PYG{k}{struct}\PYG{+w}{ }\PYG{n+nc}{list\PYGZus{}head}\PYG{+w}{ }\PYG{o}{*}\PYG{n}{head}\PYG{p}{)}
\PYG{p}{\PYGZob{}}
\PYG{+w}{    }\PYG{n}{\PYGZus{}\PYGZus{}list\PYGZus{}add}\PYG{p}{(}\PYG{n}{new}\PYG{p}{,}\PYG{+w}{ }\PYG{n}{head}\PYG{o}{\PYGZhy{}}\PYG{o}{\PYGZgt{}}\PYG{n}{prev}\PYG{p}{,}\PYG{+w}{ }\PYG{n}{head}\PYG{p}{)}\PYG{p}{;}
\PYG{p}{\PYGZcb{}}

\PYG{k}{static}\PYG{+w}{ }\PYG{n}{\PYGZus{}\PYGZus{}inline\PYGZus{}\PYGZus{}}\PYG{+w}{ }\PYG{k+kt}{void}\PYG{+w}{ }\PYG{n}{\PYGZus{}\PYGZus{}list\PYGZus{}del}\PYG{p}{(}\PYG{k}{struct}\PYG{+w}{ }\PYG{n+nc}{list\PYGZus{}head}\PYG{+w}{ }\PYG{o}{*}\PYG{n}{prev}\PYG{p}{,}
\PYG{+w}{                }\PYG{k}{struct}\PYG{+w}{ }\PYG{n+nc}{list\PYGZus{}head}\PYG{+w}{ }\PYG{o}{*}\PYG{n}{next}\PYG{p}{)}
\PYG{p}{\PYGZob{}}
\PYG{+w}{    }\PYG{n}{next}\PYG{o}{\PYGZhy{}}\PYG{o}{\PYGZgt{}}\PYG{n}{prev}\PYG{+w}{ }\PYG{o}{=}\PYG{+w}{ }\PYG{n}{prev}\PYG{p}{;}
\PYG{+w}{    }\PYG{n}{prev}\PYG{o}{\PYGZhy{}}\PYG{o}{\PYGZgt{}}\PYG{n}{next}\PYG{+w}{ }\PYG{o}{=}\PYG{+w}{ }\PYG{n}{next}\PYG{p}{;}
\PYG{p}{\PYGZcb{}}

\PYG{k}{static}\PYG{+w}{ }\PYG{n}{\PYGZus{}\PYGZus{}inline\PYGZus{}\PYGZus{}}\PYG{+w}{ }\PYG{k+kt}{void}\PYG{+w}{ }\PYG{n}{list\PYGZus{}del}\PYG{p}{(}\PYG{k}{struct}\PYG{+w}{ }\PYG{n+nc}{list\PYGZus{}head}\PYG{+w}{ }\PYG{o}{*}\PYG{n}{entry}\PYG{p}{)}
\PYG{p}{\PYGZob{}}
\PYG{+w}{    }\PYG{n}{\PYGZus{}\PYGZus{}list\PYGZus{}del}\PYG{p}{(}\PYG{n}{entry}\PYG{o}{\PYGZhy{}}\PYG{o}{\PYGZgt{}}\PYG{n}{prev}\PYG{p}{,}\PYG{+w}{ }\PYG{n}{entry}\PYG{o}{\PYGZhy{}}\PYG{o}{\PYGZgt{}}\PYG{n}{next}\PYG{p}{)}\PYG{p}{;}
\PYG{p}{\PYGZcb{}}

\PYG{k}{static}\PYG{+w}{ }\PYG{n}{\PYGZus{}\PYGZus{}inline\PYGZus{}\PYGZus{}}\PYG{+w}{ }\PYG{k+kt}{int}\PYG{+w}{ }\PYG{n}{list\PYGZus{}empty}\PYG{p}{(}\PYG{k}{struct}\PYG{+w}{ }\PYG{n+nc}{list\PYGZus{}head}\PYG{+w}{ }\PYG{o}{*}\PYG{n}{head}\PYG{p}{)}
\PYG{p}{\PYGZob{}}
\PYG{+w}{    }\PYG{k}{return}\PYG{+w}{ }\PYG{n}{head}\PYG{o}{\PYGZhy{}}\PYG{o}{\PYGZgt{}}\PYG{n}{next}\PYG{+w}{ }\PYG{o}{=}\PYG{o}{=}\PYG{+w}{ }\PYG{n}{head}\PYG{p}{;}
\PYG{p}{\PYGZcb{}}

\PYG{k}{static}\PYG{+w}{ }\PYG{n}{\PYGZus{}\PYGZus{}inline\PYGZus{}\PYGZus{}}\PYG{+w}{ }\PYG{k+kt}{void}\PYG{+w}{ }\PYG{n}{list\PYGZus{}splice}\PYG{p}{(}\PYG{k}{struct}\PYG{+w}{ }\PYG{n+nc}{list\PYGZus{}head}\PYG{+w}{ }\PYG{o}{*}\PYG{n}{list}\PYG{p}{,}
\PYG{+w}{                                   }\PYG{k}{struct}\PYG{+w}{ }\PYG{n+nc}{list\PYGZus{}head}\PYG{+w}{ }\PYG{o}{*}\PYG{n}{head}\PYG{p}{)}
\PYG{p}{\PYGZob{}}
\PYG{+w}{    }\PYG{k}{struct}\PYG{+w}{ }\PYG{n+nc}{list\PYGZus{}head}\PYG{+w}{ }\PYG{o}{*}\PYG{n}{first}\PYG{+w}{ }\PYG{o}{=}\PYG{+w}{ }\PYG{n}{list}\PYG{o}{\PYGZhy{}}\PYG{o}{\PYGZgt{}}\PYG{n}{next}\PYG{p}{;}

\PYG{+w}{    }\PYG{k}{if}\PYG{+w}{ }\PYG{p}{(}\PYG{n}{first}\PYG{+w}{ }\PYG{o}{!}\PYG{o}{=}\PYG{+w}{ }\PYG{n}{list}\PYG{p}{)}\PYG{+w}{ }\PYG{p}{\PYGZob{}}
\PYG{+w}{        }\PYG{k}{struct}\PYG{+w}{ }\PYG{n+nc}{list\PYGZus{}head}\PYG{+w}{ }\PYG{o}{*}\PYG{n}{last}\PYG{+w}{ }\PYG{o}{=}\PYG{+w}{ }\PYG{n}{list}\PYG{o}{\PYGZhy{}}\PYG{o}{\PYGZgt{}}\PYG{n}{prev}\PYG{p}{;}
\PYG{+w}{        }\PYG{k}{struct}\PYG{+w}{ }\PYG{n+nc}{list\PYGZus{}head}\PYG{+w}{ }\PYG{o}{*}\PYG{n}{at}\PYG{+w}{ }\PYG{o}{=}\PYG{+w}{ }\PYG{n}{head}\PYG{o}{\PYGZhy{}}\PYG{o}{\PYGZgt{}}\PYG{n}{next}\PYG{p}{;}

\PYG{+w}{        }\PYG{n}{first}\PYG{o}{\PYGZhy{}}\PYG{o}{\PYGZgt{}}\PYG{n}{prev}\PYG{+w}{ }\PYG{o}{=}\PYG{+w}{ }\PYG{n}{head}\PYG{p}{;}
\PYG{+w}{        }\PYG{n}{head}\PYG{o}{\PYGZhy{}}\PYG{o}{\PYGZgt{}}\PYG{n}{next}\PYG{+w}{ }\PYG{o}{=}\PYG{+w}{ }\PYG{n}{first}\PYG{p}{;}

\PYG{+w}{        }\PYG{n}{last}\PYG{o}{\PYGZhy{}}\PYG{o}{\PYGZgt{}}\PYG{n}{next}\PYG{+w}{ }\PYG{o}{=}\PYG{+w}{ }\PYG{n}{at}\PYG{p}{;}
\PYG{+w}{        }\PYG{n}{at}\PYG{o}{\PYGZhy{}}\PYG{o}{\PYGZgt{}}\PYG{n}{prev}\PYG{+w}{ }\PYG{o}{=}\PYG{+w}{ }\PYG{n}{last}\PYG{p}{;}
\PYG{+w}{    }\PYG{p}{\PYGZcb{}}
\PYG{p}{\PYGZcb{}}

\PYG{c+cp}{\PYGZsh{}}\PYG{c+cp}{define list\PYGZus{}entry(ptr, type, member) \PYGZbs{}}
\PYG{c+cp}{    ((type *)((char *)(ptr)\PYGZhy{}(unsigned long)(\PYGZam{}((type *)0)\PYGZhy{}\PYGZgt{}member)))}

\PYG{c+cp}{\PYGZsh{}}\PYG{c+cp}{define list\PYGZus{}for\PYGZus{}each(pos, head) \PYGZbs{}}
\PYG{c+cp}{    for (pos = (head)\PYGZhy{}\PYGZgt{}next; pos != (head); pos = pos\PYGZhy{}\PYGZgt{}next)}

\PYG{c+cp}{\PYGZsh{}}\PYG{c+cp}{endif }\PYG{c+cm}{/* \PYGZus{}\PYGZus{}KERNEL\PYGZus{}\PYGZus{} */}

\PYG{c+cp}{\PYGZsh{}}\PYG{c+cp}{endif}
\end{sphinxVerbatim}


\section{Arama Algoritmaları ve Sözlük Tarzı Veri Yapıları}
\label{\detokenize{kernellinkedlistandhashtable:arama-algoritmalari-ve-sozluk-tarzi-veri-yapilari}}
\sphinxAtStartPar
Arama işlemleri bir anahtar (key) eşliğinde yapılmaktadır. Anahtar bulunması istenen varlığı
temsil etmektedir. Örneğin pek çok kişinin bilgileri bir veri yapısında tutuluyor olabilir.
Biz de ismini bildiğimiz bir kişiyi bu veri yapısında arayabiliriz. Burada anahtar isimdir.
Veri yapıları dünyasında anahtar\sphinxhyphen{}değer çiftlerini tutan ve anahtar verildiğinde ona karşı
gelen değerin elde edilmesini sağlayan veri yapılarına genel olarak \sphinxstyleemphasis{sözlük (dictionary)}
tarzı veri yapıları denilmektedir.

\sphinxAtStartPar
Elemanların sıralı olmadığı liste tarzı veri yapılarında elemanların tek tek gözden
geçirilmesi yoluyla yapılan aramalara \sphinxstyleemphasis{sıralı arama (sequential search)} denilmektedir.
Sıralı arama oldukça yavaş bir yöntemdir. Sıralı aramada listedeki eleman sayısı N olmak
üzere anahtarın bulunması için ortalama N/2 kez karşılaştırmanın yapılması gerekmektedir.
Bu tür algoritmaların karmaşıklığına \sphinxstyleemphasis{Big O} notasyonunda \sphinxstyleemphasis{O(N) karmaşıklık} denildiğini
anımsayınız. Ancak ne olursa olsun eleman sayısının makul olduğu bir durumda sıralı arama
en iyi yöntem haline de gelebilmektedir. Örneğin en fazla 20 civarında elemanın bulunduğu
bir durumda bu elemanları diziye yerleştirip sıralı bir biçimde aramak en etkin yöntem
haline gelebilmektedir.

\sphinxAtStartPar
Eğer elemanlar sıralıysa ve herhangi bir elemana erişim çok hızlı (buna \sphinxstyleemphasis{rastgele erişim}
de denilmektedir) yapılabiliyorsa \sphinxstyleemphasis{ikili arama (binary search)} en iyi yöntemdir. İkili
aramanın algoritmik karmaşıklığı O(log N) biçimindedir. Aslında ikili aramanın daha genel
bir biçimine \sphinxstyleemphasis{enterpolasyon araması (interpolation search)} denilmektedir. Enterpolasyon
aramasında bölme ortadan değil daha uygun yerlerden yapılmaktadır. Ancak enterpolasyon
araması dizi dağılımının bilindiği ve özellikle de dizinin düzgün dağıldığı durumlarda
faydalı bir etki oluşturmaktadır. Diziyi önce sıraya dizip sonra ikili arama uygulamak
ise genellikle iyi bir fikir değildir. Çünkü sırayı korumak için araya eleman ekleme ve
aradan eleman silme gibi işlemlerde O(N) karmaşıklıkta kaydırmaların yapılması
gerekmektedir.


\subsection{İndeksli Arama}
\label{\detokenize{kernellinkedlistandhashtable:indeksli-arama}}
\sphinxAtStartPar
Peki ideal bir anahtar\sphinxhyphen{}değer araması nasıl olabilir? Şüphesiz ideal durumda aramanın
O(1) karmaşıklıkta yapılması istenir. Anahtarın bir \sphinxcode{\sphinxupquote{int}} değer olduğunu ve kişinin
numarasını belirttiğini düşünelim. Biz de numarasını bildiğimiz kişinin bilgilerini elde
etmek isteyelim. Örneğimizdeki kişilerin bilgileri \sphinxcode{\sphinxupquote{PERSON}} isimli bir yapıyla temsil
edilmiş olsun:

\begin{sphinxVerbatim}[commandchars=\\\{\}]
\PYG{k}{struct}\PYG{+w}{ }\PYG{n+nc}{PERSON}\PYG{+w}{ }\PYG{p}{\PYGZob{}}
\PYG{+w}{    }\PYG{p}{.}\PYG{p}{.}\PYG{p}{.}
\PYG{p}{\PYGZcb{}}\PYG{p}{;}
\end{sphinxVerbatim}

\sphinxAtStartPar
\sphinxcode{\sphinxupquote{struct PERSON}} türünden büyük bir dizi açabiliriz:

\begin{sphinxVerbatim}[commandchars=\\\{\}]
\PYG{k}{struct}\PYG{+w}{ }\PYG{n+nc}{PERSON}\PYG{+w}{ }\PYG{n}{people}\PYG{p}{[}\PYG{n}{MAX\PYGZus{}SIZE}\PYG{p}{]}\PYG{p}{;}
\end{sphinxVerbatim}

\sphinxAtStartPar
Sonra da kişilerin numaralarını indeks yaparak bu diziye yerleştirebiliriz. Örneğin
numarası 123 olan kişinin bilgilerini diziye şöyle yerleştirebiliriz:

\begin{sphinxVerbatim}[commandchars=\\\{\}]
\PYG{n}{people}\PYG{p}{[}\PYG{l+m+mi}{123}\PYG{p}{]}\PYG{+w}{ }\PYG{o}{=}\PYG{+w}{ }\PYG{n}{person\PYGZus{}info}\PYG{p}{;}
\end{sphinxVerbatim}

\sphinxAtStartPar
Artık numarası 123 olan bu kişinin bilgilerini O(1) karmaşıklıkta aşağıdaki gibi elde
edebiliriz:

\begin{sphinxVerbatim}[commandchars=\\\{\}]
\PYG{n}{person\PYGZus{}info}\PYG{+w}{ }\PYG{o}{=}\PYG{+w}{ }\PYG{n}{people}\PYG{p}{[}\PYG{l+m+mi}{123}\PYG{p}{]}\PYG{p}{;}
\end{sphinxVerbatim}

\sphinxAtStartPar
Bu yöntem ilk bakışta çok iyi bir yöntem gibi gözükse de genellikle kullanılabilir bir
yöntem değildir. Çünkü burada anahtar \sphinxcode{\sphinxupquote{int}} türdendir. Ancak uygulamalarda anahtarlar
farklı türlerden olabilmektedir. Örneğin anahtar kişinin adı soyadı olabilir. Ad ve soyad
gibi yazısal bilgiler indeks belirtmemektedir. Bu yöntemin diğer bir sakıncası da
anahtarların yüksek basamaklı sayılardan oluşabildiği durumlarda dizilerin çok fazla yer
kaplamasıdır. Örneğin TC kimlik numarasının anahtar yapılarak kişilerin bilgilerinin elde
edilmesinin istendiği bir durumu düşünelim. TC kimlik numarası 11 basamaklı bir sayıdır.
Yani skalası 100 milyarlık sınırdadır. 100 milyarlık bir yapı dizisini bu amaçla oluşturmak
mümkün olmayabilir, mümkün olsa da etkin olmayabilir. Bu yönteme \sphinxstyleemphasis{indeksli arama (index
search)} denilmektedir. Ancak çok özel durumlarda bu yöntem uygun bir yöntem olarak
kullanılabilmektedir.


\section{Hash Tabloları}
\label{\detokenize{kernellinkedlistandhashtable:hash-tablolari}}
\sphinxAtStartPar
Algoritmik aramalarda en çok kullanılan yöntemlerden biri \sphinxstyleemphasis{hash tabloları (hash tables)}
denilen yöntemdir. Hash tabloları aslında yukarıda belirttiğimiz indeksli arama ile sıralı
aramanın hibrit bir biçimi gibidir. Yöntemde ismine \sphinxstyleemphasis{hash tablosu (hash table)} denilen
makul uzunlukta bir dizi oluşturulur. Sonra anahtarlar ismine \sphinxstyleemphasis{hash fonksiyonu (hash
function)} denilen bir fonksiyona sokularak dizi indeksine dönüştürülür. Sonra da dizinin
o indeksteki elemanına başvurulur. Örneğin kişilerin bilgilerini TC kimlik numaralarına
göre saklayıp geri almak isteyelim. Hash tablomuzun uzunluğu da 1000 olsun. Hash
fonksiyonumuzun da “1000’e bölümden elde edilen kalan” değerini veren fonksiyon olduğunu
varsayalım. Bu durumda örneğin 2566198712 TC kimlik numarasına sahip kişinin bilgileri
hash tablosunun 712’nci indeksteki elemanında saklanacaktır. 72484926820 TC kimlik
numarasına sahip kişinin bilgileri de dizinin 820’nci indeksteki elemanında
saklanacaktır. Ancak farklı kişilerin TC kimlik numaraları hash fonksiyonuna sokulduğunda
aynı indeks değerleri de elde edilebilecektir. Örneğin 6238517712 TC numarasına sahip
kişi de dizinin 712’nci indeksteki elemanına yerleşmek isteyecektir. İşte hash tablosu
yönteminde bu duruma \sphinxstyleemphasis{çakışma (collision)} denilmektedir. Hash tablosu yöntemi çakışma
durumunda izlenecek stratejiye göre çeşitli alt kollara ayrılmaktadır.


\subsection{Çakışma Çözümleme Stratejileri}
\label{\detokenize{kernellinkedlistandhashtable:cakisma-cozumleme-stratejileri}}
\sphinxAtStartPar
Hash tabloları yönteminde çakışma durumunda bu sorunu çözmek için temel olarak iki alt
yöntem grubu kullanılmaktadır: \sphinxstyleemphasis{ayrı zincir oluşturma (separate chaining)} yöntemi ve
\sphinxstyleemphasis{açık adresleme (open addressing)} yöntemi. Açık adresleme yöntemi de kendi aralarında
\sphinxstyleemphasis{doğrusal yoklama (linear probing)}, \sphinxstyleemphasis{karesel yoklama (quadratic probing)}, \sphinxstyleemphasis{çift hash’leme
(double hashing)} gibi alt yöntemlere ayrılmaktadır. Ayrı zincir oluşturma ve açık
adresleme yöntemlerinin dışında başka çakışma çözümleme stratejileri de vardır. Ancak
ağırlıklı olarak bu stratejiler tercih edilmektedir.

\begin{sphinxadmonition}{note}{Not:}
\sphinxAtStartPar
Linux çekirdeklerindeki tüm hash tablolarında \sphinxstyleemphasis{ayrı zincir oluşturma (separate
chaining)} alt yöntemi tercih edilmiştir.
\end{sphinxadmonition}


\subsubsection{Ayrı Zincir Oluşturma (Separate Chaining)}
\label{\detokenize{kernellinkedlistandhashtable:ayri-zincir-olusturma-separate-chaining}}
\sphinxAtStartPar
Ayrı zincir oluşturma (separate chaining) yönteminde hash tablosu aslında bir bağlı liste
dizisi gibi oluşturulur. Yani hash tablosunun her elemanı bağlı listenin ilk elemanını
(head pointer) gösteren bir gösterici durumundadır. Anahtar hash fonksiyonuna sokulur, bir
indeks elde edilir ve o indeksteki bağlı listenin hemen önüne (ya da duruma göre arkasına)
eklenir. Eleman aranırken anahtar yine aynı hash fonksiyonuna sokulur ve dizinin ilgili
indeksindeki bağlı listede sıralı arama yapılır.

\sphinxAtStartPar
Hash tablolarına eleman insert etmek O(1) karmaşıklıktadır. Tabii burada kullanılacak hash
fonksiyonu da önemlidir. Küçük döngüler içeren hash fonksiyonları O(1) karmaşıklığı
yükseltmemektedir. Elemanın silinmesi de O(1) karmaşıklıkta yapılabilmektedir. Eleman
aramanın O(1) karmaşıklıkta yapılabilmesi için bağlı listelerdeki zincir uzunluklarının
uzun olmaması gerekir. Örneğin yukarıda 10 kadar eleman için en hızlı arama yönteminin
sıralı arama olduğunu söylemiştik. Bu koşul sağlandığında sıralı aramanın O(1)
karmaşıklıkta olduğu söylenebilir. O halde eğer zincirlerdeki ortalama eleman makul bir
düzeyde tutulursa arama işlemi de O(1) karmaşıklıkta yapılabilecektir. Ancak hash tablosu
küçük fakat tabloya eklenecek eleman fazla ise bu durumda hash tablosu yöntemi artık sıralı
arama yöntemine benzer hale gelir. Yani bu yöntemin \sphinxstyleemphasis{en kötü durumdaki (worst case)}
karmaşıklığının O(N) olduğu söylenebilir.

\sphinxAtStartPar
Tabii hash tablosu yöntemini kullanan kişiler sistem hakkında bazı ön bilgilere sahip olursa
tabloyu uygun bir büyüklükte oluşturabilirler. İşletim sistemi gibi yüksek performans
isteyen sistemlerde hash tablolarının zincirlerinin ortalama 1 civarında tutulması uygun
olabilmektedir. Hash tabloları da duruma göre büyütülebilmektedir. Ancak büyütmenin önemli
bir zaman maliyeti vardır. Yeni bir hash tablosunun tahsis edilmesi; eski tablodaki
elemanların yeniden hash’lenerek yeni tabloya yerleştirilmesi uzun zaman alan bir işlemdir.
İşletim sistemlerinin çekirdeklerinde bu biçimde uzun zaman alacak işlemler tercih edilmez.
Bu nedenle Linux çekirdeğindeki hash tabloları büyütülmemektedir.

\sphinxAtStartPar
Hash tablolarında tablo elemanlarına (zincirlere değil) İngilizce \sphinxstyleemphasis{bucket (kova)},
tablodaki toplam eleman sayısının bucket sayısına bölümüne de \sphinxstyleemphasis{yükleme faktörü (load
factor)} denilmektedir. İdeal yükleme faktörünün \textless{}= 1 olduğunu söyleyebiliriz.

\sphinxAtStartPar
Sözlük tarzı veri yapılarında genel olarak aynı anahtara ilişkin birden fazla anahtar\sphinxhyphen{}değer
çifti veri yapısına yerleştirilememektedir. (Bazı kütüphanelerde buna izin verilebilmektedir.)
Eğer aynı anahtara ilişkin yeni bir değer insert edilmeye çalışılırsa eski değer yeni
değerle yer değiştirilmektedir. Bazı tasarımlar ise aynı anahtara ilişkin insert yapmayı
engellemektedir. Yani bu tasarımlarda yalnızca olmayan elemanlar tabloya insert
edilebilmektedir. Linux çekirdeğindeki hash tablolarında anahtarlar birbirinden farklı olmak
zorundadır.


\subsubsection{Hash Fonksiyonu Tasarımı}
\label{\detokenize{kernellinkedlistandhashtable:hash-fonksiyonu-tasarimi}}
\sphinxAtStartPar
Peki hash tablolarında kullanılacak iyi bir hash fonksiyonu nasıl olmalıdır? İyi bir hash
fonksiyonunun \sphinxstyleemphasis{hızlı} olması gerekir. Çünkü insert gibi arama gibi işlemlerde hash
fonksiyonu kullanılacaktır. İyi bir hash fonksiyonunun \sphinxstyleemphasis{anahtarlar yanlı bile olsa} tabloya
onları iyi bir biçimde yaydırabilmesi gerekir. Örneğin aslında sayısal anahtarlar için
“bölümden elde edilen kalan” iyi bir hash fonksiyonu değildir.

\sphinxAtStartPar
Hash tablolarında tablonun asal sayı uzunluğunda olması hash fonksiyonlarının daha iyi
yaydırmasına yardımcı olmaktadır. (Örneğin tablo uzunluğu için 100 yerine 101 değeri tercih
edilebilmektedir.) Hash fonksiyonları \sphinxstyleemphasis{sayıyı indekse dönüştüren} ve \sphinxstyleemphasis{yazıyı indekse
dönüştüren} fonksiyonlar biçiminde oluşturulabilir. Hash tablolarının 2’nin kuvveti
uzunluğunda alınması hash fonksiyonlarının daha hızlı çalışmasına da yol açabilmektedir.
(Örneğin bölümden elde edilen kalan yerine bit düzeyinde öteleme işlemleri hız kazancı
sağlayabilmektedir.) Linux çekirdeklerinde hash tabloları için genellikle sayfa katlarında
(yani 4096’nın katlarında) alanlar tahsis edilmektedir.


\subsubsection{Açık Adresleme ve Doğrusal Yoklama}
\label{\detokenize{kernellinkedlistandhashtable:acik-adresleme-ve-dogrusal-yoklama}}
\sphinxAtStartPar
Her ne kadar Linux çekirdeğindeki hash tablolarında \sphinxstyleemphasis{ayrı zincir oluşturma (separate
chaining)} stratejisi kullanılıyorsa da burada \sphinxstyleemphasis{açık adresleme (open addressing)} stratejisi
hakkında da küçük bir açıklama yapmak istiyoruz.

\sphinxAtStartPar
Açık adresleme yöntemi de kendi arasında \sphinxstyleemphasis{yoklama (probing)} biçimine göre çeşitli alt
yöntemlere ayrılmaktadır. Açık adreslemenin en yaygın ve basit biçimi \sphinxstyleemphasis{doğrusal yoklama
(linear probing)} denilen biçimidir.

\sphinxAtStartPar
Doğrusal yoklama oldukça basit bir fikre dayanmaktadır. Bu yöntemde yine bir hash tablosu
oluşturulur. Ancak hash tablosunda bağlı listelerin adresleri tutulmaz; bizzat değerlerin
kendisi tutulur. Tabloya eleman ekleneceği zaman yine anahtardan bir hash değeri elde edilir.
Doğrudan değer tablonun hash ile elde edilen indeksine yerleştirilir. Başka bir anahtar aynı
hash değerini verdiğinde (yani çakışma durumu oluştuğunda) o indeksten itibaren boş yer
bulunana kadar yan yana indekslere sırasıyla bakılır. Örneğin hash olarak 123 değerini elde
etmiş olalım. Tablonun 123’üncü elemanının dolu olduğunu düşünelim. Bu durumda 124’üncü
elemanına bakarız. O da doluysa 125’inci elemanına bakarız. Ta ki boş bir indeks bulunana
kadar. Değeri ilk boş indekse yerleştiririz.

\sphinxAtStartPar
Bu yöntemde arama işlemi de benzer biçimde yapılmaktadır. Yani aranacak elemanın hash değeri
elde edilir. O indekse başvurulur. Anahtar o indekste değilse anahtar bulunana kadar ya da
boş bir kova (bucket) görülene kadar yan yana diğer indekslere bakılır.

\sphinxAtStartPar
Anahtara dayalı eleman silme de benzer biçimde yapılmaktadır. Ancak eleman silindiğinde
ilgili kovanın (bucket) boşaltılması arama işlemlerinde sorunlara yol açabilecektir. Burada
yöntemlerden biri silinen elemana ilişkin kovanın boş yapılmayıp silinmenin özel bir değerle
belirtilmesidir. Örneğin her kova için bir durum bayrağı tutulabilir. Bu durum bayrağı
ilgili kovanın \sphinxstyleemphasis{dolu} olduğunu, \sphinxstyleemphasis{boş} olduğunu ya da \sphinxstyleemphasis{silinmiş} olduğunu belirtebilir.
Böylece arama sırasında \sphinxstyleemphasis{silinmiş} kovalar görüldüğünde durulmaz. İlk boş kova görüldüğünde
durulur. Tabii silinmiş kovalara yeni elemanlar eklenebilir.


\section{Linux Çekirdeğinde Hash Tablosu Gerçekleştirimi}
\label{\detokenize{kernellinkedlistandhashtable:linux-cekirdeginde-hash-tablosu-gerceklestirimi}}
\sphinxAtStartPar
Şimdi de Linux çekirdeğindeki hash tablosu gerçekleştirimi üzerinde duralım. Linux kaynak
kodlarında hash tablosuna ilişkin yapılar ve fonksiyonlar bağlı listelere ilişkin yapıların
ve fonksiyonların bulunduğu başlık dosyasında tanımlanmıştır. Yani hash tabloları için ayrı
bir başlık dosyası oluşturulmamıştır. RCU’suz hash tablolarının gerçekleştirimi
\sphinxcode{\sphinxupquote{include/linux/list.h}} dosyası içerisinde, RCU’lu hash tablolarının gerçekleştirimi ise
\sphinxcode{\sphinxupquote{include/linux/rculist.h}} dosyası içerisinde bulunmaktadır. Ancak her iki gerçekleştirim
de aynı yapıları kullanmaktadır.


\subsection{hlist\_head ve hlist\_node Yapıları}
\label{\detokenize{kernellinkedlistandhashtable:hlist-head-ve-hlist-node-yapilari}}
\sphinxAtStartPar
Ayrı zincir oluşturma yönteminde hash tablosundaki kovaların (buckets) aslında bağlı
listelerin ilk düğümünün yerini tutan göstericilerden oluştuğunu belirtmiştik. Linux
çekirdeklerinde bağlı listenin kök düğümü de eleman düğümleri de \sphinxcode{\sphinxupquote{list\_head}} yapısıyla
temsil ediliyordu. Ancak ayrı zincir oluşturmalı hash tablolarındaki kovalarda (buckets)
bağlı listenin son düğümünün yerinin tutulmasına gerek yoktur. Elemanlar hemen listenin
başına eklenebilir. Arama da listenin başından itibaren yapılabilir. Tabii performans
bakımından bağlı liste düğümlerinin yine çift bağlı (doubly linked) olması gerekir. Çünkü
adresini bildiğimiz bir düğümü hash tablosundan kolaylıkla çıkartabiliriz. O halde çekirdek
tablolarında \sphinxstyleemphasis{kök düğümde tek gösterici, eleman düğümlerinde ise çift gösterici}
bulunmalıdır.

\sphinxAtStartPar
Çekirdekte kullanılan hash tablosu gerçekleştiriminde kök düğüm \sphinxcode{\sphinxupquote{hlist\_head}} yapısıyla
temsil edilmiştir:

\begin{sphinxVerbatim}[commandchars=\\\{\}]
\PYG{k}{struct}\PYG{+w}{ }\PYG{n+nc}{hlist\PYGZus{}head}\PYG{+w}{ }\PYG{p}{\PYGZob{}}
\PYG{+w}{    }\PYG{k}{struct}\PYG{+w}{ }\PYG{n+nc}{hlist\PYGZus{}node}\PYG{+w}{ }\PYG{o}{*}\PYG{n}{first}\PYG{p}{;}
\PYG{p}{\PYGZcb{}}\PYG{p}{;}
\end{sphinxVerbatim}

\sphinxAtStartPar
Görüldüğü gibi yapıda tek bir gösterici vardır. O da zincirdeki ilk elemanı göstermektedir.
Zincirdeki bağlı listenin düğümleri de \sphinxcode{\sphinxupquote{hlist\_node}} yapısıyla temsil edilmiştir:

\begin{sphinxVerbatim}[commandchars=\\\{\}]
\PYG{k}{struct}\PYG{+w}{ }\PYG{n+nc}{hlist\PYGZus{}node}\PYG{+w}{ }\PYG{p}{\PYGZob{}}
\PYG{+w}{    }\PYG{k}{struct}\PYG{+w}{ }\PYG{n+nc}{hlist\PYGZus{}node}\PYG{+w}{ }\PYG{o}{*}\PYG{n}{next}\PYG{p}{,}\PYG{+w}{ }\PYG{o}{*}\PYG{o}{*}\PYG{n}{pprev}\PYG{p}{;}
\PYG{p}{\PYGZcb{}}\PYG{p}{;}
\end{sphinxVerbatim}

\sphinxAtStartPar
Buradaki düğüm yapısını listelerdeki düğüm yapısı ile karşılaştırınız:

\begin{sphinxVerbatim}[commandchars=\\\{\}]
\PYG{k}{struct}\PYG{+w}{ }\PYG{n+nc}{list\PYGZus{}head}\PYG{+w}{ }\PYG{p}{\PYGZob{}}
\PYG{+w}{    }\PYG{k}{struct}\PYG{+w}{ }\PYG{n+nc}{list\PYGZus{}head}\PYG{+w}{ }\PYG{o}{*}\PYG{n}{next}\PYG{p}{,}\PYG{+w}{ }\PYG{o}{*}\PYG{n}{prev}\PYG{p}{;}
\PYG{p}{\PYGZcb{}}\PYG{p}{;}
\end{sphinxVerbatim}

\sphinxAtStartPar
Hash tablosundaki düğümlerin \sphinxcode{\sphinxupquote{pprev}} elemanı önceki düğümün adresini değil önceki
düğümdeki \sphinxcode{\sphinxupquote{next}} elemanının adresini göstermektedir. Bu nedenle \sphinxcode{\sphinxupquote{pprev}} göstericiyi
gösteren göstericidir. Hash tablolarındaki düğümlerde geri gitmek için bir neden yoktur.
Ancak eleman silme durumunda bu tasarım önceki elemanın \sphinxcode{\sphinxupquote{next}} göstericisinin daha kolay
güncellenmesine yol açmaktadır. Örneğin \sphinxcode{\sphinxupquote{p}} göstericisi bir \sphinxcode{\sphinxupquote{hlist\_node}} düğümünü
gösteriyor olsun. Biz de bu düğümü silecek olalım. Burada önceki düğümün \sphinxcode{\sphinxupquote{next}} elemanının
sileceğimiz düğümün \sphinxcode{\sphinxupquote{next}} elemanındaki düğümü göstermesi sağlanmalıdır. Bu işlem pratik
olarak şöyle yapılabilmektedir:

\begin{sphinxVerbatim}[commandchars=\\\{\}]
\PYG{o}{*}\PYG{n}{p}\PYG{o}{\PYGZhy{}}\PYG{o}{\PYGZgt{}}\PYG{n}{pprev}\PYG{+w}{ }\PYG{o}{=}\PYG{+w}{ }\PYG{n}{p}\PYG{o}{\PYGZhy{}}\PYG{o}{\PYGZgt{}}\PYG{n}{next}\PYG{p}{;}
\end{sphinxVerbatim}

\sphinxAtStartPar
Halbuki düğümler \sphinxcode{\sphinxupquote{list\_head}} yapısıyla temsil ediliyor olsaydı bu işlem ancak şöyle
yapılabilirdi:

\begin{sphinxVerbatim}[commandchars=\\\{\}]
\PYG{n}{p}\PYG{o}{\PYGZhy{}}\PYG{o}{\PYGZgt{}}\PYG{n}{prev}\PYG{o}{\PYGZhy{}}\PYG{o}{\PYGZgt{}}\PYG{n}{next}\PYG{+w}{ }\PYG{o}{=}\PYG{+w}{ }\PYG{n}{p}\PYG{o}{\PYGZhy{}}\PYG{o}{\PYGZgt{}}\PYG{n}{next}\PYG{p}{;}
\end{sphinxVerbatim}

\sphinxAtStartPar
Görüldüğü gibi bu güncellemede fazladan bir işlem yapılmaktadır. \sphinxcode{\sphinxupquote{hlist\_node}} tasarımının
diğer önemli bir faydası da bu tasarımda kök düğüm için özel bir işlemin yapılmasına gerek
kalmamasıdır. Yani listeye ilk kez eleman eklerken \sphinxcode{\sphinxupquote{*p\sphinxhyphen{}\textgreater{}pprev}} kök düğümün \sphinxcode{\sphinxupquote{next}}
göstericisi haline gelecektir. Heterojen yapılarda bu tür güncellemelerin yapılması daha
fazla çabayı gerektirmektedir.

\sphinxAtStartPar
Peki neden bütün çift bağlı listelerde bu teknik kullanılmıyor? \sphinxcode{\sphinxupquote{hlist\_node}} yapısında
olduğu gibi \sphinxcode{\sphinxupquote{prev}} göstericisi önceki düğümün başlangıç adresini göstermek yerine neden
önceki düğümün \sphinxcode{\sphinxupquote{next}} göstericisini göstermiyor? İşte geriye doğru ilerlemenin gerekli
olabildiği durumlarda bu tasarımda geriye gidiş zorlaşmaktadır. Ancak yukarıda da
belirttiğimiz gibi hash tablolarındaki zincirlerde zaten geriye gitmenin bir anlamı yoktur.


\subsection{Hash Tablosunun Oluşturulması}
\label{\detokenize{kernellinkedlistandhashtable:hash-tablosunun-olusturulmasi}}
\sphinxAtStartPar
Peki biz yukarıdaki \sphinxcode{\sphinxupquote{hlist\_head}} ve \sphinxcode{\sphinxupquote{hlist\_node}} yapılarını kullanarak hash tablosunu
nasıl oluşturabiliriz? İşte yapılacak ilk şey \sphinxcode{\sphinxupquote{hlist\_head}} yapısı türünden bir dizi tahsis
etmektir. Örneğin biz bu işlemi kullanıcı modunda aşağıdaki gibi simüle edebiliriz:

\begin{sphinxVerbatim}[commandchars=\\\{\}]
\PYG{c+cp}{\PYGZsh{}}\PYG{c+cp}{include}\PYG{+w}{ }\PYG{c+cpf}{\PYGZlt{}stdio.h\PYGZgt{}}
\PYG{c+cp}{\PYGZsh{}}\PYG{c+cp}{include}\PYG{+w}{ }\PYG{c+cpf}{\PYGZlt{}stdlib.h\PYGZgt{}}

\PYG{c+cp}{\PYGZsh{}}\PYG{c+cp}{define TABLE\PYGZus{}SIZE    4096}

\PYG{k}{struct}\PYG{+w}{ }\PYG{n+nc}{hlist\PYGZus{}head}\PYG{+w}{ }\PYG{p}{\PYGZob{}}
\PYG{+w}{    }\PYG{k}{struct}\PYG{+w}{ }\PYG{n+nc}{hlist\PYGZus{}node}\PYG{+w}{ }\PYG{o}{*}\PYG{n}{first}\PYG{p}{;}
\PYG{p}{\PYGZcb{}}\PYG{p}{;}

\PYG{k}{struct}\PYG{+w}{ }\PYG{n+nc}{hlist\PYGZus{}node}\PYG{+w}{ }\PYG{p}{\PYGZob{}}
\PYG{+w}{    }\PYG{k}{struct}\PYG{+w}{ }\PYG{n+nc}{hlist\PYGZus{}node}\PYG{+w}{ }\PYG{o}{*}\PYG{n}{next}\PYG{p}{,}\PYG{+w}{ }\PYG{o}{*}\PYG{o}{*}\PYG{n}{pprev}\PYG{p}{;}
\PYG{p}{\PYGZcb{}}\PYG{p}{;}

\PYG{k+kt}{int}\PYG{+w}{ }\PYG{n+nf}{main}\PYG{p}{(}\PYG{k+kt}{void}\PYG{p}{)}
\PYG{p}{\PYGZob{}}
\PYG{+w}{    }\PYG{k}{struct}\PYG{+w}{ }\PYG{n+nc}{hlist\PYGZus{}head}\PYG{+w}{ }\PYG{o}{*}\PYG{n}{hash\PYGZus{}table}\PYG{p}{;}

\PYG{+w}{    }\PYG{k}{if}\PYG{+w}{ }\PYG{p}{(}\PYG{p}{(}\PYG{n}{hash\PYGZus{}table}\PYG{+w}{ }\PYG{o}{=}\PYG{+w}{ }\PYG{p}{(}\PYG{k}{struct}\PYG{+w}{ }\PYG{n+nc}{hlist\PYGZus{}head}\PYG{+w}{ }\PYG{o}{*}\PYG{p}{)}\PYG{n}{malloc}\PYG{p}{(}
\PYG{+w}{            }\PYG{k}{sizeof}\PYG{p}{(}\PYG{k}{struct}\PYG{+w}{ }\PYG{n+nc}{hlist\PYGZus{}head}\PYG{p}{)}\PYG{+w}{ }\PYG{o}{*}\PYG{+w}{ }\PYG{n}{TABLE\PYGZus{}SIZE}\PYG{p}{)}\PYG{p}{)}\PYG{+w}{ }\PYG{o}{=}\PYG{o}{=}\PYG{+w}{ }\PYG{n+nb}{NULL}\PYG{p}{)}\PYG{+w}{ }\PYG{p}{\PYGZob{}}
\PYG{+w}{        }\PYG{n}{fprintf}\PYG{p}{(}\PYG{n}{stderr}\PYG{p}{,}\PYG{+w}{ }\PYG{l+s}{\PYGZdq{}}\PYG{l+s}{cannot allocate memory!...}\PYG{l+s+se}{\PYGZbs{}n}\PYG{l+s}{\PYGZdq{}}\PYG{p}{)}\PYG{p}{;}
\PYG{+w}{        }\PYG{n}{exit}\PYG{p}{(}\PYG{n}{EXIT\PYGZus{}FAILURE}\PYG{p}{)}\PYG{p}{;}
\PYG{+w}{    }\PYG{p}{\PYGZcb{}}

\PYG{+w}{    }\PYG{c+cm}{/* ... */}

\PYG{+w}{    }\PYG{n}{free}\PYG{p}{(}\PYG{n}{hash\PYGZus{}table}\PYG{p}{)}\PYG{p}{;}

\PYG{+w}{    }\PYG{k}{return}\PYG{+w}{ }\PYG{l+m+mi}{0}\PYG{p}{;}
\PYG{p}{\PYGZcb{}}
\end{sphinxVerbatim}

\sphinxAtStartPar
Çekirdekteki \sphinxcode{\sphinxupquote{list.h}} dosyası içerisinde hash tablosuna eleman eklemek için çeşitli basit
fonksiyonlar bulundurulmuştur. Tabii \sphinxcode{\sphinxupquote{hlist\_node}} düğümleri de aslında başka yapıların
elemanı durumunda olacaktır. Yine asıl yapı nesnesinin adresi \sphinxcode{\sphinxupquote{container\_of}} makrosuyla
elde edilecektir.

\sphinxAtStartPar
Başlangıçta \sphinxcode{\sphinxupquote{hlist\_head}} yapısındaki \sphinxcode{\sphinxupquote{first}} elemanı \sphinxcode{\sphinxupquote{NULL}} değerinde olmalıdır.
Bunun için \sphinxcode{\sphinxupquote{list.h}} içerisinde makrolar bulundurulmuştur:

\begin{sphinxVerbatim}[commandchars=\\\{\}]
\PYG{c+cp}{\PYGZsh{}}\PYG{c+cp}{define HLIST\PYGZus{}HEAD\PYGZus{}INIT \PYGZob{} .first = NULL \PYGZcb{}}
\PYG{c+cp}{\PYGZsh{}}\PYG{c+cp}{define HLIST\PYGZus{}HEAD(name) struct hlist\PYGZus{}head name = \PYGZob{} .first = NULL \PYGZcb{}}
\PYG{c+cp}{\PYGZsh{}}\PYG{c+cp}{define INIT\PYGZus{}HLIST\PYGZus{}HEAD(ptr) ((ptr)\PYGZhy{}\PYGZgt{}first = NULL)}
\end{sphinxVerbatim}

\sphinxAtStartPar
Örneğin:

\begin{sphinxVerbatim}[commandchars=\\\{\}]
\PYG{k}{for}\PYG{+w}{ }\PYG{p}{(}\PYG{k+kt}{int}\PYG{+w}{ }\PYG{n}{i}\PYG{+w}{ }\PYG{o}{=}\PYG{+w}{ }\PYG{l+m+mi}{0}\PYG{p}{;}\PYG{+w}{ }\PYG{n}{i}\PYG{+w}{ }\PYG{o}{\PYGZlt{}}\PYG{+w}{ }\PYG{n}{TABLE\PYGZus{}SIZE}\PYG{p}{;}\PYG{+w}{ }\PYG{o}{+}\PYG{o}{+}\PYG{n}{i}\PYG{p}{)}
\PYG{+w}{    }\PYG{n}{INIT\PYGZus{}HLIST\PYGZus{}HEAD}\PYG{p}{(}\PYG{o}{\PYGZam{}}\PYG{n}{hash\PYGZus{}table}\PYG{p}{[}\PYG{n}{i}\PYG{p}{]}\PYG{p}{)}\PYG{p}{;}
\end{sphinxVerbatim}

\sphinxAtStartPar
Zincirlerdeki son düğümün \sphinxcode{\sphinxupquote{next}} elemanı da \sphinxcode{\sphinxupquote{NULL}} değerindedir. Böylece arama
\sphinxcode{\sphinxupquote{NULL}} görmeyene kadar ilerlenerek yapılabilmektedir.

\sphinxAtStartPar
Tablodaki zincirlerin önüne eleman eklemek için \sphinxcode{\sphinxupquote{hlist\_add\_head}} fonksiyonu
bulundurulmuştur:

\begin{sphinxVerbatim}[commandchars=\\\{\}]
\PYG{k}{static}\PYG{+w}{ }\PYG{k+kr}{inline}\PYG{+w}{ }\PYG{k+kt}{void}\PYG{+w}{ }\PYG{n+nf}{hlist\PYGZus{}add\PYGZus{}head}\PYG{p}{(}\PYG{k}{struct}\PYG{+w}{ }\PYG{n+nc}{hlist\PYGZus{}node}\PYG{+w}{ }\PYG{o}{*}\PYG{n}{n}\PYG{p}{,}\PYG{+w}{ }\PYG{k}{struct}\PYG{+w}{ }\PYG{n+nc}{hlist\PYGZus{}head}\PYG{+w}{ }\PYG{o}{*}\PYG{n}{h}\PYG{p}{)}
\PYG{p}{\PYGZob{}}
\PYG{+w}{    }\PYG{k}{struct}\PYG{+w}{ }\PYG{n+nc}{hlist\PYGZus{}node}\PYG{+w}{ }\PYG{o}{*}\PYG{n}{first}\PYG{+w}{ }\PYG{o}{=}\PYG{+w}{ }\PYG{n}{h}\PYG{o}{\PYGZhy{}}\PYG{o}{\PYGZgt{}}\PYG{n}{first}\PYG{p}{;}
\PYG{+w}{    }\PYG{n}{WRITE\PYGZus{}ONCE}\PYG{p}{(}\PYG{n}{n}\PYG{o}{\PYGZhy{}}\PYG{o}{\PYGZgt{}}\PYG{n}{next}\PYG{p}{,}\PYG{+w}{ }\PYG{n}{first}\PYG{p}{)}\PYG{p}{;}
\PYG{+w}{    }\PYG{k}{if}\PYG{+w}{ }\PYG{p}{(}\PYG{n}{first}\PYG{p}{)}
\PYG{+w}{        }\PYG{n}{WRITE\PYGZus{}ONCE}\PYG{p}{(}\PYG{n}{first}\PYG{o}{\PYGZhy{}}\PYG{o}{\PYGZgt{}}\PYG{n}{pprev}\PYG{p}{,}\PYG{+w}{ }\PYG{o}{\PYGZam{}}\PYG{n}{n}\PYG{o}{\PYGZhy{}}\PYG{o}{\PYGZgt{}}\PYG{n}{next}\PYG{p}{)}\PYG{p}{;}
\PYG{+w}{    }\PYG{n}{WRITE\PYGZus{}ONCE}\PYG{p}{(}\PYG{n}{h}\PYG{o}{\PYGZhy{}}\PYG{o}{\PYGZgt{}}\PYG{n}{first}\PYG{p}{,}\PYG{+w}{ }\PYG{n}{n}\PYG{p}{)}\PYG{p}{;}
\PYG{+w}{    }\PYG{n}{WRITE\PYGZus{}ONCE}\PYG{p}{(}\PYG{n}{n}\PYG{o}{\PYGZhy{}}\PYG{o}{\PYGZgt{}}\PYG{n}{pprev}\PYG{p}{,}\PYG{+w}{ }\PYG{o}{\PYGZam{}}\PYG{n}{h}\PYG{o}{\PYGZhy{}}\PYG{o}{\PYGZgt{}}\PYG{n}{first}\PYG{p}{)}\PYG{p}{;}
\PYG{p}{\PYGZcb{}}
\end{sphinxVerbatim}

\sphinxAtStartPar
Fonksiyonun birinci parametresi eklenecek yeni düğümün adresini, ikinci parametresi ise kök
düğüm nesnesini belirtmektedir. Buradaki \sphinxcode{\sphinxupquote{WRITE\_ONCE}} atamayı volatile erişimle
yapmaktadır. \sphinxcode{\sphinxupquote{WRITE\_ONCE(a, b)}} çağrısını \sphinxcode{\sphinxupquote{a = b}} gibi düşünebilirsiniz.


\subsection{Düğüm Silme}
\label{\detokenize{kernellinkedlistandhashtable:id1}}
\sphinxAtStartPar
Hash tablosundan bir düğüm silmek için \sphinxcode{\sphinxupquote{hlist\_del}} fonksiyonu kullanılmaktadır. Fonksiyon
\sphinxcode{\sphinxupquote{list.h}} içerisinde şöyle tanımlanmıştır:

\begin{sphinxVerbatim}[commandchars=\\\{\}]
\PYG{k}{static}\PYG{+w}{ }\PYG{k+kr}{inline}\PYG{+w}{ }\PYG{k+kt}{void}\PYG{+w}{ }\PYG{n+nf}{hlist\PYGZus{}del}\PYG{p}{(}\PYG{k}{struct}\PYG{+w}{ }\PYG{n+nc}{hlist\PYGZus{}node}\PYG{+w}{ }\PYG{o}{*}\PYG{n}{n}\PYG{p}{)}
\PYG{p}{\PYGZob{}}
\PYG{+w}{    }\PYG{n}{\PYGZus{}\PYGZus{}hlist\PYGZus{}del}\PYG{p}{(}\PYG{n}{n}\PYG{p}{)}\PYG{p}{;}
\PYG{+w}{    }\PYG{n}{n}\PYG{o}{\PYGZhy{}}\PYG{o}{\PYGZgt{}}\PYG{n}{next}\PYG{+w}{ }\PYG{o}{=}\PYG{+w}{ }\PYG{n}{LIST\PYGZus{}POISON1}\PYG{p}{;}
\PYG{+w}{    }\PYG{n}{n}\PYG{o}{\PYGZhy{}}\PYG{o}{\PYGZgt{}}\PYG{n}{pprev}\PYG{+w}{ }\PYG{o}{=}\PYG{+w}{ }\PYG{n}{LIST\PYGZus{}POISON2}\PYG{p}{;}
\PYG{p}{\PYGZcb{}}
\end{sphinxVerbatim}

\sphinxAtStartPar
Burada asıl silme işlemini yapan fonksiyon \sphinxcode{\sphinxupquote{\_\_hlist\_del}} fonksiyonudur. Düğüm silindikten
sonra silinen düğümün \sphinxcode{\sphinxupquote{next}} ve \sphinxcode{\sphinxupquote{pprev}} göstericilerine güvenlik amacıyla özel değerler
atandığını görüyorsunuz. Bu özel değerler geçerli adresler belirtmemektedir. Eğer silinen
düğüm yanlışlıkla kullanılırsa \sphinxstyleemphasis{page fault} oluşmasına yol açacaktır. \sphinxcode{\sphinxupquote{\_\_hlist\_del}}
fonksiyonu da şöyle tanımlanmıştır:

\begin{sphinxVerbatim}[commandchars=\\\{\}]
\PYG{k}{static}\PYG{+w}{ }\PYG{k+kr}{inline}\PYG{+w}{ }\PYG{k+kt}{void}\PYG{+w}{ }\PYG{n+nf}{\PYGZus{}\PYGZus{}hlist\PYGZus{}del}\PYG{p}{(}\PYG{k}{struct}\PYG{+w}{ }\PYG{n+nc}{hlist\PYGZus{}node}\PYG{+w}{ }\PYG{o}{*}\PYG{n}{n}\PYG{p}{)}
\PYG{p}{\PYGZob{}}
\PYG{+w}{    }\PYG{k}{struct}\PYG{+w}{ }\PYG{n+nc}{hlist\PYGZus{}node}\PYG{+w}{ }\PYG{o}{*}\PYG{n}{next}\PYG{+w}{ }\PYG{o}{=}\PYG{+w}{ }\PYG{n}{n}\PYG{o}{\PYGZhy{}}\PYG{o}{\PYGZgt{}}\PYG{n}{next}\PYG{p}{;}
\PYG{+w}{    }\PYG{k}{struct}\PYG{+w}{ }\PYG{n+nc}{hlist\PYGZus{}node}\PYG{+w}{ }\PYG{o}{*}\PYG{o}{*}\PYG{n}{pprev}\PYG{+w}{ }\PYG{o}{=}\PYG{+w}{ }\PYG{n}{n}\PYG{o}{\PYGZhy{}}\PYG{o}{\PYGZgt{}}\PYG{n}{pprev}\PYG{p}{;}

\PYG{+w}{    }\PYG{n}{WRITE\PYGZus{}ONCE}\PYG{p}{(}\PYG{o}{*}\PYG{n}{pprev}\PYG{p}{,}\PYG{+w}{ }\PYG{n}{next}\PYG{p}{)}\PYG{p}{;}
\PYG{+w}{    }\PYG{k}{if}\PYG{+w}{ }\PYG{p}{(}\PYG{n}{next}\PYG{p}{)}
\PYG{+w}{        }\PYG{n}{WRITE\PYGZus{}ONCE}\PYG{p}{(}\PYG{n}{next}\PYG{o}{\PYGZhy{}}\PYG{o}{\PYGZgt{}}\PYG{n}{pprev}\PYG{p}{,}\PYG{+w}{ }\PYG{n}{pprev}\PYG{p}{)}\PYG{p}{;}
\PYG{p}{\PYGZcb{}}
\end{sphinxVerbatim}

\sphinxAtStartPar
Zincirdeki son düğümün \sphinxcode{\sphinxupquote{next}} göstericisinde \sphinxcode{\sphinxupquote{NULL}} adres bulunmalıdır. Fonksiyonun
içerisinde bu durumun kontrol edildiğine ve ilk düğümün silinmesinin özel bir durum
oluşturmadığına dikkat ediniz.


\subsection{Ek Insert Fonksiyonları}
\label{\detokenize{kernellinkedlistandhashtable:ek-insert-fonksiyonlari}}
\sphinxAtStartPar
Çok fazla gereksinim olmasa da zincirdeki belli bir düğümün önüne ve arkasına düğüm insert
eden \sphinxcode{\sphinxupquote{hlist\_add\_before}} ve \sphinxcode{\sphinxupquote{hlist\_add\_behind}} fonksiyonları da bulundurulmuştur:

\begin{sphinxVerbatim}[commandchars=\\\{\}]
\PYG{k}{static}\PYG{+w}{ }\PYG{k+kr}{inline}\PYG{+w}{ }\PYG{k+kt}{void}\PYG{+w}{ }\PYG{n+nf}{hlist\PYGZus{}add\PYGZus{}before}\PYG{p}{(}\PYG{k}{struct}\PYG{+w}{ }\PYG{n+nc}{hlist\PYGZus{}node}\PYG{+w}{ }\PYG{o}{*}\PYG{n}{n}\PYG{p}{,}
\PYG{+w}{                                     }\PYG{k}{struct}\PYG{+w}{ }\PYG{n+nc}{hlist\PYGZus{}node}\PYG{+w}{ }\PYG{o}{*}\PYG{n}{next}\PYG{p}{)}
\PYG{p}{\PYGZob{}}
\PYG{+w}{    }\PYG{n}{WRITE\PYGZus{}ONCE}\PYG{p}{(}\PYG{n}{n}\PYG{o}{\PYGZhy{}}\PYG{o}{\PYGZgt{}}\PYG{n}{pprev}\PYG{p}{,}\PYG{+w}{ }\PYG{n}{next}\PYG{o}{\PYGZhy{}}\PYG{o}{\PYGZgt{}}\PYG{n}{pprev}\PYG{p}{)}\PYG{p}{;}
\PYG{+w}{    }\PYG{n}{WRITE\PYGZus{}ONCE}\PYG{p}{(}\PYG{n}{n}\PYG{o}{\PYGZhy{}}\PYG{o}{\PYGZgt{}}\PYG{n}{next}\PYG{p}{,}\PYG{+w}{ }\PYG{n}{next}\PYG{p}{)}\PYG{p}{;}
\PYG{+w}{    }\PYG{n}{WRITE\PYGZus{}ONCE}\PYG{p}{(}\PYG{n}{next}\PYG{o}{\PYGZhy{}}\PYG{o}{\PYGZgt{}}\PYG{n}{pprev}\PYG{p}{,}\PYG{+w}{ }\PYG{o}{\PYGZam{}}\PYG{n}{n}\PYG{o}{\PYGZhy{}}\PYG{o}{\PYGZgt{}}\PYG{n}{next}\PYG{p}{)}\PYG{p}{;}
\PYG{+w}{    }\PYG{n}{WRITE\PYGZus{}ONCE}\PYG{p}{(}\PYG{o}{*}\PYG{p}{(}\PYG{n}{n}\PYG{o}{\PYGZhy{}}\PYG{o}{\PYGZgt{}}\PYG{n}{pprev}\PYG{p}{)}\PYG{p}{,}\PYG{+w}{ }\PYG{n}{n}\PYG{p}{)}\PYG{p}{;}
\PYG{p}{\PYGZcb{}}

\PYG{k}{static}\PYG{+w}{ }\PYG{k+kr}{inline}\PYG{+w}{ }\PYG{k+kt}{void}\PYG{+w}{ }\PYG{n+nf}{hlist\PYGZus{}add\PYGZus{}behind}\PYG{p}{(}\PYG{k}{struct}\PYG{+w}{ }\PYG{n+nc}{hlist\PYGZus{}node}\PYG{+w}{ }\PYG{o}{*}\PYG{n}{n}\PYG{p}{,}
\PYG{+w}{                                     }\PYG{k}{struct}\PYG{+w}{ }\PYG{n+nc}{hlist\PYGZus{}node}\PYG{+w}{ }\PYG{o}{*}\PYG{n}{prev}\PYG{p}{)}
\PYG{p}{\PYGZob{}}
\PYG{+w}{    }\PYG{n}{WRITE\PYGZus{}ONCE}\PYG{p}{(}\PYG{n}{n}\PYG{o}{\PYGZhy{}}\PYG{o}{\PYGZgt{}}\PYG{n}{next}\PYG{p}{,}\PYG{+w}{ }\PYG{n}{prev}\PYG{o}{\PYGZhy{}}\PYG{o}{\PYGZgt{}}\PYG{n}{next}\PYG{p}{)}\PYG{p}{;}
\PYG{+w}{    }\PYG{n}{WRITE\PYGZus{}ONCE}\PYG{p}{(}\PYG{n}{prev}\PYG{o}{\PYGZhy{}}\PYG{o}{\PYGZgt{}}\PYG{n}{next}\PYG{p}{,}\PYG{+w}{ }\PYG{n}{n}\PYG{p}{)}\PYG{p}{;}
\PYG{+w}{    }\PYG{n}{WRITE\PYGZus{}ONCE}\PYG{p}{(}\PYG{n}{n}\PYG{o}{\PYGZhy{}}\PYG{o}{\PYGZgt{}}\PYG{n}{pprev}\PYG{p}{,}\PYG{+w}{ }\PYG{o}{\PYGZam{}}\PYG{n}{prev}\PYG{o}{\PYGZhy{}}\PYG{o}{\PYGZgt{}}\PYG{n}{next}\PYG{p}{)}\PYG{p}{;}

\PYG{+w}{    }\PYG{k}{if}\PYG{+w}{ }\PYG{p}{(}\PYG{n}{n}\PYG{o}{\PYGZhy{}}\PYG{o}{\PYGZgt{}}\PYG{n}{next}\PYG{p}{)}
\PYG{+w}{        }\PYG{n}{WRITE\PYGZus{}ONCE}\PYG{p}{(}\PYG{n}{n}\PYG{o}{\PYGZhy{}}\PYG{o}{\PYGZgt{}}\PYG{n}{next}\PYG{o}{\PYGZhy{}}\PYG{o}{\PYGZgt{}}\PYG{n}{pprev}\PYG{p}{,}\PYG{+w}{ }\PYG{o}{\PYGZam{}}\PYG{n}{n}\PYG{o}{\PYGZhy{}}\PYG{o}{\PYGZgt{}}\PYG{n}{next}\PYG{p}{)}\PYG{p}{;}
\PYG{p}{\PYGZcb{}}
\end{sphinxVerbatim}


\subsection{Döngü Makroları}
\label{\detokenize{kernellinkedlistandhashtable:dongu-makrolari}}
\sphinxAtStartPar
Hash tablosundaki bir \sphinxcode{\sphinxupquote{hlist\_node}} adresi bilindiğinde onun içinde bulunduğu asıl yapının
adresinin \sphinxcode{\sphinxupquote{container\_of}} makrosu ile elde edilebildiğini biliyorsunuz. Ancak tıpkı bağlı
listelerde olduğu gibi hash tablolarında da bunun için \sphinxcode{\sphinxupquote{container\_of}} makrosu ile aynı
işlemi yapan bir \sphinxcode{\sphinxupquote{entry}} makrosu bulundurulmuştur:

\begin{sphinxVerbatim}[commandchars=\\\{\}]
\PYG{c+cp}{\PYGZsh{}}\PYG{c+cp}{define hlist\PYGZus{}entry(ptr, type, member) container\PYGZus{}of(ptr, type, member)}
\end{sphinxVerbatim}

\sphinxAtStartPar
Hash tablosundaki bir zinciri dolaşmak için \sphinxcode{\sphinxupquote{hlist\_for\_each}} döngü makrosu
kullanılmaktadır:

\begin{sphinxVerbatim}[commandchars=\\\{\}]
\PYG{c+cp}{\PYGZsh{}}\PYG{c+cp}{define hlist\PYGZus{}for\PYGZus{}each(pos, head) \PYGZbs{}}
\PYG{c+cp}{    for (pos = (head)\PYGZhy{}\PYGZgt{}first; pos ; pos = pos\PYGZhy{}\PYGZgt{}next)}
\end{sphinxVerbatim}

\sphinxAtStartPar
Makronun ilk parametresi \sphinxcode{\sphinxupquote{hlist\_node}} türünden bir göstericiyi, ikinci parametresi ise
bağlı liste zincirinin başlangıç düğümüne ilişkin \sphinxcode{\sphinxupquote{hlist\_head}} nesnesinin adresini
almaktadır. Bu döngü makrosunun \sphinxcode{\sphinxupquote{next}} elemanı \sphinxcode{\sphinxupquote{NULL}} olmayana kadar ilerleme
sağladığına dikkat ediniz. Döngü her yinelendikçe birinci parametreye girilen göstericinin
içerisine zincirdeki sonraki düğümün adresi yerleştirilmektedir. Bu döngü makrosuyla zincir
dolaşılırken döngünün her yinelenmesinde asıl yapı nesnesinin değil onun içerisindeki
\sphinxcode{\sphinxupquote{hlist\_node}} nesnesinin adresinin elde edildiğine da dikkat ediniz. Bu adresten hareketle
\sphinxcode{\sphinxupquote{container\_of}} ya da \sphinxcode{\sphinxupquote{hlist\_entry}} makrolarıyla asıl nesnenin başlangıç adresinin elde
edilmesi gerekir. İşte bu iki işlemi aynı anda yapan \sphinxcode{\sphinxupquote{hlist\_for\_each\_entry}} isimli bir
döngü makrosu da bulundurulmuştur:

\begin{sphinxVerbatim}[commandchars=\\\{\}]
\PYG{c+cp}{\PYGZsh{}}\PYG{c+cp}{define hlist\PYGZus{}for\PYGZus{}each\PYGZus{}entry(pos, head, member)                           \PYGZbs{}}
\PYG{c+cp}{for (pos = hlist\PYGZus{}entry((head)\PYGZhy{}\PYGZgt{}first, typeof(*(pos)), member);            \PYGZbs{}}
\PYG{c+cp}{    pos;                                                                   \PYGZbs{}}
\PYG{c+cp}{    pos = hlist\PYGZus{}entry((pos)\PYGZhy{}\PYGZgt{}member.next, typeof(*(pos)), member))}
\end{sphinxVerbatim}

\sphinxAtStartPar
Bu makronun artık birinci parametresi doğrudan asıl yapı türünden bir göstericiyi, ikinci
ve üçüncü parametreler sırasıyla bağlı listenin \sphinxcode{\sphinxupquote{hlist\_head}} adresini ve asıl yapıdaki
bağ elemanının ismini almaktadır.

\sphinxAtStartPar
Yukarıdaki makroda bir noktaya dikkatinizi çekmek istiyoruz. Hash tablosundaki
\sphinxcode{\sphinxupquote{hlist\_head}} kök düğümünün \sphinxcode{\sphinxupquote{first}} elemanında \sphinxcode{\sphinxupquote{NULL}} adres varsa yukarıdaki
\sphinxcode{\sphinxupquote{hlist\_for\_each\_entry}} makrosu tanımsız davranışa yol açar. İşte eğer \sphinxcode{\sphinxupquote{first}} elemanı
\sphinxcode{\sphinxupquote{NULL}} ise bu durumu ele alıp dolaşımı sonlandıran \sphinxcode{\sphinxupquote{hlist\_for\_each\_entry\_safe}} döngü
makrosu da bulundurulmuştur:

\begin{sphinxVerbatim}[commandchars=\\\{\}]
\PYG{c+cp}{\PYGZsh{}}\PYG{c+cp}{define hlist\PYGZus{}for\PYGZus{}each\PYGZus{}entry\PYGZus{}safe(pos, n, head, member)                   \PYGZbs{}}
\PYG{c+cp}{for (pos = hlist\PYGZus{}entry\PYGZus{}safe((head)\PYGZhy{}\PYGZgt{}first, typeof(*pos), member);         \PYGZbs{}}
\PYG{c+cp}{    pos \PYGZam{}\PYGZam{} (\PYGZob{} n = pos\PYGZhy{}\PYGZgt{}member.next; 1; \PYGZcb{});                                 \PYGZbs{}}
\PYG{c+cp}{    pos = hlist\PYGZus{}entry\PYGZus{}safe(n, typeof(*pos), member))}
\end{sphinxVerbatim}

\sphinxAtStartPar
Makronun birinci parametresi asıl yapı nesnesi türünden göstericiyi, ikinci parametresi
\sphinxcode{\sphinxupquote{hlist\_node}} türünden göstericiyi, üçüncü parametresi \sphinxcode{\sphinxupquote{hlist\_head}} türünden kök düğümün
adresini, dördüncü parametresi de asıl yapıdaki bağ düğümünün ismini almaktadır. Bu
makronun \sphinxcode{\sphinxupquote{hlist\_for\_each\_entry}} makrosundan tek farkı zincir boşsa bir soruna yol açmadan
döngünün sonlanmasını sağlamasıdır. Döngü makrosunun içerisinde \sphinxcode{\sphinxupquote{hlist\_entry\_safe}}
makrosunun kullanıldığına dikkat ediniz. Bu makro zincirin sonundaki \sphinxcode{\sphinxupquote{NULL}} adresi de
dikkate almaktadır:

\begin{sphinxVerbatim}[commandchars=\\\{\}]
\PYG{c+cp}{\PYGZsh{}}\PYG{c+cp}{define hlist\PYGZus{}entry\PYGZus{}safe(ptr, type, member)                   \PYGZbs{}}
\PYG{c+cp}{(\PYGZob{} typeof(ptr) \PYGZus{}\PYGZus{}\PYGZus{}\PYGZus{}ptr = (ptr);                               \PYGZbs{}}
\PYG{c+cp}{     \PYGZus{}\PYGZus{}\PYGZus{}\PYGZus{}ptr ? hlist\PYGZus{}entry(\PYGZus{}\PYGZus{}\PYGZus{}\PYGZus{}ptr, type, member) : NULL;     \PYGZbs{}}
\PYG{c+cp}{\PYGZcb{})}
\end{sphinxVerbatim}

\sphinxAtStartPar
\sphinxcode{\sphinxupquote{list.h}} dosyası içerisinde hash tablolarına ilişkin başka yararlı fonksiyonlar da vardır.
Bu fonksiyonları ileride gerektiğinde açıklayacağız. Bunları siz de inceleyebilirsiniz.

\sphinxAtStartPar
Bir hash tablosunun tamamen yok edilmesi için yalnızca hash tablosunun değil onun tüm
zincirlerdeki elemanlarının da serbest bırakılması gerekir. Çekirdekte genel olarak böyle
bir gereksinim yoktur.


\subsection{Örnek Uygulama}
\label{\detokenize{kernellinkedlistandhashtable:ornek-uygulama}}
\sphinxAtStartPar
Aşağıda \sphinxcode{\sphinxupquote{list.h}} içerisindeki hash tablosu işlemleri için kullanıcı modunda
çalıştırılabilecek tam bir örnek verilmiştir.

\begin{sphinxVerbatim}[commandchars=\\\{\}]
\PYG{c+cp}{\PYGZsh{}}\PYG{c+cp}{include}\PYG{+w}{ }\PYG{c+cpf}{\PYGZlt{}stdio.h\PYGZgt{}}
\PYG{c+cp}{\PYGZsh{}}\PYG{c+cp}{include}\PYG{+w}{ }\PYG{c+cpf}{\PYGZlt{}stddef.h\PYGZgt{}}
\PYG{c+cp}{\PYGZsh{}}\PYG{c+cp}{include}\PYG{+w}{ }\PYG{c+cpf}{\PYGZlt{}stdlib.h\PYGZgt{}}
\PYG{c+cp}{\PYGZsh{}}\PYG{c+cp}{include}\PYG{+w}{ }\PYG{c+cpf}{\PYGZlt{}string.h\PYGZgt{}}
\PYG{c+cp}{\PYGZsh{}}\PYG{c+cp}{include}\PYG{+w}{ }\PYG{c+cpf}{\PYGZlt{}time.h\PYGZgt{}}

\PYG{c+cp}{\PYGZsh{}}\PYG{c+cp}{define TABLE\PYGZus{}SIZE    4096}

\PYG{k}{struct}\PYG{+w}{ }\PYG{n+nc}{hlist\PYGZus{}head}\PYG{+w}{ }\PYG{p}{\PYGZob{}}
\PYG{+w}{    }\PYG{k}{struct}\PYG{+w}{ }\PYG{n+nc}{hlist\PYGZus{}node}\PYG{+w}{ }\PYG{o}{*}\PYG{n}{first}\PYG{p}{;}
\PYG{p}{\PYGZcb{}}\PYG{p}{;}

\PYG{k}{struct}\PYG{+w}{ }\PYG{n+nc}{hlist\PYGZus{}node}\PYG{+w}{ }\PYG{p}{\PYGZob{}}
\PYG{+w}{    }\PYG{k}{struct}\PYG{+w}{ }\PYG{n+nc}{hlist\PYGZus{}node}\PYG{+w}{ }\PYG{o}{*}\PYG{n}{next}\PYG{p}{,}\PYG{+w}{ }\PYG{o}{*}\PYG{o}{*}\PYG{n}{pprev}\PYG{p}{;}
\PYG{p}{\PYGZcb{}}\PYG{p}{;}

\PYG{c+cp}{\PYGZsh{}}\PYG{c+cp}{define HLIST\PYGZus{}HEAD\PYGZus{}INIT \PYGZob{} .first = NULL \PYGZcb{}}
\PYG{c+cp}{\PYGZsh{}}\PYG{c+cp}{define HLIST\PYGZus{}HEAD(name) struct hlist\PYGZus{}head name = \PYGZob{} .first = NULL \PYGZcb{}}
\PYG{c+cp}{\PYGZsh{}}\PYG{c+cp}{define INIT\PYGZus{}HLIST\PYGZus{}HEAD(ptr) ((ptr)\PYGZhy{}\PYGZgt{}first = NULL)}

\PYG{c+cp}{\PYGZsh{}}\PYG{c+cp}{define WRITE\PYGZus{}ONCE(a, b)     ((a) = (b))   }\PYG{c+cm}{/* bize özgü */}
\PYG{c+cp}{\PYGZsh{}}\PYG{c+cp}{define LIST\PYGZus{}POISON1         (struct hlist\PYGZus{}node *)0x00100100}
\PYG{c+cp}{\PYGZsh{}}\PYG{c+cp}{define LIST\PYGZus{}POISON2         (struct hlist\PYGZus{}node **)0x00200200}

\PYG{k}{static}\PYG{+w}{ }\PYG{k+kr}{inline}\PYG{+w}{ }\PYG{k+kt}{void}\PYG{+w}{ }\PYG{n+nf}{hlist\PYGZus{}add\PYGZus{}head}\PYG{p}{(}\PYG{k}{struct}\PYG{+w}{ }\PYG{n+nc}{hlist\PYGZus{}node}\PYG{+w}{ }\PYG{o}{*}\PYG{n}{n}\PYG{p}{,}
\PYG{+w}{                                   }\PYG{k}{struct}\PYG{+w}{ }\PYG{n+nc}{hlist\PYGZus{}head}\PYG{+w}{ }\PYG{o}{*}\PYG{n}{h}\PYG{p}{)}
\PYG{p}{\PYGZob{}}
\PYG{+w}{    }\PYG{k}{struct}\PYG{+w}{ }\PYG{n+nc}{hlist\PYGZus{}node}\PYG{+w}{ }\PYG{o}{*}\PYG{n}{first}\PYG{+w}{ }\PYG{o}{=}\PYG{+w}{ }\PYG{n}{h}\PYG{o}{\PYGZhy{}}\PYG{o}{\PYGZgt{}}\PYG{n}{first}\PYG{p}{;}
\PYG{+w}{    }\PYG{n}{WRITE\PYGZus{}ONCE}\PYG{p}{(}\PYG{n}{n}\PYG{o}{\PYGZhy{}}\PYG{o}{\PYGZgt{}}\PYG{n}{next}\PYG{p}{,}\PYG{+w}{ }\PYG{n}{first}\PYG{p}{)}\PYG{p}{;}
\PYG{+w}{    }\PYG{k}{if}\PYG{+w}{ }\PYG{p}{(}\PYG{n}{first}\PYG{p}{)}
\PYG{+w}{        }\PYG{n}{WRITE\PYGZus{}ONCE}\PYG{p}{(}\PYG{n}{first}\PYG{o}{\PYGZhy{}}\PYG{o}{\PYGZgt{}}\PYG{n}{pprev}\PYG{p}{,}\PYG{+w}{ }\PYG{o}{\PYGZam{}}\PYG{n}{n}\PYG{o}{\PYGZhy{}}\PYG{o}{\PYGZgt{}}\PYG{n}{next}\PYG{p}{)}\PYG{p}{;}
\PYG{+w}{    }\PYG{n}{WRITE\PYGZus{}ONCE}\PYG{p}{(}\PYG{n}{h}\PYG{o}{\PYGZhy{}}\PYG{o}{\PYGZgt{}}\PYG{n}{first}\PYG{p}{,}\PYG{+w}{ }\PYG{n}{n}\PYG{p}{)}\PYG{p}{;}
\PYG{+w}{    }\PYG{n}{WRITE\PYGZus{}ONCE}\PYG{p}{(}\PYG{n}{n}\PYG{o}{\PYGZhy{}}\PYG{o}{\PYGZgt{}}\PYG{n}{pprev}\PYG{p}{,}\PYG{+w}{ }\PYG{o}{\PYGZam{}}\PYG{n}{h}\PYG{o}{\PYGZhy{}}\PYG{o}{\PYGZgt{}}\PYG{n}{first}\PYG{p}{)}\PYG{p}{;}
\PYG{p}{\PYGZcb{}}

\PYG{k}{static}\PYG{+w}{ }\PYG{k+kr}{inline}\PYG{+w}{ }\PYG{k+kt}{void}\PYG{+w}{ }\PYG{n+nf}{\PYGZus{}\PYGZus{}hlist\PYGZus{}del}\PYG{p}{(}\PYG{k}{struct}\PYG{+w}{ }\PYG{n+nc}{hlist\PYGZus{}node}\PYG{+w}{ }\PYG{o}{*}\PYG{n}{n}\PYG{p}{)}
\PYG{p}{\PYGZob{}}
\PYG{+w}{    }\PYG{k}{struct}\PYG{+w}{ }\PYG{n+nc}{hlist\PYGZus{}node}\PYG{+w}{ }\PYG{o}{*}\PYG{n}{next}\PYG{+w}{ }\PYG{o}{=}\PYG{+w}{ }\PYG{n}{n}\PYG{o}{\PYGZhy{}}\PYG{o}{\PYGZgt{}}\PYG{n}{next}\PYG{p}{;}
\PYG{+w}{    }\PYG{k}{struct}\PYG{+w}{ }\PYG{n+nc}{hlist\PYGZus{}node}\PYG{+w}{ }\PYG{o}{*}\PYG{o}{*}\PYG{n}{pprev}\PYG{+w}{ }\PYG{o}{=}\PYG{+w}{ }\PYG{n}{n}\PYG{o}{\PYGZhy{}}\PYG{o}{\PYGZgt{}}\PYG{n}{pprev}\PYG{p}{;}
\PYG{+w}{    }\PYG{n}{WRITE\PYGZus{}ONCE}\PYG{p}{(}\PYG{o}{*}\PYG{n}{pprev}\PYG{p}{,}\PYG{+w}{ }\PYG{n}{next}\PYG{p}{)}\PYG{p}{;}
\PYG{+w}{    }\PYG{k}{if}\PYG{+w}{ }\PYG{p}{(}\PYG{n}{next}\PYG{p}{)}
\PYG{+w}{        }\PYG{n}{WRITE\PYGZus{}ONCE}\PYG{p}{(}\PYG{n}{next}\PYG{o}{\PYGZhy{}}\PYG{o}{\PYGZgt{}}\PYG{n}{pprev}\PYG{p}{,}\PYG{+w}{ }\PYG{n}{pprev}\PYG{p}{)}\PYG{p}{;}
\PYG{p}{\PYGZcb{}}

\PYG{k}{static}\PYG{+w}{ }\PYG{k+kr}{inline}\PYG{+w}{ }\PYG{k+kt}{void}\PYG{+w}{ }\PYG{n+nf}{hlist\PYGZus{}del}\PYG{p}{(}\PYG{k}{struct}\PYG{+w}{ }\PYG{n+nc}{hlist\PYGZus{}node}\PYG{+w}{ }\PYG{o}{*}\PYG{n}{n}\PYG{p}{)}
\PYG{p}{\PYGZob{}}
\PYG{+w}{    }\PYG{n}{\PYGZus{}\PYGZus{}hlist\PYGZus{}del}\PYG{p}{(}\PYG{n}{n}\PYG{p}{)}\PYG{p}{;}
\PYG{+w}{    }\PYG{n}{n}\PYG{o}{\PYGZhy{}}\PYG{o}{\PYGZgt{}}\PYG{n}{next}\PYG{+w}{ }\PYG{o}{=}\PYG{+w}{ }\PYG{n}{LIST\PYGZus{}POISON1}\PYG{p}{;}
\PYG{+w}{    }\PYG{n}{n}\PYG{o}{\PYGZhy{}}\PYG{o}{\PYGZgt{}}\PYG{n}{pprev}\PYG{+w}{ }\PYG{o}{=}\PYG{+w}{ }\PYG{n}{LIST\PYGZus{}POISON2}\PYG{p}{;}
\PYG{p}{\PYGZcb{}}

\PYG{c+cp}{\PYGZsh{}}\PYG{c+cp}{define container\PYGZus{}of(ptr, type, member) (\PYGZob{}          \PYGZbs{}}
\PYG{c+cp}{    void *\PYGZus{}\PYGZus{}mptr = (void *)(ptr);                   \PYGZbs{}}
\PYG{c+cp}{    ((type *)(\PYGZus{}\PYGZus{}mptr \PYGZhy{} offsetof(type, member))); \PYGZcb{})}

\PYG{c+cp}{\PYGZsh{}}\PYG{c+cp}{define hlist\PYGZus{}entry(ptr, type, member) \PYGZbs{}}
\PYG{c+cp}{    container\PYGZus{}of(ptr, type, member)}

\PYG{c+cp}{\PYGZsh{}}\PYG{c+cp}{define hlist\PYGZus{}entry\PYGZus{}safe(ptr, type, member) \PYGZbs{}}
\PYG{c+cp}{    (\PYGZob{} typeof(ptr) \PYGZus{}\PYGZus{}\PYGZus{}\PYGZus{}ptr = (ptr); \PYGZbs{}}
\PYG{c+cp}{       \PYGZus{}\PYGZus{}\PYGZus{}\PYGZus{}ptr ? hlist\PYGZus{}entry(\PYGZus{}\PYGZus{}\PYGZus{}\PYGZus{}ptr, type, member) : NULL; \PYGZbs{}}
\PYG{c+cp}{    \PYGZcb{})}

\PYG{c+cp}{\PYGZsh{}}\PYG{c+cp}{define hlist\PYGZus{}for\PYGZus{}each(pos, head) \PYGZbs{}}
\PYG{c+cp}{    for (pos = (head)\PYGZhy{}\PYGZgt{}first; pos ; pos = pos\PYGZhy{}\PYGZgt{}next)}

\PYG{c+cp}{\PYGZsh{}}\PYG{c+cp}{define hlist\PYGZus{}for\PYGZus{}each\PYGZus{}safe(pos, n, head) \PYGZbs{}}
\PYG{c+cp}{    for (pos = (head)\PYGZhy{}\PYGZgt{}first; pos \PYGZam{}\PYGZam{} (\PYGZob{} n = pos\PYGZhy{}\PYGZgt{}next; 1; \PYGZcb{}); \PYGZbs{}}
\PYG{c+cp}{        pos = n)}

\PYG{c+cp}{\PYGZsh{}}\PYG{c+cp}{define hlist\PYGZus{}for\PYGZus{}each\PYGZus{}entry(pos, head, member)                    \PYGZbs{}}
\PYG{c+cp}{    for (pos = hlist\PYGZus{}entry((head)\PYGZhy{}\PYGZgt{}first, typeof(*(pos)), member); \PYGZbs{}}
\PYG{c+cp}{        pos;                                                        \PYGZbs{}}
\PYG{c+cp}{        pos = hlist\PYGZus{}entry((pos)\PYGZhy{}\PYGZgt{}member.next, typeof(*(pos)), member))}

\PYG{c+cp}{\PYGZsh{}}\PYG{c+cp}{define hlist\PYGZus{}for\PYGZus{}each\PYGZus{}entry\PYGZus{}safe(pos, n, head, member)             \PYGZbs{}}
\PYG{c+cp}{    for (pos = hlist\PYGZus{}entry\PYGZus{}safe((head)\PYGZhy{}\PYGZgt{}first, typeof(*pos), member);\PYGZbs{}}
\PYG{c+cp}{        pos \PYGZam{}\PYGZam{} (\PYGZob{} n = pos\PYGZhy{}\PYGZgt{}member.next; 1; \PYGZcb{});                       \PYGZbs{}}
\PYG{c+cp}{        pos = hlist\PYGZus{}entry\PYGZus{}safe(n, typeof(*pos), member))}

\PYG{c+cm}{/* Test code */}

\PYG{k}{struct}\PYG{+w}{ }\PYG{n+nc}{PERSON}\PYG{+w}{ }\PYG{p}{\PYGZob{}}
\PYG{+w}{    }\PYG{k+kt}{char}\PYG{+w}{ }\PYG{n}{name}\PYG{p}{[}\PYG{l+m+mi}{32}\PYG{p}{]}\PYG{p}{;}
\PYG{+w}{    }\PYG{k+kt}{int}\PYG{+w}{ }\PYG{n}{no}\PYG{p}{;}
\PYG{+w}{    }\PYG{k}{struct}\PYG{+w}{ }\PYG{n+nc}{hlist\PYGZus{}node}\PYG{+w}{ }\PYG{n}{hlink}\PYG{p}{;}
\PYG{p}{\PYGZcb{}}\PYG{p}{;}

\PYG{k+kt}{void}\PYG{+w}{ }\PYG{n+nf}{set\PYGZus{}random\PYGZus{}person}\PYG{p}{(}\PYG{k}{struct}\PYG{+w}{ }\PYG{n+nc}{PERSON}\PYG{+w}{ }\PYG{o}{*}\PYG{n}{per}\PYG{p}{)}\PYG{p}{;}
\PYG{k+kt}{unsigned}\PYG{+w}{ }\PYG{k+kt}{int}\PYG{+w}{ }\PYG{n+nf}{hash\PYGZus{}func}\PYG{p}{(}\PYG{k+kt}{unsigned}\PYG{+w}{ }\PYG{k+kt}{int}\PYG{+w}{ }\PYG{n}{key}\PYG{p}{)}\PYG{p}{;}

\PYG{k+kt}{int}\PYG{+w}{ }\PYG{n+nf}{main}\PYG{p}{(}\PYG{k+kt}{void}\PYG{p}{)}
\PYG{p}{\PYGZob{}}
\PYG{+w}{    }\PYG{k}{struct}\PYG{+w}{ }\PYG{n+nc}{hlist\PYGZus{}head}\PYG{+w}{ }\PYG{o}{*}\PYG{n}{hash\PYGZus{}table}\PYG{p}{;}
\PYG{+w}{    }\PYG{k}{struct}\PYG{+w}{ }\PYG{n+nc}{PERSON}\PYG{+w}{ }\PYG{o}{*}\PYG{n}{per}\PYG{p}{;}
\PYG{+w}{    }\PYG{k+kt}{unsigned}\PYG{+w}{ }\PYG{k+kt}{int}\PYG{+w}{ }\PYG{n}{hash}\PYG{p}{;}
\PYG{+w}{    }\PYG{k+kt}{int}\PYG{+w}{ }\PYG{n}{no}\PYG{p}{;}

\PYG{+w}{    }\PYG{n}{srand}\PYG{p}{(}\PYG{n}{time}\PYG{p}{(}\PYG{n+nb}{NULL}\PYG{p}{)}\PYG{p}{)}\PYG{p}{;}

\PYG{+w}{    }\PYG{k}{if}\PYG{+w}{ }\PYG{p}{(}\PYG{p}{(}\PYG{n}{hash\PYGZus{}table}\PYG{+w}{ }\PYG{o}{=}\PYG{+w}{ }\PYG{p}{(}\PYG{k}{struct}\PYG{+w}{ }\PYG{n+nc}{hlist\PYGZus{}head}\PYG{+w}{ }\PYG{o}{*}\PYG{p}{)}\PYG{n}{malloc}\PYG{p}{(}
\PYG{+w}{            }\PYG{k}{sizeof}\PYG{p}{(}\PYG{k}{struct}\PYG{+w}{ }\PYG{n+nc}{hlist\PYGZus{}head}\PYG{p}{)}\PYG{+w}{ }\PYG{o}{*}\PYG{+w}{ }\PYG{n}{TABLE\PYGZus{}SIZE}\PYG{p}{)}\PYG{p}{)}\PYG{+w}{ }\PYG{o}{=}\PYG{o}{=}\PYG{+w}{ }\PYG{n+nb}{NULL}\PYG{p}{)}\PYG{+w}{ }\PYG{p}{\PYGZob{}}
\PYG{+w}{        }\PYG{n}{fprintf}\PYG{p}{(}\PYG{n}{stderr}\PYG{p}{,}\PYG{+w}{ }\PYG{l+s}{\PYGZdq{}}\PYG{l+s}{cannot allocate memory!...}\PYG{l+s+se}{\PYGZbs{}n}\PYG{l+s}{\PYGZdq{}}\PYG{p}{)}\PYG{p}{;}
\PYG{+w}{        }\PYG{n}{exit}\PYG{p}{(}\PYG{n}{EXIT\PYGZus{}FAILURE}\PYG{p}{)}\PYG{p}{;}
\PYG{+w}{    }\PYG{p}{\PYGZcb{}}

\PYG{+w}{    }\PYG{k}{for}\PYG{+w}{ }\PYG{p}{(}\PYG{k+kt}{int}\PYG{+w}{ }\PYG{n}{i}\PYG{+w}{ }\PYG{o}{=}\PYG{+w}{ }\PYG{l+m+mi}{0}\PYG{p}{;}\PYG{+w}{ }\PYG{n}{i}\PYG{+w}{ }\PYG{o}{\PYGZlt{}}\PYG{+w}{ }\PYG{n}{TABLE\PYGZus{}SIZE}\PYG{p}{;}\PYG{+w}{ }\PYG{o}{+}\PYG{o}{+}\PYG{n}{i}\PYG{p}{)}
\PYG{+w}{        }\PYG{n}{INIT\PYGZus{}HLIST\PYGZus{}HEAD}\PYG{p}{(}\PYG{o}{\PYGZam{}}\PYG{n}{hash\PYGZus{}table}\PYG{p}{[}\PYG{n}{i}\PYG{p}{]}\PYG{p}{)}\PYG{p}{;}

\PYG{+w}{    }\PYG{k}{for}\PYG{+w}{ }\PYG{p}{(}\PYG{k+kt}{int}\PYG{+w}{ }\PYG{n}{i}\PYG{+w}{ }\PYG{o}{=}\PYG{+w}{ }\PYG{l+m+mi}{0}\PYG{p}{;}\PYG{+w}{ }\PYG{n}{i}\PYG{+w}{ }\PYG{o}{\PYGZlt{}}\PYG{+w}{ }\PYG{l+m+mi}{100}\PYG{p}{;}\PYG{+w}{ }\PYG{o}{+}\PYG{o}{+}\PYG{n}{i}\PYG{p}{)}\PYG{+w}{ }\PYG{p}{\PYGZob{}}
\PYG{+w}{        }\PYG{k}{if}\PYG{+w}{ }\PYG{p}{(}\PYG{p}{(}\PYG{n}{per}\PYG{+w}{ }\PYG{o}{=}\PYG{+w}{ }\PYG{p}{(}\PYG{k}{struct}\PYG{+w}{ }\PYG{n+nc}{PERSON}\PYG{+w}{ }\PYG{o}{*}\PYG{p}{)}\PYG{n}{malloc}\PYG{p}{(}\PYG{k}{sizeof}\PYG{p}{(}\PYG{k}{struct}\PYG{+w}{ }\PYG{n+nc}{PERSON}\PYG{p}{)}\PYG{p}{)}\PYG{p}{)}\PYG{+w}{ }\PYG{o}{=}\PYG{o}{=}\PYG{+w}{ }\PYG{n+nb}{NULL}\PYG{p}{)}\PYG{+w}{ }\PYG{p}{\PYGZob{}}
\PYG{+w}{            }\PYG{n}{fprintf}\PYG{p}{(}\PYG{n}{stderr}\PYG{p}{,}\PYG{+w}{ }\PYG{l+s}{\PYGZdq{}}\PYG{l+s}{cannot allocate memory!...}\PYG{l+s+se}{\PYGZbs{}n}\PYG{l+s}{\PYGZdq{}}\PYG{p}{)}\PYG{p}{;}
\PYG{+w}{            }\PYG{n}{exit}\PYG{p}{(}\PYG{n}{EXIT\PYGZus{}FAILURE}\PYG{p}{)}\PYG{p}{;}
\PYG{+w}{        }\PYG{p}{\PYGZcb{}}
\PYG{+w}{        }\PYG{k}{if}\PYG{+w}{ }\PYG{p}{(}\PYG{n}{i}\PYG{+w}{ }\PYG{o}{=}\PYG{o}{=}\PYG{+w}{ }\PYG{l+m+mi}{50}\PYG{p}{)}\PYG{+w}{ }\PYG{p}{\PYGZob{}}
\PYG{+w}{            }\PYG{n}{strcpy}\PYG{p}{(}\PYG{n}{per}\PYG{o}{\PYGZhy{}}\PYG{o}{\PYGZgt{}}\PYG{n}{name}\PYG{p}{,}\PYG{+w}{ }\PYG{l+s}{\PYGZdq{}}\PYG{l+s}{ALI SERCE}\PYG{l+s}{\PYGZdq{}}\PYG{p}{)}\PYG{p}{;}
\PYG{+w}{            }\PYG{n}{per}\PYG{o}{\PYGZhy{}}\PYG{o}{\PYGZgt{}}\PYG{n}{no}\PYG{+w}{ }\PYG{o}{=}\PYG{+w}{ }\PYG{l+m+mi}{12345678}\PYG{p}{;}
\PYG{+w}{        }\PYG{p}{\PYGZcb{}}
\PYG{+w}{        }\PYG{k}{else}
\PYG{+w}{            }\PYG{n}{set\PYGZus{}random\PYGZus{}person}\PYG{p}{(}\PYG{n}{per}\PYG{p}{)}\PYG{p}{;}
\PYG{+w}{        }\PYG{n}{hash}\PYG{+w}{ }\PYG{o}{=}\PYG{+w}{ }\PYG{n}{hash\PYGZus{}func}\PYG{p}{(}\PYG{n}{per}\PYG{o}{\PYGZhy{}}\PYG{o}{\PYGZgt{}}\PYG{n}{no}\PYG{p}{)}\PYG{p}{;}
\PYG{+w}{        }\PYG{n}{hlist\PYGZus{}add\PYGZus{}head}\PYG{p}{(}\PYG{o}{\PYGZam{}}\PYG{n}{per}\PYG{o}{\PYGZhy{}}\PYG{o}{\PYGZgt{}}\PYG{n}{hlink}\PYG{p}{,}\PYG{+w}{ }\PYG{o}{\PYGZam{}}\PYG{n}{hash\PYGZus{}table}\PYG{p}{[}\PYG{n}{hash}\PYG{p}{]}\PYG{p}{)}\PYG{p}{;}
\PYG{+w}{    }\PYG{p}{\PYGZcb{}}

\PYG{+w}{    }\PYG{c+cm}{/* hlist\PYGZus{}for\PYGZus{}each ile arama */}
\PYG{+w}{    }\PYG{p}{\PYGZob{}}
\PYG{+w}{        }\PYG{k}{struct}\PYG{+w}{ }\PYG{n+nc}{hlist\PYGZus{}node}\PYG{+w}{ }\PYG{o}{*}\PYG{n}{node}\PYG{p}{;}
\PYG{+w}{        }\PYG{k}{struct}\PYG{+w}{ }\PYG{n+nc}{PERSON}\PYG{+w}{ }\PYG{o}{*}\PYG{n}{per\PYGZus{}find}\PYG{p}{;}

\PYG{+w}{        }\PYG{n}{printf}\PYG{p}{(}\PYG{l+s}{\PYGZdq{}}\PYG{l+s}{Person no:}\PYG{l+s}{\PYGZdq{}}\PYG{p}{)}\PYG{p}{;}
\PYG{+w}{        }\PYG{n}{scanf}\PYG{p}{(}\PYG{l+s}{\PYGZdq{}}\PYG{l+s}{\PYGZpc{}d}\PYG{l+s}{\PYGZdq{}}\PYG{p}{,}\PYG{+w}{ }\PYG{o}{\PYGZam{}}\PYG{n}{no}\PYG{p}{)}\PYG{p}{;}

\PYG{+w}{        }\PYG{n}{hash}\PYG{+w}{ }\PYG{o}{=}\PYG{+w}{ }\PYG{n}{hash\PYGZus{}func}\PYG{p}{(}\PYG{n}{no}\PYG{p}{)}\PYG{p}{;}

\PYG{+w}{        }\PYG{n}{hlist\PYGZus{}for\PYGZus{}each}\PYG{p}{(}\PYG{n}{node}\PYG{p}{,}\PYG{+w}{ }\PYG{o}{\PYGZam{}}\PYG{n}{hash\PYGZus{}table}\PYG{p}{[}\PYG{n}{hash}\PYG{p}{]}\PYG{p}{)}\PYG{+w}{ }\PYG{p}{\PYGZob{}}
\PYG{+w}{            }\PYG{n}{per\PYGZus{}find}\PYG{+w}{ }\PYG{o}{=}\PYG{+w}{ }\PYG{n}{hlist\PYGZus{}entry}\PYG{p}{(}\PYG{n}{node}\PYG{p}{,}\PYG{+w}{ }\PYG{k}{struct}\PYG{+w}{ }\PYG{n+nc}{PERSON}\PYG{p}{,}\PYG{+w}{ }\PYG{n}{hlink}\PYG{p}{)}\PYG{p}{;}
\PYG{+w}{            }\PYG{k}{if}\PYG{+w}{ }\PYG{p}{(}\PYG{n}{per\PYGZus{}find}\PYG{o}{\PYGZhy{}}\PYG{o}{\PYGZgt{}}\PYG{n}{no}\PYG{+w}{ }\PYG{o}{=}\PYG{o}{=}\PYG{+w}{ }\PYG{n}{no}\PYG{p}{)}\PYG{+w}{ }\PYG{p}{\PYGZob{}}
\PYG{+w}{                }\PYG{n}{printf}\PYG{p}{(}\PYG{l+s}{\PYGZdq{}}\PYG{l+s}{Found: \PYGZpc{}s, \PYGZpc{}d}\PYG{l+s+se}{\PYGZbs{}n}\PYG{l+s}{\PYGZdq{}}\PYG{p}{,}\PYG{+w}{ }\PYG{n}{per\PYGZus{}find}\PYG{o}{\PYGZhy{}}\PYG{o}{\PYGZgt{}}\PYG{n}{name}\PYG{p}{,}\PYG{+w}{ }\PYG{n}{per\PYGZus{}find}\PYG{o}{\PYGZhy{}}\PYG{o}{\PYGZgt{}}\PYG{n}{no}\PYG{p}{)}\PYG{p}{;}
\PYG{+w}{                }\PYG{k}{break}\PYG{p}{;}
\PYG{+w}{            }\PYG{p}{\PYGZcb{}}
\PYG{+w}{        }\PYG{p}{\PYGZcb{}}
\PYG{+w}{        }\PYG{k}{if}\PYG{+w}{ }\PYG{p}{(}\PYG{n}{node}\PYG{+w}{ }\PYG{o}{=}\PYG{o}{=}\PYG{+w}{ }\PYG{n+nb}{NULL}\PYG{p}{)}
\PYG{+w}{            }\PYG{n}{printf}\PYG{p}{(}\PYG{l+s}{\PYGZdq{}}\PYG{l+s}{cannot find...}\PYG{l+s+se}{\PYGZbs{}n}\PYG{l+s}{\PYGZdq{}}\PYG{p}{)}\PYG{p}{;}
\PYG{+w}{    }\PYG{p}{\PYGZcb{}}

\PYG{+w}{    }\PYG{c+cm}{/* hlist\PYGZus{}for\PYGZus{}each\PYGZus{}entry\PYGZus{}safe ile arama */}
\PYG{+w}{    }\PYG{p}{\PYGZob{}}
\PYG{+w}{        }\PYG{k}{struct}\PYG{+w}{ }\PYG{n+nc}{PERSON}\PYG{+w}{ }\PYG{o}{*}\PYG{n}{per\PYGZus{}find}\PYG{p}{;}
\PYG{+w}{        }\PYG{k}{struct}\PYG{+w}{ }\PYG{n+nc}{hlist\PYGZus{}node}\PYG{+w}{ }\PYG{o}{*}\PYG{n}{node}\PYG{p}{;}

\PYG{+w}{        }\PYG{n}{hash}\PYG{+w}{ }\PYG{o}{=}\PYG{+w}{ }\PYG{n}{hash\PYGZus{}func}\PYG{p}{(}\PYG{n}{no}\PYG{p}{)}\PYG{p}{;}

\PYG{+w}{        }\PYG{n}{hlist\PYGZus{}for\PYGZus{}each\PYGZus{}entry\PYGZus{}safe}\PYG{p}{(}\PYG{n}{per\PYGZus{}find}\PYG{p}{,}\PYG{+w}{ }\PYG{n}{node}\PYG{p}{,}\PYG{+w}{ }\PYG{o}{\PYGZam{}}\PYG{n}{hash\PYGZus{}table}\PYG{p}{[}\PYG{n}{hash}\PYG{p}{]}\PYG{p}{,}\PYG{+w}{ }\PYG{n}{hlink}\PYG{p}{)}\PYG{+w}{ }\PYG{p}{\PYGZob{}}
\PYG{+w}{            }\PYG{k}{if}\PYG{+w}{ }\PYG{p}{(}\PYG{n}{per\PYGZus{}find}\PYG{o}{\PYGZhy{}}\PYG{o}{\PYGZgt{}}\PYG{n}{no}\PYG{+w}{ }\PYG{o}{=}\PYG{o}{=}\PYG{+w}{ }\PYG{n}{no}\PYG{p}{)}\PYG{+w}{ }\PYG{p}{\PYGZob{}}
\PYG{+w}{                }\PYG{n}{printf}\PYG{p}{(}\PYG{l+s}{\PYGZdq{}}\PYG{l+s}{Found: \PYGZpc{}s, \PYGZpc{}d}\PYG{l+s+se}{\PYGZbs{}n}\PYG{l+s}{\PYGZdq{}}\PYG{p}{,}\PYG{+w}{ }\PYG{n}{per\PYGZus{}find}\PYG{o}{\PYGZhy{}}\PYG{o}{\PYGZgt{}}\PYG{n}{name}\PYG{p}{,}\PYG{+w}{ }\PYG{n}{per\PYGZus{}find}\PYG{o}{\PYGZhy{}}\PYG{o}{\PYGZgt{}}\PYG{n}{no}\PYG{p}{)}\PYG{p}{;}
\PYG{+w}{                }\PYG{k}{break}\PYG{p}{;}
\PYG{+w}{            }\PYG{p}{\PYGZcb{}}
\PYG{+w}{        }\PYG{p}{\PYGZcb{}}
\PYG{+w}{        }\PYG{k}{if}\PYG{+w}{ }\PYG{p}{(}\PYG{n}{per\PYGZus{}find}\PYG{+w}{ }\PYG{o}{=}\PYG{o}{=}\PYG{+w}{ }\PYG{n+nb}{NULL}\PYG{p}{)}
\PYG{+w}{            }\PYG{n}{printf}\PYG{p}{(}\PYG{l+s}{\PYGZdq{}}\PYG{l+s}{cannot find...}\PYG{l+s+se}{\PYGZbs{}n}\PYG{l+s}{\PYGZdq{}}\PYG{p}{)}\PYG{p}{;}
\PYG{+w}{    }\PYG{p}{\PYGZcb{}}

\PYG{+w}{    }\PYG{c+cm}{/* Tüm tabloyu serbest bırakma */}
\PYG{+w}{    }\PYG{p}{\PYGZob{}}
\PYG{+w}{        }\PYG{k}{struct}\PYG{+w}{ }\PYG{n+nc}{PERSON}\PYG{+w}{ }\PYG{o}{*}\PYG{n}{per}\PYG{p}{,}\PYG{+w}{ }\PYG{o}{*}\PYG{n}{temp}\PYG{p}{;}
\PYG{+w}{        }\PYG{k}{struct}\PYG{+w}{ }\PYG{n+nc}{hlist\PYGZus{}node}\PYG{+w}{ }\PYG{o}{*}\PYG{n}{node}\PYG{p}{;}

\PYG{+w}{        }\PYG{k}{for}\PYG{+w}{ }\PYG{p}{(}\PYG{k+kt}{int}\PYG{+w}{ }\PYG{n}{i}\PYG{+w}{ }\PYG{o}{=}\PYG{+w}{ }\PYG{l+m+mi}{0}\PYG{p}{;}\PYG{+w}{ }\PYG{n}{i}\PYG{+w}{ }\PYG{o}{\PYGZlt{}}\PYG{+w}{ }\PYG{n}{TABLE\PYGZus{}SIZE}\PYG{p}{;}\PYG{+w}{ }\PYG{o}{+}\PYG{o}{+}\PYG{n}{i}\PYG{p}{)}\PYG{+w}{ }\PYG{p}{\PYGZob{}}
\PYG{+w}{            }\PYG{n}{temp}\PYG{+w}{ }\PYG{o}{=}\PYG{+w}{ }\PYG{n+nb}{NULL}\PYG{p}{;}
\PYG{+w}{            }\PYG{n}{hlist\PYGZus{}for\PYGZus{}each\PYGZus{}entry\PYGZus{}safe}\PYG{p}{(}\PYG{n}{per}\PYG{p}{,}\PYG{+w}{ }\PYG{n}{node}\PYG{p}{,}\PYG{+w}{ }\PYG{o}{\PYGZam{}}\PYG{n}{hash\PYGZus{}table}\PYG{p}{[}\PYG{n}{i}\PYG{p}{]}\PYG{p}{,}\PYG{+w}{ }\PYG{n}{hlink}\PYG{p}{)}\PYG{+w}{ }\PYG{p}{\PYGZob{}}
\PYG{+w}{                }\PYG{n}{free}\PYG{p}{(}\PYG{n}{temp}\PYG{p}{)}\PYG{p}{;}
\PYG{+w}{                }\PYG{n}{temp}\PYG{+w}{ }\PYG{o}{=}\PYG{+w}{ }\PYG{n}{per}\PYG{p}{;}
\PYG{+w}{            }\PYG{p}{\PYGZcb{}}
\PYG{+w}{            }\PYG{n}{free}\PYG{p}{(}\PYG{n}{temp}\PYG{p}{)}\PYG{p}{;}
\PYG{+w}{        }\PYG{p}{\PYGZcb{}}
\PYG{+w}{        }\PYG{n}{free}\PYG{p}{(}\PYG{n}{hash\PYGZus{}table}\PYG{p}{)}\PYG{p}{;}
\PYG{+w}{    }\PYG{p}{\PYGZcb{}}

\PYG{+w}{    }\PYG{k}{return}\PYG{+w}{ }\PYG{l+m+mi}{0}\PYG{p}{;}
\PYG{p}{\PYGZcb{}}

\PYG{k+kt}{void}\PYG{+w}{ }\PYG{n+nf}{set\PYGZus{}random\PYGZus{}person}\PYG{p}{(}\PYG{k}{struct}\PYG{+w}{ }\PYG{n+nc}{PERSON}\PYG{+w}{ }\PYG{o}{*}\PYG{n}{per}\PYG{p}{)}
\PYG{p}{\PYGZob{}}
\PYG{+w}{    }\PYG{k+kt}{int}\PYG{+w}{ }\PYG{n}{i}\PYG{p}{;}

\PYG{+w}{    }\PYG{k}{for}\PYG{+w}{ }\PYG{p}{(}\PYG{n}{i}\PYG{+w}{ }\PYG{o}{=}\PYG{+w}{ }\PYG{l+m+mi}{0}\PYG{p}{;}\PYG{+w}{ }\PYG{n}{i}\PYG{+w}{ }\PYG{o}{\PYGZlt{}}\PYG{+w}{ }\PYG{l+m+mi}{31}\PYG{p}{;}\PYG{+w}{ }\PYG{o}{+}\PYG{o}{+}\PYG{n}{i}\PYG{p}{)}
\PYG{+w}{        }\PYG{n}{per}\PYG{o}{\PYGZhy{}}\PYG{o}{\PYGZgt{}}\PYG{n}{name}\PYG{p}{[}\PYG{n}{i}\PYG{p}{]}\PYG{+w}{ }\PYG{o}{=}\PYG{+w}{ }\PYG{n}{rand}\PYG{p}{(}\PYG{p}{)}\PYG{+w}{ }\PYG{o}{\PYGZpc{}}\PYG{+w}{ }\PYG{l+m+mi}{26}\PYG{+w}{ }\PYG{o}{+}\PYG{+w}{ }\PYG{l+s+sc}{\PYGZsq{}}\PYG{l+s+sc}{A}\PYG{l+s+sc}{\PYGZsq{}}\PYG{p}{;}
\PYG{+w}{    }\PYG{n}{per}\PYG{o}{\PYGZhy{}}\PYG{o}{\PYGZgt{}}\PYG{n}{name}\PYG{p}{[}\PYG{n}{i}\PYG{p}{]}\PYG{+w}{ }\PYG{o}{=}\PYG{+w}{ }\PYG{l+s+sc}{\PYGZsq{}}\PYG{l+s+sc}{\PYGZbs{}0}\PYG{l+s+sc}{\PYGZsq{}}\PYG{p}{;}

\PYG{+w}{    }\PYG{n}{per}\PYG{o}{\PYGZhy{}}\PYG{o}{\PYGZgt{}}\PYG{n}{no}\PYG{+w}{ }\PYG{o}{=}\PYG{+w}{ }\PYG{n}{rand}\PYG{p}{(}\PYG{p}{)}\PYG{+w}{ }\PYG{o}{\PYGZpc{}}\PYG{+w}{ }\PYG{l+m+mi}{1000000}\PYG{p}{;}
\PYG{p}{\PYGZcb{}}

\PYG{k+kt}{unsigned}\PYG{+w}{ }\PYG{k+kt}{int}\PYG{+w}{ }\PYG{n+nf}{hash\PYGZus{}func}\PYG{p}{(}\PYG{k+kt}{unsigned}\PYG{+w}{ }\PYG{k+kt}{int}\PYG{+w}{ }\PYG{n}{key}\PYG{p}{)}
\PYG{p}{\PYGZob{}}
\PYG{+w}{    }\PYG{n}{key}\PYG{+w}{ }\PYG{o}{=}\PYG{+w}{ }\PYG{p}{(}\PYG{n}{key}\PYG{+w}{ }\PYG{o}{\PYGZca{}}\PYG{+w}{ }\PYG{l+m+mi}{61}\PYG{p}{)}\PYG{+w}{ }\PYG{o}{\PYGZca{}}\PYG{+w}{ }\PYG{p}{(}\PYG{n}{key}\PYG{+w}{ }\PYG{o}{\PYGZgt{}}\PYG{o}{\PYGZgt{}}\PYG{+w}{ }\PYG{l+m+mi}{16}\PYG{p}{)}\PYG{p}{;}
\PYG{+w}{    }\PYG{n}{key}\PYG{+w}{ }\PYG{o}{=}\PYG{+w}{ }\PYG{n}{key}\PYG{+w}{ }\PYG{o}{+}\PYG{+w}{ }\PYG{p}{(}\PYG{n}{key}\PYG{+w}{ }\PYG{o}{\PYGZlt{}}\PYG{o}{\PYGZlt{}}\PYG{+w}{ }\PYG{l+m+mi}{3}\PYG{p}{)}\PYG{p}{;}
\PYG{+w}{    }\PYG{n}{key}\PYG{+w}{ }\PYG{o}{=}\PYG{+w}{ }\PYG{n}{key}\PYG{+w}{ }\PYG{o}{\PYGZca{}}\PYG{+w}{ }\PYG{p}{(}\PYG{n}{key}\PYG{+w}{ }\PYG{o}{\PYGZgt{}}\PYG{o}{\PYGZgt{}}\PYG{+w}{ }\PYG{l+m+mi}{4}\PYG{p}{)}\PYG{p}{;}
\PYG{+w}{    }\PYG{n}{key}\PYG{+w}{ }\PYG{o}{=}\PYG{+w}{ }\PYG{n}{key}\PYG{+w}{ }\PYG{o}{*}\PYG{+w}{ }\PYG{l+m+mh}{0x27d4eb2d}\PYG{p}{;}
\PYG{+w}{    }\PYG{n}{key}\PYG{+w}{ }\PYG{o}{=}\PYG{+w}{ }\PYG{n}{key}\PYG{+w}{ }\PYG{o}{\PYGZca{}}\PYG{+w}{ }\PYG{p}{(}\PYG{n}{key}\PYG{+w}{ }\PYG{o}{\PYGZgt{}}\PYG{o}{\PYGZgt{}}\PYG{+w}{ }\PYG{l+m+mi}{15}\PYG{p}{)}\PYG{p}{;}

\PYG{+w}{    }\PYG{k}{return}\PYG{+w}{ }\PYG{n}{key}\PYG{+w}{ }\PYG{o}{\PYGZpc{}}\PYG{+w}{ }\PYG{n}{TABLE\PYGZus{}SIZE}\PYG{p}{;}
\PYG{p}{\PYGZcb{}}
\end{sphinxVerbatim}

\sphinxAtStartPar
Örnekte \sphinxcode{\sphinxupquote{TABLE\_SIZE}} uzunluğunda, her bir elemanı \sphinxcode{\sphinxupquote{hlist\_head}} türünden olan bir dizi
oluşturulmuştur. Sonra bu \sphinxcode{\sphinxupquote{hlist\_head}} elemanlarına ilk değer verilmiştir (yani onların
\sphinxcode{\sphinxupquote{first}} göstericilerine \sphinxcode{\sphinxupquote{NULL}} adres yerleştirilmiştir). Ondan sonra hash tablosuna
rastgele 100 eleman eklenmiştir; ancak bunun yanı sıra belli bir eleman (numarası 12345678
olan “ALI SERCE”) da listeye ayrıca eklenmiştir. Program biterken tüm hash tablosu
zincirleriyle birlikte serbest bırakılmıştır.

\begin{sphinxadmonition}{note}{Not:}
\sphinxAtStartPar
Bağlı liste düğümleri serbest bırakılırken sonraki düğümün adresini saklamak gerekir.
\sphinxcode{\sphinxupquote{NULL}} adrese \sphinxcode{\sphinxupquote{free}} uygulamanın bir soruna yol açmayacağını anımsayınız.
\end{sphinxadmonition}


\section{RCU’lu Hash Tablosu Fonksiyonları}
\label{\detokenize{kernellinkedlistandhashtable:rcu-lu-hash-tablosu-fonksiyonlari}}
\sphinxAtStartPar
Linux çekirdeğine kilitsiz (lock\sphinxhyphen{}free) RCU mekanizması eklendiğinde tıpkı bağlı listelerde
olduğu gibi hash tablolarına da yukarıdaki fonksiyonların sonu \sphinxcode{\sphinxupquote{\_rcu}} ile biten RCU uyumlu
versiyonları eklenmiştir. Bu fonksiyonlar \sphinxcode{\sphinxupquote{include/linux/rculist.h}} dosyası içerisindedir.
Bu dosyada yukarıda gördüğümüz hash tablosu fonksiyonlarının RCU mekanizmalı versiyonları
bulunmaktadır. Biz buradaki fonksiyonların yalnızca isimlerini vereceğiz. RCU mekanizması
başka bir başlık altında ele alınacaktır:

\begin{sphinxVerbatim}[commandchars=\\\{\}]
\PYG{n}{hlist\PYGZus{}add\PYGZus{}head\PYGZus{}rcu}
\PYG{n}{hlist\PYGZus{}del\PYGZus{}rcu}
\PYG{n}{hlist\PYGZus{}add\PYGZus{}before\PYGZus{}rcu}
\PYG{n}{hlist\PYGZus{}add\PYGZus{}behind\PYGZus{}rcu}
\PYG{n}{hlist\PYGZus{}for\PYGZus{}each\PYGZus{}entry\PYGZus{}rcu}
\PYG{n}{hlist\PYGZus{}for\PYGZus{}each\PYGZus{}entry\PYGZus{}srcu}\PYG{+w}{    }\PYG{c+cm}{/* safe version */}

\PYG{n}{hlist\PYGZus{}first\PYGZus{}rcu}\PYG{p}{(}\PYG{n}{head}\PYG{p}{)}
\PYG{n}{hlist\PYGZus{}next\PYGZus{}rcu}\PYG{p}{(}\PYG{n}{node}\PYG{p}{)}
\PYG{n}{hlist\PYGZus{}pprev\PYGZus{}rcu}\PYG{p}{(}\PYG{n}{node}\PYG{p}{)}
\end{sphinxVerbatim}

\sphinxAtStartPar
Bu RCU’lu fonksiyonlar tıpkı bağlı listelerde olduğu gibi birden fazla okuyan ancak tek bir
yazan taraf varsa beklemeye yol açmamaktadır. Tabii birden fazla yazan tarafın ayrıca bir
senkronizasyon nesnesiyle korunması gerekir. Bu makrolarla dolaşım yapılırken yine bağlı
listelerde olduğu gibi ilgili kod bloğunun başına ve sonuna \sphinxcode{\sphinxupquote{rcu\_read\_lock}} ve
\sphinxcode{\sphinxupquote{rcu\_read\_unlock}} çağrılarının yerleştirilmesi gerekmektedir.

\sphinxAtStartPar
Çekirdek tıpkı listelerde olduğu gibi pek çok yerde (ancak her yerde değil) hash tablolarıyla
RCU mekanizması eşliğinde işlem yapmaktadır. Çekirdeğin RCU mekanizması eşliğinde işlem
yaptığı yerlerde sizin de bu RCU’lu fonksiyonları kullanmanız gerekir.

\sphinxstepscope


\chapter{Proses Kontrol Bloğu ve task\_struct Yapısı}
\label{\detokenize{processcontrolblock:proses-kontrol-blogu-ve-task-struct-yapisi}}\label{\detokenize{processcontrolblock::doc}}
\sphinxAtStartPar
Bu bölümde proses kavramı ve proses kontrol bloğu üzerinde duracağız ve Linux çekirdeğindeki proses
kontrol bloğunu temsil eden task\_struct yapısını ele alacağız.


\section{Proses Kavramı}
\label{\detokenize{processcontrolblock:proses-kavrami}}
\sphinxAtStartPar
İşletim sistemlerinde \sphinxstyleemphasis{program} terimi çoğu kez “çalıştırılabilen bir dosyayı” ya da “bir kaynak dosyayı”
belirtmektedir. Çalışmakta olan programlara ise \sphinxstyleemphasis{proses} denilmektedir. Bir program çalıştırıldığında artık
o bir proses haline gelmektedir. Aynı programı birden fazla kez de çalıştırabiliriz. Bu durumda birbirinden
bağımsız birden fazla proses oluşacaktır.


\section{Proses Kontrol Bloğu}
\label{\detokenize{processcontrolblock:proses-kontrol-blogu}}
\sphinxAtStartPar
İşletim sistemlerinin proses yönetimlerindeki en önemli veri yapısı \sphinxstyleemphasis{proses kontrol bloğu (process control
block)} denilen veri yapısıdır. İşletim sistemleri her proses için kavramsal olarak ismine “proses kontrol
bloğu” denilen bir yapı türünden nesne oluşturmaktadır. Proseslerin yönetimi proses kontrol bloğu denilen
bu veri yapısına başvurularak yapılmaktadır. Proses kontrol bloğu terimi yerine bazı kaynaklar \sphinxstyleemphasis{proses
betimleyicisi (process descriptor)} terimini de kullanmaktadır.

\sphinxAtStartPar
Proses kontrol blokları içerisinde prosese ilişkin gerekli bütün bilgiler bulundurulmaktadır. Bu
bilgilerden bazıları şunlardır:
\begin{itemize}
\item {} 
\sphinxAtStartPar
Prosesin kullanıcı id’si, grup id’si gibi hesap (credential) bilgileri

\item {} 
\sphinxAtStartPar
Proses id’si

\item {} 
\sphinxAtStartPar
Prosesin üst prosesinin, alt proseslerinin, kardeş proseslerinin hangi prosesler olduğu bilgisi

\item {} 
\sphinxAtStartPar
Prosesin durumsal bilgisi

\item {} 
\sphinxAtStartPar
Prosesin bağlamsal geçişini (context switch) sağlama için gereksinim duyulan alanlar

\item {} 
\sphinxAtStartPar
Prosesin bellekte nereye yüklü olduğu gibi bellekle ilgili bilgileri

\item {} 
\sphinxAtStartPar
Prosesin çizelgeleyici ile ilgili olan bilgileri (örneğin çizelgeleme politikası, önceliği gibi)

\item {} 
\sphinxAtStartPar
Prosesin CPU kullanım istatistiği

\item {} 
\sphinxAtStartPar
Prosesin sinyal bilgileri

\item {} 
\sphinxAtStartPar
Prosesin proseslerarası haberleşme için gerekli olan bilgileri

\item {} 
\sphinxAtStartPar
Prosesin çalışma dizini (current working directory)

\item {} 
\sphinxAtStartPar
Prosesin açmış olduğu dosyalara ilişkin bilgiler

\end{itemize}


\section{\sphinxstyleliteralintitle{\sphinxupquote{task\_struct}} Yapısı}
\label{\detokenize{processcontrolblock:task-struct-yapisi}}
\sphinxAtStartPar
Linux çekirdeğinde proses kontrol bloğu \sphinxcode{\sphinxupquote{\textless{}include/linux/sched.h\textgreater{}}} dosyası içerisinde bulunan
\sphinxcode{\sphinxupquote{task\_struct}} isimli yapıyla temsil edilmiştir. Bu \sphinxcode{\sphinxupquote{task\_struct}} nesnesi eskiden daha az eleman
içeriyordu. Sonra çekirdek gittikçe geliştirilince dev bir yapı haline geldi. Örneğin bu \sphinxcode{\sphinxupquote{task\_struct}}
nesnesi Linus projeye ilk başlandığında (Linux 0.01) şu biçimdeydi:

\begin{sphinxVerbatim}[commandchars=\\\{\}]
\PYG{k}{struct}\PYG{+w}{ }\PYG{n+nc}{task\PYGZus{}struct}\PYG{+w}{ }\PYG{p}{\PYGZob{}}
\PYG{c+cm}{/* these are hardcoded \PYGZhy{} don\PYGZsq{}t touch */}
\PYG{+w}{    }\PYG{k+kt}{long}\PYG{+w}{ }\PYG{n}{state}\PYG{p}{;}\PYG{+w}{    }\PYG{c+cm}{/* \PYGZhy{}1 unrunnable, 0 runnable, \PYGZgt{}0 stopped */}
\PYG{+w}{    }\PYG{k+kt}{long}\PYG{+w}{ }\PYG{n}{counter}\PYG{p}{;}
\PYG{+w}{    }\PYG{k+kt}{long}\PYG{+w}{ }\PYG{n}{priority}\PYG{p}{;}
\PYG{+w}{    }\PYG{k+kt}{long}\PYG{+w}{ }\PYG{n}{signal}\PYG{p}{;}
\PYG{+w}{    }\PYG{n}{fn\PYGZus{}ptr}\PYG{+w}{ }\PYG{n}{sig\PYGZus{}restorer}\PYG{p}{;}
\PYG{+w}{    }\PYG{n}{fn\PYGZus{}ptr}\PYG{+w}{ }\PYG{n}{sig\PYGZus{}fn}\PYG{p}{[}\PYG{l+m+mi}{32}\PYG{p}{]}\PYG{p}{;}
\PYG{c+cm}{/* various fields */}
\PYG{+w}{    }\PYG{k+kt}{int}\PYG{+w}{ }\PYG{n}{exit\PYGZus{}code}\PYG{p}{;}
\PYG{+w}{    }\PYG{k+kt}{unsigned}\PYG{+w}{ }\PYG{k+kt}{long}\PYG{+w}{ }\PYG{n}{end\PYGZus{}code}\PYG{p}{,}\PYG{n}{end\PYGZus{}data}\PYG{p}{,}\PYG{n}{brk}\PYG{p}{,}\PYG{n}{start\PYGZus{}stack}\PYG{p}{;}
\PYG{+w}{    }\PYG{k+kt}{long}\PYG{+w}{ }\PYG{n}{pid}\PYG{p}{,}\PYG{n}{father}\PYG{p}{,}\PYG{n}{pgrp}\PYG{p}{,}\PYG{n}{session}\PYG{p}{,}\PYG{n}{leader}\PYG{p}{;}
\PYG{+w}{    }\PYG{k+kt}{unsigned}\PYG{+w}{ }\PYG{k+kt}{short}\PYG{+w}{ }\PYG{n}{uid}\PYG{p}{,}\PYG{n}{euid}\PYG{p}{,}\PYG{n}{suid}\PYG{p}{;}
\PYG{+w}{    }\PYG{k+kt}{unsigned}\PYG{+w}{ }\PYG{k+kt}{short}\PYG{+w}{ }\PYG{n}{gid}\PYG{p}{,}\PYG{n}{egid}\PYG{p}{,}\PYG{n}{sgid}\PYG{p}{;}
\PYG{+w}{    }\PYG{k+kt}{long}\PYG{+w}{ }\PYG{n}{alarm}\PYG{p}{;}
\PYG{+w}{    }\PYG{k+kt}{long}\PYG{+w}{ }\PYG{n}{utime}\PYG{p}{,}\PYG{n}{stime}\PYG{p}{,}\PYG{n}{cutime}\PYG{p}{,}\PYG{n}{cstime}\PYG{p}{,}\PYG{n}{start\PYGZus{}time}\PYG{p}{;}
\PYG{+w}{    }\PYG{k+kt}{unsigned}\PYG{+w}{ }\PYG{k+kt}{short}\PYG{+w}{ }\PYG{n}{used\PYGZus{}math}\PYG{p}{;}
\PYG{c+cm}{/* file system info */}
\PYG{+w}{    }\PYG{k+kt}{int}\PYG{+w}{ }\PYG{n}{tty}\PYG{p}{;}\PYG{+w}{       }\PYG{c+cm}{/* \PYGZhy{}1 if no tty, so it must be signed */}
\PYG{+w}{    }\PYG{k+kt}{unsigned}\PYG{+w}{ }\PYG{k+kt}{short}\PYG{+w}{ }\PYG{n}{umask}\PYG{p}{;}
\PYG{+w}{    }\PYG{k}{struct}\PYG{+w}{ }\PYG{n+nc}{m\PYGZus{}inode}\PYG{+w}{ }\PYG{o}{*}\PYG{+w}{ }\PYG{n}{pwd}\PYG{p}{;}
\PYG{+w}{    }\PYG{k}{struct}\PYG{+w}{ }\PYG{n+nc}{m\PYGZus{}inode}\PYG{+w}{ }\PYG{o}{*}\PYG{+w}{ }\PYG{n}{root}\PYG{p}{;}
\PYG{+w}{    }\PYG{k+kt}{unsigned}\PYG{+w}{ }\PYG{k+kt}{long}\PYG{+w}{ }\PYG{n}{close\PYGZus{}on\PYGZus{}exec}\PYG{p}{;}
\PYG{+w}{    }\PYG{k}{struct}\PYG{+w}{ }\PYG{n+nc}{file}\PYG{+w}{ }\PYG{o}{*}\PYG{+w}{ }\PYG{n}{filp}\PYG{p}{[}\PYG{n}{NR\PYGZus{}OPEN}\PYG{p}{]}\PYG{p}{;}
\PYG{c+cm}{/* ldt for this task 0 \PYGZhy{} zero 1 \PYGZhy{} cs 2 \PYGZhy{} ds\PYGZam{}ss */}
\PYG{+w}{    }\PYG{k}{struct}\PYG{+w}{ }\PYG{n+nc}{desc\PYGZus{}struct}\PYG{+w}{ }\PYG{n}{ldt}\PYG{p}{[}\PYG{l+m+mi}{3}\PYG{p}{]}\PYG{p}{;}
\PYG{c+cm}{/* tss for this task */}
\PYG{+w}{    }\PYG{k}{struct}\PYG{+w}{ }\PYG{n+nc}{tss\PYGZus{}struct}\PYG{+w}{ }\PYG{n}{tss}\PYG{p}{;}
\PYG{p}{\PYGZcb{}}\PYG{p}{;}
\end{sphinxVerbatim}

\sphinxAtStartPar
Çekirdeğin 2.4’lü versiyonunda bu yapı oldukça büyümüş, bellek, sinyal, dosya sistemi ve kimlik bilgileri
gibi pek çok alanı kapsayan kapsamlı bir yapı haline gelmiştir. Bugünkü 6’lı çekirdeklerde bu yapı birkaç
sayfa uzunluğundadır. Eskiden yapının her elemanı çekirdek derlemesine dahil ediliyordu. Ancak 2.6’lı
versiyonlardan başlanarak artık yapı elemanlarının bazıları konfigürasyona bağlı olarak yapıya dahil
edilmektedir. Bu nedenle gelişmiş çekirdeklerde yapıda aşağıdaki gibi \sphinxcode{\sphinxupquote{\#ifdef}} blokları görürseniz
şaşırmayınız:

\begin{sphinxVerbatim}[commandchars=\\\{\}]
\PYG{k}{struct}\PYG{+w}{ }\PYG{n+nc}{task\PYGZus{}struct}\PYG{+w}{ }\PYG{p}{\PYGZob{}}
\PYG{c+cp}{\PYGZsh{}}\PYG{c+cp}{ifdef CONFIG\PYGZus{}THREAD\PYGZus{}INFO\PYGZus{}IN\PYGZus{}TASK}
\PYG{+w}{    }\PYG{k}{struct}\PYG{+w}{ }\PYG{n+nc}{thread\PYGZus{}info}\PYG{+w}{    }\PYG{n}{thread\PYGZus{}info}\PYG{p}{;}
\PYG{c+cp}{\PYGZsh{}}\PYG{c+cp}{endif}
\PYG{+w}{    }\PYG{k+kt}{unsigned}\PYG{+w}{ }\PYG{k+kt}{int}\PYG{+w}{          }\PYG{n}{\PYGZus{}\PYGZus{}state}\PYG{p}{;}

\PYG{+w}{    }\PYG{c+cm}{/* saved state for \PYGZdq{}spinlock sleepers\PYGZdq{} */}
\PYG{+w}{    }\PYG{k+kt}{unsigned}\PYG{+w}{ }\PYG{k+kt}{int}\PYG{+w}{          }\PYG{n}{saved\PYGZus{}state}\PYG{p}{;}

\PYG{+w}{    }\PYG{n}{randomized\PYGZus{}struct\PYGZus{}fields\PYGZus{}start}

\PYG{+w}{    }\PYG{k+kt}{void}\PYG{+w}{                 }\PYG{o}{*}\PYG{n}{stack}\PYG{p}{;}
\PYG{+w}{    }\PYG{n}{refcount\PYGZus{}t}\PYG{+w}{            }\PYG{n}{usage}\PYG{p}{;}
\PYG{+w}{    }\PYG{k+kt}{unsigned}\PYG{+w}{ }\PYG{k+kt}{int}\PYG{+w}{          }\PYG{n}{flags}\PYG{p}{;}
\PYG{+w}{    }\PYG{p}{.}\PYG{p}{.}\PYG{p}{.}
\PYG{p}{\PYGZcb{}}\PYG{p}{;}
\end{sphinxVerbatim}

\sphinxAtStartPar
Bilindiği gibi Linux sistemlerinde yeni bir proses \sphinxstyleemphasis{fork} isimli POSIX fonksiyonuyla yaratılmaktadır.
\sphinxstyleemphasis{fork} POSIX fonksiyonu çekirdek içerisindeki \sphinxstyleemphasis{sys\_fork} isimli sistem fonksiyonunu çağırmaktadır. Yeni
çekirdeklerde bu fonksiyon da \sphinxstyleemphasis{kernel\_clone} isimli çekirdek fonksiyonunu çağırmaktadır. İşte
\sphinxcode{\sphinxupquote{task\_struct}} nesnesi bu fonksiyonlar tarafından çekirdeğin heap alanında yaratılmaktadır.


\section{Proseslerin ID Değerleri (PID)}
\label{\detokenize{processcontrolblock:proseslerin-id-degerleri-pid}}
\sphinxAtStartPar
Bilindiği gibi UNIX/Linux sistemlerinde her prosesin (yani çalışmakta olan her programın) o anda sistem
genelinde tek (unique) olan bir \sphinxstyleemphasis{proses id (process id)} değeri vardır. Proses id değeri hem çekirdek
tarafından hem de kullanıcı modundan çağrılabilen sistem fonksiyonları tarafından kullanılan bir
kavramdır. Proses id değeri bir prosesi temsil etmekte kullanılan tamsayısal bir değerdir. \sphinxstyleemphasis{ps} komutuyla
proseslere ilişkin proses id değerlerinin görüntülendiğini anımsayınız. Kullanıcı modunda \sphinxstyleemphasis{getpid} POSIX
fonksiyonu çalışmakta olan programın proses id değerini, \sphinxstyleemphasis{getppid} POSIX fonksiyonu ise üst prosesin
proses id değerini vermektedir. Tabii prosesin proses id değeri ve üst prosesin proses id değeri proses
kontrol bloğunda yani \sphinxcode{\sphinxupquote{task\_struct}} nesnesi içerisinde saklanmaktadır.


\section{Thread Kavramı}
\label{\detokenize{processcontrolblock:thread-kavrami}}
\sphinxAtStartPar
İşletim sistemlerine thread kavramı 90’lı yıllarda sokulmuştur. Ancak ilk thread denemeleri çeşitli
çalışmalar eşliğinde daha önceleri yapılmıştır. Linux işletim sistemine thread’ler ilk kez çekirdeğin
2.0 versiyonuyla eklenmiş olsa da daha sonra çeşitli düzeltmelerle thread yapısı biraz değiştirilmiştir.

\sphinxAtStartPar
Proses kavramı çalışmakta olan programın tüm bilgilerini temsil etmektedir. Halbuki thread yalnızca bir
akış belirtmektedir. Bir proses tek thread’le çalışmaya başlar. Prosesin çalışmaya başladığı thread’e
\sphinxstyleemphasis{ana thread (main thread)} de denilmektedir. Thread’ler de kullanıcı modundan sistem fonksiyonları
(\sphinxstyleemphasis{sys\_clone} sistem fonksiyonu) ile yaratılmaktadır.

\sphinxAtStartPar
UNIX türevi sistemlere thread’ler ilk kez sokulduğunda farklı varyantlar farklı thread kütüphaneleri
kullanıyordu. Sonra thread kavramı POSIX standartlarına da POSIX.1c ile 1995 yılında sokuldu. Böylece
POSIX taşınabilir thread fonksiyonları oluşturuldu. POSIX’in taşınabilir thread fonksiyonlarının
gerçekleştirimi için Linux ve diğer UNIX varyantları çekirdek içerisinde de değişiklikler yapmak zorunda
kalmıştır.

\sphinxAtStartPar
Thread’lerin yalnızca bir akış belirttiğini söylemiştik. Proses kavramı ise tüm thread’lerle birlikte
çalışmakta olan programın tüm bilgilerini temsil etmektedir. Pek çok bilgi thread’e özgü değildir,
prosese özgüdür. Örneğin bir thread’in “çalışma dizini (current working directory)” diye bir kavram
yoktur. Prosesin çalışma dizini diye bir kavram vardır. Örneğin bir thread’in kullanıcı id’si, grup id’si
diye bir kavram yoktur. Prosesin kullanıcı id’si ve grup id’si diye bir kavram vardır. Yani çalışmakta
olan programa ilişkin pek çok bilgi programın tüm thread’leri için de geçerlidir.

\sphinxAtStartPar
Aslında thread’ler aynı bellek alanını ve ortak bilgileri kullanan prosesler gibi de düşünülebilir. Zaten
Linux işletim sisteminde thread yaratmakla proses yaratmak aslında aynı çekirdek fonksiyonuyla
yapılmaktadır. Yani bir prosesin thread’lerini biz aynı bellek alanını kullanan prosesler gibi de
düşünebiliriz.

\sphinxAtStartPar
Linux işletim sisteminde thread’ler de adeta aynı bellek alanı üzerinde çalışan birer proses gibi
düşünülmüştür. Bu nedenle yalnızca prosesler için değil thread’ler için de \sphinxcode{\sphinxupquote{task\_struct}} nesneleri
oluşturulmaktadır. Tabii işletim sistemi bir prosesin thread’lerini izleyen paragraflarda göreceğimiz
biçimde bağlı listelerde tutmaktadır.


\section{Prosesler Arasındaki Altlık\sphinxhyphen{}Üstlük İlişkisi}
\label{\detokenize{processcontrolblock:prosesler-arasindaki-altlik-ustluk-iliskisi}}
\sphinxAtStartPar
UNIX/Linux sistemlerinde prosesler arasında kuvvetli bir \sphinxstyleemphasis{altlık\sphinxhyphen{}üstlük (parent\sphinxhyphen{}child)} ilişkisi vardır.
Bir proses yaratıldığında prosesin \sphinxcode{\sphinxupquote{task\_struct}} bilgilerinin çoğu üst prosesten (parent process)
alınmaktadır. Örneğin bir program başka bir programı \sphinxstyleemphasis{fork}/\sphinxstyleemphasis{exec} ile çalıştırdığında çalıştırılan
program çalıştıran programın pek çok özelliğini almaktadır. Örneğin \sphinxstyleemphasis{fork} işlemi sırasında \sphinxstyleemphasis{fork}
yapan programın kullanıcı id’si ve çalışma dizini neyse yeni yaratılan alt prosesin (child process)
kullanıcı id’si ve çalışma dizini de aynı olur. Daha teknik olarak ifade edersek \sphinxstyleemphasis{fork} işlemi sırasında
alt proses için yaratılan \sphinxcode{\sphinxupquote{task\_struct}} nesnesinin pek çok elemanı üst prosesin \sphinxcode{\sphinxupquote{task\_struct}}
nesnesinden kopyalanmaktadır.

\sphinxAtStartPar
POSIX thread sisteminde thread’ler arasında önemli bir altlık\sphinxhyphen{}üstlük ilişkisi yoktur. Yani birkaç özel
durum dışında bir thread’i hangi thread’in yarattığının önemi yoktur.


\section{\sphinxstyleliteralintitle{\sphinxupquote{task\_struct}} Yapısının Dallı Budaklı Tasarımı}
\label{\detokenize{processcontrolblock:task-struct-yapisinin-dalli-budakli-tasarimi}}
\sphinxAtStartPar
\sphinxcode{\sphinxupquote{task\_struct}} yapısının pek çok gösterici elemanı vardır. Bu gösterici elemanları başka yapıları
göstermektedir. Hatta o gösterici elemanların gösterdiği yapıların elemanları da başka yapıları
göstermektedir. O halde \sphinxcode{\sphinxupquote{task\_struct}} aslında dallı budaklı bir yapıdır. Biz bir bilginin proses
kontrol bloğunda olduğunu söylediğimiz zaman o bilginin hemen \sphinxcode{\sphinxupquote{task\_struct}} içerisindeki bir elemanda
olduğunu kastetmeyeceğiz. O bilgi \sphinxcode{\sphinxupquote{task\_struct}} içerisindeki bir elemanın gösterdiği başka bir yapının
içerisinde de olabilir. Önemli olan o bilgiye \sphinxcode{\sphinxupquote{task\_struct}} yapısından hareketle erişilebilmesidir.

\sphinxAtStartPar
Bir \sphinxcode{\sphinxupquote{task\_struct}} yapı nesnesinin ana elemanları diğer bir \sphinxcode{\sphinxupquote{task\_struct}} nesnesine kopyalanırsa
aslında asıl yapı nesnesi ile kopyalanmış olan yapı nesnesinin gösterici elemanları aynı yerleri gösterir
duruma gelir. Programlamada bu tür kopyalamaya \sphinxstyleemphasis{sığ kopyalama (shallow copy)} denildiğini anımsayınız.
Aslında \sphinxcode{\sphinxupquote{task\_struct}} yapısının bu biçimdeki dallı budaklı tasarımı üst proses\sphinxhyphen{}alt proses ilişkisinde
alt prosesin üst prosesten birtakım özellikleri alması (inherit etmesi) sürecinde faydalar sağlamaktadır.
Böylece alt prosesin \sphinxcode{\sphinxupquote{task\_struct}} nesnesine daha az eleman kopyalanmaktadır. Örneğin güncel
çekirdeklerdeki \sphinxcode{\sphinxupquote{task\_struct}} yapısının aşağıdaki kısmına dikkat ediniz:

\begin{sphinxVerbatim}[commandchars=\\\{\}]
\PYG{k}{struct}\PYG{+w}{ }\PYG{n+nc}{task\PYGZus{}struct}\PYG{+w}{ }\PYG{p}{\PYGZob{}}
\PYG{+w}{    }\PYG{c+cm}{/* ... */}

\PYG{+w}{    }\PYG{c+cm}{/* Filesystem information: */}
\PYG{+w}{    }\PYG{k}{struct}\PYG{+w}{ }\PYG{n+nc}{fs\PYGZus{}struct}\PYG{+w}{        }\PYG{o}{*}\PYG{n}{fs}\PYG{p}{;}

\PYG{+w}{    }\PYG{c+cm}{/* Open file information: */}
\PYG{+w}{    }\PYG{k}{struct}\PYG{+w}{ }\PYG{n+nc}{files\PYGZus{}struct}\PYG{+w}{     }\PYG{o}{*}\PYG{n}{files}\PYG{p}{;}

\PYG{+w}{    }\PYG{c+cm}{/* ... */}
\PYG{p}{\PYGZcb{}}\PYG{p}{;}
\end{sphinxVerbatim}

\sphinxAtStartPar
Burada aslında \sphinxcode{\sphinxupquote{fs}} göstericisi aşağıdaki gibi bir yapıyı göstermektedir:

\begin{sphinxVerbatim}[commandchars=\\\{\}]
\PYG{k}{struct}\PYG{+w}{ }\PYG{n+nc}{fs\PYGZus{}struct}\PYG{+w}{ }\PYG{p}{\PYGZob{}}
\PYG{+w}{    }\PYG{k+kt}{int}\PYG{+w}{ }\PYG{n}{users}\PYG{p}{;}
\PYG{+w}{    }\PYG{n}{spinlock\PYGZus{}t}\PYG{+w}{ }\PYG{n}{lock}\PYG{p}{;}
\PYG{+w}{    }\PYG{n}{seqcount\PYGZus{}spinlock\PYGZus{}t}\PYG{+w}{ }\PYG{n}{seq}\PYG{p}{;}
\PYG{+w}{    }\PYG{k+kt}{int}\PYG{+w}{ }\PYG{n}{umask}\PYG{p}{;}
\PYG{+w}{    }\PYG{k+kt}{int}\PYG{+w}{ }\PYG{n}{in\PYGZus{}exec}\PYG{p}{;}
\PYG{+w}{    }\PYG{k}{struct}\PYG{+w}{ }\PYG{n+nc}{path}\PYG{+w}{ }\PYG{n}{root}\PYG{p}{,}\PYG{+w}{ }\PYG{n}{pwd}\PYG{p}{;}
\PYG{p}{\PYGZcb{}}\PYG{+w}{ }\PYG{n}{\PYGZus{}\PYGZus{}randomize\PYGZus{}layout}\PYG{p}{;}
\end{sphinxVerbatim}

\sphinxAtStartPar
Prosesin kök dizininin ve çalışma dizinin burada tutulduğuna dikkat ediniz. \sphinxcode{\sphinxupquote{files}} göstericisi ise
aşağıdaki gibi bir yapıyı göstermektedir:

\begin{sphinxVerbatim}[commandchars=\\\{\}]
\PYG{k}{struct}\PYG{+w}{ }\PYG{n+nc}{files\PYGZus{}struct}\PYG{+w}{ }\PYG{p}{\PYGZob{}}
\PYG{+w}{    }\PYG{n}{atomic\PYGZus{}t}\PYG{+w}{ }\PYG{n}{count}\PYG{p}{;}
\PYG{+w}{    }\PYG{k+kt}{bool}\PYG{+w}{ }\PYG{n}{resize\PYGZus{}in\PYGZus{}progress}\PYG{p}{;}
\PYG{+w}{    }\PYG{n}{wait\PYGZus{}queue\PYGZus{}head\PYGZus{}t}\PYG{+w}{ }\PYG{n}{resize\PYGZus{}wait}\PYG{p}{;}

\PYG{+w}{    }\PYG{k}{struct}\PYG{+w}{ }\PYG{n+nc}{fdtable}\PYG{+w}{ }\PYG{n}{\PYGZus{}\PYGZus{}rcu}\PYG{+w}{ }\PYG{o}{*}\PYG{n}{fdt}\PYG{p}{;}
\PYG{+w}{    }\PYG{k}{struct}\PYG{+w}{ }\PYG{n+nc}{fdtable}\PYG{+w}{ }\PYG{n}{fdtab}\PYG{p}{;}

\PYG{+w}{    }\PYG{n}{spinlock\PYGZus{}t}\PYG{+w}{ }\PYG{n}{file\PYGZus{}lock}\PYG{+w}{ }\PYG{n}{\PYGZus{}\PYGZus{}\PYGZus{}\PYGZus{}cacheline\PYGZus{}aligned\PYGZus{}in\PYGZus{}smp}\PYG{p}{;}
\PYG{+w}{    }\PYG{k+kt}{unsigned}\PYG{+w}{ }\PYG{k+kt}{int}\PYG{+w}{ }\PYG{n}{next\PYGZus{}fd}\PYG{p}{;}
\PYG{+w}{    }\PYG{k+kt}{unsigned}\PYG{+w}{ }\PYG{k+kt}{long}\PYG{+w}{ }\PYG{n}{close\PYGZus{}on\PYGZus{}exec\PYGZus{}init}\PYG{p}{[}\PYG{l+m+mi}{1}\PYG{p}{]}\PYG{p}{;}
\PYG{+w}{    }\PYG{k+kt}{unsigned}\PYG{+w}{ }\PYG{k+kt}{long}\PYG{+w}{ }\PYG{n}{open\PYGZus{}fds\PYGZus{}init}\PYG{p}{[}\PYG{l+m+mi}{1}\PYG{p}{]}\PYG{p}{;}
\PYG{+w}{    }\PYG{k+kt}{unsigned}\PYG{+w}{ }\PYG{k+kt}{long}\PYG{+w}{ }\PYG{n}{full\PYGZus{}fds\PYGZus{}bits\PYGZus{}init}\PYG{p}{[}\PYG{l+m+mi}{1}\PYG{p}{]}\PYG{p}{;}
\PYG{+w}{    }\PYG{k}{struct}\PYG{+w}{ }\PYG{n+nc}{file}\PYG{+w}{ }\PYG{n}{\PYGZus{}\PYGZus{}rcu}\PYG{+w}{ }\PYG{o}{*}\PYG{+w}{ }\PYG{n}{fd\PYGZus{}array}\PYG{p}{[}\PYG{n}{NR\PYGZus{}OPEN\PYGZus{}DEFAULT}\PYG{p}{]}\PYG{p}{;}
\PYG{p}{\PYGZcb{}}\PYG{p}{;}
\end{sphinxVerbatim}

\sphinxAtStartPar
\sphinxstyleemphasis{fork} işlemi sırasında sığ kopyalama yapıldığından dolayı \sphinxstyleemphasis{fork} işleminden sonra üst ve alt prosesin
\sphinxcode{\sphinxupquote{task\_struct}} nesnelerinin \sphinxcode{\sphinxupquote{fs}} ve \sphinxcode{\sphinxupquote{files}} göstericileri aynı \sphinxcode{\sphinxupquote{fs\_struct}} ve \sphinxcode{\sphinxupquote{files\_struct}}
nesnelerini gösteriyor olacaktır.


\section{\sphinxstyleliteralintitle{\sphinxupquote{current}} Makrosu}
\label{\detokenize{processcontrolblock:current-makrosu}}
\sphinxAtStartPar
Thread’lerin de tıpkı prosesler gibi \sphinxcode{\sphinxupquote{task\_struct}} nesnelerine sahip olduğunu belirtmiştik. Peki
işletim sisteminin çizelgeleyici (scheduler) alt sistemi neleri çizelgelemektedir? İşte çizelgeleyici
alt sistem aslında yalnızca thread’leri yani thread’lere ilişkin \sphinxcode{\sphinxupquote{task\_struct}} nesnelerini
çizelgelemektedir. Yani çizelgeleyici alt sistemin \sphinxstyleemphasis{çalışma kuyruğunda (run queue)} yalnızca
\sphinxcode{\sphinxupquote{task\_struct}} nesnelerinin olduğunu varsayabilirsiniz. Bir proses yaratıldığında zaten o proses için
yaratılan \sphinxcode{\sphinxupquote{task\_struct}} nesnesi aynı zamanda prosesin ana thread’inin \sphinxcode{\sphinxupquote{task\_struct}} nesnesi gibidir.
Başka bir deyişle aslında bütün \sphinxcode{\sphinxupquote{task\_struct}} nesneleri birer thread belirtmektedir.

\sphinxAtStartPar
Eskiden Linux sistemleri yalnızca tek işlemciyi destekliyordu. Sonra 2.0 versiyonuyla birlikte çekirdek
zamanla birden fazla işlemciyi ya da çekirdeği (core) destekler hale geldi. İşte bir thread çalışırken
\sphinxcode{\sphinxupquote{task\_struct *}} türünden \sphinxstyleemphasis{current} isimli global değişken o anda çalışmakta olan thread’e ilişkin
\sphinxcode{\sphinxupquote{task\_struct}} nesnesini gösterecek biçimde ayarlanmaktadır. Yani bir çekirdek kodunda \sphinxstyleemphasis{current}
göstericisini görürseniz bu \sphinxstyleemphasis{current} göstericisi o anda çalışmakta olan thread’e ilişkin \sphinxcode{\sphinxupquote{task\_struct}}
nesnesini gösteriyor durumda olacaktır. Tabii bu \sphinxstyleemphasis{current} göstericisinin o anda çalışan thread’e ilişkin
\sphinxcode{\sphinxupquote{task\_struct}} nesnesini göstermesini bağlamsal geçişi (context switch) gerçekleştiren çizelgeleyici alt
sistem sağlamaktadır.

\sphinxAtStartPar
Bir süredir birden fazla işlemciyi ya da çekirdeği destekleyen Linux çekirdeklerinde artık \sphinxstyleemphasis{current}
değişkeni bir gösterici değil fonksiyon belirtmektedir. Mevcut çekirdeklerde \sphinxstyleemphasis{current} aşağıdaki gibi bir
makro biçiminde tanımlanmıştır:

\begin{sphinxVerbatim}[commandchars=\\\{\}]
\PYG{c+cp}{\PYGZsh{}}\PYG{c+cp}{define current get\PYGZus{}current()}
\end{sphinxVerbatim}

\sphinxAtStartPar
Görüldüğü gibi artık biz \sphinxstyleemphasis{current} değişkenini kullandığımızda aslında \sphinxstyleemphasis{get\_current()} fonksiyonunu
çağırıp bu fonksiyonun geri dönüş değerini kullanmış olmaktayız. Bu konunun ayrıntıları thread’in
\sphinxstyleemphasis{çekirdek stack alanı (kernel stack)} ve \sphinxstyleemphasis{işlemciye özgü global alanlar} konusuyla ilgilidir ve başka
bir bölümde ele alınacaktır.

\sphinxAtStartPar
Biz kursumuzda \sphinxstyleemphasis{current} için “current göstericisi” ya da “current makrosu” terimlerini kullanacağız.


\section{\sphinxstyleliteralintitle{\sphinxupquote{fork}} ve \sphinxstyleliteralintitle{\sphinxupquote{pthread\_create}} Çağrı Zincirleri}
\label{\detokenize{processcontrolblock:fork-ve-pthread-create-cagri-zincirleri}}
\sphinxAtStartPar
Bilindiği gibi UNIX/Linux sistemlerinde kullanıcı modunda yeni bir proses \sphinxstyleemphasis{fork} POSIX fonksiyonu ile,
yeni bir thread de \sphinxstyleemphasis{pthread\_create} fonksiyonu ile yaratılmaktadır. Yukarıda da belirttiğimiz gibi
aslında bu iki çağrı da belirli bir noktadan sonra çekirdek içerisindeki aynı fonksiyonları
çağırmaktadır:

\sphinxincludegraphics[]{graphviz-83c3f3b28900cf53342612b50317355030cacee4.pdf}

\sphinxAtStartPar
O halde aslında çekirdek gözüyle bakıldığında bir thread başka bir thread tarafından yaratılmaktadır.
Yani bu yaratımda iki thread söz konusudur: Yaratan thread ve yaratılan thread. Yaratan thread’e ilişkin
zaten \sphinxcode{\sphinxupquote{task\_struct}} nesnesi mevcuttur. O halde yaratılan thread için de bir \sphinxcode{\sphinxupquote{task\_struct}} nesnesi
oluşturulup bir biçimde bu yapılar ilişkilendirilecektir.


\section{\sphinxstyleliteralintitle{\sphinxupquote{task\_struct}} Nesneleri Arasındaki İlişkiler}
\label{\detokenize{processcontrolblock:task-struct-nesneleri-arasindaki-iliskiler}}
\sphinxAtStartPar
Şimdi de \sphinxcode{\sphinxupquote{task\_struct}} nesneleri arasındaki ilişkilerin bağlı listelerle nasıl oluşturulduğu üzerinde
duracağız. Çekirdeğin bir prosesin thread’lerine sonra da alt proseslerine nasıl eriştiğini açıklayacağız.


\subsection{Thread Listesi: \sphinxstyleliteralintitle{\sphinxupquote{thread\_head}} / \sphinxstyleliteralintitle{\sphinxupquote{thread\_node}}}
\label{\detokenize{processcontrolblock:thread-listesi-thread-head-thread-node}}
\sphinxAtStartPar
Bir prosesin thread’lerine erişim için eskiden \sphinxcode{\sphinxupquote{task\_struct}} içerisindeki \sphinxcode{\sphinxupquote{thread\_group}} isimli
döngüsel bağlı liste bağı kullanılıyordu. 4.2 çekirdeği (2015) ile birlikte bu veri yapısında bir
değişiklik yapılmıştır. Yeni sistemde prosesin thread’lerine ilişkin bağlı listenin kök düğümü prosesin
sinyal bilgilerinin bulunduğu yerde saklanmaktadır. Prosesin sinyal bilgileri \sphinxcode{\sphinxupquote{task\_struct}} içerisindeki
\sphinxcode{\sphinxupquote{signal}} göstericisinin gösterdiği yerdeki \sphinxcode{\sphinxupquote{signal\_struct}} yapısı içerisinde tutulmaktadır. İşte
\sphinxcode{\sphinxupquote{signal\_struct}} yapısının \sphinxcode{\sphinxupquote{thread\_head}} elemanı da prosesin thread’lerine ilişkin bu kök düğümü
belirtmektedir. Prosesin thread’lerinin bağlı listesi için ise \sphinxcode{\sphinxupquote{task\_struct}} içerisindeki \sphinxcode{\sphinxupquote{thread\_node}}
elemanı kullanılmaktadır. Thread listesine ilişkin ilgili elemanları şöyle betimleyebiliriz:

\begin{sphinxVerbatim}[commandchars=\\\{\}]
\PYG{k}{struct}\PYG{+w}{ }\PYG{n+nc}{signal\PYGZus{}struct}\PYG{+w}{ }\PYG{p}{\PYGZob{}}
\PYG{+w}{    }\PYG{c+cm}{/* ... */}
\PYG{+w}{    }\PYG{k}{struct}\PYG{+w}{ }\PYG{n+nc}{list\PYGZus{}head}\PYG{+w}{ }\PYG{n}{thread\PYGZus{}head}\PYG{p}{;}
\PYG{+w}{    }\PYG{c+cm}{/* ... */}
\PYG{p}{\PYGZcb{}}\PYG{p}{;}

\PYG{k}{struct}\PYG{+w}{ }\PYG{n+nc}{task\PYGZus{}struct}\PYG{+w}{ }\PYG{p}{\PYGZob{}}
\PYG{+w}{    }\PYG{c+cm}{/* ... */}
\PYG{+w}{    }\PYG{k}{struct}\PYG{+w}{ }\PYG{n+nc}{signal\PYGZus{}struct}\PYG{+w}{ }\PYG{o}{*}\PYG{n}{signal}\PYG{p}{;}
\PYG{+w}{    }\PYG{k}{struct}\PYG{+w}{ }\PYG{n+nc}{list\PYGZus{}head}\PYG{+w}{ }\PYG{n}{thread\PYGZus{}node}\PYG{p}{;}
\PYG{+w}{    }\PYG{c+cm}{/* ... */}
\PYG{p}{\PYGZcb{}}\PYG{p}{;}
\end{sphinxVerbatim}

\sphinxAtStartPar
Aşağıdaki diyagram bu bağlı listenin yapısını göstermektedir. Mavi renkteki \sphinxcode{\sphinxupquote{thread\_head}} kök düğümü
\sphinxcode{\sphinxupquote{signal\_struct}} içerisindedir; turuncu renkteki \sphinxcode{\sphinxupquote{thread\_node}} düğümleri ise her bir thread’e ait
\sphinxcode{\sphinxupquote{task\_struct}} içerisindedir:

\sphinxincludegraphics[]{graphviz-669290f226b7eeefaa77ce92d0bc4b981f46628a.pdf}

\sphinxAtStartPar
\sphinxcode{\sphinxupquote{thread\_node}} bağlı listesi çekirdeğe ilk eklendiğinde bir süre eski \sphinxcode{\sphinxupquote{thread\_group}} listesi de
muhafaza edilmişti. Sonra tamamen eski \sphinxcode{\sphinxupquote{thread\_group}} listesi \sphinxcode{\sphinxupquote{task\_struct}} içerisinden kaldırıldı.
Artık yeni çekirdeklerde prosesin thread’lerinin \sphinxcode{\sphinxupquote{task\_struct}} nesnelerini elde etmek için prosesin
\sphinxcode{\sphinxupquote{signal\_struct}} nesnesindeki \sphinxcode{\sphinxupquote{thread\_head}} kök düğümünden başlanarak bağlı listenin \sphinxcode{\sphinxupquote{thread\_node}}
düğümlerinin dolaşılması gerekmektedir.

\sphinxAtStartPar
Aslında Linux çekirdeklerinde prosesin tüm thread’lerine ilişkin \sphinxcode{\sphinxupquote{task\_struct}} içerisinde bulunan
\sphinxcode{\sphinxupquote{signal}} göstericisi aynı \sphinxcode{\sphinxupquote{signal\_struct}} nesnesini göstermektedir. Yani \sphinxcode{\sphinxupquote{thread\_head}} kök düğümüne
biz prosesin ana thread’inden erişmek zorunda değiliz. Prosesin herhangi bir thread’ine ilişkin
\sphinxcode{\sphinxupquote{task\_struct}} nesnesindeki \sphinxcode{\sphinxupquote{signal}} göstericisi yoluyla bu kök düğüme erişebiliriz.

\sphinxAtStartPar
Çekirdek içerisinde bir \sphinxcode{\sphinxupquote{task\_struct}} yapısının prosesin ana thread’ine ilişkin olup olmadığını
belirleyen \sphinxcode{\sphinxupquote{thread\_group\_leader}} isimli bir makro da vardır:

\begin{sphinxVerbatim}[commandchars=\\\{\}]
\PYG{c+cp}{\PYGZsh{}}\PYG{c+cp}{include}\PYG{+w}{ }\PYG{c+cpf}{\PYGZlt{}linux/sched/signal.h\PYGZgt{}}

\PYG{k}{static}\PYG{+w}{ }\PYG{k+kr}{inline}\PYG{+w}{ }\PYG{k+kt}{bool}\PYG{+w}{ }\PYG{n+nf}{thread\PYGZus{}group\PYGZus{}leader}\PYG{p}{(}\PYG{k}{struct}\PYG{+w}{ }\PYG{n+nc}{task\PYGZus{}struct}\PYG{+w}{ }\PYG{o}{*}\PYG{n}{p}\PYG{p}{)}
\PYG{p}{\PYGZob{}}
\PYG{+w}{    }\PYG{k}{return}\PYG{+w}{ }\PYG{n}{p}\PYG{o}{\PYGZhy{}}\PYG{o}{\PYGZgt{}}\PYG{n}{exit\PYGZus{}signal}\PYG{+w}{ }\PYG{o}{\PYGZgt{}}\PYG{o}{=}\PYG{+w}{ }\PYG{l+m+mi}{0}\PYG{p}{;}
\PYG{p}{\PYGZcb{}}
\end{sphinxVerbatim}

\sphinxAtStartPar
Tabii aslında bir \sphinxcode{\sphinxupquote{task\_struct}} nesnesinin ana thread’e ilişkin olup olmadığı \sphinxcode{\sphinxupquote{p == p\sphinxhyphen{}\textgreater{}group\_leader}}
ya da \sphinxcode{\sphinxupquote{p\sphinxhyphen{}\textgreater{}pid == p\sphinxhyphen{}\textgreater{}tgid}} işlemiyle anlaşılabilir. Ancak özel bir bilgi olarak \sphinxcode{\sphinxupquote{task\_struct}}
içerisindeki \sphinxcode{\sphinxupquote{exit\_signal}} elemanı da bu bilgiyi verebilmektedir. Bu eleman eğer thread ana thread
değilse \sphinxhyphen{}1 değerinde, ana thread ise \textgreater{}= 0 değerinde olmaktadır.

\sphinxAtStartPar
Linux çekirdeğinde terminoloji bağlamında “prosesin ana thread’i” yerine \sphinxstyleemphasis{thread grup lideri (thread group
leader)} terimi tercih edilmektedir.


\subsection{\sphinxstyleliteralintitle{\sphinxupquote{pid}} ve \sphinxstyleliteralintitle{\sphinxupquote{tgid}} Elemanları}
\label{\detokenize{processcontrolblock:pid-ve-tgid-elemanlari}}
\sphinxAtStartPar
Bilindiği gibi UNIX/Linux sistemlerinde her prosesin sistem genelinde tek (unique) olan bir \sphinxstyleemphasis{proses id
(pid)} değeri vardır. Linux çekirdeği prosesin id değerini prosesin ana thread’ine ilişkin \sphinxcode{\sphinxupquote{task\_struct}}
içerisindeki \sphinxcode{\sphinxupquote{pid}} elemanında saklamaktadır:

\begin{sphinxVerbatim}[commandchars=\\\{\}]
\PYG{k}{struct}\PYG{+w}{ }\PYG{n+nc}{task\PYGZus{}struct}\PYG{+w}{ }\PYG{p}{\PYGZob{}}
\PYG{+w}{    }\PYG{c+cm}{/* ... */}
\PYG{+w}{    }\PYG{k+kt}{pid\PYGZus{}t}\PYG{+w}{ }\PYG{n}{pid}\PYG{p}{;}
\PYG{+w}{    }\PYG{c+cm}{/* ... */}
\PYG{p}{\PYGZcb{}}\PYG{p}{;}
\end{sphinxVerbatim}

\sphinxAtStartPar
Linux’ta her thread’in \sphinxcode{\sphinxupquote{task\_struct}} nesnesinde ayrı bir \sphinxcode{\sphinxupquote{pid}} değeri vardır, fakat prosesin pid
değeri denildiğinde prosesin ana thread’inin pid değeri anlaşılmaktadır. Ana thread’in pid değeri aynı
zamanda \sphinxcode{\sphinxupquote{task\_struct}} içerisindeki \sphinxcode{\sphinxupquote{tgid}} isimli bir eleman da tutulmaktadır. \sphinxcode{\sphinxupquote{task\_struct}} nesnesi
hangi thread’e ilişkin olursa olsun \sphinxcode{\sphinxupquote{tgid}} elemanı her zaman prosesin ana thread’inin pid değerini
tutmaktadır:

\begin{sphinxVerbatim}[commandchars=\\\{\}]
\PYG{k}{struct}\PYG{+w}{ }\PYG{n+nc}{task\PYGZus{}struct}\PYG{+w}{ }\PYG{p}{\PYGZob{}}
\PYG{+w}{    }\PYG{c+cm}{/* ... */}
\PYG{+w}{    }\PYG{k+kt}{pid\PYGZus{}t}\PYG{+w}{   }\PYG{n}{pid}\PYG{p}{;}\PYG{+w}{    }\PYG{c+cm}{/* o thread\PYGZsq{}e ilişkin pid değeri, POSIX\PYGZsq{}te böyle bir kavram yok */}
\PYG{+w}{    }\PYG{k+kt}{pid\PYGZus{}t}\PYG{+w}{   }\PYG{n}{tgid}\PYG{p}{;}\PYG{+w}{   }\PYG{c+cm}{/* ana thread\PYGZsq{}e ilişkin pid değeri, getpid bu değeri veriyor */}
\PYG{+w}{    }\PYG{c+cm}{/* ... */}
\PYG{p}{\PYGZcb{}}\PYG{p}{;}
\end{sphinxVerbatim}

\sphinxAtStartPar
O halde bazı ayrıntıları da göz ardı edersek \sphinxstyleemphasis{getpid} POSIX fonksiyonunun çağırdığı \sphinxstyleemphasis{sys\_getpid} sistem
fonksiyonu prosesin proses id değerini doğrudan \sphinxstyleemphasis{current} göstericisinin gösterdiği \sphinxcode{\sphinxupquote{task\_struct}}
nesnesinin içerisindeki \sphinxcode{\sphinxupquote{tgid}} değerinden alarak vermektedir:

\sphinxincludegraphics[]{graphviz-3574f45af61fd54730ed1b282e3650f08181ffa8.pdf}

\sphinxAtStartPar
Şimdi aklınıza “neden her thread’te ayrıca ana thread’in pid değeri tgid ismiyle tutuluyor?” sorusu
gelebilir. Bunun iki nedeni vardır: Birincisi prosesin ana thread’i sonlanabilir, bu durumda ana thread’e
ilişkin \sphinxcode{\sphinxupquote{task\_struct}} nesnesine erişilemeyebilir. İkincisi ise bu yolla prosesin id değerine daha hızlı
erişimin sağlanabilmesidir.

\sphinxAtStartPar
Linux’ta thread’lere özgü pid değeri Linux’a özgü \sphinxstyleemphasis{gettid} fonksiyonu ile elde edilebilmektedir:

\begin{sphinxVerbatim}[commandchars=\\\{\}]
\PYG{c+cp}{\PYGZsh{}}\PYG{c+cp}{define \PYGZus{}GNU\PYGZus{}SOURCE}
\PYG{c+cp}{\PYGZsh{}}\PYG{c+cp}{include}\PYG{+w}{ }\PYG{c+cpf}{\PYGZlt{}unistd.h\PYGZgt{}}

\PYG{k+kt}{pid\PYGZus{}t}\PYG{+w}{ }\PYG{n+nf}{gettid}\PYG{p}{(}\PYG{k+kt}{void}\PYG{p}{)}\PYG{p}{;}
\end{sphinxVerbatim}

\begin{sphinxadmonition}{note}{Not:}
\sphinxAtStartPar
\sphinxstyleemphasis{gettid} fonksiyonunun bir POSIX fonksiyonu olmadığına dikkat ediniz.
\end{sphinxadmonition}


\subsection{\sphinxstyleliteralintitle{\sphinxupquote{group\_leader}} Göstericisi}
\label{\detokenize{processcontrolblock:group-leader-gostericisi}}
\sphinxAtStartPar
Aslında Linux çekirdeği thread’lere ilişkin \sphinxcode{\sphinxupquote{task\_struct}} nesnelerinde yalnızca prosesin ana thread’ine
ilişkin pid değerini değil, aynı zamanda ana thread’in \sphinxcode{\sphinxupquote{task\_struct}} nesnesinin adresini de
tutmaktadır. Ana thread’in \sphinxcode{\sphinxupquote{task\_struct}} nesnesinin adresi \sphinxcode{\sphinxupquote{task\_struct}} yapısının \sphinxcode{\sphinxupquote{group\_leader}}
elemanında tutulmaktadır:

\begin{sphinxVerbatim}[commandchars=\\\{\}]
\PYG{k}{struct}\PYG{+w}{ }\PYG{n+nc}{task\PYGZus{}struct}\PYG{+w}{ }\PYG{p}{\PYGZob{}}
\PYG{+w}{    }\PYG{c+cm}{/* ... */}
\PYG{+w}{    }\PYG{k+kt}{pid\PYGZus{}t}\PYG{+w}{   }\PYG{n}{pid}\PYG{p}{;}\PYG{+w}{                        }\PYG{c+cm}{/* o thread\PYGZsq{}e ilişkin pid değeri */}
\PYG{+w}{    }\PYG{k+kt}{pid\PYGZus{}t}\PYG{+w}{   }\PYG{n}{tgid}\PYG{p}{;}\PYG{+w}{                       }\PYG{c+cm}{/* ana thread\PYGZsq{}e ilişkin pid değeri */}
\PYG{+w}{    }\PYG{k}{struct}\PYG{+w}{ }\PYG{n+nc}{task\PYGZus{}struct}\PYG{+w}{ }\PYG{o}{*}\PYG{n}{group\PYGZus{}leader}\PYG{p}{;}\PYG{+w}{   }\PYG{c+cm}{/* ana thread\PYGZsq{}e ilişkin task\PYGZus{}struct nesnesinin adresi */}
\PYG{+w}{    }\PYG{c+cm}{/* ... */}
\PYG{p}{\PYGZcb{}}\PYG{p}{;}
\end{sphinxVerbatim}

\sphinxAtStartPar
Peki prosesin ana thread’i sonlandığında \sphinxcode{\sphinxupquote{group\_leader}} göstericisinin ve \sphinxcode{\sphinxupquote{signal\_struct}} yapısı
içerisindeki \sphinxcode{\sphinxupquote{list\_head}} düğümünün durumu ne olacaktır? İşte Linux çekirdeği prosesin ana thread’i
sonlandığında istisna olarak ana thread’e ilişkin \sphinxcode{\sphinxupquote{task\_struct}} nesnesini yok etmemektedir (başka bir
deyişle \sphinxstyleemphasis{hortlak (zombie)} olarak tutmaktadır). Yani bu \sphinxcode{\sphinxupquote{group\_leader}} göstericisi her zaman geçerli
bir \sphinxcode{\sphinxupquote{task\_struct}} nesnesini gösteriyor durumda olmaktadır.


\subsection{\sphinxstyleliteralintitle{\sphinxupquote{real\_parent}} ve \sphinxstyleliteralintitle{\sphinxupquote{parent}} Göstericileri}
\label{\detokenize{processcontrolblock:real-parent-ve-parent-gostericileri}}
\sphinxAtStartPar
Şimdi de Linux çekirdeğinde prosesler arasındaki altlık\sphinxhyphen{}üstlük ilişkisinin nasıl sağlandığı üzerinde
duralım. Bilindiği gibi UNIX/Linux sistemlerinde her proses başka bir prosesin thread’i tarafından \sphinxstyleemphasis{fork}
işlemi ile yaratılmaktadır. Bu durumda prosesi yaratan thread’in ilişkin olduğu prosese \sphinxstyleemphasis{üst proses
(parent process)}, yeni yaratılan prosese de \sphinxstyleemphasis{alt process (child process)} denilmektedir.

\sphinxAtStartPar
Her thread’in \sphinxcode{\sphinxupquote{task\_struct}} nesnesi içerisindeki \sphinxcode{\sphinxupquote{real\_parent}} ve \sphinxcode{\sphinxupquote{parent}} isimli göstericiler o
thread’i yaratan thread’in \sphinxcode{\sphinxupquote{task\_struct}} nesnesinin adresini tutmaktadır:

\begin{sphinxVerbatim}[commandchars=\\\{\}]
\PYG{k}{struct}\PYG{+w}{ }\PYG{n+nc}{task\PYGZus{}struct}\PYG{+w}{ }\PYG{p}{\PYGZob{}}
\PYG{+w}{    }\PYG{c+cm}{/* ... */}
\PYG{+w}{    }\PYG{k+kt}{pid\PYGZus{}t}\PYG{+w}{   }\PYG{n}{pid}\PYG{p}{;}
\PYG{+w}{    }\PYG{k+kt}{pid\PYGZus{}t}\PYG{+w}{   }\PYG{n}{tgid}\PYG{p}{;}
\PYG{+w}{    }\PYG{k}{struct}\PYG{+w}{ }\PYG{n+nc}{task\PYGZus{}struct}\PYG{+w}{ }\PYG{o}{*}\PYG{n}{group\PYGZus{}leader}\PYG{p}{;}

\PYG{+w}{    }\PYG{k}{struct}\PYG{+w}{ }\PYG{n+nc}{task\PYGZus{}struct}\PYG{+w}{ }\PYG{n}{\PYGZus{}\PYGZus{}rcu}\PYG{+w}{ }\PYG{o}{*}\PYG{n}{real\PYGZus{}parent}\PYG{p}{;}\PYG{+w}{  }\PYG{c+cm}{/* fork yapan üst prosesteki thread */}
\PYG{+w}{    }\PYG{k}{struct}\PYG{+w}{ }\PYG{n+nc}{task\PYGZus{}struct}\PYG{+w}{ }\PYG{n}{\PYGZus{}\PYGZus{}rcu}\PYG{+w}{ }\PYG{o}{*}\PYG{n}{parent}\PYG{p}{;}\PYG{+w}{       }\PYG{c+cm}{/* fork yapan üst prosesteki thread */}
\PYG{+w}{    }\PYG{c+cm}{/* ... */}
\PYG{p}{\PYGZcb{}}\PYG{p}{;}
\end{sphinxVerbatim}

\sphinxAtStartPar
Burada \sphinxcode{\sphinxupquote{real\_parent}} üst prosese ilişkin thread’in \sphinxcode{\sphinxupquote{task\_struct}} nesnesinin adresini tutmaktadır.
Normalde \sphinxcode{\sphinxupquote{real\_parent}} elemanı ile \sphinxcode{\sphinxupquote{parent}} elemanı aynı \sphinxcode{\sphinxupquote{task\_struct}} nesnesini gösterir. Ancak
seyrek durumlarda (örneğin debug işlemlerinde ve \sphinxstyleemphasis{ptrace} işlemlerinde) geçici olarak \sphinxcode{\sphinxupquote{parent}} başka
bir \sphinxcode{\sphinxupquote{task\_struct}} nesnesini de (reparenting işlemi) gösteriyor durumda olabilmektedir.

\sphinxAtStartPar
\sphinxstyleemphasis{getppid} POSIX fonksiyonu üst prosesin pid değerini vermektedir. İşte bu fonksiyonun çağırdığı
\sphinxstyleemphasis{sys\_getppid} sistem fonksiyonu \sphinxcode{\sphinxupquote{real\_parent}} elemanının gösterdiği \sphinxcode{\sphinxupquote{task\_struct}} nesnesi içerisindeki
\sphinxcode{\sphinxupquote{tgid}} değerini geri döndürmektedir.


\subsection{Alt Proses Listesi: \sphinxstyleliteralintitle{\sphinxupquote{children}} / \sphinxstyleliteralintitle{\sphinxupquote{sibling}}}
\label{\detokenize{processcontrolblock:alt-proses-listesi-children-sibling}}
\sphinxAtStartPar
Şimdi de bir prosesin alt proseslerinin nasıl bağlı listelerde tutulduğu üzerinde duralım. Örneğin bir
proses üç kez \sphinxstyleemphasis{fork} yapmış olsun. Bu durumda bu prosesin üç tane alt prosesi olacaktır. Üst prosesleri
aynı olan proseslere \sphinxstyleemphasis{kardeş prosesler (sibling processes)} de denilmektedir. İşte bir prosesin alt
prosesleri \sphinxcode{\sphinxupquote{children}} isimli kök düğümle erişilen \sphinxcode{\sphinxupquote{sibling}} bağları yoluyla bağlı listede
tutulmaktadır:

\begin{sphinxVerbatim}[commandchars=\\\{\}]
\PYG{k}{struct}\PYG{+w}{ }\PYG{n+nc}{task\PYGZus{}struct}\PYG{+w}{ }\PYG{p}{\PYGZob{}}
\PYG{+w}{    }\PYG{c+cm}{/* ... */}
\PYG{+w}{    }\PYG{k}{struct}\PYG{+w}{ }\PYG{n+nc}{list\PYGZus{}head}\PYG{+w}{ }\PYG{n}{children}\PYG{p}{;}\PYG{+w}{   }\PYG{c+cm}{/* alt proses listesinin kök düğümü */}
\PYG{+w}{    }\PYG{k}{struct}\PYG{+w}{ }\PYG{n+nc}{list\PYGZus{}head}\PYG{+w}{ }\PYG{n}{sibling}\PYG{p}{;}\PYG{+w}{    }\PYG{c+cm}{/* alt proseslerin bağlı liste bağı */}
\PYG{+w}{    }\PYG{c+cm}{/* ... */}
\PYG{p}{\PYGZcb{}}\PYG{p}{;}
\end{sphinxVerbatim}

\sphinxAtStartPar
Aşağıdaki diyagram bu ilişkiyi görsel olarak ortaya koymaktadır. Üst prosesin \sphinxcode{\sphinxupquote{children}} düğümü (kök)
mavi, alt proseslerin \sphinxcode{\sphinxupquote{sibling}} düğümleri turuncu renkte gösterilmiştir:

\sphinxincludegraphics[]{graphviz-6923b7a7c0a2c37818097122869ea9955c9b5ae5.pdf}

\sphinxAtStartPar
\sphinxcode{\sphinxupquote{sibling}} bağları her zaman prosesin ana thread’ine ilişkin (grup liderine ilişkin) \sphinxcode{\sphinxupquote{task\_struct}}
nesnelerini göstermektedir. Bir thread akışında yeni bir thread yaratıldığı zaman yaratılan thread bu
\sphinxcode{\sphinxupquote{children}}/\sphinxcode{\sphinxupquote{sibling}} listesinde bulunmaz; prosesin thread’leri yalnızca \sphinxcode{\sphinxupquote{thread\_head}}/\sphinxcode{\sphinxupquote{thread\_node}}
listesinde tutulmaktadır.

\sphinxAtStartPar
Ayrıca Linux’ta alt proses listesi ana thread’in \sphinxcode{\sphinxupquote{children}} kök düğümünden hareketle dolaşılmak
zorundadır. Alt prosesler yaratıldığında yalnızca ana thread’in \sphinxcode{\sphinxupquote{children}} kök düğümü güncellenmektedir.


\subsection{Tüm Proseslerin Listesi: \sphinxstyleliteralintitle{\sphinxupquote{tasks}}}
\label{\detokenize{processcontrolblock:tum-proseslerin-listesi-tasks}}
\sphinxAtStartPar
Çekirdek ayrıca proseslerin ana thread’lerine ilişkin \sphinxcode{\sphinxupquote{task\_struct}} nesnelerini de \sphinxcode{\sphinxupquote{task\_struct}}
içerisindeki \sphinxcode{\sphinxupquote{tasks}} isimli bir düğüm ile birbirine bağlamaktadır:

\begin{sphinxVerbatim}[commandchars=\\\{\}]
\PYG{k}{struct}\PYG{+w}{ }\PYG{n+nc}{task\PYGZus{}struct}\PYG{+w}{ }\PYG{p}{\PYGZob{}}
\PYG{+w}{    }\PYG{c+cm}{/* ... */}
\PYG{+w}{    }\PYG{k+kt}{pid\PYGZus{}t}\PYG{+w}{ }\PYG{n}{pid}\PYG{p}{;}
\PYG{+w}{    }\PYG{k+kt}{pid\PYGZus{}t}\PYG{+w}{ }\PYG{n}{tgid}\PYG{p}{;}
\PYG{+w}{    }\PYG{k}{struct}\PYG{+w}{ }\PYG{n+nc}{task\PYGZus{}struct}\PYG{+w}{ }\PYG{o}{*}\PYG{n}{group\PYGZus{}leader}\PYG{p}{;}

\PYG{+w}{    }\PYG{k}{struct}\PYG{+w}{ }\PYG{n+nc}{task\PYGZus{}struct}\PYG{+w}{ }\PYG{n}{\PYGZus{}\PYGZus{}rcu}\PYG{+w}{ }\PYG{o}{*}\PYG{n}{real\PYGZus{}parent}\PYG{p}{;}
\PYG{+w}{    }\PYG{k}{struct}\PYG{+w}{ }\PYG{n+nc}{task\PYGZus{}struct}\PYG{+w}{ }\PYG{n}{\PYGZus{}\PYGZus{}rcu}\PYG{+w}{ }\PYG{o}{*}\PYG{n}{parent}\PYG{p}{;}

\PYG{+w}{    }\PYG{k}{struct}\PYG{+w}{ }\PYG{n+nc}{list\PYGZus{}head}\PYG{+w}{ }\PYG{n}{children}\PYG{p}{;}
\PYG{+w}{    }\PYG{k}{struct}\PYG{+w}{ }\PYG{n+nc}{list\PYGZus{}head}\PYG{+w}{ }\PYG{n}{sibling}\PYG{p}{;}
\PYG{+w}{    }\PYG{k}{struct}\PYG{+w}{ }\PYG{n+nc}{list\PYGZus{}head}\PYG{+w}{ }\PYG{n}{tasks}\PYG{p}{;}\PYG{+w}{         }\PYG{c+cm}{/* ana thread\PYGZsq{}lerin tutulduğu liste bağı */}
\PYG{+w}{    }\PYG{c+cm}{/* ... */}
\PYG{p}{\PYGZcb{}}\PYG{p}{;}
\end{sphinxVerbatim}

\sphinxAtStartPar
Bu \sphinxcode{\sphinxupquote{tasks}} düğümlerine ilişkin bağlı listenin kök düğümü \sphinxcode{\sphinxupquote{init\_task}} prosesindedir. Linux’ta boot
işlemi sırasında ilk oluşturulan prosese \sphinxstyleemphasis{init\_task} denilmektedir. Bu \sphinxstyleemphasis{init\_task} prosesine ilişkin
\sphinxcode{\sphinxupquote{task\_struct}} nesnesi statik bir biçimde \sphinxcode{\sphinxupquote{init/init\_task.c}} dosyasında bulunmaktadır. \sphinxstyleemphasis{init\_task}
prosesinin pid değeri 0’dır. \sphinxstyleemphasis{init\_task} prosesi \sphinxstyleemphasis{init} prosesini yarattıktan sonra işlevini sonlandırarak
durdurulmaktadır. Ancak bu prosese ilişkin \sphinxcode{\sphinxupquote{task\_struct}} nesnesi gerçek anlamda hiçbir zaman yok
edilmemektedir. İşte \sphinxcode{\sphinxupquote{init\_task.tasks}} elemanı proseslerin ana thread’lerine ilişkin \sphinxcode{\sphinxupquote{task\_struct}}
nesnelerini dolaşmak için bir kök düğüm olarak kullanılmaktadır.

\sphinxAtStartPar
Aşağıdaki diyagram \sphinxcode{\sphinxupquote{tasks}} bağlı listesini göstermektedir. Kök düğüm \sphinxcode{\sphinxupquote{init\_task}} (PID=0) koyu mavi,
proseslerin ana thread \sphinxcode{\sphinxupquote{task\_struct}} nesneleri açık mavi renktedir:

\sphinxincludegraphics[]{graphviz-33d569426d8a8e7af78eec10ded079a2cdd0fc8e.pdf}


\subsection{Özet: \sphinxstyleliteralintitle{\sphinxupquote{task\_struct}} İlişkileri}
\label{\detokenize{processcontrolblock:ozet-task-struct-iliskileri}}
\sphinxAtStartPar
Yukarıda açıkladığımız konuyu madde madde özetleyelim:
\begin{itemize}
\item {} 
\sphinxAtStartPar
Her thread’in ayrı bir pid değeri vardır, ancak POSIX’teki \sphinxstyleemphasis{getpid} fonksiyonu ana thread’in (thread
grup liderinin) pid değerini vermektedir.

\item {} 
\sphinxAtStartPar
Bir prosesin tüm thread’lerini dolaşmak için \sphinxcode{\sphinxupquote{signal}} göstericisinin gösterdiği \sphinxcode{\sphinxupquote{signal\_struct}}
yapısı içerisindeki \sphinxcode{\sphinxupquote{thread\_head}} bağlı listesi \sphinxcode{\sphinxupquote{thread\_node}} düğümleriyle dolaşılır.

\item {} 
\sphinxAtStartPar
\sphinxcode{\sphinxupquote{task\_struct}} içerisindeki \sphinxcode{\sphinxupquote{real\_parent}} ve \sphinxcode{\sphinxupquote{parent}} göstericileri o prosesi ya da thread’i
yaratan thread’in \sphinxcode{\sphinxupquote{task\_struct}} nesnesini göstermektedir.

\item {} 
\sphinxAtStartPar
\sphinxcode{\sphinxupquote{task\_struct}} içerisindeki \sphinxcode{\sphinxupquote{tgid}} prosesin ana thread’inin pid değerini, \sphinxcode{\sphinxupquote{group\_leader}}
göstericisi ise prosesin ana thread’inin \sphinxcode{\sphinxupquote{task\_struct}} nesnesinin adresini tutmaktadır. Ana thread
sonlansa bile onun \sphinxcode{\sphinxupquote{task\_struct}} nesnesi proses sonlanana kadar muhafaza edilmektedir.

\item {} 
\sphinxAtStartPar
Bir prosesin bütün alt proseslerine ilişkin ana thread’lerin \sphinxcode{\sphinxupquote{task\_struct}} nesneleri
\sphinxcode{\sphinxupquote{children}}/\sphinxcode{\sphinxupquote{sibling}} bağlı listesinde tutulmaktadır.

\item {} 
\sphinxAtStartPar
Prosesin alt proseslerini elde etmek için prosesin ana thread’inin \sphinxcode{\sphinxupquote{children}} kök düğümünden yola
çıkmak gerekir.

\item {} 
\sphinxAtStartPar
Çekirdek yaratılmış olan tüm proseslerin ana thread’lerine ilişkin \sphinxcode{\sphinxupquote{task\_struct}} nesnelerini ayrıca
\sphinxcode{\sphinxupquote{task\_struct}} içerisindeki \sphinxcode{\sphinxupquote{tasks}} düğümüyle birbirine bağlamaktadır. Bu düğümün kökü de
\sphinxcode{\sphinxupquote{init\_task.tasks}} elemanıdır.

\end{itemize}


\subsection{Sık Sorulan Sorular}
\label{\detokenize{processcontrolblock:sik-sorulan-sorular}}
\sphinxAtStartPar
Aşağıda konuyla ilgili bazı işlemlerin bu bağlı listeler kullanılarak nasıl gerçekleştirileceğine yönelik
tipik sorular ve yanıtları verilmektedir:

\sphinxAtStartPar
\sphinxstylestrong{Soru: Bir prosesin tüm thread’leri nasıl elde edilebilir?}

\sphinxAtStartPar
Bir prosesin tüm thread’leri prosesin \sphinxcode{\sphinxupquote{signal}} yapısı içerisinde bulunan \sphinxcode{\sphinxupquote{thread\_head}} elemanı kök
yapılarak \sphinxcode{\sphinxupquote{thread\_node}} düğümlerinin dolaşılmasıyla elde edilebilir.

\sphinxAtStartPar
\sphinxstylestrong{Soru: Bir prosesin bütün alt prosesleri nasıl elde edilebilir?}

\sphinxAtStartPar
Prosesin ana thread’indeki \sphinxcode{\sphinxupquote{children}} düğümü kök yapılarak \sphinxcode{\sphinxupquote{sibling}} düğümlerinin dolaşılmasıyla
prosesin tüm alt prosesleri elde edilebilir.

\sphinxAtStartPar
\sphinxstylestrong{Soru: Bir \textasciigrave{}\textasciigrave{}task\_struct\textasciigrave{}\textasciigrave{} nesnesinin prosesin ana thread’ine ilişkin olup olmadığı nasıl anlaşılır?}

\sphinxAtStartPar
Eğer \sphinxcode{\sphinxupquote{task\_struct}} nesnesinde \sphinxcode{\sphinxupquote{pid == tgid}} ise bu \sphinxcode{\sphinxupquote{task\_struct}} nesnesi ilgili prosesin ana
thread’ine ilişkin \sphinxcode{\sphinxupquote{task\_struct}} nesnesidir. Zaten prosesin tüm thread’lerine ilişkin \sphinxcode{\sphinxupquote{task\_struct}}
nesnelerinde \sphinxcode{\sphinxupquote{group\_leader}} göstericisi ana thread’e ilişkin bu \sphinxcode{\sphinxupquote{task\_struct}} nesnesini
göstermektedir. O anda çalışmakta olan thread eğer prosesin ana thread’i ise \sphinxcode{\sphinxupquote{current == group\_leader}}
koşulunun sağlanacağına da dikkat ediniz.

\sphinxAtStartPar
\sphinxstylestrong{Soeu: Sistemdeki tüm proseslerin ana thread’lerine ilişkin \textasciigrave{}\textasciigrave{}task\_struct\textasciigrave{}\textasciigrave{} nesneleri nasıl dolaşılır?}

\sphinxAtStartPar
\sphinxcode{\sphinxupquote{init\_task}} prosesindeki \sphinxcode{\sphinxupquote{tasks}} düğümü kök düğüm yapılıp \sphinxcode{\sphinxupquote{tasks}} düğümleri dolaşılırsa tüm
proseslerin ana thread’lerine ilişkin \sphinxcode{\sphinxupquote{task\_struct}} nesneleri elde edilebilir.


\section{task\_struct Yapısına ilişkin Bağlı Listeler Üzerinde Gezinme İşlemleri}
\label{\detokenize{processcontrolblock:task-struct-yapisina-iliskin-bagli-listeler-uzerinde-gezinme-islemleri}}
\sphinxAtStartPar
Şimdi de görmüş olduğumuz \sphinxcode{\sphinxupquote{task\_struct}} yapısına ilişkin bağlı listeler üzerinde gezinme işlemlerine
örnekler verelim. Bu biçimdeki küçük testler için çekirdeği yeniden derlememize gerek yoktur. Basit bir
karakter aygıt sürücüsü yazarak da bu işlemleri yapabiliriz.


\subsection{Geliştirme Ortamının Hazırlanması}
\label{\detokenize{processcontrolblock:gelistirme-ortaminin-hazirlanmasi}}
\sphinxAtStartPar
Genellikle bir Linux sistemini yüklediğimizde zaten çekirdek modüllerini ve aygıt sürücüleri
oluşturabilmek için gereken başlık dosyaları ve diğer gerekli öğeler \sphinxcode{\sphinxupquote{/usr/src}} dizini içerisindeki
\sphinxcode{\sphinxupquote{linux\sphinxhyphen{}headers\sphinxhyphen{}\$(uname \sphinxhyphen{}r)}} dizininde yüklü biçimde bulunmaktadır. Ancak bunlar yüklü değilse Debian
tabanlı sistemlerde bunları şöyle yükleyebilirsiniz:

\begin{sphinxVerbatim}[commandchars=\\\{\}]
\PYGZdl{}\PYG{+w}{ }sudo\PYG{+w}{ }apt\PYGZhy{}get\PYG{+w}{ }install\PYG{+w}{ }linux\PYGZhy{}headers\PYGZhy{}\PYG{k}{\PYGZdl{}(}uname\PYG{+w}{ }\PYGZhy{}r\PYG{k}{)}
\end{sphinxVerbatim}


\subsection{Derleme ve Yükleme}
\label{\detokenize{processcontrolblock:derleme-ve-yukleme}}
\sphinxAtStartPar
Karakter aygıt sürücüleri bir ya da birden fazla kaynak \sphinxcode{\sphinxupquote{.c}} dosyası oluşturularak yazılabilmektedir.
Aygıt sürücülerin derlenmesi oldukça karmaşık bir süreçle gerçekleşmektedir. Bu nedenle derleme işlemi
için hazır \sphinxcode{\sphinxupquote{Makefile}} dosyaları bulundurulmuştur. Programcı kendisi bir \sphinxcode{\sphinxupquote{Makefile}} dosyası oluşturup
asıl make dosyalarını buradan devreye sokar. Aygıt sürücüler için en basit \sphinxcode{\sphinxupquote{Makefile}} dosyası aşağıdaki
gibi oluşturulabilir:

\begin{sphinxVerbatim}[commandchars=\\\{\}]
\PYG{n+nv}{obj\PYGZhy{}m}\PYG{+w}{ }\PYG{o}{+=}\PYG{+w}{ }\PYG{l+s+si}{\PYGZdl{}\PYGZob{}}\PYG{n+nv}{file}\PYG{l+s+si}{\PYGZcb{}}.o

\PYG{n+nf}{all}\PYG{o}{:}
\PYG{+w}{    }make\PYG{+w}{ }\PYGZhy{}C\PYG{+w}{ }/lib/modules/\PYG{k}{\PYGZdl{}(}shell\PYG{+w}{ }uname\PYG{+w}{ }\PYGZhy{}r\PYG{k}{)}/build\PYG{+w}{ }\PYG{n+nv}{M}\PYG{o}{=}\PYG{l+s+si}{\PYGZdl{}\PYGZob{}}\PYG{n+nv}{PWD}\PYG{l+s+si}{\PYGZcb{}}\PYG{+w}{ }modules
\PYG{n+nf}{clean}\PYG{o}{:}
\PYG{+w}{    }make\PYG{+w}{ }\PYGZhy{}C\PYG{+w}{ }/lib/modules/\PYG{k}{\PYGZdl{}(}shell\PYG{+w}{ }uname\PYG{+w}{ }\PYGZhy{}r\PYG{k}{)}/build\PYG{+w}{ }\PYG{n+nv}{M}\PYG{o}{=}\PYG{l+s+si}{\PYGZdl{}\PYGZob{}}\PYG{n+nv}{PWD}\PYG{l+s+si}{\PYGZcb{}}\PYG{+w}{ }clean
\end{sphinxVerbatim}

\sphinxAtStartPar
Bu \sphinxcode{\sphinxupquote{Makefile}} dışarıdan \sphinxcode{\sphinxupquote{file}} isimli bir argüman almaktadır. Aygıt sürücünün derlenmesi şöyle
yapılır (kaynak dosya için uzantı belirtilmediğine dikkat ediniz):

\begin{sphinxVerbatim}[commandchars=\\\{\}]
\PYGZdl{}\PYG{+w}{ }make\PYG{+w}{ }\PYG{n+nv}{file}\PYG{o}{=}mydriver
\end{sphinxVerbatim}

\sphinxAtStartPar
Aygıt sürücüler derlendikten sonra \sphinxcode{\sphinxupquote{.ko}} uzantılı bir dosya elde edilecektir. Bu dosya çekirdeğe
yüklenmelidir. Aygıt sürücü dosyalarını (\sphinxcode{\sphinxupquote{.ko}} dosyalarını) çekirdeğe yüklemek için \sphinxstyleemphasis{insmod} ya da
\sphinxstyleemphasis{modprobe} komutları kullanılmaktadır. \sphinxstyleemphasis{insmod} komutu herhangi bir bağımlılığa bakmadan doğrudan yükleme
yapar. Biz kursumuzda aygıt sürücüleri \sphinxstyleemphasis{insmod} komutuyla yükleyeceğiz. Tabii aygıt sürücüleri
yükleyebilmek için \sphinxstyleemphasis{root} önceliğine sahip olmak gerekir. Örneğin:

\begin{sphinxVerbatim}[commandchars=\\\{\}]
\PYGZdl{}\PYG{+w}{ }sudo\PYG{+w}{ }insmod\PYG{+w}{ }mydriver.ko
\end{sphinxVerbatim}

\sphinxAtStartPar
\sphinxstyleemphasis{insmod} ile yüklenmiş olan aygıt sürücü çekirdekten \sphinxstyleemphasis{rmmod} komutuyla çıkartılır:

\begin{sphinxVerbatim}[commandchars=\\\{\}]
\PYGZdl{}\PYG{+w}{ }sudo\PYG{+w}{ }rmmod\PYG{+w}{ }mydriver.ko
\end{sphinxVerbatim}

\sphinxAtStartPar
Aygıt sürücünün yüklenip yüklenmediğini anlayabilmek için \sphinxcode{\sphinxupquote{/proc/devices}} dosyasına başvurabilirsiniz.


\subsection{Çekirdek Modülü ve Aygıt Sürücü Kavramları}
\label{\detokenize{processcontrolblock:cekirdek-modulu-ve-aygit-surucu-kavramlari}}
\sphinxAtStartPar
Aygıt sürücüler çekirdek modunda çekirdeğin bir parçasıymış gibi çalışan modüllerdir. Aygıt sürücüler
içerisinde kullanıcı modundaki standart C fonksiyonları ya da POSIX fonksiyonları kullanılamaz. Aygıt
sürücüler yalnızca çekirdek içerisinde \sphinxstyleemphasis{export} edilmiş fonksiyonları kullanabilirler. Bir fonksiyon ya
da nesne çekirdek kodlarında \sphinxcode{\sphinxupquote{EXPORT\_SYMBOL}} makrosu ile export edilmektedir. Örneğin:

\begin{sphinxVerbatim}[commandchars=\\\{\}]
\PYG{k+kt}{void}\PYG{+w}{ }\PYG{n+nf}{clear\PYGZus{}nlink}\PYG{p}{(}\PYG{k}{struct}\PYG{+w}{ }\PYG{n+nc}{inode}\PYG{+w}{ }\PYG{o}{*}\PYG{n}{inode}\PYG{p}{)}
\PYG{p}{\PYGZob{}}
\PYG{+w}{    }\PYG{k}{if}\PYG{+w}{ }\PYG{p}{(}\PYG{n}{inode}\PYG{o}{\PYGZhy{}}\PYG{o}{\PYGZgt{}}\PYG{n}{i\PYGZus{}nlink}\PYG{p}{)}\PYG{+w}{ }\PYG{p}{\PYGZob{}}
\PYG{+w}{        }\PYG{n}{inode}\PYG{o}{\PYGZhy{}}\PYG{o}{\PYGZgt{}}\PYG{n}{\PYGZus{}\PYGZus{}i\PYGZus{}nlink}\PYG{+w}{ }\PYG{o}{=}\PYG{+w}{ }\PYG{l+m+mi}{0}\PYG{p}{;}
\PYG{+w}{        }\PYG{n}{atomic\PYGZus{}long\PYGZus{}inc}\PYG{p}{(}\PYG{o}{\PYGZam{}}\PYG{n}{inode}\PYG{o}{\PYGZhy{}}\PYG{o}{\PYGZgt{}}\PYG{n}{i\PYGZus{}sb}\PYG{o}{\PYGZhy{}}\PYG{o}{\PYGZgt{}}\PYG{n}{s\PYGZus{}remove\PYGZus{}count}\PYG{p}{)}\PYG{p}{;}
\PYG{+w}{    }\PYG{p}{\PYGZcb{}}
\PYG{p}{\PYGZcb{}}
\PYG{n}{EXPORT\PYGZus{}SYMBOL}\PYG{p}{(}\PYG{n}{clear\PYGZus{}nlink}\PYG{p}{)}\PYG{p}{;}
\end{sphinxVerbatim}

\sphinxAtStartPar
Aslında Linux sistemlerinde bu bağlamda birbirleriyle ilişkili iki kavram vardır: \sphinxstyleemphasis{Çekirdek modülleri
(kernel modules)} ve \sphinxstyleemphasis{aygıt sürücüler (device drivers)}. Çekirdeğin içerisine yüklenebilen modüllere
“çekirdek modülleri” denilmektedir. Eğer bir çekirdek modülüne kullanıcı modundan bir dosya (buna \sphinxstyleemphasis{aygıt
dosyası (device file)} denilmektedir) yoluyla erişilip aygıt sürücüye işlemler yaptırılabiliyorsa bu tür
çekirdek modüllerine \sphinxstyleemphasis{aygıt sürücü (device driver)} de denilmektedir. Yani her aygıt sürücü bir çekirdek
modülü belirtir ancak her çekirdek modülü bir aygıt sürücü belirtmek zorunda değildir. Sistemdeki yüklü
olan çekirdek modülleri \sphinxcode{\sphinxupquote{/proc/modules}} dosyası yoluyla, yüklü olan aygıt sürücüler ise
\sphinxcode{\sphinxupquote{/proc/devices}} dosyası yoluyla görüntülenebilmektedir.

\sphinxAtStartPar
Aşağıda iskelet bir çekirdek modülü görüyorsunuz:

\begin{sphinxVerbatim}[commandchars=\\\{\}]
\PYG{c+cp}{\PYGZsh{}}\PYG{c+cp}{include}\PYG{+w}{ }\PYG{c+cpf}{\PYGZlt{}linux/module.h\PYGZgt{}}
\PYG{c+cp}{\PYGZsh{}}\PYG{c+cp}{include}\PYG{+w}{ }\PYG{c+cpf}{\PYGZlt{}linux/kernel.h\PYGZgt{}}

\PYG{n}{MODULE\PYGZus{}LICENSE}\PYG{p}{(}\PYG{l+s}{\PYGZdq{}}\PYG{l+s}{GPL}\PYG{l+s}{\PYGZdq{}}\PYG{p}{)}\PYG{p}{;}

\PYG{k}{static}\PYG{+w}{ }\PYG{k+kt}{int}\PYG{+w}{ }\PYG{n+nf}{helloworld\PYGZus{}init}\PYG{p}{(}\PYG{k+kt}{void}\PYG{p}{)}\PYG{p}{;}
\PYG{k}{static}\PYG{+w}{ }\PYG{k+kt}{void}\PYG{+w}{ }\PYG{n+nf}{helloworld\PYGZus{}exit}\PYG{p}{(}\PYG{k+kt}{void}\PYG{p}{)}\PYG{p}{;}

\PYG{k}{static}\PYG{+w}{ }\PYG{k+kt}{int}\PYG{+w}{ }\PYG{n+nf}{helloworld\PYGZus{}init}\PYG{p}{(}\PYG{k+kt}{void}\PYG{p}{)}
\PYG{p}{\PYGZob{}}
\PYG{+w}{    }\PYG{n}{printk}\PYG{p}{(}\PYG{n}{KERN\PYGZus{}INFO}\PYG{+w}{ }\PYG{l+s}{\PYGZdq{}}\PYG{l+s}{Hello World...}\PYG{l+s+se}{\PYGZbs{}n}\PYG{l+s}{\PYGZdq{}}\PYG{p}{)}\PYG{p}{;}

\PYG{+w}{    }\PYG{k}{return}\PYG{+w}{ }\PYG{l+m+mi}{0}\PYG{p}{;}
\PYG{p}{\PYGZcb{}}

\PYG{k}{static}\PYG{+w}{ }\PYG{k+kt}{void}\PYG{+w}{ }\PYG{n+nf}{helloworld\PYGZus{}exit}\PYG{p}{(}\PYG{k+kt}{void}\PYG{p}{)}
\PYG{p}{\PYGZob{}}
\PYG{+w}{    }\PYG{n}{printk}\PYG{p}{(}\PYG{n}{KERN\PYGZus{}INFO}\PYG{+w}{ }\PYG{l+s}{\PYGZdq{}}\PYG{l+s}{Goodbye World...}\PYG{l+s+se}{\PYGZbs{}n}\PYG{l+s}{\PYGZdq{}}\PYG{p}{)}\PYG{p}{;}
\PYG{p}{\PYGZcb{}}

\PYG{n}{module\PYGZus{}init}\PYG{p}{(}\PYG{n}{helloworld\PYGZus{}init}\PYG{p}{)}\PYG{p}{;}
\PYG{n}{module\PYGZus{}exit}\PYG{p}{(}\PYG{n}{helloworld\PYGZus{}exit}\PYG{p}{)}\PYG{p}{;}
\end{sphinxVerbatim}

\sphinxAtStartPar
Bir çekirdek modülü \sphinxstyleemphasis{insmod} komutuyla yüklendiğinde onun \sphinxcode{\sphinxupquote{module\_init}} makrosuyla belirtilen fonksiyonu
çağrılmaktadır. Çekirdek modülü \sphinxstyleemphasis{rmmod} ile sistemden çıkartılırken de \sphinxcode{\sphinxupquote{module\_exit}} makrosunda
belirtilen fonksiyon çağrılmaktadır.

\sphinxAtStartPar
Çekirdek modülleri içerisinde birtakım mesajlar iletilmek isteniyorsa bu mesajlar \sphinxstyleemphasis{kernel ring buffer}
denilen log sistemine yazdırılmaktadır. Bu log sistemine yazdırılan yazılar \sphinxstyleemphasis{dmesg} komutuyla ya da
\sphinxcode{\sphinxupquote{/var/log/syslog}} dosyası ile görüntülenebilmektedir. Çekirdek modülü içerisinde bu log sistemine
mesajlar \sphinxstyleemphasis{printk} isimli çekirdek fonksiyonuyla yazılmaktadır. Örneğin:

\begin{sphinxVerbatim}[commandchars=\\\{\}]
\PYG{n}{printk}\PYG{p}{(}\PYG{n}{KERN\PYGZus{}INFO}\PYG{+w}{ }\PYG{l+s}{\PYGZdq{}}\PYG{l+s}{Goodbye World...}\PYG{l+s+se}{\PYGZbs{}n}\PYG{l+s}{\PYGZdq{}}\PYG{p}{)}\PYG{p}{;}
\end{sphinxVerbatim}

\sphinxAtStartPar
\sphinxcode{\sphinxupquote{KERN\_INFO}} mesajın türünü belirtmektedir. \sphinxcode{\sphinxupquote{KERN\_INFO}} ile diğer argüman arasında virgül atomunun
bulunmadığına dikkat ediniz. \sphinxcode{\sphinxupquote{KERN\_INFO}} aslında bir sembolik sabittir ve bir string açımı
yapmaktadır. (C’de yana yana iki string’in birleştirildiğini anımsayınız.) \sphinxcode{\sphinxupquote{KERN\_INFO}} mesaj türü ile
mesaj yazdırmanın daha basit yolu \sphinxcode{\sphinxupquote{pr\_info}} makrosunu kullanmaktadır:

\begin{sphinxVerbatim}[commandchars=\\\{\}]
\PYG{n}{pr\PYGZus{}info}\PYG{p}{(}\PYG{l+s}{\PYGZdq{}}\PYG{l+s}{Goodbye World...}\PYG{l+s+se}{\PYGZbs{}n}\PYG{l+s}{\PYGZdq{}}\PYG{p}{)}\PYG{p}{;}
\end{sphinxVerbatim}


\subsection{Aygıt Dosyaları}
\label{\detokenize{processcontrolblock:aygit-dosyalari}}
\sphinxAtStartPar
Aygıt sürücülere erişmekte kullanılan özel dosyalara \sphinxstyleemphasis{aygıt dosyaları (device files)} denilmektedir.
Aygıt dosyası \sphinxstyleemphasis{open} POSIX fonksiyonuyla açıldığında çekirdek diske değil yüklü olan aygıt sürücüye
referans etmektedir. Kullanıcı modundaki programcı aygıt dosyası üzerinde dosya işlemlerini yaptığında
aslında programın akışı çekirdek moduna geçerek aygıt sürücü içerisindeki ilgili fonksiyonlar
çalıştırılmaktadır:

\sphinxincludegraphics[]{graphviz-03f40a3a12b11b1d8024af24d85ea2d2db2fd79b.pdf}

\sphinxAtStartPar
Linux sistemlerinde aygıt sürücüler (ve dolayısıyla aygıt dosyaları) iki kısma ayrılmaktadır:
\begin{itemize}
\item {} 
\sphinxAtStartPar
\sphinxstyleemphasis{Karakter Aygıt Sürücüleri (Character Device Drivers)}

\item {} 
\sphinxAtStartPar
\sphinxstyleemphasis{Blok Aygıt Sürücüleri (Block Device Drivers)}

\end{itemize}

\sphinxAtStartPar
Blok aygıt sürücüleri \sphinxstyleemphasis{hard disk, SSD, CDROM, ramdisk} gibi blok blok transferlerin yapıldığı aygıtlar
söz konusu olduğunda kullanılmaktadır. Karakter aygıt sürücülerine ilişkin aygıt dosyaları UNIX/Linux
sistemlerinde \sphinxcode{\sphinxupquote{c}} dosya türü ile, blok aygıt sürücülerine ilişkin aygıt dosyaları ise \sphinxcode{\sphinxupquote{b}} dosya türü
ile temsil edilmektedir.

\sphinxAtStartPar
UNIX/Linux sistemlerinde aygıt dosyaları genel olarak \sphinxcode{\sphinxupquote{/dev}} dizininde bulunmaktadır. Linux
sistemlerinde \sphinxstyleemphasis{udev} denilen bir servis (daemon) ile bu dizinde aygıt dosyalarının yaratılması mümkün
hale getirilmiştir.


\subsection{Aygıt Sürücülerin Majör ve Minör Numaraları}
\label{\detokenize{processcontrolblock:aygit-suruculerin-major-ve-minor-numaralari}}
\sphinxAtStartPar
Her aygıt sürücünün \sphinxstyleemphasis{majör} ve \sphinxstyleemphasis{minör} numarası vardır. Bu majör ve minör numaralar çekirdeğin aygıt
sürücüye erişebilmesi için bir adres görevi görmektedir. Majör numara aygıt sürücünün ana kodunu temsil
eder. Minör numara ise onun örneklerini (instances) temsil etmektedir.

\sphinxAtStartPar
Aygıt dosyaları komut satırında \sphinxstyleemphasis{mknod} isimli programla yaratılmaktadır:

\begin{sphinxVerbatim}[commandchars=\\\{\}]
\PYGZdl{}\PYG{+w}{ }sudo\PYG{+w}{ }mknod\PYG{+w}{ }\PYGZhy{}m\PYG{+w}{ }\PYG{l+m}{666}\PYG{+w}{ }mydriver\PYG{+w}{ }c\PYG{+w}{ }\PYG{l+m}{250}\PYG{+w}{ }\PYG{l+m}{0}
\end{sphinxVerbatim}

\sphinxAtStartPar
Tipik olarak sistem programcısı bir \sphinxstyleemphasis{kabuk betiği (shell script)} yazarak bu işlemi otomatize eder. Bu
betik önce \sphinxstyleemphasis{insmod} ile aygıt sürücüyü yükler, \sphinxcode{\sphinxupquote{/proc/devices}} dosyasına başvurup onun majör numarasını
elde eder ve aygıt dosyasını \sphinxstyleemphasis{mknod} komutuyla dinamik bir biçimde oluşturur. Biz de denemelerimizde
böyle yapacağız. Bunun için aşağıdaki gibi bir \sphinxstyleemphasis{load} betiği kullanacağız:

\begin{sphinxVerbatim}[commandchars=\\\{\}]
\PYG{c+ch}{\PYGZsh{}!/bin/bash}
\PYG{c+c1}{\PYGZsh{} load  (bu satırı dosyaya kopyalamayınız)}

\PYG{n+nv}{module}\PYG{o}{=}\PYG{n+nv}{\PYGZdl{}1}
\PYG{n+nv}{mode}\PYG{o}{=}\PYG{l+m}{666}

/sbin/insmod\PYG{+w}{ }./\PYG{l+s+si}{\PYGZdl{}\PYGZob{}}\PYG{n+nv}{module}\PYG{l+s+si}{\PYGZcb{}}.ko\PYG{+w}{ }\PYG{l+s+si}{\PYGZdl{}\PYGZob{}}\PYG{p}{@}\PYG{p}{:}\PYG{n+nv}{2}\PYG{l+s+si}{\PYGZcb{}}\PYG{+w}{ }\PYG{o}{||}\PYG{+w}{ }\PYG{n+nb}{exit}\PYG{+w}{ }\PYG{l+m}{1}
\PYG{n+nv}{major}\PYG{o}{=}\PYG{k}{\PYGZdl{}(}awk\PYG{+w}{ }\PYG{l+s+s2}{\PYGZdq{}}\PYG{l+s+s2}{\PYGZbs{}\PYGZdl{}2 == \PYGZbs{}\PYGZdq{}}\PYG{n+nv}{\PYGZdl{}module}\PYG{l+s+s2}{\PYGZbs{}\PYGZdq{} \PYGZob{}print \PYGZbs{}\PYGZdl{}1\PYGZcb{}}\PYG{l+s+s2}{\PYGZdq{}}\PYG{+w}{ }/proc/devices\PYG{k}{)}
rm\PYG{+w}{ }\PYGZhy{}f\PYG{+w}{ }\PYG{n+nv}{\PYGZdl{}module}
mknod\PYG{+w}{ }\PYGZhy{}m\PYG{+w}{ }\PYG{n+nv}{\PYGZdl{}mode}\PYG{+w}{ }\PYG{n+nv}{\PYGZdl{}module}\PYG{+w}{ }c\PYG{+w}{ }\PYG{n+nv}{\PYGZdl{}major}\PYG{+w}{ }\PYG{l+m}{0}
\end{sphinxVerbatim}

\sphinxAtStartPar
Bu betik dosyasını oluşturduktan sonra dosyaya \sphinxcode{\sphinxupquote{x}} hakkını vermeyi unutmayınız:

\begin{sphinxVerbatim}[commandchars=\\\{\}]
\PYGZdl{}\PYG{+w}{ }chmod\PYG{+w}{ }+x\PYG{+w}{ }load
\PYGZdl{}\PYG{+w}{ }sudo\PYG{+w}{ }./load\PYG{+w}{ }mydriver
\end{sphinxVerbatim}

\sphinxAtStartPar
Aygıt sürücüyü \sphinxstyleemphasis{rmmod} ile çıkartıp ilgili aygıt dosyasını silen \sphinxstyleemphasis{unload} betiği de aşağıdaki gibi
yazılabilir:

\begin{sphinxVerbatim}[commandchars=\\\{\}]
\PYG{c+ch}{\PYGZsh{}!/bin/bash}
\PYG{c+c1}{\PYGZsh{} unload  (bu satırı dosyaya kopyalamayınız)}

\PYG{n+nv}{module}\PYG{o}{=}\PYG{n+nv}{\PYGZdl{}1}

/sbin/rmmod\PYG{+w}{ }./\PYG{l+s+si}{\PYGZdl{}\PYGZob{}}\PYG{n+nv}{module}\PYG{l+s+si}{\PYGZcb{}}.ko\PYG{+w}{ }\PYG{o}{||}\PYG{+w}{ }\PYG{n+nb}{exit}\PYG{+w}{ }\PYG{l+m}{1}
rm\PYG{+w}{ }\PYGZhy{}f\PYG{+w}{ }\PYG{n+nv}{\PYGZdl{}module}
\end{sphinxVerbatim}

\begin{sphinxVerbatim}[commandchars=\\\{\}]
\PYGZdl{}\PYG{+w}{ }chmod\PYG{+w}{ }+x\PYG{+w}{ }unload
\PYGZdl{}\PYG{+w}{ }sudo\PYG{+w}{ }./unload\PYG{+w}{ }mydriver
\end{sphinxVerbatim}


\subsection{İskelet Karakter Aygıt Sürücüsü Örneği}
\label{\detokenize{processcontrolblock:iskelet-karakter-aygit-surucusu-ornegi}}
\sphinxAtStartPar
Aşağıda iskelet bir karakter aygıt sürücüsü oluşturulmuştur. Bu karakter aygıt sürücüsü için bir ioctl
kodu da hazırlanmıştır.

\sphinxAtStartPar
\sphinxstylestrong{test\sphinxhyphen{}driver.h:}

\begin{sphinxVerbatim}[commandchars=\\\{\}]
\PYG{c+cp}{\PYGZsh{}}\PYG{c+cp}{ifndef TEST\PYGZus{}DRIVER\PYGZus{}H\PYGZus{}}
\PYG{c+cp}{\PYGZsh{}}\PYG{c+cp}{define TEST\PYGZus{}DRIVER\PYGZus{}H\PYGZus{}}

\PYG{c+cp}{\PYGZsh{}}\PYG{c+cp}{include}\PYG{+w}{ }\PYG{c+cpf}{\PYGZlt{}linux/stddef.h\PYGZgt{}}
\PYG{c+cp}{\PYGZsh{}}\PYG{c+cp}{include}\PYG{+w}{ }\PYG{c+cpf}{\PYGZlt{}linux/ioctl.h\PYGZgt{}}

\PYG{c+cp}{\PYGZsh{}}\PYG{c+cp}{define TEST\PYGZus{}DRIVER\PYGZus{}MAGIC    \PYGZsq{}t\PYGZsq{}}
\PYG{c+cp}{\PYGZsh{}}\PYG{c+cp}{define IOC\PYGZus{}TEST             \PYGZus{}IO(TEST\PYGZus{}DRIVER\PYGZus{}MAGIC, 0)}

\PYG{c+cp}{\PYGZsh{}}\PYG{c+cp}{endif}
\end{sphinxVerbatim}

\sphinxAtStartPar
\sphinxstylestrong{test\sphinxhyphen{}driver.c:}

\begin{sphinxVerbatim}[commandchars=\\\{\}]
\PYG{c+cp}{\PYGZsh{}}\PYG{c+cp}{include}\PYG{+w}{ }\PYG{c+cpf}{\PYGZlt{}linux/module.h\PYGZgt{}}
\PYG{c+cp}{\PYGZsh{}}\PYG{c+cp}{include}\PYG{+w}{ }\PYG{c+cpf}{\PYGZlt{}linux/kernel.h\PYGZgt{}}
\PYG{c+cp}{\PYGZsh{}}\PYG{c+cp}{include}\PYG{+w}{ }\PYG{c+cpf}{\PYGZlt{}linux/fs.h\PYGZgt{}}
\PYG{c+cp}{\PYGZsh{}}\PYG{c+cp}{include}\PYG{+w}{ }\PYG{c+cpf}{\PYGZlt{}linux/cdev.h\PYGZgt{}}
\PYG{c+cp}{\PYGZsh{}}\PYG{c+cp}{include}\PYG{+w}{ }\PYG{c+cpf}{\PYGZdq{}test\PYGZhy{}driver.h\PYGZdq{}}

\PYG{n}{MODULE\PYGZus{}LICENSE}\PYG{p}{(}\PYG{l+s}{\PYGZdq{}}\PYG{l+s}{GPL}\PYG{l+s}{\PYGZdq{}}\PYG{p}{)}\PYG{p}{;}
\PYG{n}{MODULE\PYGZus{}AUTHOR}\PYG{p}{(}\PYG{l+s}{\PYGZdq{}}\PYG{l+s}{Kaan Aslan}\PYG{l+s}{\PYGZdq{}}\PYG{p}{)}\PYG{p}{;}
\PYG{n}{MODULE\PYGZus{}DESCRIPTION}\PYG{p}{(}\PYG{l+s}{\PYGZdq{}}\PYG{l+s}{test\PYGZhy{}driver}\PYG{l+s}{\PYGZdq{}}\PYG{p}{)}\PYG{p}{;}

\PYG{k}{static}\PYG{+w}{ }\PYG{k+kt}{int}\PYG{+w}{ }\PYG{n+nf}{test\PYGZus{}driver\PYGZus{}open}\PYG{p}{(}\PYG{k}{struct}\PYG{+w}{ }\PYG{n+nc}{inode}\PYG{+w}{ }\PYG{o}{*}\PYG{n}{inodep}\PYG{p}{,}\PYG{+w}{ }\PYG{k}{struct}\PYG{+w}{ }\PYG{n+nc}{file}\PYG{+w}{ }\PYG{o}{*}\PYG{n}{filp}\PYG{p}{)}\PYG{p}{;}
\PYG{k}{static}\PYG{+w}{ }\PYG{k+kt}{int}\PYG{+w}{ }\PYG{n+nf}{test\PYGZus{}driver\PYGZus{}release}\PYG{p}{(}\PYG{k}{struct}\PYG{+w}{ }\PYG{n+nc}{inode}\PYG{+w}{ }\PYG{o}{*}\PYG{n}{inodep}\PYG{p}{,}\PYG{+w}{ }\PYG{k}{struct}\PYG{+w}{ }\PYG{n+nc}{file}\PYG{+w}{ }\PYG{o}{*}\PYG{n}{filp}\PYG{p}{)}\PYG{p}{;}
\PYG{k}{static}\PYG{+w}{ }\PYG{k+kt}{ssize\PYGZus{}t}\PYG{+w}{ }\PYG{n+nf}{test\PYGZus{}driver\PYGZus{}read}\PYG{p}{(}\PYG{k}{struct}\PYG{+w}{ }\PYG{n+nc}{file}\PYG{+w}{ }\PYG{o}{*}\PYG{n}{filp}\PYG{p}{,}\PYG{+w}{ }\PYG{k+kt}{char}\PYG{+w}{ }\PYG{o}{*}\PYG{n}{buf}\PYG{p}{,}\PYG{+w}{ }\PYG{k+kt}{size\PYGZus{}t}\PYG{+w}{ }\PYG{n}{size}\PYG{p}{,}\PYG{+w}{ }\PYG{n}{loff\PYGZus{}t}\PYG{+w}{ }\PYG{o}{*}\PYG{n}{off}\PYG{p}{)}\PYG{p}{;}
\PYG{k}{static}\PYG{+w}{ }\PYG{k+kt}{ssize\PYGZus{}t}\PYG{+w}{ }\PYG{n+nf}{test\PYGZus{}driver\PYGZus{}write}\PYG{p}{(}\PYG{k}{struct}\PYG{+w}{ }\PYG{n+nc}{file}\PYG{+w}{ }\PYG{o}{*}\PYG{n}{filp}\PYG{p}{,}\PYG{+w}{ }\PYG{k}{const}\PYG{+w}{ }\PYG{k+kt}{char}\PYG{+w}{ }\PYG{o}{*}\PYG{n}{buf}\PYG{p}{,}\PYG{+w}{ }\PYG{k+kt}{size\PYGZus{}t}\PYG{+w}{ }\PYG{n}{size}\PYG{p}{,}\PYG{+w}{ }\PYG{n}{loff\PYGZus{}t}\PYG{+w}{ }\PYG{o}{*}\PYG{n}{off}\PYG{p}{)}\PYG{p}{;}
\PYG{k}{static}\PYG{+w}{ }\PYG{k+kt}{long}\PYG{+w}{ }\PYG{n+nf}{test\PYGZus{}driver\PYGZus{}ioctl}\PYG{p}{(}\PYG{k}{struct}\PYG{+w}{ }\PYG{n+nc}{file}\PYG{+w}{ }\PYG{o}{*}\PYG{n}{filp}\PYG{p}{,}\PYG{+w}{ }\PYG{k+kt}{unsigned}\PYG{+w}{ }\PYG{k+kt}{int}\PYG{+w}{ }\PYG{n}{cmd}\PYG{p}{,}\PYG{+w}{ }\PYG{k+kt}{unsigned}\PYG{+w}{ }\PYG{k+kt}{long}\PYG{+w}{ }\PYG{n}{arg}\PYG{p}{)}\PYG{p}{;}
\PYG{k}{static}\PYG{+w}{ }\PYG{k+kt}{long}\PYG{+w}{ }\PYG{n+nf}{ioctl\PYGZus{}test}\PYG{p}{(}\PYG{k}{struct}\PYG{+w}{ }\PYG{n+nc}{file}\PYG{+w}{ }\PYG{o}{*}\PYG{n}{filp}\PYG{p}{,}\PYG{+w}{ }\PYG{k+kt}{unsigned}\PYG{+w}{ }\PYG{k+kt}{long}\PYG{+w}{ }\PYG{n}{arg}\PYG{p}{)}\PYG{p}{;}

\PYG{k}{static}\PYG{+w}{ }\PYG{k+kt}{dev\PYGZus{}t}\PYG{+w}{ }\PYG{n}{g\PYGZus{}dev}\PYG{p}{;}
\PYG{k}{static}\PYG{+w}{ }\PYG{k}{struct}\PYG{+w}{ }\PYG{n+nc}{cdev}\PYG{+w}{ }\PYG{n}{g\PYGZus{}cdev}\PYG{p}{;}
\PYG{k}{static}\PYG{+w}{ }\PYG{k}{struct}\PYG{+w}{ }\PYG{n+nc}{file\PYGZus{}operations}\PYG{+w}{ }\PYG{n}{g\PYGZus{}fops}\PYG{+w}{ }\PYG{o}{=}\PYG{+w}{ }\PYG{p}{\PYGZob{}}
\PYG{+w}{    }\PYG{p}{.}\PYG{n}{owner}\PYG{+w}{          }\PYG{o}{=}\PYG{+w}{ }\PYG{n}{THIS\PYGZus{}MODULE}\PYG{p}{,}
\PYG{+w}{    }\PYG{p}{.}\PYG{n}{open}\PYG{+w}{           }\PYG{o}{=}\PYG{+w}{ }\PYG{n}{test\PYGZus{}driver\PYGZus{}open}\PYG{p}{,}
\PYG{+w}{    }\PYG{p}{.}\PYG{n}{read}\PYG{+w}{           }\PYG{o}{=}\PYG{+w}{ }\PYG{n}{test\PYGZus{}driver\PYGZus{}read}\PYG{p}{,}
\PYG{+w}{    }\PYG{p}{.}\PYG{n}{write}\PYG{+w}{          }\PYG{o}{=}\PYG{+w}{ }\PYG{n}{test\PYGZus{}driver\PYGZus{}write}\PYG{p}{,}
\PYG{+w}{    }\PYG{p}{.}\PYG{n}{release}\PYG{+w}{        }\PYG{o}{=}\PYG{+w}{ }\PYG{n}{test\PYGZus{}driver\PYGZus{}release}\PYG{p}{,}
\PYG{+w}{    }\PYG{p}{.}\PYG{n}{unlocked\PYGZus{}ioctl}\PYG{+w}{ }\PYG{o}{=}\PYG{+w}{ }\PYG{n}{test\PYGZus{}driver\PYGZus{}ioctl}
\PYG{p}{\PYGZcb{}}\PYG{p}{;}

\PYG{k}{static}\PYG{+w}{ }\PYG{k+kt}{int}\PYG{+w}{ }\PYG{n}{\PYGZus{}\PYGZus{}init}\PYG{+w}{ }\PYG{n+nf}{test\PYGZus{}driver\PYGZus{}init}\PYG{p}{(}\PYG{k+kt}{void}\PYG{p}{)}
\PYG{p}{\PYGZob{}}
\PYG{+w}{    }\PYG{k+kt}{int}\PYG{+w}{ }\PYG{n}{result}\PYG{p}{;}

\PYG{+w}{    }\PYG{n}{printk}\PYG{p}{(}\PYG{n}{KERN\PYGZus{}INFO}\PYG{+w}{ }\PYG{l+s}{\PYGZdq{}}\PYG{l+s}{test\PYGZhy{}driver module initialization...}\PYG{l+s+se}{\PYGZbs{}n}\PYG{l+s}{\PYGZdq{}}\PYG{p}{)}\PYG{p}{;}

\PYG{+w}{    }\PYG{k}{if}\PYG{+w}{ }\PYG{p}{(}\PYG{p}{(}\PYG{n}{result}\PYG{+w}{ }\PYG{o}{=}\PYG{+w}{ }\PYG{n}{alloc\PYGZus{}chrdev\PYGZus{}region}\PYG{p}{(}\PYG{o}{\PYGZam{}}\PYG{n}{g\PYGZus{}dev}\PYG{p}{,}\PYG{+w}{ }\PYG{l+m+mi}{0}\PYG{p}{,}\PYG{+w}{ }\PYG{l+m+mi}{1}\PYG{p}{,}\PYG{+w}{ }\PYG{l+s}{\PYGZdq{}}\PYG{l+s}{test\PYGZhy{}driver}\PYG{l+s}{\PYGZdq{}}\PYG{p}{)}\PYG{p}{)}\PYG{+w}{ }\PYG{o}{\PYGZlt{}}\PYG{+w}{ }\PYG{l+m+mi}{0}\PYG{p}{)}\PYG{+w}{ }\PYG{p}{\PYGZob{}}
\PYG{+w}{        }\PYG{n}{printk}\PYG{p}{(}\PYG{n}{KERN\PYGZus{}INFO}\PYG{+w}{ }\PYG{l+s}{\PYGZdq{}}\PYG{l+s}{cannot alloc char driver!...}\PYG{l+s+se}{\PYGZbs{}n}\PYG{l+s}{\PYGZdq{}}\PYG{p}{)}\PYG{p}{;}
\PYG{+w}{        }\PYG{k}{return}\PYG{+w}{ }\PYG{n}{result}\PYG{p}{;}
\PYG{+w}{    }\PYG{p}{\PYGZcb{}}
\PYG{+w}{    }\PYG{n}{cdev\PYGZus{}init}\PYG{p}{(}\PYG{o}{\PYGZam{}}\PYG{n}{g\PYGZus{}cdev}\PYG{p}{,}\PYG{+w}{ }\PYG{o}{\PYGZam{}}\PYG{n}{g\PYGZus{}fops}\PYG{p}{)}\PYG{p}{;}
\PYG{+w}{    }\PYG{k}{if}\PYG{+w}{ }\PYG{p}{(}\PYG{p}{(}\PYG{n}{result}\PYG{+w}{ }\PYG{o}{=}\PYG{+w}{ }\PYG{n}{cdev\PYGZus{}add}\PYG{p}{(}\PYG{o}{\PYGZam{}}\PYG{n}{g\PYGZus{}cdev}\PYG{p}{,}\PYG{+w}{ }\PYG{n}{g\PYGZus{}dev}\PYG{p}{,}\PYG{+w}{ }\PYG{l+m+mi}{1}\PYG{p}{)}\PYG{p}{)}\PYG{+w}{ }\PYG{o}{\PYGZlt{}}\PYG{+w}{ }\PYG{l+m+mi}{0}\PYG{p}{)}\PYG{+w}{ }\PYG{p}{\PYGZob{}}
\PYG{+w}{        }\PYG{n}{unregister\PYGZus{}chrdev\PYGZus{}region}\PYG{p}{(}\PYG{n}{g\PYGZus{}dev}\PYG{p}{,}\PYG{+w}{ }\PYG{l+m+mi}{1}\PYG{p}{)}\PYG{p}{;}
\PYG{+w}{        }\PYG{n}{printk}\PYG{p}{(}\PYG{n}{KERN\PYGZus{}ERR}\PYG{+w}{ }\PYG{l+s}{\PYGZdq{}}\PYG{l+s}{cannot add device!...}\PYG{l+s+se}{\PYGZbs{}n}\PYG{l+s}{\PYGZdq{}}\PYG{p}{)}\PYG{p}{;}
\PYG{+w}{        }\PYG{k}{return}\PYG{+w}{ }\PYG{n}{result}\PYG{p}{;}
\PYG{+w}{    }\PYG{p}{\PYGZcb{}}

\PYG{+w}{    }\PYG{k}{return}\PYG{+w}{ }\PYG{l+m+mi}{0}\PYG{p}{;}
\PYG{p}{\PYGZcb{}}

\PYG{k}{static}\PYG{+w}{ }\PYG{k+kt}{void}\PYG{+w}{ }\PYG{n}{\PYGZus{}\PYGZus{}exit}\PYG{+w}{ }\PYG{n+nf}{test\PYGZus{}driver\PYGZus{}exit}\PYG{p}{(}\PYG{k+kt}{void}\PYG{p}{)}
\PYG{p}{\PYGZob{}}
\PYG{+w}{    }\PYG{n}{cdev\PYGZus{}del}\PYG{p}{(}\PYG{o}{\PYGZam{}}\PYG{n}{g\PYGZus{}cdev}\PYG{p}{)}\PYG{p}{;}
\PYG{+w}{    }\PYG{n}{unregister\PYGZus{}chrdev\PYGZus{}region}\PYG{p}{(}\PYG{n}{g\PYGZus{}dev}\PYG{p}{,}\PYG{+w}{ }\PYG{l+m+mi}{1}\PYG{p}{)}\PYG{p}{;}
\PYG{+w}{    }\PYG{n}{printk}\PYG{p}{(}\PYG{n}{KERN\PYGZus{}INFO}\PYG{+w}{ }\PYG{l+s}{\PYGZdq{}}\PYG{l+s}{test\PYGZhy{}driver module exit...}\PYG{l+s+se}{\PYGZbs{}n}\PYG{l+s}{\PYGZdq{}}\PYG{p}{)}\PYG{p}{;}
\PYG{p}{\PYGZcb{}}

\PYG{k}{static}\PYG{+w}{ }\PYG{k+kt}{int}\PYG{+w}{ }\PYG{n+nf}{test\PYGZus{}driver\PYGZus{}open}\PYG{p}{(}\PYG{k}{struct}\PYG{+w}{ }\PYG{n+nc}{inode}\PYG{+w}{ }\PYG{o}{*}\PYG{n}{inodep}\PYG{p}{,}\PYG{+w}{ }\PYG{k}{struct}\PYG{+w}{ }\PYG{n+nc}{file}\PYG{+w}{ }\PYG{o}{*}\PYG{n}{filp}\PYG{p}{)}
\PYG{p}{\PYGZob{}}
\PYG{+w}{    }\PYG{n}{printk}\PYG{p}{(}\PYG{n}{KERN\PYGZus{}INFO}\PYG{+w}{ }\PYG{l+s}{\PYGZdq{}}\PYG{l+s}{test\PYGZhy{}driver opened...}\PYG{l+s+se}{\PYGZbs{}n}\PYG{l+s}{\PYGZdq{}}\PYG{p}{)}\PYG{p}{;}

\PYG{+w}{    }\PYG{k}{return}\PYG{+w}{ }\PYG{l+m+mi}{0}\PYG{p}{;}
\PYG{p}{\PYGZcb{}}

\PYG{k}{static}\PYG{+w}{ }\PYG{k+kt}{int}\PYG{+w}{ }\PYG{n+nf}{test\PYGZus{}driver\PYGZus{}release}\PYG{p}{(}\PYG{k}{struct}\PYG{+w}{ }\PYG{n+nc}{inode}\PYG{+w}{ }\PYG{o}{*}\PYG{n}{inodep}\PYG{p}{,}\PYG{+w}{ }\PYG{k}{struct}\PYG{+w}{ }\PYG{n+nc}{file}\PYG{+w}{ }\PYG{o}{*}\PYG{n}{filp}\PYG{p}{)}
\PYG{p}{\PYGZob{}}
\PYG{+w}{    }\PYG{n}{printk}\PYG{p}{(}\PYG{n}{KERN\PYGZus{}INFO}\PYG{+w}{ }\PYG{l+s}{\PYGZdq{}}\PYG{l+s}{test\PYGZhy{}driver closed...}\PYG{l+s+se}{\PYGZbs{}n}\PYG{l+s}{\PYGZdq{}}\PYG{p}{)}\PYG{p}{;}

\PYG{+w}{    }\PYG{k}{return}\PYG{+w}{ }\PYG{l+m+mi}{0}\PYG{p}{;}
\PYG{p}{\PYGZcb{}}

\PYG{k}{static}\PYG{+w}{ }\PYG{k+kt}{ssize\PYGZus{}t}\PYG{+w}{ }\PYG{n+nf}{test\PYGZus{}driver\PYGZus{}read}\PYG{p}{(}\PYG{k}{struct}\PYG{+w}{ }\PYG{n+nc}{file}\PYG{+w}{ }\PYG{o}{*}\PYG{n}{filp}\PYG{p}{,}\PYG{+w}{ }\PYG{k+kt}{char}\PYG{+w}{ }\PYG{o}{*}\PYG{n}{buf}\PYG{p}{,}\PYG{+w}{ }\PYG{k+kt}{size\PYGZus{}t}\PYG{+w}{ }\PYG{n}{size}\PYG{p}{,}\PYG{+w}{ }\PYG{n}{loff\PYGZus{}t}\PYG{+w}{ }\PYG{o}{*}\PYG{n}{off}\PYG{p}{)}
\PYG{p}{\PYGZob{}}
\PYG{+w}{    }\PYG{n}{printk}\PYG{p}{(}\PYG{n}{KERN\PYGZus{}INFO}\PYG{+w}{ }\PYG{l+s}{\PYGZdq{}}\PYG{l+s}{test\PYGZhy{}driver read...}\PYG{l+s+se}{\PYGZbs{}n}\PYG{l+s}{\PYGZdq{}}\PYG{p}{)}\PYG{p}{;}

\PYG{+w}{    }\PYG{k}{return}\PYG{+w}{ }\PYG{l+m+mi}{0}\PYG{p}{;}
\PYG{p}{\PYGZcb{}}

\PYG{k}{static}\PYG{+w}{ }\PYG{k+kt}{ssize\PYGZus{}t}\PYG{+w}{ }\PYG{n+nf}{test\PYGZus{}driver\PYGZus{}write}\PYG{p}{(}\PYG{k}{struct}\PYG{+w}{ }\PYG{n+nc}{file}\PYG{+w}{ }\PYG{o}{*}\PYG{n}{filp}\PYG{p}{,}\PYG{+w}{ }\PYG{k}{const}\PYG{+w}{ }\PYG{k+kt}{char}\PYG{+w}{ }\PYG{o}{*}\PYG{n}{buf}\PYG{p}{,}\PYG{+w}{ }\PYG{k+kt}{size\PYGZus{}t}\PYG{+w}{ }\PYG{n}{size}\PYG{p}{,}\PYG{+w}{ }\PYG{n}{loff\PYGZus{}t}\PYG{+w}{ }\PYG{o}{*}\PYG{n}{off}\PYG{p}{)}
\PYG{p}{\PYGZob{}}
\PYG{+w}{    }\PYG{n}{printk}\PYG{p}{(}\PYG{n}{KERN\PYGZus{}INFO}\PYG{+w}{ }\PYG{l+s}{\PYGZdq{}}\PYG{l+s}{test\PYGZhy{}driver write...}\PYG{l+s+se}{\PYGZbs{}n}\PYG{l+s}{\PYGZdq{}}\PYG{p}{)}\PYG{p}{;}

\PYG{+w}{    }\PYG{k}{return}\PYG{+w}{ }\PYG{l+m+mi}{0}\PYG{p}{;}
\PYG{p}{\PYGZcb{}}

\PYG{k}{static}\PYG{+w}{ }\PYG{k+kt}{long}\PYG{+w}{ }\PYG{n+nf}{test\PYGZus{}driver\PYGZus{}ioctl}\PYG{p}{(}\PYG{k}{struct}\PYG{+w}{ }\PYG{n+nc}{file}\PYG{+w}{ }\PYG{o}{*}\PYG{n}{filp}\PYG{p}{,}\PYG{+w}{ }\PYG{k+kt}{unsigned}\PYG{+w}{ }\PYG{k+kt}{int}\PYG{+w}{ }\PYG{n}{cmd}\PYG{p}{,}\PYG{+w}{ }\PYG{k+kt}{unsigned}\PYG{+w}{ }\PYG{k+kt}{long}\PYG{+w}{ }\PYG{n}{arg}\PYG{p}{)}
\PYG{p}{\PYGZob{}}
\PYG{+w}{    }\PYG{k+kt}{long}\PYG{+w}{ }\PYG{n}{result}\PYG{p}{;}

\PYG{+w}{    }\PYG{n}{printk}\PYG{p}{(}\PYG{n}{KERN\PYGZus{}INFO}\PYG{+w}{ }\PYG{l+s}{\PYGZdq{}}\PYG{l+s}{test\PYGZus{}driver\PYGZus{}ioctl...}\PYG{l+s+se}{\PYGZbs{}n}\PYG{l+s}{\PYGZdq{}}\PYG{p}{)}\PYG{p}{;}

\PYG{+w}{    }\PYG{k}{switch}\PYG{+w}{ }\PYG{p}{(}\PYG{n}{cmd}\PYG{p}{)}\PYG{+w}{ }\PYG{p}{\PYGZob{}}
\PYG{+w}{        }\PYG{k}{case}\PYG{+w}{ }\PYG{n+no}{IOC\PYGZus{}TEST}\PYG{p}{:}
\PYG{+w}{            }\PYG{n}{result}\PYG{+w}{ }\PYG{o}{=}\PYG{+w}{ }\PYG{n}{ioctl\PYGZus{}test}\PYG{p}{(}\PYG{n}{filp}\PYG{p}{,}\PYG{+w}{ }\PYG{n}{arg}\PYG{p}{)}\PYG{p}{;}
\PYG{+w}{            }\PYG{k}{break}\PYG{p}{;}
\PYG{+w}{        }\PYG{k}{default}\PYG{o}{:}
\PYG{+w}{            }\PYG{n}{result}\PYG{+w}{ }\PYG{o}{=}\PYG{+w}{ }\PYG{o}{\PYGZhy{}}\PYG{n}{ENOTTY}\PYG{p}{;}
\PYG{+w}{    }\PYG{p}{\PYGZcb{}}

\PYG{+w}{    }\PYG{k}{return}\PYG{+w}{ }\PYG{n}{result}\PYG{p}{;}
\PYG{p}{\PYGZcb{}}

\PYG{k+kt}{long}\PYG{+w}{ }\PYG{n+nf}{ioctl\PYGZus{}test}\PYG{p}{(}\PYG{k}{struct}\PYG{+w}{ }\PYG{n+nc}{file}\PYG{+w}{ }\PYG{o}{*}\PYG{n}{filp}\PYG{p}{,}\PYG{+w}{ }\PYG{k+kt}{unsigned}\PYG{+w}{ }\PYG{k+kt}{long}\PYG{+w}{ }\PYG{n}{arg}\PYG{p}{)}
\PYG{p}{\PYGZob{}}
\PYG{+w}{    }\PYG{n}{printk}\PYG{p}{(}\PYG{n}{KERN\PYGZus{}INFO}\PYG{+w}{ }\PYG{l+s}{\PYGZdq{}}\PYG{l+s}{IOC\PYGZus{}test\PYGZus{}driver\PYGZus{}TEST1}\PYG{l+s+se}{\PYGZbs{}n}\PYG{l+s}{\PYGZdq{}}\PYG{p}{)}\PYG{p}{;}

\PYG{+w}{    }\PYG{k}{return}\PYG{+w}{ }\PYG{l+m+mi}{0}\PYG{p}{;}
\PYG{p}{\PYGZcb{}}

\PYG{n}{module\PYGZus{}init}\PYG{p}{(}\PYG{n}{test\PYGZus{}driver\PYGZus{}init}\PYG{p}{)}\PYG{p}{;}
\PYG{n}{module\PYGZus{}exit}\PYG{p}{(}\PYG{n}{test\PYGZus{}driver\PYGZus{}exit}\PYG{p}{)}\PYG{p}{;}
\end{sphinxVerbatim}

\sphinxAtStartPar
Aygıt sürücüyü şöyle derleyebilirsiniz:

\begin{sphinxVerbatim}[commandchars=\\\{\}]
\PYGZdl{}\PYG{+w}{ }make\PYG{+w}{ }\PYG{n+nv}{file}\PYG{o}{=}test\PYGZhy{}driver
\PYGZdl{}\PYG{+w}{ }sudo\PYG{+w}{ }./load\PYG{+w}{ }test\PYGZhy{}driver
\end{sphinxVerbatim}

\sphinxAtStartPar
Aşağıdaki kullanıcı modu programı aygıt sürücüyü açıp ioctl kodunu çalıştırmaktadır:

\sphinxAtStartPar
\sphinxstylestrong{app.c:}

\begin{sphinxVerbatim}[commandchars=\\\{\}]
\PYG{c+cp}{\PYGZsh{}}\PYG{c+cp}{include}\PYG{+w}{ }\PYG{c+cpf}{\PYGZlt{}stdio.h\PYGZgt{}}
\PYG{c+cp}{\PYGZsh{}}\PYG{c+cp}{include}\PYG{+w}{ }\PYG{c+cpf}{\PYGZlt{}stdlib.h\PYGZgt{}}
\PYG{c+cp}{\PYGZsh{}}\PYG{c+cp}{include}\PYG{+w}{ }\PYG{c+cpf}{\PYGZlt{}fcntl.h\PYGZgt{}}
\PYG{c+cp}{\PYGZsh{}}\PYG{c+cp}{include}\PYG{+w}{ }\PYG{c+cpf}{\PYGZlt{}unistd.h\PYGZgt{}}
\PYG{c+cp}{\PYGZsh{}}\PYG{c+cp}{include}\PYG{+w}{ }\PYG{c+cpf}{\PYGZlt{}sys/ioctl.h\PYGZgt{}}
\PYG{c+cp}{\PYGZsh{}}\PYG{c+cp}{include}\PYG{+w}{ }\PYG{c+cpf}{\PYGZdq{}test\PYGZhy{}driver.h\PYGZdq{}}

\PYG{k+kt}{void}\PYG{+w}{ }\PYG{n+nf}{exit\PYGZus{}sys}\PYG{p}{(}\PYG{k}{const}\PYG{+w}{ }\PYG{k+kt}{char}\PYG{+w}{ }\PYG{o}{*}\PYG{n}{msg}\PYG{p}{)}\PYG{p}{;}

\PYG{k+kt}{int}\PYG{+w}{ }\PYG{n+nf}{main}\PYG{p}{(}\PYG{k+kt}{void}\PYG{p}{)}
\PYG{p}{\PYGZob{}}
\PYG{+w}{    }\PYG{k+kt}{int}\PYG{+w}{ }\PYG{n}{fd}\PYG{p}{;}

\PYG{+w}{    }\PYG{k}{if}\PYG{+w}{ }\PYG{p}{(}\PYG{p}{(}\PYG{n}{fd}\PYG{+w}{ }\PYG{o}{=}\PYG{+w}{ }\PYG{n}{open}\PYG{p}{(}\PYG{l+s}{\PYGZdq{}}\PYG{l+s}{test\PYGZhy{}driver}\PYG{l+s}{\PYGZdq{}}\PYG{p}{,}\PYG{+w}{ }\PYG{n}{O\PYGZus{}RDONLY}\PYG{p}{)}\PYG{p}{)}\PYG{+w}{ }\PYG{o}{=}\PYG{o}{=}\PYG{+w}{ }\PYG{l+m+mi}{\PYGZhy{}1}\PYG{p}{)}
\PYG{+w}{        }\PYG{n}{exit\PYGZus{}sys}\PYG{p}{(}\PYG{l+s}{\PYGZdq{}}\PYG{l+s}{open}\PYG{l+s}{\PYGZdq{}}\PYG{p}{)}\PYG{p}{;}

\PYG{+w}{    }\PYG{k}{if}\PYG{+w}{ }\PYG{p}{(}\PYG{n}{ioctl}\PYG{p}{(}\PYG{n}{fd}\PYG{p}{,}\PYG{+w}{ }\PYG{n}{IOC\PYGZus{}TEST}\PYG{p}{)}\PYG{+w}{ }\PYG{o}{=}\PYG{o}{=}\PYG{+w}{ }\PYG{l+m+mi}{\PYGZhy{}1}\PYG{p}{)}
\PYG{+w}{        }\PYG{n}{exit\PYGZus{}sys}\PYG{p}{(}\PYG{l+s}{\PYGZdq{}}\PYG{l+s}{ioctl}\PYG{l+s}{\PYGZdq{}}\PYG{p}{)}\PYG{p}{;}

\PYG{+w}{    }\PYG{n}{close}\PYG{p}{(}\PYG{n}{fd}\PYG{p}{)}\PYG{p}{;}

\PYG{+w}{    }\PYG{k}{return}\PYG{+w}{ }\PYG{l+m+mi}{0}\PYG{p}{;}
\PYG{p}{\PYGZcb{}}

\PYG{k+kt}{void}\PYG{+w}{ }\PYG{n+nf}{exit\PYGZus{}sys}\PYG{p}{(}\PYG{k}{const}\PYG{+w}{ }\PYG{k+kt}{char}\PYG{+w}{ }\PYG{o}{*}\PYG{n}{msg}\PYG{p}{)}
\PYG{p}{\PYGZob{}}
\PYG{+w}{    }\PYG{n}{perror}\PYG{p}{(}\PYG{n}{msg}\PYG{p}{)}\PYG{p}{;}
\PYG{+w}{    }\PYG{n}{exit}\PYG{p}{(}\PYG{n}{EXIT\PYGZus{}FAILURE}\PYG{p}{)}\PYG{p}{;}
\PYG{p}{\PYGZcb{}}
\end{sphinxVerbatim}

\begin{sphinxVerbatim}[commandchars=\\\{\}]
\PYGZdl{}\PYG{+w}{ }gcc\PYG{+w}{ }\PYGZhy{}Wall\PYG{+w}{ }\PYGZhy{}o\PYG{+w}{ }app\PYG{+w}{ }app.c
\PYGZdl{}\PYG{+w}{ }./app
\end{sphinxVerbatim}

\sphinxAtStartPar
\sphinxstyleemphasis{app} programını çalıştırdıktan sonra \sphinxstyleemphasis{dmesg} komutunu kullandığımızda log mesajları aşağıdaki gibi
görülmelidir:

\begin{sphinxVerbatim}[commandchars=\\\{\}]
[  431.787375] test\PYGZhy{}driver module initialization...
[  450.631940] test\PYGZhy{}driver opened...
[  450.631953] test\PYGZus{}driver\PYGZus{}ioctl...
[  450.631958] IOC\PYGZus{}test\PYGZus{}driver\PYGZus{}TEST1
[  450.631967] test\PYGZhy{}driver closed...
\end{sphinxVerbatim}


\subsection{RCU Mekanizması ile \sphinxstyleliteralintitle{\sphinxupquote{task\_struct}} Listelerinin Dolaşılması}
\label{\detokenize{processcontrolblock:rcu-mekanizmasi-ile-task-struct-listelerinin-dolasilmasi}}
\sphinxAtStartPar
Biz \sphinxcode{\sphinxupquote{task\_struct}} listelerini dolaşırken çekirdek bu listelere eleman eklerse ya da bu listelerden
eleman silerse bizim dolaşım yapan kodumuzda \sphinxstyleemphasis{tanımsız davranışlar (undefined behaviours)} oluşabilir.
Bu da tüm sistemin çökmesine yol açabilir. Artık belli bir süredir çekirdek \sphinxcode{\sphinxupquote{task\_struct}} listeleri
üzerinde kilitsiz bir biçimde \sphinxstyleemphasis{RCU (Read\sphinxhyphen{}Copy\sphinxhyphen{}Update)} mekanizmasıyla işlemler yapmaktadır. Dolayısıyla
bizim de kendimizi çekirdeğin kullandığı bu RCU mekanizmasına uydurarak dolaşım yapmamız gerekir.

\sphinxAtStartPar
RCU mekanizmasının ayrıntılarını ileride ayrı başlıkta ele alacağız. Ancak bu noktada RCU mekanizması
ile \sphinxcode{\sphinxupquote{task\_struct}} listelerinin dolaşılması için bazı temel bilgileri de vereceğiz.

\sphinxAtStartPar
RCU mekanizmasıyla bağlı listeleri dolaşırken işlemin başında \sphinxstyleemphasis{rcu\_read\_lock} fonksiyonunun, işlemin
sonunda da \sphinxstyleemphasis{rcu\_read\_unlock} fonksiyonunun çağrılması gerekmektedir:

\begin{sphinxVerbatim}[commandchars=\\\{\}]
\PYG{c+cp}{\PYGZsh{}}\PYG{c+cp}{include}\PYG{+w}{ }\PYG{c+cpf}{\PYGZlt{}linux/rcupdate.h\PYGZgt{}}

\PYG{k+kt}{void}\PYG{+w}{ }\PYG{n+nf}{rcu\PYGZus{}read\PYGZus{}lock}\PYG{p}{(}\PYG{k+kt}{void}\PYG{p}{)}\PYG{p}{;}
\PYG{k+kt}{void}\PYG{+w}{ }\PYG{n+nf}{rcu\PYGZus{}read\PYGZus{}unlock}\PYG{p}{(}\PYG{k+kt}{void}\PYG{p}{)}\PYG{p}{;}
\end{sphinxVerbatim}

\sphinxAtStartPar
Bizim dolaşım kodunu bu iki çağrı arasına yerleştirmemiz gerekir. RCU mekanizması altında bağlı listelerin
bağlarına erişilirken \sphinxcode{\sphinxupquote{rcu\_dereference}} makrosu kullanılmalıdır. Örneğin sistemdeki bütün proseslerin
ana thread’lerine ilişkin \sphinxcode{\sphinxupquote{task\_struct}} nesneleri aşağıdaki gibi güvenli bir biçimde dolaşılabilir:

\begin{sphinxVerbatim}[commandchars=\\\{\}]
\PYG{k}{static}\PYG{+w}{ }\PYG{k+kt}{void}\PYG{+w}{ }\PYG{n+nf}{walk\PYGZus{}processes}\PYG{p}{(}\PYG{k+kt}{void}\PYG{p}{)}
\PYG{p}{\PYGZob{}}
\PYG{+w}{    }\PYG{k}{struct}\PYG{+w}{ }\PYG{n+nc}{task\PYGZus{}struct}\PYG{+w}{ }\PYG{o}{*}\PYG{n}{ts}\PYG{p}{;}
\PYG{+w}{    }\PYG{k}{struct}\PYG{+w}{ }\PYG{n+nc}{list\PYGZus{}head}\PYG{+w}{ }\PYG{o}{*}\PYG{n}{lh}\PYG{p}{;}

\PYG{+w}{    }\PYG{n}{rcu\PYGZus{}read\PYGZus{}lock}\PYG{p}{(}\PYG{p}{)}\PYG{p}{;}

\PYG{+w}{    }\PYG{n}{lh}\PYG{+w}{ }\PYG{o}{=}\PYG{+w}{ }\PYG{n}{rcu\PYGZus{}dereference}\PYG{p}{(}\PYG{n}{init\PYGZus{}task}\PYG{p}{.}\PYG{n}{tasks}\PYG{p}{.}\PYG{n}{next}\PYG{p}{)}\PYG{p}{;}
\PYG{+w}{    }\PYG{k}{while}\PYG{+w}{ }\PYG{p}{(}\PYG{n}{lh}\PYG{+w}{ }\PYG{o}{!}\PYG{o}{=}\PYG{+w}{ }\PYG{o}{\PYGZam{}}\PYG{n}{init\PYGZus{}task}\PYG{p}{.}\PYG{n}{tasks}\PYG{p}{)}\PYG{+w}{ }\PYG{p}{\PYGZob{}}
\PYG{+w}{        }\PYG{n}{ts}\PYG{+w}{ }\PYG{o}{=}\PYG{+w}{ }\PYG{n}{container\PYGZus{}of}\PYG{p}{(}\PYG{n}{lh}\PYG{p}{,}\PYG{+w}{ }\PYG{k}{struct}\PYG{+w}{ }\PYG{n+nc}{task\PYGZus{}struct}\PYG{p}{,}\PYG{+w}{ }\PYG{n}{tasks}\PYG{p}{)}\PYG{p}{;}
\PYG{+w}{        }\PYG{n}{printk}\PYG{p}{(}\PYG{n}{KERN\PYGZus{}INFO}\PYG{+w}{ }\PYG{l+s}{\PYGZdq{}}\PYG{l+s}{PID = \PYGZpc{}d, COMM = \PYGZpc{}s}\PYG{l+s}{\PYGZdq{}}\PYG{p}{,}\PYG{+w}{ }\PYG{n}{ts}\PYG{o}{\PYGZhy{}}\PYG{o}{\PYGZgt{}}\PYG{n}{pid}\PYG{p}{,}\PYG{+w}{ }\PYG{n}{ts}\PYG{o}{\PYGZhy{}}\PYG{o}{\PYGZgt{}}\PYG{n}{comm}\PYG{p}{)}\PYG{p}{;}
\PYG{+w}{        }\PYG{n}{lh}\PYG{+w}{ }\PYG{o}{=}\PYG{+w}{ }\PYG{n}{rcu\PYGZus{}dereference}\PYG{p}{(}\PYG{n}{ts}\PYG{o}{\PYGZhy{}}\PYG{o}{\PYGZgt{}}\PYG{n}{tasks}\PYG{p}{.}\PYG{n}{next}\PYG{p}{)}\PYG{p}{;}
\PYG{+w}{    }\PYG{p}{\PYGZcb{}}

\PYG{+w}{    }\PYG{n}{rcu\PYGZus{}read\PYGZus{}unlock}\PYG{p}{(}\PYG{p}{)}\PYG{p}{;}
\PYG{p}{\PYGZcb{}}
\end{sphinxVerbatim}

\sphinxAtStartPar
\sphinxcode{\sphinxupquote{container\_of}} makrosu ile \sphinxcode{\sphinxupquote{list\_entry}} makrosunun tamamen eşdeğer olduğunu anımsayınız.
\sphinxcode{\sphinxupquote{task\_struct}} içerisindeki \sphinxcode{\sphinxupquote{comm}} elemanında en fazla 16 karakterlik \sphinxcode{\sphinxupquote{task\_struct}} nesnesini
betimleyen bir yazı bulunmaktadır. \sphinxstyleemphasis{fork} işlemi ya da thread yaratımı sırasında bu yazı üst prosesin
\sphinxcode{\sphinxupquote{comm}} elemanından kopyalanmaktadır. \sphinxstyleemphasis{exec} işlemleri sırasında ise bu yazı çalıştırılan programın
dosya isminden hareketle değiştirilmektedir.

\sphinxAtStartPar
\sphinxcode{\sphinxupquote{\textless{}linux/sched/signal.h\textgreater{}}} dosyası içerisinde \sphinxcode{\sphinxupquote{next\_task}} isimli bir makro da vardır. Bu makro bir
\sphinxcode{\sphinxupquote{task\_struct}} nesnesinin adresini parametre olarak alıp RCU mekanizmasına uygun olarak onun
\sphinxcode{\sphinxupquote{tasks.next}} göstericisinin gösterdiği yerdeki \sphinxcode{\sphinxupquote{task\_struct}} nesnesinin adresini vermektedir:

\begin{sphinxVerbatim}[commandchars=\\\{\}]
\PYG{c+cp}{\PYGZsh{}}\PYG{c+cp}{include}\PYG{+w}{ }\PYG{c+cpf}{\PYGZlt{}linux/sched/signal.h\PYGZgt{}}

\PYG{n}{next\PYGZus{}task}\PYG{p}{(}\PYG{n}{p}\PYG{p}{)}
\end{sphinxVerbatim}

\sphinxAtStartPar
Bu makro kullanılarak \sphinxstyleemphasis{walk\_process} fonksiyonu şöyle de yazılabilir:

\begin{sphinxVerbatim}[commandchars=\\\{\}]
\PYG{k}{static}\PYG{+w}{ }\PYG{k+kt}{void}\PYG{+w}{ }\PYG{n+nf}{walk\PYGZus{}processes}\PYG{p}{(}\PYG{k+kt}{void}\PYG{p}{)}
\PYG{p}{\PYGZob{}}
\PYG{+w}{    }\PYG{k}{struct}\PYG{+w}{ }\PYG{n+nc}{task\PYGZus{}struct}\PYG{+w}{ }\PYG{o}{*}\PYG{n}{ts}\PYG{p}{;}

\PYG{+w}{    }\PYG{n}{rcu\PYGZus{}read\PYGZus{}lock}\PYG{p}{(}\PYG{p}{)}\PYG{p}{;}

\PYG{+w}{    }\PYG{n}{ts}\PYG{+w}{ }\PYG{o}{=}\PYG{+w}{ }\PYG{n}{next\PYGZus{}task}\PYG{p}{(}\PYG{o}{\PYGZam{}}\PYG{n}{init\PYGZus{}task}\PYG{p}{)}\PYG{p}{;}
\PYG{+w}{    }\PYG{k}{while}\PYG{+w}{ }\PYG{p}{(}\PYG{n}{ts}\PYG{+w}{ }\PYG{o}{!}\PYG{o}{=}\PYG{+w}{ }\PYG{o}{\PYGZam{}}\PYG{n}{init\PYGZus{}task}\PYG{p}{)}\PYG{+w}{ }\PYG{p}{\PYGZob{}}
\PYG{+w}{        }\PYG{n}{printk}\PYG{p}{(}\PYG{n}{KERN\PYGZus{}INFO}\PYG{+w}{ }\PYG{l+s}{\PYGZdq{}}\PYG{l+s}{PID = \PYGZpc{}d, COMM = \PYGZpc{}s}\PYG{l+s}{\PYGZdq{}}\PYG{p}{,}\PYG{+w}{ }\PYG{n}{ts}\PYG{o}{\PYGZhy{}}\PYG{o}{\PYGZgt{}}\PYG{n}{pid}\PYG{p}{,}\PYG{+w}{ }\PYG{n}{ts}\PYG{o}{\PYGZhy{}}\PYG{o}{\PYGZgt{}}\PYG{n}{comm}\PYG{p}{)}\PYG{p}{;}
\PYG{+w}{        }\PYG{n}{ts}\PYG{+w}{ }\PYG{o}{=}\PYG{+w}{ }\PYG{n}{next\PYGZus{}task}\PYG{p}{(}\PYG{n}{ts}\PYG{p}{)}\PYG{p}{;}
\PYG{+w}{    }\PYG{p}{\PYGZcb{}}

\PYG{+w}{    }\PYG{n}{rcu\PYGZus{}read\PYGZus{}unlock}\PYG{p}{(}\PYG{p}{)}\PYG{p}{;}
\PYG{p}{\PYGZcb{}}
\end{sphinxVerbatim}


\subsection{Tüm Prosesleri Dolaşan Çekirdek Modülü}
\label{\detokenize{processcontrolblock:tum-prosesleri-dolasan-cekirdek-modulu}}
\sphinxAtStartPar
Aşağıda tüm proseslerin ana thread’lerini dolaşan örnek bir çekirdek modülü verilmiştir.

\sphinxAtStartPar
\sphinxstylestrong{test\sphinxhyphen{}module.c (next\_task ile):}

\begin{sphinxVerbatim}[commandchars=\\\{\}]
\PYG{c+cp}{\PYGZsh{}}\PYG{c+cp}{include}\PYG{+w}{ }\PYG{c+cpf}{\PYGZlt{}linux/module.h\PYGZgt{}}
\PYG{c+cp}{\PYGZsh{}}\PYG{c+cp}{include}\PYG{+w}{ }\PYG{c+cpf}{\PYGZlt{}linux/kernel.h\PYGZgt{}}
\PYG{c+cp}{\PYGZsh{}}\PYG{c+cp}{include}\PYG{+w}{ }\PYG{c+cpf}{\PYGZlt{}linux/rcupdate.h\PYGZgt{}}

\PYG{n}{MODULE\PYGZus{}LICENSE}\PYG{p}{(}\PYG{l+s}{\PYGZdq{}}\PYG{l+s}{GPL}\PYG{l+s}{\PYGZdq{}}\PYG{p}{)}\PYG{p}{;}

\PYG{k}{static}\PYG{+w}{ }\PYG{k+kt}{int}\PYG{+w}{ }\PYG{n+nf}{test\PYGZus{}module\PYGZus{}init}\PYG{p}{(}\PYG{k+kt}{void}\PYG{p}{)}\PYG{p}{;}
\PYG{k}{static}\PYG{+w}{ }\PYG{k+kt}{void}\PYG{+w}{ }\PYG{n+nf}{test\PYGZus{}module\PYGZus{}exit}\PYG{p}{(}\PYG{k+kt}{void}\PYG{p}{)}\PYG{p}{;}
\PYG{k}{static}\PYG{+w}{ }\PYG{k+kt}{void}\PYG{+w}{ }\PYG{n+nf}{walk\PYGZus{}processes}\PYG{p}{(}\PYG{k+kt}{void}\PYG{p}{)}\PYG{p}{;}

\PYG{k}{static}\PYG{+w}{ }\PYG{k+kt}{int}\PYG{+w}{ }\PYG{n+nf}{test\PYGZus{}module\PYGZus{}init}\PYG{p}{(}\PYG{k+kt}{void}\PYG{p}{)}
\PYG{p}{\PYGZob{}}
\PYG{+w}{    }\PYG{n}{printk}\PYG{p}{(}\PYG{n}{KERN\PYGZus{}INFO}\PYG{+w}{ }\PYG{l+s}{\PYGZdq{}}\PYG{l+s}{test\PYGZus{}module\PYGZus{}init...}\PYG{l+s+se}{\PYGZbs{}n}\PYG{l+s}{\PYGZdq{}}\PYG{p}{)}\PYG{p}{;}
\PYG{+w}{    }\PYG{n}{walk\PYGZus{}processes}\PYG{p}{(}\PYG{p}{)}\PYG{p}{;}

\PYG{+w}{    }\PYG{k}{return}\PYG{+w}{ }\PYG{l+m+mi}{0}\PYG{p}{;}
\PYG{p}{\PYGZcb{}}

\PYG{k}{static}\PYG{+w}{ }\PYG{k+kt}{void}\PYG{+w}{ }\PYG{n+nf}{walk\PYGZus{}processes}\PYG{p}{(}\PYG{k+kt}{void}\PYG{p}{)}
\PYG{p}{\PYGZob{}}
\PYG{+w}{    }\PYG{k}{struct}\PYG{+w}{ }\PYG{n+nc}{task\PYGZus{}struct}\PYG{+w}{ }\PYG{o}{*}\PYG{n}{ts}\PYG{p}{;}

\PYG{+w}{    }\PYG{n}{rcu\PYGZus{}read\PYGZus{}lock}\PYG{p}{(}\PYG{p}{)}\PYG{p}{;}

\PYG{+w}{    }\PYG{n}{ts}\PYG{+w}{ }\PYG{o}{=}\PYG{+w}{ }\PYG{n}{next\PYGZus{}task}\PYG{p}{(}\PYG{o}{\PYGZam{}}\PYG{n}{init\PYGZus{}task}\PYG{p}{)}\PYG{p}{;}
\PYG{+w}{    }\PYG{k}{while}\PYG{+w}{ }\PYG{p}{(}\PYG{n}{ts}\PYG{+w}{ }\PYG{o}{!}\PYG{o}{=}\PYG{+w}{ }\PYG{o}{\PYGZam{}}\PYG{n}{init\PYGZus{}task}\PYG{p}{)}\PYG{+w}{ }\PYG{p}{\PYGZob{}}
\PYG{+w}{        }\PYG{n}{printk}\PYG{p}{(}\PYG{n}{KERN\PYGZus{}INFO}\PYG{+w}{ }\PYG{l+s}{\PYGZdq{}}\PYG{l+s}{PID = \PYGZpc{}d, COMM = \PYGZpc{}s}\PYG{l+s}{\PYGZdq{}}\PYG{p}{,}\PYG{+w}{ }\PYG{n}{ts}\PYG{o}{\PYGZhy{}}\PYG{o}{\PYGZgt{}}\PYG{n}{pid}\PYG{p}{,}\PYG{+w}{ }\PYG{n}{ts}\PYG{o}{\PYGZhy{}}\PYG{o}{\PYGZgt{}}\PYG{n}{comm}\PYG{p}{)}\PYG{p}{;}
\PYG{+w}{        }\PYG{n}{ts}\PYG{+w}{ }\PYG{o}{=}\PYG{+w}{ }\PYG{n}{next\PYGZus{}task}\PYG{p}{(}\PYG{n}{ts}\PYG{p}{)}\PYG{p}{;}
\PYG{+w}{    }\PYG{p}{\PYGZcb{}}

\PYG{+w}{    }\PYG{n}{rcu\PYGZus{}read\PYGZus{}unlock}\PYG{p}{(}\PYG{p}{)}\PYG{p}{;}
\PYG{p}{\PYGZcb{}}

\PYG{k}{static}\PYG{+w}{ }\PYG{k+kt}{void}\PYG{+w}{ }\PYG{n+nf}{test\PYGZus{}module\PYGZus{}exit}\PYG{p}{(}\PYG{k+kt}{void}\PYG{p}{)}
\PYG{p}{\PYGZob{}}
\PYG{+w}{    }\PYG{n}{printk}\PYG{p}{(}\PYG{n}{KERN\PYGZus{}INFO}\PYG{+w}{ }\PYG{l+s}{\PYGZdq{}}\PYG{l+s}{test\PYGZus{}module\PYGZus{}exit...}\PYG{l+s+se}{\PYGZbs{}n}\PYG{l+s}{\PYGZdq{}}\PYG{p}{)}\PYG{p}{;}
\PYG{p}{\PYGZcb{}}

\PYG{n}{module\PYGZus{}init}\PYG{p}{(}\PYG{n}{test\PYGZus{}module\PYGZus{}init}\PYG{p}{)}\PYG{p}{;}
\PYG{n}{module\PYGZus{}exit}\PYG{p}{(}\PYG{n}{test\PYGZus{}module\PYGZus{}exit}\PYG{p}{)}\PYG{p}{;}
\end{sphinxVerbatim}

\sphinxAtStartPar
\sphinxcode{\sphinxupquote{\textless{}linux/sched/signal.h\textgreater{}}} başlık dosyasında bulunan \sphinxcode{\sphinxupquote{for\_each\_process}} isimli döngü makrosu da
\sphinxcode{\sphinxupquote{tasks}} listesini daha kolay dolaşmak için kullanılabilmektedir. Bu makro kendi içerisinde
\sphinxcode{\sphinxupquote{rcu\_dereference}} işlemini de yapmaktadır:

\begin{sphinxVerbatim}[commandchars=\\\{\}]
\PYG{n}{rcu\PYGZus{}read\PYGZus{}lock}\PYG{p}{(}\PYG{p}{)}\PYG{p}{;}

\PYG{n}{for\PYGZus{}each\PYGZus{}process}\PYG{p}{(}\PYG{n}{ts}\PYG{p}{)}\PYG{+w}{ }\PYG{p}{\PYGZob{}}
\PYG{+w}{    }\PYG{c+cm}{/* ... */}
\PYG{p}{\PYGZcb{}}

\PYG{n}{rcu\PYGZus{}read\PYGZus{}unlock}\PYG{p}{(}\PYG{p}{)}\PYG{p}{;}
\end{sphinxVerbatim}

\sphinxAtStartPar
\sphinxstylestrong{test\sphinxhyphen{}module.c (for\_each\_process ile):}

\begin{sphinxVerbatim}[commandchars=\\\{\}]
\PYG{c+cp}{\PYGZsh{}}\PYG{c+cp}{include}\PYG{+w}{ }\PYG{c+cpf}{\PYGZlt{}linux/module.h\PYGZgt{}}
\PYG{c+cp}{\PYGZsh{}}\PYG{c+cp}{include}\PYG{+w}{ }\PYG{c+cpf}{\PYGZlt{}linux/kernel.h\PYGZgt{}}
\PYG{c+cp}{\PYGZsh{}}\PYG{c+cp}{include}\PYG{+w}{ }\PYG{c+cpf}{\PYGZlt{}linux/sched/signal.h\PYGZgt{}}

\PYG{n}{MODULE\PYGZus{}LICENSE}\PYG{p}{(}\PYG{l+s}{\PYGZdq{}}\PYG{l+s}{GPL}\PYG{l+s}{\PYGZdq{}}\PYG{p}{)}\PYG{p}{;}

\PYG{k}{static}\PYG{+w}{ }\PYG{k+kt}{int}\PYG{+w}{ }\PYG{n+nf}{test\PYGZus{}module\PYGZus{}init}\PYG{p}{(}\PYG{k+kt}{void}\PYG{p}{)}\PYG{p}{;}
\PYG{k}{static}\PYG{+w}{ }\PYG{k+kt}{void}\PYG{+w}{ }\PYG{n+nf}{test\PYGZus{}module\PYGZus{}exit}\PYG{p}{(}\PYG{k+kt}{void}\PYG{p}{)}\PYG{p}{;}
\PYG{k}{static}\PYG{+w}{ }\PYG{k+kt}{void}\PYG{+w}{ }\PYG{n+nf}{walk\PYGZus{}processes}\PYG{p}{(}\PYG{k+kt}{void}\PYG{p}{)}\PYG{p}{;}

\PYG{k}{static}\PYG{+w}{ }\PYG{k+kt}{int}\PYG{+w}{ }\PYG{n+nf}{test\PYGZus{}module\PYGZus{}init}\PYG{p}{(}\PYG{k+kt}{void}\PYG{p}{)}
\PYG{p}{\PYGZob{}}
\PYG{+w}{    }\PYG{n}{printk}\PYG{p}{(}\PYG{n}{KERN\PYGZus{}INFO}\PYG{+w}{ }\PYG{l+s}{\PYGZdq{}}\PYG{l+s}{test\PYGZus{}module\PYGZus{}init...}\PYG{l+s+se}{\PYGZbs{}n}\PYG{l+s}{\PYGZdq{}}\PYG{p}{)}\PYG{p}{;}
\PYG{+w}{    }\PYG{n}{walk\PYGZus{}processes}\PYG{p}{(}\PYG{p}{)}\PYG{p}{;}

\PYG{+w}{    }\PYG{k}{return}\PYG{+w}{ }\PYG{l+m+mi}{0}\PYG{p}{;}
\PYG{p}{\PYGZcb{}}

\PYG{k}{static}\PYG{+w}{ }\PYG{k+kt}{void}\PYG{+w}{ }\PYG{n+nf}{walk\PYGZus{}processes}\PYG{p}{(}\PYG{k+kt}{void}\PYG{p}{)}
\PYG{p}{\PYGZob{}}
\PYG{+w}{    }\PYG{k}{struct}\PYG{+w}{ }\PYG{n+nc}{task\PYGZus{}struct}\PYG{+w}{ }\PYG{o}{*}\PYG{n}{ts}\PYG{p}{;}

\PYG{+w}{    }\PYG{n}{rcu\PYGZus{}read\PYGZus{}lock}\PYG{p}{(}\PYG{p}{)}\PYG{p}{;}

\PYG{+w}{    }\PYG{n}{for\PYGZus{}each\PYGZus{}process}\PYG{p}{(}\PYG{n}{ts}\PYG{p}{)}\PYG{+w}{ }\PYG{p}{\PYGZob{}}
\PYG{+w}{        }\PYG{n}{printk}\PYG{p}{(}\PYG{n}{KERN\PYGZus{}INFO}\PYG{+w}{ }\PYG{l+s}{\PYGZdq{}}\PYG{l+s}{PID = \PYGZpc{}d, COMM = \PYGZpc{}s}\PYG{l+s}{\PYGZdq{}}\PYG{p}{,}\PYG{+w}{ }\PYG{n}{ts}\PYG{o}{\PYGZhy{}}\PYG{o}{\PYGZgt{}}\PYG{n}{pid}\PYG{p}{,}\PYG{+w}{ }\PYG{n}{ts}\PYG{o}{\PYGZhy{}}\PYG{o}{\PYGZgt{}}\PYG{n}{comm}\PYG{p}{)}\PYG{p}{;}
\PYG{+w}{    }\PYG{p}{\PYGZcb{}}

\PYG{+w}{    }\PYG{n}{rcu\PYGZus{}read\PYGZus{}unlock}\PYG{p}{(}\PYG{p}{)}\PYG{p}{;}
\PYG{p}{\PYGZcb{}}

\PYG{k}{static}\PYG{+w}{ }\PYG{k+kt}{void}\PYG{+w}{ }\PYG{n+nf}{test\PYGZus{}module\PYGZus{}exit}\PYG{p}{(}\PYG{k+kt}{void}\PYG{p}{)}
\PYG{p}{\PYGZob{}}
\PYG{+w}{    }\PYG{n}{printk}\PYG{p}{(}\PYG{n}{KERN\PYGZus{}INFO}\PYG{+w}{ }\PYG{l+s}{\PYGZdq{}}\PYG{l+s}{test\PYGZus{}module\PYGZus{}exit...}\PYG{l+s+se}{\PYGZbs{}n}\PYG{l+s}{\PYGZdq{}}\PYG{p}{)}\PYG{p}{;}
\PYG{p}{\PYGZcb{}}

\PYG{n}{module\PYGZus{}init}\PYG{p}{(}\PYG{n}{test\PYGZus{}module\PYGZus{}init}\PYG{p}{)}\PYG{p}{;}
\PYG{n}{module\PYGZus{}exit}\PYG{p}{(}\PYG{n}{test\PYGZus{}module\PYGZus{}exit}\PYG{p}{)}\PYG{p}{;}
\end{sphinxVerbatim}

\sphinxAtStartPar
Ayrıca \sphinxcode{\sphinxupquote{\textless{}linux/rculist.h\textgreater{}}} dosyasında bulunan \sphinxcode{\sphinxupquote{list\_for\_each\_entry\_rcu}} isimli döngü makrosu da
RCU mekanizmasına uygun biçimde bağlı listeleri dolaşmak için kullanılabilmektedir:

\begin{sphinxVerbatim}[commandchars=\\\{\}]
\PYG{c+cp}{\PYGZsh{}}\PYG{c+cp}{include}\PYG{+w}{ }\PYG{c+cpf}{\PYGZlt{}linux/rculist.h\PYGZgt{}}

\PYG{n}{list\PYGZus{}for\PYGZus{}each\PYGZus{}entry\PYGZus{}rcu}\PYG{p}{(}\PYG{n}{pos}\PYG{p}{,}\PYG{+w}{ }\PYG{n}{head}\PYG{p}{,}\PYG{+w}{ }\PYG{n}{member}\PYG{p}{)}
\end{sphinxVerbatim}

\sphinxAtStartPar
Bu makroyu kullanarak tüm proseslerin ana thread’lerini dolaşan fonksiyon şöyle yazılabilir:

\begin{sphinxVerbatim}[commandchars=\\\{\}]
\PYG{k}{static}\PYG{+w}{ }\PYG{k+kt}{void}\PYG{+w}{ }\PYG{n+nf}{walk\PYGZus{}processes}\PYG{p}{(}\PYG{k+kt}{void}\PYG{p}{)}
\PYG{p}{\PYGZob{}}
\PYG{+w}{    }\PYG{k}{struct}\PYG{+w}{ }\PYG{n+nc}{task\PYGZus{}struct}\PYG{+w}{ }\PYG{o}{*}\PYG{n}{ts}\PYG{p}{;}

\PYG{+w}{    }\PYG{n}{rcu\PYGZus{}read\PYGZus{}lock}\PYG{p}{(}\PYG{p}{)}\PYG{p}{;}

\PYG{+w}{    }\PYG{n}{list\PYGZus{}for\PYGZus{}each\PYGZus{}entry\PYGZus{}rcu}\PYG{p}{(}\PYG{n}{ts}\PYG{p}{,}\PYG{+w}{ }\PYG{o}{\PYGZam{}}\PYG{n}{init\PYGZus{}task}\PYG{p}{.}\PYG{n}{tasks}\PYG{p}{,}\PYG{+w}{ }\PYG{n}{tasks}\PYG{p}{)}\PYG{+w}{ }\PYG{p}{\PYGZob{}}
\PYG{+w}{        }\PYG{n}{printk}\PYG{p}{(}\PYG{n}{KERN\PYGZus{}INFO}\PYG{+w}{ }\PYG{l+s}{\PYGZdq{}}\PYG{l+s}{PID = \PYGZpc{}d, COMM = \PYGZpc{}s}\PYG{l+s}{\PYGZdq{}}\PYG{p}{,}\PYG{+w}{ }\PYG{n}{ts}\PYG{o}{\PYGZhy{}}\PYG{o}{\PYGZgt{}}\PYG{n}{pid}\PYG{p}{,}\PYG{+w}{ }\PYG{n}{ts}\PYG{o}{\PYGZhy{}}\PYG{o}{\PYGZgt{}}\PYG{n}{comm}\PYG{p}{)}\PYG{p}{;}
\PYG{+w}{    }\PYG{p}{\PYGZcb{}}

\PYG{+w}{    }\PYG{n}{rcu\PYGZus{}read\PYGZus{}unlock}\PYG{p}{(}\PYG{p}{)}\PYG{p}{;}
\PYG{p}{\PYGZcb{}}
\end{sphinxVerbatim}

\sphinxAtStartPar
Aynı başlık dosyasındaki \sphinxcode{\sphinxupquote{list\_entry\_rcu}} makrosu bir bağın adresini alarak RCU mekanizmasına uygun bir
biçimde bağın bulunduğu nesnenin adresini vermektedir. Aslında önceki paragraflarda gördüğümüz
\sphinxcode{\sphinxupquote{next\_task}} makrosu da bu makro kullanılarak yazılmıştır:

\begin{sphinxVerbatim}[commandchars=\\\{\}]
\PYG{c+cp}{\PYGZsh{}}\PYG{c+cp}{define next\PYGZus{}task(p) \PYGZbs{}}
\PYG{c+cp}{    list\PYGZus{}entry\PYGZus{}rcu((p)\PYGZhy{}\PYGZgt{}tasks.next, struct task\PYGZus{}struct, tasks)}
\end{sphinxVerbatim}


\subsection{Thread Listesini Dolaşan Aygıt Sürücü Örneği}
\label{\detokenize{processcontrolblock:thread-listesini-dolasan-aygit-surucu-ornegi}}
\sphinxAtStartPar
Şimdi de thread’lere ilişkin \sphinxcode{\sphinxupquote{task\_struct}} nesnelerinin dolaşılmasına örnekler verelim. Anımsanacağı
gibi yeni çekirdeklerde belli bir prosesin thread’lerine ilişkin \sphinxcode{\sphinxupquote{task\_struct}} nesneleri
\sphinxcode{\sphinxupquote{signal\_struct}} nesnesindeki \sphinxcode{\sphinxupquote{thread\_head}} düğümü kök alınarak \sphinxcode{\sphinxupquote{thread\_node}} düğümlerinin
dolaşılmasıyla elde ediliyordu.

\sphinxAtStartPar
Belli bir prosesin thread’lerinin dolaşılabilmesi için çekirdek modülü yetersiz kalmaktadır. Çünkü
çekirdek modüllerindeki fonksiyonlar dışarıdan prosesler tarafından çağrılamamaktadır. Bu nedenle thread
dolaşım testlerini ancak bir aygıt sürücü yazarak oluşturabiliriz. Aygıt sürücülerdeki fonksiyonlar
kullanıcı modundaki thread’ler tarafından çağrılabilmektedir. Bu fonksiyonlarda \sphinxstyleemphasis{current} makrosu
kullanıldığında ilgili fonksiyonu çağıran thread’in \sphinxcode{\sphinxupquote{task\_struct}} nesnesinin adresi elde edilecektir.

\sphinxAtStartPar
Prosesin thread’leri aşağıdaki gibi dolaşılabilir:

\begin{sphinxVerbatim}[commandchars=\\\{\}]
\PYG{k}{static}\PYG{+w}{ }\PYG{k+kt}{void}\PYG{+w}{ }\PYG{n+nf}{walk\PYGZus{}process\PYGZus{}thread}\PYG{p}{(}\PYG{k+kt}{void}\PYG{p}{)}
\PYG{p}{\PYGZob{}}
\PYG{+w}{    }\PYG{k}{struct}\PYG{+w}{ }\PYG{n+nc}{task\PYGZus{}struct}\PYG{+w}{ }\PYG{o}{*}\PYG{n}{ts}\PYG{p}{;}

\PYG{+w}{    }\PYG{n}{rcu\PYGZus{}read\PYGZus{}lock}\PYG{p}{(}\PYG{p}{)}\PYG{p}{;}

\PYG{+w}{    }\PYG{n}{list\PYGZus{}for\PYGZus{}each\PYGZus{}entry\PYGZus{}rcu}\PYG{p}{(}\PYG{n}{ts}\PYG{p}{,}\PYG{+w}{ }\PYG{o}{\PYGZam{}}\PYG{n}{current}\PYG{o}{\PYGZhy{}}\PYG{o}{\PYGZgt{}}\PYG{n}{signal}\PYG{o}{\PYGZhy{}}\PYG{o}{\PYGZgt{}}\PYG{n}{thread\PYGZus{}head}\PYG{p}{,}\PYG{+w}{ }\PYG{n}{thread\PYGZus{}node}\PYG{p}{)}\PYG{+w}{ }\PYG{p}{\PYGZob{}}
\PYG{+w}{        }\PYG{n}{printk}\PYG{p}{(}\PYG{n}{KERN\PYGZus{}INFO}\PYG{+w}{ }\PYG{l+s}{\PYGZdq{}}\PYG{l+s}{Thread PID = \PYGZpc{}d, COMM = \PYGZpc{}s}\PYG{l+s}{\PYGZdq{}}\PYG{p}{,}\PYG{+w}{ }\PYG{n}{ts}\PYG{o}{\PYGZhy{}}\PYG{o}{\PYGZgt{}}\PYG{n}{pid}\PYG{p}{,}\PYG{+w}{ }\PYG{n}{ts}\PYG{o}{\PYGZhy{}}\PYG{o}{\PYGZgt{}}\PYG{n}{comm}\PYG{p}{)}\PYG{p}{;}
\PYG{+w}{    }\PYG{p}{\PYGZcb{}}

\PYG{+w}{    }\PYG{n}{rcu\PYGZus{}read\PYGZus{}unlock}\PYG{p}{(}\PYG{p}{)}\PYG{p}{;}
\PYG{p}{\PYGZcb{}}
\end{sphinxVerbatim}

\sphinxAtStartPar
Bir \sphinxcode{\sphinxupquote{task\_struct}} adresini alıp onun \sphinxcode{\sphinxupquote{thread\_node.next}} göstericisinin gösterdiği yerdeki
\sphinxcode{\sphinxupquote{task\_struct}} adresini veren \sphinxcode{\sphinxupquote{next\_thread}} isimli bir fonksiyon da vardır:

\begin{sphinxVerbatim}[commandchars=\\\{\}]
\PYG{c+cp}{\PYGZsh{}}\PYG{c+cp}{include}\PYG{+w}{ }\PYG{c+cpf}{\PYGZlt{}linux/sched/signal.h\PYGZgt{}}

\PYG{k}{static}\PYG{+w}{ }\PYG{k+kr}{inline}\PYG{+w}{ }\PYG{k}{struct}\PYG{+w}{ }\PYG{n+nc}{task\PYGZus{}struct}\PYG{+w}{ }\PYG{o}{*}\PYG{n}{next\PYGZus{}thread}\PYG{p}{(}\PYG{k}{struct}\PYG{+w}{ }\PYG{n+nc}{task\PYGZus{}struct}\PYG{+w}{ }\PYG{o}{*}\PYG{n}{p}\PYG{p}{)}
\PYG{p}{\PYGZob{}}
\PYG{+w}{    }\PYG{k}{return}\PYG{+w}{ }\PYG{n}{\PYGZus{}\PYGZus{}next\PYGZus{}thread}\PYG{p}{(}\PYG{n}{p}\PYG{p}{)}\PYG{+w}{ }\PYG{o}{?}\PYG{o}{:}\PYG{+w}{ }\PYG{n}{p}\PYG{o}{\PYGZhy{}}\PYG{o}{\PYGZgt{}}\PYG{n}{group\PYGZus{}leader}\PYG{p}{;}
\PYG{p}{\PYGZcb{}}
\end{sphinxVerbatim}

\sphinxAtStartPar
Bu fonksiyon ile prosesin herhangi bir thread’inden başlayarak tüm thread’leri dolaşılabilir:

\begin{sphinxVerbatim}[commandchars=\\\{\}]
\PYG{k}{static}\PYG{+w}{ }\PYG{k+kt}{void}\PYG{+w}{ }\PYG{n+nf}{walk\PYGZus{}process\PYGZus{}thread}\PYG{p}{(}\PYG{k+kt}{void}\PYG{p}{)}
\PYG{p}{\PYGZob{}}
\PYG{+w}{    }\PYG{k}{struct}\PYG{+w}{ }\PYG{n+nc}{task\PYGZus{}struct}\PYG{+w}{ }\PYG{o}{*}\PYG{n}{ts}\PYG{p}{;}

\PYG{+w}{    }\PYG{n}{rcu\PYGZus{}read\PYGZus{}lock}\PYG{p}{(}\PYG{p}{)}\PYG{p}{;}

\PYG{+w}{    }\PYG{n}{ts}\PYG{+w}{ }\PYG{o}{=}\PYG{+w}{ }\PYG{n}{current}\PYG{p}{;}
\PYG{+w}{    }\PYG{k}{do}\PYG{+w}{ }\PYG{p}{\PYGZob{}}
\PYG{+w}{        }\PYG{n}{printk}\PYG{p}{(}\PYG{n}{KERN\PYGZus{}INFO}\PYG{+w}{ }\PYG{l+s}{\PYGZdq{}}\PYG{l+s}{Thread PID = \PYGZpc{}d, COMM = \PYGZpc{}s}\PYG{l+s}{\PYGZdq{}}\PYG{p}{,}\PYG{+w}{ }\PYG{n}{ts}\PYG{o}{\PYGZhy{}}\PYG{o}{\PYGZgt{}}\PYG{n}{pid}\PYG{p}{,}\PYG{+w}{ }\PYG{n}{ts}\PYG{o}{\PYGZhy{}}\PYG{o}{\PYGZgt{}}\PYG{n}{comm}\PYG{p}{)}\PYG{p}{;}
\PYG{+w}{        }\PYG{n}{ts}\PYG{+w}{ }\PYG{o}{=}\PYG{+w}{ }\PYG{n}{next\PYGZus{}thread}\PYG{p}{(}\PYG{n}{ts}\PYG{p}{)}\PYG{p}{;}
\PYG{+w}{    }\PYG{p}{\PYGZcb{}}\PYG{+w}{ }\PYG{k}{while}\PYG{+w}{ }\PYG{p}{(}\PYG{n}{ts}\PYG{+w}{ }\PYG{o}{!}\PYG{o}{=}\PYG{+w}{ }\PYG{n}{current}\PYG{p}{)}\PYG{p}{;}

\PYG{+w}{    }\PYG{n}{rcu\PYGZus{}read\PYGZus{}unlock}\PYG{p}{(}\PYG{p}{)}\PYG{p}{;}
\PYG{p}{\PYGZcb{}}
\end{sphinxVerbatim}

\sphinxAtStartPar
Aslında prosesin thread’lerini basit bir biçimde dolaşmak için \sphinxcode{\sphinxupquote{for\_each\_thread}} döngü makrosu tercih
edilmektedir. Makronun birinci parametresi prosesin herhangi bir thread’inin \sphinxcode{\sphinxupquote{task\_struct}} adresini
almaktadır. Döngü her yinelendikçe yeni bir thread’e ilişkin \sphinxcode{\sphinxupquote{task\_struct}} yapının adresi ikinci
parametreyle belirtilen göstericinin içerisine yerleştirilmektedir:

\begin{sphinxVerbatim}[commandchars=\\\{\}]
\PYG{c+cp}{\PYGZsh{}}\PYG{c+cp}{include}\PYG{+w}{ }\PYG{c+cpf}{\PYGZlt{}linux/sched/signal.h\PYGZgt{}}

\PYG{n}{for\PYGZus{}each\PYGZus{}thread}\PYG{p}{(}\PYG{n}{p}\PYG{p}{,}\PYG{+w}{ }\PYG{n}{t}\PYG{p}{)}
\end{sphinxVerbatim}

\begin{sphinxVerbatim}[commandchars=\\\{\}]
\PYG{k}{static}\PYG{+w}{ }\PYG{k+kt}{void}\PYG{+w}{ }\PYG{n+nf}{walk\PYGZus{}process\PYGZus{}thread}\PYG{p}{(}\PYG{k+kt}{void}\PYG{p}{)}
\PYG{p}{\PYGZob{}}
\PYG{+w}{    }\PYG{k}{struct}\PYG{+w}{ }\PYG{n+nc}{task\PYGZus{}struct}\PYG{+w}{ }\PYG{o}{*}\PYG{n}{ts}\PYG{p}{;}

\PYG{+w}{    }\PYG{n}{rcu\PYGZus{}read\PYGZus{}lock}\PYG{p}{(}\PYG{p}{)}\PYG{p}{;}

\PYG{+w}{    }\PYG{n}{for\PYGZus{}each\PYGZus{}thread}\PYG{p}{(}\PYG{n}{current}\PYG{p}{,}\PYG{+w}{ }\PYG{n}{ts}\PYG{p}{)}\PYG{+w}{ }\PYG{p}{\PYGZob{}}
\PYG{+w}{        }\PYG{n}{printk}\PYG{p}{(}\PYG{n}{KERN\PYGZus{}INFO}\PYG{+w}{ }\PYG{l+s}{\PYGZdq{}}\PYG{l+s}{Thread PID = \PYGZpc{}d, COMM = \PYGZpc{}s}\PYG{l+s}{\PYGZdq{}}\PYG{p}{,}\PYG{+w}{ }\PYG{n}{ts}\PYG{o}{\PYGZhy{}}\PYG{o}{\PYGZgt{}}\PYG{n}{pid}\PYG{p}{,}\PYG{+w}{ }\PYG{n}{ts}\PYG{o}{\PYGZhy{}}\PYG{o}{\PYGZgt{}}\PYG{n}{comm}\PYG{p}{)}\PYG{p}{;}
\PYG{+w}{    }\PYG{p}{\PYGZcb{}}

\PYG{+w}{    }\PYG{n}{rcu\PYGZus{}read\PYGZus{}unlock}\PYG{p}{(}\PYG{p}{)}\PYG{p}{;}
\PYG{p}{\PYGZcb{}}
\end{sphinxVerbatim}

\sphinxAtStartPar
Aşağıda prosesin thread’lerini dolaşan tam aygıt sürücü kodları ve test programı verilmiştir. Aygıt
sürücüyü yükledikten sonra \sphinxstyleemphasis{app} programı çalıştırılmalıdır. Bu program önce 10 tane thread yaratıp sonra
aygıt sürücüdeki ioctl kodunu çalıştırmaktadır.

\sphinxAtStartPar
\sphinxstylestrong{test\sphinxhyphen{}driver.c (thread dolaşımı):}

\begin{sphinxVerbatim}[commandchars=\\\{\}]
\PYG{c+cp}{\PYGZsh{}}\PYG{c+cp}{include}\PYG{+w}{ }\PYG{c+cpf}{\PYGZlt{}linux/module.h\PYGZgt{}}
\PYG{c+cp}{\PYGZsh{}}\PYG{c+cp}{include}\PYG{+w}{ }\PYG{c+cpf}{\PYGZlt{}linux/kernel.h\PYGZgt{}}
\PYG{c+cp}{\PYGZsh{}}\PYG{c+cp}{include}\PYG{+w}{ }\PYG{c+cpf}{\PYGZlt{}linux/fs.h\PYGZgt{}}
\PYG{c+cp}{\PYGZsh{}}\PYG{c+cp}{include}\PYG{+w}{ }\PYG{c+cpf}{\PYGZlt{}linux/cdev.h\PYGZgt{}}
\PYG{c+cp}{\PYGZsh{}}\PYG{c+cp}{include}\PYG{+w}{ }\PYG{c+cpf}{\PYGZdq{}test\PYGZhy{}driver.h\PYGZdq{}}

\PYG{n}{MODULE\PYGZus{}LICENSE}\PYG{p}{(}\PYG{l+s}{\PYGZdq{}}\PYG{l+s}{GPL}\PYG{l+s}{\PYGZdq{}}\PYG{p}{)}\PYG{p}{;}
\PYG{n}{MODULE\PYGZus{}AUTHOR}\PYG{p}{(}\PYG{l+s}{\PYGZdq{}}\PYG{l+s}{Kaan Aslan}\PYG{l+s}{\PYGZdq{}}\PYG{p}{)}\PYG{p}{;}
\PYG{n}{MODULE\PYGZus{}DESCRIPTION}\PYG{p}{(}\PYG{l+s}{\PYGZdq{}}\PYG{l+s}{test\PYGZhy{}driver}\PYG{l+s}{\PYGZdq{}}\PYG{p}{)}\PYG{p}{;}

\PYG{c+cm}{/* ... (open/release/read/write fonksiyonları önceki örnekle aynıdır) ... */}

\PYG{k}{static}\PYG{+w}{ }\PYG{k+kt}{void}\PYG{+w}{ }\PYG{n+nf}{walk\PYGZus{}process\PYGZus{}thread}\PYG{p}{(}\PYG{k+kt}{void}\PYG{p}{)}\PYG{p}{;}

\PYG{k+kt}{long}\PYG{+w}{ }\PYG{n+nf}{ioctl\PYGZus{}test}\PYG{p}{(}\PYG{k}{struct}\PYG{+w}{ }\PYG{n+nc}{file}\PYG{+w}{ }\PYG{o}{*}\PYG{n}{filp}\PYG{p}{,}\PYG{+w}{ }\PYG{k+kt}{unsigned}\PYG{+w}{ }\PYG{k+kt}{long}\PYG{+w}{ }\PYG{n}{arg}\PYG{p}{)}
\PYG{p}{\PYGZob{}}
\PYG{+w}{    }\PYG{n}{walk\PYGZus{}process\PYGZus{}thread}\PYG{p}{(}\PYG{p}{)}\PYG{p}{;}

\PYG{+w}{    }\PYG{k}{return}\PYG{+w}{ }\PYG{l+m+mi}{0}\PYG{p}{;}
\PYG{p}{\PYGZcb{}}

\PYG{k}{static}\PYG{+w}{ }\PYG{k+kt}{void}\PYG{+w}{ }\PYG{n+nf}{walk\PYGZus{}process\PYGZus{}thread}\PYG{p}{(}\PYG{k+kt}{void}\PYG{p}{)}
\PYG{p}{\PYGZob{}}
\PYG{+w}{    }\PYG{k}{struct}\PYG{+w}{ }\PYG{n+nc}{task\PYGZus{}struct}\PYG{+w}{ }\PYG{o}{*}\PYG{n}{ts}\PYG{p}{;}

\PYG{+w}{    }\PYG{n}{rcu\PYGZus{}read\PYGZus{}lock}\PYG{p}{(}\PYG{p}{)}\PYG{p}{;}

\PYG{+w}{    }\PYG{n}{for\PYGZus{}each\PYGZus{}thread}\PYG{p}{(}\PYG{n}{current}\PYG{p}{,}\PYG{+w}{ }\PYG{n}{ts}\PYG{p}{)}\PYG{+w}{ }\PYG{p}{\PYGZob{}}
\PYG{+w}{        }\PYG{n}{printk}\PYG{p}{(}\PYG{n}{KERN\PYGZus{}INFO}\PYG{+w}{ }\PYG{l+s}{\PYGZdq{}}\PYG{l+s}{Thread PID = \PYGZpc{}d, COMM = \PYGZpc{}s}\PYG{l+s}{\PYGZdq{}}\PYG{p}{,}\PYG{+w}{ }\PYG{n}{ts}\PYG{o}{\PYGZhy{}}\PYG{o}{\PYGZgt{}}\PYG{n}{pid}\PYG{p}{,}\PYG{+w}{ }\PYG{n}{ts}\PYG{o}{\PYGZhy{}}\PYG{o}{\PYGZgt{}}\PYG{n}{comm}\PYG{p}{)}\PYG{p}{;}
\PYG{+w}{    }\PYG{p}{\PYGZcb{}}

\PYG{+w}{    }\PYG{n}{rcu\PYGZus{}read\PYGZus{}unlock}\PYG{p}{(}\PYG{p}{)}\PYG{p}{;}
\PYG{p}{\PYGZcb{}}

\PYG{n}{module\PYGZus{}init}\PYG{p}{(}\PYG{n}{test\PYGZus{}driver\PYGZus{}init}\PYG{p}{)}\PYG{p}{;}
\PYG{n}{module\PYGZus{}exit}\PYG{p}{(}\PYG{n}{test\PYGZus{}driver\PYGZus{}exit}\PYG{p}{)}\PYG{p}{;}
\end{sphinxVerbatim}

\sphinxAtStartPar
\sphinxstylestrong{app.c (thread dolaşımı için test programı):}

\begin{sphinxVerbatim}[commandchars=\\\{\}]
\PYG{c+cp}{\PYGZsh{}}\PYG{c+cp}{define \PYGZus{}GNU\PYGZus{}SOURCE}

\PYG{c+cp}{\PYGZsh{}}\PYG{c+cp}{include}\PYG{+w}{ }\PYG{c+cpf}{\PYGZlt{}stdio.h\PYGZgt{}}
\PYG{c+cp}{\PYGZsh{}}\PYG{c+cp}{include}\PYG{+w}{ }\PYG{c+cpf}{\PYGZlt{}stdlib.h\PYGZgt{}}
\PYG{c+cp}{\PYGZsh{}}\PYG{c+cp}{include}\PYG{+w}{ }\PYG{c+cpf}{\PYGZlt{}string.h\PYGZgt{}}
\PYG{c+cp}{\PYGZsh{}}\PYG{c+cp}{include}\PYG{+w}{ }\PYG{c+cpf}{\PYGZlt{}stdint.h\PYGZgt{}}
\PYG{c+cp}{\PYGZsh{}}\PYG{c+cp}{include}\PYG{+w}{ }\PYG{c+cpf}{\PYGZlt{}fcntl.h\PYGZgt{}}
\PYG{c+cp}{\PYGZsh{}}\PYG{c+cp}{include}\PYG{+w}{ }\PYG{c+cpf}{\PYGZlt{}unistd.h\PYGZgt{}}
\PYG{c+cp}{\PYGZsh{}}\PYG{c+cp}{include}\PYG{+w}{ }\PYG{c+cpf}{\PYGZlt{}pthread.h\PYGZgt{}}
\PYG{c+cp}{\PYGZsh{}}\PYG{c+cp}{include}\PYG{+w}{ }\PYG{c+cpf}{\PYGZlt{}sys/ioctl.h\PYGZgt{}}
\PYG{c+cp}{\PYGZsh{}}\PYG{c+cp}{include}\PYG{+w}{ }\PYG{c+cpf}{\PYGZdq{}test\PYGZhy{}driver.h\PYGZdq{}}

\PYG{c+cp}{\PYGZsh{}}\PYG{c+cp}{define NTHREADS 10}

\PYG{k+kt}{void}\PYG{+w}{ }\PYG{n+nf}{exit\PYGZus{}sys}\PYG{p}{(}\PYG{k}{const}\PYG{+w}{ }\PYG{k+kt}{char}\PYG{+w}{ }\PYG{o}{*}\PYG{n}{msg}\PYG{p}{)}\PYG{p}{;}
\PYG{k+kt}{void}\PYG{+w}{ }\PYG{o}{*}\PYG{n+nf}{thread\PYGZus{}proc}\PYG{p}{(}\PYG{k+kt}{void}\PYG{+w}{ }\PYG{o}{*}\PYG{n}{param}\PYG{p}{)}\PYG{p}{;}

\PYG{k+kt}{int}\PYG{+w}{ }\PYG{n+nf}{main}\PYG{p}{(}\PYG{k+kt}{void}\PYG{p}{)}
\PYG{p}{\PYGZob{}}
\PYG{+w}{    }\PYG{k+kt}{int}\PYG{+w}{ }\PYG{n}{fd}\PYG{p}{;}
\PYG{+w}{    }\PYG{k+kt}{int}\PYG{+w}{ }\PYG{n}{result}\PYG{p}{;}
\PYG{+w}{    }\PYG{n}{pthread\PYGZus{}t}\PYG{+w}{ }\PYG{n}{tids}\PYG{p}{[}\PYG{n}{NTHREADS}\PYG{p}{]}\PYG{p}{;}

\PYG{+w}{    }\PYG{k}{for}\PYG{+w}{ }\PYG{p}{(}\PYG{k+kt}{int}\PYG{+w}{ }\PYG{n}{i}\PYG{+w}{ }\PYG{o}{=}\PYG{+w}{ }\PYG{l+m+mi}{0}\PYG{p}{;}\PYG{+w}{ }\PYG{n}{i}\PYG{+w}{ }\PYG{o}{\PYGZlt{}}\PYG{+w}{ }\PYG{n}{NTHREADS}\PYG{p}{;}\PYG{+w}{ }\PYG{o}{+}\PYG{o}{+}\PYG{n}{i}\PYG{p}{)}\PYG{+w}{ }\PYG{p}{\PYGZob{}}
\PYG{+w}{        }\PYG{k}{if}\PYG{+w}{ }\PYG{p}{(}\PYG{p}{(}\PYG{n}{result}\PYG{+w}{ }\PYG{o}{=}\PYG{+w}{ }\PYG{n}{pthread\PYGZus{}create}\PYG{p}{(}\PYG{o}{\PYGZam{}}\PYG{n}{tids}\PYG{p}{[}\PYG{n}{i}\PYG{p}{]}\PYG{p}{,}\PYG{+w}{ }\PYG{n+nb}{NULL}\PYG{p}{,}\PYG{+w}{ }\PYG{n}{thread\PYGZus{}proc}\PYG{p}{,}\PYG{+w}{ }\PYG{n+nb}{NULL}\PYG{p}{)}\PYG{p}{)}\PYG{+w}{ }\PYG{o}{!}\PYG{o}{=}\PYG{+w}{ }\PYG{l+m+mi}{0}\PYG{p}{)}\PYG{+w}{ }\PYG{p}{\PYGZob{}}
\PYG{+w}{            }\PYG{n}{fprintf}\PYG{p}{(}\PYG{n}{stderr}\PYG{p}{,}\PYG{+w}{ }\PYG{l+s}{\PYGZdq{}}\PYG{l+s}{pthread\PYGZus{}create: \PYGZpc{}s}\PYG{l+s+se}{\PYGZbs{}n}\PYG{l+s}{\PYGZdq{}}\PYG{p}{,}\PYG{+w}{ }\PYG{n}{strerror}\PYG{p}{(}\PYG{n}{result}\PYG{p}{)}\PYG{p}{)}\PYG{p}{;}
\PYG{+w}{            }\PYG{n}{exit}\PYG{p}{(}\PYG{n}{EXIT\PYGZus{}FAILURE}\PYG{p}{)}\PYG{p}{;}
\PYG{+w}{        }\PYG{p}{\PYGZcb{}}
\PYG{+w}{    }\PYG{p}{\PYGZcb{}}

\PYG{+w}{    }\PYG{k}{if}\PYG{+w}{ }\PYG{p}{(}\PYG{p}{(}\PYG{n}{fd}\PYG{+w}{ }\PYG{o}{=}\PYG{+w}{ }\PYG{n}{open}\PYG{p}{(}\PYG{l+s}{\PYGZdq{}}\PYG{l+s}{test\PYGZhy{}driver}\PYG{l+s}{\PYGZdq{}}\PYG{p}{,}\PYG{+w}{ }\PYG{n}{O\PYGZus{}RDONLY}\PYG{p}{)}\PYG{p}{)}\PYG{+w}{ }\PYG{o}{=}\PYG{o}{=}\PYG{+w}{ }\PYG{l+m+mi}{\PYGZhy{}1}\PYG{p}{)}
\PYG{+w}{        }\PYG{n}{exit\PYGZus{}sys}\PYG{p}{(}\PYG{l+s}{\PYGZdq{}}\PYG{l+s}{open}\PYG{l+s}{\PYGZdq{}}\PYG{p}{)}\PYG{p}{;}

\PYG{+w}{    }\PYG{k}{if}\PYG{+w}{ }\PYG{p}{(}\PYG{n}{ioctl}\PYG{p}{(}\PYG{n}{fd}\PYG{p}{,}\PYG{+w}{ }\PYG{n}{IOC\PYGZus{}TEST}\PYG{p}{)}\PYG{+w}{ }\PYG{o}{=}\PYG{o}{=}\PYG{+w}{ }\PYG{l+m+mi}{\PYGZhy{}1}\PYG{p}{)}
\PYG{+w}{        }\PYG{n}{exit\PYGZus{}sys}\PYG{p}{(}\PYG{l+s}{\PYGZdq{}}\PYG{l+s}{ioctl}\PYG{l+s}{\PYGZdq{}}\PYG{p}{)}\PYG{p}{;}

\PYG{+w}{    }\PYG{n}{close}\PYG{p}{(}\PYG{n}{fd}\PYG{p}{)}\PYG{p}{;}

\PYG{+w}{    }\PYG{k}{for}\PYG{+w}{ }\PYG{p}{(}\PYG{k+kt}{int}\PYG{+w}{ }\PYG{n}{i}\PYG{+w}{ }\PYG{o}{=}\PYG{+w}{ }\PYG{l+m+mi}{0}\PYG{p}{;}\PYG{+w}{ }\PYG{n}{i}\PYG{+w}{ }\PYG{o}{\PYGZlt{}}\PYG{+w}{ }\PYG{n}{NTHREADS}\PYG{p}{;}\PYG{+w}{ }\PYG{o}{+}\PYG{o}{+}\PYG{n}{i}\PYG{p}{)}
\PYG{+w}{        }\PYG{n}{pthread\PYGZus{}join}\PYG{p}{(}\PYG{n}{tids}\PYG{p}{[}\PYG{n}{i}\PYG{p}{]}\PYG{p}{,}\PYG{+w}{ }\PYG{n+nb}{NULL}\PYG{p}{)}\PYG{p}{;}

\PYG{+w}{    }\PYG{k}{return}\PYG{+w}{ }\PYG{l+m+mi}{0}\PYG{p}{;}
\PYG{p}{\PYGZcb{}}

\PYG{k+kt}{void}\PYG{+w}{ }\PYG{o}{*}\PYG{n+nf}{thread\PYGZus{}proc}\PYG{p}{(}\PYG{k+kt}{void}\PYG{+w}{ }\PYG{o}{*}\PYG{n}{param}\PYG{p}{)}
\PYG{p}{\PYGZob{}}
\PYG{+w}{    }\PYG{n}{printf}\PYG{p}{(}\PYG{l+s}{\PYGZdq{}}\PYG{l+s}{Thread PID: \PYGZpc{}jd}\PYG{l+s+se}{\PYGZbs{}n}\PYG{l+s}{\PYGZdq{}}\PYG{p}{,}\PYG{+w}{ }\PYG{p}{(}\PYG{k+kt}{intmax\PYGZus{}t}\PYG{p}{)}\PYG{n}{gettid}\PYG{p}{(}\PYG{p}{)}\PYG{p}{)}\PYG{p}{;}
\PYG{+w}{    }\PYG{n}{sleep}\PYG{p}{(}\PYG{l+m+mi}{120}\PYG{p}{)}\PYG{p}{;}

\PYG{+w}{    }\PYG{k}{return}\PYG{+w}{ }\PYG{n+nb}{NULL}\PYG{p}{;}
\PYG{p}{\PYGZcb{}}

\PYG{k+kt}{void}\PYG{+w}{ }\PYG{n+nf}{exit\PYGZus{}sys}\PYG{p}{(}\PYG{k}{const}\PYG{+w}{ }\PYG{k+kt}{char}\PYG{+w}{ }\PYG{o}{*}\PYG{n}{msg}\PYG{p}{)}
\PYG{p}{\PYGZob{}}
\PYG{+w}{    }\PYG{n}{perror}\PYG{p}{(}\PYG{n}{msg}\PYG{p}{)}\PYG{p}{;}
\PYG{+w}{    }\PYG{n}{exit}\PYG{p}{(}\PYG{n}{EXIT\PYGZus{}FAILURE}\PYG{p}{)}\PYG{p}{;}
\PYG{p}{\PYGZcb{}}
\end{sphinxVerbatim}


\section{Tüm \sphinxstyleliteralintitle{\sphinxupquote{task\_struct}} Nesnelerini Dolaşma}
\label{\detokenize{processcontrolblock:tum-task-struct-nesnelerini-dolasma}}
\sphinxAtStartPar
Peki sistemdeki bütün \sphinxcode{\sphinxupquote{task\_struct}} nesnelerini dolaşabilir miyiz? Bunu yapmanın en pratik yolu
\sphinxcode{\sphinxupquote{init\_task}} nesnesinden hareketle tüm proseslerin ana thread’lerine ilişkin \sphinxcode{\sphinxupquote{task\_struct}} nesnelerini
elde edip o ana thread’lerden faydalanarak onların diğer thread’lerine ilişkin \sphinxcode{\sphinxupquote{task\_struct}} nesnelerini
elde etmektir. Aslında bunu iç içe iki döngü oluşturarak RCU mekanizması eşliğinde yapan
\sphinxcode{\sphinxupquote{for\_each\_process\_thread}} isimli bir makro vardır:

\begin{sphinxVerbatim}[commandchars=\\\{\}]
\PYG{c+cp}{\PYGZsh{}}\PYG{c+cp}{include}\PYG{+w}{ }\PYG{c+cpf}{\PYGZlt{}linux/sched/signal.h\PYGZgt{}}

\PYG{c+cp}{\PYGZsh{}}\PYG{c+cp}{define for\PYGZus{}each\PYGZus{}process\PYGZus{}thread(p, t) \PYGZbs{}}
\PYG{c+cp}{    for\PYGZus{}each\PYGZus{}process(p) for\PYGZus{}each\PYGZus{}thread(p, t)}
\end{sphinxVerbatim}

\sphinxAtStartPar
Makroya biz iki \sphinxcode{\sphinxupquote{task\_struct}} göstericisi veririz. Makro her yinelemede prosesin ana thread’ine ilişkin
\sphinxcode{\sphinxupquote{task\_struct}} nesnesinin adresini birinci parametreyle verilen göstericinin içerisine, o prosesin
thread’lerine ilişkin \sphinxcode{\sphinxupquote{task\_struct}} nesnelerinin adresleri de ikinci parametreyle verilen göstericinin
içerisine yerleştirmektedir:

\begin{sphinxVerbatim}[commandchars=\\\{\}]
\PYG{k}{static}\PYG{+w}{ }\PYG{k+kt}{void}\PYG{+w}{ }\PYG{n+nf}{walk\PYGZus{}all\PYGZus{}threads}\PYG{p}{(}\PYG{k+kt}{void}\PYG{p}{)}
\PYG{p}{\PYGZob{}}
\PYG{+w}{    }\PYG{k}{struct}\PYG{+w}{ }\PYG{n+nc}{task\PYGZus{}struct}\PYG{+w}{ }\PYG{o}{*}\PYG{n}{tsp}\PYG{p}{;}
\PYG{+w}{    }\PYG{k}{struct}\PYG{+w}{ }\PYG{n+nc}{task\PYGZus{}struct}\PYG{+w}{ }\PYG{o}{*}\PYG{n}{tst}\PYG{p}{;}

\PYG{+w}{    }\PYG{n}{rcu\PYGZus{}read\PYGZus{}lock}\PYG{p}{(}\PYG{p}{)}\PYG{p}{;}

\PYG{+w}{    }\PYG{n}{for\PYGZus{}each\PYGZus{}process\PYGZus{}thread}\PYG{p}{(}\PYG{n}{tsp}\PYG{p}{,}\PYG{+w}{ }\PYG{n}{tst}\PYG{p}{)}\PYG{+w}{ }\PYG{p}{\PYGZob{}}
\PYG{+w}{        }\PYG{k}{if}\PYG{+w}{ }\PYG{p}{(}\PYG{n}{tsp}\PYG{+w}{ }\PYG{o}{=}\PYG{o}{=}\PYG{+w}{ }\PYG{n}{tst}\PYG{p}{)}\PYG{+w}{ }\PYG{p}{\PYGZob{}}
\PYG{+w}{            }\PYG{n}{printk}\PYG{p}{(}\PYG{n}{KERN\PYGZus{}INFO}\PYG{+w}{ }\PYG{l+s}{\PYGZdq{}}\PYG{l+s}{Process Main Thread PID = \PYGZpc{}d, COMM = \PYGZpc{}s}\PYG{l+s+se}{\PYGZbs{}n}\PYG{l+s}{\PYGZdq{}}\PYG{p}{,}
\PYG{+w}{                   }\PYG{n}{tsp}\PYG{o}{\PYGZhy{}}\PYG{o}{\PYGZgt{}}\PYG{n}{pid}\PYG{p}{,}\PYG{+w}{ }\PYG{n}{tst}\PYG{o}{\PYGZhy{}}\PYG{o}{\PYGZgt{}}\PYG{n}{comm}\PYG{p}{)}\PYG{p}{;}
\PYG{+w}{        }\PYG{p}{\PYGZcb{}}
\PYG{+w}{        }\PYG{k}{else}\PYG{+w}{ }\PYG{p}{\PYGZob{}}
\PYG{+w}{            }\PYG{n}{printk}\PYG{p}{(}\PYG{n}{KERN\PYGZus{}INFO}\PYG{+w}{ }\PYG{l+s}{\PYGZdq{}}\PYG{l+s}{    Thread PID == \PYGZpc{}d COMM = \PYGZpc{}s}\PYG{l+s+se}{\PYGZbs{}n}\PYG{l+s}{\PYGZdq{}}\PYG{p}{,}
\PYG{+w}{                   }\PYG{n}{tst}\PYG{o}{\PYGZhy{}}\PYG{o}{\PYGZgt{}}\PYG{n}{pid}\PYG{p}{,}\PYG{+w}{ }\PYG{n}{tst}\PYG{o}{\PYGZhy{}}\PYG{o}{\PYGZgt{}}\PYG{n}{comm}\PYG{p}{)}\PYG{p}{;}
\PYG{+w}{        }\PYG{p}{\PYGZcb{}}
\PYG{+w}{    }\PYG{p}{\PYGZcb{}}

\PYG{+w}{    }\PYG{n}{rcu\PYGZus{}read\PYGZus{}unlock}\PYG{p}{(}\PYG{p}{)}\PYG{p}{;}
\PYG{p}{\PYGZcb{}}
\end{sphinxVerbatim}

\sphinxAtStartPar
\sphinxstyleemphasis{dmesg} çıktısı aşağıdakine benzer biçimde elde edilecektir:

\begin{sphinxVerbatim}[commandchars=\\\{\}]
[ 6022.928044] Process Main Thread PID = 4271, COMM = kworker/0:0
[ 6022.928045] Process Main Thread PID = 4279, COMM = kworker/u259:0
[ 6022.928048] Process Main Thread PID = 4543, COMM = kworker/3:1
[ 6022.928051] Process Main Thread PID = 4565, COMM = bash
[ 6022.928053] Process Main Thread PID = 4599, COMM = app
[ 6022.928054]     Thread PID == 4600 COMM = app
[ 6022.928056]     Thread PID == 4601 COMM = app
[ 6022.928057]     Thread PID == 4602 COMM = app
[ 6022.928059]     Thread PID == 4603 COMM = app
[ 6022.928060]     Thread PID == 4604 COMM = app
[ 6022.928087] Process Main Thread PID = 4610, COMM = sudo
[ 6022.928090] Process Main Thread PID = 4612, COMM = insmod
\end{sphinxVerbatim}


\subsection{Alt Proses Listesini Dolaşma}
\label{\detokenize{processcontrolblock:alt-proses-listesini-dolasma}}
\sphinxAtStartPar
Anımsanacağı gibi prosesin alt proses listesinin kök düğümü ana prosesin ana thread’ine ilişkin
\sphinxcode{\sphinxupquote{task\_struct}} nesnesinin \sphinxcode{\sphinxupquote{children}} elemanında tutuluyordu. \sphinxcode{\sphinxupquote{children}}/\sphinxcode{\sphinxupquote{sibling}} listesinde
yalnızca alt proseslerin ana thread’lerine ilişkin \sphinxcode{\sphinxupquote{task\_struct}} nesnelerinin bulunduğunu da
anımsayınız.

\sphinxAtStartPar
Alt prosesleri dolaşan hazır bir döngü makrosu yoktur. Ancak biz yukarıdaki tekniklerle dolaşımı
yapabiliriz:

\begin{sphinxVerbatim}[commandchars=\\\{\}]
\PYG{k}{static}\PYG{+w}{ }\PYG{k+kt}{void}\PYG{+w}{ }\PYG{n+nf}{walk\PYGZus{}child\PYGZus{}processes}\PYG{p}{(}\PYG{k+kt}{void}\PYG{p}{)}
\PYG{p}{\PYGZob{}}
\PYG{+w}{    }\PYG{k}{struct}\PYG{+w}{ }\PYG{n+nc}{task\PYGZus{}struct}\PYG{+w}{ }\PYG{o}{*}\PYG{n}{ts}\PYG{p}{;}

\PYG{+w}{    }\PYG{n}{rcu\PYGZus{}read\PYGZus{}lock}\PYG{p}{(}\PYG{p}{)}\PYG{p}{;}

\PYG{+w}{    }\PYG{n}{list\PYGZus{}for\PYGZus{}each\PYGZus{}entry\PYGZus{}rcu}\PYG{p}{(}\PYG{n}{ts}\PYG{p}{,}\PYG{+w}{ }\PYG{o}{\PYGZam{}}\PYG{n}{current}\PYG{o}{\PYGZhy{}}\PYG{o}{\PYGZgt{}}\PYG{n}{children}\PYG{p}{,}\PYG{+w}{ }\PYG{n}{sibling}\PYG{p}{)}\PYG{+w}{ }\PYG{p}{\PYGZob{}}
\PYG{+w}{        }\PYG{n}{printk}\PYG{p}{(}\PYG{n}{KERN\PYGZus{}INFO}\PYG{+w}{ }\PYG{l+s}{\PYGZdq{}}\PYG{l+s}{Child Process Main Thread PID = \PYGZpc{}d, COMM = \PYGZpc{}s}\PYG{l+s+se}{\PYGZbs{}n}\PYG{l+s}{\PYGZdq{}}\PYG{p}{,}
\PYG{+w}{               }\PYG{n}{ts}\PYG{o}{\PYGZhy{}}\PYG{o}{\PYGZgt{}}\PYG{n}{pid}\PYG{p}{,}\PYG{+w}{ }\PYG{n}{ts}\PYG{o}{\PYGZhy{}}\PYG{o}{\PYGZgt{}}\PYG{n}{comm}\PYG{p}{)}\PYG{p}{;}
\PYG{+w}{    }\PYG{p}{\PYGZcb{}}

\PYG{+w}{    }\PYG{n}{rcu\PYGZus{}read\PYGZus{}unlock}\PYG{p}{(}\PYG{p}{)}\PYG{p}{;}
\PYG{p}{\PYGZcb{}}
\end{sphinxVerbatim}

\sphinxAtStartPar
Bu kodu bir karakter aygıt sürücüsünün ioctl kodu içerisinden çağırabilirsiniz. Kodu önce çeşitli alt
prosesler yarattıktan sonra aygıt sürücü üzerinde ioctl çağrısı yapan bir programla test edebilirsiniz.
Aşağıdaki test programı 10 tane alt proses ve her alt proseste 3 tane thread yaratmaktadır:

\sphinxAtStartPar
\sphinxstylestrong{app.c (alt proses dolaşımı için test programı):}

\begin{sphinxVerbatim}[commandchars=\\\{\}]
\PYG{c+cp}{\PYGZsh{}}\PYG{c+cp}{define \PYGZus{}GNU\PYGZus{}SOURCE}

\PYG{c+cp}{\PYGZsh{}}\PYG{c+cp}{include}\PYG{+w}{ }\PYG{c+cpf}{\PYGZlt{}stdio.h\PYGZgt{}}
\PYG{c+cp}{\PYGZsh{}}\PYG{c+cp}{include}\PYG{+w}{ }\PYG{c+cpf}{\PYGZlt{}stdlib.h\PYGZgt{}}
\PYG{c+cp}{\PYGZsh{}}\PYG{c+cp}{include}\PYG{+w}{ }\PYG{c+cpf}{\PYGZlt{}string.h\PYGZgt{}}
\PYG{c+cp}{\PYGZsh{}}\PYG{c+cp}{include}\PYG{+w}{ }\PYG{c+cpf}{\PYGZlt{}stdint.h\PYGZgt{}}
\PYG{c+cp}{\PYGZsh{}}\PYG{c+cp}{include}\PYG{+w}{ }\PYG{c+cpf}{\PYGZlt{}fcntl.h\PYGZgt{}}
\PYG{c+cp}{\PYGZsh{}}\PYG{c+cp}{include}\PYG{+w}{ }\PYG{c+cpf}{\PYGZlt{}unistd.h\PYGZgt{}}
\PYG{c+cp}{\PYGZsh{}}\PYG{c+cp}{include}\PYG{+w}{ }\PYG{c+cpf}{\PYGZlt{}sys/wait.h\PYGZgt{}}
\PYG{c+cp}{\PYGZsh{}}\PYG{c+cp}{include}\PYG{+w}{ }\PYG{c+cpf}{\PYGZlt{}pthread.h\PYGZgt{}}
\PYG{c+cp}{\PYGZsh{}}\PYG{c+cp}{include}\PYG{+w}{ }\PYG{c+cpf}{\PYGZlt{}sys/ioctl.h\PYGZgt{}}
\PYG{c+cp}{\PYGZsh{}}\PYG{c+cp}{include}\PYG{+w}{ }\PYG{c+cpf}{\PYGZdq{}test\PYGZhy{}driver.h\PYGZdq{}}

\PYG{c+cp}{\PYGZsh{}}\PYG{c+cp}{define NCHILDS  10}
\PYG{c+cp}{\PYGZsh{}}\PYG{c+cp}{define NTHREADS 3}

\PYG{k+kt}{void}\PYG{+w}{ }\PYG{n+nf}{exit\PYGZus{}sys}\PYG{p}{(}\PYG{k}{const}\PYG{+w}{ }\PYG{k+kt}{char}\PYG{+w}{ }\PYG{o}{*}\PYG{n}{msg}\PYG{p}{)}\PYG{p}{;}
\PYG{k+kt}{void}\PYG{+w}{ }\PYG{o}{*}\PYG{n+nf}{thread\PYGZus{}proc}\PYG{p}{(}\PYG{k+kt}{void}\PYG{+w}{ }\PYG{o}{*}\PYG{n}{param}\PYG{p}{)}\PYG{p}{;}

\PYG{k+kt}{int}\PYG{+w}{ }\PYG{n+nf}{main}\PYG{p}{(}\PYG{k+kt}{void}\PYG{p}{)}
\PYG{p}{\PYGZob{}}
\PYG{+w}{    }\PYG{k+kt}{int}\PYG{+w}{ }\PYG{n}{fd}\PYG{p}{;}
\PYG{+w}{    }\PYG{k+kt}{int}\PYG{+w}{ }\PYG{n}{result}\PYG{p}{;}
\PYG{+w}{    }\PYG{k+kt}{pid\PYGZus{}t}\PYG{+w}{ }\PYG{n}{pid}\PYG{p}{;}

\PYG{+w}{    }\PYG{k}{for}\PYG{+w}{ }\PYG{p}{(}\PYG{k+kt}{int}\PYG{+w}{ }\PYG{n}{i}\PYG{+w}{ }\PYG{o}{=}\PYG{+w}{ }\PYG{l+m+mi}{0}\PYG{p}{;}\PYG{+w}{ }\PYG{n}{i}\PYG{+w}{ }\PYG{o}{\PYGZlt{}}\PYG{+w}{ }\PYG{n}{NCHILDS}\PYG{p}{;}\PYG{+w}{ }\PYG{o}{+}\PYG{o}{+}\PYG{n}{i}\PYG{p}{)}\PYG{+w}{ }\PYG{p}{\PYGZob{}}
\PYG{+w}{        }\PYG{k}{if}\PYG{+w}{ }\PYG{p}{(}\PYG{p}{(}\PYG{n}{pid}\PYG{+w}{ }\PYG{o}{=}\PYG{+w}{ }\PYG{n}{fork}\PYG{p}{(}\PYG{p}{)}\PYG{p}{)}\PYG{+w}{ }\PYG{o}{=}\PYG{o}{=}\PYG{+w}{ }\PYG{l+m+mi}{\PYGZhy{}1}\PYG{p}{)}\PYG{+w}{ }\PYG{p}{\PYGZob{}}
\PYG{+w}{            }\PYG{n}{perror}\PYG{p}{(}\PYG{l+s}{\PYGZdq{}}\PYG{l+s}{fork}\PYG{l+s}{\PYGZdq{}}\PYG{p}{)}\PYG{p}{;}
\PYG{+w}{            }\PYG{n}{exit}\PYG{p}{(}\PYG{n}{EXIT\PYGZus{}FAILURE}\PYG{p}{)}\PYG{p}{;}
\PYG{+w}{        }\PYG{p}{\PYGZcb{}}

\PYG{+w}{        }\PYG{k}{if}\PYG{+w}{ }\PYG{p}{(}\PYG{n}{pid}\PYG{+w}{ }\PYG{o}{=}\PYG{o}{=}\PYG{+w}{ }\PYG{l+m+mi}{0}\PYG{p}{)}\PYG{+w}{ }\PYG{p}{\PYGZob{}}
\PYG{+w}{            }\PYG{n}{pthread\PYGZus{}t}\PYG{+w}{ }\PYG{n}{tids}\PYG{p}{[}\PYG{n}{NTHREADS}\PYG{p}{]}\PYG{p}{;}

\PYG{+w}{            }\PYG{k}{for}\PYG{+w}{ }\PYG{p}{(}\PYG{k+kt}{int}\PYG{+w}{ }\PYG{n}{k}\PYG{+w}{ }\PYG{o}{=}\PYG{+w}{ }\PYG{l+m+mi}{0}\PYG{p}{;}\PYG{+w}{ }\PYG{n}{k}\PYG{+w}{ }\PYG{o}{\PYGZlt{}}\PYG{+w}{ }\PYG{n}{NTHREADS}\PYG{p}{;}\PYG{+w}{ }\PYG{o}{+}\PYG{o}{+}\PYG{n}{k}\PYG{p}{)}\PYG{+w}{ }\PYG{p}{\PYGZob{}}
\PYG{+w}{                }\PYG{k}{if}\PYG{+w}{ }\PYG{p}{(}\PYG{p}{(}\PYG{n}{result}\PYG{+w}{ }\PYG{o}{=}\PYG{+w}{ }\PYG{n}{pthread\PYGZus{}create}\PYG{p}{(}\PYG{o}{\PYGZam{}}\PYG{n}{tids}\PYG{p}{[}\PYG{n}{k}\PYG{p}{]}\PYG{p}{,}\PYG{+w}{ }\PYG{n+nb}{NULL}\PYG{p}{,}\PYG{+w}{ }\PYG{n}{thread\PYGZus{}proc}\PYG{p}{,}\PYG{+w}{ }\PYG{n+nb}{NULL}\PYG{p}{)}\PYG{p}{)}\PYG{+w}{ }\PYG{o}{!}\PYG{o}{=}\PYG{+w}{ }\PYG{l+m+mi}{0}\PYG{p}{)}\PYG{+w}{ }\PYG{p}{\PYGZob{}}
\PYG{+w}{                    }\PYG{n}{fprintf}\PYG{p}{(}\PYG{n}{stderr}\PYG{p}{,}\PYG{+w}{ }\PYG{l+s}{\PYGZdq{}}\PYG{l+s}{pthread\PYGZus{}create: \PYGZpc{}s}\PYG{l+s+se}{\PYGZbs{}n}\PYG{l+s}{\PYGZdq{}}\PYG{p}{,}\PYG{+w}{ }\PYG{n}{strerror}\PYG{p}{(}\PYG{n}{result}\PYG{p}{)}\PYG{p}{)}\PYG{p}{;}
\PYG{+w}{                    }\PYG{n}{exit}\PYG{p}{(}\PYG{n}{EXIT\PYGZus{}FAILURE}\PYG{p}{)}\PYG{p}{;}
\PYG{+w}{                }\PYG{p}{\PYGZcb{}}
\PYG{+w}{            }\PYG{p}{\PYGZcb{}}

\PYG{+w}{            }\PYG{k}{for}\PYG{+w}{ }\PYG{p}{(}\PYG{k+kt}{int}\PYG{+w}{ }\PYG{n}{k}\PYG{+w}{ }\PYG{o}{=}\PYG{+w}{ }\PYG{l+m+mi}{0}\PYG{p}{;}\PYG{+w}{ }\PYG{n}{k}\PYG{+w}{ }\PYG{o}{\PYGZlt{}}\PYG{+w}{ }\PYG{n}{NTHREADS}\PYG{p}{;}\PYG{+w}{ }\PYG{o}{+}\PYG{o}{+}\PYG{n}{k}\PYG{p}{)}
\PYG{+w}{                }\PYG{n}{pthread\PYGZus{}join}\PYG{p}{(}\PYG{n}{tids}\PYG{p}{[}\PYG{n}{k}\PYG{p}{]}\PYG{p}{,}\PYG{+w}{ }\PYG{n+nb}{NULL}\PYG{p}{)}\PYG{p}{;}

\PYG{+w}{            }\PYG{n}{\PYGZus{}exit}\PYG{p}{(}\PYG{l+m+mi}{0}\PYG{p}{)}\PYG{p}{;}
\PYG{+w}{        }\PYG{p}{\PYGZcb{}}
\PYG{+w}{        }\PYG{k}{else}\PYG{+w}{ }\PYG{p}{\PYGZob{}}
\PYG{+w}{            }\PYG{n}{printf}\PYG{p}{(}\PYG{l+s}{\PYGZdq{}}\PYG{l+s}{Child PID = \PYGZpc{}jd}\PYG{l+s+se}{\PYGZbs{}n}\PYG{l+s}{\PYGZdq{}}\PYG{p}{,}\PYG{+w}{ }\PYG{p}{(}\PYG{k+kt}{intmax\PYGZus{}t}\PYG{p}{)}\PYG{n}{pid}\PYG{p}{)}\PYG{p}{;}
\PYG{+w}{        }\PYG{p}{\PYGZcb{}}
\PYG{+w}{    }\PYG{p}{\PYGZcb{}}

\PYG{+w}{    }\PYG{k}{if}\PYG{+w}{ }\PYG{p}{(}\PYG{p}{(}\PYG{n}{fd}\PYG{+w}{ }\PYG{o}{=}\PYG{+w}{ }\PYG{n}{open}\PYG{p}{(}\PYG{l+s}{\PYGZdq{}}\PYG{l+s}{test\PYGZhy{}driver}\PYG{l+s}{\PYGZdq{}}\PYG{p}{,}\PYG{+w}{ }\PYG{n}{O\PYGZus{}RDONLY}\PYG{p}{)}\PYG{p}{)}\PYG{+w}{ }\PYG{o}{=}\PYG{o}{=}\PYG{+w}{ }\PYG{l+m+mi}{\PYGZhy{}1}\PYG{p}{)}
\PYG{+w}{        }\PYG{n}{exit\PYGZus{}sys}\PYG{p}{(}\PYG{l+s}{\PYGZdq{}}\PYG{l+s}{open}\PYG{l+s}{\PYGZdq{}}\PYG{p}{)}\PYG{p}{;}

\PYG{+w}{    }\PYG{k}{if}\PYG{+w}{ }\PYG{p}{(}\PYG{n}{ioctl}\PYG{p}{(}\PYG{n}{fd}\PYG{p}{,}\PYG{+w}{ }\PYG{n}{IOC\PYGZus{}TEST}\PYG{p}{)}\PYG{+w}{ }\PYG{o}{=}\PYG{o}{=}\PYG{+w}{ }\PYG{l+m+mi}{\PYGZhy{}1}\PYG{p}{)}
\PYG{+w}{        }\PYG{n}{exit\PYGZus{}sys}\PYG{p}{(}\PYG{l+s}{\PYGZdq{}}\PYG{l+s}{ioctl}\PYG{l+s}{\PYGZdq{}}\PYG{p}{)}\PYG{p}{;}

\PYG{+w}{    }\PYG{n}{close}\PYG{p}{(}\PYG{n}{fd}\PYG{p}{)}\PYG{p}{;}

\PYG{+w}{    }\PYG{k}{for}\PYG{+w}{ }\PYG{p}{(}\PYG{k+kt}{int}\PYG{+w}{ }\PYG{n}{i}\PYG{+w}{ }\PYG{o}{=}\PYG{+w}{ }\PYG{l+m+mi}{0}\PYG{p}{;}\PYG{+w}{ }\PYG{n}{i}\PYG{+w}{ }\PYG{o}{\PYGZlt{}}\PYG{+w}{ }\PYG{n}{NCHILDS}\PYG{p}{;}\PYG{+w}{ }\PYG{o}{+}\PYG{o}{+}\PYG{n}{i}\PYG{p}{)}
\PYG{+w}{        }\PYG{k}{if}\PYG{+w}{ }\PYG{p}{(}\PYG{n}{wait}\PYG{p}{(}\PYG{n+nb}{NULL}\PYG{p}{)}\PYG{+w}{ }\PYG{o}{=}\PYG{o}{=}\PYG{+w}{ }\PYG{l+m+mi}{\PYGZhy{}1}\PYG{p}{)}\PYG{+w}{ }\PYG{p}{\PYGZob{}}
\PYG{+w}{            }\PYG{n}{perror}\PYG{p}{(}\PYG{l+s}{\PYGZdq{}}\PYG{l+s}{wait}\PYG{l+s}{\PYGZdq{}}\PYG{p}{)}\PYG{p}{;}
\PYG{+w}{            }\PYG{n}{exit}\PYG{p}{(}\PYG{n}{EXIT\PYGZus{}FAILURE}\PYG{p}{)}\PYG{p}{;}
\PYG{+w}{        }\PYG{p}{\PYGZcb{}}

\PYG{+w}{    }\PYG{k}{return}\PYG{+w}{ }\PYG{l+m+mi}{0}\PYG{p}{;}
\PYG{p}{\PYGZcb{}}

\PYG{k+kt}{void}\PYG{+w}{ }\PYG{o}{*}\PYG{n+nf}{thread\PYGZus{}proc}\PYG{p}{(}\PYG{k+kt}{void}\PYG{+w}{ }\PYG{o}{*}\PYG{n}{param}\PYG{p}{)}
\PYG{p}{\PYGZob{}}
\PYG{+w}{    }\PYG{n}{sleep}\PYG{p}{(}\PYG{l+m+mi}{60}\PYG{p}{)}\PYG{p}{;}

\PYG{+w}{    }\PYG{k}{return}\PYG{+w}{ }\PYG{n+nb}{NULL}\PYG{p}{;}
\PYG{p}{\PYGZcb{}}

\PYG{k+kt}{void}\PYG{+w}{ }\PYG{n+nf}{exit\PYGZus{}sys}\PYG{p}{(}\PYG{k}{const}\PYG{+w}{ }\PYG{k+kt}{char}\PYG{+w}{ }\PYG{o}{*}\PYG{n}{msg}\PYG{p}{)}
\PYG{p}{\PYGZob{}}
\PYG{+w}{    }\PYG{n}{perror}\PYG{p}{(}\PYG{n}{msg}\PYG{p}{)}\PYG{p}{;}
\PYG{+w}{    }\PYG{n}{exit}\PYG{p}{(}\PYG{n}{EXIT\PYGZus{}FAILURE}\PYG{p}{)}\PYG{p}{;}
\PYG{p}{\PYGZcb{}}
\end{sphinxVerbatim}

\begin{sphinxadmonition}{note}{Not:}
\sphinxAtStartPar
Çekirdeğin uyguladığı RCU mekanizması tek bir yazan ve birden fazla okuyan akışın bulunduğu durumda
beklemeyi ortadan kaldırmaktadır. Ancak veri yapısına birden fazla yazanın olması durumunda yazan
tarafların bir kilit mekanizmasıyla senkronize edilmesi gerekir. İşte güncel çekirdekler hala
\sphinxcode{\sphinxupquote{task\_struct}} bağlı listelerine yazma için \sphinxstyleemphasis{tasklist\_lock} kilidini kullanmaktadır. \sphinxstyleemphasis{tasklist\_lock}
okuma yazma kilidi artık çekirdek modülleri ve aygıt sürücüler için export edilmemektedir.
\end{sphinxadmonition}


\section{PID Tahsisatı}
\label{\detokenize{processcontrolblock:pid-tahsisati}}
\sphinxAtStartPar
Şimdi de pid değerleriyle \sphinxcode{\sphinxupquote{task\_struct}} nesneleri arasındaki ilişkiyi ele alalım. Bilindiği gibi
kullanıcı modunda prosesler pid değerleriyle temsil edilmektedir. Örneğin pid parametresi alan tipik
POSIX fonksiyonları şunlardır: \sphinxstyleemphasis{kill}, \sphinxstyleemphasis{waitpid}, \sphinxstyleemphasis{getpgid}, \sphinxstyleemphasis{setpgid}, \sphinxstyleemphasis{getsid}, \sphinxstyleemphasis{getpriority},
\sphinxstyleemphasis{setpriority}.

\sphinxAtStartPar
O halde çekirdeğin pid değerinden hareketle o pid değerine ilişkin \sphinxcode{\sphinxupquote{task\_struct}} nesnesini hızlı bir
biçimde elde etmesi gerekmektedir. Bu bağlamda çekirdek geliştiricisinin aşağıdaki iki sorunun yanıtını
biliyor olması gerekir:
\begin{enumerate}
\sphinxsetlistlabels{\arabic}{enumi}{enumii}{}{.}%
\item {} 
\sphinxAtStartPar
Yeni bir proses ya da thread yaratıldığında çekirdek o proses ya da thread’e ilişkin pid değerini
nasıl üretmektedir?

\item {} 
\sphinxAtStartPar
Çekirdek belli bir pid değerine ilişkin proses ya da thread’in \sphinxcode{\sphinxupquote{task\_struct}} nesnesine nasıl
ulaşmaktadır?

\end{enumerate}


\subsection{PID Üretim Mekanizmasının Tarihsel Gelişimi}
\label{\detokenize{processcontrolblock:pid-uretim-mekanizmasinin-tarihsel-gelisimi}}
\sphinxAtStartPar
Bir \sphinxcode{\sphinxupquote{task\_struct}} nesnesi çekirdek tarafından yaratıldığında ona atanacak pid değeri eskiden oldukça
basit bir biçimde belirleniyordu. Linux 0.01 versiyonundaki \sphinxstyleemphasis{find\_empty\_process} fonksiyonu bu bakımdan
çok ilkeldi:

\begin{sphinxVerbatim}[commandchars=\\\{\}]
\PYG{k+kt}{int}\PYG{+w}{ }\PYG{n+nf}{find\PYGZus{}empty\PYGZus{}process}\PYG{p}{(}\PYG{k+kt}{void}\PYG{p}{)}
\PYG{p}{\PYGZob{}}
\PYG{+w}{    }\PYG{k+kt}{int}\PYG{+w}{ }\PYG{n}{i}\PYG{p}{;}

\PYG{+w}{    }\PYG{n+nl}{repeat}\PYG{p}{:}
\PYG{+w}{        }\PYG{k}{if}\PYG{+w}{ }\PYG{p}{(}\PYG{p}{(}\PYG{o}{+}\PYG{o}{+}\PYG{n}{last\PYGZus{}pid}\PYG{p}{)}\PYG{+w}{ }\PYG{o}{\PYGZlt{}}\PYG{+w}{ }\PYG{l+m+mi}{0}\PYG{p}{)}\PYG{+w}{ }\PYG{n}{last\PYGZus{}pid}\PYG{+w}{ }\PYG{o}{=}\PYG{+w}{ }\PYG{l+m+mi}{1}\PYG{p}{;}
\PYG{+w}{        }\PYG{k}{for}\PYG{+w}{ }\PYG{p}{(}\PYG{n}{i}\PYG{+w}{ }\PYG{o}{=}\PYG{+w}{ }\PYG{l+m+mi}{0}\PYG{p}{;}\PYG{+w}{ }\PYG{n}{i}\PYG{+w}{ }\PYG{o}{\PYGZlt{}}\PYG{+w}{ }\PYG{n}{NR\PYGZus{}TASKS}\PYG{p}{;}\PYG{+w}{ }\PYG{n}{i}\PYG{o}{+}\PYG{o}{+}\PYG{p}{)}
\PYG{+w}{            }\PYG{k}{if}\PYG{+w}{ }\PYG{p}{(}\PYG{n}{task}\PYG{p}{[}\PYG{n}{i}\PYG{p}{]}\PYG{+w}{ }\PYG{o}{\PYGZam{}}\PYG{o}{\PYGZam{}}\PYG{+w}{ }\PYG{n}{task}\PYG{p}{[}\PYG{n}{i}\PYG{p}{]}\PYG{o}{\PYGZhy{}}\PYG{o}{\PYGZgt{}}\PYG{n}{pid}\PYG{+w}{ }\PYG{o}{=}\PYG{o}{=}\PYG{+w}{ }\PYG{n}{last\PYGZus{}pid}\PYG{p}{)}\PYG{+w}{ }\PYG{k}{goto}\PYG{+w}{ }\PYG{n}{repeat}\PYG{p}{;}
\PYG{+w}{    }\PYG{k}{for}\PYG{+w}{ }\PYG{p}{(}\PYG{n}{i}\PYG{+w}{ }\PYG{o}{=}\PYG{+w}{ }\PYG{l+m+mi}{1}\PYG{p}{;}\PYG{+w}{ }\PYG{n}{i}\PYG{+w}{ }\PYG{o}{\PYGZlt{}}\PYG{+w}{ }\PYG{n}{NR\PYGZus{}TASKS}\PYG{p}{;}\PYG{+w}{ }\PYG{n}{i}\PYG{o}{+}\PYG{o}{+}\PYG{p}{)}
\PYG{+w}{        }\PYG{k}{if}\PYG{+w}{ }\PYG{p}{(}\PYG{o}{!}\PYG{n}{task}\PYG{p}{[}\PYG{n}{i}\PYG{p}{]}\PYG{p}{)}
\PYG{+w}{            }\PYG{k}{return}\PYG{+w}{ }\PYG{n}{i}\PYG{p}{;}

\PYG{+w}{    }\PYG{k}{return}\PYG{+w}{ }\PYG{o}{\PYGZhy{}}\PYG{n}{EAGAIN}\PYG{p}{;}
\PYG{p}{\PYGZcb{}}
\end{sphinxVerbatim}

\sphinxAtStartPar
Bu ilk versiyonda gördüğünüz gibi sistemde tahsis edilebilecek en fazla \sphinxcode{\sphinxupquote{task\_struct}} nesnesi 64 kadardı
(\sphinxcode{\sphinxupquote{NR\_TASKS = 64}}). Zaten bu versiyonda 64 tane \sphinxcode{\sphinxupquote{task\_struct}} yapısı baştan statik olarak tahsis
edilmişti:

\begin{sphinxVerbatim}[commandchars=\\\{\}]
\PYG{k}{struct}\PYG{+w}{ }\PYG{n+nc}{task\PYGZus{}struct}\PYG{+w}{ }\PYG{o}{*}\PYG{n}{task}\PYG{p}{[}\PYG{n}{NR\PYGZus{}TASKS}\PYG{p}{]}\PYG{+w}{ }\PYG{o}{=}\PYG{+w}{ }\PYG{p}{\PYGZob{}}\PYG{o}{\PYGZam{}}\PYG{p}{(}\PYG{n}{init\PYGZus{}task}\PYG{p}{.}\PYG{n}{task}\PYG{p}{)}\PYG{p}{,}\PYG{+w}{ }\PYG{p}{\PYGZcb{}}\PYG{p}{;}
\end{sphinxVerbatim}

\sphinxAtStartPar
2.4 çekirdeğinde pid numarasının belirlenmesi işlemi belli bir değerden itibaren aday olan pid değerlerinin
herhangi bir \sphinxcode{\sphinxupquote{task\_struct}} nesnesi tarafından kullanılıp kullanılmadığına bakılarak yapılmıştır. Bu
versiyondaki \sphinxstyleemphasis{get\_pid} fonksiyonu tüm \sphinxcode{\sphinxupquote{task\_struct}} nesnelerini dolaşarak boş bir pid değerini doğrusal
aramayla tespit etmeye çalışmaktadır.

\sphinxAtStartPar
pid araması 2.6 çekirdekleriyle birlikte bitmap veri yapısı kullanılarak elde edilmeye başlanmıştır.
Bitmap veri yapısı bitleri tutan bir veri yapısıdır. Modern işlemcilerde belli bir adresten itibaren 0
olan ilk bitin yerini veren makine komutları bulunmaktadır. Bu veri yapısı kullanılarak boş bir pid
değeri elde eden fonksiyondaki kilit satır şudur:

\begin{sphinxVerbatim}[commandchars=\\\{\}]
\PYG{n}{offset}\PYG{+w}{ }\PYG{o}{=}\PYG{+w}{ }\PYG{n}{find\PYGZus{}next\PYGZus{}offset}\PYG{p}{(}\PYG{n}{map}\PYG{p}{,}\PYG{+w}{ }\PYG{n}{offset}\PYG{p}{)}\PYG{p}{;}
\end{sphinxVerbatim}

\sphinxAtStartPar
\sphinxstyleemphasis{find\_next\_offset} fonksiyonu bitmap’teki ilk boş olan 0 bitinin offset değerini bulan özel makine komutu
kullanılarak yazılmıştır. Kursumuzda bitmap veri yapısını ileride ayrı bir başlık altında inceleyeceğiz.

\sphinxAtStartPar
Çekirdeğin 4.15 versiyonu ile pid tahsisatında bitmap kullanımı bırakılarak radix ağaçları kullanılmaya
başlanmıştır. Güncel çekirdekler ise boş pid değerinin üretimi için artık radix ağaçlarının özel bir
biçimi olan \sphinxstyleemphasis{XArray} veri yapısını kullanmaktadır. Güncel çekirdeklerde boş pid değerini bu karmaşık
yöntemle elde eden \sphinxcode{\sphinxupquote{alloc\_pid}} isimli yüksek seviyeli bir fonksiyon bulunmaktadır:

\begin{sphinxVerbatim}[commandchars=\\\{\}]
\PYG{k}{struct}\PYG{+w}{ }\PYG{n+nc}{pid}\PYG{+w}{ }\PYG{o}{*}\PYG{n}{alloc\PYGZus{}pid}\PYG{p}{(}\PYG{k}{struct}\PYG{+w}{ }\PYG{n+nc}{pid\PYGZus{}namespace}\PYG{+w}{ }\PYG{o}{*}\PYG{n}{ns}\PYG{p}{,}\PYG{+w}{ }\PYG{k+kt}{pid\PYGZus{}t}\PYG{+w}{ }\PYG{o}{*}\PYG{n}{set\PYGZus{}tid}\PYG{p}{,}\PYG{+w}{ }\PYG{k+kt}{size\PYGZus{}t}\PYG{+w}{ }\PYG{n}{set\PYGZus{}tid\PYGZus{}size}\PYG{p}{)}\PYG{p}{;}
\end{sphinxVerbatim}


\subsection{Maksimum PID Değeri}
\label{\detokenize{processcontrolblock:maksimum-pid-degeri}}
\sphinxAtStartPar
Linux çekirdekleri maksimum pid değeri için bir üst limit kullanmaktadır. 2.6 ve günümüze kadarki
çekirdeklerde \sphinxcode{\sphinxupquote{PID\_MAX\_LIMIT}} sembolik sabiti şöyle bildirilmiştir:

\begin{sphinxVerbatim}[commandchars=\\\{\}]
\PYG{c+cp}{\PYGZsh{}}\PYG{c+cp}{define PID\PYGZus{}MAX\PYGZus{}LIMIT (IS\PYGZus{}ENABLED(CONFIG\PYGZus{}BASE\PYGZus{}SMALL) ? PAGE\PYGZus{}SIZE * 8 : \PYGZbs{}}
\PYG{c+cp}{        (sizeof(long) \PYGZgt{} 4 ? 4 * 1024 * 1024 : PID\PYGZus{}MAX\PYGZus{}DEFAULT))}
\end{sphinxVerbatim}

\sphinxAtStartPar
Görüldüğü gibi \sphinxcode{\sphinxupquote{CONFIG\_BASE\_SMALL}} konfigürasyon parametresi \sphinxcode{\sphinxupquote{n}} yapılmış ve 64 bit sistemlerde bu
limit 4 × 1024 × 1024 = 4.194.304 değerini belirtmektedir. 32 bit sistemlerde ise bu değer \sphinxcode{\sphinxupquote{PID\_MAX\_DEFAULT}}
yani 32.768’dir.

\sphinxAtStartPar
Bugün kullandığımız 64 bitlik işlemcilerdeki Linux sistemlerinde maksimum pid değeri 4.194.304, 32 bit
sistemlerde ise 32.767 biçimindedir. Bu değer \sphinxstyleemphasis{proc} dosya sisteminde \sphinxcode{\sphinxupquote{/proc/sys/kernel/pid\_max}}
dosyasında da belirtilmektedir. Programcı isterse \sphinxstyleemphasis{sysctl} komutu yoluyla sistem çalışırken bu değeri
düşürebilir ancak yükseltemez.

\sphinxAtStartPar
Sisteminizde tahsis edilebilecek maksimum \sphinxcode{\sphinxupquote{task\_struct}} nesnelerinin sayısını (yani maksimum thread
sayısını) \sphinxcode{\sphinxupquote{/proc/sys/kernel/threads\sphinxhyphen{}max}} dosyasından öğrenebilirsiniz. Bu değer sistem çalışırken de
değiştirilebilir:

\begin{sphinxVerbatim}[commandchars=\\\{\}]
\PYGZdl{}\PYG{+w}{ }\PYG{n+nb}{echo}\PYG{+w}{ }\PYG{l+m}{100000}\PYG{+w}{ }\PYG{p}{|}\PYG{+w}{ }sudo\PYG{+w}{ }tee\PYG{+w}{ }/proc/sys/kernel/threads\PYGZhy{}max
\end{sphinxVerbatim}


\section{PID’den \sphinxstyleliteralintitle{\sphinxupquote{task\_struct}} Nesnesine Hızlı Erişim}
\label{\detokenize{processcontrolblock:pid-den-task-struct-nesnesine-hizli-erisim}}
\sphinxAtStartPar
“Çekirdek belli bir pid değerine ilişkin \sphinxcode{\sphinxupquote{task\_struct}} adresini nasıl hızlı bir biçimde elde
etmektedir?” sorusuna yanıt verelim. Düz mantıkla sistemdeki bütün \sphinxcode{\sphinxupquote{task\_struct}} nesneleri dolaşılarak
elde edilebilir. Ancak böyle bir arama çekirdek için çok yavaştır. Bu tür hızlı aramalar için Linux
çekirdeğinde genel olarak hash tabloları, bazen de arama ağaçları kullanılmaktadır.

\sphinxAtStartPar
Linux’un 2.6.24 versiyonuna kadar bunun için hash tabloları kullanılıyordu. Ancak bu versiyondan itibaren
bu amaçla \sphinxstyleemphasis{radix ağaçları} da işin içine sokulmuştur. Güncel çekirdeklerdeki arama sistemi radix
ağaçlarının özel bir biçimi olan \sphinxstyleemphasis{XArray} veri yapısı ve hash tablolarının hibrit bir biçimine benzemektedir.
Biz kursumuzun bu noktasında Linux çekirdeğinde hash tablolarının gerçekleştirimi üzerinde duracağız.

\sphinxAtStartPar
Çekirdeğin ilk 0.01 versiyonunda zaten en fazla 64 tane \sphinxcode{\sphinxupquote{task\_struct}} nesnesi oluşturulabiliyordu.
Bu versiyonda pid değerine ilişkin \sphinxcode{\sphinxupquote{task\_struct}} nesnesinin bulunması için bu \sphinxcode{\sphinxupquote{task\_struct}}
nesnelerinin tutulduğu \sphinxcode{\sphinxupquote{task}} isimli global dizide sıralı arama yapılmıştır.

\sphinxAtStartPar
Çekirdeğin 2.2 versiyonlarında pid değerinden hareketle \sphinxcode{\sphinxupquote{task\_struct}} nesnesinin bulunması için hash
tablosu kullanılmıştır. Bu versiyonlarda bu işi yapan \sphinxstyleemphasis{find\_task\_by\_pid} fonksiyonu aşağıdaki gibi
yazılmıştır:

\begin{sphinxVerbatim}[commandchars=\\\{\}]
\PYG{k}{extern}\PYG{+w}{ }\PYG{n}{\PYGZus{}\PYGZus{}inline\PYGZus{}\PYGZus{}}\PYG{+w}{ }\PYG{k}{struct}\PYG{+w}{ }\PYG{n+nc}{task\PYGZus{}struct}\PYG{+w}{ }\PYG{o}{*}\PYG{n}{find\PYGZus{}task\PYGZus{}by\PYGZus{}pid}\PYG{p}{(}\PYG{k+kt}{int}\PYG{+w}{ }\PYG{n}{pid}\PYG{p}{)}
\PYG{p}{\PYGZob{}}
\PYG{+w}{    }\PYG{k}{struct}\PYG{+w}{ }\PYG{n+nc}{task\PYGZus{}struct}\PYG{+w}{ }\PYG{o}{*}\PYG{n}{p}\PYG{p}{,}\PYG{+w}{ }\PYG{o}{*}\PYG{o}{*}\PYG{n}{htable}\PYG{+w}{ }\PYG{o}{=}\PYG{+w}{ }\PYG{o}{\PYGZam{}}\PYG{n}{pidhash}\PYG{p}{[}\PYG{n}{pid\PYGZus{}hashfn}\PYG{p}{(}\PYG{n}{pid}\PYG{p}{)}\PYG{p}{]}\PYG{p}{;}

\PYG{+w}{    }\PYG{k}{for}\PYG{+w}{ }\PYG{p}{(}\PYG{n}{p}\PYG{+w}{ }\PYG{o}{=}\PYG{+w}{ }\PYG{o}{*}\PYG{n}{htable}\PYG{p}{;}\PYG{+w}{ }\PYG{n}{p}\PYG{+w}{ }\PYG{o}{\PYGZam{}}\PYG{o}{\PYGZam{}}\PYG{+w}{ }\PYG{n}{p}\PYG{o}{\PYGZhy{}}\PYG{o}{\PYGZgt{}}\PYG{n}{pid}\PYG{+w}{ }\PYG{o}{!}\PYG{o}{=}\PYG{+w}{ }\PYG{n}{pid}\PYG{p}{;}\PYG{+w}{ }\PYG{n}{p}\PYG{+w}{ }\PYG{o}{=}\PYG{+w}{ }\PYG{n}{p}\PYG{o}{\PYGZhy{}}\PYG{o}{\PYGZgt{}}\PYG{n}{pidhash\PYGZus{}next}\PYG{p}{)}
\PYG{+w}{        }\PYG{p}{;}

\PYG{+w}{    }\PYG{k}{return}\PYG{+w}{ }\PYG{n}{p}\PYG{p}{;}
\PYG{p}{\PYGZcb{}}
\end{sphinxVerbatim}

\sphinxAtStartPar
Bu versiyonlarda henüz daha sonra açıkladığımız \sphinxcode{\sphinxupquote{hlist}} hash tablosu fonksiyonları çekirdekte
bulunmuyordu. \sphinxstyleemphasis{find\_task\_by\_pid} fonksiyonunda önce pid değeri \sphinxstyleemphasis{pid\_hashfn} isimli bir hash
fonksiyonuna sokularak bir indeks değeri elde edilmiş, sonra da \sphinxcode{\sphinxupquote{pidhash}} dizisinin bu indeksteki
bağlı liste zincirinde arama yapılmıştır. \sphinxcode{\sphinxupquote{pidhash}} dizisi şöyle tanımlanmıştır:

\begin{sphinxVerbatim}[commandchars=\\\{\}]
\PYG{k}{struct}\PYG{+w}{ }\PYG{n+nc}{task\PYGZus{}struct}\PYG{+w}{ }\PYG{o}{*}\PYG{n}{pidhash}\PYG{p}{[}\PYG{n}{PIDHASH\PYGZus{}SZ}\PYG{p}{]}\PYG{p}{;}
\end{sphinxVerbatim}

\sphinxAtStartPar
Görüldüğü gibi hash tablosundaki zincirler doğrudan \sphinxcode{\sphinxupquote{task\_struct}} nesnelerini tutmaktadır. Bu
versiyonlarda bu zincirler için \sphinxcode{\sphinxupquote{task\_struct}} içerisinde iki link elemanı bulunduruluyordu:

\begin{sphinxVerbatim}[commandchars=\\\{\}]
\PYG{k}{struct}\PYG{+w}{ }\PYG{n+nc}{task\PYGZus{}struct}\PYG{+w}{ }\PYG{p}{\PYGZob{}}
\PYG{+w}{    }\PYG{c+cm}{/* ... */}
\PYG{+w}{    }\PYG{k}{struct}\PYG{+w}{ }\PYG{n+nc}{task\PYGZus{}struct}\PYG{+w}{  }\PYG{o}{*}\PYG{n}{pidhash\PYGZus{}next}\PYG{p}{;}
\PYG{+w}{    }\PYG{k}{struct}\PYG{+w}{ }\PYG{n+nc}{task\PYGZus{}struct}\PYG{+w}{ }\PYG{o}{*}\PYG{o}{*}\PYG{n}{pidhash\PYGZus{}pprev}\PYG{p}{;}
\PYG{+w}{    }\PYG{c+cm}{/* ... */}
\PYG{p}{\PYGZcb{}}\PYG{p}{;}
\end{sphinxVerbatim}

\sphinxAtStartPar
Çekirdeğin 2.4 versiyonunda da algoritmada ve yukarıdaki fonksiyonda bir değişiklik yapılmamıştır.
Yani yine bir hash tablosu eşliğinde arama yapılmaktadır.


\subsection{2.6 Versiyonu: pid Nesnesi ve Bağlı Listeler}
\label{\detokenize{processcontrolblock:versiyonu-pid-nesnesi-ve-bagli-listeler}}
\sphinxAtStartPar
Çekirdeğin 2.6’lı versiyonlarında artık her farklı pid değeri için o pid değerine ilişkin bir
\sphinxcode{\sphinxupquote{pid}} nesnesi (\sphinxcode{\sphinxupquote{struct pid}} nesnesi) oluşturulmaya başlanmıştır. Bu versiyonlarda çekirdek her
farklı pid değeri için bir \sphinxcode{\sphinxupquote{pid}} nesnesi oluşturup bu \sphinxcode{\sphinxupquote{pid}} nesnelerini de hash tablolarında
saklamaktadır. Her \sphinxcode{\sphinxupquote{pid}} nesnesi de aşağıda açıklayacağımız gibi bir grup bağlı listenin kök
düğümlerini tutmaktadır. Bu versiyonlarda bir pid değerine ilişkin \sphinxcode{\sphinxupquote{task\_struct}} nesnesini bulmak
için çekirdek önce hash tablosundan o pid değerine ilişkin \sphinxcode{\sphinxupquote{pid}} nesnesini elde etmekte, sonra o
\sphinxcode{\sphinxupquote{pid}} nesnesinin içerisindeki bağlı listelerde arama yapmaktadır:

\sphinxincludegraphics[]{graphviz-d03a3cb3856ae563a636f462c10c108a85b22dc0.pdf}

\sphinxAtStartPar
2.6’lı versiyonlardaki \sphinxcode{\sphinxupquote{pid}} yapısı şöyleydi:

\begin{sphinxVerbatim}[commandchars=\\\{\}]
\PYG{k}{struct}\PYG{+w}{ }\PYG{n+nc}{pid}\PYG{+w}{ }\PYG{p}{\PYGZob{}}
\PYG{+w}{    }\PYG{n}{atomic\PYGZus{}t}\PYG{+w}{     }\PYG{n}{count}\PYG{p}{;}
\PYG{+w}{    }\PYG{k+kt}{unsigned}\PYG{+w}{ }\PYG{k+kt}{int}\PYG{+w}{ }\PYG{n}{level}\PYG{p}{;}
\PYG{+w}{    }\PYG{c+cm}{/* lists of tasks that use this pid */}
\PYG{+w}{    }\PYG{k}{struct}\PYG{+w}{ }\PYG{n+nc}{hlist\PYGZus{}head}\PYG{+w}{ }\PYG{n}{tasks}\PYG{p}{[}\PYG{n}{PIDTYPE\PYGZus{}MAX}\PYG{p}{]}\PYG{p}{;}
\PYG{+w}{    }\PYG{k}{struct}\PYG{+w}{ }\PYG{n+nc}{rcu\PYGZus{}head}\PYG{+w}{ }\PYG{n}{rcu}\PYG{p}{;}
\PYG{+w}{    }\PYG{k}{struct}\PYG{+w}{ }\PYG{n+nc}{upid}\PYG{+w}{ }\PYG{n}{numbers}\PYG{p}{[}\PYG{l+m+mi}{1}\PYG{p}{]}\PYG{p}{;}
\PYG{p}{\PYGZcb{}}\PYG{p}{;}
\end{sphinxVerbatim}

\sphinxAtStartPar
Bu versiyonlardaki durumu şöyle özetleyebiliriz:
\begin{enumerate}
\sphinxsetlistlabels{\arabic}{enumi}{enumii}{}{.}%
\item {} 
\sphinxAtStartPar
Çekirdek her pid için bir \sphinxcode{\sphinxupquote{pid}} nesnesi (\sphinxcode{\sphinxupquote{struct pid}} nesnesi) oluşturup onu bir hash
tablosunda saklamaktadır. Yani \sphinxcode{\sphinxupquote{pid}} nesneleri için bir hash tablosu kullanılmıştır.

\item {} 
\sphinxAtStartPar
\sphinxcode{\sphinxupquote{pid}} yapısının içerisindeki \sphinxcode{\sphinxupquote{tasks}} elemanı o pid değerine ilişkin bağlı liste zincirlerinin
kök düğümlerini tutmaktadır.

\item {} 
\sphinxAtStartPar
Sistem bir pid değerine ilişkin \sphinxcode{\sphinxupquote{task\_struct}} nesnesini elde etmek için önce o pid değerine
ilişkin \sphinxcode{\sphinxupquote{pid}} nesnesini hash tablosundan bulmakta, sonra onun içerisindeki ilgili bağlı listede
arama yapmaktadır.

\end{enumerate}

\sphinxAtStartPar
2.6 çekirdekleriyle birlikte Linux’a \sphinxstyleemphasis{isim alanları (name spaces)} kavramı da sokulmaya başlanmıştır.
Bu versiyonlar artık pid değerleri için bir isim alanı da kullanmaktadır. Pid isim alanı sayesinde
sanki çekirdek birden fazla pid dünyasına sahip gibi bir etki oluşturulmaktadır. Pid isim alanları iç
içe de oluşturulabilmektedir. Artık bu versiyonlarla birlikte pid değerleri sistem genelinde tek değil
isim alanı genelinde tektir. Linux çekirdeğine eklenen bu isim alanları özelliği \sphinxstyleemphasis{docker} gibi container
teknolojilerinin gelişmesine katkı sağlamıştır.


\subsection{\sphinxstyleliteralintitle{\sphinxupquote{pid\_type}} Enum Değerleri ve Bağlı Listeler}
\label{\detokenize{processcontrolblock:pid-type-enum-degerleri-ve-bagli-listeler}}
\sphinxAtStartPar
\sphinxcode{\sphinxupquote{pid}} yapısındaki \sphinxcode{\sphinxupquote{tasks}} dizisinin \sphinxcode{\sphinxupquote{PIDTYPE\_MAX}} kadar elemana sahip olduğunu görüyorsunuz.
\sphinxcode{\sphinxupquote{PIDTYPE\_MAX}} aşağıdaki gibi bir \sphinxcode{\sphinxupquote{enum}} türünün elemanıdır:

\begin{sphinxVerbatim}[commandchars=\\\{\}]
\PYG{c+cm}{/* 2.6\PYGZsq{}lı versiyonlar */}
\PYG{k}{enum}\PYG{+w}{ }\PYG{n}{pid\PYGZus{}type}\PYG{+w}{ }\PYG{p}{\PYGZob{}}
\PYG{+w}{    }\PYG{n}{PIDTYPE\PYGZus{}PID}\PYG{p}{,}
\PYG{+w}{    }\PYG{n}{PIDTYPE\PYGZus{}PGID}\PYG{p}{,}
\PYG{+w}{    }\PYG{n}{PIDTYPE\PYGZus{}SID}\PYG{p}{,}
\PYG{+w}{    }\PYG{n}{PIDTYPE\PYGZus{}MAX}\PYG{+w}{   }\PYG{c+cm}{/* = 3 */}
\PYG{p}{\PYGZcb{}}\PYG{p}{;}

\PYG{c+cm}{/* Güncel versiyonlar */}
\PYG{k}{enum}\PYG{+w}{ }\PYG{n}{pid\PYGZus{}type}\PYG{+w}{ }\PYG{p}{\PYGZob{}}
\PYG{+w}{    }\PYG{n}{PIDTYPE\PYGZus{}PID}\PYG{p}{,}
\PYG{+w}{    }\PYG{n}{PIDTYPE\PYGZus{}TGID}\PYG{p}{,}
\PYG{+w}{    }\PYG{n}{PIDTYPE\PYGZus{}PGID}\PYG{p}{,}
\PYG{+w}{    }\PYG{n}{PIDTYPE\PYGZus{}SID}\PYG{p}{,}
\PYG{+w}{    }\PYG{n}{PIDTYPE\PYGZus{}MAX}\PYG{+w}{   }\PYG{c+cm}{/* = 4 */}
\PYG{p}{\PYGZcb{}}\PYG{p}{;}
\end{sphinxVerbatim}

\sphinxAtStartPar
Daha sonra \sphinxcode{\sphinxupquote{PIDTYPE\_TGID}} ile temsil edilen bir bağlı listenin daha eklendiğine dikkat ediniz.
Peki pid değerini saklamak için neden birden fazla bağlı liste kullanılmaktadır? İşte 2.6’lı
çekirdeklerle birlikte bu tarz aramalarda hız kazancı sağlamak için tasarım değiştirilmiştir.
Aşağıdaki diyagram her listenin hangi \sphinxcode{\sphinxupquote{task\_struct}} nesnelerini tuttuğunu göstermektedir:

\sphinxincludegraphics[]{graphviz-d2c495fb313135c1c3a5d8e6e313d7110ba3a50c.pdf}

\sphinxAtStartPar
Bu bağlı listeleri daha ayrıntılı açıklayalım:
\begin{itemize}
\item {} 
\sphinxAtStartPar
\sphinxstylestrong{PIDTYPE\_PID:} Bu bağlı liste zincirinde aslında tek bir eleman vardır. Amacımız yalnızca belli
bir pid değerine ilişkin \sphinxcode{\sphinxupquote{task\_struct}} nesnesinin adresinin bulunmasıysa hemen bu listenin ilk
elemanını alabiliriz.

\item {} 
\sphinxAtStartPar
\sphinxstylestrong{PIDTYPE\_TGID:} Bu bağlı liste belli bir prosesin thread’lerini dolaşmak için kullanılmaktadır.
Bu bağlı listedeki tüm \sphinxcode{\sphinxupquote{task\_struct}} nesneleri aynı prosesin thread’lerini oluşturmaktadır. Yani
biz çekirdeğe bir prosesin ana thread’ine ilişkin bir pid değeri verdiğimizde çekirdek bu bağlı
liste zincirini dolaşarak o prosesin tüm thread’lerinin \sphinxcode{\sphinxupquote{task\_struct}} adreslerini elde
edebilmektedir.

\item {} 
\sphinxAtStartPar
\sphinxstylestrong{PIDTYPE\_PGID:} Aynı proses grubuna ilişkin proseslerin ana thread’lerinin \sphinxcode{\sphinxupquote{task\_struct}}
nesnelerini tutmaktadır. Yani aranan pid değeri bir proses grup liderine ilişkin pid belirtiyorsa bu
zincirde o proses grubundaki tüm proseslerin ana thread’lerinin \sphinxcode{\sphinxupquote{task\_struct}} nesnelerinin
adresleri bulunmaktadır.

\item {} 
\sphinxAtStartPar
\sphinxstylestrong{PIDTYPE\_SID:} Belli bir oturuma (session) ilişkin tüm proseslerin (onların ana thread’lerinin)
\sphinxcode{\sphinxupquote{task\_struct}} nesnelerinin adreslerini tutmaktadır.

\end{itemize}

\sphinxAtStartPar
2.6 versiyonlarından önce tek bir hash tablosu vardı. Bu nedenle proses grubundaki proseslerin
\sphinxcode{\sphinxupquote{task\_struct}} nesneleri ancak tüm \sphinxcode{\sphinxupquote{task\_struct}} nesneleri dolaşılarak elde edilebiliyordu.
Halbuki 2.6’daki bu tasarımla bu biçimdeki dolaşımlar bu hash tablosu ve bağlı listeler yoluyla çok
daha hızlı yapılabilmektedir.

\sphinxAtStartPar
Aşağıda 2.6 versiyonlarında pid değerinden \sphinxcode{\sphinxupquote{task\_struct}} nesnesine ulaşım çağrı zinciri
görülmektedir:

\sphinxincludegraphics[]{graphviz-c369e3ba32b095626ac71a139e75290841d8fb34.pdf}

\sphinxAtStartPar
Hash tablosunda pid nesnesini arayan \sphinxstyleemphasis{find\_pid\_ns} fonksiyonu, \sphinxcode{\sphinxupquote{upid}} yapı nesnelerini dolaşmakta
ve hem pid değerine hem de pid isim alanına birlikte bakmaktadır:

\begin{sphinxVerbatim}[commandchars=\\\{\}]
\PYG{k}{struct}\PYG{+w}{ }\PYG{n+nc}{upid}\PYG{+w}{ }\PYG{p}{\PYGZob{}}
\PYG{+w}{    }\PYG{k+kt}{int}\PYG{+w}{ }\PYG{n}{nr}\PYG{p}{;}
\PYG{+w}{    }\PYG{k}{struct}\PYG{+w}{ }\PYG{n+nc}{pid\PYGZus{}namespace}\PYG{+w}{ }\PYG{o}{*}\PYG{n}{ns}\PYG{p}{;}
\PYG{+w}{    }\PYG{k}{struct}\PYG{+w}{ }\PYG{n+nc}{hlist\PYGZus{}node}\PYG{+w}{ }\PYG{n}{pid\PYGZus{}chain}\PYG{p}{;}
\PYG{p}{\PYGZcb{}}\PYG{p}{;}

\PYG{k}{static}\PYG{+w}{ }\PYG{k}{struct}\PYG{+w}{ }\PYG{n+nc}{hlist\PYGZus{}head}\PYG{+w}{ }\PYG{o}{*}\PYG{n}{pid\PYGZus{}hash}\PYG{p}{;}\PYG{+w}{   }\PYG{c+cm}{/* global hash tablosu */}

\PYG{k}{struct}\PYG{+w}{ }\PYG{n+nc}{pid}\PYG{+w}{ }\PYG{o}{*}\PYG{n}{find\PYGZus{}pid\PYGZus{}ns}\PYG{p}{(}\PYG{k+kt}{int}\PYG{+w}{ }\PYG{n}{nr}\PYG{p}{,}\PYG{+w}{ }\PYG{k}{struct}\PYG{+w}{ }\PYG{n+nc}{pid\PYGZus{}namespace}\PYG{+w}{ }\PYG{o}{*}\PYG{n}{ns}\PYG{p}{)}
\PYG{p}{\PYGZob{}}
\PYG{+w}{    }\PYG{k}{struct}\PYG{+w}{ }\PYG{n+nc}{hlist\PYGZus{}node}\PYG{+w}{ }\PYG{o}{*}\PYG{n}{elem}\PYG{p}{;}
\PYG{+w}{    }\PYG{k}{struct}\PYG{+w}{ }\PYG{n+nc}{upid}\PYG{+w}{ }\PYG{o}{*}\PYG{n}{pnr}\PYG{p}{;}

\PYG{+w}{    }\PYG{n}{hlist\PYGZus{}for\PYGZus{}each\PYGZus{}entry\PYGZus{}rcu}\PYG{p}{(}\PYG{n}{pnr}\PYG{p}{,}\PYG{+w}{ }\PYG{n}{elem}\PYG{p}{,}
\PYG{+w}{            }\PYG{o}{\PYGZam{}}\PYG{n}{pid\PYGZus{}hash}\PYG{p}{[}\PYG{n}{pid\PYGZus{}hashfn}\PYG{p}{(}\PYG{n}{nr}\PYG{p}{,}\PYG{+w}{ }\PYG{n}{ns}\PYG{p}{)}\PYG{p}{]}\PYG{p}{,}\PYG{+w}{ }\PYG{n}{pid\PYGZus{}chain}\PYG{p}{)}
\PYG{+w}{        }\PYG{k}{if}\PYG{+w}{ }\PYG{p}{(}\PYG{n}{pnr}\PYG{o}{\PYGZhy{}}\PYG{o}{\PYGZgt{}}\PYG{n}{nr}\PYG{+w}{ }\PYG{o}{=}\PYG{o}{=}\PYG{+w}{ }\PYG{n}{nr}\PYG{+w}{ }\PYG{o}{\PYGZam{}}\PYG{o}{\PYGZam{}}\PYG{+w}{ }\PYG{n}{pnr}\PYG{o}{\PYGZhy{}}\PYG{o}{\PYGZgt{}}\PYG{n}{ns}\PYG{+w}{ }\PYG{o}{=}\PYG{o}{=}\PYG{+w}{ }\PYG{n}{ns}\PYG{p}{)}
\PYG{+w}{            }\PYG{k}{return}\PYG{+w}{ }\PYG{n}{container\PYGZus{}of}\PYG{p}{(}\PYG{n}{pnr}\PYG{p}{,}\PYG{+w}{ }\PYG{k}{struct}\PYG{+w}{ }\PYG{n+nc}{pid}\PYG{p}{,}
\PYG{+w}{                    }\PYG{n}{numbers}\PYG{p}{[}\PYG{n}{ns}\PYG{o}{\PYGZhy{}}\PYG{o}{\PYGZgt{}}\PYG{n}{level}\PYG{p}{]}\PYG{p}{)}\PYG{p}{;}
\PYG{+w}{    }\PYG{k}{return}\PYG{+w}{ }\PYG{n+nb}{NULL}\PYG{p}{;}
\PYG{p}{\PYGZcb{}}
\end{sphinxVerbatim}

\sphinxAtStartPar
Burada \sphinxcode{\sphinxupquote{hlist\_node}} bağının \sphinxcode{\sphinxupquote{upid}} nesnesinin içerisinde olduğuna, \sphinxcode{\sphinxupquote{upid}} nesnesinin de \sphinxcode{\sphinxupquote{pid}}
nesnesinin içerisinde olduğuna dikkat ediniz.


\subsection{4.20 Versiyonu: Hash Tablosundan XArray’e Geçiş}
\label{\detokenize{processcontrolblock:versiyonu-hash-tablosundan-xarray-e-gecis}}
\sphinxAtStartPar
Çekirdek versiyonu 4.20’lere geldiğinde yukarıdaki mekanizmada önemli değişiklikler yapılmıştır. En
önemli değişiklik pid ve isim alanından hareketle \sphinxcode{\sphinxupquote{pid}} nesnesinin bulunması işleminde hash tablosu
yerine \sphinxstyleemphasis{radix ağacının} kullanılmaya başlanmasıdır. Aslında 2.5 ile birlikte çekirdekte başka alt
sistemlerde radix ağacı kullanılmaya başlanmıştı. Sonra bu radix ağaçlarının 4.20 versiyonuyla birlikte
\sphinxstyleemphasis{XArray} isimli yeni bir gerçekleştirimi yapılmıştır.

\sphinxAtStartPar
4.20 ve sonrasında \sphinxcode{\sphinxupquote{pid}} yapısındaki bağlı listelerde bir farklılık yoktur. Yalnızca pid ve isim alanı
bilgisinden hareketle \sphinxcode{\sphinxupquote{pid}} yapı nesnesinin elde edilmesinde hash tablosu yerine radix ağacı (XArray
gerçekleştirimi) kullanılmaktadır.

\sphinxAtStartPar
Bu bağlamda pid mekanizmasında radix ağacının hash tablosuna tercih edilmesinin tipik nedenleri şunlardır:
\begin{itemize}
\item {} 
\sphinxAtStartPar
Radix ağacındaki arama hash tablolarından yavaş gözükse bile aslında pratikte durum tam böyle
değildir. Buradaki radix ağacı seyrek tutulmaktadır. Bu yüzden arama hızlı yapılır.

\item {} 
\sphinxAtStartPar
pid’leri artan sırayla gezmek, boş pid bulmak (en küçük kullanılmayan pid’i seçmek) çok daha
kolaydır.

\item {} 
\sphinxAtStartPar
Ağaçta yalnızca kullanılan düğümler tutulmaktadır. Bu da bellek kazancı sağlamaktadır.

\item {} 
\sphinxAtStartPar
pid sayısı arttıkça hash tablosunun performansı düşme eğilimi gösterdiği halde radix ağacının
performansı hash tablosuna kıyasla düşmez.

\item {} 
\sphinxAtStartPar
Radix ağacı yalnızca pid araması için değil, aynı zamanda yeni yaratılan \sphinxcode{\sphinxupquote{task\_struct}} nesneleri
için pid tahsis edilmesinde de kullanılabilmektedir.

\end{itemize}

\sphinxAtStartPar
Güncel çekirdeklerde hem pid aramaları hem de boş pid değerlerinin tespit edilmesi işlemleri bu XArray
gerçekleştirimi ile sağlanmaktadır.


\subsection{Radix Ağaçlarının Kullanılması}
\label{\detokenize{processcontrolblock:radix-agaclarinin-kullanilmasi}}
\sphinxAtStartPar
Biz radix arama ağacını ve XArray gerçekleştirimini daha sonra başka bir bölümde ele alacağız. Burada
yalnızca radix arama ağacının temel mekanizması üzerinde açıklamalar yapacağız. Radix arama ağaçları
için \sphinxstyleemphasis{dijital arama ağaçları (digital search tree)}, \sphinxstyleemphasis{önek arama ağacı (prefix search tree)},
\sphinxstyleemphasis{sıkıştırılmış arama ağaçları (compressed search trie)} gibi isimler de kullanılmaktadır.

\sphinxAtStartPar
Bu arama ağaçlarında düğümlerde anahtar olarak anahtarın tamamı değil onun ön kısımları kullanılmaktadır.
Böylece ağaç bir leksikografik sıralama ağacı gibi de kullanılabilmektedir. Tipik anlatım sayısal
değerlerden oluşan ve öneklerin bit belirttiği anahtarlar üzerinde yapılmaktadır. Ağacın kökünden
girilir. Her kademede bir bit daha anahtara eklenir.

\sphinxAtStartPar
Örneğin 4 bitlik şu sayıların radix ağacına yerleştirilmek istendiğini düşünelim:
\sphinxcode{\sphinxupquote{1100, 1010, 0101, 1011, 0010, 0111, 0000, 1000}}

\sphinxAtStartPar
Ağacın her adımında yüksek anlamlı bitten başlayarak solda 0, sağda 1 dallanması yapılmaktadır.
Aşağıdaki diyagram bu sekiz anahtarın ağaca yerleşimi sonucunu göstermektedir:

\sphinxincludegraphics[]{graphviz-21dd7abf9ab48251cef00bbd686af32dd57471af.pdf}

\sphinxAtStartPar
Peki bu ağaçta arama nasıl yapılır? İşte arama için kökten girilir. Her bit pozisyonu için bir aşağıya
inilir. Örneğin \sphinxcode{\sphinxupquote{0000}} değerini arayacak olalım: ilk bit 0 olduğu için kökten sola ineriz, ikinci bit
de 0 olduğu için sola devam ederiz, üçüncü bit de 0 olduğu için sola, dördüncü bit de 0 olduğu için
sola gideriz. Artık değeri buluruz. Radix ağaçlarını enlemesine (breadth\sphinxhyphen{}first gibi) dolaştığımızda
anahtarların sıralı bir biçimde elde edilebildiğine dikkat ediniz.


\subsection{Yapraklarda Değer Saklama Yöntemi}
\label{\detokenize{processcontrolblock:yapraklarda-deger-saklama-yontemi}}
\sphinxAtStartPar
Biz yukarıdaki anlatımda düğümlerde anahtarların tutulduğunu varsaydık. Aslında düğümlerde anahtarların
tutulmasına da gerek yoktur. Zaten her kademede bir bit ilerlendiğine göre arama yaprağa gelindiğinde
sonlandırılabilir. Aşağıdaki ağaçta bit\sphinxhyphen{}bit dallanma yapılmış ve yalnızca yapraklarda değerler
tutulmuştur:

\sphinxincludegraphics[]{graphviz-9bd2a3a59306e5b814afee2c126387db6d8a638b.pdf}

\sphinxAtStartPar
Görüldüğü gibi bu tasarımda hiç anahtar saklanmadan yaprak görülene kadar ilerlenirse ya ilgili anahtar
bulunmuş olur ya da bulunmamış olur. Peki her düğümde anahtarı tutmakla yalnızca yapraklarda değeri
tutmanın hangisi daha iyi bir yöntemdir?


\begin{savenotes}\sphinxattablestart
\sphinxthistablewithglobalstyle
\centering
\begin{tabular}[t]{\X{20}{100}\X{40}{100}\X{40}{100}}
\sphinxtoprule
\begin{varwidth}[t]{\sphinxcolwidth{1}{3}}
\sphinxstyletheadfamily \sphinxAtStartPar
Yöntem
\sphinxbeforeendvarwidth
\end{varwidth}%
&\begin{varwidth}[t]{\sphinxcolwidth{1}{3}}
\sphinxstyletheadfamily \sphinxAtStartPar
Avantajları
\sphinxbeforeendvarwidth
\end{varwidth}%
&\begin{varwidth}[t]{\sphinxcolwidth{1}{3}}
\sphinxstyletheadfamily \sphinxAtStartPar
Dezavantajları
\sphinxbeforeendvarwidth
\end{varwidth}%
\\
\sphinxmidrule
\sphinxtableatstartofbodyhook\begin{varwidth}[t]{\sphinxcolwidth{1}{3}}
\sphinxAtStartPar
\sphinxstylestrong{Her düğümde anahtar saklama}
\sphinxbeforeendvarwidth
\end{varwidth}%
&\begin{varwidth}[t]{\sphinxcolwidth{1}{3}}
\sphinxAtStartPar
Erken durabilme; ara düğümde tam eşleşme bulununca yaprağa kadar inmek gerekmez.
Prefix aramaları daha kolay.
\sphinxbeforeendvarwidth
\end{varwidth}%
&\begin{varwidth}[t]{\sphinxcolwidth{1}{3}}
\sphinxAtStartPar
Veri yapısı karmaşıklaşır; bellek kullanımı artar.
\sphinxbeforeendvarwidth
\end{varwidth}%
\\
\sphinxhline\begin{varwidth}[t]{\sphinxcolwidth{1}{3}}
\sphinxAtStartPar
\sphinxstylestrong{Yalnızca yapraklarda değer saklama}
\sphinxbeforeendvarwidth
\end{varwidth}%
&\begin{varwidth}[t]{\sphinxcolwidth{1}{3}}
\sphinxAtStartPar
Basit tasarım: iç düğümler yalnızca yönlendirme bilgisi tutar. Bellek daha düzenli kullanılır.
Ekleme/silme algoritmaları daha etkin.
\sphinxbeforeendvarwidth
\end{varwidth}%
&\begin{varwidth}[t]{\sphinxcolwidth{1}{3}}
\sphinxAtStartPar
Her zaman yaprağa kadar inmek gerekir. Prefix tabanlı aramalarda ek iş yapılması gerekir.
\sphinxbeforeendvarwidth
\end{varwidth}%
\\
\sphinxbottomrule
\end{tabular}
\sphinxtableafterendhook\par
\sphinxattableend\end{savenotes}

\begin{sphinxadmonition}{note}{Not:}
\sphinxAtStartPar
Linux’un klasik radix ve XArray gerçekleştirimlerinde değerler her zaman yapraklarda tutulmaktadır.
Ara düğümlerde anahtarlar tutulmamaktadır.
\end{sphinxadmonition}


\subsection{Dallanma Faktörü: n\sphinxhyphen{}li Ağaç}
\label{\detokenize{processcontrolblock:dallanma-faktoru-n-li-agac}}
\sphinxAtStartPar
Radix ağacının ikili ağaç olması zorunlu değildir. Dallanma bit bit yapıldığında yukarıdaki örnekte
olduğu gibi radix ağacı ikili ağaç durumunda olur. Ancak dallanma n bit n bit de yapılabilir. Bu
durumda ağaç da 2$^{\text{n}}$\sphinxhyphen{}li ağaç olacaktır.

\sphinxAtStartPar
Örneğin 32 bitlik pid değerlerini böyle bir ağaçta tutmak isteyelim ve dallanmayı 6 bit 6 bit yapalım.
6 bit \(\rightarrow\) 2$^{\text{6}}$ = 64 farklı dallanmaya yol açacaktır. Bu durumda ağacımızın her düğümünde 64
gösterici bulunacaktır ve yapraklara varabilmek için en fazla 6 kademe ilerlenecektir:

\sphinxincludegraphics[]{graphviz-3a2581653bc0601e9b800b2dd99212c992373625.pdf}

\sphinxAtStartPar
Güncel sürümlerde Linux’un pid araması için kullandığı XArray ağacı bu örnekteki gibi altışar bitlidir.

\sphinxAtStartPar
Hash tablosu ile XArray/Radix ağacı arasındaki avantaj ve dezavantajları karşılaştıralım:


\begin{savenotes}\sphinxattablestart
\sphinxthistablewithglobalstyle
\centering
\begin{tabular}[t]{\X{30}{100}\X{35}{100}\X{35}{100}}
\sphinxtoprule
\begin{varwidth}[t]{\sphinxcolwidth{1}{3}}
\sphinxstyletheadfamily \sphinxAtStartPar
Kriter
\sphinxbeforeendvarwidth
\end{varwidth}%
&\begin{varwidth}[t]{\sphinxcolwidth{1}{3}}
\sphinxstyletheadfamily \sphinxAtStartPar
Hash Tablosu
\sphinxbeforeendvarwidth
\end{varwidth}%
&\begin{varwidth}[t]{\sphinxcolwidth{1}{3}}
\sphinxstyletheadfamily \sphinxAtStartPar
XArray / Radix Ağacı
\sphinxbeforeendvarwidth
\end{varwidth}%
\\
\sphinxmidrule
\sphinxtableatstartofbodyhook\begin{varwidth}[t]{\sphinxcolwidth{1}{3}}
\sphinxAtStartPar
\sphinxstylestrong{Arama karmaşıklığı}
\sphinxbeforeendvarwidth
\end{varwidth}%
&\begin{varwidth}[t]{\sphinxcolwidth{1}{3}}
\sphinxAtStartPar
Ortalama O(1); çakışmada en kötü O(N)
\sphinxbeforeendvarwidth
\end{varwidth}%
&\begin{varwidth}[t]{\sphinxcolwidth{1}{3}}
\sphinxAtStartPar
O(log$_{\text{k}}$ N), tipik derinlik 2\textendash{}3
\sphinxbeforeendvarwidth
\end{varwidth}%
\\
\sphinxhline\begin{varwidth}[t]{\sphinxcolwidth{1}{3}}
\sphinxAtStartPar
\sphinxstylestrong{Bellek kullanımı}
\sphinxbeforeendvarwidth
\end{varwidth}%
&\begin{varwidth}[t]{\sphinxcolwidth{1}{3}}
\sphinxAtStartPar
Tablo önceden sabit boyutta ayrılır; seyrek pid’lerde boşa gider
\sphinxbeforeendvarwidth
\end{varwidth}%
&\begin{varwidth}[t]{\sphinxcolwidth{1}{3}}
\sphinxAtStartPar
Yalnızca kullanılan düğümler açılır; seyrek alanda çok verimli
\sphinxbeforeendvarwidth
\end{varwidth}%
\\
\sphinxhline\begin{varwidth}[t]{\sphinxcolwidth{1}{3}}
\sphinxAtStartPar
\sphinxstylestrong{Çakışma}
\sphinxbeforeendvarwidth
\end{varwidth}%
&\begin{varwidth}[t]{\sphinxcolwidth{1}{3}}
\sphinxAtStartPar
Mümkün; performansı düşürebilir
\sphinxbeforeendvarwidth
\end{varwidth}%
&\begin{varwidth}[t]{\sphinxcolwidth{1}{3}}
\sphinxAtStartPar
Yok; her pid tek bir yolu izler
\sphinxbeforeendvarwidth
\end{varwidth}%
\\
\sphinxhline\begin{varwidth}[t]{\sphinxcolwidth{1}{3}}
\sphinxAtStartPar
\sphinxstylestrong{Ölçeklenebilirlik}
\sphinxbeforeendvarwidth
\end{varwidth}%
&\begin{varwidth}[t]{\sphinxcolwidth{1}{3}}
\sphinxAtStartPar
Tablo boyutunu baştan iyi seçmek gerekir
\sphinxbeforeendvarwidth
\end{varwidth}%
&\begin{varwidth}[t]{\sphinxcolwidth{1}{3}}
\sphinxAtStartPar
pid alanı büyüse bile uyum sağlar
\sphinxbeforeendvarwidth
\end{varwidth}%
\\
\sphinxhline\begin{varwidth}[t]{\sphinxcolwidth{1}{3}}
\sphinxAtStartPar
\sphinxstylestrong{Sıralı gezinme}
\sphinxbeforeendvarwidth
\end{varwidth}%
&\begin{varwidth}[t]{\sphinxcolwidth{1}{3}}
\sphinxAtStartPar
Doğal sıralama yok; yavaş
\sphinxbeforeendvarwidth
\end{varwidth}%
&\begin{varwidth}[t]{\sphinxcolwidth{1}{3}}
\sphinxAtStartPar
Doğası gereği sıralı; çok hızlı
\sphinxbeforeendvarwidth
\end{varwidth}%
\\
\sphinxhline\begin{varwidth}[t]{\sphinxcolwidth{1}{3}}
\sphinxAtStartPar
\sphinxstylestrong{Boş pid bulma}
\sphinxbeforeendvarwidth
\end{varwidth}%
&\begin{varwidth}[t]{\sphinxcolwidth{1}{3}}
\sphinxAtStartPar
Ayrı veri yapısı gerektirir
\sphinxbeforeendvarwidth
\end{varwidth}%
&\begin{varwidth}[t]{\sphinxcolwidth{1}{3}}
\sphinxAtStartPar
Aynı ağaç hem arama hem tahsisat için kullanılır
\sphinxbeforeendvarwidth
\end{varwidth}%
\\
\sphinxbottomrule
\end{tabular}
\sphinxtableafterendhook\par
\sphinxattableend\end{savenotes}

\sphinxAtStartPar
Özetle: tekil arama için hash tablosu genellikle biraz daha hızlıdır; ancak büyük ve seyrek pid isim
alanlarında XArray çok daha verimli bellek kullanır. Sıralı gezinme ve boş pid bulma gerektirdiğinde
XArray üstün gelir. Güncel çekirdeklerde hem pid aramaları hem de boş pid değerlerinin tespit edilmesi
işlemleri bu XArray gerçekleştirimi ile sağlanmaktadır.


\subsection{Güncel Çekirdek: pid Araması (4.20+)}
\label{\detokenize{processcontrolblock:guncel-cekirdek-pid-aramasi-4-20}}
\sphinxAtStartPar
Çekirdeğin 4.20 versiyonundan itibaren her pid isim alanı için ayrı bir XArray ağacı oluşturulmaya
başlanmıştır. Anımsanacağı gibi 2.6’lı versiyonlarda hash tablolarına geçildiğinde tüm pid isim alanları
aynı hash tablosunda bulunuyordu. Halbuki güncel çekirdeklerde bir pid nesnesi aranacaksa o nesne belli
bir pid isim alanındaki XArray ağacında aranmaktadır.

\sphinxAtStartPar
4.20 ve sonrasında \sphinxcode{\sphinxupquote{task\_struct}} yapısında \sphinxcode{\sphinxupquote{thread\_pid}} elemanı da eklenmiştir:

\begin{sphinxVerbatim}[commandchars=\\\{\}]
\PYG{k}{struct}\PYG{+w}{ }\PYG{n+nc}{task\PYGZus{}struct}\PYG{+w}{ }\PYG{p}{\PYGZob{}}
\PYG{+w}{    }\PYG{c+cm}{/* ... */}
\PYG{+w}{    }\PYG{k+kt}{pid\PYGZus{}t}\PYG{+w}{        }\PYG{n}{pid}\PYG{p}{;}
\PYG{+w}{    }\PYG{k}{struct}\PYG{+w}{ }\PYG{n+nc}{pid}\PYG{+w}{  }\PYG{o}{*}\PYG{n}{thread\PYGZus{}pid}\PYG{p}{;}\PYG{+w}{   }\PYG{c+cm}{/* pid değerine ilişkin pid nesnesinin adresi */}
\PYG{+w}{    }\PYG{c+cm}{/* ... */}
\PYG{p}{\PYGZcb{}}\PYG{p}{;}
\end{sphinxVerbatim}

\sphinxAtStartPar
Güncel \sphinxcode{\sphinxupquote{pid}} yapısı şöyledir:

\begin{sphinxVerbatim}[commandchars=\\\{\}]
\PYG{k}{struct}\PYG{+w}{ }\PYG{n+nc}{pid}\PYG{+w}{ }\PYG{p}{\PYGZob{}}
\PYG{+w}{    }\PYG{n}{refcount\PYGZus{}t}\PYG{+w}{    }\PYG{n}{count}\PYG{p}{;}
\PYG{+w}{    }\PYG{k+kt}{unsigned}\PYG{+w}{ }\PYG{k+kt}{int}\PYG{+w}{  }\PYG{n}{level}\PYG{p}{;}
\PYG{+w}{    }\PYG{n}{spinlock\PYGZus{}t}\PYG{+w}{    }\PYG{n}{lock}\PYG{p}{;}
\PYG{+w}{    }\PYG{k}{struct}\PYG{+w}{ }\PYG{p}{\PYGZob{}}
\PYG{+w}{        }\PYG{n}{u64}\PYG{+w}{           }\PYG{n}{ino}\PYG{p}{;}
\PYG{+w}{        }\PYG{k}{struct}\PYG{+w}{ }\PYG{n+nc}{rb\PYGZus{}node}\PYG{+w}{ }\PYG{n}{pidfs\PYGZus{}node}\PYG{p}{;}
\PYG{+w}{        }\PYG{k}{struct}\PYG{+w}{ }\PYG{n+nc}{dentry}\PYG{+w}{ }\PYG{o}{*}\PYG{n}{stashed}\PYG{p}{;}
\PYG{+w}{        }\PYG{k}{struct}\PYG{+w}{ }\PYG{n+nc}{pidfs\PYGZus{}attr}\PYG{+w}{ }\PYG{o}{*}\PYG{n}{attr}\PYG{p}{;}
\PYG{+w}{    }\PYG{p}{\PYGZcb{}}\PYG{p}{;}
\PYG{+w}{    }\PYG{k}{struct}\PYG{+w}{ }\PYG{n+nc}{hlist\PYGZus{}head}\PYG{+w}{   }\PYG{n}{tasks}\PYG{p}{[}\PYG{n}{PIDTYPE\PYGZus{}MAX}\PYG{p}{]}\PYG{p}{;}
\PYG{+w}{    }\PYG{k}{struct}\PYG{+w}{ }\PYG{n+nc}{hlist\PYGZus{}head}\PYG{+w}{   }\PYG{n}{inodes}\PYG{p}{;}
\PYG{+w}{    }\PYG{n}{wait\PYGZus{}queue\PYGZus{}head\PYGZus{}t}\PYG{+w}{   }\PYG{n}{wait\PYGZus{}pidfd}\PYG{p}{;}
\PYG{+w}{    }\PYG{k}{struct}\PYG{+w}{ }\PYG{n+nc}{rcu\PYGZus{}head}\PYG{+w}{     }\PYG{n}{rcu}\PYG{p}{;}
\PYG{+w}{    }\PYG{k}{struct}\PYG{+w}{ }\PYG{n+nc}{upid}\PYG{+w}{         }\PYG{n}{numbers}\PYG{p}{[}\PYG{p}{]}\PYG{p}{;}\PYG{+w}{   }\PYG{c+cm}{/* isim alanı bilgisi */}
\PYG{p}{\PYGZcb{}}\PYG{p}{;}
\end{sphinxVerbatim}

\sphinxAtStartPar
Güncel versiyonlarda \sphinxcode{\sphinxupquote{upid}} yapısı şöyledir (\sphinxcode{\sphinxupquote{hlist\_node}} bağı kaldırılmış, ns bilgisi korunmuştur):

\begin{sphinxVerbatim}[commandchars=\\\{\}]
\PYG{k}{struct}\PYG{+w}{ }\PYG{n+nc}{upid}\PYG{+w}{ }\PYG{p}{\PYGZob{}}
\PYG{+w}{    }\PYG{k+kt}{int}\PYG{+w}{                  }\PYG{n}{nr}\PYG{p}{;}
\PYG{+w}{    }\PYG{k}{struct}\PYG{+w}{ }\PYG{n+nc}{pid\PYGZus{}namespace}\PYG{+w}{ }\PYG{o}{*}\PYG{n}{ns}\PYG{p}{;}
\PYG{p}{\PYGZcb{}}\PYG{p}{;}
\end{sphinxVerbatim}

\sphinxAtStartPar
Güncel çekirdeklerde pid araması yapan yüksek seviyeli fonksiyonlardan biri \sphinxstyleemphasis{find\_vpid} fonksiyonudur:

\begin{sphinxVerbatim}[commandchars=\\\{\}]
\PYG{k}{struct}\PYG{+w}{ }\PYG{n+nc}{pid}\PYG{+w}{ }\PYG{o}{*}\PYG{n}{find\PYGZus{}vpid}\PYG{p}{(}\PYG{k+kt}{int}\PYG{+w}{ }\PYG{n}{nr}\PYG{p}{)}
\PYG{p}{\PYGZob{}}
\PYG{+w}{    }\PYG{k}{return}\PYG{+w}{ }\PYG{n}{find\PYGZus{}pid\PYGZus{}ns}\PYG{p}{(}\PYG{n}{nr}\PYG{p}{,}\PYG{+w}{ }\PYG{n}{task\PYGZus{}active\PYGZus{}pid\PYGZus{}ns}\PYG{p}{(}\PYG{n}{current}\PYG{p}{)}\PYG{p}{)}\PYG{p}{;}
\PYG{p}{\PYGZcb{}}
\PYG{n}{EXPORT\PYGZus{}SYMBOL\PYGZus{}GPL}\PYG{p}{(}\PYG{n}{find\PYGZus{}vpid}\PYG{p}{)}\PYG{p}{;}
\end{sphinxVerbatim}

\sphinxAtStartPar
Bu fonksiyon belli bir pid değerine ilişkin pid nesnesini çağrıyı yapan prosesin bulunduğu pid isim
alanı içerisindeki XArray ağacında aramaktadır. Aşağıdaki diyagram güncel çekirdekteki çağrı zincirini
göstermektedir:

\sphinxincludegraphics[]{graphviz-712ac7de98d3a8f80e1d354370c2ce1e0b095664.pdf}

\sphinxAtStartPar
Prosesin içinde bulunduğu pid isim alanının bilgisine ulaşım çağrı zinciri şöyledir:

\begin{sphinxVerbatim}[commandchars=\\\{\}]
\PYG{k}{struct}\PYG{+w}{ }\PYG{n+nc}{pid\PYGZus{}namespace}\PYG{+w}{ }\PYG{o}{*}\PYG{n}{task\PYGZus{}active\PYGZus{}pid\PYGZus{}ns}\PYG{p}{(}\PYG{k}{struct}\PYG{+w}{ }\PYG{n+nc}{task\PYGZus{}struct}\PYG{+w}{ }\PYG{o}{*}\PYG{n}{tsk}\PYG{p}{)}
\PYG{p}{\PYGZob{}}
\PYG{+w}{    }\PYG{k}{return}\PYG{+w}{ }\PYG{n}{ns\PYGZus{}of\PYGZus{}pid}\PYG{p}{(}\PYG{n}{task\PYGZus{}pid}\PYG{p}{(}\PYG{n}{tsk}\PYG{p}{)}\PYG{p}{)}\PYG{p}{;}
\PYG{p}{\PYGZcb{}}

\PYG{k}{static}\PYG{+w}{ }\PYG{k+kr}{inline}\PYG{+w}{ }\PYG{k}{struct}\PYG{+w}{ }\PYG{n+nc}{pid}\PYG{+w}{ }\PYG{o}{*}\PYG{n}{task\PYGZus{}pid}\PYG{p}{(}\PYG{k}{struct}\PYG{+w}{ }\PYG{n+nc}{task\PYGZus{}struct}\PYG{+w}{ }\PYG{o}{*}\PYG{n}{task}\PYG{p}{)}
\PYG{p}{\PYGZob{}}
\PYG{+w}{    }\PYG{k}{return}\PYG{+w}{ }\PYG{n}{task}\PYG{o}{\PYGZhy{}}\PYG{o}{\PYGZgt{}}\PYG{n}{thread\PYGZus{}pid}\PYG{p}{;}
\PYG{p}{\PYGZcb{}}

\PYG{k}{static}\PYG{+w}{ }\PYG{k+kr}{inline}\PYG{+w}{ }\PYG{k}{struct}\PYG{+w}{ }\PYG{n+nc}{pid\PYGZus{}namespace}\PYG{+w}{ }\PYG{o}{*}\PYG{n}{ns\PYGZus{}of\PYGZus{}pid}\PYG{p}{(}\PYG{k}{struct}\PYG{+w}{ }\PYG{n+nc}{pid}\PYG{+w}{ }\PYG{o}{*}\PYG{n}{pid}\PYG{p}{)}
\PYG{p}{\PYGZob{}}
\PYG{+w}{    }\PYG{k}{struct}\PYG{+w}{ }\PYG{n+nc}{pid\PYGZus{}namespace}\PYG{+w}{ }\PYG{o}{*}\PYG{n}{ns}\PYG{+w}{ }\PYG{o}{=}\PYG{+w}{ }\PYG{n+nb}{NULL}\PYG{p}{;}
\PYG{+w}{    }\PYG{k}{if}\PYG{+w}{ }\PYG{p}{(}\PYG{n}{pid}\PYG{p}{)}
\PYG{+w}{        }\PYG{n}{ns}\PYG{+w}{ }\PYG{o}{=}\PYG{+w}{ }\PYG{n}{pid}\PYG{o}{\PYGZhy{}}\PYG{o}{\PYGZgt{}}\PYG{n}{numbers}\PYG{p}{[}\PYG{n}{pid}\PYG{o}{\PYGZhy{}}\PYG{o}{\PYGZgt{}}\PYG{n}{level}\PYG{p}{]}\PYG{p}{.}\PYG{n}{ns}\PYG{p}{;}
\PYG{+w}{    }\PYG{k}{return}\PYG{+w}{ }\PYG{n}{ns}\PYG{p}{;}
\PYG{p}{\PYGZcb{}}
\end{sphinxVerbatim}

\sphinxAtStartPar
XArray ağacında arama yapan \sphinxstyleemphasis{find\_pid\_ns} fonksiyonu artık son derece sade bir görünümdedir:

\begin{sphinxVerbatim}[commandchars=\\\{\}]
\PYG{k}{struct}\PYG{+w}{ }\PYG{n+nc}{pid}\PYG{+w}{ }\PYG{o}{*}\PYG{n}{find\PYGZus{}pid\PYGZus{}ns}\PYG{p}{(}\PYG{k+kt}{int}\PYG{+w}{ }\PYG{n}{nr}\PYG{p}{,}\PYG{+w}{ }\PYG{k}{struct}\PYG{+w}{ }\PYG{n+nc}{pid\PYGZus{}namespace}\PYG{+w}{ }\PYG{o}{*}\PYG{n}{ns}\PYG{p}{)}
\PYG{p}{\PYGZob{}}
\PYG{+w}{    }\PYG{k}{return}\PYG{+w}{ }\PYG{n}{idr\PYGZus{}find}\PYG{p}{(}\PYG{o}{\PYGZam{}}\PYG{n}{ns}\PYG{o}{\PYGZhy{}}\PYG{o}{\PYGZgt{}}\PYG{n}{idr}\PYG{p}{,}\PYG{+w}{ }\PYG{n}{nr}\PYG{p}{)}\PYG{p}{;}
\PYG{p}{\PYGZcb{}}
\PYG{n}{EXPORT\PYGZus{}SYMBOL\PYGZus{}GPL}\PYG{p}{(}\PYG{n}{find\PYGZus{}pid\PYGZus{}ns}\PYG{p}{)}\PYG{p}{;}
\end{sphinxVerbatim}

\sphinxAtStartPar
Ağaçta arama yapan asıl fonksiyon \sphinxstyleemphasis{idr\_find} isimli fonksiyondur.


\subsection{XArray Gerçekleştirimi}
\label{\detokenize{processcontrolblock:xarray-gerceklestirimi}}
\sphinxAtStartPar
XArray veri yapısının aşağı seviyeli gerçekleştirimi \sphinxcode{\sphinxupquote{include/linux/xarray.h}} dosyası içerisinde
yapılmıştır. XArray ağacı bu dosyadaki \sphinxcode{\sphinxupquote{xarray}} isimli yapıyla temsil edilmiştir:

\begin{sphinxVerbatim}[commandchars=\\\{\}]
\PYG{k}{struct}\PYG{+w}{ }\PYG{n+nc}{xarray}\PYG{+w}{ }\PYG{p}{\PYGZob{}}
\PYG{+w}{    }\PYG{n}{spinlock\PYGZus{}t}\PYG{+w}{   }\PYG{n}{xa\PYGZus{}lock}\PYG{p}{;}
\PYG{+w}{    }\PYG{c+cm}{/* private: */}
\PYG{+w}{    }\PYG{n}{gfp\PYGZus{}t}\PYG{+w}{        }\PYG{n}{xa\PYGZus{}flags}\PYG{p}{;}
\PYG{+w}{    }\PYG{k+kt}{void}\PYG{+w}{ }\PYG{n}{\PYGZus{}\PYGZus{}rcu}\PYG{+w}{  }\PYG{o}{*}\PYG{n}{xa\PYGZus{}head}\PYG{p}{;}
\PYG{p}{\PYGZcb{}}\PYG{p}{;}
\end{sphinxVerbatim}

\sphinxAtStartPar
Bu ağaç üzerinde işlem yapan \sphinxcode{\sphinxupquote{xa\_}} öneki ile başlayan bir grup fonksiyon vardır:

\begin{sphinxVerbatim}[commandchars=\\\{\}]
\PYG{k+kt}{void}\PYG{+w}{  }\PYG{o}{*}\PYG{n+nf}{xa\PYGZus{}load}\PYG{p}{(}\PYG{k}{struct}\PYG{+w}{ }\PYG{n+nc}{xarray}\PYG{+w}{ }\PYG{o}{*}\PYG{p}{,}\PYG{+w}{ }\PYG{k+kt}{unsigned}\PYG{+w}{ }\PYG{k+kt}{long}\PYG{+w}{ }\PYG{n}{index}\PYG{p}{)}\PYG{p}{;}
\PYG{k+kt}{void}\PYG{+w}{  }\PYG{o}{*}\PYG{n+nf}{xa\PYGZus{}store}\PYG{p}{(}\PYG{k}{struct}\PYG{+w}{ }\PYG{n+nc}{xarray}\PYG{+w}{ }\PYG{o}{*}\PYG{p}{,}\PYG{+w}{ }\PYG{k+kt}{unsigned}\PYG{+w}{ }\PYG{k+kt}{long}\PYG{+w}{ }\PYG{n}{index}\PYG{p}{,}\PYG{+w}{ }\PYG{k+kt}{void}\PYG{+w}{ }\PYG{o}{*}\PYG{n}{entry}\PYG{p}{,}\PYG{+w}{ }\PYG{n}{gfp\PYGZus{}t}\PYG{p}{)}\PYG{p}{;}
\PYG{k+kt}{void}\PYG{+w}{  }\PYG{o}{*}\PYG{n+nf}{xa\PYGZus{}erase}\PYG{p}{(}\PYG{k}{struct}\PYG{+w}{ }\PYG{n+nc}{xarray}\PYG{+w}{ }\PYG{o}{*}\PYG{p}{,}\PYG{+w}{ }\PYG{k+kt}{unsigned}\PYG{+w}{ }\PYG{k+kt}{long}\PYG{+w}{ }\PYG{n}{index}\PYG{p}{)}\PYG{p}{;}
\PYG{k+kt}{void}\PYG{+w}{  }\PYG{o}{*}\PYG{n+nf}{xa\PYGZus{}store\PYGZus{}range}\PYG{p}{(}\PYG{k}{struct}\PYG{+w}{ }\PYG{n+nc}{xarray}\PYG{+w}{ }\PYG{o}{*}\PYG{p}{,}\PYG{+w}{ }\PYG{k+kt}{unsigned}\PYG{+w}{ }\PYG{k+kt}{long}\PYG{+w}{ }\PYG{n}{first}\PYG{p}{,}
\PYG{+w}{                      }\PYG{k+kt}{unsigned}\PYG{+w}{ }\PYG{k+kt}{long}\PYG{+w}{ }\PYG{n}{last}\PYG{p}{,}\PYG{+w}{ }\PYG{k+kt}{void}\PYG{+w}{ }\PYG{o}{*}\PYG{n}{entry}\PYG{p}{,}\PYG{+w}{ }\PYG{n}{gfp\PYGZus{}t}\PYG{p}{)}\PYG{p}{;}
\PYG{k+kt}{bool}\PYG{+w}{   }\PYG{n+nf}{xa\PYGZus{}get\PYGZus{}mark}\PYG{p}{(}\PYG{k}{struct}\PYG{+w}{ }\PYG{n+nc}{xarray}\PYG{+w}{ }\PYG{o}{*}\PYG{p}{,}\PYG{+w}{ }\PYG{k+kt}{unsigned}\PYG{+w}{ }\PYG{k+kt}{long}\PYG{+w}{ }\PYG{n}{index}\PYG{p}{,}\PYG{+w}{ }\PYG{n}{xa\PYGZus{}mark\PYGZus{}t}\PYG{p}{)}\PYG{p}{;}
\PYG{k+kt}{void}\PYG{+w}{   }\PYG{n+nf}{xa\PYGZus{}set\PYGZus{}mark}\PYG{p}{(}\PYG{k}{struct}\PYG{+w}{ }\PYG{n+nc}{xarray}\PYG{+w}{ }\PYG{o}{*}\PYG{p}{,}\PYG{+w}{ }\PYG{k+kt}{unsigned}\PYG{+w}{ }\PYG{k+kt}{long}\PYG{+w}{ }\PYG{n}{index}\PYG{p}{,}\PYG{+w}{ }\PYG{n}{xa\PYGZus{}mark\PYGZus{}t}\PYG{p}{)}\PYG{p}{;}
\PYG{k+kt}{void}\PYG{+w}{   }\PYG{n+nf}{xa\PYGZus{}clear\PYGZus{}mark}\PYG{p}{(}\PYG{k}{struct}\PYG{+w}{ }\PYG{n+nc}{xarray}\PYG{+w}{ }\PYG{o}{*}\PYG{p}{,}\PYG{+w}{ }\PYG{k+kt}{unsigned}\PYG{+w}{ }\PYG{k+kt}{long}\PYG{+w}{ }\PYG{n}{index}\PYG{p}{,}\PYG{+w}{ }\PYG{n}{xa\PYGZus{}mark\PYGZus{}t}\PYG{p}{)}\PYG{p}{;}
\PYG{k+kt}{void}\PYG{+w}{  }\PYG{o}{*}\PYG{n+nf}{xa\PYGZus{}find}\PYG{p}{(}\PYG{k}{struct}\PYG{+w}{ }\PYG{n+nc}{xarray}\PYG{+w}{ }\PYG{o}{*}\PYG{p}{,}\PYG{+w}{ }\PYG{k+kt}{unsigned}\PYG{+w}{ }\PYG{k+kt}{long}\PYG{+w}{ }\PYG{o}{*}\PYG{n}{index}\PYG{p}{,}
\PYG{+w}{               }\PYG{k+kt}{unsigned}\PYG{+w}{ }\PYG{k+kt}{long}\PYG{+w}{ }\PYG{n}{max}\PYG{p}{,}\PYG{+w}{ }\PYG{n}{xa\PYGZus{}mark\PYGZus{}t}\PYG{p}{)}\PYG{p}{;}
\PYG{k+kt}{void}\PYG{+w}{  }\PYG{o}{*}\PYG{n+nf}{xa\PYGZus{}find\PYGZus{}after}\PYG{p}{(}\PYG{k}{struct}\PYG{+w}{ }\PYG{n+nc}{xarray}\PYG{+w}{ }\PYG{o}{*}\PYG{p}{,}\PYG{+w}{ }\PYG{k+kt}{unsigned}\PYG{+w}{ }\PYG{k+kt}{long}\PYG{+w}{ }\PYG{o}{*}\PYG{n}{index}\PYG{p}{,}
\PYG{+w}{                     }\PYG{k+kt}{unsigned}\PYG{+w}{ }\PYG{k+kt}{long}\PYG{+w}{ }\PYG{n}{max}\PYG{p}{,}\PYG{+w}{ }\PYG{n}{xa\PYGZus{}mark\PYGZus{}t}\PYG{p}{)}\PYG{p}{;}
\PYG{k+kt}{unsigned}\PYG{+w}{ }\PYG{k+kt}{int}\PYG{+w}{ }\PYG{n+nf}{xa\PYGZus{}extract}\PYG{p}{(}\PYG{k}{struct}\PYG{+w}{ }\PYG{n+nc}{xarray}\PYG{+w}{ }\PYG{o}{*}\PYG{p}{,}\PYG{+w}{ }\PYG{k+kt}{void}\PYG{+w}{ }\PYG{o}{*}\PYG{o}{*}\PYG{n}{dst}\PYG{p}{,}
\PYG{+w}{                        }\PYG{k+kt}{unsigned}\PYG{+w}{ }\PYG{k+kt}{long}\PYG{+w}{ }\PYG{n}{start}\PYG{p}{,}\PYG{+w}{ }\PYG{k+kt}{unsigned}\PYG{+w}{ }\PYG{k+kt}{long}\PYG{+w}{ }\PYG{n}{max}\PYG{p}{,}
\PYG{+w}{                        }\PYG{k+kt}{unsigned}\PYG{+w}{ }\PYG{k+kt}{int}\PYG{+w}{ }\PYG{n}{n}\PYG{p}{,}\PYG{+w}{ }\PYG{n}{xa\PYGZus{}mark\PYGZus{}t}\PYG{p}{)}\PYG{p}{;}
\PYG{k+kt}{void}\PYG{+w}{   }\PYG{n+nf}{xa\PYGZus{}destroy}\PYG{p}{(}\PYG{k}{struct}\PYG{+w}{ }\PYG{n+nc}{xarray}\PYG{+w}{ }\PYG{o}{*}\PYG{p}{)}\PYG{p}{;}
\end{sphinxVerbatim}

\sphinxAtStartPar
Bu fonksiyonların tanımlamaları \sphinxcode{\sphinxupquote{lib/xarray.c}} dosyasında yapılmıştır.

\sphinxAtStartPar
Güncel çekirdeklerin yukarıda belirttiğimiz aşağı seviyeli fonksiyonlarının yanı sıra aynı zamanda
bunları kullanan yüksek seviyeli \sphinxstyleemphasis{IDR} fonksiyonları da bulunmaktadır. Çekirdek kodları incelendiğinde
genellikle bu yüksek seviyeli \sphinxcode{\sphinxupquote{idr\_}} önekli fonksiyonların kullanıldığı görülmektedir. Aşağıdaki
diyagram üç katman arasındaki ilişkiyi göstermektedir:

\sphinxincludegraphics[]{graphviz-a1e484d869b51762e6b9cdfd85b8d93d483aa988.pdf}

\sphinxAtStartPar
Boş pid değerinin elde edilmesi de aynı XArray altyapısı üzerinden yapılmaktadır:

\begin{sphinxVerbatim}[commandchars=\\\{\}]
\PYG{c+cm}{/* Boş pid tahsisatı çağrı zinciri */}
\PYG{n}{alloc\PYGZus{}pid}\PYG{+w}{ }\PYG{o}{\PYGZhy{}}\PYG{o}{\PYGZhy{}}\PYG{o}{\PYGZgt{}}\PYG{+w}{ }\PYG{n}{idr\PYGZus{}alloc\PYGZus{}cyclic}\PYG{+w}{ }\PYG{o}{\PYGZhy{}}\PYG{o}{\PYGZhy{}}\PYG{o}{\PYGZgt{}}\PYG{+w}{ }\PYG{n}{idr\PYGZus{}alloc\PYGZus{}u32}\PYG{+w}{ }\PYG{o}{\PYGZhy{}}\PYG{o}{\PYGZhy{}}\PYG{o}{\PYGZgt{}}\PYG{+w}{ }\PYG{n}{idr\PYGZus{}get\PYGZus{}free}
\end{sphinxVerbatim}

\sphinxAtStartPar
XArray gerçekleştirimi hakkında ileride başka konularda ek bilgiler vereceğiz.

\sphinxstepscope


\chapter{Linux Çekirdeğinde Dosya Sistemi \sphinxhyphen{} I. Bölüm}
\label{\detokenize{filesystem1:linux-cekirdeginde-dosya-sistemi-i-bolum}}\label{\detokenize{filesystem1:dosya-sistemi-1}}\label{\detokenize{filesystem1::doc}}
\sphinxAtStartPar
Önceki bölümlerde \sphinxcode{\sphinxupquote{task\_struct}} nesnelerinin birbirleriyle bağlantısını ve pid değerleri ile
\sphinxcode{\sphinxupquote{task\_struct}} nesneleri arasındaki ilişkileri ele aldık. Şimdi dikkatimizi bir süre dosya sistemine
ilişkin çekirdek veri yapıları üzerine yönelteceğiz. Bilindiği gibi UNIX/Linux sistemlerinde pek çok
kavram kullanıcıya bir dosya gibi gösterilmektedir. Bu bölümde belli bir derinliğe kadar çekirdeğin
dosya işlemleri için oluşturduğu organizasyon üzerinde duracağız. Daha sonra başka bir bölümde
dosya sistemine ilişkin aşağı seviyeli ayrıntıları ele alacağız.


\section{Giriş}
\label{\detokenize{filesystem1:giris}}
\sphinxAtStartPar
İşletim sistemlerinin \sphinxstyleemphasis{dosya sistemi (file system)} denilen alt sistemlerinin iki tarafı vardır:
\begin{enumerate}
\sphinxsetlistlabels{\arabic}{enumi}{enumii}{}{.}%
\item {} 
\sphinxAtStartPar
\sphinxstylestrong{Disk tarafı}

\item {} 
\sphinxAtStartPar
\sphinxstylestrong{Bellek tarafı}

\end{enumerate}

\sphinxAtStartPar
Dosya bilgileri disk üzerindeki bloklarda tutulmaktadır. (Bu bloklara Microsoft dünyasında
\sphinxstyleemphasis{cluster} da denilmektedir.) Hangi dosyaların diskin hangi bloklarında tutulduğu, dosyaların
metadata bilgilerinin diskte nasıl saklandığı gibi belirlemeler dosya sisteminin disk tarafını;
diskteki dosya sisteminin çekirdekteki temsilinin oluşturulması ve işletim sisteminin açılan
dosyalar için yaptığı düzenlemeler ise dosya sisteminin bellek tarafını oluşturmaktadır.


\section{Disk Aktarımına İlişkin Temel Bilgiler}
\label{\detokenize{filesystem1:disk-aktarimina-iliskin-temel-bilgiler}}
\sphinxAtStartPar
Biz kursumuzda “disk” terimini ikinci bellekleri belirten genel bir terim olarak kullanacağız.
Bir süre önceye kadar disk olarak ağırlıklı biçimde \sphinxstyleemphasis{hard disk} dnilen elektromekanik birimler
kullanılıyordu. Ancak bir süredir artık disk olarak yarı iletken teknolojiler kullanılarak oluşturulmuş
\sphinxstyleemphasis{SSD (Solid State Disk)} denilen diskler kullanılmaktadır. Bugün ağırlıklı olarak kullandığımız SSD disklerin
herhangi bir mekanik parçası yoktur. SSD’ler \sphinxstyleemphasis{NAND Flash} denilen bellek teknolojisini kullanmaktadır.
SSD’ler hard disklere göre oldukça hızlıdır. Ancak onların en önemli handikapları belli bir yazma ömrünün
olmasıdır. SSD’lerde aynı bölgeye belli sayıdan daha fazla yazma yapıldığında artık SSD’nin o bölgesi
bozulabilmektedir. Tabii teknoloji bu bakımdan da ilerleme içerisindedir. SSD teknolojisi ile USB yuvalarına
taktığımız flash belleklerin teknolojisi birbirine benzemektedir.


\subsection{Sektörler ve Disk Denetleyicisi}
\label{\detokenize{filesystem1:sektorler-ve-disk-denetleyicisi}}
\sphinxAtStartPar
Kullandığımız disk birimi ister \sphinxstyleemphasis{hard disk} olsun isterse SSD olsun disk ile bilgisayarımızın
RAM’i arasındaki transferler \sphinxstyleemphasis{sektör (sector)} denilen byte blokları düzeyinde yapılmaktadır. Sektör
bir diskten okunabilecek ya da bir diske yazılabilecek en küçük birimdir. Bir sektör tipik olarak
512 byte’tır. Diskte byte düzeyinde erişim yoktur. Sektörel erişim vardır. Örneğin diskteki bir
sektörde bulunan bir byte üzerinde değişiklik ancak şöyle yapılabilmektedir: Önce o byte’ın
içinde bulunduğu sektör RAM’e okunur. Sonra byte RAM üzerinde değiştirilir. Sonra aynı sektör
yeniden diske yazılır.

\sphinxAtStartPar
Diskteki her sektörün ilk sektör 0 olmak üzere bir mantıksal numarası vardır. Hard disklerde
ardışıl numaralı mantıksal sektörler disk üzerinde de fiziksel olarak peşi sıra bulunmaktadır.
Mekanik hard disklerde bilgiler \sphinxstyleemphasis{track} denilen yollara yazılmaktadır. Ardışıl sektörler aynı
track’te bulunurlar. Dolayısıyla hard disklerde diskin kafası bir kez konumlandırıldığında
ardışıl sektörlere daha hızlı okuma yazma yapılabilmektedir. SSD’ler mekanik öğe barındırmadığı
için rastgele erişimlidir. Yani her sektörden okuma aynı hızda yapılmakta ve her sektöre yazma da
aynı hızda yapılmaktadır.

\sphinxAtStartPar
Modern bilgisayar sistemlerinde disk birimine doğrudan erişilmez. Disk erişimlerinde bu işleme
aracılık eden ismine \sphinxstyleemphasis{disk denetleyicisi (disk controller)} denen yerel bir işlemciden
faydalanılmaktadır. Yani sistem programcıları ya da işletim sistemlerini yazanlar disk
denetleyicisini programlar, disk denetleyicisi isteği elektriksel olarak disk birimine iletir,
okuma yazma işlemleri de disk birimi tarafından yapılır:

\sphinxincludegraphics[]{graphviz-b805a3ca799722a88d44016727257daf15ddca43.pdf}

\sphinxAtStartPar
Bugünkü masaüstü bilgisayarlarımızda \sphinxstyleemphasis{SATA} ve \sphinxstyleemphasis{NVMe} en çok kullanılan disk denetleyicileridir.


\subsection{DMA (Direct Memory Access)}
\label{\detokenize{filesystem1:dma-direct-memory-access}}
\sphinxAtStartPar
Peki işletim sistemi tarafından disk denetleyicisi “falanca sektörleri oku” ya da “falanca
sektörlere yaz” biçiminde  programlandıktan sonra aktarım nasıl yapılmaktadır?  Aktarım CPU
tarafından tek tek byte’ların denetleyiciden alınarak RAM’e yerleştirilmesi yoluyla
yapılmamaktadır. (Çok eskiden ilk PC mimarilerinde aktarım böyle de yapılabiliyordu.) Çünkü
CPU’nun bu işle meşgul olması önemli bir zaman kaybı oluşturmaktadır. Bu tür disk ile RAM
arasındaki aktarımlar için \sphinxstyleemphasis{DMA (Direct Memory Access)} denilen yardımcı denetleyiciler
kullanılmaktadır.

\sphinxAtStartPar
Tipik olarak CPU’da çalışan kod (yani işletim sistemi) disk denetleyicisine transfer isteğini
ve aktarımda kullanılacak bellek alanlarının adresini bildirir. Disk denetleyicisi de disk
birimini ve DMA’yı elektriksel düzeyde programlayarak aktarımın DMA üzerinden doğrudan RAM’e
yapılmasını sağlar. Aktarım sırasında artık CPU bu işle meşgul olmaz, işletim sistemi de
CPU’yu başka bir thread’i çalıştırması için \sphinxstyleemphasis{bağlamsal geçişe (task switch)} sokar. Tabii
aktarım işlemi bittiğinde disk denetleyicisi CPU’yu bir donanım kesmesi yoluyla durumdan
haberdar etmektedir.

\sphinxAtStartPar
Yani disk ile RAM arasındaki aktarım işlemleri tipik olarak şöyle yapılmaktadır:
\begin{enumerate}
\sphinxsetlistlabels{\arabic}{enumi}{enumii}{}{.}%
\item {} 
\sphinxAtStartPar
İşletim sistemi disk denetleyicisine aktarılacak sektörlere ilişkin bilgileri ve transfer
adreslerini CPU yoluyla elektriksel olarak iletir.

\item {} 
\sphinxAtStartPar
Okuma söz konusuysa disk denetleyicisi disk birimine elektriksel düzeyde komutlar göndererek
sektörlerin okunmasını ve DMA yoluyla bunların RAM’de uygun yerlere aktarılmasını sağlar.
Eğer yazma söz konusuysa RAM’de belirtilen adresteki bilgiler yine DMA yoluyla disk birimine
iletilerek yazma gerçekleştirilir.

\item {} 
\sphinxAtStartPar
Aktarım işlemi bittiğinde disk denetleyicisi bir donanım kesmesi yoluyla CPU’yu durumdan
haberdar eder.

\item {} 
\sphinxAtStartPar
İşletim sistemi aktarım için gereken kodları çalıştırdıktan sonra aktarım bitene kadar meşgul
bir döngüde beklemez. Başka thread’ler çalıştırılabiliyorsa \sphinxstyleemphasis{bağlamsal geçiş (context switch)}
oluşturarak CPU’nun boş biçimde beklemesinin önüne geçer.

\end{enumerate}

\sphinxAtStartPar
Bu süreci aşağıdaki diyagram özetlemektedir:

\sphinxincludegraphics[]{graphviz-bb623a5be421a0c57c65160a2fdd09b9e85dd5fb.pdf}

\sphinxAtStartPar
Eskiden Intel tabanlı PC mimarisinde ISA bus kullanıldığı zamanlarda tek bir merkezi DMA
denetleyicisi (Intel 8237) vardı. Ancak daha sonra PCI bus kullanılmaya başlanmasıyla birlikte
artık transfer yapabilen her donanım birimi kendi DMA denetleyicisini de içermeye başladı.
Bugün Intel tabanlı ve Apple Silicon tabanlı bilgisayar mimarilerinde disk denetleyicisi kendi
içerisindeki DMA denetleyicisini programlayarak transferi gerçekleştirmektedir. Disk
denetleyicilerinin programlanması ise artık uzunca bir süredir \sphinxstyleemphasis{bellekten tabanlı IO
(memory\sphinxhyphen{}mapped IO)} tekniği ile yapılmaktadır.

\sphinxAtStartPar
Hard disklerde disk birimi içerisinde önbellekler (cache) de bulundurulmaktadır. Böylece disk denetleyicisi aynı
sektörleri disk biriminden istediği zaman disk birimi eğer ilgili sektörler önbellek içerisindeyse hiç kafa
hareketleri yapmadan onları doğrudan önbellekten verebilmektedir. Bugünlerde örneğin 1 TB’lık hard disklerde 64MB,
128, 256 MB civarında önbellekler kullanılmaktadır. Yalnızca hard disklede değil SSD’lerde de bir önbellek sistemi
vardır. SSD’lerdeki önbellek sistemi özellikle yazma işlemlerinde hız kazancı sağlamakta ve aynı sektörlere sürekli
yazım yapıldığında o bölgenin yıpranmasını (wear leveling) engellemektedir. Tabii bu öönbellek sistemleri tamamen
disk birimleri tarafından içsel olarak (built\sphinxhyphen{}in) işletilmektedir. Bu önbellek sistemleri işletim sistemleri tarafından
erişilebilir değildir.


\subsection{Blok Kavramı}
\label{\detokenize{filesystem1:blok-kavrami}}
\sphinxAtStartPar
Yukarıda da belirttiğimiz gibi bir disk biriminde transfer edilecek en küçük birime sektör denilmektedir. Bir sektör
tipik olarak 512 byte uzunluğundadır. Ancak aslında sektör uzunlukları da disk üreticilerine bağlı olarak değişebilmektedir.
512 byte sektör uzunlukları bugün için standart bir uzunluktur. Tabii zaman geçtikçe diskler büyüdüğü için sektör
uzunluklarının da büyüyebileceğini söylemek istiyoruz. Nitekim 4K uzunluğunda sektörlere sahip olan diskler özellikle
büyük sistemlerde gittikçe yaygınlaşmaktadır. Disk birimi her sektöre ilk sektör 0 olmak üzere mantıksal bir numara
vermektedir. Yani adeta disk üzerindeki her sektörün bir adresi vardır. Disk denetleyicisi disk birimine transfer
edilecek sektörlerin numaralarını elektriksel düzeyde iletmektedir. (Bu biçimde mantıksal sektör numaraları kullanılmadan
önce 80’lerde ve 90’ların ilk yarısında sektörlerin yerleri “fiziksel koordinat sistemi” denilen “hangi yüz (head)”,
“hangi track”, “hangi sektör dilimi” biçiminde üç parametreyle belirtiliyordu.)

\sphinxAtStartPar
Sektör kavramı aslında dosya sistemleri için küçük bir depolama birimidir. İşletim sistemleri
bir dosyanın parçası olabilecek en küçük disk alanı için sektör yerine \sphinxstyleemphasis{blok (block)} ya da
\sphinxstyleemphasis{cluster} denilen daha büyük birimleri kullanmaktadır. Blok terimi daha çok UNIX/Linux
sistemlerinde kullanılmaktadır. Microsoft ise blok yerine \sphinxstyleemphasis{cluster} terimini kullanmaktadır.
Bir blok ardışıl n tane sektörden oluşmaktadır. Uygulamada bu n değeri 2’nin bir kuvveti olur. Ardışıllık hard disklerde
önemli bir unsurdur. Çünkü hard disklerde en önemli zaman kaybı mekanik bir birim olan disk kafasının track hizasına
çekilmesinde yaşanmaktadır. Disk kafası track hizasına çekildiğinde disk dönerken artık ardışıl sektörler hiç kafa
hareketi yapılmadan okunup yazılabilmektedir. Peki neden işletim sistemi dosyalar söz konusu olduğunda bir dosyanın
parçası olabilecek en küçük birim için sektör değil de ardışıl n tane sektör kullanmaktadır? İşte bunun birkaç nedeni
vardır:
\begin{enumerate}
\sphinxsetlistlabels{\arabic}{enumi}{enumii}{}{.}%
\item {} 
\sphinxAtStartPar
Dosyaların parçaları disk üzerinde ardışıl yerlerde olmak zorunda değildir. Eğer dosyalar
çok fazla parçadan oluşursa hard disklerde (ve kısmen de olsa SSD’lerde de) bu parçalar disk
üzerinde daha fazla yayılmış olur, bunlara erişmek için gereken zaman artar.

\item {} 
\sphinxAtStartPar
Eğer dosyanın parçaları sektör gibi küçük birimlerden oluşsaydı bu parçaların diskteki
yerlerine ilişkin metadata tabloları büyürdü. Bu da hem disk alanını hem de işletim sisteminin
bellekte yaptığı düzenlemede alan verimsizliği oluştururdu.

\item {} 
\sphinxAtStartPar
CPU’ların kullandığı sayfalama mekanizmasında genellikle 4K uzunluklar kullanılmaktadır.
Dosya parçalarının 4K uzunluğun katlarında olması dosya sistemi ile sayfalama sistemi arasında
daha iyi bir uyumun ortaya çıkmasına yol açmaktadır.

\end{enumerate}

\sphinxAtStartPar
Peki bu durumda işletim sistemleri blok denilen dosyanın parçası olabilecek en küçük birim için hangi uzunluğu
kullanmaktadır? İşte genellikle bu karar disk formatlanırken diskin (disk bölümünün) büyüklüğüne bakılarak verilmektedir.
Dosyaların son bloklarında kalan kullanılmayan alanların oluşmasına \sphinxstyleemphasis{içsel bölünme (internal fragmentation)}
denilmektedir. Küçük disklerde (disk bölümlerinde) içsel bölünmenin etkisi daha büyük olacağından blokların 1K gibi
küçük uzunluklarda alınması uygun olabilir. Ancak orta büyüklükte disklerde içsel bölünmenin etkisi göreli olarak
azalacağı için bloklar 4k gibi bir değerde seçilebilmektedir. Büyük disklerde ise 8K, 16K blok büyüklükleri tercih
edilmektedir. Aslında blok büyüklükleri ilgili disk bölümü formatlanırken (Linux sistemlerinde \sphinxcode{\sphinxupquote{mkfs.xxx}} programlarıyla
formatlama yapılmaktadır) belirlenmektedir. Yani kullanıcı isterse kendisi bu programda kendi tercih ettiği blok
uzunluğunu kullanabilir. Ancak kullanıcılar genellikle böyle bir belirleme yapmazlar. Bu durumda bu programlar disk
bölümünün büyüklüğüne bağlı olarak yukarıda açıkladığımız gibi uygun bir blok büyüklüğünü seçerler.


\subsection{Mount İşlemi}
\label{\detokenize{filesystem1:mount-islemi}}
\sphinxAtStartPar
UNIX/Linux sistemlerinde dosya sistemi için tek bir kök vardır. Blok aygıtları (örneğin hard
diskler, flash bellekler vb.) belli bir dizine mount edilmektedir. Mount işlemi bir dizin üzerine
uygulanır; mount işlemi sonucunda o dizinin içeriği görünmez, artık mount edilen dosya sisteminin
kök dizini mount dizininde gözükür. Dolayısıyla bu sistemlerde farklı dizinler farklı blok
büyüklüklerine ilişkin dosya sistemlerinin içerisinde olabilmektedir.

\sphinxAtStartPar
Anımsanacağı gibi \sphinxcode{\sphinxupquote{stat}} POSIX fonksiyonu ya da komut satırından uygulanan \sphinxstyleemphasis{stat} komutu belli
bir dosyanın bilgilerini verirken o dosyanın içinde bulunduğu dosya sisteminin blok uzunluğunu
da vermektedir. Linux sistemlerinde bu bilgi doğrudan dosya sistemine ilişkin blok aygıtı
üzerinde \sphinxcode{\sphinxupquote{dumpe2fs}} programıyla da elde edilebilmektedir.

\sphinxAtStartPar
Windows sistemlerinde de \sphinxstyleemphasis{cluster} adı altında blok sistemi kullanılmaktadır. O sistemlerde blok
uzunluklarını \sphinxstyleemphasis{chkdsk} programı ile ya da \sphinxstyleemphasis{fsutil} programı ile komut satırından elde
edebilirsiniz.


\subsection{Blok Numaralandırması}
\label{\detokenize{filesystem1:blok-numaralandirmasi}}
\sphinxAtStartPar
İşletim sistemleri işlemlerini kolaylaştırmak için her bloğa bir numara da vermektedir. Örneğin
bir bloğun 4K (tipik olarak 8 sektör) olduğunu düşünelim. İşletim sistemi için ilgili disk (aslında
disk bölümü ya da genel olarak blok aygıtı da diyebiliriz) bloklardan oluşmaktadır. Örneğin
diskin (disk bölümünün) ilk 8 sektörü artık 0’ıncı bloktur. Sonraki 8 sektör 1’inci bloktur.

\sphinxAtStartPar
İşletim sistemi içsel olarak artık ilgili diski bloklardan oluşan ve her bloğun bir numarasının
olduğu mantıksal bir depolama alanı gibi ele almaktadır. Yani işletim sistemi için yalnızca
sektörlerin değil aynı zamanda dosya sistemine ilişkin blokların da numaraları vardır.


\subsection{Sayfa Önbelleği (Page Cache)}
\label{\detokenize{filesystem1:sayfa-onbellegi-page-cache}}
\sphinxAtStartPar
İşletim sistemleri son okunan ya da yazılan disk bloklarını RAM’de bir önbellek sisteminde
saklamaktadır. Bu önbellek sistemine genel olarak \sphinxstyleemphasis{işletim sisteminin disk önbellek sistemi}
denilmektedir. Linux dünyasında eskiden bu önbellek sistemine \sphinxstyleemphasis{buffer cache} deniliyordu. Sonra
bu önbellek sistemi iyileştirildi ve ismi \sphinxstyleemphasis{page cache} olarak değiştirildi.

\sphinxAtStartPar
İşletim sistemlerinin bu disk önbellek sistemleri disk erişimini ciddi boyutta azaltmakta ve
sistem performansı üzerinde en önemli olumlu etkilerden birini oluşturmaktadır. Eğer işletim
sistemlerinde böyle bir disk önbellek sistemi olmasaydı sistemler çok yavaş çalışırdı.


\section{Dosyalardan Okuma ve Yazma İşlemleri}
\label{\detokenize{filesystem1:dosyalardan-okuma-ve-yazma-islemleri}}
\sphinxAtStartPar
Biz UNIX/Linux sistemlerinde bir dosyadan okuma yapmak için ya da bir dosyaya yazma yapmak için
\sphinxcode{\sphinxupquote{read}}/\sphinxcode{\sphinxupquote{write}} POSIX fonksiyonlarını kullanmış olalım. Bu fonksiyonlar çekirdek içerisindeki
\sphinxcode{\sphinxupquote{sys\_read}} ve \sphinxcode{\sphinxupquote{sys\_write}} sistem fonksiyonlarını çalıştırmaktadır.

\sphinxAtStartPar
İşletim sistemi bellek tarafında yaptığı organizasyonla okunacak ya da yazılacak bilginin ilgili
blok aygıtının (disk bölümünün) kaç numaralı bloğuna ve sektörüne ilişkin olduğunu
belirleyebilmektedir. Ancak \sphinxcode{\sphinxupquote{sys\_read}} ve \sphinxcode{\sphinxupquote{sys\_write}} gibi sistem fonksiyonları hemen diske
yönelmez. Bu fonksiyonlar önce dosyanın ilgili bölümünün RAM’de oluşturulmuş bir sayfa önbelleği
(\sphinxstyleemphasis{page cache}) içerisinde olup olmadığına bakmaktadır. Eğer ilgili bölüm bu önbellek sisteminin
içerisinde varsa bu fonksiyonlar diske hiç erişmeden dolayısıyla da hiç bloke olmadan bu okuma
yazma işlemini gerçekleştirmektedir.

\sphinxAtStartPar
\sphinxcode{\sphinxupquote{read}}/\sphinxcode{\sphinxupquote{write}} POSIX fonksiyonları çağrıldığında yapılan işlemler aşağıdaki diyagramda
özetlenmiştir:

\sphinxincludegraphics[]{graphviz-8027075cd695b166407e55ea6ae99e013eea4f82.pdf}

\sphinxAtStartPar
Peki dosyadaki okunacak ya da yazılacak kısım sayfa önbelleğinde (\sphinxstyleemphasis{page cache}) yoksa gerçek
transfer nasıl yapılmaktadır? Linux sistemlerinde bu transferlerin yapıldığı birime \sphinxstyleemphasis{blok aygıt
sürücüleri (block device drivers)} denilmektedir.

\sphinxAtStartPar
Bir Linux sistemi kurulduğunda zaten temel disk denetleyicileri üzerinden transfer yapabilen blok
aygıt sürücüleri çekirdeğin içerisine gömülmüş durumda olur. Ancak sistem programcısının kendisi
de blok aygıt sürücüleri yazabilir. Örneğin bir gömülü Linux sisteminde yeni bir SD kart birimi
için bir blok aygıt sürücüsü yazmak zorunda kalabilirsiniz.

\sphinxAtStartPar
İşletim sistemi bu IO isteklerini hemen blok aygıt sürücüsüne göndermez. Çünkü çok sayıda
farklı proses aynı disk sektörlerini okuyacak ya da o sektörlere yazacak olabilir. İşletim
sistemi önce istekleri sıraya dizer, mümkünse birleştirir, bu biçimdeki iyileştirme işleminden
sonra istekleri blok aygıt sürücüsüne gönderir. Bu sürece \sphinxstyleemphasis{IO çizelgelemesi (IO scheduling)}
denilmektedir.

\sphinxAtStartPar
O halde bir dosya okuması ya da yazması sonucunda gelişen olayları şöyle özetleyebiliriz:
\begin{enumerate}
\sphinxsetlistlabels{\arabic}{enumi}{enumii}{}{.}%
\item {} 
\sphinxAtStartPar
Kullanıcı modunda çalışan program (yani proses) \sphinxcode{\sphinxupquote{read}} ya da \sphinxcode{\sphinxupquote{write}} POSIX fonksiyonlarını
çağırır. (UNIX/Linux sistemlerindeki C derleyicilerinin standart dosya fonksiyonları da eğer
okunacak ya da yazılacak kısım kendi tamponlarında yoksa zaten bu POSIX fonksiyonlarını
çağırmaktadır.)

\item {} 
\sphinxAtStartPar
\sphinxcode{\sphinxupquote{read}} ve \sphinxcode{\sphinxupquote{write}} POSIX fonksiyonları Linux’ta \sphinxcode{\sphinxupquote{sys\_read}} ve \sphinxcode{\sphinxupquote{sys\_write}} isimli sistem
fonksiyonlarını çağırır. Artık akış kullanıcı modundan (user mode) çekirdek moduna (kernel
mode) geçmiştir.

\item {} 
\sphinxAtStartPar
\sphinxcode{\sphinxupquote{sys\_read}} ve \sphinxcode{\sphinxupquote{sys\_write}} fonksiyonları önce okunacak ya da yazılacak yerin Linux’un disk
önbellek sistemi olan sayfa önbelleğinde (\sphinxstyleemphasis{page cache}, eski ismiyle \sphinxstyleemphasis{buffer cache}) olup
olmadığına bakar. Eğer ilgili disk blokları sayfa önbelleğinde varsa akış hiç bloke olmadan
sayfa önbelleği içerisinden karşılanır. Aksi hâlde Linux çekirdeği isteği \sphinxstyleemphasis{IO çizelgeleyicisi
(IO scheduler)} denilen çekirdek birimine iletir; \sphinxcode{\sphinxupquote{read}}/\sphinxcode{\sphinxupquote{write}} çağrısını yapan thread
bloke edilir.

\item {} 
\sphinxAtStartPar
IO çizelgeleyicisi istekleri çizelgeler. Okuma işlemi söz konusuysa sayfa önbelleğinde
transfer edilecek önbellek bloklarını tahsis eder. Yazma işlemi söz konusuysa diske transfer
edilecek önbellek bloklarını belirler. Blok aygıt sürücüsüne transfer edilecek sektörleri ve
transfere ilişkin bellek adreslerini iletir.

\item {} 
\sphinxAtStartPar
Gerçek transfer işlemi blok aygıt sürücüsü tarafından yapılmaktadır. İşletim sistemi blok
aygıt sürücüsüne hangi sektörlerin sayfa önbelleğindeki hangi adreslere (ya da tersi yönde)
transfer edileceğini bir kuyruk sistemi yardımıyla iletmektedir.

\item {} 
\sphinxAtStartPar
Blok aygıt sürücüsü diskten istenen sektörleri sayfa önbelleği içerisinde belirtilen adrese
ya da sayfa önbelleğindeki bilgileri diskin belirtilen sektörlerine transfer eder.

\item {} 
\sphinxAtStartPar
Artık okuma söz konusuysa okunan bilgi sayfa önbelleği içerisindedir. İşlemi başlatan
thread’in blokesi çözülür. \sphinxcode{\sphinxupquote{sys\_read}} sistem fonksiyonu bunu sayfa önbelleği içerisinden
programcının kullanıcı modundaki adresine kopyalar.

\end{enumerate}

\sphinxAtStartPar
Bu süreci aşağıdaki diyagram özetlemektedir:

\sphinxincludegraphics[]{graphviz-f82b3dc117927b6fbff691376f5395e08d140542.pdf}

\sphinxAtStartPar
Linux sistemlerinde yukarıda özetlediğimiz olaylar silsilesi zaman içerisinde değişikliklere
uğratılarak ve sürekli geliştirilerek bugünkü durumuna getirilmiştir.


\subsection{Yazma İşleminin Ayrıntıları}
\label{\detokenize{filesystem1:yazma-isleminin-ayrintilari}}
\sphinxAtStartPar
Buradaki süreçte yazma olayı söz konusu olduğunda bazı ayrıntılar da devreye girmektedir.
Kullanıcı modundaki program \sphinxcode{\sphinxupquote{write}} POSIX fonksiyonunu çağırıp bu fonksiyon da \sphinxcode{\sphinxupquote{sys\_write}}
fonksiyonunu çağırdığında bu sistem fonksiyonu yazılmak istenen bilgiler diske yazılana kadar
\sphinxcode{\sphinxupquote{write}} işlemini yapan thread’i bloke etmez. Yazma işlemi her zaman Linux’un RAM’deki sayfa
önbelleğine yapılmaktadır. \sphinxcode{\sphinxupquote{sys\_write}} fonksiyonu yazmayı sayfa önbelleği içerisine yaptıktan
sonra hemen \sphinxstyleemphasis{başarılı} olarak geri dönmektedir.

\sphinxAtStartPar
Sistem programlama terminolojisinde IO işlemlerinde \sphinxstyleemphasis{senkron} terimi \sphinxstyleemphasis{fonksiyon geri döndüğünde
tüm işlemlerin yapılıp bitmiş olması} anlamına gelmektedir. \sphinxstyleemphasis{Asenkron} terimi ise \sphinxstyleemphasis{işlemin
başlatılması, fonksiyonun geri dönmesi ancak işlemin aslında arka planda devam etmesi} anlamına
gelmektedir. Görüldüğü gibi modern işletim sistemlerinde diske yazma işlemi aslında \sphinxstyleemphasis{senkron}
bir işlem değildir.

\sphinxAtStartPar
Ancak bunun bir istisnası vardır. Eğer bir dosya \sphinxcode{\sphinxupquote{O\_DSYNC}} ya da \sphinxcode{\sphinxupquote{O\_SYNC}} bayraklarıyla
açılmışsa o dosyaya yapılan yazma işlemleri aygıta aktarılana kadar thread \sphinxcode{\sphinxupquote{write}} fonksiyonunda
bloke edilmektedir. Yani bu bayraklar yazma işlemlerinin senkron yapılmasını sağlamaktadır.


\subsection{Gecikmeli Yazım ve Flush Thread’leri}
\label{\detokenize{filesystem1:gecikmeli-yazim-ve-flush-thread-leri}}
\sphinxAtStartPar
Peki yazma işleminde sayfa önbelleğine yazılan bilgiler çekirdek tarafından ne zaman gerçek
aygıta aktarılmaktadır? İşletim sistemi kasten bu tür durumlarda araya belli bir gecikme
koymaktadır. Böylece peşi sıra yapılan \sphinxcode{\sphinxupquote{write}} işlemlerinin tek tek gereksiz biçimde aygıta
yapılması engellenir, bunlar biriktirilerek ve çizelgelenerek blok aygıt sürücüsüne aktarılır.
Bu biçimdeki aktarmaya \sphinxstyleemphasis{gecikmeli yazım (delayed write)} da denilmektedir.

\sphinxAtStartPar
Buradaki süreci aşağıdaki şekille özetleyebiliriz:

\sphinxincludegraphics[]{graphviz-73df1d47451c50e1e304783fb01f87558cdfd1d8.pdf}

\sphinxAtStartPar
Modern Linux sistemlerinde kirlenmiş sayfaların \sphinxstyleemphasis{flush} edilmesi \sphinxstyleemphasis{çekirdek thread’leri (kernel
threads)} tarafından yapılmaktadır. Örneğin Linux sistemlerinde bu işlemlerden \sphinxstyleemphasis{flush} isimli
çekirdek thread’leri sorumludur. Eskiden Linux çekirdeklerinin 2.6.32 versiyonuna kadar bu
işlemler \sphinxstyleemphasis{pdflush} isimli tek bir çekirdek thread tarafından yapılıyordu. Bu versiyondan sonra
artık her blok aygıt sürücüsü için ayrı bir flush thread’i oluşturulmaya başlandı. Bu thread’leri
komut satırında şu şekilde görüntüleyebilirsiniz:

\begin{sphinxVerbatim}[commandchars=\\\{\}]
\PYGZdl{}\PYG{+w}{ }ps\PYG{+w}{ }\PYGZhy{}aux\PYG{+w}{ }\PYG{p}{|}\PYG{+w}{ }grep\PYG{+w}{ }flush
\end{sphinxVerbatim}

\sphinxAtStartPar
flush thread’leri arka planda sürekli olarak sayfa önbelleğini izler. Orada \sphinxstyleemphasis{kirlenmiş (dirty)}
olan sektörleri ilgili blok aygıt sürücüsüne gönderir.


\subsection{Flush Thread Parametreleri}
\label{\detokenize{filesystem1:flush-thread-parametreleri}}
\sphinxAtStartPar
Modern Linux sistemleri için yazma gecikmesi ortalama 5 saniye civarındadır. Ancak bu değerler
değiştirilebilmektedir. Flush thread’lerinin parametreleri aşağıdaki tabloda özetlenmiştir:


\begin{savenotes}\sphinxattablestart
\sphinxthistablewithglobalstyle
\centering
\sphinxcapstartof{table}
\sphinxthecaptionisattop
\sphinxcaption{Flush Thread Parametreleri}\label{\detokenize{filesystem1:id1}}
\sphinxaftertopcaption
\begin{tabular}[t]{\X{35}{100}\X{20}{100}\X{45}{100}}
\sphinxtoprule
\begin{varwidth}[t]{\sphinxcolwidth{1}{3}}
\sphinxstyletheadfamily \sphinxAtStartPar
Parametre
\sphinxbeforeendvarwidth
\end{varwidth}%
&\begin{varwidth}[t]{\sphinxcolwidth{1}{3}}
\sphinxstyletheadfamily \sphinxAtStartPar
Varsayılan Değer
\sphinxbeforeendvarwidth
\end{varwidth}%
&\begin{varwidth}[t]{\sphinxcolwidth{1}{3}}
\sphinxstyletheadfamily \sphinxAtStartPar
Anlamı
\sphinxbeforeendvarwidth
\end{varwidth}%
\\
\sphinxmidrule
\sphinxtableatstartofbodyhook\begin{varwidth}[t]{\sphinxcolwidth{1}{3}}
\sphinxAtStartPar
\sphinxcode{\sphinxupquote{dirty\_writeback\_centisecs}}
\sphinxbeforeendvarwidth
\end{varwidth}%
&\begin{varwidth}[t]{\sphinxcolwidth{1}{3}}
\sphinxAtStartPar
500
\sphinxbeforeendvarwidth
\end{varwidth}%
&\begin{varwidth}[t]{\sphinxcolwidth{1}{3}}
\sphinxAtStartPar
Flusher thread’in periyodik olarak çalıştığı aralık (santi saniye cinsinden).
500 cs = 5 saniye. Bu aralıkta çekirdek \sphinxstyleemphasis{dirty} sayfaları kontrol eder.
\sphinxbeforeendvarwidth
\end{varwidth}%
\\
\sphinxhline\begin{varwidth}[t]{\sphinxcolwidth{1}{3}}
\sphinxAtStartPar
\sphinxcode{\sphinxupquote{dirty\_expire\_centisecs}}
\sphinxbeforeendvarwidth
\end{varwidth}%
&\begin{varwidth}[t]{\sphinxcolwidth{1}{3}}
\sphinxAtStartPar
3000
\sphinxbeforeendvarwidth
\end{varwidth}%
&\begin{varwidth}[t]{\sphinxcolwidth{1}{3}}
\sphinxAtStartPar
Bir \sphinxstyleemphasis{dirty} sayfa en fazla bu kadar süre (santi saniye cinsinden) RAM’de kalabilir.
3000 cs = 30 saniye sonra \sphinxstyleemphasis{süresi dolmuş} sayılır ve flush edilir.
\sphinxbeforeendvarwidth
\end{varwidth}%
\\
\sphinxhline\begin{varwidth}[t]{\sphinxcolwidth{1}{3}}
\sphinxAtStartPar
\sphinxcode{\sphinxupquote{dirty\_ratio}} / \sphinxcode{\sphinxupquote{dirty\_background\_ratio}}
\sphinxbeforeendvarwidth
\end{varwidth}%
&\begin{varwidth}[t]{\sphinxcolwidth{1}{3}}
\sphinxAtStartPar
\%20 / \%10 civarı
\sphinxbeforeendvarwidth
\end{varwidth}%
&\begin{varwidth}[t]{\sphinxcolwidth{1}{3}}
\sphinxAtStartPar
RAM’in ne kadarı \sphinxstyleemphasis{dirty} sayfalarla dolarsa flush işleminin başlatılacağını belirler
(bellek baskısı durumunda zaman beklenmez).
\sphinxbeforeendvarwidth
\end{varwidth}%
\\
\sphinxbottomrule
\end{tabular}
\sphinxtableafterendhook\par
\sphinxattableend\end{savenotes}

\sphinxAtStartPar
Bu değerler proc dosya sisteminden görüntülenebilmektedir:

\begin{sphinxVerbatim}[commandchars=\\\{\}]
\PYGZdl{}\PYG{+w}{ }cat\PYG{+w}{ }/proc/sys/vm/dirty\PYGZus{}writeback\PYGZus{}centisecs
\PYGZdl{}\PYG{+w}{ }cat\PYG{+w}{ }/proc/sys/vm/dirty\PYGZus{}expire\PYGZus{}centisecs
\end{sphinxVerbatim}

\sphinxAtStartPar
\sphinxstyleemphasis{sysctl} komutu ile de bu değerler değiştirilebilmektedir:

\begin{sphinxVerbatim}[commandchars=\\\{\}]
\PYGZdl{}\PYG{+w}{ }sudo\PYG{+w}{ }sysctl\PYG{+w}{ }\PYGZhy{}w\PYG{+w}{ }vm.dirty\PYGZus{}writeback\PYGZus{}centisecs\PYG{o}{=}\PYG{l+m}{100}
\end{sphinxVerbatim}

\sphinxAtStartPar
\sphinxstyleemphasis{sysctl} komutu zaten kendi içerisinde \sphinxcode{\sphinxupquote{/proc/sys}} dizinindeki dosyalar üzerinde güncelleme
işlemleri yapmaktadır.

\sphinxAtStartPar
flush thread’lerinin çalışması daha ayrıntılı olarak \sphinxstyleemphasis{sayfa önbelleği (page cache)} konusunun
ele alındığı bölümde açıklanacaktır.


\subsection{Gecikmeli Yazımın Gerekçeleri}
\label{\detokenize{filesystem1:gecikmeli-yazimin-gerekceleri}}
\sphinxAtStartPar
\sphinxstyleemphasis{Gecikmeli yazım (delayed write)} işleminin gerekçeleri nelerdir? En önemli gerekçe peşi sıra
yapılan yazma işlemlerinin tek hamlede aygıta yansıtılmasıdır. Bu sayede yazma işlemini yapan
thread bloke olmaz ve toplamda bu işlemler paralel yürütüldüğü için sistem performansı yükselir.
Aynı zamanda flash belleklerde ve SSD’lerde bu gecikme sürekli yazım sonucunda oluşan belleğin
eskimesini de kısmen engellemektedir.


\subsection{Kirlenmiş Sayfaların Erken Temizlenmesi}
\label{\detokenize{filesystem1:kirlenmis-sayfalarin-erken-temizlenmesi}}
\sphinxAtStartPar
Sayfa önbelleğinde kirlenen sayfalar bazı durumlarda işletim sisteminin \sphinxstyleemphasis{tazeleyici (flusher)}
thread’lerini beklemeden de diske aktarılabilmektedir. Örneğin bir dosya kapatıldığında artık
bu işlem arka planda dosyanın kirlenmiş sayfalarının da diske yazılmasına yol açmaktadır.


\subsection{sync, fsync ve İlgili Açış Bayrakları}
\label{\detokenize{filesystem1:sync-fsync-ve-ilgili-acis-bayraklari}}
\sphinxAtStartPar
UNIX/Linux sistemlerinin çoğunda bazı özel sistem fonksiyonları yoluyla ya da \sphinxcode{\sphinxupquote{open}}
fonksiyonundaki bayraklarla da bu duruma müdahale edilebilmektedir.

\sphinxAtStartPar
\sphinxcode{\sphinxupquote{sync}} POSIX fonksiyonu çağrıldığında o anda dosya sistemine ilişkin \sphinxstyleemphasis{kirlenmiş (dirty olan)}
sayfaların hepsi flush edilmektedir:

\begin{sphinxVerbatim}[commandchars=\\\{\}]
\PYG{c+cp}{\PYGZsh{}}\PYG{c+cp}{include}\PYG{+w}{ }\PYG{c+cpf}{\PYGZlt{}unistd.h\PYGZgt{}}

\PYG{k+kt}{void}\PYG{+w}{ }\PYG{n+nf}{sync}\PYG{p}{(}\PYG{k+kt}{void}\PYG{p}{)}\PYG{p}{;}
\end{sphinxVerbatim}

\sphinxAtStartPar
\sphinxcode{\sphinxupquote{sync}} fonksiyonu asenkron biçimde çalışmaktadır. Yani fonksiyon geri döndüğünde tüm blokların
flush edilmiş olma garantisi yoktur. Aynı zamanda Linux sistemlerde \sphinxstyleemphasis{sync} isimli bir kabuk komutu
da bulunmaktadır. Bu komut \sphinxcode{\sphinxupquote{sync}} fonksiyonunu çağırmaktadır.

\sphinxAtStartPar
\sphinxcode{\sphinxupquote{fsync}} POSIX fonksiyonu ise belli bir dosyaya ilişkin kirlenmiş sayfaların flush edilmesi için
kullanılmaktadır:

\begin{sphinxVerbatim}[commandchars=\\\{\}]
\PYG{c+cp}{\PYGZsh{}}\PYG{c+cp}{include}\PYG{+w}{ }\PYG{c+cpf}{\PYGZlt{}unistd.h\PYGZgt{}}

\PYG{k+kt}{int}\PYG{+w}{ }\PYG{n+nf}{fsync}\PYG{p}{(}\PYG{k+kt}{int}\PYG{+w}{ }\PYG{n}{fildes}\PYG{p}{)}\PYG{p}{;}
\end{sphinxVerbatim}

\sphinxAtStartPar
\sphinxcode{\sphinxupquote{fsync}} fonksiyonu senkronize (synchronous) çalışmaktadır. Yani fonksiyon geri döndüğünde sayfa
önbelleğindeki kirlenmiş sayfaların flush edilmiş olması garanti edilmektedir.

\sphinxAtStartPar
Bir dosya açılırken \sphinxcode{\sphinxupquote{open}} POSIX fonksiyonunda kullanılan konuyla ilgili üç bayrak vardır:
\begin{description}
\sphinxlineitem{\sphinxcode{\sphinxupquote{O\_DSYNC}} \sphinxstylestrong{Bayrağı}}
\sphinxAtStartPar
Bu bayrak POSIX’in “Base Definitions” bölümündeki “Synchronized I/O Data Integrity Completion”
başlığında açıklanan yazma koşullarının sağlanacağını belirtmektedir. Bu bayrak kullanıldığında
aşağıdaki iki durumun çekirdek tarafından sağlanması garanti edilmektedir:
\begin{itemize}
\item {} 
\sphinxAtStartPar
Dosyaya yazdırılan bilgilerin \sphinxcode{\sphinxupquote{write}} fonksiyonu geri döndüğünde hedefe transfer edilmiş
olması.

\item {} 
\sphinxAtStartPar
Yazılan bilginin dosyadan okunabilmesi için gereken metadata bilgilerinin hedefe transfer
edilmiş olması. (Tüm metadata bilgilerinin hedefe transfer edilmiş olması gerekmemektedir.)

\end{itemize}

\sphinxlineitem{\sphinxcode{\sphinxupquote{O\_SYNC}} \sphinxstylestrong{Bayrağı}}
\sphinxAtStartPar
Bu bayrak POSIX’in “Base Definitions” bölümündeki “Synchronized I/O File Integrity Completion”
başlığında açıklanan koşulların sağlanacağını belirtmektedir. \sphinxcode{\sphinxupquote{O\_SYNC}} bayrağı \sphinxcode{\sphinxupquote{O\_DSYNC}}
bayrağını kapsamaktadır. Fakat bu bayrak \sphinxcode{\sphinxupquote{write}} fonksiyonu geri dönmeden önce tüm metadata
bilgilerinin hedefe transfer edilmiş olmasını zorunlu tutmaktadır.

\sphinxlineitem{\sphinxcode{\sphinxupquote{O\_RSYNC}}}
\sphinxAtStartPar
Bu okuma işlemi ile ilgilidir. Tek başına değil \sphinxcode{\sphinxupquote{O\_DSYNC}} ya da \sphinxcode{\sphinxupquote{O\_SYNC}} bayraklarıyla
birlikte kullanılır. Eğer \sphinxcode{\sphinxupquote{O\_RSYNC}} bayrağı \sphinxcode{\sphinxupquote{O\_DSYNC}} bayrağı ile birlikte kullanılırsa
\sphinxcode{\sphinxupquote{read}} işlemini etkileyecek olan daha önce yapılmış \sphinxcode{\sphinxupquote{write}} işlemleri varsa \sphinxcode{\sphinxupquote{read}}
fonksiyonu geri dönmeden önce bu \sphinxcode{\sphinxupquote{write}} işlemleri için \sphinxcode{\sphinxupquote{O\_DSYNC}} bayrağında belirtilen
semantik uygulanmaktadır. Eğer bu bayrak \sphinxcode{\sphinxupquote{O\_SYNC}} ile birlikte kullanılırsa \sphinxcode{\sphinxupquote{read}} işlemini
etkileyecek olan daha önce yapılmış \sphinxcode{\sphinxupquote{write}} işlemleri varsa \sphinxcode{\sphinxupquote{read}} fonksiyonu geri dönmeden
önce bu \sphinxcode{\sphinxupquote{write}} işlemleri için \sphinxcode{\sphinxupquote{O\_SYNC}} bayrağında belirtilen semantik uygulanmaktadır.

\end{description}


\subsection{C Standart Kütüphanesinin Rolü}
\label{\detokenize{filesystem1:c-standart-kutuphanesinin-rolu}}
\sphinxAtStartPar
Peki C’nin standart dosya fonksiyonları bu süreçte nerede yer almaktadır? C’nin dosya fonksiyonları
aslında neticede POSIX fonksiyonlarını çağırmaktadır. Ancak C’nin standart dosya fonksiyonları
işletim sisteminin okuma yazma fonksiyonlarını daha az çağırmak için \sphinxstyleemphasis{kullanıcı alanında (user
space)} her dosya için bir önbellek de oluşturmaktadır. Bu önbellek sistemine genellikle önbellek
yerine \sphinxstyleemphasis{tamponlama (buffering)} sistemi, burada kullanılan önbelleğe de \sphinxstyleemphasis{tampon (buffer)}
denilmektedir.

\sphinxAtStartPar
Örneğin Linux’ta biz C’nin \sphinxcode{\sphinxupquote{getc}} gibi dosya fonksiyonunu çağırmış olalım. Standart C
kütüphanesi \sphinxcode{\sphinxupquote{fgetc}} ile 1 byte okumak istediğimizde \sphinxcode{\sphinxupquote{read}} POSIX fonksiyonu ile 1 byte
okumamaktadır. \sphinxcode{\sphinxupquote{read}} POSIX fonksiyonu ile \sphinxcode{\sphinxupquote{\textless{}stdio.h\textgreater{}}} dosyasında belirtilen \sphinxcode{\sphinxupquote{BUFSIZ}} kadar
byte’ı bir tampona okumakta ve oradan 1 byte’ı programcıya vermektedir. Bu nedenle C’nin standart
fonksiyonlarına \sphinxstyleemphasis{tamponlı IO (buffered IO)} fonksiyonları da denilmektedir.

\sphinxAtStartPar
O hâlde C’nin standart dosya fonksiyonlarıyla yapılan tipik bir okuma işlemi şöyle
gerçekleşmektedir:

\begin{sphinxVerbatim}[commandchars=\\\{\}]
C\PYGZsq{}deki okuma fonksiyonu
      ↓
read POSIX fonksiyonu
      ↓
sys\PYGZus{}read sistem fonksiyonu
      ↓
...
\end{sphinxVerbatim}

\sphinxAtStartPar
C’deki bu tamponlama yani önbellek mekanizması C’ye özgü değildir. Diğer programlama dillerinin
de standart kütüphanelerinde benzer biçimde tamponlamalar yapılmaktadır. Örneğin C++’taki
\sphinxcode{\sphinxupquote{\textless{}iostream\textgreater{}}} kütüphanesi, C\#’ta kullanılan .NET kütüphaneleri, Java’da kullanılan temel
kütüphanelerde hep kullanıcı alanında C’de olduğu gibi tamponlama yapmaktadır. Ancak bunların
hepsi neticede Linux sistemlerinde POSIX fonksiyonlarını, onlar da sistem fonksiyonlarını
çağırmaktadır.

\sphinxAtStartPar
POSIX dosya fonksiyonlarının Linux’taki işletim sisteminin sistem fonksiyonlarını çağırdığını
belirtmiştik. İşletim sisteminin sistem fonksiyonları bilgi önbellekte olsa bile belli bir
yavaşlık oluşturmaktadır. Programın akışının kullanıcı modundan çekirdek moduna geçirilmesi ve
akışın ilgili sistem fonksiyonuna aktarılması göreli bir zaman kaybına yol açmaktadır. Bu nedenle
bu kütüphanelerin kullanıcı modunda tamponlama yapması önemli olmaktadır.


\section{Çekirdekteki Dosya İşlemi Veri Yapıları}
\label{\detokenize{filesystem1:cekirdekteki-dosya-islemi-veri-yapilari}}
\sphinxAtStartPar
Şimdiye kadar blok ve sektör düzeyinde okuma yazmaların kabaca nasıl gerçekleştirildiğini
açıkladık. Ancak çekirdeğin açık dosyalar için oluşturduğu organizasyon hakkında bilgi vermedik.
Şimdi sürecin bu yönü üzerinde duracağız.

\sphinxAtStartPar
Anımsanacağı gibi UNIX/Linux sistemlerinde dosyalar \sphinxcode{\sphinxupquote{open}} isimli POSIX fonksiyonuyla
açılmaktadır. \sphinxcode{\sphinxupquote{open}} POSIX fonksiyonu başarı durumunda ismine \sphinxstyleemphasis{dosya betimleyicisi (file
descriptor)} denilen bir handle değeri vermektedir. \sphinxcode{\sphinxupquote{read}}, \sphinxcode{\sphinxupquote{write}}, \sphinxcode{\sphinxupquote{lseek}}, \sphinxcode{\sphinxupquote{close}}
gibi POSIX’in diğer dosya fonksiyonları bu dosya betimleyicisini parametre olarak alıp hangi
dosya üzerinde işlem yapılacağını bu betimleyiciden hareketle belirlemektedir.

\sphinxAtStartPar
\sphinxcode{\sphinxupquote{open}}, \sphinxcode{\sphinxupquote{read}}, \sphinxcode{\sphinxupquote{write}}, \sphinxcode{\sphinxupquote{lseek}}, \sphinxcode{\sphinxupquote{close}} gibi POSIX’in temel dosya fonksiyonları
Linux sistemlerinde aslında neredeyse doğrudan Linux’un ilgili sistem fonksiyonlarını
çağırmaktadır:

\begin{sphinxVerbatim}[commandchars=\\\{\}]
open  \PYGZhy{}\PYGZhy{}\PYGZhy{}\PYGZgt{} sys\PYGZus{}open
read  \PYGZhy{}\PYGZhy{}\PYGZhy{}\PYGZgt{} sys\PYGZus{}read
write \PYGZhy{}\PYGZhy{}\PYGZhy{}\PYGZgt{} sys\PYGZus{}write
lseek \PYGZhy{}\PYGZhy{}\PYGZhy{}\PYGZgt{} sys\PYGZus{}lseek
close \PYGZhy{}\PYGZhy{}\PYGZhy{}\PYGZgt{} sys\PYGZus{}close
\end{sphinxVerbatim}

\sphinxAtStartPar
Bu nedenle birtakım ayrıntıları da göz ardı edersek biz Linux sistemlerinde dosya işlemlerini
yapan temel POSIX fonksiyonlarının aslında doğrudan \sphinxcode{\sphinxupquote{sys\_xxx}} sistem fonksiyonlarını çağırdığını
varsayabiliriz.

\sphinxAtStartPar
UNIX/Linux sistemlerinde kullanılan standart C kütüphaneleri aynı zamanda POSIX fonksiyonlarını
da içermektedir. Bilindiği gibi bugün masaüstü Linux sistemlerinde en fazla kullanılan standart C
kütüphanesi GNU’nun \sphinxstyleemphasis{libc} kütüphanesidir. Eğer standart C fonksiyonlarının ve POSIX
fonksiyonlarının nasıl yazıldığını merak ediyorsanız gömülü Linux sistemleri için daha minimalist
biçimde yazılmış olan kütüphanelerin kaynak kodlarını inceleyebilirsiniz. Bu iş için iki alternatif
\sphinxstyleemphasis{musl} ve \sphinxstyleemphasis{uclibc} kütüphaneleridir. \sphinxstyleemphasis{uclibc} kütüphanesine \sphinxstyleemphasis{Mikro C kütüphanesi} de denilmektedir.
Bu kütüphanelerin kaynak kodlarını \sphinxhref{https://elixir.bootlin.com}{elixir.bootlin.com} sitesinden
inceleyebilirsiniz. Bu site yalnızca Linux çekirdekleri için değil başka projeler için de kodlar
üzerinde gezinme olanağı sunmaktadır. Sadeliği nedeniyle \sphinxstyleemphasis{musl} kütüphanesini incelemenizi salık
veririz. Kütüphanenin kodları üzerinde gezinebilmek için aşağıdaki bağlantıdan
faydalanabilirsiniz:

\sphinxAtStartPar
\sphinxurl{https://elixir.bootlin.com/musl/v1.2.5/source}


\section{task\_struct İçerisindeki Dosya Sistemine İlişkin Veri Yapıları}
\label{\detokenize{filesystem1:task-struct-icerisindeki-dosya-sistemine-iliskin-veri-yapilari}}
\sphinxAtStartPar
\sphinxcode{\sphinxupquote{task\_struct}} yapısı (proses kontrol bloğu) içerisinde proseslere ilişkin dosya işlemleri için
kullanılan iki önemli eleman bulunmaktadır:

\begin{sphinxVerbatim}[commandchars=\\\{\}]
\PYG{k}{struct}\PYG{+w}{ }\PYG{n+nc}{task\PYGZus{}struct}\PYG{+w}{ }\PYG{p}{\PYGZob{}}
\PYG{+w}{    }\PYG{c+cm}{/* ... */}

\PYG{+w}{    }\PYG{c+cm}{/* Filesystem information: */}
\PYG{+w}{    }\PYG{k}{struct}\PYG{+w}{ }\PYG{n+nc}{fs\PYGZus{}struct}\PYG{+w}{        }\PYG{o}{*}\PYG{n}{fs}\PYG{p}{;}

\PYG{+w}{    }\PYG{c+cm}{/* Open file information: */}
\PYG{+w}{    }\PYG{k}{struct}\PYG{+w}{ }\PYG{n+nc}{files\PYGZus{}struct}\PYG{+w}{     }\PYG{o}{*}\PYG{n}{files}\PYG{p}{;}

\PYG{+w}{    }\PYG{c+cm}{/* ... */}
\PYG{p}{\PYGZcb{}}\PYG{p}{;}
\end{sphinxVerbatim}

\sphinxAtStartPar
Bu iki eleman çok uzun süredir \sphinxcode{\sphinxupquote{task\_struct}} yapısı içerisinde bulunmaktadır. Ancak buradaki
\sphinxcode{\sphinxupquote{fs\_struct}} ve \sphinxcode{\sphinxupquote{files\_struct}} yapılarının içeriğinde çekirdeğin versiyonları ilerledikçe
çeşitli değişiklikler de yapılmıştır.


\subsection{fs\_struct Yapısı}
\label{\detokenize{filesystem1:fs-struct-yapisi}}
\sphinxAtStartPar
Yuradaki \sphinxcode{\sphinxupquote{fs\_struct}} yapısı açık dosyalara ilişkin yapılan organizasyonla ilgili değildir.
Prosesin kök dizini ve çalışma dizini gibi dosya sistemine ilişkin proses bilgileri burada
tutulmaktadır. Mevcut çekirdeklerde \sphinxcode{\sphinxupquote{fs\_struct}} yapısı \sphinxcode{\sphinxupquote{include/linux/fs\_struct.h}} dosyası
içerisinde şöyle bildirilmiştir:

\begin{sphinxVerbatim}[commandchars=\\\{\}]
\PYG{k}{struct}\PYG{+w}{ }\PYG{n+nc}{fs\PYGZus{}struct}\PYG{+w}{ }\PYG{p}{\PYGZob{}}
\PYG{+w}{    }\PYG{k+kt}{int}\PYG{+w}{ }\PYG{n}{users}\PYG{p}{;}
\PYG{+w}{    }\PYG{n}{spinlock\PYGZus{}t}\PYG{+w}{ }\PYG{n}{lock}\PYG{p}{;}
\PYG{+w}{    }\PYG{n}{seqcount\PYGZus{}spinlock\PYGZus{}t}\PYG{+w}{ }\PYG{n}{seq}\PYG{p}{;}
\PYG{+w}{    }\PYG{k+kt}{int}\PYG{+w}{ }\PYG{n}{umask}\PYG{p}{;}
\PYG{+w}{    }\PYG{k+kt}{int}\PYG{+w}{ }\PYG{n}{in\PYGZus{}exec}\PYG{p}{;}
\PYG{+w}{    }\PYG{k}{struct}\PYG{+w}{ }\PYG{n+nc}{path}\PYG{+w}{ }\PYG{n}{root}\PYG{p}{,}\PYG{+w}{ }\PYG{n}{pwd}\PYG{p}{;}
\PYG{p}{\PYGZcb{}}\PYG{+w}{ }\PYG{n}{\PYGZus{}\PYGZus{}randomize\PYGZus{}layout}\PYG{p}{;}
\end{sphinxVerbatim}

\sphinxAtStartPar
Buradaki \sphinxcode{\sphinxupquote{root}} ve \sphinxcode{\sphinxupquote{pwd}} elemanları sırasıyla prosesin kök dizinini ve çalışma dizinini
(current working directory) tutmaktadır. \sphinxcode{\sphinxupquote{umask}} elemanı ise prosesin umask değerini
tutmaktadır. Buradaki \sphinxcode{\sphinxupquote{path}} yapısı da şöyle bildirilmiştir:

\begin{sphinxVerbatim}[commandchars=\\\{\}]
\PYG{k}{struct}\PYG{+w}{ }\PYG{n+nc}{path}\PYG{+w}{ }\PYG{p}{\PYGZob{}}
\PYG{+w}{    }\PYG{k}{struct}\PYG{+w}{ }\PYG{n+nc}{vfsmount}\PYG{+w}{ }\PYG{o}{*}\PYG{n}{mnt}\PYG{p}{;}
\PYG{+w}{    }\PYG{k}{struct}\PYG{+w}{ }\PYG{n+nc}{dentry}\PYG{+w}{ }\PYG{o}{*}\PYG{n}{dentry}\PYG{p}{;}
\PYG{p}{\PYGZcb{}}\PYG{+w}{ }\PYG{n}{\PYGZus{}\PYGZus{}randomize\PYGZus{}layout}\PYG{p}{;}
\end{sphinxVerbatim}

\sphinxAtStartPar
\sphinxcode{\sphinxupquote{vfsmount}} ve \sphinxcode{\sphinxupquote{dentry}} yapıları çeşitli başka bölümlerde ele alınacaktır. Ayrıntıları göz
ardı edersek \sphinxcode{\sphinxupquote{fs\_struct}} yapısındaki en önemli elemanlar prosesin kök dizininin yeri, prosesin
çalışma dizininin yeri ve prosesin umask değeridir.

\sphinxAtStartPar
Bu yapı zaman içerisinde de geliştirilmiştir. Çekirdeğin 0.01 versiyonunda bu bilgiler doğrudan
\sphinxcode{\sphinxupquote{task\_struct}} içerisinde bulunmaktaydı:

\begin{sphinxVerbatim}[commandchars=\\\{\}]
\PYG{k}{struct}\PYG{+w}{ }\PYG{n+nc}{task\PYGZus{}struct}\PYG{+w}{ }\PYG{p}{\PYGZob{}}
\PYG{+w}{    }\PYG{c+cm}{/* ... */}

\PYG{+w}{    }\PYG{k+kt}{unsigned}\PYG{+w}{ }\PYG{k+kt}{short}\PYG{+w}{ }\PYG{n}{umask}\PYG{p}{;}
\PYG{+w}{    }\PYG{k}{struct}\PYG{+w}{ }\PYG{n+nc}{m\PYGZus{}inode}\PYG{+w}{ }\PYG{o}{*}\PYG{+w}{ }\PYG{n}{pwd}\PYG{p}{;}
\PYG{+w}{    }\PYG{k}{struct}\PYG{+w}{ }\PYG{n+nc}{m\PYGZus{}inode}\PYG{+w}{ }\PYG{o}{*}\PYG{+w}{ }\PYG{n}{root}\PYG{p}{;}
\PYG{+w}{    }\PYG{k+kt}{unsigned}\PYG{+w}{ }\PYG{k+kt}{long}\PYG{+w}{ }\PYG{n}{close\PYGZus{}on\PYGZus{}exec}\PYG{p}{;}

\PYG{+w}{    }\PYG{c+cm}{/* ... */}
\PYG{p}{\PYGZcb{}}\PYG{p}{;}
\end{sphinxVerbatim}


\section{files\_struct Yapısı}
\label{\detokenize{filesystem1:files-struct-yapisi}}
\sphinxAtStartPar
\sphinxcode{\sphinxupquote{task\_struct}} içerisindeki \sphinxcode{\sphinxupquote{files}} isimli gösterici prosesin açmış olduğu dosyalara ilişkin
bilgilerin tutulduğu \sphinxcode{\sphinxupquote{files\_struct}} türünden yapı nesnesini göstermektedir. \sphinxcode{\sphinxupquote{files\_struct}}
yapısı da zaman içerisinde değişiklere uğratılmıştır. Güncel çekirdeklerde bu yapı
\sphinxcode{\sphinxupquote{include/linux/fdtable.h}} dosyası içerisinde şöyle bildirilmiştir:

\begin{sphinxVerbatim}[commandchars=\\\{\}]
\PYG{k}{struct}\PYG{+w}{ }\PYG{n+nc}{files\PYGZus{}struct}\PYG{+w}{ }\PYG{p}{\PYGZob{}}
\PYG{c+cm}{/*}
\PYG{c+cm}{ * read mostly part}
\PYG{c+cm}{ */}
\PYG{+w}{    }\PYG{n}{atomic\PYGZus{}t}\PYG{+w}{ }\PYG{n}{count}\PYG{p}{;}
\PYG{+w}{    }\PYG{k+kt}{bool}\PYG{+w}{ }\PYG{n}{resize\PYGZus{}in\PYGZus{}progress}\PYG{p}{;}
\PYG{+w}{    }\PYG{n}{wait\PYGZus{}queue\PYGZus{}head\PYGZus{}t}\PYG{+w}{ }\PYG{n}{resize\PYGZus{}wait}\PYG{p}{;}

\PYG{+w}{    }\PYG{k}{struct}\PYG{+w}{ }\PYG{n+nc}{fdtable}\PYG{+w}{ }\PYG{n}{\PYGZus{}\PYGZus{}rcu}\PYG{+w}{ }\PYG{o}{*}\PYG{n}{fdt}\PYG{p}{;}
\PYG{+w}{    }\PYG{k}{struct}\PYG{+w}{ }\PYG{n+nc}{fdtable}\PYG{+w}{ }\PYG{n}{fdtab}\PYG{p}{;}
\PYG{c+cm}{/*}
\PYG{c+cm}{ * written part on a separate cache line in SMP}
\PYG{c+cm}{ */}
\PYG{+w}{    }\PYG{n}{spinlock\PYGZus{}t}\PYG{+w}{ }\PYG{n}{file\PYGZus{}lock}\PYG{+w}{ }\PYG{n}{\PYGZus{}\PYGZus{}\PYGZus{}\PYGZus{}cacheline\PYGZus{}aligned\PYGZus{}in\PYGZus{}smp}\PYG{p}{;}
\PYG{+w}{    }\PYG{k+kt}{unsigned}\PYG{+w}{ }\PYG{k+kt}{int}\PYG{+w}{ }\PYG{n}{next\PYGZus{}fd}\PYG{p}{;}
\PYG{+w}{    }\PYG{k+kt}{unsigned}\PYG{+w}{ }\PYG{k+kt}{long}\PYG{+w}{ }\PYG{n}{close\PYGZus{}on\PYGZus{}exec\PYGZus{}init}\PYG{p}{[}\PYG{l+m+mi}{1}\PYG{p}{]}\PYG{p}{;}
\PYG{+w}{    }\PYG{k+kt}{unsigned}\PYG{+w}{ }\PYG{k+kt}{long}\PYG{+w}{ }\PYG{n}{open\PYGZus{}fds\PYGZus{}init}\PYG{p}{[}\PYG{l+m+mi}{1}\PYG{p}{]}\PYG{p}{;}
\PYG{+w}{    }\PYG{k+kt}{unsigned}\PYG{+w}{ }\PYG{k+kt}{long}\PYG{+w}{ }\PYG{n}{full\PYGZus{}fds\PYGZus{}bits\PYGZus{}init}\PYG{p}{[}\PYG{l+m+mi}{1}\PYG{p}{]}\PYG{p}{;}
\PYG{+w}{    }\PYG{k}{struct}\PYG{+w}{ }\PYG{n+nc}{file}\PYG{+w}{ }\PYG{n}{\PYGZus{}\PYGZus{}rcu}\PYG{+w}{ }\PYG{o}{*}\PYG{+w}{ }\PYG{n}{fd\PYGZus{}array}\PYG{p}{[}\PYG{n}{NR\PYGZus{}OPEN\PYGZus{}DEFAULT}\PYG{p}{]}\PYG{p}{;}
\PYG{p}{\PYGZcb{}}\PYG{p}{;}
\end{sphinxVerbatim}

\sphinxAtStartPar
Bu yapı versiyonlar arasında önemli farklılıklar göstermiştir. 2.4 çekirdeğindeki hali:

\begin{sphinxVerbatim}[commandchars=\\\{\}]
\PYG{k}{struct}\PYG{+w}{ }\PYG{n+nc}{files\PYGZus{}struct}\PYG{+w}{ }\PYG{p}{\PYGZob{}}
\PYG{+w}{    }\PYG{n}{atomic\PYGZus{}t}\PYG{+w}{ }\PYG{n}{count}\PYG{p}{;}
\PYG{+w}{    }\PYG{n}{rwlock\PYGZus{}t}\PYG{+w}{ }\PYG{n}{file\PYGZus{}lock}\PYG{p}{;}
\PYG{+w}{    }\PYG{k+kt}{int}\PYG{+w}{ }\PYG{n}{max\PYGZus{}fds}\PYG{p}{;}
\PYG{+w}{    }\PYG{k+kt}{int}\PYG{+w}{ }\PYG{n}{max\PYGZus{}fdset}\PYG{p}{;}
\PYG{+w}{    }\PYG{k+kt}{int}\PYG{+w}{ }\PYG{n}{next\PYGZus{}fd}\PYG{p}{;}
\PYG{+w}{    }\PYG{k}{struct}\PYG{+w}{ }\PYG{n+nc}{file}\PYG{+w}{ }\PYG{o}{*}\PYG{o}{*}\PYG{+w}{ }\PYG{n}{fd}\PYG{p}{;}\PYG{+w}{          }\PYG{c+cm}{/* current fd array */}
\PYG{+w}{    }\PYG{n}{fd\PYGZus{}set}\PYG{+w}{ }\PYG{o}{*}\PYG{n}{close\PYGZus{}on\PYGZus{}exec}\PYG{p}{;}
\PYG{+w}{    }\PYG{n}{fd\PYGZus{}set}\PYG{+w}{ }\PYG{o}{*}\PYG{n}{open\PYGZus{}fds}\PYG{p}{;}
\PYG{+w}{    }\PYG{n}{fd\PYGZus{}set}\PYG{+w}{ }\PYG{n}{close\PYGZus{}on\PYGZus{}exec\PYGZus{}init}\PYG{p}{;}
\PYG{+w}{    }\PYG{n}{fd\PYGZus{}set}\PYG{+w}{ }\PYG{n}{open\PYGZus{}fds\PYGZus{}init}\PYG{p}{;}
\PYG{+w}{    }\PYG{k}{struct}\PYG{+w}{ }\PYG{n+nc}{file}\PYG{+w}{ }\PYG{o}{*}\PYG{+w}{ }\PYG{n}{fd\PYGZus{}array}\PYG{p}{[}\PYG{n}{NR\PYGZus{}OPEN\PYGZus{}DEFAULT}\PYG{p}{]}\PYG{p}{;}
\PYG{p}{\PYGZcb{}}\PYG{p}{;}
\end{sphinxVerbatim}

\sphinxAtStartPar
Çekirdeğin 0.01 versiyonunda bu bilgiler doğrudan \sphinxcode{\sphinxupquote{task\_struct}} içerisinde bulunuyordu:

\begin{sphinxVerbatim}[commandchars=\\\{\}]
\PYG{k}{struct}\PYG{+w}{ }\PYG{n+nc}{task\PYGZus{}struct}\PYG{+w}{ }\PYG{p}{\PYGZob{}}
\PYG{+w}{    }\PYG{c+cm}{/* ... */}

\PYG{+w}{    }\PYG{k+kt}{unsigned}\PYG{+w}{ }\PYG{k+kt}{short}\PYG{+w}{ }\PYG{n}{umask}\PYG{p}{;}
\PYG{+w}{    }\PYG{k}{struct}\PYG{+w}{ }\PYG{n+nc}{m\PYGZus{}inode}\PYG{+w}{ }\PYG{o}{*}\PYG{+w}{ }\PYG{n}{pwd}\PYG{p}{;}
\PYG{+w}{    }\PYG{k}{struct}\PYG{+w}{ }\PYG{n+nc}{m\PYGZus{}inode}\PYG{+w}{ }\PYG{o}{*}\PYG{+w}{ }\PYG{n}{root}\PYG{p}{;}
\PYG{+w}{    }\PYG{k+kt}{unsigned}\PYG{+w}{ }\PYG{k+kt}{long}\PYG{+w}{ }\PYG{n}{close\PYGZus{}on\PYGZus{}exec}\PYG{p}{;}

\PYG{+w}{    }\PYG{c+cm}{/* ... */}
\PYG{p}{\PYGZcb{}}\PYG{p}{;}
\end{sphinxVerbatim}


\section{Dosya Nesnesi (struct file)}
\label{\detokenize{filesystem1:dosya-nesnesi-struct-file}}
\sphinxAtStartPar
Linux’ta ne zaman \sphinxcode{\sphinxupquote{open}} POSIX fonksiyonuyla bir dosya açılsa \sphinxcode{\sphinxupquote{sys\_open}} sistem fonksiyonu
açılan dosya için \sphinxcode{\sphinxupquote{file}} isimli (\sphinxcode{\sphinxupquote{struct file}} türünden) bir yapı nesnesini tahsis edip
dosya işlemleri için gereken bilgileri bu yapı nesnesinin içerisine yerleştirmektedir.
\sphinxcode{\sphinxupquote{sys\_read}}, \sphinxcode{\sphinxupquote{sys\_write}}, \sphinxcode{\sphinxupquote{sys\_lseek}}, \sphinxcode{\sphinxupquote{sys\_close}} gibi sistem fonksiyonları da dosya
üzerinde işlem yapabilmek için bu \sphinxcode{\sphinxupquote{file}} yapısındaki bilgileri kullanmaktadır. İşletim
sistemlerinde bu amaçla kullanılan nesnelere \sphinxstyleemphasis{dosya nesnesi (file object)} de denilmektedir.

\sphinxAtStartPar
Biz aşağıdaki gibi bir dosya açmış olalım:

\begin{sphinxVerbatim}[commandchars=\\\{\}]
\PYG{n}{fd}\PYG{+w}{ }\PYG{o}{=}\PYG{+w}{ }\PYG{n}{open}\PYG{p}{(}\PYG{p}{.}\PYG{p}{.}\PYG{p}{.}\PYG{p}{)}\PYG{p}{;}
\end{sphinxVerbatim}

\sphinxAtStartPar
\sphinxcode{\sphinxupquote{sys\_open}} sistem fonksiyonu açılmak istenen dosyanın diskteki yerini ve metadata bilgilerini
bulur, o bilgilerden hareketle bir dosya nesnesi (file object) oluşturur. O dosya nesnesinin
adresini de \sphinxcode{\sphinxupquote{files\_struct}} nesnesinin içerisine yerleştirir. Böylece \sphinxcode{\sphinxupquote{sys\_read}},
\sphinxcode{\sphinxupquote{sys\_write}}, \sphinxcode{\sphinxupquote{sys\_lseek}}, \sphinxcode{\sphinxupquote{sys\_close}} gibi sistem fonksiyonları \sphinxcode{\sphinxupquote{task\_struct}} nesnesinden
hareketle bu dosya nesnesine erişebilmektedir.

\sphinxAtStartPar
Güncel çekirdeklerde \sphinxcode{\sphinxupquote{file}} yapısı \sphinxcode{\sphinxupquote{include/linux/fs.h}} dosyasının içerisinde şöyle
bildirilmiştir:

\begin{sphinxVerbatim}[commandchars=\\\{\}]
\PYG{k}{struct}\PYG{+w}{ }\PYG{n+nc}{file}\PYG{+w}{ }\PYG{p}{\PYGZob{}}
\PYG{+w}{    }\PYG{n}{spinlock\PYGZus{}t}\PYG{+w}{                   }\PYG{n}{f\PYGZus{}lock}\PYG{p}{;}
\PYG{+w}{    }\PYG{n}{fmode\PYGZus{}t}\PYG{+w}{                      }\PYG{n}{f\PYGZus{}mode}\PYG{p}{;}
\PYG{+w}{    }\PYG{k}{const}\PYG{+w}{ }\PYG{k}{struct}\PYG{+w}{ }\PYG{n+nc}{file\PYGZus{}operations}\PYG{+w}{ }\PYG{o}{*}\PYG{n}{f\PYGZus{}op}\PYG{p}{;}
\PYG{+w}{    }\PYG{k}{struct}\PYG{+w}{ }\PYG{n+nc}{address\PYGZus{}space}\PYG{+w}{        }\PYG{o}{*}\PYG{n}{f\PYGZus{}mapping}\PYG{p}{;}
\PYG{+w}{    }\PYG{k+kt}{void}\PYG{+w}{                        }\PYG{o}{*}\PYG{n}{private\PYGZus{}data}\PYG{p}{;}
\PYG{+w}{    }\PYG{k}{struct}\PYG{+w}{ }\PYG{n+nc}{inode}\PYG{+w}{                }\PYG{o}{*}\PYG{n}{f\PYGZus{}inode}\PYG{p}{;}
\PYG{+w}{    }\PYG{k+kt}{unsigned}\PYG{+w}{ }\PYG{k+kt}{int}\PYG{+w}{                 }\PYG{n}{f\PYGZus{}flags}\PYG{p}{;}
\PYG{+w}{    }\PYG{k+kt}{unsigned}\PYG{+w}{ }\PYG{k+kt}{int}\PYG{+w}{                 }\PYG{n}{f\PYGZus{}iocb\PYGZus{}flags}\PYG{p}{;}
\PYG{+w}{    }\PYG{k}{const}\PYG{+w}{ }\PYG{k}{struct}\PYG{+w}{ }\PYG{n+nc}{cred}\PYG{+w}{           }\PYG{o}{*}\PYG{n}{f\PYGZus{}cred}\PYG{p}{;}
\PYG{+w}{    }\PYG{k}{struct}\PYG{+w}{ }\PYG{n+nc}{fown\PYGZus{}struct}\PYG{+w}{          }\PYG{o}{*}\PYG{n}{f\PYGZus{}owner}\PYG{p}{;}
\PYG{+w}{    }\PYG{c+cm}{/* \PYGZhy{}\PYGZhy{}\PYGZhy{} cacheline 1 boundary (64 bytes) \PYGZhy{}\PYGZhy{}\PYGZhy{} */}
\PYG{+w}{    }\PYG{k}{struct}\PYG{+w}{ }\PYG{n+nc}{path}\PYG{+w}{                  }\PYG{n}{f\PYGZus{}path}\PYG{p}{;}
\PYG{+w}{    }\PYG{k}{union}\PYG{+w}{ }\PYG{p}{\PYGZob{}}
\PYG{+w}{        }\PYG{k}{struct}\PYG{+w}{ }\PYG{n+nc}{mutex}\PYG{+w}{             }\PYG{n}{f\PYGZus{}pos\PYGZus{}lock}\PYG{p}{;}
\PYG{+w}{        }\PYG{n}{u64}\PYG{+w}{                      }\PYG{n}{f\PYGZus{}pipe}\PYG{p}{;}
\PYG{+w}{    }\PYG{p}{\PYGZcb{}}\PYG{p}{;}
\PYG{+w}{    }\PYG{n}{loff\PYGZus{}t}\PYG{+w}{                       }\PYG{n}{f\PYGZus{}pos}\PYG{p}{;}
\PYG{c+cp}{\PYGZsh{}}\PYG{c+cp}{ifdef CONFIG\PYGZus{}SECURITY}
\PYG{+w}{    }\PYG{k+kt}{void}\PYG{+w}{                        }\PYG{o}{*}\PYG{n}{f\PYGZus{}security}\PYG{p}{;}
\PYG{c+cp}{\PYGZsh{}}\PYG{c+cp}{endif}
\PYG{+w}{    }\PYG{c+cm}{/* \PYGZhy{}\PYGZhy{}\PYGZhy{} cacheline 2 boundary (128 bytes) \PYGZhy{}\PYGZhy{}\PYGZhy{} */}
\PYG{+w}{    }\PYG{n}{errseq\PYGZus{}t}\PYG{+w}{                     }\PYG{n}{f\PYGZus{}wb\PYGZus{}err}\PYG{p}{;}
\PYG{+w}{    }\PYG{n}{errseq\PYGZus{}t}\PYG{+w}{                     }\PYG{n}{f\PYGZus{}sb\PYGZus{}err}\PYG{p}{;}
\PYG{c+cp}{\PYGZsh{}}\PYG{c+cp}{ifdef CONFIG\PYGZus{}EPOLL}
\PYG{+w}{    }\PYG{k}{struct}\PYG{+w}{ }\PYG{n+nc}{hlist\PYGZus{}head}\PYG{+w}{           }\PYG{o}{*}\PYG{n}{f\PYGZus{}ep}\PYG{p}{;}
\PYG{c+cp}{\PYGZsh{}}\PYG{c+cp}{endif}
\PYG{+w}{    }\PYG{k}{union}\PYG{+w}{ }\PYG{p}{\PYGZob{}}
\PYG{+w}{        }\PYG{k}{struct}\PYG{+w}{ }\PYG{n+nc}{callback\PYGZus{}head}\PYG{+w}{     }\PYG{n}{f\PYGZus{}task\PYGZus{}work}\PYG{p}{;}
\PYG{+w}{        }\PYG{k}{struct}\PYG{+w}{ }\PYG{n+nc}{llist\PYGZus{}node}\PYG{+w}{        }\PYG{n}{f\PYGZus{}llist}\PYG{p}{;}
\PYG{+w}{        }\PYG{k}{struct}\PYG{+w}{ }\PYG{n+nc}{file\PYGZus{}ra\PYGZus{}state}\PYG{+w}{     }\PYG{n}{f\PYGZus{}ra}\PYG{p}{;}
\PYG{+w}{        }\PYG{n}{freeptr\PYGZus{}t}\PYG{+w}{                }\PYG{n}{f\PYGZus{}freeptr}\PYG{p}{;}
\PYG{+w}{    }\PYG{p}{\PYGZcb{}}\PYG{p}{;}
\PYG{+w}{    }\PYG{n}{file\PYGZus{}ref\PYGZus{}t}\PYG{+w}{                   }\PYG{n}{f\PYGZus{}ref}\PYG{p}{;}
\PYG{+w}{    }\PYG{c+cm}{/* \PYGZhy{}\PYGZhy{}\PYGZhy{} cacheline 3 boundary (192 bytes) \PYGZhy{}\PYGZhy{}\PYGZhy{} */}
\PYG{p}{\PYGZcb{}}\PYG{+w}{ }\PYG{n}{\PYGZus{}\PYGZus{}randomize\PYGZus{}layout}
\PYG{n}{\PYGZus{}\PYGZus{}attribute\PYGZus{}\PYGZus{}}\PYG{p}{(}\PYG{p}{(}\PYG{n}{aligned}\PYG{p}{(}\PYG{l+m+mi}{4}\PYG{p}{)}\PYG{p}{)}\PYG{p}{)}\PYG{p}{;}
\end{sphinxVerbatim}


\section{Dosya Betimleyici Tablosu}
\label{\detokenize{filesystem1:dosya-betimleyici-tablosu}}
\sphinxAtStartPar
Şimdi bir dosya \sphinxcode{\sphinxupquote{sys\_open}} sistem fonksiyonuyla açıldığında dosya nesnesinin ve dosya
betimleyicisinin nasıl saklandığını açıklayalım. Önce çekirdeğin 0.01 versiyonunu inceleyelim.
Bu ilkel versiyonda henüz \sphinxcode{\sphinxupquote{files\_struct}} biçiminde bir yapı yoktu. Açık dosya bilgileri doğrudan
\sphinxcode{\sphinxupquote{task\_struct}} içerisinde bulunan aşağıdaki elemanlarda saklanıyordu:

\begin{sphinxVerbatim}[commandchars=\\\{\}]
\PYG{k}{struct}\PYG{+w}{ }\PYG{n+nc}{task\PYGZus{}struct}\PYG{+w}{ }\PYG{p}{\PYGZob{}}
\PYG{+w}{    }\PYG{c+cm}{/* ... */}

\PYG{+w}{    }\PYG{k+kt}{unsigned}\PYG{+w}{ }\PYG{k+kt}{long}\PYG{+w}{ }\PYG{n}{close\PYGZus{}on\PYGZus{}exec}\PYG{p}{;}
\PYG{+w}{    }\PYG{k}{struct}\PYG{+w}{ }\PYG{n+nc}{file}\PYG{+w}{ }\PYG{o}{*}\PYG{n}{filp}\PYG{p}{[}\PYG{n}{NR\PYGZus{}OPEN}\PYG{p}{]}\PYG{p}{;}

\PYG{+w}{    }\PYG{c+cm}{/* ... */}
\PYG{p}{\PYGZcb{}}\PYG{p}{;}
\end{sphinxVerbatim}

\sphinxAtStartPar
Burada \sphinxcode{\sphinxupquote{filp}} isimli dizinin \sphinxcode{\sphinxupquote{struct file *}} türünden olduğuna dikkat ediniz. Yani \sphinxcode{\sphinxupquote{filp}}
dizisi file nesnelerinin adreslerini tutan bir gösterici dizisidir. Bu versiyonda \sphinxcode{\sphinxupquote{NR\_OPEN}}
şöyle tanımlanmıştır:

\begin{sphinxVerbatim}[commandchars=\\\{\}]
\PYG{c+cp}{\PYGZsh{}}\PYG{c+cp}{define NR\PYGZus{}OPEN 20}
\end{sphinxVerbatim}

\sphinxAtStartPar
Dolayısıyla bu ilkel versiyonda bir proses en fazla 20 dosyayı açık durumda tutabiliyordu.
UNIX/Linux dünyasında dosya nesnelerinin adreslerini tutan bu gösterici dizilerine \sphinxstyleemphasis{dosya
betimleyici tablosu (file descriptor table)} denilmektedir.

\sphinxincludegraphics[]{graphviz-dfe8286a0b4b844362f9d759a0f69989cb8a5721.pdf}

\sphinxAtStartPar
\sphinxcode{\sphinxupquote{open}} POSIX fonksiyonunun (yani \sphinxcode{\sphinxupquote{sys\_open}} sistem fonksiyonunun) verdiği \sphinxstyleemphasis{dosya betimleyicisi
(file descriptor)} aslında dosya betimleyici tablosundaki dizide bir indeks belirtmektedir.
\sphinxcode{\sphinxupquote{open}} POSIX fonksiyonunun dosya betimleyici tablosundaki en düşük boş indeksi vereceği POSIX
standartlarında garanti altına alınmıştır.

\sphinxAtStartPar
Dosya betimleyici tablosunun \sphinxstyleemphasis{prosese özgü} olduğuna dikkat ediniz. Bir proseste açılmış olan
dosyaya ilişkin dosya nesnesinin adresi o prosesteki dosya betimleyici tablosuna yazılmaktadır.
Dosya betimleyicileri sistem genelinde bir değer belirtmemektedir; dosya betimleyici değerleri
yalnızca ilgili proses için anlamlıdır.

\sphinxAtStartPar
Yukarıdaki 0.01 versiyonunda konuyla ilgili \sphinxcode{\sphinxupquote{unsigned long}} türden \sphinxcode{\sphinxupquote{close\_on\_exec}} isimli bir
elemanın da bulunduğunu görüyorsunuz. Bu elemanın her biti bir betimleyicinin \sphinxstyleemphasis{close\sphinxhyphen{}on\sphinxhyphen{}exec}
durumunu belirtmektedir. Söz konusu bit 1 ise ilgili betimleyici exec işlemleri sırasında close
edilir, 0 ise close edilmez. POSIX standartlarında bir dosya açıldığında close\sphinxhyphen{}on\sphinxhyphen{}exec bayrağının
varsayılan durumda 0 olduğu belirtilmiştir.


\subsection{İlk Boş Betimleyicinin Bulunması}
\label{\detokenize{filesystem1:ilk-bos-betimleyicinin-bulunmasi}}
\sphinxAtStartPar
\sphinxcode{\sphinxupquote{sys\_open}} sistem fonksiyonu öncelikle dosya betimleyici tablosundaki ilk boş betimleyiciyi
bulmaya çalışır. 0.01 versiyonunda boş betimleyici şöyle bulunmuştur:

\begin{sphinxVerbatim}[commandchars=\\\{\}]
\PYG{k+kt}{int}\PYG{+w}{ }\PYG{n+nf}{sys\PYGZus{}open}\PYG{p}{(}\PYG{k}{const}\PYG{+w}{ }\PYG{k+kt}{char}\PYG{+w}{ }\PYG{o}{*}\PYG{+w}{ }\PYG{n}{filename}\PYG{p}{,}\PYG{+w}{ }\PYG{k+kt}{int}\PYG{+w}{ }\PYG{n}{flag}\PYG{p}{,}\PYG{+w}{ }\PYG{k+kt}{int}\PYG{+w}{ }\PYG{n}{mode}\PYG{p}{)}\PYG{+w}{ }\PYG{p}{\PYGZob{}}
\PYG{+w}{    }\PYG{c+cm}{/* ... */}

\PYG{+w}{    }\PYG{k}{for}\PYG{p}{(}\PYG{n}{fd}\PYG{o}{=}\PYG{l+m+mi}{0}\PYG{+w}{ }\PYG{p}{;}\PYG{+w}{ }\PYG{n}{fd}\PYG{o}{\PYGZlt{}}\PYG{n}{NR\PYGZus{}OPEN}\PYG{+w}{ }\PYG{p}{;}\PYG{+w}{ }\PYG{n}{fd}\PYG{o}{+}\PYG{o}{+}\PYG{p}{)}
\PYG{+w}{        }\PYG{k}{if}\PYG{+w}{ }\PYG{p}{(}\PYG{o}{!}\PYG{n}{current}\PYG{o}{\PYGZhy{}}\PYG{o}{\PYGZgt{}}\PYG{n}{filp}\PYG{p}{[}\PYG{n}{fd}\PYG{p}{]}\PYG{p}{)}
\PYG{+w}{            }\PYG{k}{break}\PYG{p}{;}

\PYG{+w}{    }\PYG{c+cm}{/* ... */}
\PYG{p}{\PYGZcb{}}
\end{sphinxVerbatim}

\sphinxAtStartPar
Görüldüğü gibi bu ilkel versiyonda dosya betimleyici tablosu üzerinde tek tek sıralı arama
yapılmış, ilk boş betimleyici elde edilmiştir. Ancak uzunca bir süredir proseslerin varsayılan
dosya betimleyici tablolarının varsayılan uzunlukları 1024’tür ve bu uzunluk da büyütülebilmektedir.


\subsection{Bitmap Tabanlı Hızlı Arama}
\label{\detokenize{filesystem1:bitmap-tabanli-hizli-arama}}
\sphinxAtStartPar
1024 elemanlı bir tabloda sıralı arama ile ilk boş elemanın bulunması yavaş bir işlemdir. Bu
nedenle bir süre sonra Linux çekirdeklerinde bu arama işlemi bit düzeyinde aramayla
hızlandırılmıştır.

\sphinxAtStartPar
Bit düzeyinde arama yönteminde dosya betimleyici tablosunun uzunluğu kadar bit dizisi
oluşturulur. Bu bit dizisindeki 0 olan bitler betimleyici tablosundaki boş elemanları, 1 olan
bitler dolu elemanları belirtmektedir.

\sphinxAtStartPar
Güncel çekirdeklerde iki düzey bitmap kullanılmaktadır:
\begin{itemize}
\item {} 
\sphinxAtStartPar
\sphinxstylestrong{Birinci düzey bitmap (open\_fds):} Tüm dosya betimleyicilerinin dolu/boş olduğunu tutar.

\item {} 
\sphinxAtStartPar
\sphinxstylestrong{İkinci düzey bitmap (full\_fds\_bits):} Birinci düzey bitmap’teki tüm bitleri 0 olmayan
(yani tamamen dolu) \sphinxcode{\sphinxupquote{unsigned long}} elemanların indeksini hızlıca bulmak için kullanılır.

\end{itemize}

\sphinxAtStartPar
Bu hızlandırma mantığında:
\begin{enumerate}
\sphinxsetlistlabels{\arabic}{enumi}{enumii}{}{.}%
\item {} 
\sphinxAtStartPar
Çekirdek hemen ikinci düzey bitmap’e yönelmez. Önce \sphinxcode{\sphinxupquote{open\_fds}} dizisinde tek bir
\sphinxcode{\sphinxupquote{unsigned long}} elemanda hızlı bir arama yapar.

\item {} 
\sphinxAtStartPar
Eğer bu arama başarısız olursa \sphinxcode{\sphinxupquote{full\_fds\_bits}}’te ilk 0 olan bitin indeksi elde edilir;
birinci düzey bitmap’te yalnızca bu indeksteki \sphinxcode{\sphinxupquote{unsigned long}} dizi elemanında arama yapılır.

\end{enumerate}

\sphinxAtStartPar
Güncel çekirdeklerdeki \sphinxcode{\sphinxupquote{sys\_open}} sistem fonksiyonundan başlanarak ilk boş dosya betimleyicisinin
bulunması için yapılan çağrılar şöyledir:

\begin{sphinxVerbatim}[commandchars=\\\{\}]
sys\PYGZus{}open
  \(\rightarrow\) do\PYGZus{}sys\PYGZus{}open
    \(\rightarrow\) do\PYGZus{}sys\PYGZus{}openat2
      \(\rightarrow\) \PYGZus{}\PYGZus{}get\PYGZus{}unused\PYGZus{}fd\PYGZus{}flags
        \(\rightarrow\) alloc\PYGZus{}fd
          \(\rightarrow\) find\PYGZus{}next\PYGZus{}fd
            \(\rightarrow\) find\PYGZus{}next\PYGZus{}zero\PYGZus{}bit
\end{sphinxVerbatim}

\sphinxAtStartPar
Güncel çekirdeklerdeki \sphinxcode{\sphinxupquote{find\_next\_fd}} fonksiyonu şöyle yazılmıştır:

\begin{sphinxVerbatim}[commandchars=\\\{\}]
\PYG{k}{static}\PYG{+w}{ }\PYG{k+kt}{unsigned}\PYG{+w}{ }\PYG{k+kt}{int}\PYG{+w}{ }\PYG{n+nf}{find\PYGZus{}next\PYGZus{}fd}\PYG{p}{(}\PYG{k}{struct}\PYG{+w}{ }\PYG{n+nc}{fdtable}\PYG{+w}{ }\PYG{o}{*}\PYG{n}{fdt}\PYG{p}{,}\PYG{+w}{ }\PYG{k+kt}{unsigned}\PYG{+w}{ }\PYG{k+kt}{int}\PYG{+w}{ }\PYG{n}{start}\PYG{p}{)}
\PYG{p}{\PYGZob{}}
\PYG{+w}{    }\PYG{k+kt}{unsigned}\PYG{+w}{ }\PYG{k+kt}{int}\PYG{+w}{ }\PYG{n}{maxfd}\PYG{+w}{ }\PYG{o}{=}\PYG{+w}{ }\PYG{n}{fdt}\PYG{o}{\PYGZhy{}}\PYG{o}{\PYGZgt{}}\PYG{n}{max\PYGZus{}fds}\PYG{p}{;}
\PYG{+w}{    }\PYG{k+kt}{unsigned}\PYG{+w}{ }\PYG{k+kt}{int}\PYG{+w}{ }\PYG{n}{maxbit}\PYG{+w}{ }\PYG{o}{=}\PYG{+w}{ }\PYG{n}{maxfd}\PYG{+w}{ }\PYG{o}{/}\PYG{+w}{ }\PYG{n}{BITS\PYGZus{}PER\PYGZus{}LONG}\PYG{p}{;}
\PYG{+w}{    }\PYG{k+kt}{unsigned}\PYG{+w}{ }\PYG{k+kt}{int}\PYG{+w}{ }\PYG{n}{bitbit}\PYG{+w}{ }\PYG{o}{=}\PYG{+w}{ }\PYG{n}{start}\PYG{+w}{ }\PYG{o}{/}\PYG{+w}{ }\PYG{n}{BITS\PYGZus{}PER\PYGZus{}LONG}\PYG{p}{;}
\PYG{+w}{    }\PYG{k+kt}{unsigned}\PYG{+w}{ }\PYG{k+kt}{int}\PYG{+w}{ }\PYG{n}{bit}\PYG{p}{;}

\PYG{+w}{    }\PYG{c+cm}{/*}
\PYG{c+cm}{     * Try to avoid looking at the second level bitmap}
\PYG{c+cm}{     */}
\PYG{+w}{    }\PYG{n}{bit}\PYG{+w}{ }\PYG{o}{=}\PYG{+w}{ }\PYG{n}{find\PYGZus{}next\PYGZus{}zero\PYGZus{}bit}\PYG{p}{(}\PYG{o}{\PYGZam{}}\PYG{n}{fdt}\PYG{o}{\PYGZhy{}}\PYG{o}{\PYGZgt{}}\PYG{n}{open\PYGZus{}fds}\PYG{p}{[}\PYG{n}{bitbit}\PYG{p}{]}\PYG{p}{,}\PYG{+w}{ }\PYG{n}{BITS\PYGZus{}PER\PYGZus{}LONG}\PYG{p}{,}
\PYG{+w}{                }\PYG{n}{start}\PYG{+w}{ }\PYG{o}{\PYGZam{}}\PYG{+w}{ }\PYG{p}{(}\PYG{n}{BITS\PYGZus{}PER\PYGZus{}LONG}\PYG{+w}{ }\PYG{o}{\PYGZhy{}}\PYG{+w}{ }\PYG{l+m+mi}{1}\PYG{p}{)}\PYG{p}{)}\PYG{p}{;}
\PYG{+w}{    }\PYG{k}{if}\PYG{+w}{ }\PYG{p}{(}\PYG{n}{bit}\PYG{+w}{ }\PYG{o}{\PYGZlt{}}\PYG{+w}{ }\PYG{n}{BITS\PYGZus{}PER\PYGZus{}LONG}\PYG{p}{)}
\PYG{+w}{        }\PYG{k}{return}\PYG{+w}{ }\PYG{n}{bit}\PYG{+w}{ }\PYG{o}{+}\PYG{+w}{ }\PYG{n}{bitbit}\PYG{+w}{ }\PYG{o}{*}\PYG{+w}{ }\PYG{n}{BITS\PYGZus{}PER\PYGZus{}LONG}\PYG{p}{;}

\PYG{+w}{    }\PYG{n}{bitbit}\PYG{+w}{ }\PYG{o}{=}\PYG{+w}{ }\PYG{n}{find\PYGZus{}next\PYGZus{}zero\PYGZus{}bit}\PYG{p}{(}\PYG{n}{fdt}\PYG{o}{\PYGZhy{}}\PYG{o}{\PYGZgt{}}\PYG{n}{full\PYGZus{}fds\PYGZus{}bits}\PYG{p}{,}\PYG{+w}{ }\PYG{n}{maxbit}\PYG{p}{,}\PYG{+w}{ }\PYG{n}{bitbit}\PYG{p}{)}\PYG{+w}{ }\PYG{o}{*}\PYG{+w}{ }\PYG{n}{BITS\PYGZus{}PER\PYGZus{}LONG}\PYG{p}{;}
\PYG{+w}{    }\PYG{k}{if}\PYG{+w}{ }\PYG{p}{(}\PYG{n}{bitbit}\PYG{+w}{ }\PYG{o}{\PYGZgt{}}\PYG{o}{=}\PYG{+w}{ }\PYG{n}{maxfd}\PYG{p}{)}
\PYG{+w}{        }\PYG{k}{return}\PYG{+w}{ }\PYG{n}{maxfd}\PYG{p}{;}
\PYG{+w}{    }\PYG{k}{if}\PYG{+w}{ }\PYG{p}{(}\PYG{n}{bitbit}\PYG{+w}{ }\PYG{o}{\PYGZgt{}}\PYG{+w}{ }\PYG{n}{start}\PYG{p}{)}
\PYG{+w}{        }\PYG{n}{start}\PYG{+w}{ }\PYG{o}{=}\PYG{+w}{ }\PYG{n}{bitbit}\PYG{p}{;}
\PYG{+w}{    }\PYG{k}{return}\PYG{+w}{ }\PYG{n}{find\PYGZus{}next\PYGZus{}zero\PYGZus{}bit}\PYG{p}{(}\PYG{n}{fdt}\PYG{o}{\PYGZhy{}}\PYG{o}{\PYGZgt{}}\PYG{n}{open\PYGZus{}fds}\PYG{p}{,}\PYG{+w}{ }\PYG{n}{maxfd}\PYG{p}{,}\PYG{+w}{ }\PYG{n}{start}\PYG{p}{)}\PYG{p}{;}
\PYG{p}{\PYGZcb{}}
\end{sphinxVerbatim}

\sphinxAtStartPar
Makine dili düzeyinde aramanın \sphinxcode{\sphinxupquote{ffz}} fonksiyonunda ve \sphinxcode{\sphinxupquote{\_\_ffs}} fonksiyonunda yapıldığına dikkat
ediniz.

\sphinxAtStartPar
\sphinxcode{\sphinxupquote{files\_struct}} yapısının \sphinxcode{\sphinxupquote{next\_fd}} elemanı aramanın başlatılacağı yeri belirtmektedir. Yani
\sphinxcode{\sphinxupquote{next\_fd}}’nin belirttiği değerden küçük tüm dosya betimleyicileri doludur. Dolayısıyla arama
\sphinxcode{\sphinxupquote{full\_fds\_bits}} dizisinin başından başlatılmamaktadır.


\section{2.6 Çekirdeğinde files\_struct ve fdtable}
\label{\detokenize{filesystem1:cekirdeginde-files-struct-ve-fdtable}}
\sphinxAtStartPar
Çekirdeğin 2.6 versiyonlarına gelindiğinde \sphinxcode{\sphinxupquote{files\_struct}} yapısının içerisi \sphinxcode{\sphinxupquote{fdtable}} isimli
bir yapı ile biraz daha derli toplu fakat biraz daha karmaşık hale getirilmiştir. 2.6’lı
versiyonlardaki \sphinxcode{\sphinxupquote{files\_struct}} yapısı şöyledir:

\begin{sphinxVerbatim}[commandchars=\\\{\}]
\PYG{k}{struct}\PYG{+w}{ }\PYG{n+nc}{files\PYGZus{}struct}\PYG{+w}{ }\PYG{p}{\PYGZob{}}
\PYG{c+cm}{/*}
\PYG{c+cm}{ * read mostly part}
\PYG{c+cm}{ */}
\PYG{+w}{    }\PYG{n}{atomic\PYGZus{}t}\PYG{+w}{ }\PYG{n}{count}\PYG{p}{;}
\PYG{+w}{    }\PYG{k}{struct}\PYG{+w}{ }\PYG{n+nc}{fdtable}\PYG{+w}{ }\PYG{n}{\PYGZus{}\PYGZus{}rcu}\PYG{+w}{ }\PYG{o}{*}\PYG{n}{fdt}\PYG{p}{;}
\PYG{+w}{    }\PYG{k}{struct}\PYG{+w}{ }\PYG{n+nc}{fdtable}\PYG{+w}{ }\PYG{n}{fdtab}\PYG{p}{;}
\PYG{c+cm}{/*}
\PYG{c+cm}{ * written part on a separate cache line in SMP}
\PYG{c+cm}{ */}
\PYG{+w}{    }\PYG{n}{spinlock\PYGZus{}t}\PYG{+w}{ }\PYG{n}{file\PYGZus{}lock}\PYG{+w}{ }\PYG{n}{\PYGZus{}\PYGZus{}\PYGZus{}\PYGZus{}cacheline\PYGZus{}aligned\PYGZus{}in\PYGZus{}smp}\PYG{p}{;}
\PYG{+w}{    }\PYG{k+kt}{int}\PYG{+w}{ }\PYG{n}{next\PYGZus{}fd}\PYG{p}{;}
\PYG{+w}{    }\PYG{k}{struct}\PYG{+w}{ }\PYG{n+nc}{embedded\PYGZus{}fd\PYGZus{}set}\PYG{+w}{ }\PYG{n}{close\PYGZus{}on\PYGZus{}exec\PYGZus{}init}\PYG{p}{;}
\PYG{+w}{    }\PYG{k}{struct}\PYG{+w}{ }\PYG{n+nc}{embedded\PYGZus{}fd\PYGZus{}set}\PYG{+w}{ }\PYG{n}{open\PYGZus{}fds\PYGZus{}init}\PYG{p}{;}
\PYG{+w}{    }\PYG{k}{struct}\PYG{+w}{ }\PYG{n+nc}{file}\PYG{+w}{ }\PYG{n}{\PYGZus{}\PYGZus{}rcu}\PYG{+w}{ }\PYG{o}{*}\PYG{+w}{ }\PYG{n}{fd\PYGZus{}array}\PYG{p}{[}\PYG{n}{NR\PYGZus{}OPEN\PYGZus{}DEFAULT}\PYG{p}{]}\PYG{p}{;}
\PYG{p}{\PYGZcb{}}\PYG{p}{;}
\end{sphinxVerbatim}

\sphinxAtStartPar
\sphinxcode{\sphinxupquote{fdtable}} yapısı da şöyledir:

\begin{sphinxVerbatim}[commandchars=\\\{\}]
\PYG{k}{struct}\PYG{+w}{ }\PYG{n+nc}{fdtable}\PYG{+w}{ }\PYG{p}{\PYGZob{}}
\PYG{+w}{    }\PYG{k+kt}{unsigned}\PYG{+w}{ }\PYG{k+kt}{int}\PYG{+w}{ }\PYG{n}{max\PYGZus{}fds}\PYG{p}{;}
\PYG{+w}{    }\PYG{k}{struct}\PYG{+w}{ }\PYG{n+nc}{file}\PYG{+w}{ }\PYG{n}{\PYGZus{}\PYGZus{}rcu}\PYG{+w}{ }\PYG{o}{*}\PYG{o}{*}\PYG{n}{fd}\PYG{p}{;}\PYG{+w}{      }\PYG{c+cm}{/* current fd array */}
\PYG{+w}{    }\PYG{n}{fd\PYGZus{}set}\PYG{+w}{ }\PYG{o}{*}\PYG{n}{close\PYGZus{}on\PYGZus{}exec}\PYG{p}{;}
\PYG{+w}{    }\PYG{n}{fd\PYGZus{}set}\PYG{+w}{ }\PYG{o}{*}\PYG{n}{open\PYGZus{}fds}\PYG{p}{;}
\PYG{+w}{    }\PYG{k}{struct}\PYG{+w}{ }\PYG{n+nc}{rcu\PYGZus{}head}\PYG{+w}{ }\PYG{n}{rcu}\PYG{p}{;}
\PYG{+w}{    }\PYG{k}{struct}\PYG{+w}{ }\PYG{n+nc}{fdtable}\PYG{+w}{ }\PYG{o}{*}\PYG{n}{next}\PYG{p}{;}
\PYG{p}{\PYGZcb{}}\PYG{p}{;}
\end{sphinxVerbatim}

\sphinxAtStartPar
Bu versiyonlarda çekirdek her zaman \sphinxcode{\sphinxupquote{fdt}} göstericisinin gösterdiği yerden işlemine
başlamaktadır. \sphinxcode{\sphinxupquote{fdt}} göstericisi işin başında yapı içerisindeki \sphinxcode{\sphinxupquote{fdtab}} yapı nesnesini
göstermektedir. \sphinxcode{\sphinxupquote{fdtab}} yapı nesnesinin içerisinde de \sphinxcode{\sphinxupquote{fd}}, \sphinxcode{\sphinxupquote{close\_on\_exec}}, \sphinxcode{\sphinxupquote{open\_fds}}
göstericileri vardır. Bu göstericiler de işin başında \sphinxcode{\sphinxupquote{files\_struct}} içerisindeki \sphinxcode{\sphinxupquote{fd\_array}},
\sphinxcode{\sphinxupquote{close\_on\_exec\_init}} ve \sphinxcode{\sphinxupquote{open\_fds\_init}} elemanlarını göstermektedir.

\sphinxAtStartPar
Bu versiyonlarda \sphinxcode{\sphinxupquote{current}} göstericisinden hareketle \sphinxcode{\sphinxupquote{fdx}} betimleyicisinin gösterdiği
yerdeki dosya nesnesine \sphinxcode{\sphinxupquote{current\sphinxhyphen{}\textgreater{}files\sphinxhyphen{}\textgreater{}fdt\sphinxhyphen{}\textgreater{}fd{[}fdx{]}}} ifadesiyle erişilebilir.

\sphinxAtStartPar
2.6’lı çekirdeklerde de dosya betimleyicisinden hareketle dosya nesnesine erişimi sağlayan
\sphinxcode{\sphinxupquote{fget}} ve \sphinxcode{\sphinxupquote{fput}} fonksiyonları bulunmaktadır:

\begin{sphinxVerbatim}[commandchars=\\\{\}]
\PYG{k}{struct}\PYG{+w}{ }\PYG{n+nc}{file}\PYG{+w}{ }\PYG{o}{*}\PYG{n}{fget}\PYG{p}{(}\PYG{k+kt}{unsigned}\PYG{+w}{ }\PYG{k+kt}{int}\PYG{+w}{ }\PYG{n}{fd}\PYG{p}{)}
\PYG{p}{\PYGZob{}}
\PYG{+w}{    }\PYG{k}{return}\PYG{+w}{ }\PYG{n}{\PYGZus{}\PYGZus{}fget}\PYG{p}{(}\PYG{n}{fd}\PYG{p}{,}\PYG{+w}{ }\PYG{n}{FMODE\PYGZus{}PATH}\PYG{p}{)}\PYG{p}{;}
\PYG{p}{\PYGZcb{}}
\PYG{n}{EXPORT\PYGZus{}SYMBOL}\PYG{p}{(}\PYG{n}{fget}\PYG{p}{)}\PYG{p}{;}
\end{sphinxVerbatim}

\begin{sphinxVerbatim}[commandchars=\\\{\}]
\PYG{k+kt}{void}\PYG{+w}{ }\PYG{n+nf}{fput}\PYG{p}{(}\PYG{k}{struct}\PYG{+w}{ }\PYG{n+nc}{file}\PYG{+w}{ }\PYG{o}{*}\PYG{n}{file}\PYG{p}{)}
\PYG{p}{\PYGZob{}}
\PYG{+w}{    }\PYG{k}{if}\PYG{+w}{ }\PYG{p}{(}\PYG{n}{atomic\PYGZus{}long\PYGZus{}dec\PYGZus{}and\PYGZus{}test}\PYG{p}{(}\PYG{o}{\PYGZam{}}\PYG{n}{file}\PYG{o}{\PYGZhy{}}\PYG{o}{\PYGZgt{}}\PYG{n}{f\PYGZus{}count}\PYG{p}{)}\PYG{p}{)}
\PYG{+w}{        }\PYG{n}{\PYGZus{}\PYGZus{}fput}\PYG{p}{(}\PYG{n}{file}\PYG{p}{)}\PYG{p}{;}
\PYG{p}{\PYGZcb{}}
\end{sphinxVerbatim}

\sphinxAtStartPar
\sphinxcode{\sphinxupquote{fget}} fonksiyonuyla elde edilen dosya nesnesi \sphinxcode{\sphinxupquote{fput}} fonksiyonuyla geri bırakılmaktadır.


\section{Güncel Çekirdeklerdeki Dosya Veri Yapıları}
\label{\detokenize{filesystem1:guncel-cekirdeklerdeki-dosya-veri-yapilari}}
\sphinxAtStartPar
Belli bir zamandan sonra artık bit dizisi oluşturmak için \sphinxcode{\sphinxupquote{fd\_set}} yapısının kullanılmasından
vazgeçilmiştir. Güncel çekirdeklerdeki \sphinxcode{\sphinxupquote{files\_struct}} yapısı şöyledir:

\begin{sphinxVerbatim}[commandchars=\\\{\}]
\PYG{k}{struct}\PYG{+w}{ }\PYG{n+nc}{files\PYGZus{}struct}\PYG{+w}{ }\PYG{p}{\PYGZob{}}
\PYG{c+cm}{/*}
\PYG{c+cm}{ * read mostly part}
\PYG{c+cm}{ */}
\PYG{+w}{    }\PYG{n}{atomic\PYGZus{}t}\PYG{+w}{ }\PYG{n}{count}\PYG{p}{;}
\PYG{+w}{    }\PYG{k+kt}{bool}\PYG{+w}{ }\PYG{n}{resize\PYGZus{}in\PYGZus{}progress}\PYG{p}{;}
\PYG{+w}{    }\PYG{n}{wait\PYGZus{}queue\PYGZus{}head\PYGZus{}t}\PYG{+w}{ }\PYG{n}{resize\PYGZus{}wait}\PYG{p}{;}

\PYG{+w}{    }\PYG{k}{struct}\PYG{+w}{ }\PYG{n+nc}{fdtable}\PYG{+w}{ }\PYG{n}{\PYGZus{}\PYGZus{}rcu}\PYG{+w}{ }\PYG{o}{*}\PYG{n}{fdt}\PYG{p}{;}
\PYG{+w}{    }\PYG{k}{struct}\PYG{+w}{ }\PYG{n+nc}{fdtable}\PYG{+w}{ }\PYG{n}{fdtab}\PYG{p}{;}
\PYG{c+cm}{/*}
\PYG{c+cm}{ * written part on a separate cache line in SMP}
\PYG{c+cm}{ */}
\PYG{+w}{    }\PYG{n}{spinlock\PYGZus{}t}\PYG{+w}{ }\PYG{n}{file\PYGZus{}lock}\PYG{+w}{ }\PYG{n}{\PYGZus{}\PYGZus{}\PYGZus{}\PYGZus{}cacheline\PYGZus{}aligned\PYGZus{}in\PYGZus{}smp}\PYG{p}{;}
\PYG{+w}{    }\PYG{k+kt}{unsigned}\PYG{+w}{ }\PYG{k+kt}{int}\PYG{+w}{ }\PYG{n}{next\PYGZus{}fd}\PYG{p}{;}
\PYG{+w}{    }\PYG{k+kt}{unsigned}\PYG{+w}{ }\PYG{k+kt}{long}\PYG{+w}{ }\PYG{n}{close\PYGZus{}on\PYGZus{}exec\PYGZus{}init}\PYG{p}{[}\PYG{l+m+mi}{1}\PYG{p}{]}\PYG{p}{;}
\PYG{+w}{    }\PYG{k+kt}{unsigned}\PYG{+w}{ }\PYG{k+kt}{long}\PYG{+w}{ }\PYG{n}{open\PYGZus{}fds\PYGZus{}init}\PYG{p}{[}\PYG{l+m+mi}{1}\PYG{p}{]}\PYG{p}{;}
\PYG{+w}{    }\PYG{k+kt}{unsigned}\PYG{+w}{ }\PYG{k+kt}{long}\PYG{+w}{ }\PYG{n}{full\PYGZus{}fds\PYGZus{}bits\PYGZus{}init}\PYG{p}{[}\PYG{l+m+mi}{1}\PYG{p}{]}\PYG{p}{;}
\PYG{+w}{    }\PYG{k}{struct}\PYG{+w}{ }\PYG{n+nc}{file}\PYG{+w}{ }\PYG{n}{\PYGZus{}\PYGZus{}rcu}\PYG{+w}{ }\PYG{o}{*}\PYG{+w}{ }\PYG{n}{fd\PYGZus{}array}\PYG{p}{[}\PYG{n}{NR\PYGZus{}OPEN\PYGZus{}DEFAULT}\PYG{p}{]}\PYG{p}{;}
\PYG{p}{\PYGZcb{}}\PYG{p}{;}
\end{sphinxVerbatim}

\sphinxAtStartPar
\sphinxcode{\sphinxupquote{fdtable}} yapısı da şöyledir:

\begin{sphinxVerbatim}[commandchars=\\\{\}]
\PYG{k}{struct}\PYG{+w}{ }\PYG{n+nc}{fdtable}\PYG{+w}{ }\PYG{p}{\PYGZob{}}
\PYG{+w}{    }\PYG{k+kt}{unsigned}\PYG{+w}{ }\PYG{k+kt}{int}\PYG{+w}{ }\PYG{n}{max\PYGZus{}fds}\PYG{p}{;}
\PYG{+w}{    }\PYG{k}{struct}\PYG{+w}{ }\PYG{n+nc}{file}\PYG{+w}{ }\PYG{n}{\PYGZus{}\PYGZus{}rcu}\PYG{+w}{ }\PYG{o}{*}\PYG{o}{*}\PYG{n}{fd}\PYG{p}{;}\PYG{+w}{      }\PYG{c+cm}{/* current fd array */}
\PYG{+w}{    }\PYG{k+kt}{unsigned}\PYG{+w}{ }\PYG{k+kt}{long}\PYG{+w}{ }\PYG{o}{*}\PYG{n}{close\PYGZus{}on\PYGZus{}exec}\PYG{p}{;}
\PYG{+w}{    }\PYG{k+kt}{unsigned}\PYG{+w}{ }\PYG{k+kt}{long}\PYG{+w}{ }\PYG{o}{*}\PYG{n}{open\PYGZus{}fds}\PYG{p}{;}
\PYG{+w}{    }\PYG{k+kt}{unsigned}\PYG{+w}{ }\PYG{k+kt}{long}\PYG{+w}{ }\PYG{o}{*}\PYG{n}{full\PYGZus{}fds\PYGZus{}bits}\PYG{p}{;}
\PYG{+w}{    }\PYG{k}{struct}\PYG{+w}{ }\PYG{n+nc}{rcu\PYGZus{}head}\PYG{+w}{ }\PYG{n}{rcu}\PYG{p}{;}
\PYG{p}{\PYGZcb{}}\PYG{p}{;}
\end{sphinxVerbatim}

\sphinxAtStartPar
Görüldüğü gibi artık bit dizileri \sphinxcode{\sphinxupquote{fd\_set}} yerine doğrudan \sphinxcode{\sphinxupquote{unsigned long}} türden bir dizi
biçiminde oluşturulmaktadır.

\sphinxAtStartPar
Güncel versiyonlarda \sphinxcode{\sphinxupquote{fcheck}} biçiminde bir makro ve \sphinxcode{\sphinxupquote{fcheck\_files}} isimli bir fonksiyon
yoktur. Ancak yine güncel versiyonlarda dosya betimleyicisi yoluyla dosya nesnesine erişimi
referans sayacını artırarak yapan \sphinxcode{\sphinxupquote{fget}} fonksiyonu ve sayacı eksilterek nesneyi bırakan
\sphinxcode{\sphinxupquote{fput}} fonksiyonu bulunmaktadır:

\begin{sphinxVerbatim}[commandchars=\\\{\}]
\PYG{k}{struct}\PYG{+w}{ }\PYG{n+nc}{file}\PYG{+w}{ }\PYG{o}{*}\PYG{n}{fget}\PYG{p}{(}\PYG{k+kt}{unsigned}\PYG{+w}{ }\PYG{k+kt}{int}\PYG{+w}{ }\PYG{n}{fd}\PYG{p}{)}
\PYG{p}{\PYGZob{}}
\PYG{+w}{    }\PYG{k}{return}\PYG{+w}{ }\PYG{n}{\PYGZus{}\PYGZus{}fget}\PYG{p}{(}\PYG{n}{fd}\PYG{p}{,}\PYG{+w}{ }\PYG{n}{FMODE\PYGZus{}PATH}\PYG{p}{)}\PYG{p}{;}
\PYG{p}{\PYGZcb{}}
\PYG{n}{EXPORT\PYGZus{}SYMBOL}\PYG{p}{(}\PYG{n}{fget}\PYG{p}{)}\PYG{p}{;}
\end{sphinxVerbatim}

\begin{sphinxVerbatim}[commandchars=\\\{\}]
\PYG{k+kt}{void}\PYG{+w}{ }\PYG{n+nf}{fput}\PYG{p}{(}\PYG{k}{struct}\PYG{+w}{ }\PYG{n+nc}{file}\PYG{+w}{ }\PYG{o}{*}\PYG{n}{file}\PYG{p}{)}
\PYG{p}{\PYGZob{}}
\PYG{+w}{    }\PYG{k}{if}\PYG{+w}{ }\PYG{p}{(}\PYG{n}{unlikely}\PYG{p}{(}\PYG{n}{file\PYGZus{}ref\PYGZus{}put}\PYG{p}{(}\PYG{o}{\PYGZam{}}\PYG{n}{file}\PYG{o}{\PYGZhy{}}\PYG{o}{\PYGZgt{}}\PYG{n}{f\PYGZus{}ref}\PYG{p}{)}\PYG{p}{)}\PYG{p}{)}
\PYG{+w}{        }\PYG{n}{\PYGZus{}\PYGZus{}fput\PYGZus{}deferred}\PYG{p}{(}\PYG{n}{file}\PYG{p}{)}\PYG{p}{;}
\PYG{p}{\PYGZcb{}}
\PYG{n}{EXPORT\PYGZus{}SYMBOL}\PYG{p}{(}\PYG{n}{fput}\PYG{p}{)}\PYG{p}{;}
\end{sphinxVerbatim}


\section{Thread’ler, Fork ve Dosya Betimleyicileri}
\label{\detokenize{filesystem1:thread-ler-fork-ve-dosya-betimleyicileri}}
\sphinxAtStartPar
POSIX sistemlerinde dolayısıyla da Linux çekirdeğinde thread’lerin ayrı dosya betimleyici
tabloları yoktur. Dosya betimleyicileri ve dosya betimleyici tablosu prosese özgüdür. Bir thread
yaratıldığında thread’e ilişkin \sphinxcode{\sphinxupquote{task\_struct}} nesnesinin \sphinxcode{\sphinxupquote{fs}} ve \sphinxcode{\sphinxupquote{files}} gibi elemanları,
onu yaratan thread’in \sphinxcode{\sphinxupquote{task\_struct}} nesnesinden sığ kopyalamayla kopyalanmaktadır. Dolayısıyla
aslında prosesin bütün thread’leri açık dosyalara ilişkin aynı veri yapısı nesnelerini
kullanmaktadır:

\begin{sphinxVerbatim}[commandchars=\\\{\}]
\PYG{k}{struct}\PYG{+w}{ }\PYG{n+nc}{task\PYGZus{}struct}\PYG{+w}{ }\PYG{p}{\PYGZob{}}
\PYG{+w}{    }\PYG{c+cm}{/* ... */}

\PYG{+w}{    }\PYG{c+cm}{/* Filesystem information: */}
\PYG{+w}{    }\PYG{k}{struct}\PYG{+w}{ }\PYG{n+nc}{fs\PYGZus{}struct}\PYG{+w}{        }\PYG{o}{*}\PYG{n}{fs}\PYG{p}{;}

\PYG{+w}{    }\PYG{c+cm}{/* Open file information: */}
\PYG{+w}{    }\PYG{k}{struct}\PYG{+w}{ }\PYG{n+nc}{files\PYGZus{}struct}\PYG{+w}{     }\PYG{o}{*}\PYG{n}{files}\PYG{p}{;}

\PYG{+w}{    }\PYG{c+cm}{/* ... */}
\PYG{p}{\PYGZcb{}}\PYG{p}{;}
\end{sphinxVerbatim}

\sphinxAtStartPar
Yeni \sphinxcode{\sphinxupquote{task\_struct}} nesnesi yaratılıp diğer \sphinxcode{\sphinxupquote{task\_struct}} nesnesinden sığ kopyalama
yapıldığında \sphinxcode{\sphinxupquote{files}} göstericisinin de aslında aynı \sphinxcode{\sphinxupquote{files\_struct}} nesnesini göstereceğine
dikkat ediniz. Yani toplamda aslında proses için tek bir \sphinxcode{\sphinxupquote{files\_struct}} nesnesi bulunmaktadır.

\sphinxAtStartPar
\sphinxcode{\sphinxupquote{fork}} işlemi sırasında ise tamamen alt proses için yeni bir \sphinxcode{\sphinxupquote{files\_struct}} nesnesi ve yeni
bir dosya betimleyici tablosu (fd dizisi) yaratılmaktadır. Ancak üst prosesin dosya betimleyici
tablosundaki adresler yeni yaratılan alt prosesteki dosya betimleyici tablosuna kopyalanmaktadır.
Böylece üst prosesin dosya betimleyici tablosunun aynı numaralı betimleyicileriyle alt prosesin
dosya betimleyici tablosunun aynı numaralı betimleyicileri aynı dosya nesnelerini gösteriyor
durumda olur. Tabii artık üst ve alt proseslerin yeni açacağı dosyalar onlara özgü olacaktır.
Paylaşılan dosya nesneleri yalnızca \sphinxcode{\sphinxupquote{fork}} öncesinde açılmış olanlardır.


\section{Dosya Sistemi Temel Nesneleri}
\label{\detokenize{filesystem1:dosya-sistemi-temel-nesneleri}}
\sphinxAtStartPar
Şimdiye kadar açık dosyalara ilişkin çekirdeğin oluşturduğu veri yapıları hakkında şu bilgileri
edindik:
\begin{itemize}
\item {} 
\sphinxAtStartPar
Açık dosyalara ilişkin \sphinxcode{\sphinxupquote{task\_struct}} içerisindeki veri yapıları.

\item {} 
\sphinxAtStartPar
Dosya betimleyicilerinin anlamı ve dosya betimleyicisi yoluyla dosya nesnelerine nasıl
erişildiği.

\item {} 
\sphinxAtStartPar
En düşük boş betimleyicinin elde edilme biçimi.

\end{itemize}

\sphinxAtStartPar
Şimdi dosya sisteminin diğer önemli veri yapıları üzerinde duracağız.


\subsection{Dosya Nesnesi, inode Nesnesi ve dentry Nesnesi}
\label{\detokenize{filesystem1:dosya-nesnesi-inode-nesnesi-ve-dentry-nesnesi}}
\sphinxAtStartPar
Çekirdek açık bir dosya üzerinde read/write gibi işlemleri yaparken dosyanın son okunma tarihi,
son güncelleme tarihi, dosya uzunluğu gibi bilgileri de güncelleyeceğine göre bu bilgilere nasıl
erişmektedir? Dosyaların uzunluk, erişim hakları, tarih\sphinxhyphen{}zaman bilgisi gibi metadata bilgileri
diskte tutulmaktadır. (Ext dosya sistemlerinde bu bilgilerin diskte tutulduğu yere inode blok
denilmektedir.) Çekirdek dosya açılırken dosyanın bu metadata bilgilerini diskten okuyarak
bellekte \sphinxcode{\sphinxupquote{inode}} isimli bir yapı nesnesinin içerisine yerleştirmektedir. Biz dosyaların
metadata bilgilerinin tutulduğu bu \sphinxstyleemphasis{inode nesnesi} türünden nesnelere \sphinxstyleemphasis{inode nesnesi} de
diyeceğiz.

\sphinxAtStartPar
Peki bir dosya açıldığında dosyanın dosya sistemindeki yeri çekirdek tarafından nasıl
saklanmaktadır? Çekirdek her dosya açıldığında o dosyanın dizin girişi bilgilerini (yani dosya
sisteminde nerede olduğu bilgisini ve bazı diğer bilgileri) ismine \sphinxstyleemphasis{dentry} denilen bir yapı
nesnesine yerleştirmektedir.

\sphinxAtStartPar
Dosya sistemi nesnelerini şöyle özetleyebiliriz:


\begin{savenotes}\sphinxattablestart
\sphinxthistablewithglobalstyle
\centering
\begin{tabular}[t]{\X{30}{100}\X{70}{100}}
\sphinxtoprule
\begin{varwidth}[t]{\sphinxcolwidth{1}{2}}
\sphinxstyletheadfamily \sphinxAtStartPar
Nesne
\sphinxbeforeendvarwidth
\end{varwidth}%
&\begin{varwidth}[t]{\sphinxcolwidth{1}{2}}
\sphinxstyletheadfamily \sphinxAtStartPar
Açıklama
\sphinxbeforeendvarwidth
\end{varwidth}%
\\
\sphinxmidrule
\sphinxtableatstartofbodyhook\begin{varwidth}[t]{\sphinxcolwidth{1}{2}}
\sphinxAtStartPar
\sphinxstylestrong{Dosya Nesnesi} (\sphinxcode{\sphinxupquote{struct file}})
\sphinxbeforeendvarwidth
\end{varwidth}%
&\begin{varwidth}[t]{\sphinxcolwidth{1}{2}}
\sphinxAtStartPar
Açılmış dosyalar üzerinde işlem yapmak için gereken tüm bilgilerin tutulduğu nesne.
\sphinxbeforeendvarwidth
\end{varwidth}%
\\
\sphinxhline\begin{varwidth}[t]{\sphinxcolwidth{1}{2}}
\sphinxAtStartPar
\sphinxstylestrong{inode Nesnesi} (\sphinxcode{\sphinxupquote{struct inode}})
\sphinxbeforeendvarwidth
\end{varwidth}%
&\begin{varwidth}[t]{\sphinxcolwidth{1}{2}}
\sphinxAtStartPar
Dosyanın diskteki metadata bilgilerini tutan nesne.
\sphinxbeforeendvarwidth
\end{varwidth}%
\\
\sphinxhline\begin{varwidth}[t]{\sphinxcolwidth{1}{2}}
\sphinxAtStartPar
\sphinxstylestrong{dentry Nesnesi} (\sphinxcode{\sphinxupquote{struct dentry}})
\sphinxbeforeendvarwidth
\end{varwidth}%
&\begin{varwidth}[t]{\sphinxcolwidth{1}{2}}
\sphinxAtStartPar
Dosyanın dosya sistemi üzerindeki yerini ve buna ilişkin bazı bilgileri tutan nesne.
\sphinxbeforeendvarwidth
\end{varwidth}%
\\
\sphinxbottomrule
\end{tabular}
\sphinxtableafterendhook\par
\sphinxattableend\end{savenotes}


\subsection{Nesneler Arası İlişki}
\label{\detokenize{filesystem1:nesneler-arasi-iliski}}
\sphinxAtStartPar
Uzunca bir süre (2.6’ya kadar ve 2.6’lı versiyonlar da dahil olmak üzere) dosya nesnesinin
(\sphinxcode{\sphinxupquote{file}} yapısının) içerisinde dentry nesnesinin adresi, dentry nesnesinin içerisinde de o dizin
girişinin inode nesnesinin adresi tutuluyordu. Daha sonraları dosya nesnesinin içerisinde de
doğrudan inode nesnesinin adresi tutulmaya başlanmıştır:

\sphinxincludegraphics[]{graphviz-2aacf48e0dc5b01b34d9248e84725ceab0748778.pdf}

\sphinxAtStartPar
UNIX/Linux sistemlerinde bir dosya sistemi kök dizinde bir yere mount edilebilmektedir. Çekirdeğin
bazı durumlarda bir dosyanın hangi dosya sisteminin içerisinde bulunduğunu anlaması da
gerekebilmektedir. 2.6 çekirdeklerinden itibaren de dosyanın dentry nesnesinin adresiyle
vfsmount bilgileri \sphinxcode{\sphinxupquote{path}} isimli bir yapıya yerleştirilmiş ve \sphinxcode{\sphinxupquote{file}} yapısının içerisindeki
\sphinxcode{\sphinxupquote{f\_path}} elemanında tutulmaya başlanmıştır:

\begin{sphinxVerbatim}[commandchars=\\\{\}]
\PYG{k}{struct}\PYG{+w}{ }\PYG{n+nc}{file}\PYG{+w}{ }\PYG{p}{\PYGZob{}}
\PYG{+w}{    }\PYG{c+cm}{/* ... */}

\PYG{+w}{    }\PYG{n}{fmode\PYGZus{}t}\PYG{+w}{             }\PYG{n}{f\PYGZus{}mode}\PYG{p}{;}
\PYG{+w}{    }\PYG{k+kt}{unsigned}\PYG{+w}{ }\PYG{k+kt}{int}\PYG{+w}{        }\PYG{n}{f\PYGZus{}flags}\PYG{p}{;}
\PYG{+w}{    }\PYG{n}{loff\PYGZus{}t}\PYG{+w}{              }\PYG{n}{f\PYGZus{}pos}\PYG{p}{;}
\PYG{+w}{    }\PYG{n}{atomic\PYGZus{}long\PYGZus{}t}\PYG{+w}{       }\PYG{n}{f\PYGZus{}count}\PYG{p}{;}\PYG{+w}{    }\PYG{c+cm}{/* güncel çekirdeklerde: file\PYGZus{}ref\PYGZus{}t f\PYGZus{}ref */}
\PYG{+w}{    }\PYG{k}{struct}\PYG{+w}{ }\PYG{n+nc}{path}\PYG{+w}{         }\PYG{n}{f\PYGZus{}path}\PYG{p}{;}
\PYG{+w}{    }\PYG{k}{struct}\PYG{+w}{ }\PYG{n+nc}{inode}\PYG{+w}{       }\PYG{o}{*}\PYG{n}{f\PYGZus{}inode}\PYG{p}{;}

\PYG{+w}{    }\PYG{c+cm}{/* ... */}
\PYG{p}{\PYGZcb{}}\PYG{p}{;}

\PYG{k}{struct}\PYG{+w}{ }\PYG{n+nc}{path}\PYG{+w}{ }\PYG{p}{\PYGZob{}}
\PYG{+w}{    }\PYG{k}{struct}\PYG{+w}{ }\PYG{n+nc}{vfsmount}\PYG{+w}{ }\PYG{o}{*}\PYG{n}{mnt}\PYG{p}{;}
\PYG{+w}{    }\PYG{k}{struct}\PYG{+w}{ }\PYG{n+nc}{dentry}\PYG{+w}{ }\PYG{o}{*}\PYG{n}{dentry}\PYG{p}{;}
\PYG{p}{\PYGZcb{}}\PYG{+w}{ }\PYG{n}{\PYGZus{}\PYGZus{}randomize\PYGZus{}layout}\PYG{p}{;}
\end{sphinxVerbatim}


\subsection{Dosya Nesnesi Elemanları}
\label{\detokenize{filesystem1:dosya-nesnesi-elemanlari}}
\sphinxAtStartPar
Açık bir dosyanın \sphinxcode{\sphinxupquote{open}} POSIX fonksiyonunda kullanılan açış bayrakları (\sphinxcode{\sphinxupquote{O\_}} ile başlayan
bayraklar) \sphinxcode{\sphinxupquote{file}} yapısının \sphinxcode{\sphinxupquote{f\_flags}} elemanında saklanmaktadır. Örneğin:

\begin{sphinxVerbatim}[commandchars=\\\{\}]
\PYG{n}{fd}\PYG{+w}{ }\PYG{o}{=}\PYG{+w}{ }\PYG{n}{open}\PYG{p}{(}\PYG{l+s}{\PYGZdq{}}\PYG{l+s}{test.txt}\PYG{l+s}{\PYGZdq{}}\PYG{p}{,}\PYG{+w}{ }\PYG{n}{O\PYGZus{}RDWR}\PYG{o}{|}\PYG{n}{O\PYGZus{}APPEND}\PYG{p}{)}\PYG{p}{;}
\end{sphinxVerbatim}

\sphinxAtStartPar
Buradaki \sphinxcode{\sphinxupquote{O\_RDWR}} ve \sphinxcode{\sphinxupquote{O\_APPEND}} bayrakları \sphinxcode{\sphinxupquote{file}} yapısının \sphinxcode{\sphinxupquote{f\_flags}} elemanında
saklanmaktadır. Bu \sphinxcode{\sphinxupquote{f\_flags}} elemanının daha çabuk işleme sokulabilecek yeniden düzenlenmiş
hali yapının \sphinxcode{\sphinxupquote{f\_mode}} elemanında saklanmaktadır.

\sphinxAtStartPar
Okuma yazma işlemlerinin \sphinxstyleemphasis{dosya göstericisi (file pointer)} denilen bir offset’ten itibaren
yapıldığını anımsayınız. Dosya göstericisinin konumu da \sphinxcode{\sphinxupquote{file}} yapısının \sphinxcode{\sphinxupquote{f\_pos}} elemanında
tutulmaktadır.

\sphinxAtStartPar
Dosya nesnelerini birden fazla betimleyici gösterebilmektedir. Örneğin \sphinxcode{\sphinxupquote{dup}} ve \sphinxcode{\sphinxupquote{dup2}} POSIX
fonksiyonları aynı dosya nesnesini gösteren farklı bir betimleyicinin oluşturulmasına yol
açmaktadır. Bu durumda \sphinxcode{\sphinxupquote{close}} fonksiyonu ile dosya kapatıldığında dosya hemen silinmez. Çünkü
onu kullanan başka betimleyiciler de bulunuyor olabilir. İşte \sphinxcode{\sphinxupquote{file}} yapısının içerisindeki
referans sayacı (uzun süre \sphinxcode{\sphinxupquote{f\_count}}, güncel çekirdeklerde \sphinxcode{\sphinxupquote{f\_ref}} ismiyle) o dosya
nesnesinin kaç betimleyici tarafından gösterildiği bilgisini tutmaktadır. Her betimleyici
\sphinxcode{\sphinxupquote{close}} fonksiyonu ile kapatıldığında bu sayaç 1 eksiltilir; sayaç 0’a düştüğünde dosya nesnesi
silinir.


\section{inode ve dentry Önbellekleri}
\label{\detokenize{filesystem1:inode-ve-dentry-onbellekleri}}
\sphinxAtStartPar
Her açılan yeni dosya için çekirdek yeni bir dosya nesnesi yaratmaktadır. Ancak aynı dosyalar
açıldığında bu dosyalar için tek bir dentry ve inode nesnesi oluşturmaktadır. Bunun için \sphinxcode{\sphinxupquote{open}}
fonksiyonuyla bir dosya açıldığında çekirdeğin bu dosyaya ilişkin dentry nesnesi ve inode nesnesi
daha önce yaratılmış mı diye bir araştırma yapması gerekmektedir.

\sphinxAtStartPar
İşte çekirdek dentry ve inode nesneleri için ayrı önbellek (cache) sistemleri oluşturmaktadır.
Bunlara Linux sistemlerinde \sphinxstyleemphasis{dentry önbelleği (dentry cache)} ve \sphinxstyleemphasis{inode önbelleği (inode cache)}
denilmektedir. Bir dosya açıldığında çekirdek o dosyaya ilişkin dentry ve inode nesneleri zaten
bu önbellek sistemlerinde varsa hiç diske gitmeden doğrudan bu önbelleklerden onları alıp
kullanmaktadır.

\sphinxAtStartPar
Bir dosyayı onu açmış olan bütün prosesler kapatmış olsa bile o dosyaya ilişkin dentry ve inode
nesneleri bu önbellek sistemlerinde kalmaya devam edebilir. Çünkü dosyalar (özellikle bazı
merkezi dosyalar) farklı prosesler tarafından defalarca açılıp kullanılabilmektedir. (Örneğin
\sphinxcode{\sphinxupquote{/etc/passwd}} dosyası pek çok proses tarafından dolaylı bir biçimde açılıp kullanılabilmektedir.)
Bu tür durumlarda onların bu önbellekler içerisinde biriktirilmesi sistem performansını oldukça
iyileştirmektedir.

\sphinxAtStartPar
Linux’taki bu tür önbellek sistemlerinde önbellekten çıkarma için genel olarak \sphinxstyleemphasis{LRU (Least
Recently Used)} denilen \sphinxstyleemphasis{önbellek yer değiştirme (cache replacement) algoritması} işletilmektedir.
Yani önbellekten toplamda en az kullanılanlar değil son zamanlarda en az kullanılanlar
çıkartılmaktadır.

\sphinxAtStartPar
Bu noktada \sphinxstyleemphasis{dentry önbelleği} ve \sphinxstyleemphasis{inode önbelleği} ile \sphinxstyleemphasis{sayfa önbelleği (page cache)} arasındaki
ilişkiye de değinelim. Sayfa önbelleği disk blokları temelinde organize edilen ve her türlü disk
bloğunun transfer edilmesinde kullanılan aşağı seviyeli bir önbellek sistemidir. Halbuki dentry
ve inode önbellekleri yalnızca dentry ve inode elemanlarından oluşan önbelleklerdir. Bir dosya
açıldığında eğer o dosyaya ilişkin dentry ve inode nesneleri dentry ve inode önbelleklerinde
yoksa diskten elde edilmektedir. Tabii bu amaçla disk okuması yapılmadan önce bu disk bloklarının
sayfa önbelleğinde olup olmadığına da bakılmaktadır.


\section{open() Fonksiyonunun İşleyişi}
\label{\detokenize{filesystem1:open-fonksiyonunun-isleyisi}}
\sphinxAtStartPar
Geldiğimiz noktaya kadarki konuları dikkate aldığımızda \sphinxcode{\sphinxupquote{open}} fonksiyonuyla bir dosyanın
açılması durumunda sırasıyla şunların yapıldığını söyleyebiliriz:
\begin{enumerate}
\sphinxsetlistlabels{\arabic}{enumi}{enumii}{}{.}%
\item {} 
\sphinxAtStartPar
Prosesin dosya betimleyici tablosunda boş bir betimleyici hızlı bir biçimde bulunur.

\item {} 
\sphinxAtStartPar
Bir dosya nesnesi (\sphinxcode{\sphinxupquote{struct file}} nesnesi) yaratılır ve elemanlarına gerekli ilkdeğerler
verilir.

\item {} 
\sphinxAtStartPar
Dosyaya ilişkin dentry nesnesi ve inode nesnesi dentry önbelleğinde ve inode önbelleğinde
yoksa diskte arama yapılarak yaratılır ve bunların adresleri dosya nesnesine yerleştirilir.

\item {} 
\sphinxAtStartPar
Dosya nesnesinin adresi dosya betimleyici tablosuna (fd dizisine) yerleştirilir ve yerleştirilen
dizi indeksi dosya betimleyicisi olarak geri döndürülür.

\end{enumerate}


\section{Prosesin Kök ve Çalışma Dizini}
\label{\detokenize{filesystem1:prosesin-kok-ve-calisma-dizini}}
\sphinxAtStartPar
Bir prosesin kök dizininin ve çalışma dizininin proses kontrol bloğunda yani \sphinxcode{\sphinxupquote{task\_struct}}
nesnesinde tutulduğunu belirtmiştik. Güncel çekirdeklerde bu bilgi şöyle tutulmaktadır:

\begin{sphinxVerbatim}[commandchars=\\\{\}]
\PYG{k}{struct}\PYG{+w}{ }\PYG{n+nc}{task\PYGZus{}struct}\PYG{+w}{ }\PYG{p}{\PYGZob{}}
\PYG{+w}{    }\PYG{c+cm}{/* ... */}

\PYG{+w}{    }\PYG{k}{struct}\PYG{+w}{ }\PYG{n+nc}{fs\PYGZus{}struct}\PYG{+w}{ }\PYG{o}{*}\PYG{n}{fs}\PYG{p}{;}

\PYG{+w}{    }\PYG{c+cm}{/* ... */}
\PYG{p}{\PYGZcb{}}\PYG{p}{;}

\PYG{k}{struct}\PYG{+w}{ }\PYG{n+nc}{fs\PYGZus{}struct}\PYG{+w}{ }\PYG{p}{\PYGZob{}}
\PYG{+w}{    }\PYG{c+cm}{/* ... */}
\PYG{+w}{    }\PYG{k}{struct}\PYG{+w}{ }\PYG{n+nc}{path}\PYG{+w}{ }\PYG{n}{root}\PYG{p}{,}\PYG{+w}{ }\PYG{n}{pwd}\PYG{p}{;}

\PYG{+w}{    }\PYG{c+cm}{/* ... */}
\PYG{p}{\PYGZcb{}}\PYG{+w}{ }\PYG{n}{\PYGZus{}\PYGZus{}randomize\PYGZus{}layout}\PYG{p}{;}

\PYG{k}{struct}\PYG{+w}{ }\PYG{n+nc}{path}\PYG{+w}{ }\PYG{p}{\PYGZob{}}
\PYG{+w}{    }\PYG{k}{struct}\PYG{+w}{ }\PYG{n+nc}{vfsmount}\PYG{+w}{ }\PYG{o}{*}\PYG{n}{mnt}\PYG{p}{;}
\PYG{+w}{    }\PYG{k}{struct}\PYG{+w}{ }\PYG{n+nc}{dentry}\PYG{+w}{ }\PYG{o}{*}\PYG{n}{dentry}\PYG{p}{;}
\PYG{p}{\PYGZcb{}}\PYG{+w}{ }\PYG{n}{\PYGZus{}\PYGZus{}randomize\PYGZus{}layout}\PYG{p}{;}
\end{sphinxVerbatim}

\sphinxAtStartPar
Görüldüğü gibi çekirdek prosesin kök dizinin ve çalışma dizininin bilgisini yol ifadesi biçiminde
değil dentry nesnesi biçiminde tutmaktadır.


\section{Dosya Sistemi Veri Yapılarına İlişkin Testler}
\label{\detokenize{filesystem1:dosya-sistemi-veri-yapilarina-iliskin-testler}}
\sphinxAtStartPar
Şimdi dosya sisteminin temel veri yapıları üzerinde görmüş olduğumuz konulara ilişkin bazı
testler yapalım. Bu testleri yapmak için bir aygıt sürücü yazarak ioctl yoluyla gerçekleştireceğiz.
(Kursumuzda kullandığımız Linux makinedeki çekirdek sürümü 6.9.2’dir.)


\subsection{ioctl\_test1 ve ioctl\_test2: Dosya Göstericisini Okuma}
\label{\detokenize{filesystem1:ioctl-test1-ve-ioctl-test2-dosya-gostericisini-okuma}}
\sphinxAtStartPar
İlk örneğimizde bir dosya betimleyicisine ilişkin dosya nesnesini elde ederek onun \sphinxcode{\sphinxupquote{f\_pos}}
elemanını yazdırmaya çalışalım. Dosya nesnesindeki \sphinxcode{\sphinxupquote{f\_pos}} elemanının dosya göstericisinin
değerini tuttuğunu anımsayınız.

\sphinxAtStartPar
\sphinxstylestrong{Düşük seviyeli erişim (ioctl\_test1):}

\begin{sphinxVerbatim}[commandchars=\\\{\}]
\PYG{k}{static}\PYG{+w}{ }\PYG{k+kt}{long}\PYG{+w}{ }\PYG{n+nf}{ioctl\PYGZus{}test1}\PYG{p}{(}\PYG{k}{struct}\PYG{+w}{ }\PYG{n+nc}{file}\PYG{+w}{ }\PYG{o}{*}\PYG{n}{filp}\PYG{p}{,}\PYG{+w}{ }\PYG{k+kt}{unsigned}\PYG{+w}{ }\PYG{k+kt}{long}\PYG{+w}{ }\PYG{n}{arg}\PYG{p}{)}
\PYG{p}{\PYGZob{}}
\PYG{+w}{    }\PYG{k}{struct}\PYG{+w}{ }\PYG{n+nc}{file}\PYG{+w}{ }\PYG{o}{*}\PYG{n}{f}\PYG{p}{;}
\PYG{+w}{    }\PYG{k+kt}{int}\PYG{+w}{ }\PYG{n}{fd}\PYG{+w}{ }\PYG{o}{=}\PYG{+w}{ }\PYG{p}{(}\PYG{k+kt}{int}\PYG{p}{)}\PYG{+w}{ }\PYG{n}{arg}\PYG{p}{;}

\PYG{+w}{    }\PYG{k}{if}\PYG{+w}{ }\PYG{p}{(}\PYG{n}{fd}\PYG{+w}{ }\PYG{o}{\PYGZlt{}}\PYG{+w}{ }\PYG{l+m+mi}{0}\PYG{+w}{ }\PYG{o}{|}\PYG{o}{|}\PYG{+w}{ }\PYG{n}{fd}\PYG{+w}{ }\PYG{o}{\PYGZgt{}}\PYG{o}{=}\PYG{+w}{ }\PYG{n}{current}\PYG{o}{\PYGZhy{}}\PYG{o}{\PYGZgt{}}\PYG{n}{files}\PYG{o}{\PYGZhy{}}\PYG{o}{\PYGZgt{}}\PYG{n}{fdt}\PYG{o}{\PYGZhy{}}\PYG{o}{\PYGZgt{}}\PYG{n}{max\PYGZus{}fds}\PYG{p}{)}\PYG{+w}{ }\PYG{p}{\PYGZob{}}
\PYG{+w}{        }\PYG{n}{printk}\PYG{p}{(}\PYG{n}{KERN\PYGZus{}INFO}\PYG{+w}{ }\PYG{l+s}{\PYGZdq{}}\PYG{l+s}{file descriptor is not valid!..}\PYG{l+s+se}{\PYGZbs{}n}\PYG{l+s}{\PYGZdq{}}\PYG{p}{)}\PYG{p}{;}
\PYG{+w}{        }\PYG{k}{return}\PYG{+w}{ }\PYG{o}{\PYGZhy{}}\PYG{n}{EBADF}\PYG{p}{;}
\PYG{+w}{    }\PYG{p}{\PYGZcb{}}

\PYG{+w}{    }\PYG{n}{f}\PYG{+w}{ }\PYG{o}{=}\PYG{+w}{ }\PYG{n}{current}\PYG{o}{\PYGZhy{}}\PYG{o}{\PYGZgt{}}\PYG{n}{files}\PYG{o}{\PYGZhy{}}\PYG{o}{\PYGZgt{}}\PYG{n}{fdt}\PYG{o}{\PYGZhy{}}\PYG{o}{\PYGZgt{}}\PYG{n}{fd}\PYG{p}{[}\PYG{n}{fd}\PYG{p}{]}\PYG{p}{;}
\PYG{+w}{    }\PYG{k}{if}\PYG{+w}{ }\PYG{p}{(}\PYG{n}{f}\PYG{+w}{ }\PYG{o}{=}\PYG{o}{=}\PYG{+w}{ }\PYG{n+nb}{NULL}\PYG{p}{)}\PYG{+w}{ }\PYG{p}{\PYGZob{}}
\PYG{+w}{        }\PYG{n}{printk}\PYG{p}{(}\PYG{n}{KERN\PYGZus{}INFO}\PYG{+w}{ }\PYG{l+s}{\PYGZdq{}}\PYG{l+s}{file descriptor is not valid!..}\PYG{l+s+se}{\PYGZbs{}n}\PYG{l+s}{\PYGZdq{}}\PYG{p}{)}\PYG{p}{;}
\PYG{+w}{        }\PYG{k}{return}\PYG{+w}{ }\PYG{o}{\PYGZhy{}}\PYG{n}{EBADF}\PYG{p}{;}
\PYG{+w}{    }\PYG{p}{\PYGZcb{}}

\PYG{+w}{    }\PYG{n}{printk}\PYG{p}{(}\PYG{n}{KERN\PYGZus{}INFO}\PYG{+w}{ }\PYG{l+s}{\PYGZdq{}}\PYG{l+s}{File pointer offset: \PYGZpc{}lld}\PYG{l+s+se}{\PYGZbs{}n}\PYG{l+s}{\PYGZdq{}}\PYG{p}{,}\PYG{+w}{ }\PYG{p}{(}\PYG{k+kt}{long}\PYG{+w}{ }\PYG{k+kt}{long}\PYG{p}{)}\PYG{n}{f}\PYG{o}{\PYGZhy{}}\PYG{o}{\PYGZgt{}}\PYG{n}{f\PYGZus{}pos}\PYG{p}{)}\PYG{p}{;}

\PYG{+w}{    }\PYG{k}{return}\PYG{+w}{ }\PYG{l+m+mi}{0}\PYG{p}{;}
\PYG{p}{\PYGZcb{}}
\end{sphinxVerbatim}

\sphinxAtStartPar
Burada dosya betimleyicisinden hareketle dosya nesnesine \sphinxcode{\sphinxupquote{current\sphinxhyphen{}\textgreater{}files\sphinxhyphen{}\textgreater{}fdt\sphinxhyphen{}\textgreater{}fd{[}fd{]}}} ifadesiyle
erişilmiştir. Ancak bu erişim aslında güvenli değildir. Çünkü tam bu erişim yapılırken dosya
betimleyici tablosu büyütülürse ya da bu betimleyici üzerinde işlem yapılırsa kodumuz kararsız
bir durumla karşı karşıya kalır. Genel olarak çekirdeğin uyguladığı senkronizasyon mekanizmasına
uygun hareket etmek gerekir.

\sphinxAtStartPar
Çekirdek nesneleri üzerinde işlem yapacak kişi nesnenin referans sayacını artırmazsa çekirdek onu
boşaltabilir. Linux çekirdeklerinde sayacı artırarak nesneyi elde eden fonksiyonlar \sphinxstyleemphasis{get} soneki
ile, sayacı azaltarak nesneyi bırakan fonksiyonlar ise \sphinxstyleemphasis{put} soneki ile isimlendirilmiştir.

\sphinxAtStartPar
\sphinxstylestrong{fget/fput kullanımı (ioctl\_test2):}

\sphinxAtStartPar
Yukarıdaki testi daha güvenli bir biçimde \sphinxcode{\sphinxupquote{fget}} ve \sphinxcode{\sphinxupquote{fput}} fonksiyonlarını kullanarak da
yapabiliriz:

\begin{sphinxVerbatim}[commandchars=\\\{\}]
\PYG{k}{static}\PYG{+w}{ }\PYG{k+kt}{long}\PYG{+w}{ }\PYG{n+nf}{ioctl\PYGZus{}test2}\PYG{p}{(}\PYG{k}{struct}\PYG{+w}{ }\PYG{n+nc}{file}\PYG{+w}{ }\PYG{o}{*}\PYG{n}{filp}\PYG{p}{,}\PYG{+w}{ }\PYG{k+kt}{unsigned}\PYG{+w}{ }\PYG{k+kt}{long}\PYG{+w}{ }\PYG{n}{arg}\PYG{p}{)}
\PYG{p}{\PYGZob{}}
\PYG{+w}{    }\PYG{k}{struct}\PYG{+w}{ }\PYG{n+nc}{file}\PYG{+w}{ }\PYG{o}{*}\PYG{n}{f}\PYG{p}{;}
\PYG{+w}{    }\PYG{k+kt}{int}\PYG{+w}{ }\PYG{n}{fd}\PYG{+w}{ }\PYG{o}{=}\PYG{+w}{ }\PYG{p}{(}\PYG{k+kt}{int}\PYG{p}{)}\PYG{+w}{ }\PYG{n}{arg}\PYG{p}{;}

\PYG{+w}{    }\PYG{k}{if}\PYG{+w}{ }\PYG{p}{(}\PYG{p}{(}\PYG{n}{f}\PYG{+w}{ }\PYG{o}{=}\PYG{+w}{ }\PYG{n}{fget}\PYG{p}{(}\PYG{n}{fd}\PYG{p}{)}\PYG{p}{)}\PYG{+w}{ }\PYG{o}{=}\PYG{o}{=}\PYG{+w}{ }\PYG{n+nb}{NULL}\PYG{p}{)}\PYG{+w}{ }\PYG{p}{\PYGZob{}}
\PYG{+w}{        }\PYG{n}{printk}\PYG{p}{(}\PYG{n}{KERN\PYGZus{}INFO}\PYG{+w}{ }\PYG{l+s}{\PYGZdq{}}\PYG{l+s}{file descriptor is not valid!..}\PYG{l+s+se}{\PYGZbs{}n}\PYG{l+s}{\PYGZdq{}}\PYG{p}{)}\PYG{p}{;}
\PYG{+w}{        }\PYG{k}{return}\PYG{+w}{ }\PYG{o}{\PYGZhy{}}\PYG{n}{EBADF}\PYG{p}{;}
\PYG{+w}{    }\PYG{p}{\PYGZcb{}}

\PYG{+w}{    }\PYG{n}{printk}\PYG{p}{(}\PYG{n}{KERN\PYGZus{}INFO}\PYG{+w}{ }\PYG{l+s}{\PYGZdq{}}\PYG{l+s}{File pointer offset: \PYGZpc{}ld}\PYG{l+s+se}{\PYGZbs{}n}\PYG{l+s}{\PYGZdq{}}\PYG{p}{,}\PYG{+w}{ }\PYG{p}{(}\PYG{k+kt}{long}\PYG{p}{)}\PYG{n}{f}\PYG{o}{\PYGZhy{}}\PYG{o}{\PYGZgt{}}\PYG{n}{f\PYGZus{}pos}\PYG{p}{)}\PYG{p}{;}

\PYG{+w}{    }\PYG{n}{fput}\PYG{p}{(}\PYG{n}{f}\PYG{p}{)}\PYG{p}{;}

\PYG{+w}{    }\PYG{k}{return}\PYG{+w}{ }\PYG{l+m+mi}{0}\PYG{p}{;}
\PYG{p}{\PYGZcb{}}
\end{sphinxVerbatim}

\sphinxAtStartPar
\sphinxcode{\sphinxupquote{fget}} isimli yüksek seviyeli çekirdek fonksiyonu dosya betimleyicisinden hareketle bize dosya
nesnesini RCU mekanizmasını kullanarak dosya nesnesinin sayacını da artırarak vermektedir. \sphinxcode{\sphinxupquote{fput}}
fonksiyonu da sayacı azaltarak dosya nesnesini geri bırakmaktadır. Bu fonksiyonların prototipleri:

\begin{sphinxVerbatim}[commandchars=\\\{\}]
\PYG{k}{struct}\PYG{+w}{ }\PYG{n+nc}{file}\PYG{+w}{ }\PYG{o}{*}\PYG{n}{fget}\PYG{p}{(}\PYG{k+kt}{unsigned}\PYG{+w}{ }\PYG{k+kt}{int}\PYG{+w}{ }\PYG{n}{fd}\PYG{p}{)}\PYG{p}{;}
\PYG{k+kt}{void}\PYG{+w}{ }\PYG{n+nf}{fput}\PYG{p}{(}\PYG{k}{struct}\PYG{+w}{ }\PYG{n+nc}{file}\PYG{+w}{ }\PYG{o}{*}\PYG{p}{)}\PYG{p}{;}
\end{sphinxVerbatim}

\sphinxAtStartPar
Bu fonksiyonların prototipleri \sphinxcode{\sphinxupquote{include/linux/file.h}} dosyası içerisindedir. Güncel çekirdeklerde
\sphinxcode{\sphinxupquote{fget}} fonksiyonunun tanımlaması \sphinxcode{\sphinxupquote{fs/file.c}} dosyası içerisinde, \sphinxcode{\sphinxupquote{fput}} fonksiyonunun
tanımlaması ise \sphinxcode{\sphinxupquote{fs/file\_table.c}} içerisinde yapılmıştır. \sphinxcode{\sphinxupquote{fget}} ve \sphinxcode{\sphinxupquote{fput}} fonksiyonları
export edildiği için aygıt sürücüler içerisinde de kullanılabilmektedir.

\sphinxAtStartPar
Güncel çekirdeklere \sphinxcode{\sphinxupquote{fget}} ve \sphinxcode{\sphinxupquote{fput}} işlemlerini daha etkin gerçekleştirmek için \sphinxcode{\sphinxupquote{fdget}} ve
\sphinxcode{\sphinxupquote{fdput}} isimli fonksiyonlar da eklenmiştir. Her ne kadar yeni çekirdeklerde \sphinxcode{\sphinxupquote{fdget}} ve
\sphinxcode{\sphinxupquote{fdput}} daha hızlı çalışıyorsa da \sphinxcode{\sphinxupquote{fget}} ve \sphinxcode{\sphinxupquote{fput}} basit arayüzü ve kullanım kolaylığı
nedeniyle tercih edilebilir.


\subsection{dup Gerçekleştirimi: ioctl\_test3}
\label{\detokenize{filesystem1:dup-gerceklestirimi-ioctl-test3}}
\sphinxAtStartPar
Şimdi de \sphinxcode{\sphinxupquote{dup}} POSIX fonksiyonunun işlevini yerine getiren (yani \sphinxcode{\sphinxupquote{sys\_dup}} sistem fonksiyonu
gibi çalışan) bir fonksiyonu kendimiz yazalım. Bunun için bir ioctl kodu oluşturabiliriz. Bu ioctl
kodu kullanıcı modundan çağrılırken \sphinxcode{\sphinxupquote{ioctl}} fonksiyonunun üçüncü parametresine aşağıdaki gibi
bir yapı nesnesinin adresini geçebiliriz:

\begin{sphinxVerbatim}[commandchars=\\\{\}]
\PYG{k}{struct}\PYG{+w}{ }\PYG{n+nc}{FDDUP}\PYG{+w}{ }\PYG{p}{\PYGZob{}}
\PYG{+w}{    }\PYG{k+kt}{int}\PYG{+w}{ }\PYG{n}{fd}\PYG{p}{;}
\PYG{+w}{    }\PYG{k+kt}{int}\PYG{+w}{ }\PYG{n}{fd\PYGZus{}dup}\PYG{p}{;}
\PYG{p}{\PYGZcb{}}\PYG{p}{;}
\end{sphinxVerbatim}

\sphinxAtStartPar
Burada yapının \sphinxcode{\sphinxupquote{fd}} elemanı çiftlenecek dosya betimleyicisini belirtmektedir. \sphinxcode{\sphinxupquote{fd\_dup}}
elemanı ise çiftleme sonucunda elde edilecek yeni betimleyiciyi belirtmektedir. Bu elemanların
\sphinxcode{\sphinxupquote{dup}} fonksiyonu ile ilişkisini şöyle temsil edebiliriz:

\begin{sphinxVerbatim}[commandchars=\\\{\}]
\PYG{n}{fd\PYGZus{}dup}\PYG{+w}{ }\PYG{o}{=}\PYG{+w}{ }\PYG{n}{dup}\PYG{p}{(}\PYG{n}{fd}\PYG{p}{)}\PYG{p}{;}
\end{sphinxVerbatim}

\sphinxAtStartPar
Dosya betimleyicisini çiftleyen ioctl kodu test amaçlı olarak şöyle yazılabilir:

\begin{sphinxVerbatim}[commandchars=\\\{\}]
\PYG{k}{static}\PYG{+w}{ }\PYG{k+kt}{long}\PYG{+w}{ }\PYG{n+nf}{ioctl\PYGZus{}test3}\PYG{p}{(}\PYG{k}{struct}\PYG{+w}{ }\PYG{n+nc}{file}\PYG{+w}{ }\PYG{o}{*}\PYG{n}{filp}\PYG{p}{,}\PYG{+w}{ }\PYG{k+kt}{unsigned}\PYG{+w}{ }\PYG{k+kt}{long}\PYG{+w}{ }\PYG{n}{arg}\PYG{p}{)}
\PYG{p}{\PYGZob{}}
\PYG{+w}{    }\PYG{k}{struct}\PYG{+w}{ }\PYG{n+nc}{file}\PYG{+w}{ }\PYG{o}{*}\PYG{n}{f}\PYG{p}{;}
\PYG{+w}{    }\PYG{k}{struct}\PYG{+w}{ }\PYG{n+nc}{fdtable}\PYG{+w}{ }\PYG{o}{*}\PYG{n}{fdt}\PYG{p}{;}
\PYG{+w}{    }\PYG{k}{struct}\PYG{+w}{ }\PYG{n+nc}{file}\PYG{+w}{ }\PYG{o}{*}\PYG{o}{*}\PYG{n}{fd\PYGZus{}table}\PYG{p}{;}
\PYG{+w}{    }\PYG{k}{struct}\PYG{+w}{ }\PYG{n+nc}{FDDUP}\PYG{+w}{ }\PYG{n}{fddup}\PYG{p}{;}
\PYG{+w}{    }\PYG{k+kt}{int}\PYG{+w}{ }\PYG{n}{i}\PYG{p}{;}

\PYG{+w}{    }\PYG{n}{f}\PYG{+w}{ }\PYG{o}{=}\PYG{+w}{ }\PYG{n+nb}{NULL}\PYG{p}{;}

\PYG{+w}{    }\PYG{k}{if}\PYG{+w}{ }\PYG{p}{(}\PYG{n}{copy\PYGZus{}from\PYGZus{}user}\PYG{p}{(}\PYG{o}{\PYGZam{}}\PYG{n}{fddup}\PYG{p}{,}\PYG{+w}{ }\PYG{p}{(}\PYG{k}{struct}\PYG{+w}{ }\PYG{n+nc}{FDDUP}\PYG{+w}{ }\PYG{o}{*}\PYG{p}{)}\PYG{n}{arg}\PYG{p}{,}\PYG{+w}{ }\PYG{k}{sizeof}\PYG{p}{(}\PYG{n}{fddup}\PYG{p}{)}\PYG{p}{)}\PYG{+w}{ }\PYG{o}{!}\PYG{o}{=}\PYG{+w}{ }\PYG{l+m+mi}{0}\PYG{p}{)}
\PYG{+w}{        }\PYG{k}{return}\PYG{+w}{ }\PYG{o}{\PYGZhy{}}\PYG{n}{EFAULT}\PYG{p}{;}

\PYG{+w}{    }\PYG{n}{fdt}\PYG{+w}{ }\PYG{o}{=}\PYG{+w}{ }\PYG{n}{current}\PYG{o}{\PYGZhy{}}\PYG{o}{\PYGZgt{}}\PYG{n}{files}\PYG{o}{\PYGZhy{}}\PYG{o}{\PYGZgt{}}\PYG{n}{fdt}\PYG{p}{;}
\PYG{+w}{    }\PYG{n}{fd\PYGZus{}table}\PYG{+w}{ }\PYG{o}{=}\PYG{+w}{ }\PYG{n}{fdt}\PYG{o}{\PYGZhy{}}\PYG{o}{\PYGZgt{}}\PYG{n}{fd}\PYG{p}{;}

\PYG{+w}{    }\PYG{k}{if}\PYG{+w}{ }\PYG{p}{(}\PYG{n}{fddup}\PYG{p}{.}\PYG{n}{fd}\PYG{+w}{ }\PYG{o}{\PYGZlt{}}\PYG{+w}{ }\PYG{l+m+mi}{0}\PYG{+w}{ }\PYG{o}{|}\PYG{o}{|}\PYG{+w}{ }\PYG{n}{fddup}\PYG{p}{.}\PYG{n}{fd}\PYG{+w}{ }\PYG{o}{\PYGZgt{}}\PYG{o}{=}\PYG{+w}{ }\PYG{n}{fdt}\PYG{o}{\PYGZhy{}}\PYG{o}{\PYGZgt{}}\PYG{n}{max\PYGZus{}fds}\PYG{p}{)}
\PYG{+w}{        }\PYG{k}{return}\PYG{+w}{ }\PYG{o}{\PYGZhy{}}\PYG{n}{EBADF}\PYG{p}{;}
\PYG{+w}{    }\PYG{k}{if}\PYG{+w}{ }\PYG{p}{(}\PYG{n}{fd\PYGZus{}table}\PYG{p}{[}\PYG{n}{fddup}\PYG{p}{.}\PYG{n}{fd}\PYG{p}{]}\PYG{+w}{ }\PYG{o}{=}\PYG{o}{=}\PYG{+w}{ }\PYG{n+nb}{NULL}\PYG{p}{)}
\PYG{+w}{        }\PYG{k}{return}\PYG{+w}{ }\PYG{o}{\PYGZhy{}}\PYG{n}{EBADF}\PYG{p}{;}

\PYG{+w}{    }\PYG{k}{for}\PYG{+w}{ }\PYG{p}{(}\PYG{n}{i}\PYG{+w}{ }\PYG{o}{=}\PYG{+w}{ }\PYG{l+m+mi}{0}\PYG{p}{;}\PYG{+w}{ }\PYG{n}{i}\PYG{+w}{ }\PYG{o}{\PYGZlt{}}\PYG{+w}{ }\PYG{n}{fdt}\PYG{o}{\PYGZhy{}}\PYG{o}{\PYGZgt{}}\PYG{n}{max\PYGZus{}fds}\PYG{p}{;}\PYG{+w}{ }\PYG{o}{+}\PYG{o}{+}\PYG{n}{i}\PYG{p}{)}
\PYG{+w}{        }\PYG{k}{if}\PYG{+w}{ }\PYG{p}{(}\PYG{n}{fd\PYGZus{}table}\PYG{p}{[}\PYG{n}{i}\PYG{p}{]}\PYG{+w}{ }\PYG{o}{=}\PYG{o}{=}\PYG{+w}{ }\PYG{n+nb}{NULL}\PYG{p}{)}\PYG{+w}{ }\PYG{p}{\PYGZob{}}
\PYG{+w}{            }\PYG{n}{f}\PYG{+w}{ }\PYG{o}{=}\PYG{+w}{ }\PYG{n}{fd\PYGZus{}table}\PYG{p}{[}\PYG{n}{fddup}\PYG{p}{.}\PYG{n}{fd}\PYG{p}{]}\PYG{p}{;}
\PYG{+w}{            }\PYG{n}{fd\PYGZus{}table}\PYG{p}{[}\PYG{n}{i}\PYG{p}{]}\PYG{+w}{ }\PYG{o}{=}\PYG{+w}{ }\PYG{n}{f}\PYG{p}{;}
\PYG{+w}{            }\PYG{o}{+}\PYG{o}{+}\PYG{n}{f}\PYG{o}{\PYGZhy{}}\PYG{o}{\PYGZgt{}}\PYG{n}{f\PYGZus{}count}\PYG{p}{.}\PYG{n}{counter}\PYG{p}{;}\PYG{+w}{     }\PYG{c+cm}{/* atomic\PYGZus{}long\PYGZus{}inc(\PYGZam{}f\PYGZhy{}\PYGZgt{}f\PYGZus{}count); */}
\PYG{+w}{            }\PYG{n}{fddup}\PYG{p}{.}\PYG{n}{fd\PYGZus{}dup}\PYG{+w}{ }\PYG{o}{=}\PYG{+w}{ }\PYG{n}{i}\PYG{p}{;}
\PYG{+w}{            }\PYG{k}{break}\PYG{p}{;}
\PYG{+w}{        }\PYG{p}{\PYGZcb{}}
\PYG{+w}{    }\PYG{k}{if}\PYG{+w}{ }\PYG{p}{(}\PYG{n}{f}\PYG{+w}{ }\PYG{o}{=}\PYG{o}{=}\PYG{+w}{ }\PYG{n+nb}{NULL}\PYG{p}{)}
\PYG{+w}{        }\PYG{k}{return}\PYG{+w}{ }\PYG{o}{\PYGZhy{}}\PYG{n}{EMFILE}\PYG{p}{;}

\PYG{+w}{    }\PYG{k}{if}\PYG{+w}{ }\PYG{p}{(}\PYG{n}{copy\PYGZus{}to\PYGZus{}user}\PYG{p}{(}\PYG{p}{(}\PYG{k}{struct}\PYG{+w}{ }\PYG{n+nc}{FDDUP}\PYG{+w}{ }\PYG{o}{*}\PYG{p}{)}\PYG{n}{arg}\PYG{p}{,}\PYG{+w}{ }\PYG{o}{\PYGZam{}}\PYG{n}{fddup}\PYG{p}{,}\PYG{+w}{ }\PYG{k}{sizeof}\PYG{p}{(}\PYG{n}{fddup}\PYG{p}{)}\PYG{p}{)}\PYG{+w}{ }\PYG{o}{!}\PYG{o}{=}\PYG{+w}{ }\PYG{l+m+mi}{0}\PYG{p}{)}
\PYG{+w}{        }\PYG{k}{return}\PYG{+w}{ }\PYG{o}{\PYGZhy{}}\PYG{n}{EFAULT}\PYG{p}{;}

\PYG{+w}{    }\PYG{k}{return}\PYG{+w}{ }\PYG{l+m+mi}{0}\PYG{p}{;}
\PYG{p}{\PYGZcb{}}
\end{sphinxVerbatim}

\sphinxAtStartPar
Kursumuzun yapıldığı 6.9.2 çekirdeğinde dosya nesnesinin sayacının \sphinxcode{\sphinxupquote{file}} yapısının içerisinde
şöyle tutulduğunu anımsatalım:

\begin{sphinxVerbatim}[commandchars=\\\{\}]
\PYG{k}{struct}\PYG{+w}{ }\PYG{n+nc}{file}\PYG{+w}{ }\PYG{p}{\PYGZob{}}
\PYG{+w}{    }\PYG{c+cm}{/* ... */}

\PYG{+w}{    }\PYG{n}{atomic\PYGZus{}long\PYGZus{}t}\PYG{+w}{       }\PYG{n}{f\PYGZus{}count}\PYG{p}{;}

\PYG{+w}{    }\PYG{c+cm}{/* ... */}
\PYG{p}{\PYGZcb{}}\PYG{p}{;}
\end{sphinxVerbatim}

\sphinxAtStartPar
Burada \sphinxcode{\sphinxupquote{atomic\_long\_t}} türü bir yapı biçiminde bildirilmiştir:

\begin{sphinxVerbatim}[commandchars=\\\{\}]
\PYG{k}{typedef}\PYG{+w}{ }\PYG{k}{struct}\PYG{+w}{ }\PYG{p}{\PYGZob{}}
\PYG{+w}{    }\PYG{k+kt}{long}\PYG{+w}{ }\PYG{n}{counter}\PYG{p}{;}
\PYG{p}{\PYGZcb{}}\PYG{+w}{ }\PYG{n}{atomic\PYGZus{}long\PYGZus{}t}\PYG{p}{;}
\end{sphinxVerbatim}

\sphinxAtStartPar
Bu \sphinxcode{\sphinxupquote{atomic\_long\_t}} türüyle belirtilen değerler üzerinde işlemler atomik düzeyde özel
fonksiyonlarla yapılmaktadır. Bu konuyla ilgili önemli birkaç fonksiyon:

\begin{sphinxVerbatim}[commandchars=\\\{\}]
\PYG{k+kt}{void}\PYG{+w}{ }\PYG{n+nf}{atomic\PYGZus{}long\PYGZus{}set}\PYG{p}{(}\PYG{n}{atomic\PYGZus{}long\PYGZus{}t}\PYG{+w}{ }\PYG{o}{*}\PYG{n}{v}\PYG{p}{,}\PYG{+w}{ }\PYG{k+kt}{long}\PYG{+w}{ }\PYG{n}{i}\PYG{p}{)}\PYG{p}{;}
\PYG{k+kt}{long}\PYG{+w}{ }\PYG{n+nf}{atomic\PYGZus{}long\PYGZus{}read}\PYG{p}{(}\PYG{k}{const}\PYG{+w}{ }\PYG{n}{atomic\PYGZus{}long\PYGZus{}t}\PYG{+w}{ }\PYG{o}{*}\PYG{n}{v}\PYG{p}{)}\PYG{p}{;}
\PYG{k+kt}{void}\PYG{+w}{ }\PYG{n+nf}{atomic\PYGZus{}long\PYGZus{}inc}\PYG{p}{(}\PYG{n}{atomic\PYGZus{}long\PYGZus{}t}\PYG{+w}{ }\PYG{o}{*}\PYG{n}{v}\PYG{p}{)}\PYG{p}{;}
\PYG{k+kt}{void}\PYG{+w}{ }\PYG{n+nf}{atomic\PYGZus{}long\PYGZus{}dec}\PYG{p}{(}\PYG{n}{atomic\PYGZus{}long\PYGZus{}t}\PYG{+w}{ }\PYG{o}{*}\PYG{n}{v}\PYG{p}{)}\PYG{p}{;}
\end{sphinxVerbatim}

\sphinxAtStartPar
\sphinxcode{\sphinxupquote{atomic\_t}} ve \sphinxcode{\sphinxupquote{atomic\_long\_t}} türleri hakkında ayrıntıları başka bir başlıkta ele alacağız.

\sphinxAtStartPar
Fonksiyonumuzdaki döngüde aslında mantıksal bir kusur da vardır. Biz döngüde yalnızca o andaki
dosya betimleyici tablosunda arama yaptık. Halbuki çekirdek aslında başlangıçta küçük bir dosya
betimleyici tablosunu tutarken daha sonra bunu gerektiğinde proses limitinin izin verdiği kadar
yükseltebilmektedir.


\subsection{Yüksek Seviyeli Fonksiyonlarla dup: ioctl\_test4}
\label{\detokenize{filesystem1:yuksek-seviyeli-fonksiyonlarla-dup-ioctl-test4}}
\sphinxAtStartPar
Testi yaptığımız makinedeki Linux çekirdeğinde anımsayacağınız gibi ilk boş betimleyiciyi hızlı
bir biçimde bulabilmek için iki düzey bitmap kullanılıyordu. Bu versiyonlarda ilk boş
betimleyiciyi bulan \sphinxcode{\sphinxupquote{get\_unused\_fd\_flags}} isimli yüksek seviyeli bir çekirdek fonksiyonu
bulunmaktadır. Bu fonksiyon export edildiği için aygıt sürücülerden de kullanılabilmektedir:

\begin{sphinxVerbatim}[commandchars=\\\{\}]
\PYG{k+kt}{int}\PYG{+w}{ }\PYG{n+nf}{get\PYGZus{}unused\PYGZus{}fd\PYGZus{}flags}\PYG{p}{(}\PYG{k+kt}{unsigned}\PYG{+w}{ }\PYG{n}{flags}\PYG{p}{)}
\PYG{p}{\PYGZob{}}
\PYG{+w}{    }\PYG{k}{return}\PYG{+w}{ }\PYG{n}{\PYGZus{}\PYGZus{}get\PYGZus{}unused\PYGZus{}fd\PYGZus{}flags}\PYG{p}{(}\PYG{n}{flags}\PYG{p}{,}\PYG{+w}{ }\PYG{n}{rlimit}\PYG{p}{(}\PYG{n}{RLIMIT\PYGZus{}NOFILE}\PYG{p}{)}\PYG{p}{)}\PYG{p}{;}
\PYG{p}{\PYGZcb{}}
\PYG{n}{EXPORT\PYGZus{}SYMBOL}\PYG{p}{(}\PYG{n}{get\PYGZus{}unused\PYGZus{}fd\PYGZus{}flags}\PYG{p}{)}\PYG{p}{;}
\end{sphinxVerbatim}

\sphinxAtStartPar
\sphinxcode{\sphinxupquote{get\_unused\_fd\_flags}} fonksiyonu aynı zamanda gerektiğinde dosya betimleyici tablosunu büyütme
işlemini de kendisi yapmaktadır. \sphinxcode{\sphinxupquote{flags}} parametresi ya \sphinxcode{\sphinxupquote{O\_CLOEXEC}} biçiminde ya da \sphinxcode{\sphinxupquote{0}}
biçiminde geçilebilmektedir.

\sphinxAtStartPar
Ayrıca RCU mekanizması eşliğinde dosya nesnesini (\sphinxcode{\sphinxupquote{struct file}} nesnesini) dosya betimleyici
tablosuna yerleştiren export edilmiş \sphinxcode{\sphinxupquote{fd\_install}} isimli yüksek seviyeli bir çekirdek fonksiyonu
da bulunmaktadır:

\begin{sphinxVerbatim}[commandchars=\\\{\}]
\PYG{k+kt}{void}\PYG{+w}{ }\PYG{n+nf}{fd\PYGZus{}install}\PYG{p}{(}\PYG{k+kt}{unsigned}\PYG{+w}{ }\PYG{k+kt}{int}\PYG{+w}{ }\PYG{n}{fd}\PYG{p}{,}\PYG{+w}{ }\PYG{k}{struct}\PYG{+w}{ }\PYG{n+nc}{file}\PYG{+w}{ }\PYG{o}{*}\PYG{n}{file}\PYG{p}{)}\PYG{p}{;}
\end{sphinxVerbatim}

\sphinxAtStartPar
Böylece yukarıdaki fonksiyonu şu hale getirebiliriz:

\begin{sphinxVerbatim}[commandchars=\\\{\}]
\PYG{k}{static}\PYG{+w}{ }\PYG{k+kt}{long}\PYG{+w}{ }\PYG{n+nf}{ioctl\PYGZus{}test4}\PYG{p}{(}\PYG{k}{struct}\PYG{+w}{ }\PYG{n+nc}{file}\PYG{+w}{ }\PYG{o}{*}\PYG{n}{filp}\PYG{p}{,}\PYG{+w}{ }\PYG{k+kt}{unsigned}\PYG{+w}{ }\PYG{k+kt}{long}\PYG{+w}{ }\PYG{n}{arg}\PYG{p}{)}
\PYG{p}{\PYGZob{}}
\PYG{+w}{    }\PYG{k}{struct}\PYG{+w}{ }\PYG{n+nc}{file}\PYG{+w}{ }\PYG{o}{*}\PYG{n}{f}\PYG{p}{;}
\PYG{+w}{    }\PYG{k}{struct}\PYG{+w}{ }\PYG{n+nc}{FDDUP}\PYG{+w}{ }\PYG{n}{fddup}\PYG{p}{;}

\PYG{+w}{    }\PYG{k}{if}\PYG{+w}{ }\PYG{p}{(}\PYG{n}{copy\PYGZus{}from\PYGZus{}user}\PYG{p}{(}\PYG{o}{\PYGZam{}}\PYG{n}{fddup}\PYG{p}{,}\PYG{+w}{ }\PYG{p}{(}\PYG{k}{struct}\PYG{+w}{ }\PYG{n+nc}{FDDUP}\PYG{+w}{ }\PYG{o}{*}\PYG{p}{)}\PYG{n}{arg}\PYG{p}{,}\PYG{+w}{ }\PYG{k}{sizeof}\PYG{p}{(}\PYG{n}{fddup}\PYG{p}{)}\PYG{p}{)}\PYG{+w}{ }\PYG{o}{!}\PYG{o}{=}\PYG{+w}{ }\PYG{l+m+mi}{0}\PYG{p}{)}
\PYG{+w}{        }\PYG{k}{return}\PYG{+w}{ }\PYG{o}{\PYGZhy{}}\PYG{n}{EFAULT}\PYG{p}{;}

\PYG{+w}{    }\PYG{k}{if}\PYG{+w}{ }\PYG{p}{(}\PYG{p}{(}\PYG{n}{f}\PYG{+w}{ }\PYG{o}{=}\PYG{+w}{ }\PYG{n}{fget}\PYG{p}{(}\PYG{n}{fddup}\PYG{p}{.}\PYG{n}{fd}\PYG{p}{)}\PYG{p}{)}\PYG{+w}{ }\PYG{o}{=}\PYG{o}{=}\PYG{+w}{ }\PYG{n+nb}{NULL}\PYG{p}{)}
\PYG{+w}{        }\PYG{k}{return}\PYG{+w}{ }\PYG{o}{\PYGZhy{}}\PYG{n}{EBADF}\PYG{p}{;}

\PYG{+w}{    }\PYG{k}{if}\PYG{+w}{ }\PYG{p}{(}\PYG{p}{(}\PYG{n}{fddup}\PYG{p}{.}\PYG{n}{fd\PYGZus{}dup}\PYG{+w}{ }\PYG{o}{=}\PYG{+w}{ }\PYG{n}{get\PYGZus{}unused\PYGZus{}fd\PYGZus{}flags}\PYG{p}{(}\PYG{l+m+mi}{0}\PYG{p}{)}\PYG{p}{)}\PYG{+w}{ }\PYG{o}{\PYGZlt{}}\PYG{+w}{ }\PYG{l+m+mi}{0}\PYG{p}{)}\PYG{+w}{ }\PYG{p}{\PYGZob{}}
\PYG{+w}{        }\PYG{n}{fput}\PYG{p}{(}\PYG{n}{f}\PYG{p}{)}\PYG{p}{;}
\PYG{+w}{        }\PYG{k}{return}\PYG{+w}{ }\PYG{n}{fddup}\PYG{p}{.}\PYG{n}{fd\PYGZus{}dup}\PYG{p}{;}
\PYG{+w}{    }\PYG{p}{\PYGZcb{}}

\PYG{+w}{    }\PYG{k}{if}\PYG{+w}{ }\PYG{p}{(}\PYG{n}{copy\PYGZus{}to\PYGZus{}user}\PYG{p}{(}\PYG{p}{(}\PYG{k}{struct}\PYG{+w}{ }\PYG{n+nc}{FDDUP}\PYG{+w}{ }\PYG{o}{*}\PYG{p}{)}\PYG{n}{arg}\PYG{p}{,}\PYG{+w}{ }\PYG{o}{\PYGZam{}}\PYG{n}{fddup}\PYG{p}{,}\PYG{+w}{ }\PYG{k}{sizeof}\PYG{p}{(}\PYG{n}{fddup}\PYG{p}{)}\PYG{p}{)}\PYG{+w}{ }\PYG{o}{!}\PYG{o}{=}\PYG{+w}{ }\PYG{l+m+mi}{0}\PYG{p}{)}\PYG{+w}{ }\PYG{p}{\PYGZob{}}
\PYG{+w}{        }\PYG{n}{fput}\PYG{p}{(}\PYG{n}{f}\PYG{p}{)}\PYG{p}{;}
\PYG{+w}{        }\PYG{k}{return}\PYG{+w}{ }\PYG{o}{\PYGZhy{}}\PYG{n}{EFAULT}\PYG{p}{;}
\PYG{+w}{    }\PYG{p}{\PYGZcb{}}

\PYG{+w}{    }\PYG{n}{fd\PYGZus{}install}\PYG{p}{(}\PYG{n}{fddup}\PYG{p}{.}\PYG{n}{fd\PYGZus{}dup}\PYG{p}{,}\PYG{+w}{ }\PYG{n}{f}\PYG{p}{)}\PYG{p}{;}

\PYG{+w}{    }\PYG{k}{return}\PYG{+w}{ }\PYG{l+m+mi}{0}\PYG{p}{;}
\PYG{p}{\PYGZcb{}}
\end{sphinxVerbatim}


\subsection{Tam Sürücü Kodu}
\label{\detokenize{filesystem1:tam-surucu-kodu}}
\sphinxAtStartPar
Yukarıdaki test işlemlerine ilişkin tam aygıt sürücü kodu aşağıda verilmiştir.

\sphinxAtStartPar
\sphinxstylestrong{test\sphinxhyphen{}driver.h}

\begin{sphinxVerbatim}[commandchars=\\\{\}]
\PYG{c+cp}{\PYGZsh{}}\PYG{c+cp}{ifndef TEST\PYGZus{}DRIVER\PYGZus{}H\PYGZus{}}
\PYG{c+cp}{\PYGZsh{}}\PYG{c+cp}{define TEST\PYGZus{}DRIVER\PYGZus{}H\PYGZus{}}

\PYG{c+cp}{\PYGZsh{}}\PYG{c+cp}{include}\PYG{+w}{ }\PYG{c+cpf}{\PYGZlt{}linux/stddef.h\PYGZgt{}}
\PYG{c+cp}{\PYGZsh{}}\PYG{c+cp}{include}\PYG{+w}{ }\PYG{c+cpf}{\PYGZlt{}linux/ioctl.h\PYGZgt{}}

\PYG{k}{struct}\PYG{+w}{ }\PYG{n+nc}{FDDUP}\PYG{+w}{ }\PYG{p}{\PYGZob{}}
\PYG{+w}{    }\PYG{k+kt}{int}\PYG{+w}{ }\PYG{n}{fd}\PYG{p}{;}
\PYG{+w}{    }\PYG{k+kt}{int}\PYG{+w}{ }\PYG{n}{fd\PYGZus{}dup}\PYG{p}{;}
\PYG{p}{\PYGZcb{}}\PYG{p}{;}

\PYG{c+cp}{\PYGZsh{}}\PYG{c+cp}{define TEST\PYGZus{}DRIVER\PYGZus{}MAGIC       \PYGZsq{}t\PYGZsq{}}
\PYG{c+cp}{\PYGZsh{}}\PYG{c+cp}{define IOC\PYGZus{}TEST1               \PYGZus{}IOR(TEST\PYGZus{}DRIVER\PYGZus{}MAGIC, 0, int)}
\PYG{c+cp}{\PYGZsh{}}\PYG{c+cp}{define IOC\PYGZus{}TEST2               \PYGZus{}IOR(TEST\PYGZus{}DRIVER\PYGZus{}MAGIC, 1, int)}
\PYG{c+cp}{\PYGZsh{}}\PYG{c+cp}{define IOC\PYGZus{}TEST3               \PYGZus{}IOWR(TEST\PYGZus{}DRIVER\PYGZus{}MAGIC, struct FDDUP)}
\PYG{c+cp}{\PYGZsh{}}\PYG{c+cp}{define IOC\PYGZus{}TEST4               \PYGZus{}IOWR(TEST\PYGZus{}DRIVER\PYGZus{}MAGIC, 3, struct FDDUP)}

\PYG{c+cp}{\PYGZsh{}}\PYG{c+cp}{endif}
\end{sphinxVerbatim}

\sphinxAtStartPar
\sphinxstylestrong{test\sphinxhyphen{}driver.c}

\begin{sphinxVerbatim}[commandchars=\\\{\}]
\PYG{c+cp}{\PYGZsh{}}\PYG{c+cp}{include}\PYG{+w}{ }\PYG{c+cpf}{\PYGZlt{}linux/module.h\PYGZgt{}}
\PYG{c+cp}{\PYGZsh{}}\PYG{c+cp}{include}\PYG{+w}{ }\PYG{c+cpf}{\PYGZlt{}linux/kernel.h\PYGZgt{}}
\PYG{c+cp}{\PYGZsh{}}\PYG{c+cp}{include}\PYG{+w}{ }\PYG{c+cpf}{\PYGZlt{}linux/fs.h\PYGZgt{}}
\PYG{c+cp}{\PYGZsh{}}\PYG{c+cp}{include}\PYG{+w}{ }\PYG{c+cpf}{\PYGZlt{}linux/cdev.h\PYGZgt{}}
\PYG{c+cp}{\PYGZsh{}}\PYG{c+cp}{include}\PYG{+w}{ }\PYG{c+cpf}{\PYGZlt{}linux/fdtable.h\PYGZgt{}}
\PYG{c+cp}{\PYGZsh{}}\PYG{c+cp}{include}\PYG{+w}{ }\PYG{c+cpf}{\PYGZlt{}linux/file.h\PYGZgt{}}
\PYG{c+cp}{\PYGZsh{}}\PYG{c+cp}{include}\PYG{+w}{ }\PYG{c+cpf}{\PYGZdq{}test\PYGZhy{}driver.h\PYGZdq{}}

\PYG{n}{MODULE\PYGZus{}LICENSE}\PYG{p}{(}\PYG{l+s}{\PYGZdq{}}\PYG{l+s}{GPL}\PYG{l+s}{\PYGZdq{}}\PYG{p}{)}\PYG{p}{;}
\PYG{n}{MODULE\PYGZus{}AUTHOR}\PYG{p}{(}\PYG{l+s}{\PYGZdq{}}\PYG{l+s}{Kaan Aslan}\PYG{l+s}{\PYGZdq{}}\PYG{p}{)}\PYG{p}{;}
\PYG{n}{MODULE\PYGZus{}DESCRIPTION}\PYG{p}{(}\PYG{l+s}{\PYGZdq{}}\PYG{l+s}{test\PYGZhy{}driver}\PYG{l+s}{\PYGZdq{}}\PYG{p}{)}\PYG{p}{;}

\PYG{k}{static}\PYG{+w}{ }\PYG{k+kt}{int}\PYG{+w}{ }\PYG{n+nf}{test\PYGZus{}driver\PYGZus{}open}\PYG{p}{(}\PYG{k}{struct}\PYG{+w}{ }\PYG{n+nc}{inode}\PYG{+w}{ }\PYG{o}{*}\PYG{n}{inodep}\PYG{p}{,}\PYG{+w}{ }\PYG{k}{struct}\PYG{+w}{ }\PYG{n+nc}{file}\PYG{+w}{ }\PYG{o}{*}\PYG{n}{filp}\PYG{p}{)}\PYG{p}{;}
\PYG{k}{static}\PYG{+w}{ }\PYG{k+kt}{int}\PYG{+w}{ }\PYG{n+nf}{test\PYGZus{}driver\PYGZus{}release}\PYG{p}{(}\PYG{k}{struct}\PYG{+w}{ }\PYG{n+nc}{inode}\PYG{+w}{ }\PYG{o}{*}\PYG{n}{inodep}\PYG{p}{,}\PYG{+w}{ }\PYG{k}{struct}\PYG{+w}{ }\PYG{n+nc}{file}\PYG{+w}{ }\PYG{o}{*}\PYG{n}{filp}\PYG{p}{)}\PYG{p}{;}
\PYG{k}{static}\PYG{+w}{ }\PYG{k+kt}{ssize\PYGZus{}t}\PYG{+w}{ }\PYG{n+nf}{test\PYGZus{}driver\PYGZus{}read}\PYG{p}{(}\PYG{k}{struct}\PYG{+w}{ }\PYG{n+nc}{file}\PYG{+w}{ }\PYG{o}{*}\PYG{n}{filp}\PYG{p}{,}\PYG{+w}{ }\PYG{k+kt}{char}\PYG{+w}{ }\PYG{o}{*}\PYG{n}{buf}\PYG{p}{,}\PYG{+w}{ }\PYG{k+kt}{size\PYGZus{}t}\PYG{+w}{ }\PYG{n}{size}\PYG{p}{,}\PYG{+w}{ }\PYG{n}{loff\PYGZus{}t}\PYG{+w}{ }\PYG{o}{*}\PYG{n}{off}\PYG{p}{)}\PYG{p}{;}
\PYG{k}{static}\PYG{+w}{ }\PYG{k+kt}{ssize\PYGZus{}t}\PYG{+w}{ }\PYG{n+nf}{test\PYGZus{}driver\PYGZus{}write}\PYG{p}{(}\PYG{k}{struct}\PYG{+w}{ }\PYG{n+nc}{file}\PYG{+w}{ }\PYG{o}{*}\PYG{n}{filp}\PYG{p}{,}\PYG{+w}{ }\PYG{k}{const}\PYG{+w}{ }\PYG{k+kt}{char}\PYG{+w}{ }\PYG{o}{*}\PYG{n}{buf}\PYG{p}{,}\PYG{+w}{ }\PYG{k+kt}{size\PYGZus{}t}\PYG{+w}{ }\PYG{n}{size}\PYG{p}{,}\PYG{+w}{ }\PYG{n}{loff\PYGZus{}t}\PYG{+w}{ }\PYG{o}{*}\PYG{n}{off}\PYG{p}{)}\PYG{p}{;}
\PYG{k}{static}\PYG{+w}{ }\PYG{k+kt}{long}\PYG{+w}{ }\PYG{n+nf}{test\PYGZus{}driver\PYGZus{}ioctl}\PYG{p}{(}\PYG{k}{struct}\PYG{+w}{ }\PYG{n+nc}{file}\PYG{+w}{ }\PYG{o}{*}\PYG{n}{filp}\PYG{p}{,}\PYG{+w}{ }\PYG{k+kt}{unsigned}\PYG{+w}{ }\PYG{k+kt}{int}\PYG{+w}{ }\PYG{n}{cmd}\PYG{p}{,}\PYG{+w}{ }\PYG{k+kt}{unsigned}\PYG{+w}{ }\PYG{k+kt}{long}\PYG{+w}{ }\PYG{n}{arg}\PYG{p}{)}\PYG{p}{;}

\PYG{k}{static}\PYG{+w}{ }\PYG{k+kt}{long}\PYG{+w}{ }\PYG{n+nf}{ioctl\PYGZus{}test1}\PYG{p}{(}\PYG{k}{struct}\PYG{+w}{ }\PYG{n+nc}{file}\PYG{+w}{ }\PYG{o}{*}\PYG{n}{filp}\PYG{p}{,}\PYG{+w}{ }\PYG{k+kt}{unsigned}\PYG{+w}{ }\PYG{k+kt}{long}\PYG{+w}{ }\PYG{n}{arg}\PYG{p}{)}\PYG{p}{;}
\PYG{k}{static}\PYG{+w}{ }\PYG{k+kt}{long}\PYG{+w}{ }\PYG{n+nf}{ioctl\PYGZus{}test2}\PYG{p}{(}\PYG{k}{struct}\PYG{+w}{ }\PYG{n+nc}{file}\PYG{+w}{ }\PYG{o}{*}\PYG{n}{filp}\PYG{p}{,}\PYG{+w}{ }\PYG{k+kt}{unsigned}\PYG{+w}{ }\PYG{k+kt}{long}\PYG{+w}{ }\PYG{n}{arg}\PYG{p}{)}\PYG{p}{;}
\PYG{k}{static}\PYG{+w}{ }\PYG{k+kt}{long}\PYG{+w}{ }\PYG{n+nf}{ioctl\PYGZus{}test3}\PYG{p}{(}\PYG{k}{struct}\PYG{+w}{ }\PYG{n+nc}{file}\PYG{+w}{ }\PYG{o}{*}\PYG{n}{filp}\PYG{p}{,}\PYG{+w}{ }\PYG{k+kt}{unsigned}\PYG{+w}{ }\PYG{k+kt}{long}\PYG{+w}{ }\PYG{n}{arg}\PYG{p}{)}\PYG{p}{;}
\PYG{k}{static}\PYG{+w}{ }\PYG{k+kt}{long}\PYG{+w}{ }\PYG{n+nf}{ioctl\PYGZus{}test4}\PYG{p}{(}\PYG{k}{struct}\PYG{+w}{ }\PYG{n+nc}{file}\PYG{+w}{ }\PYG{o}{*}\PYG{n}{filp}\PYG{p}{,}\PYG{+w}{ }\PYG{k+kt}{unsigned}\PYG{+w}{ }\PYG{k+kt}{long}\PYG{+w}{ }\PYG{n}{arg}\PYG{p}{)}\PYG{p}{;}

\PYG{k}{static}\PYG{+w}{ }\PYG{k+kt}{dev\PYGZus{}t}\PYG{+w}{ }\PYG{n}{g\PYGZus{}dev}\PYG{p}{;}
\PYG{k}{static}\PYG{+w}{ }\PYG{k}{struct}\PYG{+w}{ }\PYG{n+nc}{cdev}\PYG{+w}{ }\PYG{n}{g\PYGZus{}cdev}\PYG{p}{;}
\PYG{k}{static}\PYG{+w}{ }\PYG{k}{struct}\PYG{+w}{ }\PYG{n+nc}{file\PYGZus{}operations}\PYG{+w}{ }\PYG{n}{g\PYGZus{}fops}\PYG{+w}{ }\PYG{o}{=}\PYG{+w}{ }\PYG{p}{\PYGZob{}}
\PYG{+w}{    }\PYG{p}{.}\PYG{n}{owner}\PYG{+w}{          }\PYG{o}{=}\PYG{+w}{ }\PYG{n}{THIS\PYGZus{}MODULE}\PYG{p}{,}
\PYG{+w}{    }\PYG{p}{.}\PYG{n}{open}\PYG{+w}{           }\PYG{o}{=}\PYG{+w}{ }\PYG{n}{test\PYGZus{}driver\PYGZus{}open}\PYG{p}{,}
\PYG{+w}{    }\PYG{p}{.}\PYG{n}{read}\PYG{+w}{           }\PYG{o}{=}\PYG{+w}{ }\PYG{n}{test\PYGZus{}driver\PYGZus{}read}\PYG{p}{,}
\PYG{+w}{    }\PYG{p}{.}\PYG{n}{write}\PYG{+w}{          }\PYG{o}{=}\PYG{+w}{ }\PYG{n}{test\PYGZus{}driver\PYGZus{}write}\PYG{p}{,}
\PYG{+w}{    }\PYG{p}{.}\PYG{n}{release}\PYG{+w}{        }\PYG{o}{=}\PYG{+w}{ }\PYG{n}{test\PYGZus{}driver\PYGZus{}release}\PYG{p}{,}
\PYG{+w}{    }\PYG{p}{.}\PYG{n}{unlocked\PYGZus{}ioctl}\PYG{+w}{ }\PYG{o}{=}\PYG{+w}{ }\PYG{n}{test\PYGZus{}driver\PYGZus{}ioctl}
\PYG{p}{\PYGZcb{}}\PYG{p}{;}

\PYG{k}{static}\PYG{+w}{ }\PYG{k+kt}{int}\PYG{+w}{ }\PYG{n}{\PYGZus{}\PYGZus{}init}\PYG{+w}{ }\PYG{n+nf}{test\PYGZus{}driver\PYGZus{}init}\PYG{p}{(}\PYG{k+kt}{void}\PYG{p}{)}
\PYG{p}{\PYGZob{}}
\PYG{+w}{    }\PYG{k+kt}{int}\PYG{+w}{ }\PYG{n}{result}\PYG{p}{;}

\PYG{+w}{    }\PYG{n}{printk}\PYG{p}{(}\PYG{n}{KERN\PYGZus{}INFO}\PYG{+w}{ }\PYG{l+s}{\PYGZdq{}}\PYG{l+s}{test\PYGZhy{}driver module initialization...}\PYG{l+s+se}{\PYGZbs{}n}\PYG{l+s}{\PYGZdq{}}\PYG{p}{)}\PYG{p}{;}

\PYG{+w}{    }\PYG{k}{if}\PYG{+w}{ }\PYG{p}{(}\PYG{p}{(}\PYG{n}{result}\PYG{+w}{ }\PYG{o}{=}\PYG{+w}{ }\PYG{n}{alloc\PYGZus{}chrdev\PYGZus{}region}\PYG{p}{(}\PYG{o}{\PYGZam{}}\PYG{n}{g\PYGZus{}dev}\PYG{p}{,}\PYG{+w}{ }\PYG{l+m+mi}{0}\PYG{p}{,}\PYG{+w}{ }\PYG{l+m+mi}{1}\PYG{p}{,}\PYG{+w}{ }\PYG{l+s}{\PYGZdq{}}\PYG{l+s}{test\PYGZhy{}driver}\PYG{l+s}{\PYGZdq{}}\PYG{p}{)}\PYG{p}{)}\PYG{+w}{ }\PYG{o}{\PYGZlt{}}\PYG{+w}{ }\PYG{l+m+mi}{0}\PYG{p}{)}\PYG{+w}{ }\PYG{p}{\PYGZob{}}
\PYG{+w}{        }\PYG{n}{printk}\PYG{p}{(}\PYG{n}{KERN\PYGZus{}INFO}\PYG{+w}{ }\PYG{l+s}{\PYGZdq{}}\PYG{l+s}{cannot alloc char driver!...}\PYG{l+s+se}{\PYGZbs{}n}\PYG{l+s}{\PYGZdq{}}\PYG{p}{)}\PYG{p}{;}
\PYG{+w}{        }\PYG{k}{return}\PYG{+w}{ }\PYG{n}{result}\PYG{p}{;}
\PYG{+w}{    }\PYG{p}{\PYGZcb{}}
\PYG{+w}{    }\PYG{n}{cdev\PYGZus{}init}\PYG{p}{(}\PYG{o}{\PYGZam{}}\PYG{n}{g\PYGZus{}cdev}\PYG{p}{,}\PYG{+w}{ }\PYG{o}{\PYGZam{}}\PYG{n}{g\PYGZus{}fops}\PYG{p}{)}\PYG{p}{;}
\PYG{+w}{    }\PYG{k}{if}\PYG{+w}{ }\PYG{p}{(}\PYG{p}{(}\PYG{n}{result}\PYG{+w}{ }\PYG{o}{=}\PYG{+w}{ }\PYG{n}{cdev\PYGZus{}add}\PYG{p}{(}\PYG{o}{\PYGZam{}}\PYG{n}{g\PYGZus{}cdev}\PYG{p}{,}\PYG{+w}{ }\PYG{n}{g\PYGZus{}dev}\PYG{p}{,}\PYG{+w}{ }\PYG{l+m+mi}{1}\PYG{p}{)}\PYG{p}{)}\PYG{+w}{ }\PYG{o}{\PYGZlt{}}\PYG{+w}{ }\PYG{l+m+mi}{0}\PYG{p}{)}\PYG{+w}{ }\PYG{p}{\PYGZob{}}
\PYG{+w}{        }\PYG{n}{unregister\PYGZus{}chrdev\PYGZus{}region}\PYG{p}{(}\PYG{n}{g\PYGZus{}dev}\PYG{p}{,}\PYG{+w}{ }\PYG{l+m+mi}{1}\PYG{p}{)}\PYG{p}{;}
\PYG{+w}{        }\PYG{n}{printk}\PYG{p}{(}\PYG{n}{KERN\PYGZus{}ERR}\PYG{+w}{ }\PYG{l+s}{\PYGZdq{}}\PYG{l+s}{cannot add device!...}\PYG{l+s+se}{\PYGZbs{}n}\PYG{l+s}{\PYGZdq{}}\PYG{p}{)}\PYG{p}{;}
\PYG{+w}{        }\PYG{k}{return}\PYG{+w}{ }\PYG{n}{result}\PYG{p}{;}
\PYG{+w}{    }\PYG{p}{\PYGZcb{}}

\PYG{+w}{    }\PYG{k}{return}\PYG{+w}{ }\PYG{l+m+mi}{0}\PYG{p}{;}
\PYG{p}{\PYGZcb{}}

\PYG{k}{static}\PYG{+w}{ }\PYG{k+kt}{void}\PYG{+w}{ }\PYG{n}{\PYGZus{}\PYGZus{}exit}\PYG{+w}{ }\PYG{n+nf}{test\PYGZus{}driver\PYGZus{}exit}\PYG{p}{(}\PYG{k+kt}{void}\PYG{p}{)}
\PYG{p}{\PYGZob{}}
\PYG{+w}{    }\PYG{n}{cdev\PYGZus{}del}\PYG{p}{(}\PYG{o}{\PYGZam{}}\PYG{n}{g\PYGZus{}cdev}\PYG{p}{)}\PYG{p}{;}
\PYG{+w}{    }\PYG{n}{unregister\PYGZus{}chrdev\PYGZus{}region}\PYG{p}{(}\PYG{n}{g\PYGZus{}dev}\PYG{p}{,}\PYG{+w}{ }\PYG{l+m+mi}{1}\PYG{p}{)}\PYG{p}{;}

\PYG{+w}{    }\PYG{n}{printk}\PYG{p}{(}\PYG{n}{KERN\PYGZus{}INFO}\PYG{+w}{ }\PYG{l+s}{\PYGZdq{}}\PYG{l+s}{test\PYGZhy{}driver module exit...}\PYG{l+s+se}{\PYGZbs{}n}\PYG{l+s}{\PYGZdq{}}\PYG{p}{)}\PYG{p}{;}
\PYG{p}{\PYGZcb{}}

\PYG{k}{static}\PYG{+w}{ }\PYG{k+kt}{int}\PYG{+w}{ }\PYG{n+nf}{test\PYGZus{}driver\PYGZus{}open}\PYG{p}{(}\PYG{k}{struct}\PYG{+w}{ }\PYG{n+nc}{inode}\PYG{+w}{ }\PYG{o}{*}\PYG{n}{inodep}\PYG{p}{,}\PYG{+w}{ }\PYG{k}{struct}\PYG{+w}{ }\PYG{n+nc}{file}\PYG{+w}{ }\PYG{o}{*}\PYG{n}{filp}\PYG{p}{)}
\PYG{p}{\PYGZob{}}
\PYG{+w}{    }\PYG{n}{printk}\PYG{p}{(}\PYG{n}{KERN\PYGZus{}INFO}\PYG{+w}{ }\PYG{l+s}{\PYGZdq{}}\PYG{l+s}{test\PYGZhy{}driver opened...}\PYG{l+s+se}{\PYGZbs{}n}\PYG{l+s}{\PYGZdq{}}\PYG{p}{)}\PYG{p}{;}
\PYG{+w}{    }\PYG{k}{return}\PYG{+w}{ }\PYG{l+m+mi}{0}\PYG{p}{;}
\PYG{p}{\PYGZcb{}}

\PYG{k}{static}\PYG{+w}{ }\PYG{k+kt}{int}\PYG{+w}{ }\PYG{n+nf}{test\PYGZus{}driver\PYGZus{}release}\PYG{p}{(}\PYG{k}{struct}\PYG{+w}{ }\PYG{n+nc}{inode}\PYG{+w}{ }\PYG{o}{*}\PYG{n}{inodep}\PYG{p}{,}\PYG{+w}{ }\PYG{k}{struct}\PYG{+w}{ }\PYG{n+nc}{file}\PYG{+w}{ }\PYG{o}{*}\PYG{n}{filp}\PYG{p}{)}
\PYG{p}{\PYGZob{}}
\PYG{+w}{    }\PYG{n}{printk}\PYG{p}{(}\PYG{n}{KERN\PYGZus{}INFO}\PYG{+w}{ }\PYG{l+s}{\PYGZdq{}}\PYG{l+s}{test\PYGZhy{}driver closed...}\PYG{l+s+se}{\PYGZbs{}n}\PYG{l+s}{\PYGZdq{}}\PYG{p}{)}\PYG{p}{;}
\PYG{+w}{    }\PYG{k}{return}\PYG{+w}{ }\PYG{l+m+mi}{0}\PYG{p}{;}
\PYG{p}{\PYGZcb{}}

\PYG{k}{static}\PYG{+w}{ }\PYG{k+kt}{ssize\PYGZus{}t}\PYG{+w}{ }\PYG{n+nf}{test\PYGZus{}driver\PYGZus{}read}\PYG{p}{(}\PYG{k}{struct}\PYG{+w}{ }\PYG{n+nc}{file}\PYG{+w}{ }\PYG{o}{*}\PYG{n}{filp}\PYG{p}{,}\PYG{+w}{ }\PYG{k+kt}{char}\PYG{+w}{ }\PYG{o}{*}\PYG{n}{buf}\PYG{p}{,}\PYG{+w}{ }\PYG{k+kt}{size\PYGZus{}t}\PYG{+w}{ }\PYG{n}{size}\PYG{p}{,}\PYG{+w}{ }\PYG{n}{loff\PYGZus{}t}\PYG{+w}{ }\PYG{o}{*}\PYG{n}{off}\PYG{p}{)}
\PYG{p}{\PYGZob{}}
\PYG{+w}{    }\PYG{n}{printk}\PYG{p}{(}\PYG{n}{KERN\PYGZus{}INFO}\PYG{+w}{ }\PYG{l+s}{\PYGZdq{}}\PYG{l+s}{test\PYGZhy{}driver read...}\PYG{l+s+se}{\PYGZbs{}n}\PYG{l+s}{\PYGZdq{}}\PYG{p}{)}\PYG{p}{;}
\PYG{+w}{    }\PYG{k}{return}\PYG{+w}{ }\PYG{l+m+mi}{0}\PYG{p}{;}
\PYG{p}{\PYGZcb{}}

\PYG{k}{static}\PYG{+w}{ }\PYG{k+kt}{ssize\PYGZus{}t}\PYG{+w}{ }\PYG{n+nf}{test\PYGZus{}driver\PYGZus{}write}\PYG{p}{(}\PYG{k}{struct}\PYG{+w}{ }\PYG{n+nc}{file}\PYG{+w}{ }\PYG{o}{*}\PYG{n}{filp}\PYG{p}{,}\PYG{+w}{ }\PYG{k}{const}\PYG{+w}{ }\PYG{k+kt}{char}\PYG{+w}{ }\PYG{o}{*}\PYG{n}{buf}\PYG{p}{,}\PYG{+w}{ }\PYG{k+kt}{size\PYGZus{}t}\PYG{+w}{ }\PYG{n}{size}\PYG{p}{,}\PYG{+w}{ }\PYG{n}{loff\PYGZus{}t}\PYG{+w}{ }\PYG{o}{*}\PYG{n}{off}\PYG{p}{)}
\PYG{p}{\PYGZob{}}
\PYG{+w}{    }\PYG{n}{printk}\PYG{p}{(}\PYG{n}{KERN\PYGZus{}INFO}\PYG{+w}{ }\PYG{l+s}{\PYGZdq{}}\PYG{l+s}{test\PYGZhy{}driver write...}\PYG{l+s+se}{\PYGZbs{}n}\PYG{l+s}{\PYGZdq{}}\PYG{p}{)}\PYG{p}{;}
\PYG{+w}{    }\PYG{k}{return}\PYG{+w}{ }\PYG{l+m+mi}{0}\PYG{p}{;}
\PYG{p}{\PYGZcb{}}

\PYG{k}{static}\PYG{+w}{ }\PYG{k+kt}{long}\PYG{+w}{ }\PYG{n+nf}{test\PYGZus{}driver\PYGZus{}ioctl}\PYG{p}{(}\PYG{k}{struct}\PYG{+w}{ }\PYG{n+nc}{file}\PYG{+w}{ }\PYG{o}{*}\PYG{n}{filp}\PYG{p}{,}\PYG{+w}{ }\PYG{k+kt}{unsigned}\PYG{+w}{ }\PYG{k+kt}{int}\PYG{+w}{ }\PYG{n}{cmd}\PYG{p}{,}\PYG{+w}{ }\PYG{k+kt}{unsigned}\PYG{+w}{ }\PYG{k+kt}{long}\PYG{+w}{ }\PYG{n}{arg}\PYG{p}{)}
\PYG{p}{\PYGZob{}}
\PYG{+w}{    }\PYG{k+kt}{long}\PYG{+w}{ }\PYG{n}{result}\PYG{p}{;}

\PYG{+w}{    }\PYG{n}{printk}\PYG{p}{(}\PYG{n}{KERN\PYGZus{}INFO}\PYG{+w}{ }\PYG{l+s}{\PYGZdq{}}\PYG{l+s}{test\PYGZus{}driver\PYGZus{}ioctl...}\PYG{l+s+se}{\PYGZbs{}n}\PYG{l+s}{\PYGZdq{}}\PYG{p}{)}\PYG{p}{;}

\PYG{+w}{    }\PYG{k}{switch}\PYG{+w}{ }\PYG{p}{(}\PYG{n}{cmd}\PYG{p}{)}\PYG{+w}{ }\PYG{p}{\PYGZob{}}
\PYG{+w}{        }\PYG{k}{case}\PYG{+w}{ }\PYG{n+no}{IOC\PYGZus{}TEST1}\PYG{p}{:}
\PYG{+w}{            }\PYG{n}{result}\PYG{+w}{ }\PYG{o}{=}\PYG{+w}{ }\PYG{n}{ioctl\PYGZus{}test1}\PYG{p}{(}\PYG{n}{filp}\PYG{p}{,}\PYG{+w}{ }\PYG{n}{arg}\PYG{p}{)}\PYG{p}{;}
\PYG{+w}{            }\PYG{k}{break}\PYG{p}{;}
\PYG{+w}{        }\PYG{k}{case}\PYG{+w}{ }\PYG{n+no}{IOC\PYGZus{}TEST2}\PYG{p}{:}
\PYG{+w}{            }\PYG{n}{result}\PYG{+w}{ }\PYG{o}{=}\PYG{+w}{ }\PYG{n}{ioctl\PYGZus{}test2}\PYG{p}{(}\PYG{n}{filp}\PYG{p}{,}\PYG{+w}{ }\PYG{n}{arg}\PYG{p}{)}\PYG{p}{;}
\PYG{+w}{            }\PYG{k}{break}\PYG{p}{;}
\PYG{+w}{        }\PYG{k}{case}\PYG{+w}{ }\PYG{n+no}{IOC\PYGZus{}TEST3}\PYG{p}{:}
\PYG{+w}{            }\PYG{n}{result}\PYG{+w}{ }\PYG{o}{=}\PYG{+w}{ }\PYG{n}{ioctl\PYGZus{}test3}\PYG{p}{(}\PYG{n}{filp}\PYG{p}{,}\PYG{+w}{ }\PYG{n}{arg}\PYG{p}{)}\PYG{p}{;}
\PYG{+w}{            }\PYG{k}{break}\PYG{p}{;}
\PYG{+w}{        }\PYG{k}{case}\PYG{+w}{ }\PYG{n+no}{IOC\PYGZus{}TEST4}\PYG{p}{:}
\PYG{+w}{            }\PYG{n}{result}\PYG{+w}{ }\PYG{o}{=}\PYG{+w}{ }\PYG{n}{ioctl\PYGZus{}test4}\PYG{p}{(}\PYG{n}{filp}\PYG{p}{,}\PYG{+w}{ }\PYG{n}{arg}\PYG{p}{)}\PYG{p}{;}
\PYG{+w}{            }\PYG{k}{break}\PYG{p}{;}
\PYG{+w}{        }\PYG{k}{default}\PYG{o}{:}
\PYG{+w}{            }\PYG{n}{result}\PYG{+w}{ }\PYG{o}{=}\PYG{+w}{ }\PYG{o}{\PYGZhy{}}\PYG{n}{ENOTTY}\PYG{p}{;}
\PYG{+w}{    }\PYG{p}{\PYGZcb{}}

\PYG{+w}{    }\PYG{k}{return}\PYG{+w}{ }\PYG{n}{result}\PYG{p}{;}
\PYG{p}{\PYGZcb{}}

\PYG{k}{static}\PYG{+w}{ }\PYG{k+kt}{long}\PYG{+w}{ }\PYG{n+nf}{ioctl\PYGZus{}test1}\PYG{p}{(}\PYG{k}{struct}\PYG{+w}{ }\PYG{n+nc}{file}\PYG{+w}{ }\PYG{o}{*}\PYG{n}{filp}\PYG{p}{,}\PYG{+w}{ }\PYG{k+kt}{unsigned}\PYG{+w}{ }\PYG{k+kt}{long}\PYG{+w}{ }\PYG{n}{arg}\PYG{p}{)}
\PYG{p}{\PYGZob{}}
\PYG{+w}{    }\PYG{k}{struct}\PYG{+w}{ }\PYG{n+nc}{file}\PYG{+w}{ }\PYG{o}{*}\PYG{n}{f}\PYG{p}{;}
\PYG{+w}{    }\PYG{k+kt}{int}\PYG{+w}{ }\PYG{n}{fd}\PYG{+w}{ }\PYG{o}{=}\PYG{+w}{ }\PYG{p}{(}\PYG{k+kt}{int}\PYG{p}{)}\PYG{n}{arg}\PYG{p}{;}

\PYG{+w}{    }\PYG{k}{if}\PYG{+w}{ }\PYG{p}{(}\PYG{n}{fd}\PYG{+w}{ }\PYG{o}{\PYGZlt{}}\PYG{+w}{ }\PYG{l+m+mi}{0}\PYG{+w}{ }\PYG{o}{|}\PYG{o}{|}\PYG{+w}{ }\PYG{n}{fd}\PYG{+w}{ }\PYG{o}{\PYGZgt{}}\PYG{o}{=}\PYG{+w}{ }\PYG{n}{current}\PYG{o}{\PYGZhy{}}\PYG{o}{\PYGZgt{}}\PYG{n}{files}\PYG{o}{\PYGZhy{}}\PYG{o}{\PYGZgt{}}\PYG{n}{fdt}\PYG{o}{\PYGZhy{}}\PYG{o}{\PYGZgt{}}\PYG{n}{max\PYGZus{}fds}\PYG{p}{)}\PYG{+w}{ }\PYG{p}{\PYGZob{}}
\PYG{+w}{        }\PYG{n}{printk}\PYG{p}{(}\PYG{n}{KERN\PYGZus{}INFO}\PYG{+w}{ }\PYG{l+s}{\PYGZdq{}}\PYG{l+s}{file descriptor is not valid!..}\PYG{l+s+se}{\PYGZbs{}n}\PYG{l+s}{\PYGZdq{}}\PYG{p}{)}\PYG{p}{;}
\PYG{+w}{        }\PYG{k}{return}\PYG{+w}{ }\PYG{o}{\PYGZhy{}}\PYG{n}{EBADF}\PYG{p}{;}
\PYG{+w}{    }\PYG{p}{\PYGZcb{}}

\PYG{+w}{    }\PYG{n}{f}\PYG{+w}{ }\PYG{o}{=}\PYG{+w}{ }\PYG{n}{current}\PYG{o}{\PYGZhy{}}\PYG{o}{\PYGZgt{}}\PYG{n}{files}\PYG{o}{\PYGZhy{}}\PYG{o}{\PYGZgt{}}\PYG{n}{fdt}\PYG{o}{\PYGZhy{}}\PYG{o}{\PYGZgt{}}\PYG{n}{fd}\PYG{p}{[}\PYG{n}{fd}\PYG{p}{]}\PYG{p}{;}
\PYG{+w}{    }\PYG{k}{if}\PYG{+w}{ }\PYG{p}{(}\PYG{n}{f}\PYG{+w}{ }\PYG{o}{=}\PYG{o}{=}\PYG{+w}{ }\PYG{n+nb}{NULL}\PYG{p}{)}\PYG{+w}{ }\PYG{p}{\PYGZob{}}
\PYG{+w}{        }\PYG{n}{printk}\PYG{p}{(}\PYG{n}{KERN\PYGZus{}INFO}\PYG{+w}{ }\PYG{l+s}{\PYGZdq{}}\PYG{l+s}{file descriptor is not valid!..}\PYG{l+s+se}{\PYGZbs{}n}\PYG{l+s}{\PYGZdq{}}\PYG{p}{)}\PYG{p}{;}
\PYG{+w}{        }\PYG{k}{return}\PYG{+w}{ }\PYG{o}{\PYGZhy{}}\PYG{n}{EBADF}\PYG{p}{;}
\PYG{+w}{    }\PYG{p}{\PYGZcb{}}

\PYG{+w}{    }\PYG{n}{printk}\PYG{p}{(}\PYG{n}{KERN\PYGZus{}INFO}\PYG{+w}{ }\PYG{l+s}{\PYGZdq{}}\PYG{l+s}{File pointer offset: \PYGZpc{}ld}\PYG{l+s+se}{\PYGZbs{}n}\PYG{l+s}{\PYGZdq{}}\PYG{p}{,}\PYG{+w}{ }\PYG{p}{(}\PYG{k+kt}{long}\PYG{p}{)}\PYG{n}{f}\PYG{o}{\PYGZhy{}}\PYG{o}{\PYGZgt{}}\PYG{n}{f\PYGZus{}pos}\PYG{p}{)}\PYG{p}{;}
\PYG{+w}{    }\PYG{k}{return}\PYG{+w}{ }\PYG{l+m+mi}{0}\PYG{p}{;}
\PYG{p}{\PYGZcb{}}

\PYG{k}{static}\PYG{+w}{ }\PYG{k+kt}{long}\PYG{+w}{ }\PYG{n+nf}{ioctl\PYGZus{}test2}\PYG{p}{(}\PYG{k}{struct}\PYG{+w}{ }\PYG{n+nc}{file}\PYG{+w}{ }\PYG{o}{*}\PYG{n}{filp}\PYG{p}{,}\PYG{+w}{ }\PYG{k+kt}{unsigned}\PYG{+w}{ }\PYG{k+kt}{long}\PYG{+w}{ }\PYG{n}{arg}\PYG{p}{)}
\PYG{p}{\PYGZob{}}
\PYG{+w}{    }\PYG{k}{struct}\PYG{+w}{ }\PYG{n+nc}{file}\PYG{+w}{ }\PYG{o}{*}\PYG{n}{f}\PYG{p}{;}
\PYG{+w}{    }\PYG{k+kt}{int}\PYG{+w}{ }\PYG{n}{fd}\PYG{+w}{ }\PYG{o}{=}\PYG{+w}{ }\PYG{p}{(}\PYG{k+kt}{int}\PYG{p}{)}\PYG{n}{arg}\PYG{p}{;}

\PYG{+w}{    }\PYG{k}{if}\PYG{+w}{ }\PYG{p}{(}\PYG{p}{(}\PYG{n}{f}\PYG{+w}{ }\PYG{o}{=}\PYG{+w}{ }\PYG{n}{fget}\PYG{p}{(}\PYG{n}{fd}\PYG{p}{)}\PYG{p}{)}\PYG{+w}{ }\PYG{o}{=}\PYG{o}{=}\PYG{+w}{ }\PYG{n+nb}{NULL}\PYG{p}{)}\PYG{+w}{ }\PYG{p}{\PYGZob{}}
\PYG{+w}{        }\PYG{n}{printk}\PYG{p}{(}\PYG{n}{KERN\PYGZus{}INFO}\PYG{+w}{ }\PYG{l+s}{\PYGZdq{}}\PYG{l+s}{file descriptor is not valid!..}\PYG{l+s+se}{\PYGZbs{}n}\PYG{l+s}{\PYGZdq{}}\PYG{p}{)}\PYG{p}{;}
\PYG{+w}{        }\PYG{k}{return}\PYG{+w}{ }\PYG{o}{\PYGZhy{}}\PYG{n}{EBADF}\PYG{p}{;}
\PYG{+w}{    }\PYG{p}{\PYGZcb{}}

\PYG{+w}{    }\PYG{n}{printk}\PYG{p}{(}\PYG{n}{KERN\PYGZus{}INFO}\PYG{+w}{ }\PYG{l+s}{\PYGZdq{}}\PYG{l+s}{File pointer offset: \PYGZpc{}ld}\PYG{l+s+se}{\PYGZbs{}n}\PYG{l+s}{\PYGZdq{}}\PYG{p}{,}\PYG{+w}{ }\PYG{p}{(}\PYG{k+kt}{long}\PYG{p}{)}\PYG{n}{f}\PYG{o}{\PYGZhy{}}\PYG{o}{\PYGZgt{}}\PYG{n}{f\PYGZus{}pos}\PYG{p}{)}\PYG{p}{;}
\PYG{+w}{    }\PYG{n}{fput}\PYG{p}{(}\PYG{n}{f}\PYG{p}{)}\PYG{p}{;}
\PYG{+w}{    }\PYG{k}{return}\PYG{+w}{ }\PYG{l+m+mi}{0}\PYG{p}{;}
\PYG{p}{\PYGZcb{}}

\PYG{k}{static}\PYG{+w}{ }\PYG{k+kt}{long}\PYG{+w}{ }\PYG{n+nf}{ioctl\PYGZus{}test3}\PYG{p}{(}\PYG{k}{struct}\PYG{+w}{ }\PYG{n+nc}{file}\PYG{+w}{ }\PYG{o}{*}\PYG{n}{filp}\PYG{p}{,}\PYG{+w}{ }\PYG{k+kt}{unsigned}\PYG{+w}{ }\PYG{k+kt}{long}\PYG{+w}{ }\PYG{n}{arg}\PYG{p}{)}
\PYG{p}{\PYGZob{}}
\PYG{+w}{    }\PYG{k}{struct}\PYG{+w}{ }\PYG{n+nc}{file}\PYG{+w}{ }\PYG{o}{*}\PYG{n}{f}\PYG{p}{;}
\PYG{+w}{    }\PYG{k}{struct}\PYG{+w}{ }\PYG{n+nc}{fdtable}\PYG{+w}{ }\PYG{o}{*}\PYG{n}{fdt}\PYG{p}{;}
\PYG{+w}{    }\PYG{k}{struct}\PYG{+w}{ }\PYG{n+nc}{file}\PYG{+w}{ }\PYG{o}{*}\PYG{o}{*}\PYG{n}{fd\PYGZus{}table}\PYG{p}{;}
\PYG{+w}{    }\PYG{k}{struct}\PYG{+w}{ }\PYG{n+nc}{FDDUP}\PYG{+w}{ }\PYG{n}{fddup}\PYG{p}{;}
\PYG{+w}{    }\PYG{k+kt}{int}\PYG{+w}{ }\PYG{n}{i}\PYG{p}{;}

\PYG{+w}{    }\PYG{n}{f}\PYG{+w}{ }\PYG{o}{=}\PYG{+w}{ }\PYG{n+nb}{NULL}\PYG{p}{;}

\PYG{+w}{    }\PYG{k}{if}\PYG{+w}{ }\PYG{p}{(}\PYG{n}{copy\PYGZus{}from\PYGZus{}user}\PYG{p}{(}\PYG{o}{\PYGZam{}}\PYG{n}{fddup}\PYG{p}{,}\PYG{+w}{ }\PYG{p}{(}\PYG{k}{struct}\PYG{+w}{ }\PYG{n+nc}{FDDUP}\PYG{+w}{ }\PYG{o}{*}\PYG{p}{)}\PYG{n}{arg}\PYG{p}{,}\PYG{+w}{ }\PYG{k}{sizeof}\PYG{p}{(}\PYG{n}{fddup}\PYG{p}{)}\PYG{p}{)}\PYG{+w}{ }\PYG{o}{!}\PYG{o}{=}\PYG{+w}{ }\PYG{l+m+mi}{0}\PYG{p}{)}
\PYG{+w}{        }\PYG{k}{return}\PYG{+w}{ }\PYG{o}{\PYGZhy{}}\PYG{n}{EFAULT}\PYG{p}{;}

\PYG{+w}{    }\PYG{n}{fdt}\PYG{+w}{ }\PYG{o}{=}\PYG{+w}{ }\PYG{n}{current}\PYG{o}{\PYGZhy{}}\PYG{o}{\PYGZgt{}}\PYG{n}{files}\PYG{o}{\PYGZhy{}}\PYG{o}{\PYGZgt{}}\PYG{n}{fdt}\PYG{p}{;}
\PYG{+w}{    }\PYG{n}{fd\PYGZus{}table}\PYG{+w}{ }\PYG{o}{=}\PYG{+w}{ }\PYG{n}{fdt}\PYG{o}{\PYGZhy{}}\PYG{o}{\PYGZgt{}}\PYG{n}{fd}\PYG{p}{;}

\PYG{+w}{    }\PYG{k}{if}\PYG{+w}{ }\PYG{p}{(}\PYG{n}{fddup}\PYG{p}{.}\PYG{n}{fd}\PYG{+w}{ }\PYG{o}{\PYGZlt{}}\PYG{+w}{ }\PYG{l+m+mi}{0}\PYG{+w}{ }\PYG{o}{|}\PYG{o}{|}\PYG{+w}{ }\PYG{n}{fddup}\PYG{p}{.}\PYG{n}{fd}\PYG{+w}{ }\PYG{o}{\PYGZgt{}}\PYG{o}{=}\PYG{+w}{ }\PYG{n}{fdt}\PYG{o}{\PYGZhy{}}\PYG{o}{\PYGZgt{}}\PYG{n}{max\PYGZus{}fds}\PYG{p}{)}
\PYG{+w}{        }\PYG{k}{return}\PYG{+w}{ }\PYG{o}{\PYGZhy{}}\PYG{n}{EBADF}\PYG{p}{;}
\PYG{+w}{    }\PYG{k}{if}\PYG{+w}{ }\PYG{p}{(}\PYG{n}{fd\PYGZus{}table}\PYG{p}{[}\PYG{n}{fddup}\PYG{p}{.}\PYG{n}{fd}\PYG{p}{]}\PYG{+w}{ }\PYG{o}{=}\PYG{o}{=}\PYG{+w}{ }\PYG{n+nb}{NULL}\PYG{p}{)}
\PYG{+w}{        }\PYG{k}{return}\PYG{+w}{ }\PYG{o}{\PYGZhy{}}\PYG{n}{EBADF}\PYG{p}{;}

\PYG{+w}{    }\PYG{k}{for}\PYG{+w}{ }\PYG{p}{(}\PYG{n}{i}\PYG{+w}{ }\PYG{o}{=}\PYG{+w}{ }\PYG{l+m+mi}{0}\PYG{p}{;}\PYG{+w}{ }\PYG{n}{i}\PYG{+w}{ }\PYG{o}{\PYGZlt{}}\PYG{+w}{ }\PYG{n}{fdt}\PYG{o}{\PYGZhy{}}\PYG{o}{\PYGZgt{}}\PYG{n}{max\PYGZus{}fds}\PYG{p}{;}\PYG{+w}{ }\PYG{o}{+}\PYG{o}{+}\PYG{n}{i}\PYG{p}{)}
\PYG{+w}{        }\PYG{k}{if}\PYG{+w}{ }\PYG{p}{(}\PYG{n}{fd\PYGZus{}table}\PYG{p}{[}\PYG{n}{i}\PYG{p}{]}\PYG{+w}{ }\PYG{o}{=}\PYG{o}{=}\PYG{+w}{ }\PYG{n+nb}{NULL}\PYG{p}{)}\PYG{+w}{ }\PYG{p}{\PYGZob{}}
\PYG{+w}{            }\PYG{n}{f}\PYG{+w}{ }\PYG{o}{=}\PYG{+w}{ }\PYG{n}{fd\PYGZus{}table}\PYG{p}{[}\PYG{n}{fddup}\PYG{p}{.}\PYG{n}{fd}\PYG{p}{]}\PYG{p}{;}
\PYG{+w}{            }\PYG{n}{fd\PYGZus{}table}\PYG{p}{[}\PYG{n}{i}\PYG{p}{]}\PYG{+w}{ }\PYG{o}{=}\PYG{+w}{ }\PYG{n}{f}\PYG{p}{;}
\PYG{+w}{            }\PYG{o}{+}\PYG{o}{+}\PYG{n}{f}\PYG{o}{\PYGZhy{}}\PYG{o}{\PYGZgt{}}\PYG{n}{f\PYGZus{}count}\PYG{p}{.}\PYG{n}{counter}\PYG{p}{;}\PYG{+w}{     }\PYG{c+cm}{/* atomic\PYGZus{}long\PYGZus{}inc(\PYGZam{}f\PYGZhy{}\PYGZgt{}f\PYGZus{}count); */}
\PYG{+w}{            }\PYG{n}{fddup}\PYG{p}{.}\PYG{n}{fd\PYGZus{}dup}\PYG{+w}{ }\PYG{o}{=}\PYG{+w}{ }\PYG{n}{i}\PYG{p}{;}
\PYG{+w}{            }\PYG{k}{break}\PYG{p}{;}
\PYG{+w}{        }\PYG{p}{\PYGZcb{}}
\PYG{+w}{    }\PYG{k}{if}\PYG{+w}{ }\PYG{p}{(}\PYG{n}{f}\PYG{+w}{ }\PYG{o}{=}\PYG{o}{=}\PYG{+w}{ }\PYG{n+nb}{NULL}\PYG{p}{)}
\PYG{+w}{        }\PYG{k}{return}\PYG{+w}{ }\PYG{o}{\PYGZhy{}}\PYG{n}{EMFILE}\PYG{p}{;}

\PYG{+w}{    }\PYG{k}{if}\PYG{+w}{ }\PYG{p}{(}\PYG{n}{copy\PYGZus{}to\PYGZus{}user}\PYG{p}{(}\PYG{p}{(}\PYG{k}{struct}\PYG{+w}{ }\PYG{n+nc}{FDDUP}\PYG{+w}{ }\PYG{o}{*}\PYG{p}{)}\PYG{n}{arg}\PYG{p}{,}\PYG{+w}{ }\PYG{o}{\PYGZam{}}\PYG{n}{fddup}\PYG{p}{,}\PYG{+w}{ }\PYG{k}{sizeof}\PYG{p}{(}\PYG{n}{fddup}\PYG{p}{)}\PYG{p}{)}\PYG{+w}{ }\PYG{o}{!}\PYG{o}{=}\PYG{+w}{ }\PYG{l+m+mi}{0}\PYG{p}{)}
\PYG{+w}{        }\PYG{k}{return}\PYG{+w}{ }\PYG{o}{\PYGZhy{}}\PYG{n}{EFAULT}\PYG{p}{;}

\PYG{+w}{    }\PYG{k}{return}\PYG{+w}{ }\PYG{l+m+mi}{0}\PYG{p}{;}
\PYG{p}{\PYGZcb{}}

\PYG{k}{static}\PYG{+w}{ }\PYG{k+kt}{long}\PYG{+w}{ }\PYG{n+nf}{ioctl\PYGZus{}test4}\PYG{p}{(}\PYG{k}{struct}\PYG{+w}{ }\PYG{n+nc}{file}\PYG{+w}{ }\PYG{o}{*}\PYG{n}{filp}\PYG{p}{,}\PYG{+w}{ }\PYG{k+kt}{unsigned}\PYG{+w}{ }\PYG{k+kt}{long}\PYG{+w}{ }\PYG{n}{arg}\PYG{p}{)}
\PYG{p}{\PYGZob{}}
\PYG{+w}{    }\PYG{k}{struct}\PYG{+w}{ }\PYG{n+nc}{file}\PYG{+w}{ }\PYG{o}{*}\PYG{n}{f}\PYG{p}{;}
\PYG{+w}{    }\PYG{k}{struct}\PYG{+w}{ }\PYG{n+nc}{FDDUP}\PYG{+w}{ }\PYG{n}{fddup}\PYG{p}{;}

\PYG{+w}{    }\PYG{k}{if}\PYG{+w}{ }\PYG{p}{(}\PYG{n}{copy\PYGZus{}from\PYGZus{}user}\PYG{p}{(}\PYG{o}{\PYGZam{}}\PYG{n}{fddup}\PYG{p}{,}\PYG{+w}{ }\PYG{p}{(}\PYG{k}{struct}\PYG{+w}{ }\PYG{n+nc}{FDDUP}\PYG{+w}{ }\PYG{o}{*}\PYG{p}{)}\PYG{n}{arg}\PYG{p}{,}\PYG{+w}{ }\PYG{k}{sizeof}\PYG{p}{(}\PYG{n}{fddup}\PYG{p}{)}\PYG{p}{)}\PYG{+w}{ }\PYG{o}{!}\PYG{o}{=}\PYG{+w}{ }\PYG{l+m+mi}{0}\PYG{p}{)}
\PYG{+w}{        }\PYG{k}{return}\PYG{+w}{ }\PYG{o}{\PYGZhy{}}\PYG{n}{EFAULT}\PYG{p}{;}

\PYG{+w}{    }\PYG{k}{if}\PYG{+w}{ }\PYG{p}{(}\PYG{p}{(}\PYG{n}{f}\PYG{+w}{ }\PYG{o}{=}\PYG{+w}{ }\PYG{n}{fget}\PYG{p}{(}\PYG{n}{fddup}\PYG{p}{.}\PYG{n}{fd}\PYG{p}{)}\PYG{p}{)}\PYG{+w}{ }\PYG{o}{=}\PYG{o}{=}\PYG{+w}{ }\PYG{n+nb}{NULL}\PYG{p}{)}
\PYG{+w}{        }\PYG{k}{return}\PYG{+w}{ }\PYG{o}{\PYGZhy{}}\PYG{n}{EBADF}\PYG{p}{;}

\PYG{+w}{    }\PYG{k}{if}\PYG{+w}{ }\PYG{p}{(}\PYG{p}{(}\PYG{n}{fddup}\PYG{p}{.}\PYG{n}{fd\PYGZus{}dup}\PYG{+w}{ }\PYG{o}{=}\PYG{+w}{ }\PYG{n}{get\PYGZus{}unused\PYGZus{}fd\PYGZus{}flags}\PYG{p}{(}\PYG{l+m+mi}{0}\PYG{p}{)}\PYG{p}{)}\PYG{+w}{ }\PYG{o}{\PYGZlt{}}\PYG{+w}{ }\PYG{l+m+mi}{0}\PYG{p}{)}\PYG{+w}{ }\PYG{p}{\PYGZob{}}
\PYG{+w}{        }\PYG{n}{fput}\PYG{p}{(}\PYG{n}{f}\PYG{p}{)}\PYG{p}{;}
\PYG{+w}{        }\PYG{k}{return}\PYG{+w}{ }\PYG{n}{fddup}\PYG{p}{.}\PYG{n}{fd\PYGZus{}dup}\PYG{p}{;}
\PYG{+w}{    }\PYG{p}{\PYGZcb{}}

\PYG{+w}{    }\PYG{k}{if}\PYG{+w}{ }\PYG{p}{(}\PYG{n}{copy\PYGZus{}to\PYGZus{}user}\PYG{p}{(}\PYG{p}{(}\PYG{k}{struct}\PYG{+w}{ }\PYG{n+nc}{FDDUP}\PYG{+w}{ }\PYG{o}{*}\PYG{p}{)}\PYG{n}{arg}\PYG{p}{,}\PYG{+w}{ }\PYG{o}{\PYGZam{}}\PYG{n}{fddup}\PYG{p}{,}\PYG{+w}{ }\PYG{k}{sizeof}\PYG{p}{(}\PYG{n}{fddup}\PYG{p}{)}\PYG{p}{)}\PYG{+w}{ }\PYG{o}{!}\PYG{o}{=}\PYG{+w}{ }\PYG{l+m+mi}{0}\PYG{p}{)}\PYG{+w}{ }\PYG{p}{\PYGZob{}}
\PYG{+w}{        }\PYG{n}{fput}\PYG{p}{(}\PYG{n}{f}\PYG{p}{)}\PYG{p}{;}
\PYG{+w}{        }\PYG{k}{return}\PYG{+w}{ }\PYG{o}{\PYGZhy{}}\PYG{n}{EFAULT}\PYG{p}{;}
\PYG{+w}{    }\PYG{p}{\PYGZcb{}}

\PYG{+w}{    }\PYG{n}{fd\PYGZus{}install}\PYG{p}{(}\PYG{n}{fddup}\PYG{p}{.}\PYG{n}{fd\PYGZus{}dup}\PYG{p}{,}\PYG{+w}{ }\PYG{n}{f}\PYG{p}{)}\PYG{p}{;}

\PYG{+w}{    }\PYG{k}{return}\PYG{+w}{ }\PYG{l+m+mi}{0}\PYG{p}{;}
\PYG{p}{\PYGZcb{}}

\PYG{n}{module\PYGZus{}init}\PYG{p}{(}\PYG{n}{test\PYGZus{}driver\PYGZus{}init}\PYG{p}{)}\PYG{p}{;}
\PYG{n}{module\PYGZus{}exit}\PYG{p}{(}\PYG{n}{test\PYGZus{}driver\PYGZus{}exit}\PYG{p}{)}\PYG{p}{;}
\end{sphinxVerbatim}

\sphinxAtStartPar
\sphinxstylestrong{Makefile}

\begin{sphinxVerbatim}[commandchars=\\\{\}]
\PYG{n+nv}{obj\PYGZhy{}m}\PYG{+w}{ }\PYG{o}{+=}\PYG{+w}{ }\PYG{l+s+si}{\PYGZdl{}\PYGZob{}}\PYG{n+nv}{file}\PYG{l+s+si}{\PYGZcb{}}.o

\PYG{n+nf}{all}\PYG{o}{:}
\PYG{+w}{     }make\PYG{+w}{ }\PYGZhy{}C\PYG{+w}{ }/lib/modules/\PYG{k}{\PYGZdl{}(}shell\PYG{+w}{ }uname\PYG{+w}{ }\PYGZhy{}r\PYG{k}{)}/build\PYG{+w}{ }\PYG{n+nv}{M}\PYG{o}{=}\PYG{l+s+si}{\PYGZdl{}\PYGZob{}}\PYG{n+nv}{PWD}\PYG{l+s+si}{\PYGZcb{}}\PYG{+w}{ }modules
\PYG{n+nf}{clean}\PYG{o}{:}
\PYG{+w}{     }make\PYG{+w}{ }\PYGZhy{}C\PYG{+w}{ }/lib/modules/\PYG{k}{\PYGZdl{}(}shell\PYG{+w}{ }uname\PYG{+w}{ }\PYGZhy{}r\PYG{k}{)}/build\PYG{+w}{ }\PYG{n+nv}{M}\PYG{o}{=}\PYG{l+s+si}{\PYGZdl{}\PYGZob{}}\PYG{n+nv}{PWD}\PYG{l+s+si}{\PYGZcb{}}\PYG{+w}{ }clean
\end{sphinxVerbatim}

\sphinxAtStartPar
\sphinxstylestrong{load} (yükleyici betik)

\begin{sphinxVerbatim}[commandchars=\\\{\}]
\PYG{c+ch}{\PYGZsh{}!/bin/bash}

\PYG{n+nv}{module}\PYG{o}{=}\PYG{n+nv}{\PYGZdl{}1}
\PYG{n+nv}{mode}\PYG{o}{=}\PYG{l+m}{666}

/sbin/insmod\PYG{+w}{ }./\PYG{l+s+si}{\PYGZdl{}\PYGZob{}}\PYG{n+nv}{module}\PYG{l+s+si}{\PYGZcb{}}.ko\PYG{+w}{ }\PYG{l+s+si}{\PYGZdl{}\PYGZob{}}\PYG{p}{@}\PYG{p}{:}\PYG{n+nv}{2}\PYG{l+s+si}{\PYGZcb{}}\PYG{+w}{ }\PYG{o}{||}\PYG{+w}{ }\PYG{n+nb}{exit}\PYG{+w}{ }\PYG{l+m}{1}
\PYG{n+nv}{major}\PYG{o}{=}\PYG{k}{\PYGZdl{}(}awk\PYG{+w}{ }\PYG{l+s+s2}{\PYGZdq{}}\PYG{l+s+s2}{\PYGZbs{}\PYGZdl{}2 == \PYGZbs{}\PYGZdq{}}\PYG{n+nv}{\PYGZdl{}module}\PYG{l+s+s2}{\PYGZbs{}\PYGZdq{} \PYGZob{}print \PYGZbs{}\PYGZdl{}1\PYGZcb{}}\PYG{l+s+s2}{\PYGZdq{}}\PYG{+w}{ }/proc/devices\PYG{k}{)}
rm\PYG{+w}{ }\PYGZhy{}f\PYG{+w}{ }\PYG{n+nv}{\PYGZdl{}module}
mknod\PYG{+w}{ }\PYGZhy{}m\PYG{+w}{ }\PYG{n+nv}{\PYGZdl{}mode}\PYG{+w}{ }\PYG{n+nv}{\PYGZdl{}module}\PYG{+w}{ }c\PYG{+w}{ }\PYG{n+nv}{\PYGZdl{}major}\PYG{+w}{ }\PYG{l+m}{0}
\end{sphinxVerbatim}

\sphinxAtStartPar
\sphinxstylestrong{unload} (kaldırıcı betik)

\begin{sphinxVerbatim}[commandchars=\\\{\}]
\PYG{c+ch}{\PYGZsh{}!/bin/bash}

\PYG{n+nv}{module}\PYG{o}{=}\PYG{n+nv}{\PYGZdl{}1}

/sbin/rmmod\PYG{+w}{ }./\PYG{l+s+si}{\PYGZdl{}\PYGZob{}}\PYG{n+nv}{module}\PYG{l+s+si}{\PYGZcb{}}.ko\PYG{+w}{ }\PYG{o}{||}\PYG{+w}{ }\PYG{n+nb}{exit}\PYG{+w}{ }\PYG{l+m}{1}
rm\PYG{+w}{ }\PYGZhy{}f\PYG{+w}{ }\PYG{n+nv}{\PYGZdl{}module}
\end{sphinxVerbatim}

\sphinxAtStartPar
\sphinxstylestrong{app.c} (test uygulaması)

\begin{sphinxVerbatim}[commandchars=\\\{\}]
\PYG{c+cp}{\PYGZsh{}}\PYG{c+cp}{include}\PYG{+w}{ }\PYG{c+cpf}{\PYGZlt{}stdio.h\PYGZgt{}}
\PYG{c+cp}{\PYGZsh{}}\PYG{c+cp}{include}\PYG{+w}{ }\PYG{c+cpf}{\PYGZlt{}stdlib.h\PYGZgt{}}
\PYG{c+cp}{\PYGZsh{}}\PYG{c+cp}{include}\PYG{+w}{ }\PYG{c+cpf}{\PYGZlt{}stdint.h\PYGZgt{}}
\PYG{c+cp}{\PYGZsh{}}\PYG{c+cp}{include}\PYG{+w}{ }\PYG{c+cpf}{\PYGZlt{}fcntl.h\PYGZgt{}}
\PYG{c+cp}{\PYGZsh{}}\PYG{c+cp}{include}\PYG{+w}{ }\PYG{c+cpf}{\PYGZlt{}unistd.h\PYGZgt{}}
\PYG{c+cp}{\PYGZsh{}}\PYG{c+cp}{include}\PYG{+w}{ }\PYG{c+cpf}{\PYGZlt{}sys/ioctl.h\PYGZgt{}}
\PYG{c+cp}{\PYGZsh{}}\PYG{c+cp}{include}\PYG{+w}{ }\PYG{c+cpf}{\PYGZdq{}test\PYGZhy{}driver.h\PYGZdq{}}

\PYG{k+kt}{void}\PYG{+w}{ }\PYG{n+nf}{exit\PYGZus{}sys}\PYG{p}{(}\PYG{k}{const}\PYG{+w}{ }\PYG{k+kt}{char}\PYG{+w}{ }\PYG{o}{*}\PYG{n}{msg}\PYG{p}{)}\PYG{p}{;}

\PYG{k+kt}{int}\PYG{+w}{ }\PYG{n+nf}{main}\PYG{p}{(}\PYG{k+kt}{int}\PYG{+w}{ }\PYG{n}{argc}\PYG{p}{,}\PYG{+w}{ }\PYG{k+kt}{char}\PYG{+w}{ }\PYG{o}{*}\PYG{n}{argv}\PYG{p}{[}\PYG{p}{]}\PYG{p}{)}
\PYG{p}{\PYGZob{}}
\PYG{+w}{    }\PYG{k+kt}{int}\PYG{+w}{ }\PYG{n}{fd\PYGZus{}dev}\PYG{p}{;}
\PYG{+w}{    }\PYG{k+kt}{int}\PYG{+w}{ }\PYG{n}{fd}\PYG{p}{;}
\PYG{+w}{    }\PYG{k}{struct}\PYG{+w}{ }\PYG{n+nc}{FDDUP}\PYG{+w}{ }\PYG{n}{fddup}\PYG{p}{;}
\PYG{+w}{    }\PYG{k+kt}{off\PYGZus{}t}\PYG{+w}{ }\PYG{n}{pos}\PYG{p}{;}

\PYG{+w}{    }\PYG{k}{if}\PYG{+w}{ }\PYG{p}{(}\PYG{n}{argc}\PYG{+w}{ }\PYG{o}{!}\PYG{o}{=}\PYG{+w}{ }\PYG{l+m+mi}{2}\PYG{p}{)}\PYG{+w}{ }\PYG{p}{\PYGZob{}}
\PYG{+w}{        }\PYG{n}{fprintf}\PYG{p}{(}\PYG{n}{stderr}\PYG{p}{,}\PYG{+w}{ }\PYG{l+s}{\PYGZdq{}}\PYG{l+s}{wrong number of arguments!..}\PYG{l+s+se}{\PYGZbs{}n}\PYG{l+s}{\PYGZdq{}}\PYG{p}{)}\PYG{p}{;}
\PYG{+w}{        }\PYG{n}{exit}\PYG{p}{(}\PYG{n}{EXIT\PYGZus{}FAILURE}\PYG{p}{)}\PYG{p}{;}
\PYG{+w}{    }\PYG{p}{\PYGZcb{}}

\PYG{+w}{    }\PYG{k}{if}\PYG{+w}{ }\PYG{p}{(}\PYG{p}{(}\PYG{n}{fd\PYGZus{}dev}\PYG{+w}{ }\PYG{o}{=}\PYG{+w}{ }\PYG{n}{open}\PYG{p}{(}\PYG{l+s}{\PYGZdq{}}\PYG{l+s}{test\PYGZhy{}driver}\PYG{l+s}{\PYGZdq{}}\PYG{p}{,}\PYG{+w}{ }\PYG{n}{O\PYGZus{}RDONLY}\PYG{p}{)}\PYG{p}{)}\PYG{+w}{ }\PYG{o}{=}\PYG{o}{=}\PYG{+w}{ }\PYG{l+m+mi}{\PYGZhy{}1}\PYG{p}{)}
\PYG{+w}{        }\PYG{n}{exit\PYGZus{}sys}\PYG{p}{(}\PYG{l+s}{\PYGZdq{}}\PYG{l+s}{open}\PYG{l+s}{\PYGZdq{}}\PYG{p}{)}\PYG{p}{;}

\PYG{+w}{    }\PYG{k}{if}\PYG{+w}{ }\PYG{p}{(}\PYG{p}{(}\PYG{n}{fd}\PYG{+w}{ }\PYG{o}{=}\PYG{+w}{ }\PYG{n}{open}\PYG{p}{(}\PYG{n}{argv}\PYG{p}{[}\PYG{l+m+mi}{1}\PYG{p}{]}\PYG{p}{,}\PYG{+w}{ }\PYG{n}{O\PYGZus{}RDONLY}\PYG{p}{)}\PYG{p}{)}\PYG{+w}{ }\PYG{o}{=}\PYG{o}{=}\PYG{+w}{ }\PYG{l+m+mi}{\PYGZhy{}1}\PYG{p}{)}
\PYG{+w}{        }\PYG{n}{exit\PYGZus{}sys}\PYG{p}{(}\PYG{l+s}{\PYGZdq{}}\PYG{l+s}{open}\PYG{l+s}{\PYGZdq{}}\PYG{p}{)}\PYG{p}{;}
\PYG{+w}{    }\PYG{n}{lseek}\PYG{p}{(}\PYG{n}{fd}\PYG{p}{,}\PYG{+w}{ }\PYG{l+m+mi}{100}\PYG{p}{,}\PYG{+w}{ }\PYG{l+m+mi}{0}\PYG{p}{)}\PYG{p}{;}

\PYG{+w}{    }\PYG{n}{fddup}\PYG{p}{.}\PYG{n}{fd}\PYG{+w}{ }\PYG{o}{=}\PYG{+w}{ }\PYG{n}{fd}\PYG{p}{;}
\PYG{+w}{    }\PYG{k}{if}\PYG{+w}{ }\PYG{p}{(}\PYG{n}{ioctl}\PYG{p}{(}\PYG{n}{fd\PYGZus{}dev}\PYG{p}{,}\PYG{+w}{ }\PYG{n}{IOC\PYGZus{}TEST4}\PYG{p}{,}\PYG{+w}{ }\PYG{o}{\PYGZam{}}\PYG{n}{fddup}\PYG{p}{)}\PYG{+w}{ }\PYG{o}{=}\PYG{o}{=}\PYG{+w}{ }\PYG{l+m+mi}{\PYGZhy{}1}\PYG{p}{)}
\PYG{+w}{        }\PYG{n}{exit\PYGZus{}sys}\PYG{p}{(}\PYG{l+s}{\PYGZdq{}}\PYG{l+s}{ioctl}\PYG{l+s}{\PYGZdq{}}\PYG{p}{)}\PYG{p}{;}

\PYG{+w}{    }\PYG{n}{pos}\PYG{+w}{ }\PYG{o}{=}\PYG{+w}{ }\PYG{n}{lseek}\PYG{p}{(}\PYG{n}{fddup}\PYG{p}{.}\PYG{n}{fd\PYGZus{}dup}\PYG{p}{,}\PYG{+w}{ }\PYG{l+m+mi}{0}\PYG{p}{,}\PYG{+w}{ }\PYG{l+m+mi}{1}\PYG{p}{)}\PYG{p}{;}
\PYG{+w}{    }\PYG{n}{printf}\PYG{p}{(}\PYG{l+s}{\PYGZdq{}}\PYG{l+s}{\PYGZpc{}jd}\PYG{l+s+se}{\PYGZbs{}n}\PYG{l+s}{\PYGZdq{}}\PYG{p}{,}\PYG{+w}{ }\PYG{p}{(}\PYG{k+kt}{intmax\PYGZus{}t}\PYG{p}{)}\PYG{n}{pos}\PYG{p}{)}\PYG{p}{;}

\PYG{+w}{    }\PYG{n}{close}\PYG{p}{(}\PYG{n}{fd}\PYG{p}{)}\PYG{p}{;}
\PYG{+w}{    }\PYG{n}{close}\PYG{p}{(}\PYG{n}{fddup}\PYG{p}{.}\PYG{n}{fd}\PYG{p}{)}\PYG{p}{;}
\PYG{+w}{    }\PYG{n}{close}\PYG{p}{(}\PYG{n}{fd\PYGZus{}dev}\PYG{p}{)}\PYG{p}{;}

\PYG{+w}{    }\PYG{k}{return}\PYG{+w}{ }\PYG{l+m+mi}{0}\PYG{p}{;}
\PYG{p}{\PYGZcb{}}

\PYG{k+kt}{void}\PYG{+w}{ }\PYG{n+nf}{exit\PYGZus{}sys}\PYG{p}{(}\PYG{k}{const}\PYG{+w}{ }\PYG{k+kt}{char}\PYG{+w}{ }\PYG{o}{*}\PYG{n}{msg}\PYG{p}{)}
\PYG{p}{\PYGZob{}}
\PYG{+w}{    }\PYG{n}{perror}\PYG{p}{(}\PYG{n}{msg}\PYG{p}{)}\PYG{p}{;}
\PYG{+w}{    }\PYG{n}{exit}\PYG{p}{(}\PYG{n}{EXIT\PYGZus{}FAILURE}\PYG{p}{)}\PYG{p}{;}
\PYG{p}{\PYGZcb{}}
\end{sphinxVerbatim}

\begin{sphinxadmonition}{note}{Not:}
\sphinxAtStartPar
Bu bölümde kullanılan grafiklerin oluşturulabilmesi için Sphinx yapılandırma dosyasında
(\sphinxcode{\sphinxupquote{conf.py}}) \sphinxcode{\sphinxupquote{sphinx.ext.graphviz}} uzantısının etkin olması gerekmektedir:

\begin{sphinxVerbatim}[commandchars=\\\{\}]
\PYG{n}{extensions} \PYG{o}{=} \PYG{p}{[}
    \PYG{o}{.}\PYG{o}{.}\PYG{o}{.}
    \PYG{l+s+s1}{\PYGZsq{}}\PYG{l+s+s1}{sphinx.ext.graphviz}\PYG{l+s+s1}{\PYGZsq{}}\PYG{p}{,}
\PYG{p}{]}
\end{sphinxVerbatim}
\end{sphinxadmonition}



\renewcommand{\indexname}{Dizin}
\printindex
\end{document}